\documentclass[11pt,a4paper]{article}
\usepackage[T1]{fontenc}
\usepackage{lmodern}
\usepackage{amsmath,amssymb,amsthm,mathtools,mathrsfs}
\usepackage[margin=25mm]{geometry}
\usepackage{microtype,booktabs,longtable,array,enumitem,hyperref,xurl,graphicx}
\hypersetup{colorlinks=true,linkcolor=blue,citecolor=blue,urlcolor=blue,
 pdftitle={LM4: Finite group actions, loop structures, and profinite extensions},
 pdfauthor={AI-produced by GPT-6 Astra (OpenAI, ChatGPT Pro mode)}}
\setlength{\emergencystretch}{1.5em}
\setlength{\parskip}{3pt}
\newcommand{\Z}{\mathbb Z}\newcommand{\Q}{\mathbb Q}
\newcommand{\R}{\mathbb R}\newcommand{\C}{\mathbb C}\newcommand{\F}{\mathbb F}
\DeclareMathOperator{\Aut}{Aut}\DeclareMathOperator{\Out}{Out}
\DeclareMathOperator{\Inn}{Inn}\DeclareMathOperator{\Hom}{Hom}
\DeclareMathOperator{\id}{id}\DeclareMathOperator{\im}{im}
\DeclareMathOperator{\coker}{coker}\DeclareMathOperator{\Spec}{Spec}
\DeclareMathOperator{\Gal}{Gal}\DeclareMathOperator{\Map}{Map}
\DeclareMathOperator{\Homeo}{Homeo}\DeclareMathOperator{\Diff}{Diff}
\DeclareMathOperator{\rank}{rank}\DeclareMathOperator{\Sec}{Sec}
\DeclareMathOperator{\GL}{GL}\DeclareMathOperator{\PGL}{PGL}
\newenvironment{claim}[3]{\par\medskip\phantomsection\label{cl:#1}%
 \noindent\textbf{#1. [#2] #3.}\ \ignorespaces}{\par\medskip}
\newcommand{\passage}[1]{\par\noindent{\footnotesize\emph{Manuscript:} #1.}\par}
\begin{document}
\begin{center}
{\Large\bfseries Finite group actions, loop structures,\par
and profinite extensions\par}\smallskip
{\large A mathematical development of \emph{La Longue Marche}: LM4\par}\medskip
\end{center}
\noindent AI-produced by GPT-6 Astra (OpenAI, ChatGPT Pro mode). No priority claimed.
\par\noindent 13 November 2026

\noindent\textbf{Working manuscript: INCOMPLETE.}
The numbered statements distinguish proved assertions, complete counterexamples,
source theorems with their reductions, and the exact obligations still under
investigation. This working record makes no assertion of completion of the task.

\smallskip
\noindent\textbf{Notation and its manuscript concordance.}
Local generators, base points, coefficients and comparison maps are defined at
their use. The table fixes the global symbols and the meaning of each variant.

{\small\renewcommand{\arraystretch}{1.13}\noindent
\begin{tabular}{@{}p{.16\textwidth}p{.49\textwidth}p{.29\textwidth}@{}}
\toprule
Our notation & Meaning and Grothendieck's notation & Introducing passage / scans\\
\midrule
$X,S,U=X-S$ & Compact surface, marked set, punctured surface; $X,S,U$ & \S12: 91--92; \S19: 204\\
$(g,\nu),I$ & Genus, number and index set of punctures; $(g,\nu),I$ & \S13 bis: 115--116\\
$\pi,\pi_0,P$ & Discrete surface group, chosen discretization, its completion; $\pi,\pi_0,\pi$ & \S12: 91; \S28: 352\\
$L_i,\kappa_i$ & Indexed peripheral subgroup and exterior source map; $L_i,\kappa_i$ & \S12: 91--99\\
$T=T_\pi$ & Covariant orientation line; dual of the earlier cohomological line & 91; 116; 130\\
$A,A^0$ & Surface homeomorphism group and identity component; $A_{g,\nu},A^0$ & \S19: 204--205\\
$M,\Gamma_{g,\nu}$ & Geometric mapping class group; $\mathcal T_{g,\nu},\Gamma_{g,\nu}$ & 193--203; 204--205\\
$!,+$ & Pointwise marked-point preservation; orientation preservation & 204--205\\
$G,\Gamma,E$ & Finite action, operator group, lift extension; his $\Gamma,G,E$ as locally specified & \S15: 137--139; 307--312\\
$\Phi,\mu$ & Maximal partial-section data and torsion multiplier; $\Phi,\mu$ & 150--160; 197--203\\
$SE,S\Gamma$ & Special actual and exterior automorphism groups; $SE,S\Gamma$ & \S25: 296--300\\
$B,\overline M$ & $\Out_{\rm lac}(P)$ and closure of a specified geometric group; $\widehat{\widehat{\mathcal T}},\widehat{\mathcal T}$ where identified & 318--320; 352--358\\
$\operatorname{Discr},\operatorname{Prediscr}$ & Actual, outer, or closed-stabilizer cosets as explicitly specified & \S28: 352--358\\
$GT_0,\lambda,f$ & The normalized centralizer of $S_3$, multiplier and commutator parameter & Two-group problem: 312, 316--317\\
\bottomrule
\end{tabular}}

\clearpage
\section*{Reading the proof record}
\addcontentsline{toc}{section}{Reading the proof record}
We work in the order A--F of the supplied task. A group with loops retains its
indexed exterior peripheral maps and its orientation line; this includes the
exceptional cases where the image subgroups alone lose that data. All
homomorphisms between profinite or pro-$p$ groups are continuous. A section
retains its prescribed quotient labels, and its equivalence is conjugacy by
the named kernel. Every passage citation gives the manuscript section, folio,
and scan; a folio followed by ``verso'' names its reverse side. A transcription
alternative is retained when mathematical evidence bears on it. Our corrections
are assertions about the displayed mathematics, not declarations that a
doubtful reading has been settled.

The results register contains the full numbered assertions and dependencies.
The readings register records the editorial consequences. The report gives
the obligation-by-obligation state, source checks and computation provenance.
An ``Open'' label records an unfinished obligation in this working manuscript;
it does not replace a requested proof. A ``Known'' label includes the
derivation from the precise source theorem. The code uses exact finite
arithmetic, and the proofs identify what the computations certify and what
requires a separate infinite-group argument.

\tableofcontents
\clearpage

\section{Groups with loops and their rigidity}
\label{sec:part-a}
We use $[x,y]=xyx^{-1}y^{-1}$. A peripheral subgroup is indexed even when two indices give the same subgroup. The notation $\Out_{\mathrm{lac}}$ means preservation of the indexed peripheral conjugacy classes, with their permitted permutations; in hyperbolic type these permutations are determined by the automorphism. We specify the additional data needed in the exceptional types below. All profinite homomorphisms and automorphisms in this section are continuous.

\begin{claim}{A.1}{Proved}{The presentation, ranks, and exceptional types}
Let $X$ be a connected closed oriented topological surface of genus $g$, let $S$ consist of $\nu$ distinct points, and put $U=X\setminus S$. A discrete group with loops of type $(g,\nu)$ consists of a group $\pi$, a finite index set $I$ of cardinality $\nu$, a free rank-one integral orientation module $T$, and exterior homomorphisms $\kappa_i:T\to\pi$. Here an exterior homomorphism is a conjugacy class, under conjugation in the target, of actual homomorphisms. After choosing a generator of $T$, representatives $c_i=\kappa_i(1)$, an ordering of $I$, and the handle generators, these data admit the presentation
\[
 \pi_{g,\nu}=\left\langle a_1,b_1,\ldots,a_g,b_g,c_1,\ldots,c_\nu\ \middle|\
 \prod_{j=1}^{g}[a_j,b_j]c_1\cdots c_\nu=1\right\rangle,
 \qquad L_i=\langle c_i\rangle.
\]
The profinite version uses a free rank-one $\widehat\Z$-module $T$, continuous exterior homomorphisms, and the corresponding presentation in profinite groups, namely the quotient of the free profinite group by the closed normal subgroup generated by this relator. When $I$ is empty one additionally identifies $T$ with the dual of $H^2(\pi;\Z)$, or its continuous profinite counterpart. This is the covariant orientation convention of scan 116. More precisely, an earlier cohomological marking $\kappa:H^2(\pi;\Z)\xrightarrow{\sim}T_{\mathrm{coh}}$ on scan 91 gives, by dualization, $\kappa^\vee:T_{\mathrm{coh}}^\vee\xrightarrow{\sim}H_2(\pi;\Z)$; the covariant module used here is $T=T_{\mathrm{coh}}^\vee$. The same typed dual map specifies the profinite convention. For an isomorphism of groups with loops the data are an isomorphism $f:\pi\to\pi_1$, a bijection $s:I\to I_1$, and an isomorphism $t:T\to T_1$ such that $f\kappa_i$ is conjugate to $\kappa_{s(i)}t$ for every $i$, with the corresponding dual-cohomology compatibility when $I$ is empty. An exterior group with loops regards two such $f$ as equivalent when they differ by target inner conjugation, with the same $s,t$. Thus its isomorphism sets are explicitly the quotient of the preceding isomorphism sets by $\Inn(\pi_1)$. In all nontrivial admitted types $L_i=\kappa_i(T)$ recovers its source by A.2; the one-puncture sphere requires the original source map.
For $\nu>0$ the group is free of rank $2g+\nu-1$, and its abelianization has that rank. For $\nu=0$ the abelianization has rank $2g$. The abelian types are exactly
\[
 (0,0),\quad(0,1),\quad(0,2),\quad(1,0),
\]
with groups $1,1,\Z,\Z^2$, respectively. The same list holds for their profinite completions. The peripheral maps from the source orientation group $\Z$ to $\pi$ are injective except at $(0,1)$, among the types admitted in scans 91--92. The closed sphere is explicitly excluded there.
\end{claim}
\passage{12; 44--48; 91--100; 13 bis; 56--57; 115--118}

\begin{proof}
Eliminate $c_\nu$ when $\nu>0$; the remaining generators have no relation. Abelianizing the displayed presentation leaves only $c_1+\cdots+c_\nu=0$, giving both ranks. A free group of rank at least two is nonabelian, as witnessed by its quotient $S_3$, generated by $(12)$ and $(123)$. A closed genus-$g$ group for $g\ge2$ maps onto the free group on $a_1,\ldots,a_g$ by killing all the $b_j$. The same finite quotient witnesses noncommutativity after completion. For $\nu\ge2$, each peripheral generator belongs to a free basis obtained by eliminating a different peripheral generator. For $\nu=1$, $g\ge1$, the nonempty cyclically reduced word $\prod[a_j,b_j]$ is nontrivial and has infinite order. Its image has infinite order in every pro-$p$ completion: map the first pair into the upper unitriangular $3\times3$ matrices over $\Z_p$, and all other pairs to $1$; the commutator is $1+E_{13}$. This also proves that its procyclic closure in the full completion is $\widehat\Z$, since every prime component injects. At $(0,1)$ the peripheral word is $1$, so the source map $\Z\to1$ is not injective. These calculations prove the correction $\nu\le2$ to the bracket $\nu\le3$ on scan 116. This is mathematical evidence of an error, not a changed reading of the visible numeral.
\end{proof}
\noindent\emph{Manuscript:} \S12, folios 44--48, scans 91--100; \S13 bis, folios 56--57, scans 115--118.

\begin{claim}{A.2}{Proved}{Canonical comparison of the orientation modules}
Assume $\nu>0$ and $\pi\ne1$. There is a canonical exact sequence
\[
0\longrightarrow T_\pi\longrightarrow\bigoplus_{i\in S}L_i
 \longrightarrow\pi^{\mathrm{ab}}\longrightarrow H_1(X;\Z)\longrightarrow0,
\qquad
T_\pi=\ker\left(\bigoplus_iL_i\longrightarrow\pi^{\mathrm{ab}}\right).
\]
Each projection $T_\pi\to L_i$ is an isomorphism of rank-one integral modules. If $M\subset U$ is obtained by removing small open discs at the punctures, the boundary morphism and evaluation give canonical isomorphisms
\[
 H_2(M,\partial M;\Z)\xrightarrow{\sim}T_\pi,
 \qquad
 T_\pi\xrightarrow{\sim}\Hom\bigl(H^2_c(U;\Z),\Z\bigr).
\]
For $\nu=0$, $g\ge1$, the corresponding formula is
$T_\pi=\Hom(H^2(\pi;\Z),\Z)=H_2(X;\Z)$.
Filling the discs identifies both orientation lines with $H_2(X;\Z)$, naturally for homeomorphisms of the marked surface. All these comparisons persist with $\widehat\Z$ and continuous relative group cohomology. At $(0,1)$ the kernel formed from the image subgroup $L_1=1$ is zero, although $H^2_c(\R^2;\Z)=\Z$; the source peripheral morphism is therefore essential. At $(0,0)$ group cohomology $H^2(1;\Z)=0$ cannot supply the surface orientation line.
\end{claim}
\passage{13 bis; 56--57; 116--118; 14; 64 verso; 130}

\begin{proof}
The relevant portion of the integral homology sequence of the compact pair is
\[
0\to H_2(M,\partial M)\xrightarrow{\partial}
H_1(\partial M)\to H_1(M)\to H_1(M,\partial M).
\]
The initial zero follows from $H_2(M)=0$. A relative fundamental class has boundary the sum of the consistently oriented boundary circles. Identifying each boundary circle with its peripheral subgroup, the boundary morphism identifies $H_2(M,\partial M)$ with the stated kernel and has degree $1$ on each factor. The cokernel is $H_1(X)$ by attaching the discs and using the homology sequence. Universal coefficients for the pair, whose homology is free, identify $H^2(M,\partial M)$ with the dual of $H_2(M,\partial M)$; excision identifies the former with $H^2_c(U)$. Let $D\subset X$ be the union of the filled discs. The canonical isomorphisms $H_2(X)\xrightarrow{\sim}H_2(X,D)\xrightarrow{\sim}H_2(M,\partial M)$ are, respectively, the pair map and excision. Their inverse is the filling comparison, sending a relative fundamental class to the closed fundamental class and proving compatibility without a chosen sign. Conjugating a peripheral subgroup changes its identification by an inner normalizer action, which is trivial by A.3; hence the sequence is independent of paths to the boundary. For the annulus the same conclusion follows directly because its group is abelian.

For clarity, the profinite compact-support notation is the relative continuous cochain complex
\[
 C^*_{\mathrm{rel}}(\widehat\pi;(\widehat L_i);\Z/n)
 =\operatorname{Cone}\left(C^*_{\mathrm{cont}}(\widehat\pi;\Z/n)
 \longrightarrow\bigoplus_iC^*_{\mathrm{cont}}(\widehat L_i;\Z/n)\right)[-1].
\]
For $\nu>0$, $\widehat\pi$ has cohomological dimension one and each $\widehat L_i\simeq\widehat\Z$; thus its degree-two group is the cokernel of restriction on $H^1$. In the displayed presentation the restriction matrix is the transpose of the integral peripheral inclusion matrix. Its cokernel is canonically $\Hom(T_\pi,\Z/n)$. Taking the inverse limit over $n$ gives the claimed duality. For a closed surface of genus at least two, the continuous duality pairing identifies $H^2(\widehat\pi;\widehat\Z)$ with the completed surface orientation cohomology; this is the surface Poincare duality statement in \cite[Proposition 1.3]{A_Mochizuki2007}. At genus one, for each prime $p$ the Koszul resolution of the trivial module over $\Z_p[[\Z_p^2]]$ has differentials given by the two generators minus $1$. Applying trivial coefficients $\Z/p^r$ makes both differentials zero, so the cup product identifies degree two with the exterior square of degree one, exactly as for the discrete torus. The other prime factors have zero positive-degree cohomology with these coefficients. The Chinese remainder theorem gives the comparison for every finite coefficient $\Z/n$. The transition maps on degree one are surjective, so their first derived inverse limit vanishes; consequently passing to inverse limits introduces no extra degree-two term. The canonical map from the discrete presentation to its completion induces the identity on the surface relation and on every peripheral source; consequently the finite-coefficient comparisons commute with filling. The two trivial-group exceptions follow by the explicit cohomology calculations in the statement.
\end{proof}
\noindent\emph{Manuscript:} \S13 bis, folios 56--57, scans 116--118; \S14, folio 64 verso, scan 130. We read $H^2_!$ on scan 130 as the relative complex just specified, rather than as ordinary cohomology of the underlying abstract profinite group.

\begin{claim}{A.3}{Proved}{Discrete peripheral rigidity in every hyperbolic type}
For $2g+\nu>2$, the group $\pi=\pi_{g,\nu}$ has trivial centre and $C_\pi(H)=1$ for every finite-index subgroup $H$. For every peripheral subgroup and every integer $n\ge1$,
\[
 N_\pi(L_i)=C_\pi(L_i)=C_\pi(L_i^n)=L_i.
\]
Distinct peripheral indices give distinct, noncommensurable conjugacy classes.
\end{claim}
\passage{12; 47 verso--48; 98--99; 13 bis; 57; 117--118}

\begin{proof}
For $\nu>0$, the free action of a free group on its Cayley tree shows that a nonidentity element has an axis and that its centralizer is the infinite cyclic group of translations generated by its primitive root. An element normalizing this cyclic group preserves the axis. A reversal of the axis fixes a vertex or an edge midpoint, which is incompatible with the free, inversion-free action; the normalizer is the same cyclic group. For $\nu\ge2$ each peripheral element is primitive in a free basis. For $\nu=1$, the cyclic word $\prod[a_j,b_j]$ contains the positive letter $a_1$ exactly once. A proper cyclic power has the number of occurrences of every signed letter divisible by its exponent, so this word is not a proper power. This proves the peripheral formulas including the marginal case of scan 117.

For the closed case, the hyperbolic polygon realization of \cite[Theorem 1.2]{A_FarbMargalit2012} embeds $\pi_g$ as a torsion-free cocompact group of orientation-preserving isometries of the hyperbolic plane. Every nonidentity element is hyperbolic. Its centralizer preserves its axis and acts discretely by translations, hence is cyclic. An axis reversal would be an order-two elliptic element and is absent. A finite-index subgroup is a surface group of genus $1+[\pi:H](g-1)\ge2$. In the punctured case the covering has free rank $1+[\pi:H](2g+\nu-2)\ge2$. Thus a finite-index subgroup cannot lie in a cyclic centralizer. This proves $C_\pi(H)=1$, including $H=\pi$.

For distinct punctures with $\nu\ge3$, their primitive abelianization vectors span distinct lines; for $\nu=2$, $g\ge1$, choose the free basis containing $c_1$. The cyclically reduced lengths of $c_1$ and $c_2=(\prod[a_j,b_j]c_1)^{-1}$ are $1$ and $4g+1$, and their abelianizations force any commensurability to match equal absolute exponents. They cannot be conjugate. This proves the final assertion; at $\nu=1$ it is vacuous.
\end{proof}
\noindent\emph{Manuscript:} \S12, folios 47 verso--48, scans 98--99; \S13 bis, folio 57, scans 117--118.

\begin{claim}{A.4}{Proved}{Exact exceptional-case table}
The discrete and profinite exceptional groups have the following properties; in the profinite column replace $\Z$ by $\widehat\Z$.
\[
\begin{array}{c|c|c|c|c}
(g,\nu)&\pi& Z(\pi)&C_\pi(H)\ (H\text{ finite index/open})&\text{peripherals}\\\hline
(0,0)&1&1&1&\text{none}\\
(0,1)&1&1&1&L_1=1\\
(0,2)&\Z&\Z&\Z&L_1=L_2=\pi\\
(1,0)&\Z^2&\Z^2&\Z^2&\text{none}
\end{array}
\]
Consequently the peripheral normalizer and power-centralizer equalities hold in the plane and annulus, while distinctness of the two peripheral classes fails on the annulus. The nonzero vector $1$ on the annulus and $(1,0)$ on the torus witness failure of both centre and finite-index-centralizer assertions. The plane's failure is the noninjective source map, not a nontrivial centre. For the standard base-element formula of scan 501, the annulus satisfies the equality since its basis element generates its whole group. On the torus $C(\langle(1,0)\rangle^n)=\Z^2$ for every $n\ge1$, and $(0,1)$ witnesses its failure to equal $\langle(1,0)\rangle$. The same vector works after completion. The sphere and plane have no nontrivial base elements to test.
\end{claim}
\passage{20; 111--115 verso; 222--230; 12; 47 verso; 98; 13 bis; 56 verso; 116}

\begin{proof}
A.1 identifies the groups and peripheral images. Every subgroup of an abelian group is centralized by the entire group. The assertions now follow by direct substitution. For the profinite version the listed elements remain nonzero, as their reductions modulo $2$ show.
\end{proof}
\noindent\emph{Manuscript:} \S20, folios 111--115 verso, scans 222--230; omissions on scans 98 and 116.

\begin{claim}{A.5}{Proved}{Peripheral decomposition groups and transporters}
Let $1\to\pi\to E\xrightarrow{p}\Gamma\to1$ arise from an action on a hyperbolic marked surface. Choose one peripheral subgroup $L_s$ and the corresponding peripheral lift group $E_s\subset E$; thus $E_s\cap\pi=L_s$ and $p(E_s)=\Gamma_s$, the stabilizer of that puncture. Define
\[
\operatorname{Trans}_E(L_s^d,L_s)
 =\{e\in E:eL_s^de^{-1}\subseteq L_s\}\qquad(d\ge1).
\]
Then $N_E(L_s)=\operatorname{Trans}_E(L_s^d,L_s)=E_s$. The same assertions hold for the corresponding profinite extension and closed peripheral groups. On the annulus these assertions fail for $z\mapsto1/z$: with $E=\Z\rtimes C_2$ and the involution acting by inversion, $E_s=\Z$ but $N_E(L_s)=E$ and every transporter is $E$. More generally for every annulus action both normalizer and transporter equal $E$, while $E_s$ is the inverse image of the end stabilizer. Thus their equality with $E_s$ holds exactly when the action fixes both ends. On the one-puncture sphere the peripheral mixed group maps onto $E=\Gamma$ with kernel $\Z$; using its image, both equalities hold, whereas identifying that mixed group itself with a subgroup is invalid. The closed sphere and torus have no peripheral statement.
\end{claim}
\passage{9; 33; 69; 12; 48; 99}

\begin{proof}
A normalizer element maps to $\Gamma_s$, since the peripheral conjugacy classes are distinct. Choose $e_s\in E_s$ with the same image. Then $e_s^{-1}e$ belongs to $N_\pi(L_s)=L_s$, proving the normalizer equality. For the transporter assertion, $eL_s^de^{-1}\subset L_s$ implies $L_s^d\subset e^{-1}L_se$. The latter is abelian, so it lies in $C_\pi(L_s^d)=L_s$, by A.3 or A.7. In the discrete case a peripheral generator cannot be a proper power, so the inclusion of peripheral groups $e^{-1}L_se\subset L_s$ is an equality. In the profinite case use compactness of $E$: a closed subgroup $L$ of a compact group cannot be properly contained in an inner conjugate. Indeed $aLa^{-1}\subset L$ gives $a^kLa^{-k}\subset L$ for all $k\ge1$, and a net of positive powers of $a$ converges to $a^{-1}$ inside its procyclic closure. Passing to the limit gives the reverse conjugation inclusion and hence equality. Thus the transporter normalizes $L_s$ in both cases. This argument only uses the exact power-centralizer statement and avoids replacing the one-sided edge statement in \cite[Proposition 1.2(i)]{A_Mochizuki2007} by a stronger unstated assertion about chosen subgroup representatives. The peripheral lift groups for the action on a punctured disc map onto the puncture stabilizer by construction, so the preceding choice of $e_s$ is legitimate. For the annulus the involution swaps the two punctures, whose inertia subgroups both equal $\pi$. Its lift is the explicit element in the normalizer outside $E_s$. Replacing $\Z$ with $\widehat\Z$ gives the completed counterexample.
\end{proof}
\noindent\emph{Manuscript:} \S9, folio 33, scan 69; \S12, folio 48, scan 99. In the plane, the map from the peripheral mixed group to $E$ has a kernel $\Z$; writing that mixed group as a subgroup would already discard essential data.

\begin{claim}{A.6}{Proved}{Centre-free extension reconstruction and its abelian correction}
For a centre-free discrete group $\pi$ and a homomorphism $\rho:\Gamma\to\Out(\pi)$, the extension inducing $\rho$ is, up to a unique isomorphism inducing the identity on its prescribed kernel and quotient,
\[
 E_\rho=\{(a,\gamma)\in\Aut(\pi)\times\Gamma:[a]=\rho(\gamma)\}.
\]
The same formula holds for a topologically finitely generated centre-free profinite $\pi$ and continuous $\rho$, with the profinite topology. For $A=\Z$ or $\Z^2$ and a prescribed action $\rho:\Gamma\to\Aut(A)$, extension classes with their kernel and quotient identified are instead classified by $H^2(\Gamma;A_\rho)$. Their automorphisms inducing the identity on both ends form $Z^1(\Gamma;A_\rho)$. For finite $\Gamma$ this statement holds with $A=\widehat\Z$ or $\widehat\Z^2$ and continuous cochains. Thus both the annulus and the torus require the abelian correction.
\end{claim}
\passage{12; 47 verso; 98; 17; 77--80; 155--160}

\begin{proof}
An extension $E$ maps to $E_\rho$ by $e\mapsto(\operatorname{conj}_e,p(e))$. Its kernel is $Z(\pi)=1$. Given $(a,\gamma)$, choose $e_0$ above $\gamma$. The automorphism $a\operatorname{conj}_{e_0}^{-1}$ is inner, implemented by a unique element of $\pi$, so the map is onto. An isomorphism inducing the identity on both ends fixes both components of this formula and is unique. In the profinite case this continuous bijection between compact Hausdorff groups is an isomorphism.

For abelian $A$, choose a set-theoretic section $t$ of the quotient with $t(1)=1$ and write $t(g)t(h)=c(g,h)t(gh)$. Associativity is exactly
$c(g,h)+c(gh,k)=\rho(g)c(h,k)+c(g,hk)$. Changing $t$ by a function $b:\Gamma\to A$ adds its coboundary. Conversely this identity makes $(a,g)(b,h)=(a+\rho(g)b+c(g,h),gh)$ associative. An automorphism fixing both ends has the form $(a,g)\mapsto(a+b(g),g)$, where the homomorphism equation says $b$ is a $1$-cocycle. For finite $\Gamma$ all choices of section are continuous. For example $\Gamma=C_2$, $A=\Z$ with trivial action has the two extension classes $\Z\times C_2$ and $\Z$, with kernel $\Z\xrightarrow{\times2}\Z$ in the latter. They induce the same outer action and are not isomorphic as extensions; reduction modulo $2$ detects the different value of $t(\sigma)^2$.
\end{proof}
\noindent\emph{Manuscript:} \S12, folio 47 verso, scan 98; \S17, folios 77--80, scans 155--160.

\begin{claim}{A.7}{Known}{The full profinite rigidity statements, with reductions}
For the full profinite completion $\Pi=\widehat\pi_{g,\nu}$ of a hyperbolic type, $\Pi$ is torsion-free, $C_\Pi(H)=1$ for every open subgroup $H$, and
\[
N_\Pi(\widehat L_i)=C_\Pi(\widehat L_i^n)=\widehat L_i
\quad(n\in\Z_{>0}).
\]
Distinct peripheral classes are noncommensurable. The same conclusions hold for the pro-$p$ completions, for every prime $p$, including $2$ and $3$. The four exceptional types have exactly the completed table of A.4, and A.5 supplies the extension and transporter assertions.
\end{claim}
\passage{13 bis; 57 verso; 118; 35; 253--254; 499--501}

\begin{proof}
Take the full formation of finite groups, respectively the formation of finite $p$-groups, in \cite[Remark 1.2.2 and Proposition 1.4]{A_MochizukiTamagawa2008}. A hyperbolic complex curve with the prescribed topological presentation identifies its geometric group with $\Pi$, respectively the pro-$p$ completion. Proposition 1.4 asserts precisely the vanishing of centralizers of open subgroups, with no properness assumption, so taking $H=\Pi$ gives centre-freeness. Remark 1.2.2 gives finite $\ell$-cohomological dimension for every relevant prime, hence excludes an element of order $\ell$: such an element would give a closed $C_\ell$ subgroup, whose $\ell$-cohomological dimension is infinite. Every nontrivial finite-order element has a prime-order power; this proves torsion-freeness.

For peripherals, \cite[Lemma 1.3.7]{A_Mochizuki2004} gives $\operatorname{Comm}_\Pi(\widehat L_i)=\widehat L_i$. A normalizer is contained in that commensurator. An element centralizing $\widehat L_i^n$ conjugates the finite-index subgroup $n\widehat L_i$ to itself, so it too belongs to the commensurator. Conversely the procyclic group $\widehat L_i$ centralizes every power. This proves both equalities. In the semi-graph of the smooth pointed curve, peripherals are the cusp edge groups. Proposition 1.2(i) of \cite{A_Mochizuki2007} says that an intersection open in one edge group identifies the underlying edge; hence distinct punctures cannot have commensurable inertia groups. This gives the needed distinction of peripheral indices; A.5 proves its transporter assertion separately using power-centralizers and compactness. No claim about centralizers of arbitrary non-peripheral elements is needed here.

For $\nu\ge2$ there is also a direct source for the power-centralizer assertion: each peripheral generator can be made a free basis element, and \cite[Lemma (2.1.2)]{A_Nakamura1994} gives the centralizer of every nonidentity element of its procyclic closure. For $\nu=1$ the commensurator argument covers precisely the case that this basis argument does not cover. For every $\nu\ge1$, each handle generator $a_j,b_j$ belongs to the free basis obtained by eliminating $c_\nu$. Applying that same lemma proves equation (4) of scan 501 for all these base elements, including the once-punctured case. The basis condition is material if the phrase ``systeme de generateurs on scan 500 is read literally without a cardinality restriction: in a free profinite group on $a,b$, the redundant generating list $a,b,a^2$ has $C(a^2)=\overline{\langle a\rangle}\supsetneq\overline{\langle a^2\rangle}$. The last index is $2$, detected in abelianization modulo $2$. For the pro-$p$ version use $a,b,a^p$ and index $p$. The stated source theorem applies to a member of a free basis; it does not supply the assertion for every redundant generating list.
\end{proof}
\noindent\emph{Manuscript:} \S13 bis, folio 57 verso, scan 118; \S35, folios 253--254, scans 499--501.

\begin{claim}{A.8}{Proved}{Lemme 5 from open-subgroup centralizers alone}
Let $P$ be a discrete or profinite group for which the centralizer of every finite-index, respectively open, subgroup is trivial. Let $H\subset P$ be such a subgroup and let $u\in\Aut(P)$ preserve $H$. An automorphism restricting to the identity on $H$ is the identity on $P$. An automorphism whose restriction equals $\operatorname{conj}_h$ for some $h\in H$ equals $\operatorname{conj}_h$ on $P$.
\end{claim}
\passage{35; 252 verso--254; 498--501}

\begin{proof}
Take the finite-index normal core $N=\bigcap_{x\in P}xHx^{-1}\subset H$. If $u|_H=\id$, then for $x\in P$ and $n\in N$,
\[
 u(x)nu(x)^{-1}=u(xnx^{-1})=xnx^{-1}.
\]
Thus $x^{-1}u(x)\in C_P(N)=1$ for every $x$. For the second assertion apply the first to $\operatorname{conj}_{h^{-1}}u$. In the profinite case the core is an intersection of finitely many open subgroups and is open. This proof uses neither peripheral generators nor their power-centralizers.
\end{proof}
\noindent\emph{Manuscript:} \S35, folio 252 verso, scan 498, Lemme 5; folios 253--254, scans 500--501, proposed proof.

\begin{claim}{A.9}{Proved}{Faithfulness ascends along a finite etale cover}
Let $k$ be a characteristic-zero field, $U/k$ a hyperbolic geometrically connected smooth curve, and $V\to U$ a finite etale geometrically connected covering defined over $k$. Put $\Gamma=\operatorname{Gal}(\overline k/k)$ and choose compatible geometric base points. Faithfulness of $\Gamma\to\Out(\pi_1(U_{\overline k}))$ implies faithfulness of $\Gamma\to\Out(\pi_1(V_{\overline k}))$.
\end{claim}
\passage{35; 253; 499--500}

\begin{proof}
The covering gives an open embedding $\Pi_V\hookrightarrow\Pi_U$ of arithmetic fundamental groups, compatible with their surjections onto the same $\Gamma$. Their geometric kernels are an open subgroup $H\subset P$. Let $\gamma\in\Gamma$ act trivially in $\Out(H)$ and choose a lift $e\in\Pi_V$. Conjugation by $e$ preserves $H$ and restricts there to conjugation by some $h\in H$. By A.7 and A.8 its action on $P$ is also conjugation by $h$. Thus $\gamma$ acts trivially in $\Out(P)$ and equals $1$. A change of compatible base points conjugates these homomorphisms and does not change their kernels. The proof does not use the structure of finite subgroups of $\Gamma$.
\end{proof}
\noindent\emph{Manuscript:} \S35, folio 253, scans 499--500, Corollaire 2.

\begin{claim}{A.10}{Known}{Closed-surface base elements: two separate source reductions}
Let $\Pi=\widehat\pi_g$ with $g\ge2$, let $x$ be any of the standard surface generators $a_i,b_i$, and let $\Lambda=\overline{\langle x\rangle}$. Then $\Lambda\simeq\widehat\Z$ and, for every positive integer $n$,
\[
 C_\Pi(\Lambda^n)=N_\Pi(\Lambda^n)=\Lambda.
\]
\end{claim}
\passage{35; 254; 501}

\begin{proof}
The standard generator maps to a primitive vector of $\Pi^{\mathrm{ab}}=\widehat\Z^{2g}$; hence its procyclic closure injects there and is $\widehat\Z$. For the first source reduction, $x$ is represented by a nonseparating simple closed curve, and is therefore algebraically simple in the terminology of \cite[Definition 1.3]{A_BoggiZalesskii2019}. Proposition 3.6 of that paper gives $N_\Pi(\Lambda^n)=\Lambda$. Centralizers lie in normalizers and contain $\Lambda$, proving the centralizer equality.

For the second reduction, pinch that curve on the genus-$g$ surface. The resulting stable one-node curve has a normalization of genus $g-1$ with two marked inverse branches; at $g=2$ this is a twice-marked genus-one component, still hyperbolic. Its associated semi-graph of anabelioids has one edge group whose image in the smoothing's profinite group is exactly $\Lambda$, up to the chosen base-path conjugation. Proposition 1.2(ii) of \cite{A_Mochizuki2007} makes this edge group commensurably terminal. An element normalizing $\Lambda^n$ commensurates $\Lambda$, proving the same normalizer equality. These are independent applications of the two stated source results, and neither replaces a closed surface group by a free group.
\end{proof}
\noindent\emph{Manuscript:} \S35, folio 254, scan 501, equation (4).

\begin{claim}{A.11}{Proved}{What the abstract profinite presentation specifies}
The full profinite presentation required by scan 318, equation (4), already determines a profinite group with loops isomorphic to $\widehat\pi_{g,\nu}$ with its standard loops. Hence A.7 and A.10 apply to every abstract object satisfying that presentation axiom. Merely requiring generators to satisfy the displayed equation would not suffice: $P=\widehat\Z^2$, $c_1=(1,0)$, $c_2=(0,1)$, $c_3=(-1,-1)$ satisfies the $(0,3)$ relation, has three injective procyclic peripheral subgroups and the rank-one orientation kernel, but has centre $P$.
\end{claim}
\passage{26; 161; 318--319}

\begin{proof}
The universal property gives a continuous homomorphism from the quotient of the free profinite group by the closed normal closure of the relator to $P$. The assertion that this is a defining relation makes this homomorphism an isomorphism; its inverse sends the specified base to the standard base. This gives exactly the comparison map in scan 319, equation (5). For the weaker condition the explicit abelian example has $c_1+c_2+c_3=0$ and generates $P$, but has the additional relation $[c_1,c_2]=1$. Its peripheral inclusion matrix has kernel generated by $(1,1,1)$, and its centre contains $(1,0)\ne0$. The finite reductions and the rank calculation are checked by the accompanying code.
\end{proof}
\noindent\emph{Manuscript:} \S26, folio 161, scans 318--319, equations (4)--(5). The words ``relation generatrice'' are indispensable; weakening them changes the objects by the explicit quotient above.

\begin{claim}{A.12}{Proved}{The discrete filling isomorphism and all compatibility maps}
Let $U=X\setminus S$ be hyperbolic, choose $b\in U$, and put $U'=U\setminus\{b\}$, $\pi'=\pi_1(U')$, and $\pi=\pi'/\langle\!\langle L_b\rangle\!\rangle$. Choose the peripheral representative $L_b\subset\pi'$. The normalization of representatives defines an isomorphism of discrete groups
\[
 d_b:\Out_{\mathrm{lac}}(\pi';b)\xrightarrow{\sim}\Aut_{\mathrm{lac}}(\pi).
\]
On both sides old punctures may be permuted and orientation may be reversed. It identifies the pure orientation-preserving subgroups, whose common index is $2\nu!$, where $\nu=|S|$. It commutes with forgetting the distinguished puncture, with the maps to $\Out_{\mathrm{lac}}(\pi)$, with orientation characters, with the permutation maps to $\mathfrak S_\nu$, and with point-pushing and inner conjugation on their kernels.
\end{claim}
\passage{18; 93--94; 186--189}

\begin{proof}
An outer class fixing $[L_b]$ has a representative $\alpha$ preserving $L_b$ itself: conjugate any representative by an element taking its peripheral image back to $L_b$. Two normalized representatives differ by conjugation by an element of $N_{\pi'}(L_b)=L_b$, by A.3. Such conjugation becomes the identity after killing $L_b$. Thus $\alpha$ induces a well-defined actual automorphism $d_b([\alpha])$ of the specified quotient $\pi$. Products of normalized representatives are normalized, so $d_b$ is a homomorphism. Its definition also proves compatibility with the permutation of the remaining loops and with the orientation module, using A.2's filling comparison.

We give the group-theoretic identification of this homomorphism with the topological one. The Dehn--Nielsen--Baer theorem identifies the extended mapping class group of a hyperbolic punctured surface with the outer automorphisms preserving its peripheral conjugacy classes, and in the closed case with $\Out(\pi)$; the precise versions are \cite[Theorems 8.1 and 8.8]{A_FarbMargalit2012}. The evaluation map of the homeomorphism group of $(X,S)$ at $b$ is a locally trivial fibration over $U$: over a sufficiently small disc, move $b$ to a nearby point by a homeomorphism supported in that disc, which constructs a local section. Its connecting map pushes $b$ round a loop. The identity component of this homeomorphism group is contractible in hyperbolic type, by the surface homeomorphism-group theorem in \cite[Lemma 2.1 and Theorem 1.1]{A_Yagasaki2000}; hence the homotopy sequence gives
\[
1\to\pi\xrightarrow{\mathrm{push}}\Out_{\mathrm{lac}}(\pi';b)
\xrightarrow{\mathrm{fill}}\Out_{\mathrm{lac}}(\pi)\to1.
\]
One may alternatively use the extended Birman exact sequence, \cite[Theorem 4.6 and the immediately following extended-group remark]{A_FarbMargalit2012}. Surjectivity follows directly by moving the image of $b$ back to $b$ along an arc avoiding $S$.

Choose the convention that the point-push associated to a based loop $\gamma$ moves the base-point track in the direction for which the induced based map is $\eta\mapsto\gamma\eta\gamma^{-1}$. This can also be checked on the boundary of the homotopy square traversed by $\eta$ under the isotopy. It proves $d_b\circ\mathrm{push}=\operatorname{conj}$. The filled mapping class gives the same outer automorphism as $d_b$ by construction. Thus the last sequence maps to
\[
1\to\pi\xrightarrow{\operatorname{conj}}\Aut_{\mathrm{lac}}(\pi)
\to\Out_{\mathrm{lac}}(\pi)\to1
\]
by the identity on the end groups. Injectivity and surjectivity of $d_b$ follow by the elementary exact-sequence argument: a kernel element belongs to point-pushing and maps trivially by the injective inner map, and any target element can first be matched in the outer quotient and then in the kernel.

Finally every permutation of old punctures is realized by half-twists supported in small discs along arcs, and an orientation-reversing homeomorphism can be chosen to fix the finite marked set pointwise. Hence the map to $\{\pm1\}\times\mathfrak S_\nu$ is onto; its kernel is precisely the pure orientation-preserving group. Its index is $2\nu!$. This also identifies the subgroups denoted $\mathcal T^{00}$ on scan 188 and proves the compatibility checks left on scan 189.
\end{proof}
\noindent\emph{Manuscript:} \S18, folios 93--94, scans 186--189. On scan 187 the inequality must be $2g+\nu'\ge4$ when $\nu'$ counts punctures before filling; the meaningful alternative $\le4$ contradicts both examples and the stated hyperbolicity requirement. The mathematical argument supports the first alternative without settling the handwriting.

\begin{claim}{A.13}{Proved}{Filling into the two abelian targets}
Exactly one of the drilled and filled groups is abelian in the two cases $(0,3)\to(0,2)$ and $(1,1)\to(1,0)$. In the first, the homomorphism on automorphisms of the underlying groups with peripheral subgroup classes is
\[
 C_2\times C_2=\Out_{\mathrm{lac}}(F_2;[\langle xy\rangle])
 \longrightarrow\Aut(\Z)=C_2,
\]
with kernel of order two. Its nontrivial element has representative $x\mapsto y^{-1}$, $y\mapsto x^{-1}$. Its restriction to the positive-orientation subgroup is an isomorphism. Retaining the two indexed peripheral source maps and their orientation line on the annulus enlarges the target automorphism group to $C_2\times C_2$ and restores the full isomorphism. In the second case the filling map is the isomorphism $\operatorname{GL}_2(\Z)\to\operatorname{GL}_2(\Z)$, with positive subgroup $\operatorname{SL}_2(\Z)$.
\end{claim}
\passage{18; 94--97; 189--194}

\begin{proof}
A.1 lists the two cases. For a three-punctured sphere the extended mapping class group is $\mathfrak S_3\times C_2$: a homeomorphism fixing the three punctures is isotopic, in each orientation class, to the identity or a fixed reflection, by cutting along three joining arcs and applying the disc isotopy. The stabilizer of the distinguished puncture is $C_2\times C_2$. On $F_2=\langle x,y\rangle$, the old-puncture transposition $x\leftrightarrow y$ and the reflection $x\mapsto x^{-1},y\mapsto y^{-1}$ both induce inversion after imposing $xy=1$. Their product is the displayed nontrivial kernel element. It takes the distinguished loop to its inverse and reverses orientation. Direct free-word and finite-group calculations appear in the attached program.

For an annulus with ordered end sources $\kappa_1(1)=t$ and $\kappa_2(1)=t^{-1}$, an automorphism includes a sign $\epsilon\in\{\pm1\}$ on $T$ and a permutation $s\in\mathfrak S_2$. Compatibility forces the underlying map $t\mapsto t^{\epsilon\operatorname{sgn}(s)}$. Thus there are four automorphisms of this full data, and the map forgetting $s,\epsilon$ is exactly the preceding quotient. This calculation specifies the information lost by retaining only subgroup images.

For the torus, every matrix in $\operatorname{GL}_2(\Z)$ is realized by a linear torus homeomorphism fixing the origin. A based torus homeomorphism inducing the identity matrix is isotopic relative to that point to the identity. Indeed the restriction of evaluation to the translation subgroup $\R^2/\Z^2$ is the identity of the torus, so evaluation is surjective on fundamental groups and the point-pushing kernel is zero. The unbased kernel is the identity component: isotope the images of two primitive curves intersecting once back to those curves, use the identity matrix to remove their remaining twists, and cut along them to a disc, where the Alexander isotopy finishes the argument. Removing the fixed point identifies this mapping class group with the peripheral-preserving outer group of the once-punctured torus by Dehn--Nielsen--Baer. Filling induces abelianization, hence the asserted identity map on matrices. Its orientation sign is the determinant.
\end{proof}
\noindent\emph{Manuscript:} \S18, folios 94--97, scans 189--194. The abelian cases in the margin of scan 194 refer to the filled types; the corresponding drilled groups are both $F_2$.

\begin{claim}{A.14}{Known}{Primitive loops in finite-step solvable quotients}
Let $\Sigma$ be a nonempty set of primes, $F$ a free pro-$\Sigma$ group of finite rank at least two, and $F^{[j]}$ its closed derived series. Put $F_m=F/F^{[m]}$, $m\ge2$. For any element $x$ of a free basis and any positive integer $n$,
\[
 C_{F_m}(x^n)=N_{F_m}(\overline{\langle x\rangle})
 =\overline{\langle x\rangle}\simeq\Z_\Sigma,
 \qquad \Z_\Sigma=\prod_{p\in\Sigma}\Z_p.
\]
Every open subgroup of $F_m$ has trivial centralizer. For $\nu\ge2$, all peripheral loops of the pro-$\Sigma$ group of type $(g,\nu)$ satisfy these formulas in every $F_m$. Their indexed conjugacy classes are pairwise distinct in all hyperbolic types.
\end{claim}
\passage{13 bis; 57--57 verso; 117--118; strengthened PartA(c)}

\begin{proof}
The power-centralizer formula is \cite[Theorem 1.5, finite-rank case]{A_Yamaguchi2026}; the source explicitly treats every nonzero integer power, including powers divisible by $2$ or $3$. We checked its finite-rank proof: Proposition 1.3 uses the finite semidirect product $(\Z/N)[G/\langle\bar x\rangle]\rtimes(G\times\langle\bar x\rangle)$ to force the centralizer into $\langle x\rangle F_m^{[m-1]}$; the coefficient at the distinguished coset forces the first $G$ coordinate into $\langle\bar x\rangle$. At $m=2$, the Blanchfield--Lyndon injection of $F'/F''$ into $\Z_\Sigma[[F^{\mathrm{ab}}]]^r$ and injectivity of multiplication by $\bar x^n-1$ remove the intersection with $F_m'$. That injectivity is proved by passing from $C_M$ to $C_{kM}$: invariant coefficients downstairs are divisible by every $\Sigma$-integer $k$, so are zero. The finite-rank induction passes to finite abelian covers where a positive power of $x$ is a free basis element. These are exactly the finite-rank hypotheses used here; no infinite-rank inverse-limit assertion is used.

If $u$ normalizes $\langle x\rangle$, write $uxu^{-1}=x^\lambda$, $\lambda\in\Z_\Sigma^\times$. The primitive abelianization coordinate gives $\lambda=1$, so the normalizer lies in the already computed centralizer. An open subgroup contains positive powers of two distinct free generators; their cyclic centralizers intersect trivially because their abelianization coordinates are independent. This proves the open-centralizer assertion directly.

Each peripheral generator for $\nu\ge2$ belongs to a free basis after eliminating another boundary generator. For $\nu\ge3$ distinct peripheral abelianization lines separate the conjugacy classes. When $\nu=2$ and $g\ge1$, kill $c_1$ and all but one handle pair, mapping onto a free metabelian group on $a,b$. The image of $c_1$ is $1$ and that of $c_2$ is $[a,b]^{-1}\ne1$; the latter has infinite order in the upper unitriangular $3\times3$ group over $\Z_\Sigma$. Hence they cannot be conjugate, or commensurable with finite indices. These maps factor through every $F_m$, $m\ge2$.
\end{proof}
\noindent\emph{Manuscript:} \S13 bis, scans 117--118, and the strengthened finite-step task in Part A(c).

\begin{claim}{A.15}{Proved}{The compact intersection calculation used in descent}
Let $(H_\alpha)$ be a decreasing directed family of closed subgroups of a profinite group $P$, and put $K=\bigcap_\alpha H_\alpha$. Then
\[
 K^{[j]}=\bigcap_\alpha H_\alpha^{[j]}\quad(j\ge0).
\]
For a closed subgroup $L\subset P$ and decreasing closed normal subgroups $N_\alpha$ of $P$, one also has $\bigcap_\alpha LN_\alpha=L(\bigcap_\alpha N_\alpha)$.
\end{claim}
\passage{Auxiliary compact-group lemma for finite-step PartA(c)}

\begin{proof}
For any open normal $V\subset P$, the images of $H_\alpha$ in the finite group $P/V$ eventually stabilize. Their eventual image equals that of $K$: each compatible coset fibre is a decreasing family of nonempty compact sets. Taking a fixed number of derived subgroups commutes with homomorphic images, and taking the inverse limit of all finite images recovers the closed derived subgroup. This proves the first equality. For the second equality, an element $z$ of every $LN_\alpha$ gives nonempty closed sets $L\cap zN_\alpha$; compactness makes their intersection nonempty. A member of that intersection gives the required single factorization of $z$.
\end{proof}

\begin{claim}{A.16}{Proved}{The one-puncture solvable threshold, with a finite certificate}
Let $F$ be the pro-$\Sigma$ completion of $\pi_{g,1}$, $g\ge1$, let $c=\prod_{j=1}^g[a_j,b_j]$ (the inverse peripheral generator), and write $c$ for its image in $F_m=F/F^{[m]}$. For every positive integer $n$,
\[
\begin{array}{ll}
m=2:& C_{F_2}(c^n)=N_{F_2}(\overline{\langle c\rangle})=F_2',\\[2pt]
m\ge3:&C_{F_m}(c^n)=N_{F_m}(\overline{\langle c\rangle})
 =\overline{\langle c\rangle}.
\end{array}
\]
In the first line the inclusion $\overline{\langle c\rangle}\subset F_2'$ is strict. Thus peripheral rigidity fails at the smallest solvable depth $m=2$ and is restored at every $m\ge3$, for the full profinite and every pro-$p$ version.
\end{claim}
\passage{13 bis; 57 verso; 118; strengthened PartA(c)}

\begin{proof}
Let $H\subset F$ run over open normal subgroups containing $F'$, with $|F/H|\ge3$. They form a family whose intersection is $F'$. The corresponding finite abelian covering of the once-punctured surface has exactly $d=|F/H|$ punctures, because $c\in F'$ has trivial image in the deck group. Choose the cusp above the original based peripheral loop. Its loop is $c$ itself; the other loops are its conjugates by deck representatives. The submodule of $H^{\mathrm{ab}}$ generated by these cusps is
\[
 \Z_\Sigma^{F/H}\big/\Z_\Sigma\!\left(\sum_{q\in F/H}e_q\right),
\]
with the deck group acting by left translation. This follows either from the relative surface homology sequence of A.2 or directly from the free-group peripheral presentation. For $d\ge3$ the lines generated by distinct $e_q$ in this module are distinct; moreover $ne_q=ae_1$, for $0\ne n\in\Z$ and $a\in\Z_\Sigma$, forces $q=1$ and $a=n$. To check this last assertion, lift the equality to the free module and compare a coordinate different from both $1$ and $q$; its coefficient forces the possible multiple of $\sum e_q$ to be zero, after which the coordinates at $q$ and $1$ give the result.

For $m\ge2$, $F^{[m]}\subset H'$, so the natural map $H\to H^{\mathrm{ab}}$ factors through the image of $H$ in $F_m$. Let $u\in F$ represent an element of $C_{F_m}(c^n)$. The relation $uc^nu^{-1}=c^n$ modulo $F^{[m]}$ implies $ne_{\bar u}=ne_1$ in the cusp submodule of $H^{\mathrm{ab}}$. Hence $u\in H$ for every such $H$, so its image belongs to $F_m'$. The same argument for a normalizer, using $uc u^{-1}=c^\lambda$ modulo $F^{[m]}$, $\lambda\in\Z_\Sigma^\times$, gives $u\in F'$ and $\lambda=1$. The subgroup $\overline{\langle c\rangle}$ is indeed $\Z_\Sigma$: the unitriangular $3\times3$ quotient used in A.1 proves injectivity of all its prime components. Thus its normalizer equals its centralizer.

When $m=2$, the derived subgroup is abelian, contains $c$, and centralizes it. This gives equality with $F_2'$. Now let $m\ge3$ and let $uF^{[m]}$ centralize $c^n$. We have proved $u\in F'\subset H$. Since $F'\subset H$,
$F^{[m]}=(F')^{[m-1]}\subset H^{[m-1]}$; therefore there is a well-defined continuous epimorphism
\[
 H/F^{[m]}\longrightarrow H/H^{[m-1]}.
\]
The group $H$ is free pro-$\Sigma$, and $c$ is a member of a free basis because this cover has at least two punctures. By A.14, applied at depth $m-1\ge2$, the image of $u$ in this quotient belongs to $\overline{\langle c\rangle}$. Consequently
\[
 u\in\bigcap_H\overline{\langle c\rangle}H^{[m-1]}
 =\overline{\langle c\rangle}\left(\bigcap_H H^{[m-1]}\right)
 =\overline{\langle c\rangle}(F')^{[m-1]}
 =\overline{\langle c\rangle}F^{[m]},
\]
where both interchanges are proved in A.15. This proves the second line.

For strictness in the first line, set $a=a_1$, $b=b_1$, and kill the remaining handle pairs. In $\operatorname{UT}_4(\F_p)$ take
\[
 x=1+E_{12}+E_{23},\qquad y=1+E_{34}.
\]
Their group is metabelian; direct multiplication gives $[x,y]=1+E_{24}$ and $x[x,y]x^{-1}=1+E_{24}+E_{14}$. The second matrix is not a power of the first. Thus $aca^{-1}\in F_2'$ centralizes $c$ but does not belong to its procyclic closure. For every $p\in\Sigma$ this quotient factors through $F_2$ and proves strictness; the attached program independently enumerates and verifies the cases $p=2,3,5$. In particular there is no excluded residue-characteristic caveat.
\end{proof}
\noindent\emph{Manuscript:} \S13 bis, scan 118, in the strengthened finite-step question of Part A(c). This assertion does not assert a new literature priority; the searched sources and the proof audit are listed in the report.

\begin{claim}{A.17}{Proved}{Centre and torsion in all finite-step surface quotients}
Let $P$ be the full profinite or pro-$\Sigma$ completion of a hyperbolic surface group, open or closed, and let $P_m=P/P^{[m]}$. For every $m\ge2$, $P_m$ is torsion-free and centre-free. This verifies the form of \cite[Theorem 2.9 and Corollary 2.10]{A_Yamaguchi2026} needed here by the following argument from finite covers and transfer.
\end{claim}
\passage{Strengthened PartA(c); Yamaguchi2026 Theorem2.9 and Corollary2.10}

\begin{proof}
Every open subgroup of $P$ is the pro-$\Sigma$ group of a finite covering surface, so its abelianization is a free $\Z_\Sigma$-module. A closed subgroup $K$ is the intersection of its open neighbourhood subgroups; by A.15 its abelianization is the inverse limit of their abelianizations. More explicitly, a compatible system in those abelianizations lifts to a compatible coset system in the compact group, and the kernel is the intersection of the closed derived subgroups. Thus $K^{\mathrm{ab}}$ is torsion-free. Applying this to each $P^{[j]}$ and inducting in
\[
1\to(P^{[m-1]})^{\mathrm{ab}}\to P_m\to P_{m-1}\to1
\]
proves torsion-freeness, starting with the torsion-free abelianization.

We also need the faithful action of a finite deck group $Q=H/N$ on $N^{\mathrm{ab}}$, for open $N\trianglelefteq H\subset P$. Choose $\ell\in\Sigma$. If a subgroup $Q_0$ acts trivially on $H^1$ of the covering surface with $\Q_\ell$ coefficients, transfer gives an isomorphism of that $H^1$ with the $H^1$ of the quotient by $Q_0$. The two surfaces are both proper or both nonproper, and their Euler characteristics are consequently equal. Multiplicativity of Euler characteristic gives $\chi(Y)=|Q_0|\chi(Y/Q_0)$. Both are negative, so $Q_0=1$. The comparison with $N^{\mathrm{ab}}\otimes_{\Z_\Sigma}\Q_\ell$ is the canonical duality between first homology and first cohomology. Thus the deck action on $N^{\mathrm{ab}}$ is faithful.

Put $K=P^{[m-1]}$, $A=K/K'$, and $Q=P/K'$. For an open normal subgroup $N\subset Q$ containing $A$, its preimage $\widetilde N\subset P$ contains $K$, hence $K'\subset\widetilde N'$. Therefore $N^{\mathrm{ab}}=\widetilde N^{\mathrm{ab}}$, and the preceding deck faithfulness implies that an element of $Z(Q)$ belongs to every such $N$. Their intersection is $A$.

We finish by proving $A^{P/K}=0$. For any open normal $V\subset P$ containing $K$, write $j:V^{\mathrm{ab}}\to P^{\mathrm{ab}}$ for inclusion and $t:P^{\mathrm{ab}}\to V^{\mathrm{ab}}$ for transfer. The composite $tj$ on the $P/V$-invariants is multiplication by $[P:V]$, which is injective because $V^{\mathrm{ab}}$ is torsion-free. Hence inclusion is injective on those invariants. Passing to the inverse limit over $V$ shows that $K^{\mathrm{ab},P/K}\to P^{\mathrm{ab}}$ is injective. But $K\subset P'$, so this map is zero. Therefore $A^{P/K}=0$. Since $Z(Q)\subset A$, the centre is zero, as claimed. This proof checks precisely the transfer and finite-cover steps used in the cited unrefereed result.
\end{proof}

\begin{claim}{A.18}{Proved}{The remaining affine solvable rigidity list and a closed bound}
For all hyperbolic types with $\nu>0$, and all $m\ge2$, the maximal $m$-step-solvable full profinite or pro-$\Sigma$ quotient has trivial centre and trivial centralizers of open subgroups, and is torsion-free. Peripheral conjugacy classes remain distinct. For $\nu\ge2$ all peripheral normalizer and integer-power-centralizer equalities hold at every $m\ge2$; for $\nu=1$ their exact values are those of A.16. In an extension obtained by quotienting a geometric action extension by the $m$th closed derived subgroup of its kernel, the equality $N_E(L_s)=E_s$ holds in the former cases and at $m\ge3$ in the latter. At $\nu=1,m=2$ its correct value is $N_E(L_s)=F_2'E_s$, with the intersection $F_2'\cap E_s=L_s$. In all these cases every transporter $\operatorname{Trans}_E(L_s^d,L_s)$, $d\ge1$, equals the same displayed normalizer. Since the kernels are centre-free by A.17, the extension reconstruction of A.6 also applies to all these solvable quotients. Standard free-basis generators satisfy the power-centralizer formula at every $m\ge2$ by A.14.
For a closed hyperbolic group $P$ and an open subgroup $H\subset P_m$, the proven universal bound is
\[
 C_{P_m}(H)\subseteq P_m^{[m-1]}.
\]
\end{claim}
\passage{13 bis; 57 verso; 118; 35; 252 verso--254; 498--501; strengthened PartA(c)}

\begin{proof}
The first assertions are A.14, A.16, and A.17. For the extension calculation, the quotient map is injective on the peripheral group, since the peripheral generator has infinite order in $P_2$ by A.1. The image of a geometric peripheral lift group still maps onto the puncture stabilizer. The kernel of the normalizer map to that stabilizer is the normalizer calculated in A.14 or A.16, so the same lifting argument as A.5 proves the equalities. In the metabelian one-puncture case there is just one puncture, so the stabilizer is the whole acting group and multiplication by $E_s$ gives the asserted $F_2'E_s$.

It remains to justify the one-sided transporter equality in these quotients. Suppose $eL_i^de^{-1}\subset L_i$. The action sends $L_i$ to a conjugate of some $L_j$ and its chosen generator to a conjugate of $c_j^\lambda$, with $\lambda\in\Z_\Sigma^\times$. For $\nu\ge3$, the primitive abelianization lines imply $j=i$. For $\nu=2$, $g\ge1$, the quotient killing $c_i$ and all but one handle pair sends $c_j$ to a nonzero commutator of infinite order in $\operatorname{UT}_3(\Z_\Sigma)$; thus $j\ne i$ would contradict the given containment. Choose an element of $E_i$ with the same image as $e$ in the acting group and divide it out. The resulting element $u$ lies in $P_m$ and satisfies $uc_i^du^{-1}=c_i^a$ for some $a\in\Z_\Sigma$. The primitive abelianization coordinate forces $a=d$, so A.14 puts $u$ in $L_i$. This proves the equality for $\nu\ge2$. For $\nu=1$, the acting group already fixes the only puncture. After the same division by $E_i$, apply the finite abelian cover calculation in A.16 to $uc^du^{-1}=c^a$. It forces $a=d$ and reduces the transporter to the integer-power centralizer calculated there. This gives $E_i$ at $m\ge3$ and $F_2'E_i$ at $m=2$.

The closed bound recorded here is strengthened to triviality in A.19 below. For this bound, put $A=P_m^{[m-1]}$ and choose any open normal $N\subset P_m$ containing $A$. The group $N^{\mathrm{ab}}$ is the torsion-free abelianization of its preimage in $P$, as in A.17. If $z$ centralizes $H$, its conjugation action fixes the image of $H\cap N$ in $N^{\mathrm{ab}}$. This image has finite index, so some nonzero integer multiple of every element is fixed. Torsion-freeness then makes the action trivial on all of $N^{\mathrm{ab}}$. Deck faithfulness from A.17 implies $z\in N$. Intersecting all these $N$ yields $z\in A$.
\end{proof}
\noindent\emph{Manuscript:} \S13 bis, scan 118; \S35, scans 498--501; additional finite-step task A(c).

\begin{claim}{A.19}{Proved}{Closed finite-step power centralizers and slimness}
Let $P$ be the pro-$\Sigma$ completion of a closed oriented surface group of genus $g\ge2$, let $m\ge2$, and put $P_m=P/P^{[m]}$. Let $x$ be represented by a nonseparating simple closed curve, in particular by any standard handle generator. Then for every positive integer $n$,
\[
 C_{P_m}(x^n)=N_{P_m}(\overline{\langle x^n\rangle})
 =\overline{\langle x\rangle}\simeq\Z_\Sigma.
\]
Consequently $C_{P_m}(H)=1$ for every open subgroup $H$. This includes the full profinite and every pro-$p$ quotient, without a restriction at $p=2$ or $3$.
\end{claim}
\passage{35; 254; 501; strengthened PartA(c)}

\begin{proof}
We first give the completed module construction required for the metabelian case. Let $F$ be a free pro-$\Sigma$ group on $z_1,\ldots,z_r$, let $A=F/F'$, and put $R=\Z_\Sigma[[A]]$. The continuous homomorphism
\[
 \mu:F\longrightarrow R^r\rtimes A,
 \qquad z_i\longmapsto(e_i,\bar z_i),
\]
exists by freeness; $A$ acts by multiplication on the coefficient module, so the product in the target is $(v,a)(w,b)=(v+aw,ab)$. Its kernel is exactly $F''$. The inclusion $F''\subset\ker\mu$ follows because the target is metabelian. To prove the reverse inclusion without assuming any exactness result about completed chain complexes, take an arbitrary finite metabelian quotient $G$ of $F$. Write $K=G'$, $Q=G/K$, and choose a section $t:Q\to G$ with $t(1)=1$. There is an injective homomorphism
\[
 G\longrightarrow K^Q\rtimes Q,
 \qquad g\longmapsto(f_g,\bar g),\qquad
 f_g(q)=t(q)g\,t(q\bar g)^{-1}.
\]
Indeed $f_{gh}(q)=f_g(q)f_h(q\bar g)$, and, when $\bar g=1$, evaluation at $q=1$ recovers $g$. Use additive notation for the finite abelian group $K$. The finite module $K^Q$ is an $R$-module through $A\to Q$ and right translation on $Q$. Sending $e_i$ to the function coordinate of the image of $z_i$ gives a continuous $R$-linear map $R^r\to K^Q$. The resulting homomorphism of semidirect products makes the homomorphism from $F$ factor through $\mu$, since it agrees on every free generator. Hence $\ker\mu$ maps trivially to every finite metabelian quotient of $F$. These quotients separate $F/F''$, proving the assertion. In particular the function coordinate gives an injective $R$-linear morphism
\[
 F'/F''\lhook\joinrel\longrightarrow R^r.
\]
Its image is a closed $R$-submodule: conjugation gives the $A$-action, integral powers give the scalar action, and continuity extends these to the completed group ring.

Apply this construction with $r=2g$ and
\[
 P=F/\langle\!\langle r_0\rangle\!\rangle,
 \qquad r_0=\prod_{i=1}^g[a_i,b_i],
 \qquad A=P^{\mathrm{ab}}=\Z_\Sigma^{2g}.
\]
The function coordinate of $r_0$ is
\[
 v=(1-\bar b_1,\bar a_1-1,\ldots,1-\bar b_g,\bar a_g-1)\in R^{2g}.
\]
This is obtained by multiplying the four factors of each commutator in $R^{2g}\rtimes A$. Since $r_0\in F'$, its closed normal closure maps in $F'/F''$ onto exactly $Rv$: its conjugates give the $A$-multiples of $v$, their products give their integral sums, and their closure is $Rv$. The last module is closed because it is the image of the compact ring $R$. Thus the natural quotient maps induce the injective comparison
\[
 B:=P'/P''\lhook\joinrel\longrightarrow R^{2g}/Rv.                 \tag{A.19.1}
\]
Indeed $P''$ is the image of $F''$ and the kernel in $F'/F''$ is precisely the just calculated module. This establishes the needed completed Fox-module comparison directly.

Choose a geometric basis with $x=a_1$, and set $h=\bar a_1^n-1$. We use two elementary facts about these completed group rings. First, multiplication by $h$ on $R$ is injective. Second, multiplication by $\bar a_2-1$ on $R/hR$ is injective. Here
\[
 R/hR\xrightarrow{\sim}
 \Z_\Sigma[[A/\overline{\langle\bar a_1^n\rangle}]]                \tag{A.19.2}
\]
is the map induced by the group quotient. Its kernel is the closed ideal generated by $h$; that ideal equals $hR$ since $hR$ is compact. To prove the injectivity assertions, one can use the following coefficient calculation. In a completed group ring of $\Z_\Sigma\times D$, consider an element invariant under the chosen infinite cyclic coordinate, or under its $n$th power. At finite coefficient level $\Z/k$ and finite coordinate level $C_M$, enlarge the coordinate level to $C_{kM}$, taking $M$ divisible by the $\Sigma$-part of $n$ in the second case. Invariance makes the coefficients constant on each orbit. The projection down to $C_M$ sums $k$ equal coefficients on each orbit fibre and therefore is zero modulo $k$. This works with every finite quotient of the remaining factor $D$. All finite images of the original element vanish, so the element is zero. Apply this argument first to the $a_1$ coordinate of $R$, and then to the independent $a_2$ coordinate in the group of (A.19.2). It applies to arbitrary $\Sigma$, because $k$ and the chosen coordinate orders have prime factors in $\Sigma$.

These ring facts show that multiplication by $h$ is injective on $R^{2g}/Rv$. If $hz=vt$, inspect the $b_2$ coordinate, which is $h z_{b_2}=(\bar a_2-1)t$. Reducing modulo $h$ and using the second injectivity gives $t=ht'$ for some $t'\in R$. The first injectivity then gives $z=vt'$ in every coordinate. Consequently (A.19.1) implies
\[
 B^{\langle x^n\rangle}=0.                                      \tag{A.19.3}
\]

For completeness, we next constrain the abelianization of an arbitrary centralizer element, at every depth $m\ge2$. Put $V=\Z_\Sigma^{2g}$ with its ordered handle basis and put $W=\bigwedge^2V/\Z_\Sigma\omega$, where $\omega=\sum_i a_i\wedge b_i$. For every $\Sigma$-integer $k$, form the finite class-two group with underlying set $(V/kV)\times(W/kW)$ and product
\[
 (q,w)(q',w')=(q+q',w+w'+\beta(q,q')),
 \qquad
 \beta(q,q')=\sum_{i<j}q_iq'_j\,z_i\wedge z_j.
\]
Its commutator is $(0,q\wedge q')$. The standard surface relation maps to $\omega=0$, so sending the handle generators to their coordinate vectors defines a quotient homomorphism from $P$. The finite group is metabelian and a $\Sigma$-group; the homomorphism therefore factors through every $P_m$, $m\ge2$. If $u\in P_m$ centralizes $a_1^n$ and has abelianization $q$, these homomorphisms give
$n(q\wedge a_1)=0$ in $W$. The module $W$ is free because $\omega$ is primitive, so multiplication by the nonzero integer $n$ is injective. Write $q\wedge a_1=t\omega$. The coefficient of $a_2\wedge b_2$ on the left is zero; on the right it is $t$. Hence $t=0$, and the remaining coordinates give $q=\lambda a_1$ for some $\lambda\in\Z_\Sigma$. The procyclic group generated by $a_1$ injects into $V$, so dividing $u$ by $a_1^\lambda$ puts it in $P_m'$ without changing its centralization property. At $m=2$, (A.19.3) makes this remaining element trivial, proving the centralizer formula.

We now induct on $m$, simultaneously for every genus at least two and every nonseparating simple closed curve. Suppose $m\ge3$, and divide a centralizer element as just explained so that it is represented by $u\in P'$. Let $H$ run over all open normal subgroups of $P$ containing $P'$, and let $d$ be the order of $x$ in $P/H$. The corresponding finite abelian covering of the closed surface has genus at least two. The based lift of the simple curve $x$ is represented by $x^d$, traversed once. This lift is a simple nonseparating curve: its integral homology class pushes forward to $d[x]\ne0$, whereas a separating curve has zero homology class. It is therefore a standard handle generator for some geometric basis of the covering surface group $H$.

Since $P'\subset H$, we have $P^{[m]}\subset H^{[m-1]}$. The quotient morphism
\[
 H/P^{[m]}\longrightarrow H/H^{[m-1]}
\]
is consequently defined. The image of $u$ centralizes $(x^d)^n$ in its target. By the induction hypothesis at depth $m-1$, that image belongs to $\overline{\langle x^d\rangle}$. Hence
\[
 u\in\bigcap_H\overline{\langle x\rangle}H^{[m-1]}
 =\overline{\langle x\rangle}\left(\bigcap_HH\right)^{[m-1]}
 =\overline{\langle x\rangle}P^{[m]},
\]
by A.15 and $\bigcap_HH=P'$. The intersection of $\overline{\langle x\rangle}$ with $P'$ is trivial, again by its primitive abelianization coordinate. Since $u\in P'$, its image in $P_m$ is trivial. This completes the induction.

If an element normalizes $\overline{\langle x^n\rangle}$, its conjugation sends $x^n$ to $x^{n\lambda}$ for a unit $\lambda\in\Z_\Sigma^\times$. Abelianization gives $n\lambda=n$, whence $\lambda=1$, so the normalizer equals the calculated centralizer. Finally an open subgroup contains positive integer powers of $a_1$ and $b_1$. Their calculated centralizers have trivial intersection because they inject into distinct independent abelianization coordinates. This proves slimness.
\end{proof}
\noindent\emph{Manuscript:} \S35, scan 501, equation (4), and the finite-step strengthening in Part A(c). The proof establishes the exact module comparison and all closed-group descent maps used; it does not rely on the unknown general slimness statement mentioned before Proposition 2.6 in \cite{A_Yamaguchi2026}, nor claim a literature priority.

\begin{claim}{A.20}{Proved}{The typed completion comparison for filling}
For the two hyperbolic groups in A.12, put $\Pi'=\widehat{\pi'}$ and $\Pi=\widehat\pi$. There is a commutative diagram of profinite groups
\[
\begin{array}{ccc}
 \widehat{\Out_{\mathrm{lac}}(\pi';b)}
 &\xrightarrow{\ \widehat{d_b}\ }&
 \widehat{\Aut_{\mathrm{lac}}(\pi)}\\
 \big\downarrow c_{\mathrm{out}}&&\big\downarrow c_{\mathrm{aut}}\\
 \Out_{\mathrm{lac}}(\Pi';b)
 &\xrightarrow{\ d_b^{\mathrm{cont}}\ }&
 \Aut_{\mathrm{lac}}(\Pi).
\end{array}
\]
The upper horizontal map is an isomorphism. The lower horizontal map is the normalized filling homomorphism on all continuous loop automorphisms. The vertical maps are induced by extending discrete automorphisms to group completions. The same construction with pro-$p$ completions of the surface groups still has the full profinite completions of the discrete automorphism groups in its upper row. These assertions specify the precise consequence of the discrete filling theorem; no isomorphism assertion for the lower horizontal or vertical arrows is used.
\end{claim}
\passage{18; 93--94; 186--189; comparison with profinite PartA(c)}

\begin{proof}
The presentation gives $\Pi'/\overline{\langle\!\langle L_b\rangle\!\rangle}=\Pi$. By A.7, $N_{\Pi'}(L_b)=L_b$. Normalize an outer representative to preserve $L_b$, then pass to this quotient. Two normalized representatives differ by conjugation by an element of $L_b$, which dies in the quotient. This proves that $d_b^{\mathrm{cont}}$ is a well-defined homomorphism. It is continuous: the closed subgroup of $\Aut(\Pi')$ preserving $L_b$ is compact, its quotient onto the indicated outer group has kernel the inner automorphisms by $L_b$, and the continuous quotient action factors through this quotient. The same proof works for the pro-$p$ groups using A.7.

Every discrete automorphism extends uniquely to an automorphism of its completion; extensions of inner automorphisms remain inner. Since the resulting target automorphism groups are profinite, the extension maps themselves extend to the full profinite completions of the discrete automorphism groups. Normalizing a discrete representative and then completing gives the same quotient automorphism as completing and then normalizing. Thus the square commutes on the dense discrete subgroup and therefore everywhere. A.12 implies that its upper arrow is an isomorphism. The quotient maps $\widehat\pi\twoheadrightarrow\pi^{(p)}$ and $\widehat{\pi'}\twoheadrightarrow(\pi')^{(p)}$ have characteristic kernels and commute with filling, giving the further comparison from the full square to its pro-$p$ version.

For odd $p$, replacing the upper groups by their pro-$p$ completions would in general fail to define these same maps: an orientation-reversing involution is killed by every homomorphism to a pro-$p$ group, whereas its action on the pro-$p$ orientation line is multiplication by $-1\ne1$. This is an explicit obstruction to that proposed factorization, rather than an identification of two differently typed completion operations.
\end{proof}

\section{Finite actions, sections, and rotation data}
\label{sec:partB}

In this section $U=X\setminus S$ is an orientable connected surface of finite type,
unless another hypothesis is explicitly stated. A universal covering
$p:\widetilde U\to U$ is fixed, and $\pi=\operatorname{Deck}(p)$ acts on the left.
For a faithful action of a finite group $G$ on $U$, we write $E$ for the group of
all lifts of its elements. Thus $1\to\pi\to E\to G\to1$ is an extension of
\emph{discrete} groups. A partial section means a finite subgroup of $E$ projecting
isomorphically onto a subgroup of $G$; its equivalence relation throughout this
section is conjugacy by $\pi$. We use $G^+$ for the orientation-preserving subgroup.
The claims below also specify the corrections required when $G^+=1$.

\subsection{The plane lemma and the section correspondence}

\begin{claim}{B.1}{Proved}{Finite plane actions and simultaneous local models}
Every faithful action of a finite group $H$ on $\R^2$ by homeomorphisms is
conjugate to an orthogonal action. When $H^+\ne1$, $H^+$ is cyclic, $H$ is cyclic
or dihedral, and $(\R^2)^H$ is a single point. A nontrivial orientation-reversing
involution has a fixed line. Every stabilizer in a finite action on a surface
with universal cover $\R^2$ has an invariant disc on which the whole stabilizer
acts by an orthogonal action. These conclusions supply the two local assertions
used in \S15, folios 70--74, scans 141--149, and \S33, scan 440.
\end{claim}
\passage{15; 70-74; 141-149; 33; 222; 440}

\begin{proof}
We use precisely Kolev's Theorem 1.2: a compact subgroup of
$\operatorname{Homeo}(S^2)$ is conjugate to a closed subgroup of $O(3)$
\cite{bKolev2006}. Each homeomorphism of $\R^2$ extends to its one-point
compactification, and the finite group fixes the added point $\infty$.
After the stated conjugacy the group fixes a point $q$ of the round sphere.
Its action on $q^\perp$ gives an injective homomorphism into $O(2)$.
Stereographic projection from $q$ is equivariant and identifies the action on
$S^2\setminus\{q\}$ with that linear action. Finite subgroups of $SO(2)$ are the
cyclic groups of roots of unity: an element of least positive angle generates
the group, by division with remainder of angles. An orientation-reversing
orthogonal matrix in dimension two is a reflection and conjugates every rotation
to its inverse. This proves the cyclic/dihedral assertion and all assertions about
fixed sets.

For the local statement, let $H$ stabilize $x\in U$ and choose
$\widetilde x\in p^{-1}(x)$. For each $h\in H$ there is a unique lift fixing
$\widetilde x$. Uniqueness makes the assignment a homomorphism, and it is
injective because its projection to $U$ is the original action. The preceding
plane argument linearizes these lifts simultaneously. Pull back a sufficiently
small round disc about the common fixed point; choose it inside an evenly
covered neighbourhood of $\widetilde x$. Its projection is the required
$H$-invariant disc. The same assertion for a sphere follows directly from
Kolev's theorem. In particular the fixed locus of an anti-involution is locally
a diameter, with the involution interchanging its two sides. No differentiability
of the original action has entered the argument. The single-element versions
are Constantin--Kolev, Lemma 2.5, Theorems 3.1 and 4.1, and Corollary 4.2
\cite{bConstantinKolev1994}.
\end{proof}

\begin{claim}{B.2}{Refuted}{The omitted orientation hypothesis in scans 146--147}
The displayed lemma on scan 147, taken without its preceding orientation
hypothesis, is false. Its correct form is: an orientation-preserving finite-order
homeomorphism $u$ of $\C$ commuting with $T(z)=z+1$ and satisfying $u(0)=0$ is
the identity. The quotient $\C/\Z$ provides a complete proof along the route
proposed in the manuscript.
\end{claim}
\passage{15; 72-73; 146-147}

\begin{proof}
Complex conjugation is a counterexample to the unqualified sentence: it has
order two, fixes zero, and commutes with $T$. For the corrected statement,
$u(k)=k$ for every integer $k$. B.1 gives at most one fixed point for a
nonidentity orientation-preserving periodic plane homeomorphism, so $u=1$.

Here is the quotient argument. Commutation with $T$ gives a finite-order
homeomorphism $\bar u$ of the annulus $A=\C/\Z$, inducing $+1$ on
$\pi_1(A)=\Z$. It preserves the orientation of $A$. The two end loops, oriented
as boundary loops of a compact subannulus, represent opposite homology classes.
Consequently an orientation-preserving map inducing $+1$ on $H_1(A;\Z)$ fixes
each end. Filling the ends extends $\bar u$ to an orientation-preserving periodic
homeomorphism of $S^2$ fixing both added points and $[0]$. By
Constantin--Kolev, Theorem 4.1, a nonidentity such map has exactly two fixed
points. Hence $\bar u=1$. Its lift $u$ is a deck transformation $T^k$; $u(0)=0$
forces $k=0$. On scan 146 the equation $G=G^\circ$ is exactly what makes the
lifted action orientation preserving. The mathematical repair therefore restores
an inherited hypothesis; it does not justify a spectral-sequence argument.
\end{proof}

\begin{claim}{B.3}{Proved}{Sections and isolated fixed points}
Let $\widetilde U\cong\R^2$. A faithful finite action on $U$ with $G^+\ne1$
has a bijection
\[
 U^G\longrightarrow \operatorname{Sec}(E\to G)/\pi.
\]
For every section subgroup $K$, $C_\pi(K)=1$. Existence of a section forces $G$
to be cyclic or dihedral. More generally, every nontrivial finite subgroup of
$E$ whose orientation-preserving part is nontrivial is contained in exactly one
maximal finite subgroup of $E$. For an orientation-preserving action the map
from points with nontrivial stabilizer to maximal partial-section classes is a
$G$-equivariant bijection. These are the theorem and corollary of \S15, folios
73--75, scans 148--151, with their stated orientation hypotheses.
\end{claim}
\passage{15; 73-75; 148-151}

\begin{proof}
For $x\in U^G$, the lifts of $G$ fixing a chosen point over $x$ form a section.
Changing that point conjugates the section by the corresponding element of
$\pi$. Conversely a section $K$ is a finite faithful group of plane
homeomorphisms. B.1 supplies a unique fixed point $\widetilde x$. Its image is
fixed by $G$, and the section fixing $\widetilde x$ is $K$. These constructions
are inverse: conjugate sections have deck-translated fixed points and hence
the same image in $U$. An element of $C_\pi(K)$ preserves the unique fixed point.
A deck transformation with a fixed point is the identity. This proves the
centralizer assertion. B.1 gives the possible finite groups.

Let $K\le E$ have nontrivial positive part and fix $\widetilde x$. The stabilizer
$E_{\widetilde x}$ is finite, since it intersects $\pi$ trivially and injects into
the finite group $G$. Any finite subgroup containing $K$ has a fixed point by
B.1; that point must be $\widetilde x$. It is therefore contained in
$E_{\widetilde x}$. This proves uniqueness as an actual subgroup, before any
conjugacy quotient. In the orientation-preserving case every nontrivial partial
section has this property. Applying the same construction to each stabilizer
$G_x$ gives the final bijection and its equivariance. The argument also applies
to a nonorientable $U$ whenever the finite lifted group has nontrivial positive
part; in particular the cardinality condition $|G|\ge3$ on scan 149 ensures it.
\end{proof}

\begin{claim}{B.4}{Proved}{Reflection sections, ovals, and punctured ovals}
Let $\sigma$ be an orientation-reversing involution of $U$ with
$\widetilde U\cong\R^2$, and let $E\to C_2$ be its lift extension. There is a
bijection
\[
 \pi_0(U^\sigma)\longrightarrow\operatorname{Sec}(E\to C_2)/\pi.
\]
For the section $s$ associated with a component $Z$, inclusion induces an
isomorphism $\pi_1(Z)\cong C_\pi(s)$. Thus the centralizer is $1$ for a line and
$\Z$ for an oval. This includes punctures fixed or permuted by $\sigma$ and proves
\S33, folios 222--225, scans 440--445, in the required noncompact cases.
\end{claim}
\passage{33; 222-225; 440-445}

\begin{proof}
Each order-two lift $s$ is a reflection of the plane by B.1. Its fixed set
$L$ is a line. The restriction $p:L\to U^\sigma$ is locally a homeomorphism
onto the real locus, by the local model in B.1. Its image is the entire component
$Z$ containing any of its points. To verify this without a properness assumption,
lift a path in $Z$ starting at $p(\widetilde x)$: applying $s$ gives another lift
with the same initial point, so uniqueness of path lifting fixes the lifted path.
The same argument for a homotopy shows that $L\to Z$ is its universal covering.
Every section has such a fixed line, so every section class arises. Two fixed
lines projecting to the same component contain lifts of the same point; a deck
transformation taking one lift to the other conjugates the associated involutions.
The two lifts fixing the resulting point agree, because lifts of a specified
base homeomorphism are unique once the value at one point is fixed.

A deck transformation preserving $L$ normalizes $\langle s\rangle$ and hence
centralizes $s$. Conversely a deck transformation centralizing $s$ preserves
$L$. Therefore the deck group of $L\to Z$ is exactly $C_\pi(s)$. Since $L\cong\R$,
this proves both the asserted isomorphism and injectivity of $\pi_1(Z)\to\pi$.
A finite-type real locus is a finite union of circles and open intervals after
its real punctures are removed. Each interval is homeomorphic to $\R$.
This argument is also the explicit reduction of Mochizuki's Theorem 3.6 and
Corollary 3.7 \cite{bMochizuki2003} to Grothendieck's formulation: his full
centralizer is $\pi_1(Z)\times C_2$, and intersecting it with $\pi$ gives the
centralizer used here.
\end{proof}

For $G=1$ the section set has one element and its centralizer is $\pi$; the
comparison is the bijection $\pi_0(U)\to\{*\}$. The map from individual points
is not injective. The hypothesis on the universal cover has content: on $S^2$
a nontrivial rotation has two fixed points while its lift extension, with kernel
$1$, has one section; the antipodal involution has no fixed point and still has
one section. These examples explain exactly why the same formula does not extend
to that excluded surface.

\begin{claim}{B.5}{Proved}{Recovering the incidence and cyclic order of real punctures}
For a hyperbolic $U=X\setminus S$ with anti-involution $\sigma$, its discrete
extension together with its indexed peripheral subgroups determines the following
finite graph. Its vertices are the section classes of B.4, its edges are the real
punctures, and each edge is attached to its two adjacent real half-branches.
Every punctured oval gives one cycle, allowing a loop or two parallel edges.
The graph recovers which punctures lie on each oval and their cyclic order up
to reversal. Unpunctured ovals are the vertices with infinite cyclic centralizer.
The claim uses the extension and peripheral incidence maps; a naked set of
section classes alone has no such incidence information.
\end{claim}
\passage{33; 221-225; 437,440-445}

\begin{proof}
A puncture is real exactly when its $C_2$-stabilizer is all of $C_2$. For such a
puncture its peripheral extension is
\[
 E_i=\langle t,s\mid s^2=1,;sts^{-1}=t^{-1}\rangle,
 \qquad L_i=\langle t\rangle.
\]
Indeed the local anti-involution of the filled disc is reflection. On the
logarithmic covering $z=\exp(2\pi i w)$ of the punctured disc its two types of
order-two lifts are $w\mapsto-\overline w$ and
$w\mapsto1-\overline w$. Algebraically the reflections $t^k s$ have two
$L_i$-conjugacy classes, distinguished by the parity of $k$, since conjugation
by $t^m$ adds $2m$. The two inclusions of these classes into the section set of
$E$ are precisely the two real half-branches. They can have the same image.
Attach the edge $i$ using these two maps. Walking through the intervals of a
punctured circle gives a cycle whose edge labels are its punctures. Conversely
these cycles determine the partition of the punctures and their circular order,
with the orientation of each circle unchosen. B.4 identifies every remaining
vertex with an unpunctured oval. A nonreal pair of punctures has trivial
stabilizer and supplies no edge. Finally any two configurations on circles with
the same circular order are related by a homeomorphism: map successive closed
arcs monotonically, and extend over collars. Thus this recovery specifies the
topological position of the punctures, including the single-puncture and
single-interval case, without specifying their analytic positions.
\end{proof}

\subsection{Finite automorphisms and the withdrawn assertions}

\begin{claim}{B.6}{Proved}{The fundamental loop lemma of \S18}
Let $\pi=\pi_1(U)$ be hyperbolic and discrete, with its indexed peripheral
structure and orientation module. A finite-order orientation-preserving
automorphism of $\pi$ fixing a peripheral index and normalizing its subgroup
$L_i$ is the identity. Consequently the peripheral extension attached to a
nontrivial finite-order positive outer automorphism fixing $i$ cannot split.
This proves the lemma and corollary of scans 193--194 without adding realizability
to the hypotheses.
\end{claim}
\passage{18; 96-98; 193-196}

\begin{proof}
The based Dehn--Nielsen--Baer correspondence identifies the given finite cyclic
group with a finite subgroup of the mapping class group with one additional
marked base point. Kerckhoff's realization theorem realizes this group by
homeomorphisms fixing that point and preserving the marked punctures. The
original paper explicitly treats cusps in \S IV, pp.~255--256, in addition to
the orientation-reversing convention on p.~235 \cite{C_Kerckhoff1983}; hence the
punctured case used here is included. The section determined by the fixed base
point identifies the induced genuine action on $\pi$ with the given one.

An orientation-preserving automorphism fixing $i$ acts as $+1$ on $L_i\cong\Z$,
so it fixes a generator $\ell$ of $L_i$. A nontrivial finite positive cyclic group
of plane lifts has one fixed point by B.1. Commutation with the nontrivial deck
transformation $\ell$ would give the second fixed point $\ell\widetilde x$,
a contradiction. Thus the automorphism is the identity. A splitting of the
peripheral extension gives a finite-order lift normalizing $L_i$; the result
just proved makes its outer class trivial.

The margin on scan 194 writes the \emph{filled} type as $(g,\nu)$ and the
acting group as type $(g,\nu+1)$. For $(g,\nu)=(0,2)$ or $(1,0)$ the acting
group is still a hyperbolic free group of rank two, so the same proof applies.
For a literal annulus, a positive automorphism fixing one of its two indexed
ends acts trivially on $\pi=\Z$. Forgetting the indices permits the positive
end interchange, which induces $-1$ on $\Z$ while normalizing its common
peripheral subgroup. For a literal unpunctured torus there is no peripheral
index, so the displayed loop lemma has no instance. These are the precise
comparisons required by the marginal remark, also reflected in A.13.
\end{proof}

\begin{claim}{B.7}{Refuted}{The unrestricted fixed-subgroup question, with its replacement}
The assertion $\pi^G=1$ for every finite $G\le\Aut_{\rm lac}(\pi)$ is false.
For a hyperbolic discrete $\pi$, a nontrivial finite genuine loop action has
fixed subgroup $1$ when its positive part is nontrivial. Its remaining possibility
is a reflection action of $C_2$; its fixed subgroup is $1$ or $\Z$, as in B.4.
For any finite-type faithful geometric action there are finitely many classes
of full and partial sections. This answers questions (c) and (e), scans 197--198.
\end{claim}
\passage{18; 98-99; 197-198}

\begin{proof}
A finite subgroup of $\Aut_{\rm lac}(\pi)$ intersects $\Inn(\pi)$ trivially:
$\pi$ is torsion-free and centre-free. Its outer image is therefore finite and
is realized by Kerckhoff. The canonical centre-free extension identifies the
given genuine action with a section of its geometric lift extension. B.3 and
B.4 now give the complete list of fixed subgroups. An explicit counterexample
is the double of an orientable one-holed torus, a genus-two surface, with the
reflection interchanging the two halves. Its fixed oval is essential and a
base point on it gives a genuine order-two automorphism with fixed subgroup
$\Z$.

For finiteness, compactify $U$ by its finitely many ends. Every nonidentity
positive periodic homeomorphism has isolated fixed points by B.1. The resulting
closed discrete subset of the compact surface is finite. A reflection has a
compact one-dimensional fixed submanifold, hence finitely many components;
removing finitely many punctures preserves finiteness. Apply B.3 or B.4 to each
of the finitely many subgroups of $G$. This proves the asserted finiteness.
For the abelian groups the genuine actions can be read directly: a nonidentity
positive finite-order element of $SL_2(\Z)$ fixes no nonzero lattice vector,
whereas a determinant $-1$ involution fixes a primitive rank-one sublattice.
For $\Z$, the nonidentity automorphism is $-1$ and its fixed subgroup is zero.
The trivial group of operators fixes all of $\pi$ in every type.
\end{proof}

\begin{claim}{B.8}{Refuted}{Unique maximal sections fail for reflections, already in type $(0,3)$}
The orientation-free extension of question (d), scan 198, and of the expectation
on scan 155 is false. Put
\[
 E=\langle s_0,s_1,s_2\mid s_j^2=1, (s_0s_1)^2=1, (s_1s_2)^3=1\rangle.
\]
There is a geometric extension $1\to\pi_{0,3}\to E\to S_3\times C_2\to1$.
The reflection subgroup $\langle s_1\rangle$ is contained in two distinct maximal
finite subgroups $\langle s_0,s_1\rangle\cong C_2^2$ and
$\langle s_1,s_2\rangle\cong S_3$, which are not $\pi_{0,3}$-conjugate.
The valid uniqueness statement is B.3, requiring a nontrivial positive part.
\end{claim}
\passage{15; 77; 155; 18; 99-101; 198,201-203}

\begin{proof}
Take a hyperbolic triangle with angles $\pi/2,\pi/3,0$ and reflect in its three
sides. Reflected triangles tile the plane; the local tilings at the two finite
vertices give exactly the displayed relations. Equivalently its group is the
amalgam $C_2^2*_{C_2}S_3$, with common subgroup generated by $s_1$.
Write $S_3=\langle r,a\mid r^3=a^2=1, ara=r^{-1}\rangle$, and let $\tau$ generate
the central $C_2$. Define the surjection
\[
 s_0\mapsto\tau,\qquad s_1\mapsto a\tau,\qquad s_2\mapsto ra\tau.
\]
It is injective on both vertex groups. A finite subgroup of this amalgam fixes
a vertex of its Bass--Serre tree: the finite orbit spans a finite tree, whose
centre is fixed (there are no edge inversions). Thus the kernel is torsion-free.
The character sending each $s_j$ to $-1$ factors through the exponent of $\tau$,
so the kernel preserves orientation. The group Euler characteristic is
$1/4+1/6-1/2=-1/12$, and the index is $12$; hence the kernel has Euler
characteristic $-1$. The stabilizer of the cusp is $\langle s_0,s_2\rangle$;
its image has order four. The kernel therefore has $12/4=3$ cusps, and its
quotient surface has genus zero. This constructs the asserted geometric extension.

Each vertex group is maximal finite, again by the tree action: a finite overgroup
would fix a vertex, and a vertex group of order greater than the edge group fixes
only its vertex. Their intersection is $\langle s_1\rangle$ by the normal form
for the amalgam. Their orders differ, and their generated group is all of the
infinite amalgam. This proves every claimed failure. The finite quotient, its
vertex and cusp images, and the Euler-characteristic calculation are checked in
\texttt{code/check\_part\_b.py}. Grothendieck's proof on scans 201--203 applies
when the finite overgroups are positive cyclic groups: they centralize their
intersection, and B.3 makes the kernel of that centralizer map trivial. A
reflection has a line of fixed points, so that argument supplies neither of
these requirements in the generalization just refuted.
\end{proof}

\begin{claim}{B.9}{Proved}{The valid normalizer assertion}
For a hyperbolic discrete group with loops, a nontrivial finite subgroup
$K\le\Aut^+_{\rm lac}(\pi)$ is cyclic and
\[
 N_{\Aut^+_{\rm lac}(\pi)}(K)=C_{\Aut^+_{\rm lac}(\pi)}(K).
\]
An orientation-reversing element of its full normalizer acts on $K$ by inversion.
Every finite positive overgroup of $K$ is cyclic. These are the valid assertions
within the passage labelled ``Ultra fort!'' on scans 198--199.
\end{claim}
\passage{18; 99-100; 198-200}

\begin{proof}
Realize the outer image of $K$ and use its given section. B.3 makes $K$ cyclic.
The boundary rotation character of B.12 is injective on $K$ and is natural for
loop automorphisms with their action on the orientation module. A positive
normalizer therefore preserves every value of that character, so fixes every
element of $K$. A negative normalizer changes its sign, so inverts every element.
The same realization argument proves cyclicity of each finite positive overgroup.
This proof compares a normalizing automorphism on the canonical oriented boundary;
it makes no assertion that an infinite normalizer acts as a finite isometry group.
\end{proof}

\begin{claim}{B.10}{Refuted}{The withdrawn uniqueness of topological roots}
The lemma marked ``faux!'' on scan 200 is false even for the punctured torus:
two different finite-order homeomorphisms fixing the base point can have the
same nontrivial square and induce the same automorphism of $\pi_1$.
\end{claim}
\passage{18; 100; 200}

\begin{proof}
On $U=\R^2/\Z^2\setminus\{0\}$ choose $p=(1/2,1/2)$ and let
$r(x,y)=(-y,x)$, $h=r^2$. Choose a small disc $D$ whose four $r$-translates are
disjoint and avoid both $0$ and $p$. Choose a nonidentity homeomorphism $f_0$
supported in $D$ and isotopic to the identity there relative to its boundary;
for example use a radial bump rotation in disc coordinates. Put
$f=f_0(hf_0h^{-1})$. Its two disjoint supports show that $fh=hf$; the same
supports show that $fr\ne rf$, since $r$ moves their union to the disjoint pair
$rD,r^3D$. Hence $v=frf^{-1}\ne r$, but $v^2=h$ and $v(p)=p$.
The isotopy of $f$ fixes $0$ and $p$, so $f_*=1$ on $\pi_1(U,p)$ and therefore
$v_*=r_*$. This verifies the exact data of the withdrawn statement. What survives
is uniqueness up to conjugacy by the identity component in the realization
comparison, not equality of the homeomorphisms.
\end{proof}

\begin{claim}{B.11}{Refuted}{Torsion in the centralizer need not be a subgroup}
The assertion following the normalizer assertion on scan 199, that the torsion
elements of the centralizer form one largest finite cyclic overgroup, is false
for the hyperbolic type $(1,1)$.
\end{claim}
\passage{18; 99-100; 198-200}

\begin{proof}
Keep $U,p,r,h$ as in B.10. The essential simple curve
$c=\{x+y=1/2\pmod1\}$ avoids $0,p$ and is invariant under $h$.
The based mapping class $T_c$ commutes with $h$, since conjugation of a positive
Dehn twist by an orientation-preserving map is the positive twist about the
image curve. Thus the two genuine automorphisms induced by $r$ and
$T_crT_c^{-1}$ have order four and common square $h_*$; both centralize $h_*$.
On $H_1(U;\Z)$, with a consistent choice of twist sign, their matrices are
\[
 S=\begin{pmatrix}0&-1\\1&0\end{pmatrix},\qquad
 TST^{-1}=\begin{pmatrix}-2&-5\\1&2\end{pmatrix},\qquad
 T=\begin{pmatrix}2&1\\-1&0\end{pmatrix}.
\]
Both square to $-I$, while their product has trace $-6$ and determinant one.
Its eigenvalues are $-3\pm\sqrt8$, so it has infinite order. The corresponding
product of automorphisms has infinite order as well. These integral equalities
are checked by the supplied program. Consequently the centralizer's torsion
set is not closed under multiplication, and there is no finite subgroup
containing all its finite overgroups of $\langle h_*\rangle$. This coexists with
the exact equality of normalizer and centralizer in B.9.
\end{proof}

\subsection{The multiplier and the abelian surfaces}

\begin{claim}{B.12}{Proved}{An intrinsic discrete formula for the multiplier}
For a geometric finite extension of a hyperbolic discrete surface group,
$E$ with its peripheral structure and orientation module $T$ determines
$\mu_E$ on all torsion elements. On $E^+$ it is the rotation number of its
canonical oriented boundary action. More explicitly, let
$\widetilde{\operatorname{Homeo}}^+(S^1)\to\operatorname{Homeo}^+(S^1)$ be the
central extension by the positively oriented deck translation $z$. For
$u\in E^+$ of order $n$ and any lift $\widetilde u$, write
$\widetilde u^{,n}=z^k$. Then
\[
 \mu_E(u)=\frac{k}{n}\pmod{\Z}\quad\hbox{in }T\otimes\Q/\Z.
\]
For a nonidentity orientation-reversing torsion element we put
$\mu_E(u)$ equal to the unique nonzero element of order two. With this
convention the formula is natural under isomorphisms preserving the loops and
transporting $T$, satisfies $\mu(u^m)=m\mu(u)$, and has the same order as $u$.
This answers the discrete intrinsicity question of \S15, folios 77--79,
scans 155--160.
\end{claim}
\passage{15; 77-79; 155-160}

\begin{proof}
Choose an oriented finite-area hyperbolic uniformization of $U$. Its boundary
at infinity is a circle. We explain the ambiguity in this choice. A marking
isomorphism of surface groups preserving their peripheral subgroups and the
orientation is induced by an orientation-preserving homeomorphism, by the
Dehn--Nielsen--Baer correspondence. We can choose that marking quasiconformal:
on a compact core smooth it, and on each cusp use standard cylindrical
coordinates and extend with bounded dilatation. Its lift extends to an oriented
boundary homeomorphism. This extension is uniquely determined by its
$\pi$-equivariance. Indeed it must carry the attracting and repelling endpoints
of every hyperbolic element to those of the corresponding element. These
endpoints are dense in the boundary circle of a finite-area surface group.
For completeness, density follows from its full limit set: a boundary arc
meets translates of a fundamental ideal polygon arbitrarily near infinity;
conjugating either of two hyperbolic elements with disjoint endpoints by the
corresponding translates gives endpoints in that arc. Thus two equivariant
boundary maps agree on a dense set and hence everywhere.

For each automorphism of $\pi$ preserving the peripheral structure,
Dehn--Nielsen--Baer provides the boundary map conjugating its deck action.
The preceding uniqueness makes composition exact and makes this circle,
together with its $E$-action, independent of the uniformization up to a unique
$\pi$-equivariant boundary isomorphism. The boundary orientation is recorded by
$T$, so changing orientation acts by $-1$. Realize the finite action and choose
an invariant conformal structure, using B.1's cyclic and dihedral local charts
on the quotient. Its hyperbolic metric makes a finite-order positive lift an
elliptic isometry. In the disc model centred at its fixed point it is the same
rotation on the tangent circle and on the circle at infinity. The displayed
formula therefore gives its local multiplier.

Replacing $\widetilde u$ by $z^j\widetilde u$ replaces $k$ by $k+nj$, so the
value is independent of the lift. An orientation-preserving finite-order circle
homeomorphism of order $n$ is conjugate to a rotation by $k/n$ with
$\gcd(k,n)=1$; this also follows by ordering the finite orbit on the circle
and extending a homeomorphism of one fundamental interval around that orbit.
This proves the order and power statements. A negative finite plane isometry
is a reflection, of order two, so its assigned value has exactly the required
power property and is unchanged by sign reversal.

The normalizer formula suggested in the margin of scan 159 is recovered as
follows. Drill the orbit of the fixed point of $u$. Above the new puncture the
positive peripheral normalizer is infinite cyclic, generated by a positively
oriented element $v$, with peripheral subgroup generated by $v^n$. A lift of
$u$ to this normalizer is $v^a$ modulo $\langle v^n\rangle$, and
$\mu(u)=a/n$. In a local coordinate this is the equality between the rotation
$z\mapsto e^{2\pi ia/n}z$ and monodromy in $z\mapsto z^n$. Harvey's Theorem 7
\cite{bHarvey1971} computes exactly this inverse-residue relation from the
elliptic generators of the orbifold epimorphism. His convention uses an inverse
deck transformation, accounting for the sign visible in his proof on
pp.~393--394. Our local-coordinate calculation also applies to cusps and to the
positive cyclic part of a dihedral stabilizer; B.1 supplies the reflection part.
\end{proof}

\begin{claim}{B.13}{Refuted}{Forgetting loops can change the multiplier}
An abstract extension, with its orientation line regarded as an additional
rank-one module, need not determine the multiplier after the peripheral
structure is forgotten. There are isomorphic extensions with kernel a free
group of rank two and quotient $C_3$ whose geometric multipliers disagree,
even after allowing either isomorphism between their orientation lines.
\end{claim}
\passage{15; 77-79; 155-160}

\begin{proof}
Let $E_0=\langle u,v\mid u^3=v^3=1\rangle=C_3*C_3$. Realize it as the
positive orbifold group of a sphere with one cusp and two cone points of order
three; choose $u,v$ as its positive local generators. The epimorphisms
\[
 q_+(u)=q_+(v)=1,\qquad q_-(u)=1,\quad q_-(v)=-1
 \quad (\text{values in }\Z/3)
\]
are injective on every finite subgroup, because finite subgroups of a free
product are conjugate into its factors (the tree argument in B.8 applies).
Their kernels are torsion-free of index three and Euler characteristic $-1$.
The cusp generator is $(uv)^{-1}$. Its image under $q_+$ is nonzero, so there
is one cusp and the kernel is a once-punctured torus group. Its image under
$q_-$ is zero, so there are three cusps and the kernel is a pair-of-pants group.
Both have free rank two.

The automorphism $\psi(u)=u,\ \psi(v)=v^{-1}$ satisfies $q_-\psi=q_+$ and
is an isomorphism of the two abstract extensions. The positive local generator
values in each realization are $\mu(u)=\mu(v)=1/3$. Thus $\psi$ preserves the
value at $u$ but changes the value at $v$ to $-1/3$. Reversing the isomorphism
of orientation lines would instead break the equality at $u$. No orientation-line
identification repairs both. The isomorphism sends the cusp word $uv$ to
$uv^{-1}$, displaying precisely which peripheral comparison has been discarded.
The finite monodromy calculation is included in the computational record.
\end{proof}

\begin{claim}{B.14}{Proved}{Finite annulus and torus actions, and their multipliers}
Every finite topological action on the torus is topologically conjugate to an
affine action on $V/\Lambda$, where $\Lambda\cong\Z^2$ and
$V=\Lambda\otimes\R$. For a fixed finite operator group $G$ such an action is
specified by a homomorphism $A:G\to GL(\Lambda)$ and a cocycle
$b\in Z^1(G,V/\Lambda)$, subject to faithfulness of
$g\mapsto(x\mapsto A_gx+b_g)$. Translation conjugacy corresponds exactly to a
coboundary, and
\[
 H^1(G,V/\Lambda)\xrightarrow{\sim}H^2(G,\Lambda)
\]
is the connecting homomorphism. Linear changes of the lattice give the further
conjugacies. Every finite annulus action has a simultaneous normal form
\[
 (x,t)\longmapsto(\epsilon_g x+a_g,\eta_g t),
 \qquad x\in\R/\Z,\quad t\in\R,\quad
 \epsilon_g,\eta_g\in\{1,-1\}.
\]
These descriptions determine the multiplier of every finite-order element in
the respective lift extensions.
\end{claim}
\passage{15; 78-79; 157-159; 17; 89-90; 179-181}

\begin{proof}
We first justify simultaneous conformalization. B.1 gives genuine cyclic or
dihedral charts for every finite stabilizer. The quotient is a compact
surface orbifold after the ends are added; its interior special points have
cyclic charts, and its mirror edges and corners have reflection and dihedral
charts. Choose a smooth metric on the quotient away from these charts, choose
round cone and mirror metrics in them, and patch in collars using an invariant
partition of unity on their finite local covers. The local charts make the
patched lifts smooth. Pullback is a smooth invariant conformal structure on
the original surface. This construction proves the required passage from the
original topological action; it does not average a merely topological metric
and call the result Riemannian.

Uniformization of the torus gives $\C/\Lambda$. Every holomorphic automorphism
lifts to an entire map $f$ with $f(z+\lambda)-f(z)\in\Lambda$. Its derivative
is periodic and hence bounded on $\C$, so it is constant; $f(z)=az+c$.
An antiholomorphic automorphism is obtained by composing such a map with
complex conjugation on its source. Hence every element is real affine on the
same lattice. The action law is precisely
$b_{gh}=b_g+A_gb_h$. Averaging cochains over the finite group gives
$H^q(G,V)=0$ for $q>0$. The exact sequence
$0\to\Lambda\to V\to V/\Lambda\to0$ therefore gives the asserted isomorphism.
In concrete terms, lifts $v_g\in V$ of $b_g$ give the extension cocycle
$v_g+A_gv_h-v_{gh}\in\Lambda$. Conjugating by translation adds a coboundary.
Moreover any topological conjugacy between two affine actions can be replaced
by an affine conjugacy. A lift of the conjugacy has form $Bx+f(x)$, with
$B\in GL_2(\Z)$ and $f$ periodic. Average its equivariance equation over the
torus with Haar measure. Affine lattice automorphisms preserve that measure,
so $Bx+\int f$ satisfies the same equation modulo $\Lambda$.

For the annulus, extend the action to its two-point sphere compactification
and apply Kolev. In the resulting conformal sphere choose a M\"obius coordinate
sending the two ends to zero and infinity. The maps are then of the forms
$z\mapsto az$, $a/z$, $a\overline z$, or $a/\overline z$. On the logarithmic
radial coordinate they act affinely with slopes $\eta=\pm1$. A finite group
of affine transformations of the real line fixes the average of any orbit;
translation to this common fixed point removes every radial constant. The
angular coordinate gives the claimed $\epsilon x+a$. This proves the displayed
simultaneous normal form. The action on $\pi_1$ is $\epsilon_g$ and its
orientation sign is $\epsilon_g\eta_g$.

For the promised computation on the torus, an element of the lift extension
is $x\mapsto Ax+v$. It has finite order exactly when $A$ has finite order,
say $n$, and $(1+A+\cdots+A^{n-1})v=0$. Averaging its orbit gives a fixed point,
and translation to that point conjugates it to $A$. A nonidentity positive
finite-order matrix has trace $-2,-1,0,1$, giving orders $2,3,4,6$, respectively.
For order two its multiplier is $1/2$. For the other orders it is $+1/n$ or
$-1/n$ according to the sign of $\det(w,Aw)$ for any nonzero $w$; this sign
is independent of $w$, as is seen after averaging an inner product and writing
$A$ as a rotation. A negative finite-order matrix is a reflection and receives
$1/2$. A lift with $A=I$ has finite order only when it is the identity.
For the annulus a lift is $(x,t)\mapsto(\epsilon x+a+k,\eta t)$. With
$\epsilon=-1$ its square is the identity; with $\epsilon=1$ finite order
requires $a+k=0$, leaving the identity or $t\mapsto-t$. Thus every nonidentity
torsion lift on an annulus has order two and multiplier $1/2$.
\end{proof}

\begin{claim}{B.15}{Refuted}{The ``only counterexample'' of scan 179}
The assertion that a faithful homotopically trivial finite annulus action
failing to be a translation action must be the single anti-involution
$z\mapsto1/\overline z$ is false. For the two non-simply-connected abelian
types, the exact replacement is: homotopically trivial torus actions are
translation actions; homotopically trivial annulus actions are the finite
subgroups of $SO(2)\times C_2$ in the normal form of B.14. Those contained in
$SO(2)$ are translations, while those with nontrivial second projection include
two different types of anti-involution and groups of arbitrarily large order.
\end{claim}
\passage{17; 89-90; 179-181}

\begin{proof}
On the torus, homotopic triviality is $A_g=I$ for every $g$, so B.14 proves the
translation statement and the implication doubted on scans 180--181.
On the annulus it is $\epsilon_g=1$, which gives the asserted subgroup.
For each $n\ge1$ the maps
\[
 r(z)=e^{2\pi i/n}z,\qquad s(z)=1/\overline z
\]
generate a faithful $C_n\times C_2$ action, and every element induces the
identity on $\pi_1(\C^*)$. The fixed circle of $s$ excludes a translation
realization: a nonidentity translation in any group is fixed-point free.
In addition $s_-(z)=-1/\overline z$ is an involution inducing $+1$ on $\pi_1$,
but it has no fixed point, since its fixed-point equation is $|z|^2=-1$.
It cannot be conjugate to $s$, and it cannot be a translation of a connected
topological group: every such translation is isotopic to the identity along
a path in that group and hence preserves orientation. These formulas give the required full
counterexamples, not merely a gap in the argument.

More explicitly let the rotation subgroup have order $n$ and generator
$(x,t)\mapsto(x+1/n,t)$. A second-coset generator can be chosen as
$(x,t)\mapsto(x+m/(2n),-t)$, with $m=0$ or $1$ after multiplying by a rotation.
Its square is the $m$th rotation. This exhausts the finite subgroups projecting
onto $C_2$, and includes both the reflection and glide cases. For the omitted
simply connected surfaces, finite plane actions are the orthogonal actions of
B.1 and are all homotopically trivial; positive sphere actions are likewise
homotopic to the identity. Nontrivial rotations fixing a point already refute
any extension of the translation formulation to those types.
\end{proof}

\subsection{Anti-involutions and the later tube construction}

\begin{claim}{B.16}{Proved}{Complete topological classification of anti-involutions}
Conjugacy classes of anti-involutions of a closed oriented genus-$g$ surface
are classified by the orientability of the quotient and its number $b$ of
boundary components. The possibilities are
\[
 \begin{array}{ll}
 \text{nonorientable quotient:}&0\le b\le g,\quad
                  h=g+1-b\ge1\text{ crosscaps};\\
 \text{orientable quotient:}&b\ge1,\quad g=2h+b-1,\quad h\ge0.
 \end{array}
\]
There are $(g+1)+\lfloor(g+2)/2\rfloor=\lfloor(3g+4)/2\rfloor$ classes.
For $\nu$ unlabelled invariant punctures the additional data are an integer
$r\ge0$ and an unordered list $(q_1,\ldots,q_b)$ of nonnegative integers with
$2r+\sum q_i=\nu$: $r$ is the number of interior punctures of the quotient,
and $q_i$ counts punctures on its $i$th boundary circle.
\end{claim}
\passage{33; 221-222; 437-440}

\begin{proof}
By B.1 the quotient $Y=X/\sigma$ is a compact surface with boundary
$X^\sigma$. The complement of the real locus covers its interior with degree
two; the Euler characteristic of the boundary is zero. Hence
$\chi(X)=2\chi(Y)$. If $Y$ is nonorientable with $h$ crosscaps,
$\chi(Y)=2-h-b$, yielding $g=h+b-1$; if orientable of genus $h$ it yields
$g=2h+b-1$. In the latter case $b$ must be positive, since the oriented double
of a closed orientable quotient would be disconnected. Conversely form the
orientation double over the interior of each allowed $Y$ and join the two
local sides at its boundary. The result is a connected closed orientable surface,
with the deck interchange extended to reflection along the boundary; its genus
is given by the displayed formula. The quotient recovers $Y$, so the classification
of compact surfaces proves both existence and uniqueness of these conjugacy classes.
This also proves the complete list underlying Gross--Harris, Proposition 3.1
\cite{bGrossHarris1981}, including the two corrections in the margins of scan 439.
Grothendieck's nonorientable parameter $\gamma$ is $h-1$; thus his formula
$g=\gamma+b$ agrees exactly with ours.

An anti-involution of a punctured finite-type surface extends to its end
compactification. Punctures in its quotient lie either in the interior,
with two-point inverse image, or on the boundary, with one-point inverse image.
Any two finite interior configurations of the same cardinality are related by
an isotopy supported away from the boundary. Boundary components may be
permuted, and any two unlabelled configurations of the same cardinality on one
circle are related by a circle homeomorphism, extended over a collar. These
facts prove completeness of the punctured invariants. With individually
labelled punctures one must instead retain their assignments and their circular
orders; B.5 provides the corresponding incidence reconstruction.
\end{proof}

\begin{claim}{B.17}{Proved}{Real moduli components, with the group and marking convention fixed}
Let $\mathcal M_{g,\nu}$ mean the real moduli stack of smooth proper
geometrically connected genus-$g$ curves with an \emph{unlabelled}
degree-$\nu$ reduced puncture divisor. We retain the stack and its
real forms also in the four unstable types. Its complex
mapping class group is $\mathcal T^+_{g,\nu}$, allowing permutations of the
punctures. Write $p_b(N)$ for the number of partitions of $N$ into at most
$b$ parts, with $p_0(0)=1$ and $p_0(N)=0$ for $N>0$. Then
\[
 \#\pi_0\bigl(\mathcal M_{g,\nu}(\R)\bigr)
 =\sum_{b=0}^{g}\ \sum_{r=0}^{\lfloor\nu/2\rfloor}p_b(\nu-2r)
  +\sum_{\substack{1\le b\le g+1\\ b\equiv g+1\ (2)}}
       \ \sum_{r=0}^{\lfloor\nu/2\rfloor}p_b(\nu-2r).
\]
For closed hyperbolic surfaces this is $\lfloor(3g+4)/2\rfloor$.
In the four unstable types $(0,0)$, $(0,1)$, $(0,2)$ and $(1,0)$ the counts are,
respectively, $2,1,3,3$. In the last case no origin is prescribed.
Distinct topological anti-involution classes remain distinct in the extended
mapping class group in these hyperbolic types. The last assertion fails for
$(0,0)$, $(0,2)$, and $(1,0)$.
\end{claim}
\passage{33; 222; 439-440}

\begin{proof}
C.5 and C.6, finite-action realization and uniqueness up to the identity component, identify
conjugacy classes of anti-involutions on the fixed punctured surface with
conjugacy classes of orientation-reversing involutions in
$\mathcal T_{g,\nu}$. One can use positive conjugacy here: an orientation-reversing
conjugator can be multiplied by the target involution to become positive.
Mochizuki, Corollary 3.8, identifies these conjugacy classes with the connected
components of the real moduli \emph{stack} \cite{bMochizuki2003}. His Example 3.3
uses the full modular group of a punctured surface, which is the unlabelled
convention just specified. B.16 supplies the list. For fixed $b,r$, unordered
nonnegative lists of length $b$ summing to $\nu-2r$ are exactly the partitions
counted by $p_b(\nu-2r)$. Summing over the two quotient families proves the
formula. The supplied code enumerates these numbers for $2g+\nu\le6$ and checks
the closed formula for $0\le g\le20$; the proof holds for all $g$.

We verify the unstable component counts without using mapping-class
injectivity. A real genus-zero curve is a split or nonsplit real conic;
analytically these are the sphere with $z\mapsto\overline z$ or with
$z\mapsto-1/\overline z$. To see these are the two possibilities,
conformalize its anti-involution as in B.14 and apply the two quotient
types in B.16; uniformization identifies the resulting conformal
double with the sphere. A real point exists exactly in the first type.
There is no complex moduli parameter in genus zero. With one puncture
the point is real, and $PGL_2(\R)$ acts transitively on the real
projective line, giving one component. With two unlabelled punctures
the possibilities are two real points on the split conic, one
conjugate pair on the split conic, or one conjugate pair on the
nonsplit conic. The first is transitive under $PGL_2(\R)$, the second
under its action on the upper half-plane, and the third under the
rotations commuting with the antipodal map. These give three
components. With no puncture there are the two conic components.

For an unpointed genus-one curve, B.14 and B.16 give the three
anti-involution types with zero, one, or two real ovals. Each type
has a connected parameter space of complex structures. Explicit
representatives are
\[
 \begin{array}{c|c|c}
 b&\Lambda_t&\sigma(z)\\\hline
 0&\Z+it\Z&\overline z+\tfrac12\\
 1&\Z+(\tfrac12+it)\Z&\overline z\\
 2&\Z+it\Z&\overline z
 \end{array}\qquad(t>0).
\]
To obtain this list from an invariant lattice, rotate a primitive
real period to $1$. Conjugation of a second period forces its real
part modulo $\Z$ to be $0$ or $1/2$. The translation cocycle of
B.14 is then removed by a change of origin, except for the
nonzero half-period class in the rectangular case. That class gives
the first row. Thus every real genus-one curve occurs in one of
these connected families. The number of real ovals is locally
constant in a smooth proper real family: local trivialization can
be chosen invariant under the involution by an invariant horizontal
metric and its path lifting. Distinct rows consequently belong to
different components. This proves the number three, with no origin
added to the moduli problem. Direct substitution of the four
unstable types into the displayed partition formula gives exactly
these counts.

The mapping-class exceptions can be checked directly. On $S^2$ reflection and the antipodal
map have different fixed sets, but both represent the orientation-reversing
mapping class. On $\R^2/\Z^2$ the maps $(x,y)\mapsto(x,-y)$ and
$(x,y)\mapsto(x+1/2,-y)$ have respectively two fixed circles and no fixed
points, but both induce $\operatorname{diag}(1,-1)$ and hence the same mapping
class. On the annulus, the reflection and glide anti-involutions of B.15 also have
the same mapping class: the path $(x,t)\mapsto(x+s/2,-t)$ joins them through
homeomorphisms preserving the two ends as a set. These examples establish the exact failure of the unrestricted
injectivity hoped for on scan 440. For labelled moduli spaces or a coarse
moduli space, the quotient and its real locus have to be specified separately;
the displayed formula is for the stack and puncture convention stated above.
\end{proof}

\begin{claim}{B.18}{Proved}{The two centralizer corollaries in both quotient families}
Let $g\ge2$, and let an anti-involution of $X_g$ have $b>0$ ovals, with outer
class $\sigma$. There are exactly $b$ classes of order-two genuine
loop automorphisms lifting $\sigma$, under conjugacy by $\pi_g$.
For each lift $u$, $\pi_g^u\cong\Z$. The positive centralizer of $\sigma$ acts
transitively on those $b$ classes. More precisely, writing
$A=\Aut_{\rm lac}(\pi_g)$ and $O=\Out_{\rm lac}(\pi_g)$, there is an exact sequence
\[
 1\longrightarrow\pi_g^u\longrightarrow C_{A^+}(u)
 \longrightarrow\operatorname{Stab}_{C_{O^+}(\sigma)}([u])\longrightarrow1,
\]
and the last group has index $b$ in $C_{O^+}(\sigma)$. The corresponding full
centralizer image also has index $b$. This proves Corollaries 1 and 2,
folios 225--226, scans 445--448, for both orientable and nonorientable quotients.
\end{claim}
\passage{33; 225-226; 445-448}

\begin{proof}
B.4 proves the count and the fixed subgroup. Any permutation of the boundary
components of a compact connected quotient surface is realized by a homeomorphism:
join two components by an embedded arc, interchange them in a neighbourhood of
that arc and the two collars, and use transpositions. Each quotient homeomorphism
lifts to the oriented double; among its two lifts choose the orientation-preserving
one. It commutes with the reflection. Consequently the positive centralizer is
transitive on the ovals, and hence on the section classes.

For exactness, the centre of $\pi_g$ is trivial, so $\Inn(\pi_g)$ is identified
with $\pi_g$. Its intersection with $C_A(u)$ is exactly $\pi_g^u$.
An outer class $\alpha$ centralizing $\sigma$ preserves $[u]$ precisely when
one, hence every, lift $\widetilde\alpha$ satisfies
$\widetilde\alpha u\widetilde\alpha^{-1}=vuv^{-1}$ for an inner automorphism
$v$. Then $v^{-1}\widetilde\alpha$ centralizes $u$, proving surjectivity onto
the displayed stabilizer. Orbit--stabilizer gives index $b$ in both versions.
For $b=0$ there are no order-two genuine lifts and these assertions about a
chosen lift have no instance. With punctures the same exact sequence holds,
but the orbit size records the boundary decorations from B.16; it need not be
the total number of ovals when different ovals carry different puncture data.
\end{proof}

\begin{claim}{B.19}{Proved}{The tube twists, their maximality property, and their normalizer}
For $g\ge2$, the elements $\rho_j$ of \S33, scans 453--470, are full Dehn twists
about a cut system $c_1,\ldots,c_g$: disjoint nonseparating curves with connected
complement. They generate $H\cong\Z^g$. This is not a maximal free abelian
subgroup of $\mathcal T_g$: it is contained in a multitwist subgroup of rank
$3g-3$. It has the precise maximality property of being a full multitwist group
of largest rank whose supporting multicurve has connected complement.
Furthermore
\[
 N_{\mathcal T_g}(H)=\operatorname{Stab}_{\mathcal T_g}(\{c_1,\ldots,c_g\}),
 \qquad N_{\mathcal T_g}(H)/C_{\mathcal T_g}(H)\cong S_g\times C_2.
\]
The action on $H$ permutes its basis and multiplies every basis vector by the
same orientation sign. Cutting and capping gives an exact sequence
\[
 1\longrightarrow H\longrightarrow N_{\mathcal T_g}(H)
 \longrightarrow\mathcal T^{\rm pairs}_{0,2g}\longrightarrow1,
\]
where the last extended sphere mapping class group preserves the partition of
its $2g$ punctures into the $g$ unordered gluing pairs.
\end{claim}
\passage{33; 229-237; 453-470}

\begin{proof}
In the coordinates of scan 461 the map is
$\rho_j(z,t)=(-z e^{\pi i t},t)$ on $S^1\times[-1,1]$ and is the identity outside
the tube. Its angular displacement increases from $0$ to $2\pi$, so its relative
boundary mapping class is one full Dehn twist, with sign fixed by that coordinate
orientation. The complement of the tubes in his construction is a sphere with
$2g$ boundary circles, establishing the cut-system assertion.

Twists in disjoint collars commute. They are independent: cut the surface along
the $c_j$, choose transverse curves dual to these cuts, and measure the integer
winding of the resulting arcs when the collars are glued back. A product
$\prod T_{c_j}^{n_j}$ changes the $j$th such winding by $n_j$; an isotopy cannot
change these relative-end winding differences. This proves $H\cong\Z^g$.
Complete the $2g$-holed sphere to a pants decomposition by adding $2g-3$
essential disjoint curves. The same winding argument gives a rank-$3g-3$
multitwist group containing $H$. In contrast, cutting $k$ disjoint curves while
keeping the complement connected reduces the genus to $g-k$, so $k\le g$.
This proves the stated maximality among connected-complement multitwist systems.

For the normalizer, conjugation satisfies
$fT_cf^{-1}=T_{f(c)}^{\epsilon(f)}$. A single twist can equal a product of twists
about the $c_j$ only by being one of these twists to the corresponding power.
Here is the needed support argument. Such a twist commutes with every $T_{c_j}$.
The elementary intersection formula
$i(T_d(c),c)=i(c,d)^2$ for a single twist forces $d$ to be disjoint from each
$c_j$; otherwise its conjugation changes $T_{c_j}$. Add $d$ to the disjoint
system if it is not already there and apply the preceding independent-winding
argument. The resulting nontrivial relation is impossible. Thus a normalizer
permutes precisely the $c_j$, proving the first equality. Conversely every such
permutation of curves normalizes $H$. Every permutation of the handles is
realized positively, and a reflection of the standard cut-system model reverses
the global orientation while preserving each curve. This gives the full image
$S_g\times C_2$. Its kernel is exactly the centralizer.

A representative preserving the multicurve can be isotoped to preserve its
collars. Its restriction to the cut surface, after capping each boundary to a
puncture, preserves the gluing pairs. Conversely a homeomorphism of the capped
sphere preserving those pairs can be made radial in small puncture discs and
then extended over the paired collars. Two such extensions differ by integer
collar twists. An isotopy on the capped sphere lifts after allowing exactly
these twists, by the isotopy extension theorem on collars. This proves both
surjectivity and the stated kernel, and types the comparison with the
complement suggested on scan 470.
\end{proof}

\begin{claim}{B.20}{Refuted}{Presentation (52) and its exact correction}
For $g\ge2$, presentation (52) on scan 469 does not present the stated
$\Z^g$-by-finite group. Put $r=\varepsilon_0\varepsilon_1$ and
$\rho_j=r^j\tau'\tau r^{-j}$, $0\le j<g$. The printed presentation gives
\[
 F_g\rtimes\bigl(D_{4g}\times\langle\tau\rangle\bigr),
\]
where $D_{4g}=\langle r,\varepsilon_1\mid r^{2g}=\varepsilon_1^2=1,
\varepsilon_1r\varepsilon_1=r^{-1}\rangle$ permutes the free generators through
its action on $\Z/g$, and $\tau$ inverts each free generator.
Adding the relations $[\rho_i,\rho_j]=1$ for $i<j$ gives exactly
$\Z^g\rtimes(D_{4g}\times C_2)$ with the manuscript's specified action.
\end{claim}
\passage{33; 237; 469}

\begin{proof}
The relations say that $\tau$ commutes with $D_{4g}$, that $\tau'$ commutes
with $\varepsilon_1$ and $r^g$, and that both $\tau,\tau'$ are involutions.
Set $\rho_0=\tau'\tau$. Then $\tau\rho_0\tau=\rho_0^{-1}$,
$\varepsilon_1\rho_0\varepsilon_1=\rho_0$, and
$r^g\rho_0r^{-g}=\rho_0$. Its conjugates are indexed by the $g$ cosets of
$\langle\varepsilon_1,r^g\rangle$ in $D_{4g}$, with no further relation among
them. To check the last assertion directly, form the displayed semidirect
product with a free basis $\rho_0,\ldots,\rho_{g-1}$ and define
$\tau'=\rho_0\tau$ in it. All printed relations hold. The reverse map sends
these free generators to $r^j\tau'\tau r^{-j}$; the two maps agree inversely
on all generators. This is a proof by explicit mutually inverse homomorphisms.
Abelianizing only the free factor adds exactly the stated commutators and gives
the corrected group.

There is also a finite certificate. For $g=2$ set
$\varepsilon_0=(13)$, $\varepsilon_1=\tau=1$, $\tau'=(12)$ in $S_3$.
All nine displayed relators evaluate to $1$, but
$\rho_0=(12)$ and $\rho_1=(23)$ do not commute. The program checks every one
of these evaluations. Relations (33), scan 454, already assert commutation
for the geometric twists; they are exactly the missing input to (52).
\end{proof}

\begin{claim}{B.21}{Proved}{The one-holed-torus subgroups and the proved part of the analogy}
For $g\ge2$, choose a one-holed torus $W_j$ containing $c_j$ and a transverse curve $d_j$
with $i(c_j,d_j)=1$. Its group fixing its boundary pointwise embeds by extension
by the identity into $\mathcal T_g^+$. It is
\[
 B_3=\langle A_j,B_j\mid A_jB_jA_j=B_jA_jB_j\rangle,
\]
and capping the boundary gives
\[
 1\longrightarrow\langle T_{\partial W_j}\rangle\cong\Z
 \longrightarrow B_3\longrightarrow SL_2(\Z)\longrightarrow1,
 \qquad (A_jB_j)^6=T_{\partial W_j}^{\pm1}.
\]
This extension has no section. For disjoint chosen $W_j$, their supported
subgroups commute and, together with $H$, preserve their boundary multicurve;
they do not generate $\mathcal T_g^+$. On explicit generators the tube subgroup
and these two twists act by
\[
 \begin{array}{lll}
 \rho_j:&a_j\mapsto a_j,&b_j\mapsto b_ja_j,\\
 B_j:&a_j\mapsto a_jb_j^{-1},&b_j\mapsto b_j,
 \end{array}
\]
fixing the other handle generators of
$\pi_g=\langle a_j,b_j\mid\prod_j[a_j,b_j]=1\rangle$.
\end{claim}
\passage{33; 236-237; 467-470}

\begin{proof}
Choose the based handle loops so that their commutator is the boundary loop.
The displayed substitutions preserve $[a_j,b_j]$ exactly, and are the effect
of crossing the appropriate twist collar. Twists in curves meeting once obey
the braid relation: in these generators substitution gives $ABA=BAB$.
The sixfold product $(AB)^6$ conjugates both generators by
$[a,b]^{-1}$, which is the boundary twist with this collar convention.
Both word equalities are checked by free reduction in
\texttt{code/free\_words.py}; the stated substitutions also give their direct
verification. Filling the boundary gives the standard two elementary matrices
generating $SL_2(\Z)$. Their action on primitive slopes, reduced by the Euclidean
algorithm, shows that the filled mapping class group is exactly $SL_2(\Z)$.
The kernel of filling is the boundary winding: cut along two based generating
arcs, isotope a kernel representative to fix them, and read its remaining
integer twist in the boundary collar. The kernel is therefore the infinite
cyclic group generated by the boundary twist. The braid group on three strands
has the same central extension: with $x=ABA,y=AB$ it is
$\langle x,y\mid x^2=y^3\rangle$, and its quotient by $y^6$ is $SL_2(\Z)$.
The comparison sends the kernel generator to the boundary twist, so it is an
isomorphism of these extensions.

Extension by the identity is injective because an isotopy to the identity on
the closed surface would fix the essential boundary and, after collar adjustment,
give an isotopy on the one-holed torus; a nonzero boundary twist remains
nontrivial by the winding argument in B.19. The braid group is torsion-free:
$\langle x,y\mid x^2=y^3\rangle$ is the amalgam of two infinite cyclic groups
along their subgroups of indices two and three; a finite subgroup fixes a vertex
in its tree and is conjugate into an infinite cyclic group. A section from
$SL_2(\Z)$ would inject its order-four element, contradicting this fact.

Disjoint supports commute. Each supported homeomorphism fixes the chosen
boundary curves, and so does every generator of $H$. An essential curve crossing
one of those boundaries has a twist which does not preserve it, by the
intersection formula in B.19. Hence the generated subgroup is proper.
The exact analogy proved here is the homomorphism to the usual rank-two
integral representation and its central boundary-winding extension, with braid
and disjoint-support relations. For overlapping choices of the strips one
must specify their intersections before computing the remaining relations;
no unspecified choice of strips supplies additional relations by itself.
\end{proof}

\subsection{Computations and the limits of the requested table}

\begin{claim}{B.22}{Refuted}{A finite subgroup alone does not index the requested abelian table}
For all types in Part B(f), a table indexed only by a finite subgroup of the
mapping class group cannot assign unique splitting numbers or $\Phi$.
The type $(1,0)$ and the single subgroup
$\langle\operatorname{diag}(1,-1)\rangle$ already have two different answers.
The additional extension datum is precisely the class in $H^2(G,\pi)$ in the
abelian cases of B.14.
\end{claim}
\passage{15; 78-79; 157-159; prompt B.f}

\begin{proof}
On the torus compare
$\sigma_0(x,y)=(x,-y)$ and
$\sigma_1(x,y)=(x+1/2,-y)$. Both are faithful involutions with the same
homomorphism to $GL_2(\Z)$, but the first fixes two circles and the second fixes
nothing. The corresponding extensions, with commuting lattice generators
$a,b$, have relations
\[
 \widetilde\sigma a\widetilde\sigma^{-1}=a,\qquad
 \widetilde\sigma b\widetilde\sigma^{-1}=b^{-1},\qquad
 \widetilde\sigma^2=a^\epsilon,\quad \epsilon=0\text{ or }1.
\]
The square of $a^mb^n\widetilde\sigma$ is $a^{2m+\epsilon}$. For $\epsilon=1$
there is no order-two lift. For $\epsilon=0$ such lifts require $m=0$,
and conjugation by $b$ changes $n$ by two. Thus there are exactly two section
classes in the first case and zero in the second. These are the two parity
classes of the extension cocycle. B.14 proves that retaining the extension
class gives the correct affine classification, with all translation conjugacies
accounted for. Thus the counterexample identifies the missing index and its
exact replacement.
\end{proof}

\begin{claim}{B.23}{Proved}{Complete finite-subgroup table in type $(0,3)$}
The ten conjugacy classes in the extended group
$\mathcal T_{0,3}\cong S_3\times C_2$, and their discrete section invariants,
are the following. In the table $F$ counts full-section classes, $P$ counts all
nontrivial partial-section classes with their subgroup images distinguished,
$M=|\Phi|$ counts maximal partial-section classes, and $O$ counts the group
orbits on $\Phi$.
\[
\begin{array}{lrrrrr}
 G&|G|&F&P&M&O\\\hline
 1&1&1&0&0&0\\
 C_2\text{ (positive)}&2&1&1&1&1\\
 C_2\text{ (horizontal mirror)}&2&3&3&3&3\\
 C_2\text{ (vertical mirror)}&2&1&1&1&1\\
 C_3&3&2&2&2&2\\
 C_2^2&4&1&6&3&2\\
 C_6&6&0&5&5&2\\
 S_3\text{ (positive)}&6&0&5&5&2\\
 S_3\text{ (mixed)}&6&2&7&2&2\\
 S_3\times C_2&12&0&16&5&2
\end{array}
\]
The attached CSV gives every maximal stabilizer; the JSON table gives every
partial section and all of its multiplier values. In particular it checks the
reflection failure of maximal-section uniqueness from B.8 inside this whole list.
\end{claim}
\passage{15; 75-77; 150-155; prompt B.f}

\begin{proof}
Place the three punctures at the vertices of an equatorial equilateral triangle.
The positive rotation group of that triangle is $S_3$: it is generated by the
$120$-degree vertical rotation $r$ and a half-turn $a$ around an axis through a
puncture. Reflection $\tau$ in the equator commutes with both. The mapping class
group map to permutations and orientation is an isomorphism. For its kernel,
fix three arcs forming a tree through the punctures by isotopy; its complement
is a disc, and Alexander's disc isotopy makes the remaining map isotopic to the
identity. This proves completeness of the ambient finite group and its stated
realization.

The program represents this group exactly by pairs $(p,e)\in S_3\times C_2$.
Starting with the trivial subgroup it repeatedly adjoins every group element
and takes closure, until no new subgroup occurs. Every subgroup is generated
by its elements, so this terminating procedure is exhaustive. It finds sixteen
subgroups and ten conjugacy classes, represented in the displayed order by
\[
 1,\ \langle a\rangle,\ \langle\tau\rangle,\
 \langle a\tau\rangle,\ \langle r\rangle,\ \langle a,\tau\rangle,\
 \langle r,\tau\rangle,\ \langle r,a\rangle,\
 \langle r,a\tau\rangle,\ \langle r,a,\tau\rangle.
\]
Here is the independent geometric input to its section calculation. Denote the
poles by $N,S$, and the three surviving opposite equatorial axis points by
$Q_i$; the other three axis points are the removed punctures. The fixed set of
$\langle r\rangle$ is $\{N,S\}$, and that of a positive half-turn is its single
surviving $Q_i$. The horizontal mirror has three open intervals, each containing
one $Q_i$. A vertical mirror has one real line after its single puncture is
removed, and this line contains $N,S,Q_i$. Intersect these fixed sets to find the
fixed sets of every subgroup. B.3 and B.4 turn their components into exactly
its section classes. A partial section on a mirror interval is contained in a
point stabilizer precisely when that point lies on the interval; this determines
all inclusions and hence all maximal classes. The code records and checks these
inclusions and the conjugation action of every element.

Every reflection and positive half-turn has multiplier $1/2$. For $r$, the
values at $N,S$ are respectively $1/3,-1/3$, and its powers give all remaining
values. This computes all entries in the attached JSON. The geometry also
verifies that all the centralizers occurring for nontrivial sections here are
trivial: every mirror component is a line, and every positive section fixes an
isolated point. The absence of a common finite overgroup for the two vertex
stabilizers along a vertical mirror is retained in the inclusion data.
\end{proof}

\begin{claim}{B.24}{Proved}{Completed finite-subgroup catalogue and exact indexing}
Every finite-subgroup conjugacy class in the requested range
$2g+\nu\le6$ has now been enumerated.  For the hyperbolic types,
the numbers of classes in the extended mapping class group are
\[
\begin{array}{c|r|l}
 (g,\nu)&\#\text{ classes}&\text{proofs and executed tables}\\\hline
 (0,3),(0,4),(0,5),(0,6)&10,30,20,57&B.23,B.25,B.26\\
 (1,1),(1,2),(1,3),(1,4)&13,40,38,115&B.25,B.27\\
 (2,0),(2,1),(2,2)&92,32,121&B.28,B.29\\
 (3,0)&292&B.33,B.37\text{--}B.41.
\end{array}                                             \tag{B.24.1}
\]
Each table includes the full and partial section classes, multipliers,
centralizers, extension incidences, maximal section classes $\Phi$,
and their group orbits and stabilizers.  In the four exceptional
abelian types, the finite mapping-class subgroup counts are
respectively $2,2,5,13$ for sphere, plane, annulus and torus.  Their
injectively realized finite-action counts are $3,2,7,17$; the exact
covering extension is retained as part of the section index in B.36,
as required by the explicit failures of extension-independent indexing
in B.22.  The families with a nontrivial kernel on passage to mapping
classes are described in B.14.
The exact tube words, stabilizer action, strip relations and arbitrary
transversal obstruction formerly required here are proved and executed
in B.30--B.32.
\end{claim}
\passage{\S{}18; folios 96--101; scans 193--203; \S{}33; folios 222--237; scans 440--470; Part B(e) and B(f) of the task}

\begin{proof}
The first three rows of (B.24.1) are precisely the complete catalogues
cited there.  We spell out the closed genus-three reconciliation to
record the coverage and the absence of double counting.
The positive portion has $55$ classes by B.33.  For a mixed finite
action let $H$ be its positive subgroup.  The positive quotient has
genus zero, one, two or three, by the positive classification.
The genus-zero case gives either a mirror disc ($145$ classes,
B.37) or a projective plane ($11$ classes, B.39).  The genus-one case
gives an annulus or M\"obius band ($57$ classes, B.40), or a Klein
bottle ($6$ classes, B.41).  The genus-two case is the free positive
$C_2$ action and gives $12$ mixed classes in B.38.  The genus-three
case has trivial positive part and gives its six negative involution
classes in B.38.  Thus the number of mixed classes is
\[
                   145+11+57+6+12+6=237,
                     \qquad55+237=292.                \tag{B.24.2}
\]
Every conjugacy descends to homeomorphisms of both the positive
quotient and the full quotient orbifold.  The positive quotient genus,
and the full quotient's boundary and orientability, separate the
listed portions.  Within each portion the cited proofs calculate
exactly the geometric kernel-normalizer orbits, so these counts are
irreducible counts of surface-action conjugacy classes.

B.34 supplies an independent hyperelliptic construction and is
counted through its explicit overlaps, not added a second time.
All $79$ of its rows are identified by monodromy isomorphisms:
$18$ with positive B.33 rows, $57$ with B.37 rows, and four with B.39
rows.  The executed program
\path{reconcile_genus3_catalogue.py} records these $79$ maps and
checks the corresponding section counts.  Its unified
\path{type_3_0_complete_catalogue.csv} has $292$ different
class identifiers, and \path{type_3_0_complete_sections.json}
contains the full finite group and section model for every row.
The exceptional replacement in B.36 is complete with its stated
realization index; it retains, in particular, the sphere antipodal
action's algebraic splitting without identifying it with a geometric
fixed point.  This proves the claimed coverage with exactly the
comparisons and indices used by the computations.
\end{proof}

\begin{claim}{B.25}{Proved}{Complete executed catalogues in types $(1,1)$ and $(0,4)$}
The extended mapping class groups of types $(1,1)$ and $(0,4)$ have,
respectively, $13$ and $30$ conjugacy classes of finite subgroups, including
the trivial subgroup.  The tables below give every class and the numbers of
full splittings, nontrivial partial splittings, maximal partial section
classes $\Phi$, and group orbits on $\Phi$.  The accompanying exact-arithmetic
files give every section, its image subgroup, its multiplier on every
element, its centralizer, all inclusion relations, the group stabilizers,
and the puncture stabilizers.  These are exhaustive catalogues for the two
specified surface types.
\end{claim}
\passage{\S{}15; folios 73--75; scans 148--151; \S{}15; folios 77--79; scans 155--160; \S{}18; folios 96--101; scans 193--203; Part B(f) of the task; no further manuscript assertion}

\begin{proof}
Put
\[
 S=\begin{pmatrix}0&-1\\1&0\end{pmatrix},\qquad
 R=\begin{pmatrix}0&-1\\1&1\end{pmatrix},\qquad
 F=\begin{pmatrix}0&1\\1&0\end{pmatrix}.
\]
For the punctured torus use their linear actions on
$\R^2/\Z^2\setminus\{0\}$.  The homology map identifies its extended
mapping class group with $\mathrm{GL}_2(\Z)$: surjectivity is supplied by
linear maps, and a map acting trivially on homology can be isotoped to fix
two transverse basis curves, after which the complementary punctured disc
is handled by the Alexander isotopy.  This is the punctured version of
\cite[Theorem 2.5 and pp.~54--55]{A_FarbMargalit2012}.

For the four-punctured sphere use the pillowcase
\[
 \bigl(\R^2/\Z^2\setminus (\tfrac12\Z^2/\Z^2)\bigr)/\{p\sim-p\}.
\]
Write $U(p)=p+(1/2,0)$ and $V(p)=p+(0,1/2)$.  The extended mapping
class group is
$V_4\rtimes\mathrm{PGL}_2(\Z)$ with these affine representatives.
The map to $\mathrm{PGL}_2(\Z)$ is obtained by lifting through the double
branched torus cover and acting on torus homology.  Its kernel consists of
the four half translations: a kernel element fixes all slope classes;
after a half translation it fixes the four punctures and two filling
curves, so the Alexander isotopy on their complementary discs kills it.
Linear maps split the homomorphism and act on half translations by
reduction modulo two.  This is the explicit morphism of
\cite[Proposition 2.7 and the extended statement on p.~57]{A_FarbMargalit2012}.

Here are finite amalgams that make enumeration complete.  The matrices
$S,R$ generate $\mathrm{SL}_2(\Z)$ by the Euclidean algorithm, since
$S^{-1}R=\left(\begin{smallmatrix}1&1\\0&1\end{smallmatrix}\right)$.
Modulo $\{\pm I\}$ their images have orders two and three, with no
further relation: on the real projective line, $\bar S$ sends the negative
half-line to the positive half-line, and both nonidentity powers of
$\bar R$ send the positive half-line into the negative half-line.  The
free-product ping-pong argument gives
$\mathrm{PSL}_2(\Z)=C_2*C_3$.
Conjugation by $F$ fixes $\bar S$ and inverts $\bar R$.  Consequently
\[
 \mathrm{PGL}_2(\Z)=V_4*_{C_2}S_3,
 \qquad
 \mathrm{GL}_2(\Z)=\langle S,F\rangle
       *_{\langle-I,F\rangle}\langle R,F\rangle
       =D_8*_{V_4}D_{12},
\]
where $D_k$ has order $k$.  The second identity follows also by lifting
the first through their common central subgroup $\langle-I\rangle$;
the normal form for the amalgam leaves no additional kernel.
Adding the normal group of half translations gives the pillowcase amalgam
\[
 (V_4\rtimes V_4)*_{V_4\rtimes C_2}(V_4\rtimes S_3),
\]
whose vertex groups have orders $16$ and $24$, and whose edge group has
order $8$.  The two vertex groups are $D_8\times C_2$ and $S_4$.

A finite subgroup of an amalgam fixes a vertex in its Bass--Serre tree.
Indeed the convex hull of a finite orbit is a finite tree; its centre is
fixed, and the action preserves the two vertex types, so an edge centre
also gives fixed endpoints.  Two vertex subgroups are conjugate in the
amalgam precisely through the equivalence generated by conjugacy in
each vertex group and identification of edge subgroups.  To prove the
nontrivial direction, a conjugacy makes the subgroup fix both endpoints
of the intervening path; it fixes every edge on that path, which
decomposes the conjugacy into those successive vertex and edge steps.
Thus finite enumeration in the two vertices, followed by these exact
fusion steps, is a complete conjugacy calculation in the infinite group.
The program finds $10$ and $16$ actual subgroups in the torus vertices,
and $35$ and $30$ in the pillowcase vertices.  Fusion leaves exactly
$13$ and $30$ classes.  All subgroup lists are saturated under adjoining
an arbitrary vertex-group element, so the enumeration has no search bound
that could omit a subgroup.

We next justify the computation of sections.  Let
$Q=(\tfrac1{12}\Z^2)/\Z^2$.  Every isolated fixed point needed for
either model lies in $Q$.  For the torus it solves $(A-I)p=0$; for the
pillowcase it solves
\[
 (A-\epsilon I)p=-t/2\pmod{\Z^2},\qquad
 \epsilon\in\{1,-1\},\quad t\in(\Z/2)^2.
\]
For every positive matrix in the finite vertices, the inverse matrices
and the particular solutions of the nonsingular equations have
denominators dividing $12$.  The program certifies this by exact
fractions for every matrix and translation, rather than inferring it
from a grid count.  The singular case for a nonidentity positive element
is a nontrivial translation and contributes no solutions to that sign.

For a reflection, the same equations are affine circle equations.  The
greatest common divisor of the nonzero matrix entries is one or two, so
every component meets $Q$.  On each such circle the primitive integral
kernel direction, divided by $12$, joins consecutive sampled points.
In the pillowcase the plus and minus circles meet only at two-torsion
points, all of which are included in $Q$.  Quotienting the sampled graph
by $p\sim-p$ therefore gives a subdivision of the true fixed circles.
Deleting the prescribed punctures gives exactly their connected
components.  A component adjacent to a deleted point is a line; the
others are ovals.  All rotations and all punctures are retained as
vertices, so no incidence needed for restriction of sections is lost.

B.3 and B.4 identify these fixed components with section classes.
For a partial section with a rotation use its fixed point; for a
reflection use its fixed line or oval.  Inclusion of image subgroups and
containment of the fixed point in the fixed component determine exactly
when one section extends another.  Applying each affine element gives
the exact action on this finite section poset.  This computes $\Phi$,
its orbits and its stabilizers without replacing pointwise section data
by conjugacy of image subgroups in the quotient.  Centralizers are $1$
for isolated points and lines and $\Z$ for ovals, by B.3--B.4.

At a torus fixed point the local positive derivative is $A$.  At a
nonbranch pillowcase fixed point it is $\epsilon A$, with the unique
sign specified by $Ap+t/2=\epsilon p$.  For these matrices the traces
$-2,-1,0,1$ give orders $2,3,4,6$; the half-turn has multiplier $1/2$,
and otherwise the sign of $\det(e_1,Ae_1)$ selects $+1/n$ or $-1/n$.
Reflections have multiplier $1/2$.  At a removed pillowcase branch
point the quotient coordinate squares a torus coordinate, so the
positive rotation angle doubles.  The program records these puncture
data separately.  It checks the order and power identities for every
element of every section and checks that every section containing a
rotation has exactly one maximal overgroup.  The reflection classes
with several maximal extensions are retained explicitly.

Finally, C.5--C.6 show that the chosen finite actions compute the action
attached to each mapping class subgroup up to conjugacy, so these
fixed-component computations depend on the listed subgroup class, not
on an unrecorded realization choice.
\end{proof}

In the tables $f$ is the number of full splitting classes, $p$ the number
of nontrivial partial splitting classes with image subgroup retained,
$\phi=|\Phi|$, and $o=|\Phi/\Gamma|$.  Juxtaposition in the generator
column means composition of the matrices or affine maps defined above.
For $(0,4)$ the letters $S,R,F$ mean their pillowcase classes.
The complete output files are
\texttt{type\_1\_1\_finite\_subgroups.csv},
\texttt{type\_0\_4\_finite\_subgroups.csv}, and the corresponding
\texttt{partial\_sections.json} files, produced by
\texttt{catalogue\_low\_complexity.py}.

\begin{center}Type $(1,1)$\end{center}
\begin{longtable}{lllrrrrr}
\toprule
Class & Group & Generators & $|\Gamma^+|$ & $f$ & $p$ & $\phi$ & $o$\\
\midrule\endhead
T01 & $1$ & $1$ & 1 & 1 & 0 & 0 & 0\\
T02 & $C_2$ & $SF$ & 1 & 2 & 2 & 2 & 2\\
T03 & $C_2$ & $SSF$ & 1 & 1 & 1 & 1 & 1\\
T04 & $C_2$ & $SS$ & 2 & 3 & 3 & 3 & 3\\
T05 & $C_3$ & $RR$ & 3 & 2 & 2 & 2 & 2\\
T06 & $C_2^2$ & $SS,SF$ & 2 & 3 & 10 & 3 & 3\\
T07 & $C_2^2$ & $SS,SSF$ & 2 & 1 & 6 & 3 & 2\\
T08 & $C_4$ & $S$ & 4 & 1 & 4 & 3 & 2\\
T09 & $S_3$ & $RR,RRF$ & 3 & 2 & 7 & 2 & 2\\
T10 & $S_3$ & $RR,RF$ & 3 & 0 & 5 & 5 & 2\\
T11 & $C_6$ & $R$ & 6 & 0 & 5 & 5 & 2\\
T12 & $D_8$ & $S,SF$ & 4 & 1 & 15 & 3 & 2\\
T13 & $D_{12}$ & $R,RRF$ & 6 & 0 & 16 & 5 & 2\\
\bottomrule
\end{longtable}

\begin{center}Type $(0,4)$\end{center}
\begin{longtable}{lllrrrrr}
\toprule
Class & Group & Generators & $|\Gamma^+|$ & $f$ & $p$ & $\phi$ & $o$\\
\midrule\endhead
S01 & $1$ & $1$ & 1 & 1 & 0 & 0 & 0\\
S02 & $C_2$ & $SF$ & 1 & 4 & 4 & 4 & 4\\
S03 & $C_2$ & $SFV$ & 1 & 1 & 1 & 1 & 1\\
S04 & $C_2$ & $SFUV$ & 1 & 0 & 0 & 0 & 0\\
S05 & $C_2$ & $F$ & 1 & 2 & 2 & 2 & 2\\
S06 & $C_2$ & $V$ & 2 & 2 & 2 & 2 & 2\\
S07 & $C_2$ & $S$ & 2 & 0 & 0 & 0 & 0\\
S08 & $C_3$ & $RR$ & 3 & 1 & 1 & 1 & 1\\
S09 & $C_2^2$ & $V,SF$ & 2 & 2 & 9 & 4 & 3\\
S10 & $C_2^2$ & $V,SFU$ & 2 & 0 & 3 & 3 & 2\\
S11 & $C_2^2$ & $UV,SF$ & 2 & 0 & 6 & 6 & 3\\
S12 & $C_2^2$ & $UV,SFV$ & 2 & 2 & 6 & 2 & 2\\
S13 & $C_2^2$ & $UV,F$ & 2 & 2 & 8 & 2 & 2\\
S14 & $C_2^2$ & $SF,F$ & 2 & 0 & 6 & 6 & 3\\
S15 & $C_2^2$ & $SFUV,F$ & 2 & 0 & 2 & 2 & 1\\
S16 & $C_4$ & $FU$ & 2 & 0 & 2 & 2 & 1\\
S17 & $C_2^2$ & $V,U$ & 4 & 0 & 6 & 6 & 3\\
S18 & $C_2^2$ & $UV,S$ & 4 & 0 & 2 & 2 & 1\\
S19 & $C_4$ & $SU$ & 4 & 2 & 4 & 2 & 2\\
S20 & $S_3$ & $RR,FR$ & 3 & 1 & 8 & 4 & 2\\
S21 & $C_2^3$ & $V,U,SF$ & 4 & 0 & 18 & 6 & 3\\
S22 & $C_2^3$ & $UV,SF,F$ & 4 & 0 & 12 & 6 & 2\\
S23 & $C_4\times C_2$ & $FU,SU$ & 4 & 0 & 8 & 6 & 2\\
S24 & $D_8$ & $FU,V$ & 4 & 0 & 12 & 6 & 2\\
S25 & $D_8$ & $SU,SFV$ & 4 & 2 & 16 & 2 & 2\\
S26 & $D_8$ & $FU,SFV$ & 4 & 0 & 6 & 2 & 1\\
S27 & $D_8$ & $SU,V$ & 8 & 0 & 8 & 6 & 2\\
S28 & $A_4$ & $RR,RUR$ & 12 & 0 & 10 & 10 & 2\\
S29 & $D_8\times C_2$ & $FU,SU,V$ & 8 & 0 & 28 & 6 & 2\\
S30 & $S_4$ & $FUR,URF$ & 12 & 0 & 32 & 10 & 2\\
\bottomrule
\end{longtable}

\begin{claim}{B.26}{Proved}{Complete executed spherical catalogues through six punctures}
For $\nu=3,4,5,6$, the extended mapping class group of the sphere with
$\nu$ unlabelled punctures has respectively
\[
                         10,\quad30,\quad20,\quad57
\]
conjugacy classes of finite subgroups, including the trivial subgroup.
The exact finite parameterization below gives every class once.  The
executed files \texttt{type\_0\_\nu\_orbifold\_catalogue.csv} and
\texttt{type\_0\_\nu\_orbifold\_sections.json}, for each of these four
values of $\nu$, give all partial-section classes, all their inclusion
relations, $\Phi$, its group action and stabilizers, the centralizers,
the multiplier on each section element, and the puncture stabilizers.
Thus these files settle the whole computation requested in B(f) for
the hyperbolic spherical types with $2g+\nu\le6$.
\end{claim}
\passage{\S{}15; folios 73--75; scans 148--151; \S{}15; folios 77--79; scans 155--160; \S{}18; folios 96--101; scans 193--203; \S{}33; folios 222--225; scans 440--445; Part B(f) of the task; no further manuscript assertion}

\begin{proof}
Extend an action over the punctures.  By B.1 and the sphere theorem of
\cite[Theorem 1.2]{bKolev2006}, a finite topological action on $S^2$ is
conjugate to an orthogonal action.  Write $Q=S^2/G$ with its local
orbifold groups.  A point has cyclic positive stabilizer, or a dihedral
stabilizer with mirrors; the mirror-only stabilizer has order two.
The covering formula is
\[
 \chi_{\rm orb}(Q)=\chi(|Q|)
  -\sum_{j=1}^r(1-1/m_j)
  -\frac12\sum_{j=1}^k(1-1/n_j)=\frac2N,
 \qquad N=|G|.                                      \tag{B.26.1}
\]
Here the $m_j$ are interior cone orders and the $n_j$ are corner orders.
One proves this formula directly by a triangulation whose vertices
include all cone and corner points: count each cell in the covering
with weight equal to the reciprocal of its stabilizer order.  Mirrors
are boundary strata; they introduce no additional term after the
boundary cells and their endpoints are counted together.

Since (B.26.1) is positive, the underlying compact surface is $S^2$,
$\mathbb{RP}^2$, or a disc.  The sphere has at most three cones.  A
single cone is impossible, since its local generator would be trivial
in the orbifold presentation.  Two cones have equal orders, since the
two local generators are inverses and both local cyclic groups inject.
For three cones the inequality $1/m_1+1/m_2+1/m_3>1$ gives precisely
$(2,2,n),(2,3,3),(2,3,4),(2,3,5)$.
The projective plane has at most one cone.  For a disc, positivity gives
$2r+k\le3$.  A single corner without a cone is impossible; two corners
without a cone have equal orders, by the same local-injectivity argument.
Three corners give the same four triples.  One cone and one corner
give $2/m+1/n>1$, hence $(m,n)=(2,n)$ or $(3,2)$.
Consequently the following table is exhaustive.  In its notation
$S(m_1,\ldots,m_r)$ is a sphere with cones, $P(m)$ is a projective
plane with one cone, $D(m_1,\ldots,m_r;n_1,\ldots,n_k)$ is a mirror
disc with the specified interior cones and cyclically ordered corners,
and an empty list means no such singularities.
\[
\begin{array}{c|c@{\qquad}c|c}
Q&N&Q&N\\\hline
S()&1&P()&2\\
S(n,n)&n&P(n)&2n\\
S(2,2,n)&2n&D(;)&2\\
S(2,3,3)&12&D(;n,n)&2n\\
S(2,3,4)&24&D(n;)&2n\\
S(2,3,5)&60&D(;2,2,n)&4n\\
&&D(2;n)&4n\\
&&D(;2,3,3)&24\\
&&D(;2,3,4)&48\\
&&D(;2,3,5)&120\\
&&D(3;2)&24
\end{array}                                                   \tag{B.26.2}
\]
Every integer $n$ in this table is at least two.  For $3\le\nu\le6$
only $2\le n\le6$ can occur.  Indeed, in every family depending on
$n$, the union of all point orbits of size strictly smaller than $n$
has at most two points; an invariant set of at least three punctures
therefore contains an orbit of size at least $n$.  This assertion also
follows immediately by dividing $N$ by each local group order in
(B.26.2).  In the two icosahedral rows every point orbit has at least
twelve points.  Those rows cannot occur in the present range.  These
are proved finite bounds, rather than bounds inferred from the output.

We now give the exact decoration rule used in the program.  A cone
orbit may be removed or retained; put $c_j\in\{0,1\}$ accordingly.
A corner orbit similarly has $d_j\in\{0,1\}$.  For each open mirror
edge let $e_i\ge0$ be the number of removed orbits on that edge; a
cornerless mirror circle is one edge for this purpose.  Finally let
$a\ge0$ be the number of regular interior orbits removed.  The single
equation is
\[
 \sum_jc_j\frac N{m_j}
 +\sum_jd_j\frac N{2n_j}
 +\sum_i e_i\frac N2+aN=\nu.                       \tag{B.26.3}
\]
On a sphere, permutations of cones of equal order identify decorations.
On a mirror disc, identify them under precisely the rotations and
reflections of the cyclic corner list that preserve its order labels.
The corresponding edge counts move with their edges.  For two corners
this rule includes the independent exchanges of the two corners and
the two edges.  For one or zero corners it introduces no further
identification of the single edge count.  There is at most one interior
cone in every remaining disc or projective-plane row.

For clarity, this rule also proves completeness of the conjugacy
classification.  A homeomorphism of an interval moves any two
unlabelled finite subsets of the same size into one another, and the
same assertion holds in the connected regular interior of each quotient
surface.  Such moves extend to the quotient surface by disjoint small
discs and strips.  The only remaining data are therefore the selected
singular strata and the cyclic mirror-edge counts in (B.26.3), modulo
the explicitly stated symmetries.  Conversely every quotient
homeomorphism preserves those strata, their orders, and their cyclic
incidence.  Since $S^2$ is the universal orbifold cover, a homeomorphism
between the decorated quotient orbifolds lifts to a conjugacy between
their covering actions; this follows by lifting paths and their local
cone or mirror charts.  A conjugacy of actions descends to such a
quotient homeomorphism.  Thus no extra equivalence of decorations is
possible.  Finally C.5 and C.6 identify these topological action classes
with conjugacy classes of finite subgroups in the extended mapping
class group of the punctured sphere.  The induced action there is
faithful: a nonidentity positive periodic map isotopic to the identity
would fix every puncture individually, whereas a nonidentity sphere
rotation has exactly two fixed points; a negative map cannot be
isotopic to the identity.  Here $\nu\ge3$ is essential.

Here are certificates for the finite group calculations.  The source
\texttt{catalogue\_spheres.py} constructs the cyclic and dihedral
families as regular permutation groups, the tetrahedral and octahedral
positive groups as $A_4,S_4$, their needed extensions by an explicitly
given central involution, and the full octahedral group as signed
permutations of three axes.  All generators, products, orders and signs
are recorded as finite permutations.  The program verifies each
orbifold relation, injection of every local cyclic or dihedral group,
and $N\chi_{\rm orb}(Q)=2$ using rational arithmetic.  These checks
prove that the permutation group is the entire orbifold group, as
follows.  The kernel of the displayed quotient of the orbifold
presentation produces a covering with no cone, corner or mirror
singularity because all local groups inject.  Its orientation
character is trivial, and its Euler characteristic is two.  It is
therefore a closed orientable sphere; its fundamental group is trivial.
This proves injectivity of the quotient presentation as well as its
already checked surjectivity.  In particular these certificates do
not merely exhibit a possibly proper finite quotient of an unexamined
presentation.

We describe how the section data are recovered without numerical
sampling.  Cone and corner lifts are cosets of their cyclic and
dihedral stabilizers.  In a disc presentation choose reflections
$r_0,\ldots,r_k$ with
\[
 (r_jr_{j+1})^{n_j}=1,\qquad
 r_k=e r_0e^{-1},\qquad x_1\cdots x_r e=1.            \tag{B.26.4}
\]
Indices of the corner orders are shifted here to start at zero.
The code stores $r_0,\ldots,r_{k-1}$ and derives the final reflection.
An edge lift is a coset of $\langle r_i\rangle$ and meets the adjacent
corner cosets; the closing incidence at edge zero is multiplied by
$e^{-1}$.  With no corners it instead joins the corresponding two
cosets of $\langle r_0\rangle$ by that same rule.  Insert the prescribed
removed vertices along each edge, retaining two auxiliary vertices in
each complementary segment.  These auxiliary vertices prevent a
one-edge loop from collapsing when represented as an ordinary finite
adjacency graph.  This is exactly the lifted mirror locus with a
finite subdivision.

For any reflection $r\in G$, keep the vertices whose stabilizers
contain $r$.  Before punctures are deleted the program verifies degree
two at every such vertex.  Its components are consequently the full
fixed circles of $r$.  Deleting the selected puncture vertices gives
all the fixed lines and ovals, with their exact incidence.  For a
subgroup whose positive part is nontrivial, the fixed points are
precisely the retained cone and corner lifts whose stabilizers contain
that subgroup.  A negative involution without a mirror has an empty
fixed set.  Thus B.3 and B.4 supply every section class, with centralizer
$1$ for a point or a line and $\Z$ for an oval.  A section over $H$
extends a section over $K<H$ exactly when its fixed component is
contained in the latter fixed component.  Hence the same finite
incidence calculation gives the entire partial-section poset and all
its maximal elements, including the reflection cases with several
maximal extensions.  Left multiplication on cosets gives the group
action and every stabilizer without a choice of orbit representatives.

The multiplier is exact as well.  Orient a local lift of each cone,
and at a corner use the positive generator $r_jr_{j+1}$ determined by
the oriented boundary order.  Write this generator as $u$, of order
$m$.  At the point represented by a coset $gD$, an element $h$ of its
positive stabilizer satisfies $g^{-1}hg=u^b$, and
\[
                   \mu(h)=\operatorname{sgn}(g)\frac bm
                                \pmod{\Z}.           \tag{B.26.5}
\]
Replacing $g$ by another coset representative leaves this value
unchanged: a positive local element commutes with $u$, while a
negative local element both reverses the sign and inverts $u$.
For a reflection the value is $1/2$, and for the identity it is zero.
Formula (B.26.5) is the local instance of B.12, and applies to all
sections in the files.  The code checks its order and every power
identity, as well as uniqueness of the maximal extension for every
section with nontrivial positive part.  All nontrivial subgroups of
each finite permutation group are enumerated by saturation under
adjoining each group element; this stops only when the subgroup set
is closed under every possible enlargement.

The following compact census is part of the executed output; each
entry is the number of conjugacy classes of that group order.  A dash
denotes zero.  Complete individual rows, rather than just this census,
are in the CSV files specified in the statement.
\[
\begin{array}{c|rrrrrrrrrrrrr|r}
\nu\backslash N&1&2&3&4&5&6&8&10&12&16&20&24&48&\text{total}\\\hline
3&1&3&1&1&-&3&-&-&1&-&-&-&-&10\\
4&1&6&1&11&-&1&7&-&1&1&-&1&-&30\\
5&1&4&1&4&1&3&1&3&1&-&1&-&-&20\\
6&1&7&1&14&1&9&8&1&9&1&-&4&1&57
\end{array}                                                    \tag{B.26.6}
\]
As an independent check, the program compares the multisets of all
six columns consisting of group order, positive order, full-section
count, nontrivial partial-section count, maximal-section count, and
group-orbit count, for $\nu=3,4$, against the previous permutation and
amalgam calculations in B.23 and B.25.  Both multisets agree, with
multiplicities.  The check and its inputs are attached; the completeness
proof above does not depend on that agreement.
\end{proof}

\begin{claim}{B.27}{Proved}{Complete executed punctured-torus catalogues through four punctures}
The extended mapping class groups of a torus with $\nu=1,2,3,4$
unlabelled punctures have, respectively,
\[
                         13,\quad40,\quad38,\quad115
\]
conjugacy classes of finite subgroups.  The files
\texttt{type\_1\_\nu\_affine\_catalogue.csv}, together with their
\texttt{affine\_sections.json} companions and
\texttt{torus\_affine\_normal\_forms.json}, list every class once.
They specify each group by its complete finite set of rational affine
transformations, its punctures, every full and nontrivial partial section
class, all extension incidences, $\Phi$ and its group orbits and stabilizers,
the centralizers, and $\mu$ on every section and puncture-stabilizer element.
Thus the torus portion of B(f) is complete for every hyperbolic type in the
requested range.
\end{claim}
\passage{\S{}15; folios 73--75; scans 148--151; \S{}15; folios 77--79; scans 155--160; \S{}18; folios 96--101; scans 193--203; Part B(f) of the task; no further manuscript assertion}

\begin{proof}
We give the finite classification underlying the program
\texttt{catalogue\_tori.py}; this also specifies the tables without relying
on abstract finite-group names.  By B.14 extend over the punctures and put
the action in affine form on $V/\Lambda=\R^2/\Z^2$.  Write
\[
 1\longrightarrow T\longrightarrow G\longrightarrow H\longrightarrow1,
 \qquad T\subset V/\Lambda,\quad H\subset\mathrm{GL}_2(\Z).
                                                        \tag{B.27.1}
\]
Here $T$ consists exactly of the translations in $G$.  It acts freely on
$V/\Lambda$, so each $T$-orbit on the puncture set has $|T|$ elements.
Consequently $|T|\mid\nu$ and in particular $|T|\le4$.  The complete list
of subgroups of $V/\Lambda$ of order at most four consists of the trivial
subgroup, the three cyclic subgroups of order two, the four cyclic
subgroups of order three, the six cyclic subgroups of order four, and
$(\frac12\Lambda)/\Lambda$.  Indeed a group of order at most four is
cyclic or $C_2^2$; a cyclic subgroup is generated by a point of that exact
order, and division by the number of generators gives precisely these
counts.  All fifteen groups lie in $(\frac1{12}\Lambda)/\Lambda$.
The thirteen possibilities for $H$, up to integral conjugacy, are exactly
those proved in B.25.  For each we retain only the $H$-invariant $T$.

We now enumerate the affine extensions, including the actions with no
common fixed point.  If $H$ contains a nonidentity positive element,
choose a generator $r$ of its cyclic positive subgroup.  Every affine
map with linear part $r$ has a fixed point, since $I-r$ is invertible
on $V$ and is therefore onto on $V/\Lambda$.  Translation of that point
to the origin makes this generator $(r,0)$.  A positive cyclic $H$
then gives exactly $G=\langle T,(r,0)\rangle$.  For a dihedral $H$
choose a negative generator $f$ and write its lift as $(f,b)$.  The two
relations in $H$ give exactly
\[
       (I-r)b\in T,\qquad (I+f)b\in T.                 \tag{B.27.2}
\]
Conversely these equations make the subgroup generated by
$T,(r,0),(f,b)$ a finite group with kernel $T$ and quotient $H$:
every word reduces to a translation followed by $r^j$ or $r^jf$, and
its linear part distinguishes these cosets.  Values of $b$ modulo $T$
are enough.  The finite set $(I-r)^{-1}T$ therefore supplies every
candidate, with no bound on words or finite quotients involved.
Translations of the origin preserving this normal form are exactly
$a$ with $(I-r)a\in T$; they change $b$ by $(I-f)a$ modulo $T$.

The remaining nontrivial linear groups have $H=\langle f\rangle$
of order two.  Put $W=(V/\Lambda)/T$.  Their extensions are represented
by $b\in W$ with $(I+f)b=0$, modulo $(I-f)W$.  Every such class has a
representative in $W[2]$.  To verify this assertion directly, choose a
lift $v\in V$ of $b$ for the enlarged lattice
$\Lambda_T=\{v\in V:v\bmod\Lambda\in T\}$.
Then $(I+f)v\in\Lambda_T$.  Write
$v=v_++v_-$ in the real $+1$ and $-1$ eigenspaces of $f$.
Subtracting $(I-f)(v_-/2)$ removes $v_-$, and the remaining $v_+$
satisfies $2v_+\in\Lambda_T$.  Thus the program uses precisely
\[
              2b\in T,\qquad (I+f)b\in T.             \tag{B.27.3}
\]
For two candidates $b,b'$ the translation-conjugacy test is also exact.
Let $w$ be a primitive integral row vector satisfying $wf=w$.
The image of $I-f$ on $V/\Lambda$ is the connected circle $\ker w$:
its real image is that eigenspace, and multiplication by any nonzero
integer is onto on a circle.  Hence the two candidates are equivalent
precisely when $w(b-b')\in w(T)$ in $\R/\Z$.
This includes both reflections and fixed-point-free glide reflections.

Integral changes of basis are handled without truncating an infinite
matrix group.  For $H=1$ or $\{\pm I\}$, use the generators $S,R,F$
of B.25 on the complete fifteen-element list of possible $T$'s.
Their orbit calculation is the full $\mathrm{GL}_2(\Z)$-conjugacy
calculation.  For a fixed $T$, a spanning tree of this finite orbit
and Schreier's formula give generators for its entire integral
stabilizer.  Explicitly, with a chosen representative $B_U$ carrying
$T$ to $U$, the generator associated to a matrix $A\in\{S,R,F,S^{-1},R^{-1}\}$
and the edge $U\mapsto AU$ is
\[
                       B_{AU}^{-1}AB_U.                \tag{B.27.4}
\]
For every other $H$ its integral normalizer is finite.  Here is an
exhaustiveness check for the small matrix enumeration in the code.
For a reflection conjugate to $\operatorname{diag}(-1,1)$, a commuting
integral unimodular matrix is diagonal with entries $\pm1$.
For a reflection conjugate to
$\left(\begin{smallmatrix}0&-1\\-1&0\end{smallmatrix}\right)$,
the equation is $A=\left(\begin{smallmatrix}a&b\\b&a\end{smallmatrix}\right)$,
with $(a-b)(a+b)=\pm1$, giving four solutions.  A normalizer of the
four-group generated by this reflection and $-I$ may also exchange
the two eigenspaces; the same calculation gives eight solutions.
A group with a rotation of order four has its normalizer in the eight
signed permutation matrices, since a commuting matrix has determinant
$a^2+b^2=1$.  A group with rotation of order three or six has its
normalizer in the twelve matrices for the hexagonal lattice: the
commuting equation gives the norm $a^2+ab+b^2=1$, and one reflection
supplies the other coset.  Testing these finite sets against $H$
gives the full normalizer.  Every resulting matrix has entries of
absolute value at most one, so the code's box of entries from $-2$
to $2$ contains every solution and introduces none outside the tested
normalizer equations.  Applying these matrices and the origin changes
above fuses all and only affine-conjugate candidates.
The result is $93$ affine action classes with $|T|\le4$ before choosing
punctures; some admit no invariant set of a given required size.

We explain why decorating these affine actions gives topological
conjugacy, including all permutations of the singular strata.
B.14 proves that a topological conjugacy between finite affine actions
can be replaced by an affine conjugacy.  The replacement also induces
the same permutation on the strata used here.  In detail write a lift
of the conjugacy as $h(x)=Ax+\delta(x)$ with $\delta$ periodic, and
put $\bar\delta=\int_{V/\Lambda}\delta$.  The equivariance equations
and invariance of Haar measure show that $x\mapsto Ax+\bar\delta$
has the same equivariance.  The straight interpolation of these two
lifts is an equivariant homotopy.  An isolated fixed point of a
nontrivial positive subgroup stays constant under this homotopy.
Each reflection circle stays in the same fixed-set component.
The degree on every such circle is fixed by the homology map $A$;
these circles are essential, since they lift to parallel lines.
Together with the images of the isolated vertices this determines
the permutation of the intervening open mirror intervals.  Thus the
affine normalizer gives exactly the possible stratum permutations.

To compute them, let $\mathcal V$ be the union of the fixed points of
all positive nontranslation elements.  It is a finite set.  The fixed
circles of the reflections, cut at $\mathcal V$, are open mirror
intervals or uncut mirror circles.  Take the $G$-orbits of the vertices
and of these mirror components; call their numbers $v$ and $e$.
There is one regular interior stratum in the quotient, since its
interior after removal of finitely many cone points is connected.
A vertex orbit may be removed once, a mirror-component orbit may
contain any nonnegative number of puncture orbits, and the regular
interior may contain any nonnegative number.  Thus the decoration is
\[
 (d_1,\ldots,d_v;c_1,\ldots,c_e;a),\qquad
 d_i\in\{0,1\},\ c_j,a\ge0,
\]
subject to the exact equation
\[
      \sum_i d_i|V_i|+\sum_j c_j\frac{|G|}{2}+a|G|=\nu.
                                                        \tag{B.27.5}
\]
Divide these finitely many decorations by the stratum permutations
of the affine normalizer computed above.  Counts are sufficient on
an open mirror interval or circle: an increasing interval map, or a
circle map of the prescribed degree, moves any two unlabelled finite
sets of the same size into each other.  The moves extend through
small strips in the quotient.  In the regular interior use disjoint
moving discs; its configuration space of finitely many unlabelled
points is connected.  Such quotient isotopies lift through the finite
orbifold covering.  They preserve all cone and corner charts and
therefore give equivariant homeomorphisms upstairs.  Conversely a
conjugacy preserves the displayed orbit sizes and the incidence of
these strata.  This proves completeness and injectivity of the
decorated classification.

All computations use exact integers on
$Q=(\frac1{144}\Z^2)/\Z^2$.  We record why this finite mesh loses no
component or conjugacy.  Coordinates of $T$ have denominator dividing
its order $k\le4$.  For the rotations of orders $2,3,4,6$, the inverses
of $I-r$ have denominators $2,3,2,1$, respectively.  Thus all solutions
needed for (B.27.2), all normalized origin changes, and all isolated
fixed points have denominators dividing $2k$ or $3k$.
The representatives (B.27.3) have denominator dividing $2k$.
For a reflection normalizer, the origin equation is a rank-one
integral equation with nonzero Smith invariant $1$ or $2$; every
connected solution component has a representative with denominator
dividing $4k$.  The integer $144$ is divisible by each of these
integers for $1\le k\le4$.  The code separately checks the inverse
matrix and particular-solution denominators for every positive affine
element.  On each reflection circle, moving by its primitive integral
direction divided by $144$ connects consecutive sampled points.
Every isolated intersection and every selected puncture is retained
as a vertex.  This graph is a subdivision of the actual fixed-circle
arrangement, so its components after removing punctures are exactly
the ovals and lines of the action.  The mesh has enough unused
representatives on each eligible stratum to place the at most four
punctures; the executed program explicitly verifies the prescribed
orbit counts and total cardinality for every decoration.

For each unremoved vertex $p$, enumerate all subgroups of its finite
cyclic or dihedral stabilizer that contain a positive nonidentity
element.  Each gives the section class with fixed point $p$.
For each reflection enumerate the components of its fixed circles
after deleting the punctures.  B.3--B.4 identify exactly these data
with all nontrivial partial-section classes, and identify their
centralizers as $1$ at points and lines and $\Z$ at ovals.
Containment of image subgroups and reverse containment of their fixed
components give all extension relations.  Applying every affine
matrix computes the group action on this poset and hence $\Phi$,
its orbit decomposition and every stabilizer.  At a positive fixed
point the local matrix is its actual linear part.  Its trace
$-2,-1,0,1$ gives order $2,3,4,6$, and the sign of
$\det(e_1,Ae_1)$ gives the sign of the angle except for a half-turn.
Reflections have multiplier $1/2$.  The program checks the order
and power identities on every element of every section.  It retains
rather than deletes the reflection section classes with several
maximal extensions.

Finally a nonidentity finite affine map acting trivially in the
punctured-torus mapping class group would have trivial homology
matrix and hence be a nonidentity translation.  It would also fix
each puncture, contradicting freeness of translations.  Thus each
action gives a faithful finite mapping class subgroup.  The realization
and uniqueness statements C.5--C.6 identify their conjugacy classes
with the topological action classes just enumerated.  The computation
returns the four counts in the statement.  For $\nu=1$, the multiset
of group orders, positive orders, full splitting counts and partial
splitting counts agrees, class for class, with the independent
amalgam calculation of B.25.  A first complete run with denominator
$72$ and the final run with denominator $144$ also give the same
four counts; the proof of completeness uses the latter mesh.
\end{proof}

\begin{claim}{B.28}{Proved}{Complete executed extended genus-two catalogue}
The extended mapping class group of the closed orientable genus-two
surface has exactly $92$ conjugacy classes of finite subgroups.  Of
these, $57$ contain the hyperelliptic involution and $35$ do not;
$21$ classes preserve orientation.  The files
\texttt{type\_2\_0\_extended\_catalogue.csv} and
\texttt{type\_2\_0\_extended\_sections.json} specify every group by
an exact finite permutation representation and give every full and
partial splitting, its multiplier and centralizer, all extension
incidences, $\Phi$, and its group orbits and stabilizers.
\end{claim}
\passage{\S{}15; folios 73--75; scans 148--151; \S{}15; folios 77--79; scans 155--160; \S{}18; folios 96--101; scans 193--203; \S{}33; folios 222--225; scans 440--445; Part B(f) of the task; no further manuscript assertion}

\begin{proof}
Let $\iota$ be the hyperelliptic mapping class.  The genus-two
Birman--Hilden homomorphism, with orientation reversal allowed, is
\[
 1\longrightarrow\langle\iota\rangle
 \longrightarrow\mathcal T_{2,0}
 \longrightarrow\mathcal T_{0,6}\longrightarrow1.       \tag{B.28.1}
\]
We recall its geometric content, to specify every comparison used here.
The involution is the deck map of the double cover of $S^2$ branched
at six points.  The covering is determined by the homomorphism from
the fundamental group of the six-times-punctured sphere to $C_2$
that sends each peripheral generator to the nonidentity element;
the product relation is respected because six is even.  Every
homeomorphism preserving the branch set preserves this homomorphism,
and hence has exactly two lifts, differing by $\iota$.
The genus-two hyperelliptic mapping class is central: the usual chain
of five nonseparating curves is invariant under the involution, their
Dehn twists generate the oriented genus-two mapping class group,
and a reflection of the standard branched-cover model commutes with
$\iota$ and supplies the negative coset.  The Birman--Hilden isotopy
property identifies the kernel in (B.28.1) with these two deck classes;
it also makes the lifting map onto.  This is the hyperelliptic
Birman--Hilden description used in B.25, in genus two instead of the
four-point pillowcase; see \cite[\S9.4.2, p.~257, with the proof in \S9.4.3]{A_FarbMargalit2012}.

For a finite $K\subset\mathcal T_{2,0}$ put
$E=K\langle\iota\rangle$ and let $H$ be its image in
$\mathcal T_{0,6}$.  The group $E$ is exactly the full inverse image
of $H$ in (B.28.1), and has order $2|H|$.  Thus B.26 supplies its
complete list of $57$ conjugacy classes.  A subgroup $K$ with this
fixed $E$ either equals $E$, or is a complement to
$\langle\iota\rangle$.  Every complement is normal in $E$ and is
uniquely the kernel of a homomorphism
\[
             \chi:E\longrightarrow\mathbf F_2,
             \qquad\chi(\iota)=1.                     \tag{B.28.2}
\]
Indeed $E=K\times\langle\iota\rangle$ for a complement, giving
this homomorphism; conversely the kernel of (B.28.2) has index two
and maps isomorphically onto $H$.  Because $\iota$ is central in the
whole mapping class group, conjugate $K$'s have conjugate associated
$E$'s.  For a fixed $E$, their further fusion is therefore exactly
the action of its normalizer on the characters (B.28.2).

Here are explicit presentations of all the finite groups $E$.
They are also used to calculate the sections, so the statement is
independent of the names assigned to abstract finite groups.  Let
$Q=S^2/H$ be one of the spherical orbifolds in B.26, with its
six-branch-point decoration.  Write $z=\iota$, $z^2=1$, and require
$z$ to commute with all generators.  An interior cone of order $m_j$
has a generator $x_j$ with
\[
                    x_j^{m_j}=z^{d_j},                 \tag{B.28.3}
\]
where $d_j=1$ exactly when its orbit is a branch orbit.  Let $a$ be
the number of regular interior branch orbits.  For a spherical
underlying quotient the remaining relation is
\[
                       \prod_jx_j=z^a.                 \tag{B.28.4}
\]
For a projective-plane quotient, with crosscap generator $u$, it is
$u^2\prod_jx_j=z^a$.

For a mirror disc, expand the cyclic corner list by inserting each
branch orbit on an open mirror edge as an extra entry $(m,d)=(1,1)$.
An original corner contributes $(m,d)$ with its original order $m$
and its branch flag $d$.  The distinction is retained as a color.
Let the resulting cyclic list be $(m_j,d_j)_{j=0}^{L-1}$.
There are reflections $R_0,\ldots,R_{L-1}$ and a boundary element
$e$, with
\[
 \begin{split}
 R_j^2&=1,\\
 \prod_i x_i\,e&=z^a,\\
 (R_jR_{j+1})^{m_j}&=z^{d_j},\\
 R_L&=eR_0e^{-1}.
 \end{split}
 \tag{B.28.5}
\]
When $L=0$ there is one reflection $R_0$ and the last relation is
$eR_0e^{-1}=R_0$.  A newly inserted mirror branch entry says exactly
$R_{j+1}=R_jz$.

These are presentations, rather than only relations satisfied by
some finite images.  To verify this, use the branched cover described
above.  Its full group of lifts of $H$ is a group of order $2|H|$.
A small positive circle about a branch point has deck monodromy $z$;
about an unbranched point it has trivial deck monodromy.  This gives
(B.28.3).  The product of the boundary loops gives (B.28.4), or its
crosscap or disc version.  At a branch point on a mirror, the local
map is $w\mapsto w^2$: the two reflections on the adjacent lifted
half-branches differ by $w\mapsto-w$.  This gives precisely the
inserted entry in (B.28.5).  At a corner the positive local cyclic
generator has its order doubled exactly when it is branched, giving
the other entries.  These local monodromies and $z$ generate the full
lift group.  Conversely quotienting the displayed presentation by
$z=1$ gives the original spherical orbifold presentation, which is
$H$ by the local-injectivity and Euler-characteristic certificate in
B.26.  Since $z$ is central of order at most two, the presented group
has at most $2|H|$ elements.  Its surjection onto the geometric lift
group proves equality and the presentation assertion.

We determine the entire normalizer action needed in (B.28.2), not
just conjugation inside the finite group.  A normalizing mapping
class descends through (B.28.1) and through the $H$-action to an
orbifold homeomorphism of $Q$ preserving its branch decoration.
Conversely any such orbifold homeomorphism lifts to $S^2$, because
$S^2$ is the universal orbifold cover of $Q$, and then to the double
branched cover, because the six peripheral mod-two monodromies are
preserved.  Thus its lifts give the whole normalizer action, with
inner automorphisms and the deck involution as the harmless lifting
ambiguities.  Both act trivially on characters to $\mathbf F_2$.

On a sphere, an orbifold homeomorphism sends each local cone
generator to a conjugate of a local generator with the same original
order and branch flag, or to its inverse when it reverses
orientation.  Characters ignore conjugation and inversion.  All
permutations among cones of the same pair $(m_j,d_j)$ are realized
by homeomorphisms of the sphere.  Point-pushing and pure braids
conjugate the local generators and therefore act trivially on these
characters.  The ordinary branch generators already equal the
central element $z$.  Consequently the exact action on (B.28.2)
is the permutation action on these colored cone coordinates.
The character equations are
\[
       m_j\chi(x_j)=d_j,\qquad
       \sum_j\chi(x_j)=a\pmod2.                       \tag{B.28.6}
\]
This includes, for example, the two elliptic complements in a
four-group lying over a sphere involution with no fixed branch
points: exchanging the two cone points exchanges those complements.
They must be fused, and the program does so.

On a mirror disc there is at most one interior cone.  The normalizer
preserves the cyclic order on its boundary, up to reversal.  It
therefore induces exactly the rotations and reversals of the
expanded colored corner list that preserve its colors.  Each such
symmetry is realized by a disc homeomorphism preserving the interior
decoration.  To see its action explicitly, extend the reflection
sequence by $R_{j+L}=eR_je^{-1}$.  A cyclic shift by $s$ sends
$R_j$ to $R_{j+s}$ and preserves $e$; a reversal sends $R_j$ to
$R_{s-j}$ and $e$ to $e^{-1}$.  The interior cone generator is sent
to a conjugate of itself or its inverse.  Modulo two these formulas
are literal permutations of the reflection coordinates and leave
the interior coordinate fixed.  They respect every relation in
(B.28.5).  Every remaining disc mapping class comes from moving the
interior marked points or from a boundary twist; its effect on local
generators is conjugation and hence is trivial on the characters.
The character equations are therefore
\[
 \begin{split}
 m_i\chi(x_i)&=d_i,\\
 \chi(e)&=a+\sum_i\chi(x_i),\\
 m_j\bigl(\chi(R_j)+\chi(R_{j+1})\bigr)&=d_j\pmod2.
 \end{split}
                                                        \tag{B.28.7}
\]
The program enumerates these bit vectors and divides by exactly the
colored dihedral permutations just proved.  For a projective-plane
quotient, the six-branch-point equation allows only the three cases
$P()$ with $a=3$, $P(2)$ with $(d,a)=(1,1)$, and $P(3)$ with
$(d,a)=(0,1)$.  Their full lifts are respectively $C_4,C_8,C_{12}$,
with $z$ the involution.  None has a character taking $z$ to one.
Thus no further complement fusion is needed in these cases.

Executing (B.28.6)--(B.28.7) gives $35$ character orbits, in addition
to the $57$ full groups.  The following order distribution is a
compact view of the complete $92$-row table:
\[
\begin{array}{c|rrrrrrrr}
 |K|&1&2&3&4&5&6&8&10\\\hline
 \#&1&7&1&13&1&11&14&3\\[2pt]
 |K|&12&16&20&24&32&48&96\\\hline
 \#&17&8&1&9&1&4&1
\end{array}                                             \tag{B.28.8}
\]
In particular its seven involution classes consist of the hyperelliptic
class, the elliptic positive class, and the five negative classes
predicted independently by B.16.

We finish by specifying the verification of all section data.
For each of the $57$ presentations, the program performs an exact
Todd--Coxeter enumeration and obtains the regular action on exactly
$2|H|$ cosets.  It evaluates every relator in these permutations and
checks that $z$ is nonidentity of order two.  In the quotient orbifold
for the genus-two action, an originally branched cone of order $m$
now has order $2m$, a regular branch orbit contributes a cone of
order two, and a branched corner has its order doubled.  Every
local cyclic or dihedral group injects into the computed finite
group.  The orbifold Euler characteristic multiplied by $2|H|$
is $-2$, and the covering subgroup is the kernel of a
monodromy homomorphism compatible with orientation.  Its torsion-free
covering surface is therefore orientable of genus two.
The program also checks that $z$ has exactly six section classes.

The cone and corner cosets, together with the lifted mirror coset
graph of B.26, give every fixed point and reflection oval.  Enumerating
subgroups of these local stabilizers and conjugating them by the
finite group is enough: every finite subgroup with a nonempty fixed
set is contained in such a local stabilizer.  The section multiplier
is read from the power of the corresponding positive local generator;
conjugating by a negative element reverses that angle.  For each
complement, retain exactly the section classes whose image subgroup
lies in its character kernel.  This restriction is exact because
the lift extension for the complement is the inverse image of that
subgroup in the full lift extension.  Recompute maximality in the
retained poset, intersect the section stabilizers with the complement,
and compute its actual orbits.  Thus no maximality or orbit count is
inferred merely by restricting the full-group summary.
All order and power identities for $\mu$ and all unique-maximal-extension
claims for sections containing a rotation are checked.  Finally the
$21$ positive rows agree with the independent primary-source
oriented genus-two catalogue used in B.33 on the complete multiset
of five invariants: group order, full sections, nontrivial partial
sections, $|\Phi|$, and $|\Phi/K|$.
\end{proof}

\begin{claim}{B.29}{Proved}{Complete executed punctured genus-two catalogues}
The extended mapping class groups of types $(2,1)$ and $(2,2)$ have
respectively $32$ and $121$ conjugacy classes of finite subgroups.
Their positive subgroups form $10$ and $23$ of these classes.
The files \texttt{type\_2\_\nu\_extended\_catalogue.csv} and
\texttt{type\_2\_\nu\_extended\_sections.json}, for $\nu=1,2$,
list every class with its exact group, puncture orbits and stabilizers,
all section classes and their incidences, $\Phi$, its group action
and stabilizers, centralizers, and all section and puncture multipliers.
\end{claim}
\passage{\S{}15; folios 73--75; scans 148--151; \S{}15; folios 77--79; scans 155--160; \S{}18; folios 96--101; scans 193--203; \S{}33; folios 222--225; scans 440--445; Part B(f) of the task; no further manuscript assertion}

\begin{proof}
Fill the punctures of a finite action on the genus-two surface.
The resulting finite subgroup in the closed mapping class group is
faithful: a finite homeomorphism isotopic to the identity on a closed
hyperbolic surface is the identity, by B.6.  Thus B.28 supplies the
complete list of closed actions that must be decorated.
For such an action $K$, keep the notation
$E=K\langle\iota\rangle$, $H=E/\langle\iota\rangle$, and
$Q=S^2/H$ from that proof.

The strata relevant to punctures are intrinsic to $K$.
Let $\mathcal V_K$ be the set of points whose stabilizer in $K$
contains a positive nonidentity element.  Take its $K$-orbits.
Cut the fixed circles of the reflections in $K$ at
$\mathcal V_K$ and take the $K$-orbits of the resulting open
intervals or uncut circles.  These are the mirror strata.
There is one connected regular interior in the quotient $S_2/K$.
A vertex orbit is selected zero or one time; a mirror stratum
receives any nonnegative number of puncture orbits; so does the
regular interior.  Their weights are, respectively,
\[
       |Kp|,\qquad |K|/2,\qquad |K|.                  \tag{B.29.1}
\]
The weighted sum must equal $\nu$.
Points fixed by $\iota$ but having no positive stabilizer in $K$
are ordinary points of these $K$-strata.  They introduce no extra
puncture label.  In particular punctures need not be preserved
by $\iota$ when $K$ is a complement in $E$.

We prove the normalizer permutations used to identify these
decorations.  Every mapping class normalizing $K$ normalizes $E$,
since $\iota$ is central in $\mathcal T_{2,0}$.  At the homeomorphism
level it may be represented by a normalizer of the chosen finite
action $E$ without changing its action on the intrinsic $K$-strata.
Indeed compare the two commuting involutions $\iota$ and
$h\iota h^{-1}$ after conjugating $K$ by $h$.  The finite actions of
the abstract extension $K\langle\iota\rangle$ have the same mapping
classes.  Uniqueness of realization in C.6 gives a conjugating
homeomorphism isotopic to the identity and commuting with $K$.
C.7 gives an isotopy through such commuting homeomorphisms.
It fixes the isolated $K$-fixed points individually and preserves
every reflection component.  Composing with it makes $h$ normalize
the chosen $E$ and leaves the required stratum permutation unchanged.
Thus it suffices, and is necessary, to compute the action of the
normalizer of $E$ that preserves $K$.

The full-group case $K=E$ is direct.  A cone in $Q$ has the color
consisting of its original order and its branch flag.
An ordinary branch orbit has its own color $(1,1)$.
The normalizer permits all permutations of identical interior
cone colors and exactly the dihedral symmetries of the expanded
colored boundary list of B.28.  Each is realized by a quotient
orbifold homeomorphism and hence lifts.  Pure interior braids
conjugate local generators, so they leave the corresponding
$E$-orbits unchanged.  These observations prove that this finite
permutation group is the whole normalizer action on the strata.

Suppose now $K=\ker\chi$ with $\chi(z)=1$.
For a local stabilizer $L\subset E$, its point orbit is $E/L$.
It is a single $K$-orbit when $\chi(L)=\mathbf F_2$.
When $\chi(L)=0$, it consists of two $K$-orbits, labelled
\[
                       b=\chi(g),\qquad gL\in E/L.    \tag{B.29.2}
\]
This is well defined precisely in the latter case.
It gives the exact orbit comparison, including the ambiguity in
choosing a lift of a quotient homeomorphism: composing with $z$
simultaneously changes every available sheet label $b$ to $b+1$.
The code includes this deck permutation.

An interior cone with generator $x$ contributes an intrinsic
$K$-vertex only when $K\cap\langle x\rangle$ is nontrivial.
For cone permutations, retain exactly the equal-color permutations
that also preserve the character values on their local generators.
On a sphere quotient, every split cone pair has an additional
independent interchange.  Here is the explicit lift proving that
assertion.  Choose a branch local generator $y$ with $\chi(y)=1$;
one exists because the branch decoration is nonempty.
An ordinary branch generator is $z$.  A branched old cone has
$y^m=z$ with $m$ odd whenever a character with $\chi(z)=1$ exists,
and therefore also has $\chi(y)=1$.
For a split cone generator $x$, we have $\chi(x)=0$.
A pure braid about these two quotient points acts by
\[
       x\longmapsto(xy)x(xy)^{-1},
       \qquad y\longmapsto xyx^{-1}.                 \tag{B.29.3}
\]
The other local generators are unchanged.
The character is preserved.  The conjugator for $x$ has character
one, so (B.29.2) interchanges its two $K$-orbits; the conjugator for
$y$ has character zero.  Thus it leaves the other split pairs
unchanged.  This proves the claimed independent interchanges.
Together with cone permutations they generate the whole stratum
action: pure braids generate the remaining mapping classes of a
punctured sphere, and their action on the two-element fibers can
only be these interchanges.

For a disc quotient there is at most one old interior cone.
An ordinary interior branch point permits exactly the same
interchange (B.29.3) of that split cone pair, supported away from
the boundary.  When there is no such interior branch point, the
disc has at most one interior marked cone; its pure braid group
introduces no additional permutation.  Boundary permutations are
computed as follows.  Let $L$ be the length of the expanded corner
list, write $R_{j+L}=eR_je^{-1}$, and put
$\epsilon=\chi(e)$.  A colored cyclic shift sends a corner index
$j$ to $s+j$ and an edge index $i$ to $s+i$.
A reversal sends them to $s-j-1$ and $s-i$, respectively.
For either image write the resulting integer as $qL+t$,
$0\le t<L$.  Its local generator is conjugated by $e^q$;
hence the sheet label changes by
\[
                          b\longmapsto b+q\epsilon.   \tag{B.29.4}
\]
Retain precisely the colored symmetries that preserve the reflection
character values, and include the global deck interchange.
These formulas give both vertex and mirror-component permutations.
For a cornerless boundary the same rule is its circular covering
with monodromy $\chi(e)$; the lifted mirror graph records whether
the two local sheet labels belong to one or two components.

There are no omitted boundary permutations.  An orbifold
homeomorphism of the disc preserves its cyclic boundary order
up to reversal.  After its colored dihedral permutation has been
removed, its remaining mapping class is an interior braid or a
boundary twist.  An interior braid has the local conjugators
already described and fixes the boundary.  A boundary twist has
the conjugator $e$, already represented by (B.29.4) and the deck
permutation.  Inner automorphisms coming from $K$ leave $K$-orbits
unchanged.  The projective-plane cases in B.28 have no complements;
their full-group decorations require only permutations of identical
interior cone colors.  Thus every case has an explicit complete
finite normalizer permutation group.

We spell out how this calculation handles smooth points of the
$K$-mirror locus that were corners for $E$.  The code starts with
the exact finite mirror coset graph of the full action $E$.
It finds the vertices with positive stabilizer in $K$, removes
only these vertices to form the open mirror strata, and computes
the fixed circles of each actual reflection in $K$.
Thus an $E$-corner that has only a reflection stabilizer in $K$
lies inside a single $K$-mirror stratum.  The sheet and boundary
formulas above are applied to an ordinary edge sample in each
such component.  They determine its whole component by
connectedness.  This is also a direct check that a puncture can
move through a noncritical Weierstrass point whenever the
$K$-stratum allows it.

Enumerate all vectors with the weights (B.29.1) summing to
$\nu=1$ or $2$, and divide by the proved finite permutation group.
These are precisely the topological action classes with punctures.
Indeed unlabelled finite subsets of an interval or circle with
the same cardinality are related by a homeomorphism of prescribed
degree, and the regular interior configuration space is connected.
Extend such moves by small discs and strips in the quotient and
lift the isotopies through the finite orbifold cover.  Conversely
an equivariant homeomorphism preserves the intrinsic strata,
their weights, and a stratum permutation in the normalizer just
computed.  C.5--C.6 then identify these action classes with the
finite-subgroup conjugacy classes in the punctured mapping class
group.

The executed program \texttt{catalogue\_genus2\_punctures.py}
returns $32$ and $121$ classes.  It places the punctures on the
exact coset graph, with one midpoint inserted in every edge before
any vertex is removed.  This subdivision guarantees that even
two punctures at formerly adjacent samples leave the intervening
open interval represented.  The program checks that an element
stabilizing an edge fixes both its endpoints, so the inserted
midpoint has exactly the edge stabilizer and creates no artificial
critical point.  It records regular puncture orbits by the free
$K$-set $K$, with trivial stabilizer.

All positive section classes from B.28 are retained exactly when
their fixed point has not been removed and their image subgroup
lies in $K$.  For every reflection, remove the punctures from its
actual fixed-circle graph and recompute the resulting ovals and
lines.  These operations give all partial-section classes by
B.3--B.4.  Their inclusion poset, maximal classes, $K$-action,
stabilizers and centralizers are then recomputed.
Local multipliers at punctures are retained from the positive
local generators or assigned $1/2$ to a reflection.
The positive portions contain $10$ and $23$ classes and agree
with the independent primary-source calculation in B.33 on the
full multiset of group order, full sections, partial sections,
$|\Phi|$, and $|\Phi/K|$.  The negative involutions provide
another independent check: there are $4$ classes with one
puncture and $11$ with two, exactly the decorations of the five
closed real types in B.16.
\end{proof}

\subsection{The full tube action and the strip subgroups}

\begin{claim}{B.30}{Proved}{Explicit outer words for the antipodally paired tube model}
For $g\ge2$, put $n=2g$ and label the equatorial holes by $j\in\Z/n$,
with $j$ glued by a tube to $j+g$. The following presentation and
substitutions compute the entire exterior action requested on scan 467:
\[
 \pi_g=\left\langle x_0,\ldots,x_{g-1},p_0,\ldots,p_{g-1}
 \ \middle|\ R=1\right\rangle,
 \qquad
 R=x_0\cdots x_{g-1}
   \prod_{j=0}^{g-1}(p_j^{-1}x_j^{-1}p_j).
\]
Extend the notation, with all subscripts read modulo $n$, by
\[
 \begin{gathered}
 x_{j+g}=p_j^{-1}x_j^{-1}p_j,\qquad p_{j+g}=p_j^{-1},\\
 u_j=(x_0\cdots x_j)^{-1},\quad u_{n-1}=1,\qquad
 q_j=u_j^{-1}p_j u_{j+g-1}.
 \end{gathered}
\]
Let $r=\varepsilon_0\varepsilon_1$ and $s=\varepsilon_1$ as in
B.20. Representatives of their outer automorphisms, and of $\tau$ and
$\rho_k$, are specified on the $2g$ displayed generators by
\[
\begin{array}{c|cc}
 &x_j&p_j\\ \hline
 r&x_{j+1}&p_{j+1}\\
 s&u_{-j}^{-1}x_{-j}u_{-j}&q_{-j}\\
 \tau&u_{j+g-1}^{-1}x_{j+g}^{-1}u_{j+g-1}&q_j^{-1}\\
 \rho_k&x_j&x_j^{\delta_{jk}}p_j.
\end{array}
\]
The map to $\operatorname{Out}(\pi_g)$ defined by these substitutions
is the geometric action of
$\Z^g\rtimes(D_{4g}\times C_2)$ of B.20. The finite factor acts by
\[
 r\rho_jr^{-1}=\rho_{j+1},\qquad
 s\rho_js^{-1}=\rho_{-j},\qquad
 \tau\rho_j\tau^{-1}=\rho_j^{-1},
\]
with indices in these three formulas taken modulo $g$.
\end{claim}
\passage{\S{}33; folios 234--236; scans 463--467}

\begin{proof}
We specify the base paths so that the formulas can be used as words,
rather than only as unbased pictures. Let $N,S$ be the north and south
poles in the sphere with the $n$ discs removed. Draw disjoint meridian
arcs $d_j:N\to S$ through the gap immediately following hole $j$.
The loop about that hole, based at $N$, is
$x_j=d_{j-1}d_j^{-1}$. Collapse the maximal-tree edge $d_{n-1}$ and
identify $N,S$ by that edge. The remaining arcs become
$u_j=(x_0\cdots x_j)^{-1}$, because the product telescopes to
$d_{n-1}d_j^{-1}$. In particular $\prod_{j=0}^{n-1}x_j=1$.

Let $p_j$ run from $N$ to the northern side of hole $j$, through the
tube to hole $j+g$, and back to $N$ on its northern side. The opposite
traversal is $p_{j+g}=p_j^{-1}$. Its boundary identification reads
$x_{j+g}=p_j^{-1}x_j^{-1}p_j$. Van Kampen therefore gives exactly the
presentation in the claim. The parallel path based at $S$, on the
southern side of the tube, is
$q_j=d_j^{-1}p_jd_{j+g-1}$ after inserting the northern paths, which
becomes the displayed formula on collapsing the tree edge.
These identities can also be taken as a two-vertex cellular description:
the sphere faces give the telescoping boundary relation and each tube
face gives the indicated boundary conjugation.

Rotation about the polar axis sends hole $j$ to $j+1$, preserves $N$,
and sends $p_j$ to $p_{j+1}$. This proves the $r$ row. The half-turn
$s$ about the equatorial axis through hole $0$ exchanges $N,S$ and
sends hole $j$ to $-j$. The southern boundary loop at hole $k$ is
$d_k^{-1}x_kd_k$, and the southern tube path is $q_k$, giving the
$s$ row. The antipodal map exchanges $N,S$, sends $j$ to $j+g$,
and reverses orientation. It therefore inverts the southern boundary
loop and reverses the corresponding tube path. The boundary-loop
formula written with $u_{j+g-1}$ equals that written with $u_{j+g}$,
since $u_{k-1}=x_ku_k$. This proves the $\tau$ row. Finally a positive
full twist of tube $k$ inserts $x_k$ into the crossing $p_k$ and
leaves all other chosen paths unchanged. This proves the $\rho_k$ row.
Each row comes from an actual homeomorphism, so it is an automorphism
of the stated group and has the prescribed orientation sign.

For $s$ and $\tau$ the chosen path from $N$ to the image of $N$ is
$d_{n-1}$. We consequently get actual automorphisms only after this
base-path choice. For example their squares are
\[
 s^2=\operatorname{Inn}(x_0),\qquad
 \tau^2=\operatorname{Inn}(x_0\cdots x_{g-1}),
\]
whereas $r^{2g}=1$ as an actual automorphism. These formulas follow
by concatenating the base path with its image: $s(d_{n-1})=d_0^{-1}$
and $\tau(d_{n-1})=d_{g-1}^{-1}$. Thus the displayed group relation
$s^2=\tau^2=1$ holds in the outer group with precisely the stated
inner ambiguity. The geometric symmetries satisfy $srs=r^{-1}$ and
commute with the antipodal involution on the surface. They preserve
and permute the tube collars; a positive symmetry preserves twist
sign and the antipodal map reverses it. This proves every relation
of the asserted semidirect product in the outer group. Changing the
maximal-tree edge conjugates each based comparison by its corresponding
base-path word and leaves this exterior action unchanged.

The accompanying programs check all relator images for $2\le g\le8$;
$r,s$ carry $R$ to a conjugate of $R$, and $\tau$ carries it to a
conjugate of $R^{-1}$. They additionally check all the finite-factor
and twist-conjugacy relations in the surface group, with the actual
inner conjugators above. For instance $srsr=1$ and
$\tau r\tau r^{-1}=\operatorname{Inn}(x_0\cdots x_g)$.
The word checker uses the symmetrized relator: its largest common
piece has length one, so the length-$4g$ presentation satisfies
$C'(1/6)$ for $g\ge2$, and the Dehn reduction is exact. These finite
certificates supplement the cellular proof for arbitrary $g$.
\end{proof}

\begin{claim}{B.31}{Proved}{The eight symmetries of one tube}
The stabilizer of tube $T_0$ in $D_{4g}\times\langle\tau\rangle$ is
$\langle r^g,s,\tau\rangle\cong C_2^3$. There are coordinates
$(\theta,t)\in(\R/2\pi\Z)\times[-1,1]$ on that tube in which its
three generators act as
\[
 r^g(\theta,t)=(-\theta,-t),\qquad
 s(\theta,t)=(\theta+\pi,t),\qquad
 \tau(\theta,t)=(\theta+\pi,-t).
\]
These formulas agree with the outer action in B.30 and with the
manuscript's free antipodal extension, after the explicit coordinate
change $\theta_{\rm new}=\theta_{\rm old}+\pi t/2$.
\end{claim}
\passage{\S{}33; folios 229--230; scans 453--454; \S{}33; folios 233--236; scans 461--467}

\begin{proof}
The stabilizer of the unordered pair $\{0,g\}$ under the polygon's
dihedral group is $\{1,r^g,s,r^gs\}$: a stabilizing rotation is either
$1$ or $r^g$, and every stabilizing reflection is one of these followed
by $s$. The independent central antipodal factor doubles it. In a
symmetric tube, $r^g$ interchanges its ends and reverses its angular
direction, while $s$ preserves both ends and rotates the cross section
by $\pi$. The displayed representatives commute, are involutions,
and have eight distinct combinations; the antipodal representative
is free. They extend the same boundary maps as the corresponding
sphere symmetries and hence give the desired tube extension.

The formula of scan 453, retained through scans 461--462, is
$\tau_{\rm old}(\theta,t)=(\theta+\pi+\pi t,-t)$.
Conjugating it by the stated change gives exactly
$(\theta+\pi,-t)$; this also transports both boundary identifications.
Thus it is an explicit isomorphism of the two actions. The full twist
is $(\theta,t)\mapsto(\theta+\pi+\pi t,t)$ in either coordinate.
Conjugation by the two positive generators preserves its class, while
conjugation by the negative generator inverts its class, agreeing with
B.30. All eight maps and their products are checked symbolically in
\texttt{code/verify\_tube\_local.py}.
\end{proof}

\begin{claim}{B.32}{Proved}{Relations for arbitrary strips and a uniform obstruction to generation}
Choose for each $c_j$ one simple transversal $d_j$, using a strip in
the complement of the other tubes, so that
$i(c_i,d_j)=\delta_{ij}$. Let $W_j$ be a regular neighbourhood of
$c_j\cup d_j$, a one-holed torus, and put
$A_j=T_{c_j}=\rho_j$, $B_j=T_{d_j}$. Its supported subgroup is the
central extension $\langle A_j,B_j\rangle\cong B_3$ of B.21, with
quotient $SL_2(\Z)$. The exact pairwise relations are
\[
 [A_i,A_j]=1,\quad [A_i,B_j]=1\ (i\ne j),\qquad
 A_jB_jA_j=B_jA_jB_j.
\]
For distinct $i,j$, the two supported subgroups commute precisely when
$i(d_i,d_j)=0$. For intersection one, $B_i,B_j$ obey the braid
relation and do not commute; for intersection at least two they freely
generate a free group of rank two. For every such choice of one strip
per tube and every $g\ge2$, these $g$ supported subgroups, together
with $H$, generate a proper subgroup of $\mathcal T_g^+$.
\end{claim}
\passage{\S{}33; folios 236--237; scans 467--470}

\begin{proof}
Each $W_j$ has genus one and one boundary component because the two
core curves meet transversely once. The extension-by-identity map and
its faithful boundary-twist kernel are the explicit maps of B.21.
The disjointness conditions give the first two commutation relations,
and the calculation on its two free generators in B.21 proves the
braid relation. For any two essential simple curves $a,b$, their
twists commute exactly when they are disjoint: conjugation sends
$T_b$ to $T_{T_a(b)}$, and
$i(T_a(b),b)=i(a,b)^2$. Hence intersecting curves cannot commute.
If $i(a,b)=1$, their regular neighbourhood is a one-holed torus, so
the same braid calculation applies.

For $k=i(a,b)\ge2$ the freedom assertion has an elementary ping-pong
proof. On isotopy classes of essential curves set
\[
 X_a=\{x:i(x,a)<i(x,b)\},\qquad
 X_b=\{x:i(x,b)<i(x,a)\}.
\]
Both are nonempty and disjoint. Counting crossings in a twist collar
and cancelling only crossings already present outside it gives
\[
 i(T_a^m(x),b)\ge |m|\,i(x,a)k-i(x,b).
\]
For $m\ne0$ and $x\in X_b$ this is strictly greater than
$i(x,a)=i(T_a^m(x),a)$, so $T_a^m(X_b)\subset X_a$.
The symmetric inclusion holds for powers of $T_b$. The ping-pong
argument on a reduced alternating word proves that it is nontrivial:
apply it to the appropriate one of $X_a,X_b$, inserting one further
nonzero power from the opposite factor when the first and last
factors agree. Thus $\langle T_a,T_b\rangle\cong\Z*\Z$.
These statements compute the required pairwise relations for arbitrary
intersections of the selected strips; no disjoint-strip convention is
being used to conclude the general generation assertion.

For that assertion work on $V=H_1(X_g;\F_2)$ with its nonsingular
intersection pairing. The $2g$ classes $[c_j],[d_j]$ form a basis:
pairing a linear relation first with the $[c_j]$ kills all coefficients
of the $[d_j]$, and pairing next with the $[d_j]$ kills the rest.
There is a unique quadratic refinement $q:V\to\F_2$ with
$q([c_j])=q([d_j])=1$. Explicitly, for this ordered basis $(v_i)$,
\[
 q\Bigl(\sum_i a_iv_i\Bigr)
 =\sum_i a_i+\sum_{i<j}a_ia_j\langle v_i,v_j\rangle.
\]
It obeys $q(v+w)=q(v)+q(w)+\langle v,w\rangle$ by expansion.
A twist about a curve of class $v$ acts by
$x\mapsto x+\langle x,v\rangle v$; substitution gives
\[
 q(T_vx)-q(x)=\langle x,v\rangle\bigl(q(v)+1\bigr).
\]
Every generator $A_j,B_j$ therefore preserves $q$, without any
restriction on the intersections among the $d_j$.

Join $c_1$ and $c_2$ by an embedded arc in their connected complement
and form a band sum $e$. This is a simple curve of nonzero class
$[e]=[c_1]+[c_2]$, so it is nonseparating. Since the two cores are
disjoint, $q([e])=0$. Taking $x=[d_1]$ gives
$\langle x,[e]\rangle=1$, so $T_e$ fails to preserve $q$.
Consequently $T_e$ belongs to $\mathcal T_g^+$ and lies outside the
group generated by the strip subgroups. This proves properness for
all choices under discussion. The quadratic calculation has been
exhaustively checked for every possible matrix of transversal
intersections modulo two in genera two, three and four; the general
identity is the displayed algebraic proof.

The rank-one comparison is thus the specified central quotient
$B_3\to SL_2(\Z)$ and its action on the homology of $W_j$, together
with the above braid and intersection relations. The subgroup they
generate fixes the explicit quadratic refinement. This supplies the
precise algebraic maps and the explicit obstruction governing the
proposed generation in scan 470.
\end{proof}

\begin{claim}{B.33}{Proved}{Executed positive catalogues in genera two and three}
For $(g,\nu)=(2,0),(2,1),(2,2),(3,0)$, the numbers of conjugacy
classes of finite orientation-preserving subgroups, allowing conjugacy
in the full extended mapping class group, are respectively
\[
                         21,\qquad10,\qquad23,\qquad55.
\]
The trivial subgroup is included. The files
\texttt{oriented\_higher\_genus\_catalogue.csv} and
\texttt{oriented\_higher\_genus\_sections.json} give each class,
its actual finite permutation group, every subgroup's number of
sections, all nontrivial partial-section classes and their inclusion
relations, the $G$-set $\Phi$, its stabilizers, and every value of
$\mu$ on every section. Punctures are unlabelled. These assertions
concern positive subgroups; they do not enumerate subgroups containing
orientation-reversing elements.
\end{claim}
\passage{\S{}15; folios73--75; scans148--151; \S{}17; folios77--79verso; scans155--160; \S{}18; folios96--101; scans193--203; Part B(f) of the task; positive higher-genus cases}

\begin{proof}
We specify the classification input and the computation using it.
Karab\'a\v{s}--Nedela--Skyvov\'a
\cite[Theorem 2.2, Corollary 2.3, Theorem 4.1 and \S8]{R-KNS}
identify the positive topological classes with the orbits of smooth
epimorphisms from the Fuchsian group
\[
 \Delta=\langle x_1,\ldots,x_r,a_1,b_1,\ldots,a_h,b_h\mid
 x_i^{m_i}=1,\ [a_1,b_1]\cdots[a_h,b_h]x_1\cdots x_r=1\rangle
\]
under target automorphisms and orientation-preserving orbifold
automorphisms. Here smooth means that every image of $x_i$ has
order exactly $m_i$. Their \S8 gives a complete genus-two list
of 21 classes and a genus-three list of 55 classes in 49
group/signature pairs, including the trivial actions. We opened
the accompanying data \cite{R-KNSData}, downloaded the two original
JSON files, and retained their attribution and SHA-256 digests.
The attached program reads those exact files. It verifies that
each generating vector generates the group of the stated order,
each elliptic generator has its specified order, every presentation
relation holds, and
\[
 |G|\left(2-2h-\sum_i(1-1/m_i)\right)=2-2g.
\]
The conversion from the source's right-acting permutations to the
left-acting permutations in our code is the isomorphism sending
each permutation to its inverse. For a handle word the two
commutator conventions are compared explicitly; replacing both
handle generators by their inverses converts $a^{-1}b^{-1}ab$
to the convention $aba^{-1}b^{-1}$. The files retain the actual
generating vector, so this convention does not affect a recorded
elliptic rotation.

The local-order checks imply that the kernel of each epimorphism
is torsion-free: every finite-order element of a Fuchsian group
is conjugate to a power of an elliptic generator. The covering
formula then identifies its quotient surface as genus $g$.
Conversely a positive finite action determines this epimorphism
by lifting paths in the quotient. The two constructions are
inverse modulo the stated source and target automorphisms.
Finally C.5--C.6 identify these action classes with conjugacy
classes of finite mapping-class subgroups. This supplies the
exact reduction from the classification input to the groups
asked for here.

We also recompute the equivalence relation whenever $h=0$,
rather than using only the published class count. For an ordered
generating vector $(u_1,\ldots,u_r)$ with product one, the braid
generator acts by
\[
 \sigma_i:(u_i,u_{i+1})\longmapsto
               (u_i u_{i+1}u_i^{-1},u_i).
\]
Use adjacent such generators when $m_i=m_{i+1}$ and all pure
braid generators, obtained by moving the $j$th entry next to the
$i$th, applying $\sigma_i^2$, and reversing the moves.
Theorem 4.1 of the cited paper proves that these transformations
induce the full orientation-preserving orbifold automorphism
action after inner automorphisms have been divided out. A
generating vector is compared modulo $\Aut(G)$ by its rooted,
edge-coloured Cayley graph: breadth-first enumeration with the
ordered generators gives a canonical transition table. Equality
of tables gives the explicit bijection of vertices preserving
every labelled multiplication edge; induction on words proves
that it is the unique target-group isomorphism carrying the
generating vectors to one another. Thus the canonicalization
does not approximate group automorphisms.

To allow conjugacy by negative homeomorphisms, we add the
reflection of the marked sphere. Its effect on the tuple is
\[
 (u_i)_{i=1}^r\longmapsto
 \left(p_{i-1}u_i^{-1}p_{i-1}^{-1}\right)_{i=1}^r,
 \qquad p_{i-1}=u_1\cdots u_{i-1}.
\]
The product of the new tuple is $(u_1\cdots u_r)^{-1}=1$,
as successive cancellation shows. The program computes the
orbits both before and after adding this transformation; it
recovers the published planar counts and finds no fusion.
All nonplanar group/signature pairs in these two genera have
only one positive class in the input classification. Reversing
orientation preserves the group and signature and hence fixes
that class. Consequently the positive counts remain 21 and 55
under conjugacy in the extended group.

For genus two we add all puncture configurations with
$\nu\leq2$. An invariant finite puncture set consists of selected
elliptic orbits and a number $a$ of ordinary orbits. Write
$c_i\in\{0,1\}$ for selection of the $i$th elliptic orbit.
The exact equation is
\[
                  \sum_i c_i\frac{|G|}{m_i}+a|G|=\nu.
                                                               \tag{B.33.1}
\]
The program enumerates every solution. It carries the bits
$c_i$ along under each braid move, exchanging the two bits
when the corresponding two elliptic entries are exchanged.
It then takes orbits of these decorated generating vectors
under the same coloured braids, target automorphisms, and
reflection. These are exactly the decorated topological
classes: an equivariant homeomorphism descends to the quotient
orbifold, and its action on the covering is the compatibility
of its induced automorphism with the epimorphism just described.
Conversely that compatibility lifts the quotient homeomorphism
to an equivariant homeomorphism. Ordinary unlabelled punctures
can be moved to any other such configuration through paths in
the regular part of the quotient; point-pushing the paths is
isotopic to the identity after these additional points are
forgotten and so preserves the original covering kernel.
Thus $a$ records all their remaining data.

There are just two nonplanar genus-two entries. For the trivial
group, there is one configuration for each $\nu$. The other
has group $C_2$ and signature $(1;2,2)$. Its two elliptic
images are the unique involution, and exchanging their small
discs preserves the monodromy after an adjustment of the two
handle paths. The four possible pairs of handle images in
$C_2^2$ are connected by point-pushing either branch point
around a handle: crossing the branch cut adds its nonzero
monodromy to the transverse handle image. These two pushes
change either handle coordinate independently. Hence the
exchange lifts and only the number of selected branch orbits
matters. These elementary moves account for all four vectors
in the supplied data and prove the nonplanar decoration step
used by the code. Equation (B.33.1) and the decorated orbits
give 10 and 23 classes at one and two punctures.

We now derive every section entry from the checked finite
groups. A retained elliptic orbit is the left $G$-set
$G/\langle u_i\rangle$. Its point at the coset
$v\langle u_i\rangle$ has stabilizer
$v\langle u_i\rangle v^{-1}$ and positive generator
$v u_i v^{-1}$. B.3 gives exactly one kernel-conjugacy class
of sections for every point fixed by a specified nontrivial
subgroup $K\leq G$. Such a subgroup is cyclic because it is
contained in the cyclic point stabilizer. Conversely every
such containment gives a section. Thus the exact count is
\[
 \#\Sec(E_K\to K)/\pi
  =\sum_{i:c_i=0}
       \#\{v\langle u_i\rangle:
                    K\leq v\langle u_i\rangle v^{-1}\}.
                                                               \tag{B.33.2}
\]
The formula is used for every subgroup, which the program
enumerates by saturation under adjoining each remaining group
element. It checks that no noncyclic subgroup has a section.
The trivial subgroup has one section, and it is not included
in the nontrivial partial-section totals.

The maximal partial sections are precisely these individual
elliptic points, with $G$ acting by left multiplication on the
cosets. A partial section lies below another exactly when
they have the same point and its subgroup is contained in the
other subgroup. This determines the whole inclusion relation
and proves uniqueness of a maximal extension in the positive
case. Each nontrivial section has centralizer one by B.3.
The total number of nontrivial partial-section classes is
independently checked against
\[
             \sum_{i:c_i=0}\frac{|G|}{m_i}\bigl(\tau(m_i)-1\bigr),
\]
where $\tau(m)$ is the number of positive divisors of $m$.
Finally, for $h=v u_i^b v^{-1}$ the value of $\mu$ at this
section is $b/m_i\pmod\Z$. This lists every power of every
local generator and is exactly the intrinsic multiplier
identified in B.12. Puncture stabilizers are the removed
cyclic groups or the trivial group for ordinary orbits.
All these data are computed with integer permutations and
rational fractions. The output has 109 individual action
rows and 59 independent planar braid-orbit checks.
\end{proof}

\begin{claim}{B.34}{Proved}{The complete genus-three subcatalogue containing a hyperelliptic involution}
There are exactly $79$ conjugacy classes of finite subgroups of the
extended mapping class group of a closed orientable genus-three surface
that contain a hyperelliptic involution.  Exactly $18$ preserve orientation
and $61$ contain orientation-reversing elements.  The latter $61$ classes
are listed below and are specified, together with every section and its
extension incidences, in
\texttt{type\_3\_0\_hyperelliptic\_mixed\_catalogue.csv} and
\texttt{type\_3\_0\_hyperelliptic\_sections.json}.
\end{claim}
\passage{\S{}18; folios 96--101; scans 193--203; \S{}33; folios 222--225; scans 440--445; Part B(f) of the task; the genus-three hyperelliptic subcatalogue is an added result, not a further manuscript assertion}

\begin{proof}
Fix a hyperelliptic involution $z$ and its branched quotient
$p:S_3\longrightarrow S^2$, with branch set $P$ of size eight.
On the unbranched complements this is the connected double cover
classified by the homomorphism
\[
 \pi_1(S^2\setminus P)\longrightarrow C_2
 \quad\hbox{sending each of the eight meridians to }1.
                                                        \tag{B.34.1}
\]
Their sum is zero modulo two.  Every homeomorphism of the sphere
preserving $P$ preserves (B.34.1), and therefore has exactly two lifts.
The Birman--Hilden isotopy property gives the exact sequence
\[
 1\longrightarrow\langle z\rangle\longrightarrow
 C_{\mathcal T_{3,0}}(z)\longrightarrow\mathcal T_{0,8}
 \longrightarrow1.                                    \tag{B.34.2}
\]
Here the left group is the centralizer of this specified involution;
the map is induced by the quotient $p$.  The lifting and isotopy argument
is precisely the fully ramified hyperelliptic case of
\cite[\S9.4]{A_FarbMargalit2012}.  Allowing orientation reversal adds the
negative lifts of the same sphere homeomorphisms, with the same kernel.

The involution $z$ is the unique hyperelliptic involution inside any
finite subgroup $K$ containing it.  To verify the assertion directly,
realize $K$ by isometries and use the resulting invariant complex
structure.  A hyperelliptic involution acts as $-1$ on $H^1(S_3;\mathbb R)$:
its invariant subspace is the pullback of $H^1(S^2;\mathbb R)=0$ under
its quotient, as follows by averaging cochains, or by the transfer
outside the branch points and filling their meridians.  If $z'$ were
another, the finite-order holomorphic map $zz'$ would act trivially on
$H^1$.  A nonidentity such map has isolated fixed points, each of
positive local Lefschetz index; a fixed-point-free map has Lefschetz
number zero.  But its Lefschetz number from homology is $2-6=-4$.
Thus $zz'=1$, proving uniqueness.
Consequently every element of $K$ centralizes $z$, and any conjugacy
between two such finite groups maps their hyperelliptic involutions
to one another.  Once $z$ is fixed, (B.34.2) therefore gives a bijection
\[
 \left\{\begin{array}{c}\text{finite subgroup conjugacy classes}\\
 \text{in }\mathcal T_{3,0}\text{ containing }z\end{array}\right\}
 \longleftrightarrow
 \left\{\begin{array}{c}\text{finite subgroup conjugacy classes}\\
 \text{in }\mathcal T_{0,8}\end{array}\right\}.
                                                        \tag{B.34.3}
\]
The inverse takes the full preimage, of order twice the sphere group.
The inclusion of every kernel element follows from the assumption
$z\in K$.  This also proves that no additional fusion occurs when the
chosen hyperelliptic involution is forgotten.

For completeness, B.26's decorated spherical-orbifold enumeration
applies verbatim to eight marked points: a selected cone orbit has
size $|H|/m$, a selected corner orbit has size $|H|/(2m)$, an orbit
on an open mirror edge has size $|H|/2$, and a regular orbit has size
$|H|$.  Their sum is eight.  The cyclic and dihedral rotation orders
are at most eight, since every non-polar point orbit has at least that
many points.  The tetrahedral and octahedral cases are included.
An icosahedral orbit has at least twelve points, so cannot occur.
The spherical classification and the homeomorphism classification
of these decorations in B.26 prove exhaustiveness and irredundancy.
Executing the integer weight equation gives $79$ decorated quotients,
$18$ spherical and $61$ with a negative element.

For each decorated quotient, the full lift is given by the exact
central extension presentation (B.28.3)--(B.28.5), now with eight
branch points.  Its presentation proof uses (B.34.1) in exactly the
same way: the displayed group surjects onto the full geometric lift
group of order $2|H|$, while its quotient by the central element $z$
is $H$, so it has at most $2|H|$ elements.  The two groups are
therefore isomorphic by that surjection.  For a projective-plane
quotient its generators include both $z$ and the crosscap element;
their relation is $u^2\prod x_i=z^a$.  In particular no cyclicity
assumption from the six-branch-point case is imposed.

The program \texttt{catalogue\_genus3\_hyperelliptic.py} was executed.
For every presentation it computes an exact regular permutation
representation of degree $2|H|$, verifies every defining relator,
checks that $z$ has order two, and checks injectivity on every local
cyclic or dihedral stabilizer.  The resulting orbifold Euler
characteristic is $-4/|K|$, so its oriented surface kernel has genus
three.  Exactly eight point sections have image $\langle z\rangle$,
providing an independent branch-number check.
The finite coset incidence graph of B.26 then computes every point
and oval fixed component and every partial section.  Each point
section has trivial centralizer in the surface group, each oval
section has infinite-cyclic centralizer, and the rotation multipliers,
extensions, maximal sections $\Phi$, group orbits and stabilizers
are recorded explicitly in the JSON file.  The calculation retains
all intersections of reflection ovals at dihedral stabilizers.
\end{proof}

In the following table, $b$ is both the G3H class suffix and the decorated sphere-eight row number
in \texttt{type\_3\_0\_hyperelliptic\_catalogue.csv}; $N$ is the group
order; $p$ counts all nontrivial partial splitting classes; and $m,o$
are respectively $|\Phi|$ and $|K\backslash\Phi|$.  The CSV gives the
cone, corner, mirror and ordinary branch decorations of every row,
and the JSON gives its finite group generators and all relators.
\begin{longtable}{r|r|r|r|r}
b&N&p&m&o\\\hline
\endhead
2&4&8&8&4\\
3&4&10&10&5\\
4&4&12&8&5\\
5&4&16&8&6\\
6&4&20&8&7\\
7&4&24&8&8\\
11&8&12&12&3\\
12&8&14&14&4\\
13&8&20&12&4\\
14&8&28&12&5\\
15&8&14&8&3\\
16&8&22&8&4\\
17&8&20&12&4\\
18&8&26&12&5\\
19&8&34&12&6\\
20&8&24&12&5\\
21&8&32&12&6\\
22&8&32&12&6\\
23&8&24&8&4\\
24&8&32&8&5\\
25&8&32&8&5\\
29&12&12&8&2\\
30&12&14&10&3\\
31&12&24&8&3\\
32&12&30&8&3\\
33&12&42&8&4\\
34&12&42&8&4\\
39&16&16&12&2\\
40&16&28&20&3\\
41&16&40&20&4\\
42&16&28&12&3\\
43&16&18&14&3\\
44&16&32&12&3\\
45&16&38&20&4\\
46&16&52&20&5\\
47&16&50&20&5\\
48&16&48&12&4\\
49&16&48&12&4\\
50&16&36&12&3\\
51&16&48&12&4\\
52&16&44&12&4\\
56&24&54&20&3\\
57&24&36&8&2\\
58&24&28&8&2\\
59&24&68&20&4\\
60&24&66&20&4\\
61&24&66&8&3\\
62&24&72&8&3\\
65&28&45&10&3\\
66&32&60&28&3\\
67&32&44&12&2\\
68&32&36&12&2\\
69&32&82&28&4\\
70&32&84&28&4\\
71&32&76&12&3\\
72&32&72&12&3\\
74&48&54&20&2\\
75&48&120&20&3\\
76&48&120&20&3\\
78&64&136&28&3\\
79&96&210&44&3\\
\end{longtable}
The mixed-class order distribution is
\[
\begin{array}{c|rrrrrrrrrr}
 |K|&4&8&12&16&24&28&32&48&64&96\\\hline
 \#&6&15&6&14&7&1&7&3&1&1.
\end{array}                                             \tag{B.34.4}
\]
The scope of (B.34.3) is the condition that the subgroup itself
contain a hyperelliptic involution.  A finite subgroup that only
centralizes an involution outside itself is represented by a
complement in a lift group; its conjugacy after forgetting the
external involution requires an additional calculation.  No such
complement is counted among these $79$ classes.

\begin{claim}{B.36}{Proved}{Complete exceptional catalogues and their realization fibers}
For the four exceptional types, the numbers of finite-subgroup
conjugacy classes in the extended mapping class group, and the
numbers of finite action classes inducing these subgroups
isomorphically, are as follows:
\[
\begin{array}{c|c|r|r}
(g,\nu)&\mathcal T_{g,\nu}&
 \text{finite subgroup classes}&\text{realized action classes}\\\hline
(0,0)&C_2&2&3\\
(0,1)&C_2&2&2\\
(0,2)&C_2^2&5&7\\
(1,0)&\mathrm{GL}_2(\Z)&13&17
\end{array}                                             \tag{B.36.1}
\]
The executed files
\path{exceptional_finite_subgroups_and_lifts.csv} and
\path{exceptional_sections.json} give all $29$ action rows over
these $22$ mapping class subgroup rows.  They retain the extension,
all full and partial section classes, $\Phi$ and its stabilizers,
centralizers, and every defined local multiplier.
The request in B(f) for one set of section data attached only to a
finite mapping class subgroup is false in these types: the displayed
realization fibers are required.  For the antipodal action on $S^2$,
the algebraic section exists and the fixed-point locus is empty;
there is no fixed-point rotation multiplier for that section.
\end{claim}
\passage{\S{}15; folios 68--75; scans 137--151; \S{}17; folios 88--90; scans 179--181; \S{}19; folios 111--115; scans 222--230; Part B(f) of the task; no further manuscript assertion}

\begin{proof}
We first identify the groups in (B.36.1) and the comparison being
classified.  Orientation is the component invariant for
$\operatorname{Homeo}(S^2)$ and for
$\operatorname{Homeo}(S^2,\{\infty\})$.
An orientation-preserving homeomorphism is isotopic to the identity:
fix a disc and its complementary disc by the Alexander isotopy,
with the marked point kept fixed in the second case.
For the annulus, the two component invariants are the sign
$\epsilon$ on its fundamental group and the sign $\eta$ of its
permutation of the two ends.  Their four representatives are
\[
       (x,t)\longmapsto(\epsilon x,\eta t),
       \qquad x\in\R/\Z,\quad t\in\R.                  \tag{B.36.2}
\]
The kernel of these invariants is connected: compactify the two
ends, isotope an arc joining them to a fixed arc, and apply the
Alexander isotopy on the complementary disc.  Winding of the arc
is absorbed by rotating a puncture neighborhood, since its end
is a puncture and has no prescribed boundary parametrization.
For the torus, homology gives
$\pi_0\operatorname{Homeo}(T^2)=\mathrm{GL}_2(\Z)$.
Surjectivity is given by linear maps.  A homologically trivial
homeomorphism is first isotoped to fix the origin, then the two
basis curves, and then their complementary disc; translations
lie in the identity component.  B.25 gives its thirteen finite
subgroup classes.

In this statement a realized action over a finite subgroup
$H\subset\pi_0\operatorname{Homeo}(U)$ means a finite subgroup
$G\subset\operatorname{Homeo}(U)$ for which the projection
$G\to\pi_0\operatorname{Homeo}(U)$ is an isomorphism onto a conjugate
of $H$.  Actions are identified by homeomorphism conjugacy.
This specifies the exact fiber of the comparison map.  The
additional finite actions having a nontrivial kernel in that
projection are the full affine families already classified in
B.14; their kernels are retained in that result and are not
silently identified with an injective realization.

On $S^2$, a negative involution is conjugate either to a reflection
or to the antipodal map.  Indeed the sphere theorem in B.1 conjugates
the finite action into $O(3)$.  An orthogonal involution has
eigenvalues $\pm1$; negative determinant leaves precisely the
patterns $(1,1,-1)$ and $(-1,-1,-1)$.  Their fixed sets are a
circle and the empty set.  Together with the trivial action this
gives three rows over the two mapping class subgroups.
On the plane, a negative involution is conjugate to a reflection
by B.1, giving two rows including the trivial action.
For both surfaces $\pi=1$, so the lift extension is simply
$E=G$.  For the negative $C_2$ there is exactly one full section,
exactly one nontrivial partial section, and one maximal class,
with centralizer $1$ in $\pi$ and stabilizer $C_2$.
For the trivial group there is one full section and no nontrivial
partial section.  These are algebraic assertions about the
extension and include the antipodal action.  In that action
$p=-p$ implies $p=0$, which has no solution on the unit sphere.
The program verifies $A^2=I$, $\det A=-1$, and $\det(I-A)=8$
for $A=-I_3$.  It therefore records zero fixed components next
to the one algebraic full section.

A reflection has the local multiplier $1/2$.  For the antipodal
section no fixed-point germ exists from which to read such a
multiplier; the table marks that entry as undefined.  It separately
records the formal value $1/2$ obtained by extending the
orientation-sign convention to every negative involution.
This convention has a stated source and is not substituted for
a nonexistent fixed-point comparison.  In particular the
circle fixed by a sphere reflection has trivial image in
$\pi_1(S^2)$, so the hyperbolic oval-centralizer formula $\Z$
does not transfer to this case.

For the annulus use the simultaneous normal form of B.14.
Let
\[
 a(x,t)=(-x,t),\qquad
 b(x,t)=(x,-t),\qquad
 b'(x,t)=(x+\tfrac12,-t),\qquad
 c(x,t)=(-x,-t).                                      \tag{B.36.3}
\]
The three nontrivial mapping class elements are represented by
$a,b,c$.  A lift of the class of $a$ or $c$ has an arbitrary
angular shift, removed by translation conjugacy.
A lift of the class of $b$ has angular shift $0$ or $1/2$,
because its square is the identity.  These give $b$ and $b'$;
their fixed sets distinguish them.
For the whole four-group, normalize its angular reflection to
$a$.  Commutation forces the radial generator's angular shift
to be $0$ or $1/2$, giving exactly
$\langle a,b\rangle$ and $\langle a,b'\rangle$.
Conjugation cannot exchange the three nontrivial mapping class
elements, since their component group is abelian.  The fixed
oval or its absence again distinguishes these two actions.
This proves the seven-row classification.

Here are all its section counts; $f,p,\phi,o$ mean full sections,
nontrivial partial sections, $|\Phi|$, and $|\Phi/G|$:
\[
\begin{array}{l|rrrr}
\text{action}&f&p&\phi&o\\\hline
1&1&0&0&0\\
\langle a\rangle&2&2&2&2\\
\langle b\rangle&1&1&1&1\\
\langle b'\rangle&0&0&0&0\\
\langle c\rangle&2&2&2&2\\
\langle a,b\rangle&2&7&2&2\\
\langle a,b'\rangle&0&4&4&2
\end{array}                                             \tag{B.36.4}
\]
To verify them, the fixed sets of $a,b,b',c$ are respectively
two lines, one oval, the empty set, and two points.
In the split four-group, each of its two global fixed points
extends the corresponding point section and the incident
reflection sections; the oval section has both maximal extensions.
In the glide four-group, the two angular-reflection lines and the
two fixed points of $ab'$ are disjoint.  They are all maximal,
and the group interchanges each pair.  Thus (B.36.4) also gives
the required partial-section distinction.
Every nonidentity torsion element in these sections has multiplier
$1/2$.  The centralizer in $\pi=\Z$ is $\Z$ for the radial oval
section and $1$ for every point or line section.

These calculations can also be read directly from the extension.
For an affine element write its chosen angular shift as
$a_g\in[0,1)$ and let $\epsilon_g$ be its action on $\pi=\Z$.
The program records the integer cocycle
\[
 c(g,h)=a_g+\epsilon_ga_h-a_{gh}
 \quad\hbox{and}\quad
 (n,g)(m,h)=(n+\epsilon_gm+c(g,h),gh).                 \tag{B.36.5}
\]
It checks the cocycle identity on every triple.
In the universal cover $\R^2$, a finite subgroup of the lift
extension fixes a point; its fixed set is affine and connected.
Passing to the deck quotient identifies its conjugacy class
with the corresponding fixed component just listed.
This proves the section calculation without using a
center-free hypothesis.

For the torus apply the exact affine classification in B.27 and
retain precisely the normal forms whose translation kernel has
order one.  The linear projection of such a finite group is
faithful, which is exactly the comparison stipulated above.
Conversely every injective realized mapping class subgroup has
this property and is included by B.14.
There are seventeen affine forms over the thirteen integral
linear classes.  In the order of the linear classes in B.27,
the fiber sizes are
\[
                   1,2,1,1,1,3,1,1,1,1,1,2,1.       \tag{B.36.6}
\]
The proof of B.27 enumerates all cocycles and all affine
conjugacies without a word bound.  Its averaging comparison in
B.14 proves that no further topological conjugacy is missing.
Run the same exact fixed-point and reflection-circle computation
with no punctures removed.  Every isolated point section has
centralizer $1$ in $\Z^2$; every reflection oval has centralizer
the primitive invariant copy of $\Z$.  The trivial section has
centralizer $\Z^2$.  Local rotation matrices give the multipliers,
and the complete section poset gives the full and partial counts,
$\Phi$, and every stabilizer.  The resulting seventeen rows and
their exact rational matrices are included in the JSON.

The executed program \texttt{catalogue\_exceptional.py} carries
out all these finite calculations, records the two annulus-end
stabilizers and the plane's puncture stabilizer, and verifies the
four mapping class counts and the four realization-fiber totals.
Consequently it supplies a complete corrected computational
answer for the exceptional types, including the precise fibers
where B.22 prevents attaching one extension or one section table
to a bare mapping class subgroup.
\end{proof}

\begin{claim}{B.37}{Proved}{Complete genus-three mirror-disc quotient catalogue}
There are exactly $145$ conjugacy classes of finite subgroups $K$ of
$\mathcal T_{3,0}$ for which $H=K\cap\mathcal T^+_{3,0}$ has spherical
quotient and the induced involution on $S_3/H$ has a fixed circle.
Exactly $57$ contain a hyperelliptic involution, and $88$ do not.
The files \path{type_3_0_reflection_quotient_catalogue.csv},
\path{type_3_0_reflection_quotient_certificates.json}, and
\path{type_3_0_reflection_quotient_sections.json} give every
class, its exact finite group, and all its section data.
\end{claim}
\passage{\S{}18; folios 96--101; scans 193--203; Part B(f) of the task; genus-three mirror-disc quotient subcatalogue}

\begin{proof}
The quotient by the induced involution is a mirror disc $Q$.  Let
$(m_1,\ldots,m_k)$ be its cyclic list of corner orders and
$(n_1,\ldots,n_l)$ its list of interior cone orders.  The orientation
double of $Q$ is the spherical orbifold with cone multiset
\[
               \{m_1,\ldots,m_k,n_1,n_1,\ldots,n_l,n_l\}.
                                                        \tag{B.37.1}
\]
The involution on a sphere with a fixed circle is conjugate to a
reflection: cut along that circle into two discs and extend a boundary
parametrization over either disc, then use the involution to define
the parametrization on the other disc.  Cone points on its circle
become the corners of $Q$, while its other cone points occur in
exchanged pairs.  Thus the corner and interior data in (B.37.1)
exhaust its quotient types.  The corner list is considered up to
cyclic permutation and reversal; the interior list is unordered.

Use the following presentation of $\Lambda=\pi_1^{\rm orb}(Q)$:
\[
 \left\langle x_1,\ldots,x_l,e,c_0,\ldots,c_k\ \left|
 \begin{array}{c}
 x_i^{n_i}=1,\quad x_1\cdots x_l e=1,\\
 c_j^2=1,\quad(c_{j-1}c_j)^{m_j}=1,\quad c_k=ec_0e^{-1}
 \end{array}\right.\right\rangle .                    \tag{B.37.2}
\]
For $k=0$ there is one reflection $c_0$, commuting with $e$.
The orientation kernel $\Delta=\Lambda^+$ is the orbifold group
of (B.37.1).  Put $t=c_0$, $Z_j=c_{j-1}c_j$,
$P_j=Z_1\cdots Z_j$, $U_i=x_{l+1-i}^{-1}$, and $V_i=tU_it$.
Schreier rewriting, or the doubled fundamental polygon, gives the
spherical generating vector
\[
 (Z_1,\ldots,Z_k,U_1,\ldots,U_l,V_l^{-1},\ldots,V_1^{-1}),
 \qquad P_k\,U_1\cdots U_l\,(V_1\cdots V_l)^{-1}=1.
                                                        \tag{B.37.3}
\]
For example $e=U_1\cdots U_l$ and
$P_k=t c_k=tet e^{-1}=(V_1\cdots V_l)e^{-1}$, which proves
the product relation with its stated order.  Conjugation by $t$
acts on these generators by
\[
 \psi(Z_j)=P_{j-1}Z_j^{-1}P_{j-1}^{-1},\qquad
 \psi(U_i)=V_i,\qquad\psi(V_i)=U_i.                    \tag{B.37.4}
\]
These formulas include every base-path conjugation.

The input \cite{R-KNSData}, whose completeness is established in
\cite[Theorem 2.2, Corollary 2.3, and the planar calculation]{R-KNS},
contains every $\operatorname{Aut}(H)$-class of smooth epimorphisms
$\theta:\Delta\to H$ in genus three.  Its $197$ vectors cover all
$49$ genus-three signature/group pairs, including those with positive
quotient genus.  Here we use precisely the spherical-quotient pairs.
The original input and its attribution are retained in
\texttt{research\_data/actions-genus3.json}; B.33 verifies its orders
and surface signatures.  A prescribed ordering (B.37.3) is obtained
from the published sorted ordering by explicit adjacent Hurwitz
moves.  This is a bijection of epimorphism sets, since each move has
its displayed inverse in the program.

For each vector the program tests whether its kernel is invariant
under (B.37.4).  More explicitly, it attempts to define
\[
              \alpha(\theta(a))=\theta(\psi(a))
                         \qquad(a\in\Delta)             \tag{B.37.5}
\]
by propagation on the finite Cayley graph.  Agreement on every edge
is exactly the condition that (B.37.5) be well defined.  Surjectivity
then makes it an automorphism and (B.37.4) gives $\alpha^2=1$.
The resulting full action group is explicitly
\[
           K=H\rtimes_\alpha\langle t\rangle,
                 \quad t^2=1.                         \tag{B.37.6}
\]
Its epimorphism from (B.37.2) sends $x_i$ to their recovered vector
entries and $c_j$ to $t\theta(P_j)$.  Its kernel is exactly
$\ker\theta\subset\Delta$.  In particular no choice of a further
extension class is omitted: the orbifold kernel determines
$\Lambda/\ker\theta$ and hence (B.37.6).  Conversely every action
in the claim supplies this invariant kernel and therefore occurs.

We now prove the precise fusion calculation.  Two epimorphisms from
the same orbifold group have the same kernel exactly when a unique
isomorphism between their finite images intertwines them: define
that isomorphism by $\theta(a)\mapsto\theta'(a)$, which is well defined
exactly when the kernels agree.  Thus the published
$\operatorname{Aut}(H)$-normalization records kernels without loss.
A conjugacy of the surface actions descends to a decorated-orbifold
homeomorphism of $Q$ carrying one kernel to the other.  Conversely
such a quotient homeomorphism lifts through the two regular
orbifold coverings because it carries their defining kernels to
one another; its lift is the required surface conjugacy.  This is
an exact equivalence, so it remains to calculate the full action of
these disc homeomorphisms on the finite kernel set.

A disc homeomorphism permutes its colored boundary corners cyclically,
or reverses their cyclic order.  After composing with that boundary
map it fixes the boundary pointwise up to isotopy.  The remaining
mapping class is a braid of the interior cone points.  To see this,
forget the interior points and isotope the boundary-fixed disc map
to the identity by the Alexander isotopy; the tracks of the points
are its braid, and an isotopy of those tracks gives precisely an
isotopy of the punctured disc map by isotopy extension.  The colored
braid subgroup is generated by pure braid generators and adjacent
exchanges among equal colors after the colors are sorted.  The
program uses the standard words
$\sigma_{j-1}\cdots\sigma_{i+1}\sigma_i^2
\sigma_{i+1}^{-1}\cdots\sigma_{j-1}^{-1}$ and those equal-color
exchanges.  On interior generators an adjacent positive move is
\[
                    (x_i,x_{i+1})\longmapsto
                    (x_ix_{i+1}x_i^{-1},x_i),           \tag{B.37.7}
\]
fixing $e$ and every boundary reflection.  The inverse is also used
when forming the pure braid words.

Here are the boundary moves, including their effect on the chosen
reflection.  Extend the sequence by $c_{j+k}=ec_je^{-1}$.  One cyclic
shift has
\[
 \begin{split}
 (Z_1,\ldots,Z_k)&\longmapsto
       (Z_2,\ldots,Z_k,eZ_1e^{-1}),\\
 (x_1,\ldots,x_l)&\longmapsto(x_1,\ldots,x_l),\\
 t&\longmapsto tZ_1,\qquad
 \alpha\longmapsto\operatorname{Inn}(Z_1^{-1})\alpha.
 \end{split}                                           \tag{B.37.8}
\]
A reversal has
\[
 \begin{split}
 (Z_1,\ldots,Z_k)&\longmapsto(Z_k^{-1},\ldots,Z_1^{-1}),\\
 (x_1,\ldots,x_l)&\longmapsto(x_l^{-1},\ldots,x_1^{-1}),\\
 t&\longmapsto tP_k,\qquad
 \alpha\longmapsto\operatorname{Inn}(P_k^{-1})\alpha.
 \end{split}                                           \tag{B.37.9}
\]
They follow by sending $c_j$ respectively to $c_{j+1}$ and
$c_{k-j}$ in (B.37.2); in the reversal $e$ goes to $e^{-1}$.
Only the shifts and shifted reversals preserving the corner-order
list are self-homeomorphisms of the chosen colored quotient.
A fixed braid restores the sorted interior order after a reversal.
Equations (B.37.7)--(B.37.9), together with the preceding braid
argument, prove both that every programmed fusion is realized and
that every actual conjugacy is included.

The program computes the finite orbits under these moves.  It
identifies two generating vectors under $\operatorname{Aut}(H)$ by
their rooted, edge-colored Cayley graphs: starting at the identity,
visit vertices in generator order and record each colored successor.
Equality defines a unique group isomorphism mapping each generator
to the corresponding generator, and such an isomorphism plainly
gives equality.  Thus this is an exact isomorphism test rather than
a coarse group invariant.  The computation gives $145$ orbits.

For each orbit the second program constructs the regular action of
(B.37.6) explicitly from
$(h,\epsilon)(h',\epsilon')=(h\alpha^\epsilon(h'),\epsilon+\epsilon')$.
It checks (B.37.2), every local stabilizer injection, and
$|K|\chi^{\rm orb}(Q)=-4$.  The resulting cover is therefore a closed
orientable genus-three surface.  The finite coset incidence graph
then gives every fixed point and oval, multiplier, section,
extension, maximal section, centralizer, orbit and stabilizer as in
B.26.  These checks passed for all $145$ classes.

Finally the overlap with B.34 is identified by actual monodromy
isomorphisms.  In each of B.34's $57$ mirror-disc groups the program
extracts the positive group and the generators $x_i$, $Z_j$ and
conjugation by $c_0$.  It applies (B.37.8)--(B.37.9) and the interior
braids to put the signature in the same ordering, and locates its
rooted colored Cayley graph in the computed kernel orbit.  The file
\texttt{genus3\_hyperelliptic\_reflection\_crosscheck.json} records
the resulting bijection with $57$ different rows and equality of
all eight reported section-count columns.  Thus these are exactly
the common classes, leaving $88$ additional ones.
\end{proof}

The $145$ classes have the following order distribution:
\[
\begin{array}{c|rrrrrrrr}
 |K|&4&6&8&12&14&16&18&24\\\hline
 \#&5&3&36&13&3&35&1&15\\[2pt]
 |K|&28&32&42&48&64&96&192&336\\\hline
 \#&1&13&1&11&2&4&1&1.
\end{array}                                             \tag{B.37.10}
\]
Together B.34 and B.37 therefore give $149$ distinct mixed classes.
The additional four are the projective-plane quotients in B.34.
The scope of B.37 requires both a spherical positive quotient and a
fixed-circle involution on it; the remaining quotient types are
listed in the accompanying computation report.

\begin{claim}{B.38}{Proved}{The free positive involution and the six negative involutions in genus three}
There are exactly twelve conjugacy classes of mixed finite subgroups
of $\mathcal T_{3,0}$ whose positive subgroup is a free involution.
Each is a four-group.  Their quotient signatures and numbers are
\[
\begin{array}{c|ccccc}
\text{underlying quotient}&S_{0,3}&S_{1,1}&N_{1,2}&N_{2,1}&N_{3,0}\\\hline
\text{number of classes}&4&1&4&2&1 ,
\end{array}                                             \tag{B.38.1}
\]
where $S_{h,b}$ is orientable of genus $h$ with $b$ mirror boundary
components, and $N_{q,b}$ is nonorientable with $q$ crosscaps and
$b$ mirror boundary components.  There are no cones or corners.
There are also exactly six mixed classes whose positive subgroup
is trivial: the negative involutions with quotient signatures
\[
 S_{0,4},\quad S_{1,2},\quad
 N_{1,3},\quad N_{2,2},\quad N_{3,1},\quad N_{4,0}.
                                                        \tag{B.38.2}
\]
All eighteen classes, their exact monodromies, complete section
posets, $\Phi$, stabilizers, centralizers and multipliers appear in
\path{type_3_0_free_positive_and_anti_catalogue.csv}
and its \texttt{sections.json} companion.
\end{claim}
\passage{\S{}15; folios 73--75; scans 148--151; \S{}15; folios 77--79; scans 155--160; \S{}18; folios 96--101; scans 193--203; \S{}33; folios 222--225; scans 440--445; Part B(f) of the task; no further manuscript assertion}

\begin{proof}
Let $X$ be the closed orientable genus-three surface.  First let
$G^+=\langle a\rangle$ have order two and act freely, and let
$G$ contain negative elements.  Then $|G|=4$.
A group of order four is cyclic or a four-group: an element of
order four gives the first case; otherwise $(xy)^2=1$ for all
$x,y$, which gives commutativity and the second case.
There is no nontrivial positive point stabilizer.
A negative point stabilizer is generated by a reflection, so
$Q=X/G$ has no cone or corner points.  Its orbifold Euler
characteristic is $-4/4=-1$.  The complete list of compact
underlying surfaces with this Euler characteristic is exactly
(B.38.1).

The cyclic possibility cannot occur.  A negative element of
order four has no fixed point, since a negative local finite
stabilizer element is a reflection of order two.  Its square
is free by assumption, so the whole action would be free.
The quotient would then be $N_{3,0}$.
Its fundamental group has the presentation
\[
 \langle u_1,u_2,u_3\mid u_1^2u_2^2u_3^2=1\rangle .
\]
An orientation-compatible monodromy to $C_4=\Z/4$ sends each
crosscap generator to an odd residue.  Each square has image
$2$, so the relator has image $6=2\pmod4$, a contradiction.
The program checks all eight assignments of odd residues.
Thus $G$ is the four-group
$V_4=\mathbf F_2^2$, with orientation given by its second
coordinate, and $a=(1,0)$.

We classify all four-group monodromies, including their full
topological equivalence relation.  The mirror orbifold group
of an orientable $S_{h,b}$ has generators
$a_i,b_i,e_j,c_j$ and relations
\[
 c_j^2=1,\qquad[e_j,c_j]=1,\qquad
 \prod_i[a_i,b_i]\prod_je_j=1.                         \tag{B.38.3}
\]
The negative $c_j$ is the reflection on the $j$th boundary,
and the positive $e_j$ traverses that boundary.
For $N_{q,b}$ replace the handles by crosscap generators
$u_i$, and replace the last relation by
\[
                         \prod_i u_i^2\prod_j e_j=1.  \tag{B.38.4}
\]
The orientation character $\varepsilon$ is zero on handles
and $e_j$, and one on crosscaps and $c_j$.
An orientation-compatible epimorphism to $V_4$ is exactly
\[
 q_\alpha:\Pi_Q\longrightarrow V_4,\qquad
 q_\alpha(x)=(\alpha(x),\varepsilon(x)),
 \quad
 \alpha\in H^1(\Pi_Q;\mathbf F_2)
       \setminus\{0,\varepsilon\}.                   \tag{B.38.5}
\]
In coordinates write the handle or crosscap values as $\eta_i$,
the boundary values as $\beta_j=\alpha(e_j)$, and the reflection
values as $\gamma_j=\alpha(c_j)$.
The only linear relation is
$\sum_j\beta_j=0$.
The two automorphisms of $V_4$ preserving orientation identify
$\alpha$ with $\alpha+\varepsilon$.

Every map (B.38.5) is injective on local reflection groups.
Its kernel is therefore torsion-free; one proves this either
from the universal hyperbolic orbifold cover or by cutting along
arcs to its reflection polygon, where every finite stabilizer
is a displayed local reflection group.
The kernel lies in the orientation kernel, and its Euler
characteristic is $4(-1)=-4$.  It is consequently the fundamental
group of a closed orientable genus-three surface.
Conversely every action under consideration yields exactly
this construction.  Two such actions are conjugate precisely
when an orbifold homeomorphism of $Q$ carries one monodromy
kernel to the other, or equivalently when it carries the
monodromy to the other after an automorphism of $V_4$.
This follows by lifting paths in the two covers, with the
local reflection charts giving the extension across mirrors.
Thus the following finite cohomology-orbit calculation is an
exact classification of actions.

We prove the entire mapping class action needed in this
calculation.  A boundary permutation permutes the pairs
$(\beta_j,\gamma_j)$.  The interior homology action preserves
the mod-two intersection pairing, its radical generated by
boundary classes, and the orientation character.
It introduces no reflection coordinate: an interior loop
is mapped to an interior loop, whereas each reflection generator
is mapped to a conjugate of the reflection at the corresponding
boundary.  Thus the $\gamma_j$ are affected only by boundary
permutations and by the global replacement
$\alpha\mapsto\alpha+\varepsilon$.

For $S_{0,3}$ the interior homology is generated by its boundary
classes, with their single sum relation, so boundary permutations
are the whole action.  When $\beta=(0,0,0)$, surjectivity leaves
one nonconstant reflection-color pattern up to permutation and
global color reversal.  When $\beta$ has two ones, distinguish
the boundary with value zero and use global color reversal to
make its reflection value zero.  The two remaining reflection
values have weights zero, one or two.  This gives $1+3=4$ classes.

For $S_{1,1}$ the boundary value is zero.
Global color reversal sets its reflection value to zero.
The remaining nonzero pair of handle values is a single orbit
under $\mathrm{Sp}_2(\mathbf F_2)$, generated by the two Dehn
twists in a pair of basis curves.  This gives one class.

For $N_{1,2}$ write $u$ for its one-sided core and $e$ for
either mod-two boundary class.  Its intersection form has
$u\cdot u=1$, with $e$ in the radical.
An automorphism preserving the boundary classes and the
orientation character sends $u$ either to $u$ or to $u+e$.
Both possibilities are realized.  Push one boundary disc once
around the one-sided core and return it as an unparametrized
disc; its boundary orientation may reverse.  A core arc crossing
the track once acquires one boundary circuit.  Thus, modulo two,
this homeomorphism sends $u$ to $u+e$ and leaves both reflection
coordinates unchanged.  The boundary exchange is realized away
from the core.  These maps generate every permitted action on
the displayed homology and boundary data.
After global color reversal sets $\eta=\alpha(u)$ to zero,
the case $\beta=(0,0)$ has two nonzero reflection patterns up
to exchange.  The case $\beta=(1,1)$ has two patterns: the
boundary push followed by global reversal interchanges
$\gamma=(0,0)$ and $(1,1)$, while $(0,1)$ and $(1,0)$ are
exchanged by the boundary permutation.  This gives four classes.
The boundary push is necessary for this count.

For $N_{2,1}$ the boundary value is zero.
Set its reflection value to zero by global reversal.
The intersection form on the two crosscap classes is the
identity matrix over $\mathbf F_2$; its orthogonal group has
exactly the identity and the exchange of those two classes.
They are realized by exchanging the two crosscaps.
The nonzero crosscap vectors therefore have two orbits,
represented by $(1,0)$ and $(1,1)$.

For $N_{3,0}$ the mod-two intersection form is the identity
matrix in three crosscap classes, and its characteristic
vector is $(1,1,1)$.  An orthonormal basis cannot contain that
characteristic vector, since it pairs to one with every vector
of odd self-intersection.  Its members must therefore be the
three standard basis vectors.  The full orthogonal group is
their permutation group, realized by crosscap permutations.
Global addition of $\varepsilon=(1,1,1)$ exchanges weights
one and two, leaving one orbit after $0,\varepsilon$ are
excluded.  This proves all five numbers in (B.38.1).

We now compute the sections explicitly.  A mirror boundary
with data $(\beta_j,\gamma_j)$ has local reflection
$c_j=(\gamma_j,1)$ and boundary monodromy
$e_j=(\beta_j,0)$.
Its preimage is described by the two cosets
$V_4/\langle c_j\rangle$; one trip around the boundary
permutes these cosets by $e_j$.
Consequently it gives two fixed ovals when $\beta_j=0$,
and one when $\beta_j=1$.
The two-oval case has an oval stabilizer of order two,
$\langle c_j\rangle$, and the one-oval case has stabilizer
$V_4=\langle c_j,e_j\rangle$.
Both are one $V_4$-orbit of section classes.
There are no point sections because the positive involution
is free, and there are no full sections because $V_4$ cannot
fix a point without giving a positive local stabilizer.
Every reflection section is maximal, has centralizer $\Z$
in the surface group, and has multiplier $1/2$ on its
nonidentity element.  Thus
\[
 |\Phi|=\#\{\text{nontrivial partial sections}\}
       =2b-\#\{j:\beta_j=1\},
 \qquad |\Phi/V_4|=b.                                 \tag{B.38.6}
\]
The program constructs the actual finite coset graph, with
three subdivisions per boundary traversal, computes its
components and the group action, and verifies every count
and stabilizer in this formula.

Finally a mixed group with trivial positive subgroup is $C_2$.
Its quotient has orbifold Euler characteristic $-2$ and no
cones or corners.  An orientable quotient must have mirror
boundary; otherwise its orientation character would be
trivial and could not give the required negative covering.
Solving $2-2h-b=-2$ with $b>0$, or $2-q-b=-2$ with $q>0$,
gives precisely (B.38.2).
For each quotient the monodromy is forced: it is the orientation
character, taking every crosscap and local reflection to the
nonidentity element and every handle and boundary traversal
to zero.  Its orientation double is connected, torsion-free,
and has Euler characteristic $-4$.  The classification of
compact quotient surfaces and path lifting prove that each
signature gives one conjugacy class and that distinct signatures
give distinct classes.
There are $b$ fixed ovals, hence $b$ full sections, $b$
nontrivial partial sections, and $b$ maximal classes, each
fixed by the whole $C_2$ and with surface-group centralizer
$\Z$.  The case $N_{4,0}$ has none.  These are the six
genus-three real types of B.16 in an executed form.

The code \texttt{catalogue\_genus3\_free\_positive.py} enumerates
every cohomology vector in (B.38.5), applies exactly the
proved generators of the equivalence relation, checks
surjectivity, all quotient relations, orientation compatibility,
and the covering Euler characteristic, and records every
orbit.  Its second calculation constructs all six orientation
doubles in (B.38.2).  This supplies the complete eighteen-row
catalogue and its section data.
\end{proof}

\begin{claim}{B.39}{Proved}{Complete genus-three projective-plane quotient catalogue}
There are exactly $11$ conjugacy classes of mixed finite genus-three
actions whose positive quotient is a sphere and whose induced
involution on that sphere is fixed-point-free.  Their full quotients
are projective planes with cone points.  Four classes are identified
with B.34 and seven are additional.  Consequently B.37 and the present
claim give exactly $156$ mixed classes with spherical positive quotient.
\end{claim}
\passage{\S{}18; folios 96--101; scans 193--203; Part B(f) of the task; genus-three projective-plane quotient subcatalogue}

\begin{proof}
A fixed-point-free sphere involution has projective-plane quotient:
its quotient is a closed surface of Euler characteristic one, and
its orientation cover is the given sphere.  Cone points are exchanged
in pairs.  Thus each multiplicity in the positive cone multiset must
be even.  Inspection of the complete input in B.33 leaves exactly
$\mathrm O3.4$, $\mathrm O3.9$, $\mathrm O3.11$, $\mathrm O3.15$,
$\mathrm O3.21$, $\mathrm O3.22$, and $\mathrm O3.25$.
Their paired cone orders are respectively
$(2,2,2,2)$, $(2,2,2)$, $(4,4)$, $(2,6)$, $(2,4)$, $(2,4)$,
and $(2,3)$.

For a projective plane with cone orders $n_1,\ldots,n_l$ use
\[
 \Lambda=\langle x_1,\ldots,x_l,d\mid
       x_i^{n_i}=1,\quad x_1\cdots x_l d^2=1\rangle .   \tag{B.39.1}
\]
Set $X_i=x_i$, $Y_i=d x_i d^{-1}$ and
$A=d^2=(X_1\cdots X_l)^{-1}$.  Its orientation double has the
spherical vector
\[
        v=(X_1,\ldots,X_l,Y_l^{-1},\ldots,Y_1^{-1}),
                \qquad\prod_iX_i=\prod_iY_i.            \tag{B.39.2}
\]
Conjugation by $d$ is
\[
       X_i\longmapsto Y_i,\qquad
       Y_i\longmapsto A X_iA^{-1}.                     \tag{B.39.3}
\]
For every published $\operatorname{Aut}(H)$-normalized vector in
these seven cases, the program tests that (B.39.3) descends to an
automorphism $\alpha$ of its finite positive image $H$.  The same
Cayley-graph propagation as in B.37 gives an exact kernel-invariance
test.  It also verifies
$\alpha^2=\operatorname{Inn}(A)$ and $\alpha(A)=A$.
The full finite image is then the set $H\times\{0,1\}$ with product
\[
 (h,\epsilon)(k,\delta)=
       \bigl(h\alpha^\epsilon(k) A^{\epsilon\delta},
                    \epsilon+\delta\bmod2\bigr).      \tag{B.39.4}
\]
The two identities just verified prove associativity; the element
$d=(1,1)$ has square $A$ and induces $\alpha$ on $H$.
Thus (B.39.4) is exactly $\Lambda/\ker\theta$, where
$\theta$ is the positive epimorphism (B.39.2).

For clarity, the full source fusion is as follows.  The colored
mapping class group of the punctured projective plane is generated
by the colored disc braids and by sliding each puncture once around
a one-sided loop.  One direct proof cuts the projective plane along
a projective line disjoint from the marked disc: a moving puncture
crossing the cut is separated off as a one-sided slide, after which
all the remaining tracks form a braid in the disc.  Equivalently,
forgetting punctures successively expresses their point-pushing
loops by peripheral loops and the one-sided generator.  The former
are the pure disc braids.  The unmarked projective-plane mapping
class group is trivial: isotope the image of a projective line to
that line, adjust its parametrization by a projective rotation, and
apply the disc Alexander isotopy after cutting.  This proves that
these point motions exhaust all mapping classes.  Color-preserving
permutations are supplied by the same equal-color braid generators
used in B.37.

The slide of the last puncture has the exact substitutions
\[
 x_i\longmapsto x_i\quad(i<l),\qquad
 x_l\longmapsto(x_l d)x_l^{-1}(x_l d)^{-1},\qquad
 d\longmapsto x_l d.                                   \tag{B.39.5}
\]
They are obtained by tracking the meridian through a neighborhood
of that puncture and the one-sided loop.  The two-cone case, including
its geometric puncture-slide interpretation, is verified independently
in \cite[Lemma 3.19]{C_GromadzkiSzepietowski2016}.  Other punctures
outside this neighborhood are fixed, so the same formula applies
to all $l$ used here.  It respects the product relation directly:
if $B=x_1\cdots x_{l-1}$ and $x=x_l$, then
\[
 B(x d x^{-1}d^{-1}x^{-1})(x d)^2=B x d^2=1.
\]
In the finite positive group its effect is
\[
 x_l\longmapsto x_l\alpha(x_l)^{-1}x_l^{-1},\qquad
 \alpha\longmapsto\operatorname{Inn}(x_l)\alpha,
 \qquad A\longmapsto x_l\alpha(x_l)A.                  \tag{B.39.6}
\]
To slide any other puncture, a specified sequence of adjacent braids
moves it to the last position, applies (B.39.6), and moves it back.
The full word preserves its color even when the intermediate braids
exchange different colors.  Equations (B.39.5)--(B.39.6) and the disc
braids therefore give exactly the source mapping-class action.
The kernel-to-cover conjugacy argument of B.37 applies without change.

Executing these finite orbits gives respectively $1,3,3,1,2,0,1$
classes in the seven signature/group pairs above.  For every class,
\texttt{catalogue\_genus3\_projective\_quotients.py} constructs the
regular representation of (B.39.4), checks every relation in
(B.39.1), injectivity of all local cyclic groups, and Euler
characteristic $-4$.  The finite point-coset calculation gives every
partial section and its multiplier, centralizer and extensions.
No negative element has a fixed point, since its image on the
positive quotient sphere is fixed-point-free.  Thus all nontrivial
sections in these rows have positive cyclic image and are point
sections; their surface-group centralizers are trivial.

Finally the four projective-plane groups in B.34 are compared using
their actual cone generators and conjugation by their crosscap
element.  After a fixed ordering braid, their doubled vectors are
identified in the finite kernel orbits by the rooted colored Cayley
graph isomorphism test.  The resulting exact identifications are
\[
 \mathrm{G3H}_{002}\leftrightarrow\mathrm{G3P}_{001},\quad
 \mathrm{G3H}_{011}\leftrightarrow\mathrm{G3P}_{002},\quad
 \mathrm{G3H}_{029}\leftrightarrow\mathrm{G3P}_{008},\quad
 \mathrm{G3H}_{039}\leftrightarrow\mathrm{G3P}_{010}.
                                                        \tag{B.39.7}
\]
All seven applicable section-count columns agree under this
bijection.  Hence precisely seven classes are additional to B.34.
The quotient involution on the positive sphere has a fixed circle
in B.37 and has empty fixed set here; a conjugacy induces a sphere
conjugacy carrying these fixed sets to each other.  Thus none of
these eleven rows can be conjugate to a B.37 row, proving the total
$145+11=156$.
\end{proof}

Here $p$ counts nontrivial partial section classes, and $m,o$ are
$|\Phi|$ and $|K\backslash\Phi|$.  The final column $i$ counts negative
involutions in the displayed finite group; all of them act freely.
\[
\begin{array}{c|r|l|r|r|r|r}
 \text{class}&|K|&\text{cone orders}&p&m&o&i\\\hline
 001&4&(2,2,2,2)&8&8&4&2\\
 002&8&(2,2,2)&12&12&3&0\\
 003&8&(2,2,2)&12&12&3&2\\
 004&8&(2,2,2)&12&12&3&4\\
 005&8&(4,4)&8&4&2&0\\
 006&8&(4,4)&8&4&2&2\\
 007&8&(4,4)&8&4&2&4\\
 008&12&(2,6)&12&8&2&2\\
 009&16&(2,4)&16&12&2&0\\
 010&16&(2,4)&16&12&2&0\\
 011&24&(2,3)&20&20&2&4
\end{array}                                             \tag{B.39.8}
\]
Every row has zero full splittings, since its full image contains a
negative element and no such element fixes a point.  The JSON
\texttt{type\_3\_0\_projective\_quotient\_sections.json} includes all
finite permutations, local generators, multipliers and section
incidences behind this table.

\begin{claim}{B.40}{Proved}{Complete annulus and M\"obius-band quotients in genus three}
There are exactly $57$ mixed finite genus-three action classes whose
full quotient is an annulus or a M\"obius band with its cone and mirror
corner structure.  There are $39$ annulus classes and $18$ M\"obius-band
classes.  Their positive quotient is a torus.  Every group, generating
vector, quotient signature and complete section poset is given in
\path{type_3_0_torus_mirrors_catalogue.csv},
\path{type_3_0_torus_mirrors_certificates.json}, and
\path{type_3_0_torus_mirrors_sections.json}.
\end{claim}
\passage{\S{}18; folios 96--101; scans 193--203; Part B(f) of the task; genus-three annulus and Mobius quotient cases}

\begin{proof}
An involution reversing the orientation of a torus has annulus,
M\"obius-band or Klein-bottle quotient.  Indeed its fixed set consists
of disjoint circles, and its quotient has Euler characteristic zero;
the orientation double is the torus.  The classification of compact
surfaces leaves the orientable annulus, the nonorientable M\"obius
band, and the closed nonorientable Klein bottle.  Their orientation
doubles give respectively two, one, and zero fixed circles.  Cone
points on a fixed circle become mirror corners, and other cone points
are paired to interior cones.  The Klein-bottle cases are B.41.

The complete positive input of B.33 leaves the following seven
signature/group pairs.  In each row $r$ is the number of positive
quotient cone points, all of order $m$; the last two columns are the
executed numbers of full annulus and M\"obius-band classes.
\[
\begin{array}{c|c|r|r|r|r}
 \text{positive row}&H&m&r&\text{annulus}&\text{M\"obius}\\\hline
 \mathrm O3.3&C_2&2&4&10&3\\
 \mathrm O3.5&C_3&3&2&5&4\\
 \mathrm O3.7&C_4&2&2&8&4\\
 \mathrm O3.8&C_2^2&2&2&9&4\\
 \mathrm O3.12&S_3&3&1&3&1\\
 \mathrm O3.17&D_8&2&1&3&1\\
 \mathrm O3.18&Q_8&2&1&1&1
\end{array}                                             \tag{B.40.1}
\]
Here $D_8$ has order eight.  For $l$ interior cones, the total number
of corners is $r-2l$.  On an annulus distribute them as
$k_1\le k_2$ with $k_1+k_2=r-2l$; on a M\"obius band place all
$k=r-2l$ on its single boundary.  All their orders are $m$, so these
integer choices already classify the colored quotient signatures.

We first give an exhaustive construction of the possible oriented
finite groups, rather than assuming a small-group name list.
Choose a negative coset representative $t$ in an index-two extension
of $H$.  Let $\alpha(h)=t h t^{-1}$ and $A=t^2$.  Its complete
extension data satisfy
\[
              \alpha^2=\operatorname{Inn}(A),\qquad
                          \alpha(A)=A.                 \tag{B.40.2}
\]
Conversely the product of (B.39.4), with these data, defines that
extension:
\[
 (h,\epsilon)(k,\delta)=
       (h\alpha^\epsilon(k)A^{\epsilon\delta},
                                  \epsilon+\delta\bmod2).
                                                        \tag{B.40.3}
\]
Every oriented group isomorphism restricts to some $\beta\in
\operatorname{Aut}(H)$ and sends $t$ to some $k t'$ in the negative
coset.  Therefore the entire isomorphism action on (B.40.2) is generated
by transport by $\beta$ and by
\[
       (\alpha,A)\longmapsto
       \bigl(\operatorname{Inn}(k)\alpha,\ k\alpha(k)A\bigr).
                                                        \tag{B.40.4}
\]
The converse follows by sending the displayed generators accordingly
in (B.40.3), so this is an exact fusion test.  The program enumerates
$\operatorname{Aut}(H)$ by assigning images to a generating set and
propagating on its complete Cayley graph, enumerates every pair
(B.40.2), and takes precisely the finite orbits (B.40.4).
For the seven $H$'s in (B.40.1), it obtains respectively
$2,2,4,3,1,4,4$ oriented extension classes.

Here are the full source presentations and enumeration conditions.
For an annulus, put $X=x_1\cdots x_l$ and use positive generators
$x_i,e_1,e_2$ with
\[
                x_i^m=1,\qquad X e_1e_2=1.             \tag{B.40.5}
\]
On its $j$th boundary there are reflections
$c_{j,0},\ldots,c_{j,k_j}$ satisfying
\[
 c_{j,s}^2=1,\quad(c_{j,s-1}c_{j,s})^m=1,\qquad
               c_{j,k_j}=e_j c_{j,0}e_j^{-1}.          \tag{B.40.6}
\]
For $k_j=0$, there is one reflection commuting with $e_j$.
For a M\"obius band there is one negative crosscap generator $d$,
one boundary with (B.40.6), and the relation
\[
                X e d^2=1,
                   \qquad e=X^{-1}d^{-2}.              \tag{B.40.7}
\]
In each finite extension the program assigns every possible positive
order-$m$ element to every $x_i$, every positive element to $e_1$
in the annulus, and every negative element to $d$ in the M\"obius
band.  It then assigns every sequence of negative involutions to
the mirrors, requiring exact order $m$ for each adjacent product
and the stated closing relation.  It retains precisely the vectors
generating the whole finite group.  Two involutions whose product
has exact order $m$ generate a dihedral group of order $2m$; thus
all local groups inject.  The orientation character takes the
specified signs on every generator, so the kernel is oriented and
contains no reflection.  The signature equation gives
$|K|\chi^{\rm orb}(Q)=-4$; its torsion-free surface kernel therefore
has genus three.  Conversely every action with this quotient yields
one of these enumerated epimorphisms.

Target automorphisms are removed by the exact rooted colored Cayley
graph test of B.37, now applied to the full NEC generating vector.
Its prescribed generator signs force any resulting isomorphism to
preserve the orientation character.  For source fusion, the following
are explicit complete generators.  The interior braids are
(B.37.7), with all other generators fixed.  A cyclic boundary move
sends $c_{j,s}$ to $c_{j,s+1}$, extending the reflection sequence
by (B.40.6), while fixing $x_i$, $e_j$, and the crosscap generator.
After $k_j$ moves it conjugates that boundary's reflection sequence
by $e_j$ and fixes the other generators, which supplies the boundary
Dehn twist where it is needed.

For the annulus, pushing the last interior point once around the
first boundary, with $a=x_l$ and $e=e_1$, has
\[
 \begin{aligned}
 a&\longmapsto (a e)a(a e)^{-1},\\
 e_1&\longmapsto a e_1a^{-1},\\
 c_{1,s}&\longmapsto a c_{1,s}a^{-1}.
 \end{aligned}                                           \tag{B.40.8}
\]
It fixes the other $x_i$, $e_2$ and the second reflection sequence.
This is the square of the Hurwitz exchange between the puncture
meridian and the first boundary loop; in particular it preserves
$a e_1$ exactly.  The interior braids conjugate this to pushes of
each other interior point.  When $k_1=k_2$, interchange of the two
boundary components is
\[
 \begin{aligned}
 e_1&\longmapsto e_1e_2e_1^{-1},&
 e_2&\longmapsto e_1,\\
 c_{1,s}&\longmapsto e_1c_{2,s}e_1^{-1},&
 c_{2,s}&\longmapsto c_{1,s},
 \end{aligned}                                           \tag{B.40.9}
\]
fixing the interior generators.  This is their orientation-preserving
boundary exchange, with the chosen connecting paths.

The remaining annulus generator reverses the orientation of its
fundamental region.  For the peripheral list
$(w_1,\ldots,w_{l+2})=(x_1,\ldots,x_l,e_1,e_2)$, write
$W_i=w_1\cdots w_i$.  It sends
\[
             w_i\longmapsto W_{i-1}w_i^{-1}W_{i-1}^{-1}.
                                                        \tag{B.40.10}
\]
For the first boundary its reflection sequence is reversed and
conjugated by $X$; for the second it is reversed and conjugated by
$X e_1$.  Reversing a cycle means
$(c_0,\ldots,c_{k-1})\mapsto(c_k,c_{k-1},\ldots,c_1)$,
with $c_k=e c_0e^{-1}$; an empty cycle retains its single reflection.
These substitutions verify (B.40.5)--(B.40.6) directly.

For the M\"obius band, a one-sided slide of the last interior point
is particularly explicit.  Set
\[
          a=x_l,\qquad a'=(a d)a^{-1}(a d)^{-1}.
\]
The substitution is
\[
 \begin{aligned}
  x_l&\longmapsto a',& d&\longmapsto a d,\\
  c_s&\longmapsto a'^{-1}a\,c_s\,a^{-1}a'.
 \end{aligned}                                           \tag{B.40.11}
\]
It fixes the other interior generators and determines $e$ by
(B.40.7).  To obtain it geometrically, regard the boundary temporarily
as the last peripheral loop in the projective-plane presentation:
exchange $(a,e)$ by its Hurwitz move, slide the now-last puncture
by (B.39.5), and undo that exchange.  This gives exactly
(B.40.11), including the displayed transport of every boundary
reflection.  Finally a reversal of the core has
\[
 \begin{aligned}
 (x_1,\ldots,x_l)&\longmapsto(x_l^{-1},\ldots,x_1^{-1}),\\
 d&\longmapsto d^{-1},\qquad
 c_s\longmapsto X c_{k-s}X^{-1}.
 \end{aligned}                                           \tag{B.40.12}
\]
Here $e$ goes to $X e^{-1}X^{-1}$, so the closing relation is
preserved.

We justify that this is the whole source action.  First prescribe
the cyclic permutations, and possible reversal or interchange, of
the marked boundary components.  The boundary moves above realize
exactly these permutations.  Upon forgetting the interior points,
a remaining annulus map is generated by its core twist, and a
remaining M\"obius-band map by reversal of its core.  Cutting along
a spanning arc gives this classification: the annulus permits the
integer winding of the arc, while cutting the M\"obius band reduces
a boundary-fixed map to a disc isotopy.  The boundary reversal is
(B.40.12).  The core twist is already among the full cyclic boundary
moves when both boundaries are marked; when a boundary is unmarked
its rotation isotopes it away in the unpunctured quotient.
The kernel of forgetting interior points consists of their braids.
An interior point in an annulus minus the other points has fundamental
group generated by the puncture meridians and the core loop; in a
M\"obius band the corresponding extra generator is one-sided.
Forgetting points one at a time, and using isotopy extension, therefore
shows that their braid groups are generated by the ordinary interior
braids and respectively (B.40.8) or (B.40.11).  This also includes
any twist that becomes trivial only after the interior points are
forgotten.  Thus every geometric source automorphism is included,
and each included automorphism is realized by the maps just described.
The kernel-to-cover lifting argument of B.37 proves that the computed
orbits are exactly the surface conjugacy classes.

Execution gives the $57$ classes in (B.40.1).  There is an independent
check on the source generators in every one-cone, no-corner case:
the complete primary automorphism lists in
\cite[Lemmas 3.7 and 3.8]{C_GromadzkiSzepietowski2016} give exactly the
same partitions of the entire target-normalized epimorphism sets.
The file \path{check_torus_mirror_generators_output.json}
records all six comparisons, comparing full orbit memberships and
not just their numbers.

The program then computes all section invariants.  Its mirror
incidence graph has a separate closed cycle for each quotient
boundary, with endpoint gluing by that boundary's actual element
$e_j$.  At every corner it uses the full local dihedral group, and
checks that the fixed subgraph of each reflection has degree two
at every vertex.  Consequently its connected components are exactly
the fixed ovals, including all their intersections with other
reflection fixed sets.  Cone and corner cosets give the isolated
fixed points.  The procedure of B.26 then supplies every section,
multiplier, centralizer, extension incidence, maximal section, group
orbit and stabilizer.  All $57$ presentations, local injections,
Euler characteristics and complete section calculations passed.
\end{proof}

The $57$ classes have orders $4,6,8,12,16$ with multiplicities
$13,9,25,4,6$, respectively.  The full quotient surfaces in this
claim have boundary, whereas those in B.41 are closed.  A conjugacy
of actions descends to a homeomorphism of their full quotient
orbifolds and therefore preserves boundary.  Their positive
quotients have genus one, whereas those in B.37 and B.39 have genus
zero and those in B.38 have genus two or three.  The induced
homeomorphism of positive quotients preserves genus.  These exact
quotient maps prove that the four portions contribute disjoint
conjugacy classes.

\begin{claim}{B.41}{Proved}{The six genus-three actions with Klein-bottle quotient}
There are exactly six mixed finite-subgroup classes in
$\mathcal T_{3,0}$ whose quotient orbifold has underlying
Klein bottle.  Their positive subgroups, full groups, cone orders,
and section counts are
\[
\begin{array}{c|c|c|rrrr}
G^+&G&\text{cone orders}&f&p&|\Phi|&|\Phi/G|\\\hline
C_2&C_4&2,2&0&4&4&2\\
C_2&C_2^2&2,2&0&4&4&2\\
C_3&C_6&3&0&2&2&1\\
C_4&C_8&2&0&4&4&1\\
C_4&C_4\times C_2&2&0&4&4&1\\
C_2^2&D_8&2&0&4&4&1
\end{array}                                             \tag{B.41.1}
\]
Here $f$ and $p$ denote full and nontrivial partial section classes,
and $D_8$ has order eight.  The files
\path{type_3_0_klein_quotient_catalogue.csv} and
\path{type_3_0_klein_quotient_sections.json} contain
every group operation, monodromy, local multiplier, section
stabilizer and centralizer.  All centralizers in the surface
group are trivial.  In the primary positive catalogue, these
rows lie over O3.3, O3.5, O3.7 and O3.8, respectively in
numbers $2,1,2,1$.
\end{claim}
\passage{\S{}15; folios 73--75; scans 148--151; \S{}15; folios 77--79; scans 155--160; \S{}18; folios 96--101; scans 193--203; Part B(f) of the task; no further manuscript assertion}

\begin{proof}
The orientation double of a Klein bottle is a torus.
An interior cone of order $m$ in the quotient has two lifts of
order $m$ in that double, and there are no mirrors or corners.
The complete positive catalogue in B.33 therefore leaves
precisely the positive actions
\[
 (C_2;2,2,2,2),\quad(C_3;3,3),\quad
 (C_4;2,2),\quad(C_2^2;2,2),                          \tag{B.41.2}
\]
with torus underlying quotient.
The other torus-positive rows have a single cone and cannot
arise by this double-cover comparison.
Thus the full groups have orders four, six or eight.

We enumerate all oriented group extensions, rather than
presupposing a list of abstract groups.
For a fixed positive group $H$ choose a negative representative
$t$, put $\alpha(h)=tht^{-1}$ and $A=t^2\in H$.
Associativity is equivalent to
\[
           \alpha^2=\operatorname{Inn}(A),
           \qquad\alpha(A)=A.                         \tag{B.41.3}
\]
Conversely the set $H\times\{0,1\}$ with normal form $ht^i$
and relations $th=\alpha(h)t$, $t^2=A$ is a group for every
pair in (B.41.3); checking the two possible negative factors
verifies associativity.  Changing $t$ to $kt$ changes the pair to
\[
 \bigl(\operatorname{Inn}(k)\alpha,\ k\alpha(k)A\bigr),
                                                        \tag{B.41.4}
\]
and changing the identification of $H$ by $\beta\in\operatorname{Aut}(H)$
changes it to
$(\beta\alpha\beta^{-1},\beta(A))$.
These are all orientation-preserving isomorphisms of the
extensions, because any such isomorphism restricts to an
automorphism of $H$ and sends a negative representative to
one negative coset representative.

For $H=C_2$ this gives $C_4$ and $C_2^2$.
For $H=C_3$, the two possible conjugation actions give $C_6$
and $S_3$.  Explicitly the identity action allows the square
of $t$ to be removed using (B.41.4), while inversion forces
$A=1$.
For $H=C_4=\langle a\rangle$, the identity action has two
square classes, $A=1$ or $a$, giving $C_4\times C_2$ and
$C_8$.  Inversion allows $A=1$ or $a^2$, giving $D_8$ and
$Q_8$.
For $H=C_2^2$, the identity action gives $C_2^3$ or
$C_4\times C_2$ according as $A=1$ or $A\ne1$.
Every nonidentity order-two automorphism of $H$ is conjugate
to a transvection.  Its fixed subgroup equals the image of
$1+\alpha$, so (B.41.4) removes either allowable square.
This gives one further oriented extension, $D_8$, with positive
subgroup a four-group.  These arguments prove all the extension
counts $2,2,4,3$ used in the computation.

First take a Klein bottle with one cone of order $m$.
Its orbifold group is
\[
 \Lambda_m=
 \langle d_1,d_2\mid(d_1^2d_2^2)^m=1\rangle.
                                                        \tag{B.41.5}
\]
Put $y=d_1$, $z=d_1d_2$.  Then $y$ is negative, $z$ is
positive, and a positive local cone generator is
\[
                     x=(yzy^{-1}z)^{-1}.              \tag{B.41.6}
\]
A smooth orientation-compatible epimorphism to an oriented
finite group $G$ is therefore exactly a pair $(Y,Z)$ such that
$Y$ is negative, $Z$ is positive, $\langle Y,Z\rangle=G$,
and the image of (B.41.6) has order $m$.
Gromadzki--Szepietowski \cite[Lemma 3.3]{C_GromadzkiSzepietowski2016}
prove that the full outer automorphism group of $\Lambda_m$
is generated by
\[
 (y,z)\longmapsto(yz,z),\qquad
 (y,z)\longmapsto(y^{-1},z),\qquad
 (y,z)\longmapsto(y,z^{-1}).                           \tag{B.41.7}
\]
Their proof identifies these maps with the Dehn twist and
puncture and crosscap slides of the one-pointed Klein bottle.
The automorphism assertion concerns $\Lambda_m$ itself and
does not impose their later nonorientable-kernel hypothesis.
We apply that assertion to the orientation-compatible
epimorphisms above; every generator in (B.41.7) preserves the
orientation character.  Quotienting the finite set of pairs
also by $\operatorname{Aut}(G)$ preserving orientation is
therefore exactly topological conjugacy of the surface actions.

For $H=C_3$ the cone order is three.
In the $S_3$ extension a negative element inverts every
positive element, so (B.41.6) is the identity.
In $C_6$ it equals $Z^{-2}$; $Z$ must generate the positive
$C_3$.  The first move in (B.41.7) runs through the entire
negative coset, and the last move exchanges the two choices
of $Z$.  Thus there is one class, with full group $C_6$.

For $H=C_4$ the cone order is two.
The negative elements of $D_8$ and $Q_8$ invert $H$, again
making (B.41.6) trivial.  In the two abelian extensions it is
$Z^{-2}$, so $Z$ must generate $C_4$.
For $C_4\times C_2$, multiplication of $Y$ by powers of $Z$
sets $Y$ to the order-two generator of the other factor.
For $C_8$, an automorphism sets the odd generator $Y$ to a
chosen generator, and inversion of $Z$ identifies its two
remaining possibilities.  Each gives exactly one class.

For $H=C_2^2$, the elementary abelian extension $C_2^3$
cannot be generated by two elements.
In the abelian extension $C_4\times C_2$, every $Z\in H$
has square one, so (B.41.6) is trivial.
Write the remaining extension as
\[
 D_8=\langle r,s\mid r^4=s^2=1,\ srs=r^{-1}\rangle,
 \qquad H=\langle r^2,s\rangle .
\]
The admissible $Z$ are $s,r^2s$, and every negative $Y$
gives the cone image $r^2$.
An orientation-preserving automorphism exchanges the two
choices of $Z$; the first move in (B.41.7) changes a negative
reflection to a negative rotation, and inversion then
identifies $r$ and $r^{-1}$.  All eight pairs are one orbit.
This gives the final row in (B.41.1).

It remains to handle the two order-two cones in the first case
of (B.41.2).  For either extension $G=C_4$ or $C_2^2$, let
$h$ be the nonidentity positive element.
Both cone generators must map to $h$.
The two crosscap generators may independently map to either
negative element.  Every one of these four assignments is
surjective and respects the relation
\[
                      d_1^2x_1d_2^2x_2=1.
                                                        \tag{B.41.8}
\]
They are all equivalent.  Here is a geometric source move with
its exact word comparison, so no unproved transitivity is needed.
In a one-crosscap neighborhood of a cone, slide that cone once
around the one-sided core.  On the free generators $(d,x)$,
with exterior boundary word $d^2x$, the move is
\[
        d\longmapsto dx,\qquad
        x\longmapsto (dx)^{-1}x^{-1}(dx).              \tag{B.41.9}
\]
It preserves $d^2x$ exactly and sends the cone generator to a
conjugate of its inverse.  It is an automorphism: first replace
the free basis $(d,x)$ by $(dx,x)$, then invert and conjugate
the second member by the first.  Cutting the crosscap
neighborhood along its core gives the stated cone slide;
the unchanged exterior boundary permits extension by the
identity to the other crosscap.
In either abelian target, where $x$ has order two, this
replaces the image of $d$ by its product with $h$ and leaves
the cone image unchanged.  Apply it independently at the two
cones to connect all four assignments.  This gives one class
for each of $C_4$ and $C_2^2$.  Their different exponents
distinguish the two full groups.

For each surviving monodromy the local cone groups inject.
The orbifold covering subgroup is therefore torsion-free,
and it lies in the orientation kernel by construction.
The Klein-bottle orbifold Euler characteristic is
$-\sum_j(1-1/m_j)$; multiplying it by $|G|$ gives $-4$ in
all six cases.  The covering surface is consequently closed,
orientable, and of genus three.  This proves existence and
faithfulness of every listed action, as well as completeness.

There are no mirrors, so no reflection sections.
For each cone generator $x_j$ of prime order $m_j=2$ or $3$,
the points of its critical orbit are exactly
$G/\langle x_j\rangle$.  At the coset represented by $g$,
the stabilizer is $g\langle x_j\rangle g^{-1}$.  Write
$\varepsilon:G\to\{\pm1\}$ for orientation.  The multiplier
map $\mu_{j,g}$ on the section represented by that point is
\[
      \mu_{j,g}(gx_j^kg^{-1})=
                  \varepsilon(g)\frac{k}{m_j}\pmod{\Z}.
                                                        \tag{B.41.10}
\]
Every such point supplies one nontrivial partial section,
which is maximal; its centralizer in the surface group is
trivial by B.3.  No whole $G$ is a local stabilizer, so there
are no full sections.  The resulting coset counts are
exactly (B.41.1).  The program records the permutations of
these cosets, their stabilizers and the multiplier of every
local element.  In the $C_6$ row the two point generators
have multipliers $1/3$ and $-1/3$ under the fixed positive
generator; the negative coset exchanges the two.

The executed standard-library program
\path{catalogue_genus3_klein.py} independently enumerates
all pairs (B.41.3), checks associativity of their finite
multiplication tables and the orientation homomorphism,
enumerates every smooth generating monodromy, and computes
its orbit under the complete source and target automorphism
actions.  It retains the unsuccessful oriented extensions
with their empty monodromy lists, so the exclusions above
are part of the reproducible computation.
\end{proof}

\section{C. Equivariant comparison, drilling, and the topology of special groups}
\label{sec:part-c}

Throughout this part, $U=X\setminus S$ is connected, $X$ is a closed
orientable surface, and $S$ is a finite set of cardinality $\nu$. We put
$A=\operatorname{Homeo}(X,S)$ with its compact-open topology and
$K=\pi_0(A)$. Thus $A$ permits both permutations of $S$ and orientation
reversal; $A^0$ fixes $S$ pointwise and preserves orientation. We call the
type hyperbolic when $2g-2+\nu>0$. A superscript $u$ means the centralizer
of the image of the specified homomorphism $u:\Gamma\to A$. In particular,
$(A^0)^u=A^0\cap A^u$; it is not, by definition, the identity component
of $A^u$. We distinguish these two groups until connectedness is proved.

\subsection{The categories and their exact fidelity conditions}

\begin{claim}{C.1}{Refuted}{The unrestricted full target has too many arrows}
The literal comparison from surfaces and homeomorphisms to the full
$2$-category of peripheral groupoid diagrams and arbitrary functors is
not locally essentially surjective, already for $\Gamma=1$ and type
$(0,3)$. The correct target for a realization biequivalence is the
$2$-groupoid of such diagrams with equivalences as $1$-arrows.
\end{claim}
\passage{\S{}23; folios 134--136; scans 267--270}

\begin{proof}
For the three-punctured sphere put $\pi=\langle x,y\rangle$ and take the
peripheral diagram
\[
 \coprod_{i=1}^3 B\Z\longrightarrow B\pi,
 \qquad 1\longmapsto x,\ y,\ (xy)^{-1}.
\]
The identity on its three-object index set, the trivial homomorphism on
each copy of $\Z$, and the trivial homomorphism $\pi\to\pi$ define a
strictly commuting endomorphism of this diagram. Its map on
$\pi_{\rm ab}=\Z^2$ is zero. A homeomorphism induces an automorphism on
this abelianization. Natural isomorphisms between functors induce inner
conjugations, which act trivially on abelianization. Consequently this
endomorphism cannot become isomorphic to the image of a homeomorphism.
This is a counterexample to the literal fullness assertion, not to its
core version.

Here is the precise core version we use. An admissible surface is a
finite disjoint union of finite-type orientable surfaces, each genus-zero
component having at least one puncture, as on scan 264, folio 133. The
source $2$-groupoid has strict $\Gamma$-actions as objects, equivariant
homeomorphisms as $1$-arrows, and endpoint-preserving homotopy classes of
equivariant isotopies as $2$-arrows. An object of the target is a diagram
of groupoids over $B\Gamma$
\[
 D_\Gamma\longrightarrow U_\Gamma\longrightarrow B\Gamma.
\]
Its homotopy fibre over the base object is a peripheral diagram
$\coprod_{i\in I}B\Z\to B\pi$ with the surface presentation
$\prod_{j=1}^{g}[a_j,b_j]\prod_{i=1}^{\nu}c_i=1$; the $i$th map sends
$1$ to $c_i$. The fibre condition includes the common orientation module
and permits the noninjective map in type $(0,1)$ and the repeated image
in type $(0,2)$. This is the condition $M_{\rm lac}$. For $I=\varnothing$
we use the closed surface group, with $g\geq1$.

A target $1$-arrow is a pair of groupoid equivalences over $B\Gamma$
together with the invertible natural transformation expressing
commutativity of the peripheral square. A $2$-arrow is a pair of natural
isomorphisms, with identity projection to $B\Gamma$, satisfying that
square's compatibility. Composition is composition of functors with
pasting of these specified natural transformations. Equivalently these
are equivalences of the corresponding opfibrations over
$[1]\times B\Gamma$ that preserve the displayed fibres. The topological
comparison takes the fundamental groupoids of the homotopy quotients of
a $\Gamma$-invariant punctured collar system and of $U$. Two collar
systems have a common smaller system; the induced equivalences and
natural transformations specify the independence of this choice.
No arbitrary functor is silently replaced by an equivalence.
\end{proof}

\begin{claim}{C.2}{Proved}{The coset category and the quotient correction}
For a hyperbolic peripheral pair $(\pi;(L_i)_{i\in I})$, the category
$\mathscr C$ with objects $u,p_i$, automorphism groups
$\Aut(u)=\pi$, $\Aut(p_i)=1$, and
$\Hom(p_i,u)=\pi/L_i$ determines the pair up to an isomorphism and
simultaneous inner conjugation. Its equivalence groupoid has components
$\Out_{\rm lac}(\pi)$ and trivial automorphism groups. For an action with
extension $1\to\pi\to E\to\Gamma\to1$, the corresponding category over
$B\Gamma$ uses the bimodule $E/L_i$, with right action of
$E_i/L_i=\Gamma_i$, not $E/E_i$.
\end{claim}
\passage{\S{}23; folios 133--136; scans 264--270}

\begin{proof}
The stabilizer in $\Aut(u)$ of a selected arrow $p_i\to u$ is $L_i$.
Changing that arrow to $aL_i$ conjugates its stabilizer by $a$. An
equivalence fixes the unique object with nontrivial noncyclic
endomorphism group and permutes the other objects. Its restriction at
$u$ is an automorphism $\alpha$ of $\pi$. Equivariance of its maps on
Hom-sets says precisely that $\alpha(L_i)$ is conjugate to
$L_{\sigma(i)}$. Such a map of transitive sets is determined by a coset
whose stabilizer has this property. The equality
$N_\pi(L_i)=L_i$ makes that coset unique. A natural transformation between
two such equivalences has its component $a\in\pi$ at $u$, its identity
components at the $p_i$, and imposes $\alpha'=\operatorname{Int}(a)\alpha$
and the corresponding multiplication of all cosets by $a$. An
automorphism of an equivalence therefore lies in $Z(\pi)=1$.
This proves the assertion about the equivalence groupoid and gives both
comparison constructions explicitly.

For the geometric exit interpretation, the space of exits from $p_i$ to
$u$ is the homotopy fibre of the link map $BL_i\to B\pi$. Its components
are $\pi/L_i$ and its fundamental groups are the kernels of
$L_i\to\pi$; its higher homotopy groups vanish. The peripheral injection
makes those components contractible. Thus the ordinary category above
has exactly these exit mapping spaces. This also identifies the exact
failure of this ordinary model for the plane: the fibre of
$B\Z\to *$ is $B\Z$, not a discrete space. More explicitly the
Grothendieck category of $BL_i\to B\pi$ has automorphism groups
$L_i,\pi$ and forward Hom-set $\pi$. Its comparison with the exit
category is identity at the generic group, trivial at the special
endomorphism group, and $\pi\to\pi/L_i$ on forward arrows. Conversely
the translation groupoid of the generic action on the exit Hom-set
is $(\pi/L_i)//\pi\simeq BL_i$, with its map to $B\pi$. Thus the
peripheral diagram is recovered by this explicit link construction;
these two displayed categories are not silently identified. On
classifying spaces the forward comparison collapses each link to its
special point, realizing the filling map
$B\pi\to B\pi\cup_{\coprod BL_i}\coprod *$.

In the action category take one special object for each $\Gamma$-orbit
of punctures, choosing $i$ in that orbit. The generic object has
automorphism group $E$, the special object has automorphism group
$\Gamma_i$, and composition
on $E/L_i$ is left multiplication by $E$ and right multiplication by
$E_i/L_i$. These actions commute. The homotopy fibre over the base object of
$B\Gamma$ has special objects indexed by $\Gamma/\Gamma_i$ for each
chosen representative, and hence recovers every upstairs puncture. An
equivalent unskeletonized model uses all punctures and the nonempty
$\Gamma$-transporter sets between those in the same orbit. The
projection of a morphism $eL_i$
to $B\Gamma$ is the image $\bar e\in\Gamma$; it is well-defined because
$L_i\subset\pi$. It is not well-defined on $E/E_i$ when
$\Gamma_i\ne1$: representatives $e$ and $ea$, $a\in E_i$, can have
different images in $\Gamma$.

For a geometric example take $U=\mathbf P^1\setminus\{0,\infty,1,-1\}$,
$\Gamma=C_2$ acting by $z\mapsto-z$, and quotient coordinate $w=z^2$.
Then $E=\pi_1(\mathbf P^1\setminus\{0,1,\infty\})=F_2$. With $a$ a
loop about $w=0$, the upstairs inertia is $L_0=\langle a^2\rangle$
and $E_0=\langle a\rangle$. The right $C_2$ action on
$E/\langle a^2\rangle$ is nontrivial, whereas replacing it by
$E/\langle a\rangle$ erases that action. This is the exact morphism
obstruction to the quotient-category inference in the prompt.
\end{proof}

\begin{claim}{C.3}{Proved}{The four levels and the coefficient ladder}
Let $A$ be a locally path-connected topological group whose components
are $1$-types. For the comparison from strict actions to weak actions in
the categorical group $\Pi_1A$, the first three fidelity conditions are,
on each comparison $M\to N$ of transport spaces,
\[
\begin{array}{c|l}
0&\pi_1M\to\pi_1N\text{ injective at every point},\\
1&\pi_0M\to\pi_0N\text{ injective and }\pi_1M\to\pi_1N\text{ isomorphic},\\
2&\pi_0M\to\pi_0N\text{ bijective and }\pi_1M\to\pi_1N\text{ isomorphic}.
\end{array}
\]
Level $3$ adds essential surjectivity on action objects. When $A^0$ is
simply connected, these give exactly scans 279--280: vanishing of
$\pi_1((A^0)^u)$; connectedness of $(A^0)^u$; conjugacy of two lifts of the
same $\Gamma\to K$ by an element of $A^0$; and existence of a lift.
\end{claim}
\passage{\S{}23; folios 136--137; scans 271--273; \S{}23; folios 140--141; scans 279--281}

\begin{proof}
The source Hom-groupoid is $\Pi_1M$. A functor of groupoids is faithful
exactly when it is injective on the automorphism group at each object.
It is fully faithful exactly when that condition is an isomorphism and
two source objects whose images are connected are already connected.
It is an equivalence exactly when it is fully faithful and every target
component is hit. Applying these three elementary tests proves the
table, including the cumulative meaning of each level.

When $\pi_1(A^0)=0$, the categorical group $\Pi_1A$ is equivalent to the
discrete group $K$. For strict homomorphisms $u,v:\Gamma\to A$ the map
of transport spaces is
\[
 \{a\in A:a u(\gamma)a^{-1}=v(\gamma)\ (\gamma\in\Gamma)\}
 \longrightarrow
 \{k\in K:k\bar u(\gamma)k^{-1}=\bar v(\gamma)\}.
\]
A nonempty source is a torsor under $A^u$. Its fibres over points of the
target are torsors under $(A^0)^u$. The first two lines of the table are
therefore the first two conditions stated. For the third, given a target
transporter $k$, choose $a\in A$ representing $k$. The two actions
$a u a^{-1}$ and $v$ have the same projection to $K$; conjugating them
by an element of $A^0$ produces the prescribed transporter. Conversely,
apply transporter surjectivity to the identity element of $K$. This
also verifies the stronger marginal correction on scan 280.

For comparison with the GALOIS ladder, first take the $1$-truncations
of the transport spaces; fixed spaces are not assumed to inherit the
ambient group's higher truncation. On a map of $1$-types $M\to N$, set-valued constant coefficients test bijectivity on $\pi_0$;
full faithfulness of restriction tests that bijectivity together with
surjectivity on $\pi_1$; equivalence tests bijectivity on $\pi_0$ and an
isomorphism on $\pi_1$. For groupoid-valued coefficients all three tests
coincide with a weak equivalence, because $\pi_iM=\pi_iN=0$ for $i\geq2$;
the same is true for $2$-type coefficients. These assertions follow by
applying the long exact sequence to the homotopy fibres: an $r$-connected
fibre gives isomorphisms through degree $r$ and surjectivity in degree
$r+1$. Thus level $2$ is the $r=0$ equivalence test on the transport
spaces. Levels $0$ and $1$ do not equal a rung in that ladder: their
$\pi_1$-injectivity and $\pi_0$-injectivity tests have been written out
above, whereas each rung of the coefficient ladder demands a component
bijection. Level $3$ concerns objects rather than this map of Hom-spaces.
This is a comparison of the specified maps, not an identification based
only on the words ``faithful'' and ``equivalence''.
\end{proof}

\begin{claim}{C.4}{Proved}{Coherent actions and rectification on spaces}
For a discrete group $P$, the space of self-homotopy equivalences of
$BP$ has components indexed by $\Out(P)$, each with homotopy type
$BZ(P)$. In particular, for the centre-free surface groups the map from
outer actions to coherent actions is an equivalence of action groupoids.
Every coherent action of a discrete group on a space has a strict
realization on a space of the same homotopy type. For $P=\Z$ or $\Z^2$,
the obstruction class in degree $3$ is zero and the coherent actions
with a specified linear action are classified by $H^2(\Gamma,P)$.
\end{claim}
\passage{\S{}23; folios 131--132; scans 261--263; \S{}23; folios 135--141; scans 268--281}

\begin{proof}
The groupoid of functors $BP\to BP$ has homomorphisms $P\to P$ as
objects and conjugating elements of $P$ as arrows. Its equivalence
subgroupoid has automorphisms as objects. The stabilizer of an
automorphism is $Z(P)$, and the orbits are $\Out(P)$. The nerve computes
the mapping space because nerves of groupoids represent $1$-types.
Delooping the resulting categorical group has $\pi_1=\Out(P)$,
$\pi_2=Z(P)$ and no higher homotopy groups. Its Postnikov class is the
usual extension obstruction in $H^3(\Out(P);Z(P))$. For $P$ abelian the
map $\Aut(P)\to\Out(P)$ is the identity and gives a strict section;
the categorical group is $BP\rtimes\Aut(P)$, so this Postnikov class
vanishes. A choice of coherence is then a $2$-cocycle with values in the
specified $\Gamma$-module $P$; changing it by a coherent equivalence
changes the cocycle by a coboundary. This proves the stated
classification without asserting a nonzero obstruction for either
abelian surface group.

For rectification, represent a coherent action by a fibration
$p:Y\to B\Gamma$ with the prescribed homotopy fibre. Pull it back along
$E\Gamma\to B\Gamma$. The pullback $Y'=E\Gamma\times_{B\Gamma}Y$ has a
strict deck action of $\Gamma$. Its projection to the contractible
space $E\Gamma$ has fibre the original homotopy type, so the inclusion
of any chosen fibre is a homotopy equivalence. Its homotopy quotient
recovers $Y\to B\Gamma$. This construction applies simultaneously to
peripheral diagrams by the diagrammatic pullback. It produces a space
model; the additional demand that the model be the original surface is
exactly the realization problem addressed below. For continuous
profinite actions the objects and the homotopy quotient must instead be
specified in a continuous/profinite homotopy theory; the preceding
construction makes no untyped transfer to that setting.
\end{proof}

\subsection{Finite actions on hyperbolic surfaces}

\begin{claim}{C.5}{Proved}{Existence, with punctures and orientation reversal}
For a hyperbolic type and every finite group $\Gamma$, each
homomorphism $\rho:\Gamma\to K$ lifts to a homomorphism
$u:\Gamma\to A$. Its image can be chosen to consist of isometries of a
complete finite-area hyperbolic metric on $U$.
\end{claim}
\passage{\S{}23; folios 136, 140--141; scans 270--271, 279--280; \S{}24; folio 147; scan 293}

\begin{proof}
Let $H=\rho(\Gamma)$, a finite subgroup of the extended mapping class
group. Kerckhoff's Theorem 5, together with the explicit extension in
\S IV, pp.~255--256, of the original paper \cite{C_Kerckhoff1983}, gives
a fixed hyperbolic structure for $H$. That extension explicitly permits
cusps; the paper's mapping-class convention explicitly includes
orientation-reversing maps. We use the original paper, rather than
inferring the punctured case from the announcement.

Choose the isometry in each class of $H$. It is unique. Indeed an
isometry isotopic to the identity has a lift to $\mathbb H^2$ commuting
with the surface deck group. That deck group contains two hyperbolic
elements with different endpoint pairs, so an isometry commuting with
both fixes at least three points of the ideal boundary and is the
identity. The same argument covers a finite-area punctured surface;
its non-elementary deck group again contains those elements. Products
of the chosen isometries represent products in $H$, and uniqueness
therefore makes the choices a homomorphism $H\to A$. Composing with
$\rho$ proves the claim. The isometries extend to the end compactification
and permute $S$, so the target really is the specified group $A$.
\end{proof}

\begin{claim}{C.6}{Proved}{Uniqueness up to the identity component}
Let $u,v:\Gamma\to A$ be homomorphisms with $\Gamma$ finite, $U$
hyperbolic, and $\bar u=\bar v:\Gamma\to K$. There is an
$h\in A^0$ such that $h u(\gamma)h^{-1}=v(\gamma)$ for every $\gamma$.
This holds for topological actions without a smoothness assumption.
\end{claim}
\passage{\S{}23; folios 136, 140--141; scans 270--271, 279--280}

\begin{proof}
We first justify the analytic model for either action. For a point
$x\in U$ and its finite stabilizer $H$, lift each element of $H$ to the
universal cover with the chosen lift of $x$ fixed. Uniqueness of a lift
with a specified value makes these lifts a finite group. Compactifying
$\widetilde U\simeq\R^2$ by one point gives a finite action on $S^2$
with two common fixed points. Kolev's compact-sphere theorem
\cite[Theorem 1.2]{bKolev2006} conjugates that action into $O(3)$.
A subgroup of $O(3)$ fixing two distinct points is contained in the
orthogonal group of the perpendicular plane; choosing a sufficiently
small round disc at the first point and pulling it back yields a common
invariant disc with an $O(2)$ action. Choosing it small enough makes its
translates by distinct cosets of $H$ disjoint. At a puncture we apply the
same finite-stabilizer disc argument on the compactification, as in the
finite-disc theorem proved from the same source.

The quotient consequently has ordinary charts, cone charts, mirror
charts, and corner charts with their finite cyclic or dihedral local
groups. A finite triangulation respecting these strata supplies a
smooth orbifold structure: smooth each triangle, use collar charts
along its edges, and use the displayed linear quotient charts at its
vertices. Pulling the charts to $X$ makes the original finite action
smooth for a compatible smooth structure. Average a metric and use its
conformal class. Uniformization of the punctured conformal surface gives
a finite-type structure for which each operator is conformal or
anticonformal. Thus we have conformal structures $c_u,c_v$ on the same
marked topological surface.

The extremal quasiconformal map $h:(U,c_u)\to(U,c_v)$ in the identity
isotopy class is unique, by the finite-type Teichm\"uller theorem
\cite[\S3(A)--(D)]{C_Bers1981}. For each $\gamma$, the map
$v(\gamma)^{-1}h u(\gamma)$ is orientation preserving: $u(\gamma)$ and
$v(\gamma)$ have the same orientation sign. Pre- and postcomposition by
these conformal or anticonformal maps leaves the maximal dilatation
unchanged. The equality of the two mapping-class homomorphisms says
that this map is isotopic to the identity, with the punctures retained.
Uniqueness therefore gives $v(\gamma)^{-1}h u(\gamma)=h$. The
quasiconformal map extends across the finite punctures; hence
$h\in A^0$ and has the asserted conjugacy property.
\end{proof}

\begin{claim}{C.7}{Proved}{The finite centralizer component, including mirrors}
For a finite action $u$ on a hyperbolic $U$, the group $(A^0)^u$ is
contractible. In particular it is the identity component of $A^u$.
\end{claim}
\passage{\S{}\S{}23--24; folios 140--144; scans 279--286}

\begin{proof}
Replace $\Gamma$ by its image $H$. The preceding local construction
makes $H$ a group of hyperbolic isometries. Let $\pi$ be the surface deck
group and let $E$ be the full group of lifts of $H$ to $\mathbb H^2$.
It is a discrete group of isometries, possibly with reflections; its
quotient $Q$ is a finite-type hyperbolic orbifold. Write $D$ for descent
of a centralizing homeomorphism to the labelled underlying surface of
$Q$, with cone points, corner points, mirror boundary, and punctures
retained.

We give the injectivity argument at the place where an
orientation-preserving-only Birman--Hilden theorem would not suffice.
For $f\in(A^0)^u$, the identity isotopy supplies a lift $\widetilde f$
that commutes with $\pi$. Since $f$ commutes with $H$, conjugation by
$\widetilde f$ induces an automorphism $\alpha$ of $E$ that is the
identity on $\pi$. For $e\in E$ and $a\in\pi$,
\[
 \alpha(e)a\alpha(e)^{-1}=\alpha(eae^{-1})=eae^{-1}.
\]
Thus $e^{-1}\alpha(e)$ centralizes $\pi$, and hence is the identity by
the distinct-endpoint argument in C.5. Consequently $\alpha=\id_E$.
The orbifold Dehn--Nielsen--Baer injectivity theorem, with its mirror
and corner structure, makes $D(f)$ isotopic to the identity through
homeomorphisms preserving those structures; the relevant precise form
is \cite[Theorem 2 and Proposition 3]{C_KolbeEvans2021}. At the algebraic
step in their Lemma 3 we have used the calculation just written:
centralizers of individual reflections need not be cyclic, but the
centralizer of the non-elementary surface subgroup is trivial. Thus
that inapplicable sentence of their proof is not being used. There is an orientation hypothesis in that source when the labelled
underlying quotient is orientable, and we verify it here. The canonical
$\pi$-equivariant lift $\widetilde f$ extends by the identity on
$\partial_\infty\mathbb H^2$. For a compact surface this follows from
bounded displacement on a compact fundamental domain. For a punctured
surface use a compact truncation and disjoint invariant horodiscs:
displacement is bounded on the truncated part, and preservation of each
end sends sufficiently deep cusp horodiscs into the corresponding cusp
horodiscs. A sequence escaping through distinct horodiscs has the same
ideal limit before and after the map, since their visual diameters tend
to zero; a sequence in one horodisc has its cusp endpoint as its only
ideal limit. These observations prove the claimed continuous extension.
Because $\alpha=\id$, $\widetilde f$ preserves each reflection axis.
It fixes that axis's two ideal endpoints individually, so preserves its
linear orientation, and, being orientation preserving on $\mathbb H^2$,
preserves its two sides. It therefore preserves every mirror chamber.
On such a chamber the quotient orientation is the one induced from
$\mathbb H^2$; hence $D(f)$ preserves it. Without mirrors an orientable
quotient is covered orientably and the same assertion is immediate.
Thus the theorem is applied inside its stated homeomorphism group.

Let $\mathcal H(Q)^0$ denote the identity component of that labelled
homeomorphism group. Every path starting at its identity lifts uniquely,
starting at the identity upstairs, to a path centralizing $H$. Lift
first on the regular covering locus. The local cyclic and dihedral
models extend the lift over cone points, corners and mirrors; continuity
of the extended lift follows in the charts from the finite quotient
map. The same construction for a parameter family proves the path and
homotopy lifting properties. The only ambiguity for an endpoint is a
central deck transformation. Such a deck transformation cannot belong
to $A^0$ unless it is the identity, by C.5's isometry argument. It follows
that descent identifies $(A^0)^u$ with $\mathcal H(Q)^0$, as soon as the
latter is simply connected; the construction is continuous in local
isotopy charts, so this is an isomorphism of topological groups.

We check that this group is contractible, including the small underlying
quotients. End-compactify the quotient and retain its finitely many
interior and boundary marked points. Yagasaki's compact-pair Lemma 2.1
\cite{A_Yagasaki2000} gives contractibility outside its explicit
exceptional pairs. A marked point on a boundary component is covered
by Lemma 2.1(ii). With no such point, the exceptional disk has at most
one interior marked point, the annulus and M\"obius band have none, the
sphere has at most two, the projective plane at most one, and the torus
and Klein bottle have none. Each of these has nonnegative orbifold
Euler characteristic: an interior cone subtracts $1-1/m$, an interior
puncture subtracts $1$, and a corner or boundary puncture would already
be a boundary marked point. But
$\chi_{\rm orb}(Q)=\chi(U)/|H|<0$. None occurs. The homeomorphism group
with labels has the same identity component as the group fixing all
these finite marks pointwise. Yagasaki's theorem therefore applies.
This proves contractibility, and also supplies the simple connectedness
used in the lifting construction.
\end{proof}

\begin{claim}{C.8}{Proved}{The finite core biequivalence and homotopy fixed points}
On hyperbolic surfaces, for every finite $\Gamma$ the corrected
comparison of C.1 has all four fidelity levels and is a biequivalence.
For a specified action $u$, the comparison
\[
 A^u\longrightarrow A^{h\Gamma}
      =\operatorname{Map}_\Gamma(E\Gamma,A),
 \qquad a\longmapsto(\text{constant map with value }a),
\]
where $\Gamma$ acts on $A$ by conjugation through $u$, is a weak homotopy
equivalence. In standard CW models it is a homotopy equivalence.
\end{claim}
\passage{\S{}23; folios 136--141; scans 270--281}

\begin{proof}
The nonequivariant peripheral comparison identifies $K$ with
$\Out_{\rm lac}(\pi)$ and has trivial $2$-automorphism groups, by the
peripheral Dehn--Nielsen--Baer theorem
\cite[Theorems 8.1 and 8.8]{A_FarbMargalit2012} and $Z(\pi)=1$.
C.7 gives levels $0$ and $1$, C.6 gives level $2$ by the transporter
argument of C.3, and C.5 gives level $3$. These arguments also treat a
nonfaithful homomorphism: its finite image is the action group, and a
finite-order homeomorphism in $A^0$ is the identity. Disconnected
admissible surfaces reduce to representatives of the finitely many
orbits of connected components and their stabilizers, with the stated
transporters providing the arrows between choices. This is the desired
biequivalence on its hyperbolic part.

The map $A\to K$ is $\Gamma$-equivariant. Its fibres are contractible by
the same homeomorphism-group theorem with $H=1$. Taking derived
homotopy fixed points gives a weak equivalence
$A^{h\Gamma}\to K^\Gamma=C_K(\bar u(\Gamma))$, the latter a discrete
space. The constant-map comparison followed by this map is exactly
$\pi_0(A^u)\to C_K(\bar u(\Gamma))$. Its surjectivity follows by applying
C.6 to $a u a^{-1}$ and $u$ for any representative $a$ of a centralizing
mapping class; its injectivity and its contractible fibres are C.7.
Thus the displayed map is a weak equivalence. The source components
have CW homotopy type by the ANR statement for quotient homeomorphism
groups in \cite{A_Yagasaki2000}; taking the usual CW model of the
derived target turns it into a homotopy equivalence. We state the model
convention because scan 281 writes ``homotopisme'' without specifying
an infinite-dimensional mapping-space model.
\end{proof}

\subsection{The exceptional types and infinite operator groups}

\begin{claim}{C.9}{Proved}{Corrected abelian tests and a torus calculation}
For $\pi_1(A^0)\ne0$, replacement of the four tests by connectedness and
uniqueness of lifts of $\Gamma\to K$ loses the coherent action data.
The valid tests are C.3 with weak transport spaces. At a specified
strict action the target automorphism space is the $1$-type
$A^{h\Gamma}$ and
\[
 \pi_1(A^{h\Gamma},1)=\pi_1(A,1)^\Gamma.
\]
For $A=\operatorname{Homeo}(T^2)$ and $u(C_2)=\{1,-1\}$, the group
$(A^0)^u$ has four components, all contractible, whereas the right-hand
invariant fundamental group is zero. Thus the naive connectedness test
fails; the corrected comparison does not fail in this example.
\end{claim}
\passage{\S{}23; folio 141; scans 280--281}

\begin{proof}
The homotopy-fixed-point groupoid of a $1$-type has loops equal to the
fixed automorphisms of its base object, proving the displayed equality.
Its components include the $H^1(\Gamma;\pi_1A)$ coherence data, with
possible obstruction for an individual component of $A$. Hence it is
not legitimate to test only $\pi_0A$ when $\pi_1A\ne0$.

For the torus, evaluation at $0$ gives
$A^0\cong T^2\times A^0_{\{0\}}$ by
$f\mapsto(f(0),t_{-f(0)}f)$. The second factor is contractible by the
homeomorphism-group theorem for a torus with a marked point. In the
centralizer of $-1$, $f(0)$ is one of the four points of $T^2[2]$.
Each is realized by translation. In the fibre $f(0)=0$, regarding the
torus as once punctured makes it hyperbolic; the fibre is contractible
by C.7. This proves the four-component assertion. Conjugation by $-1$
acts as $-I$ on $\pi_1(A^0)=\Z^2$, so its invariants vanish.
The inclusion of translations $T^2\hookrightarrow A^0$ is equivariant
and a weak equivalence; its homotopy fixed points have components
$H^1(C_2;\Z^2_{-})=(\Z/2)^2$ and zero higher homotopy groups. The four
translations just identified realize exactly these four components.

For the sphere, the peripheral groupoid diagram is empty over $*$ and
cannot encode its homotopy type; $\pi_2(S^2)=\Z$ already records the
failure of asphericity. The plane, annulus and torus have
$A^0\simeq S^1,S^1,T^2$, respectively; in the plane the peripheral
map $B\Z\to *$ retains the circle of automorphism $2$-arrows, and in the
annulus the two ends remain indexed even though their subgroup images
coincide. These exact maps explain why collapsing to a bare outer group
would erase data. The complete finite-action comparison for these models, including
its parametric centralizer proof, is C.10.
\end{proof}

\begin{claim}{C.10}{Proved}{Every finite action in the exceptional aspherical types}
For $U=\R^2$, $S^1\times\R$, or $T^2$, and every finite strict action,
the fixed-to-homotopy-fixed comparison in C.8 is a weak equivalence
and is a homotopy equivalence on CW models. The corrected peripheral
$2$-groupoid comparison has all four fidelity levels for these types.
The tests are those of C.3; ordinary outer actions alone do not retain
the required coherence classes.
\end{claim}
\passage{\S{}23; folios 136--141; scans 270--281}

\begin{proof}
We prove the parametric assertion, not merely classification of individual
actions. Write respectively
\[
 G=O(2),\qquad
 G=(\R/\Z)\rtimes(C_2\times C_2),\qquad
 G=(\R^2/\Z^2)\rtimes\operatorname{GL}_2(\Z).
\]
The annulus element $(a,\epsilon,\eta)$ sends
$(\theta,t)$ to $(\epsilon\theta+a,\eta t)$; the first sign acts on the
circle factor and the second on the two ends. In every case $G\to A$
is a weak equivalence. Here is a direct check of the identity components.
On the torus, evaluation at $0$ splits $A^0$ as $T^2$ times its
point-fixing subgroup. An identity isotopy can be made point fixing by
postcomposing at time $s$ with translation by the negative of the trace
of $0$. Thus that subgroup is its own identity component, and is
contractible by the punctured-torus case of the homeomorphism-group
theorem \cite[Theorem 1.1 and Lemma 2.1]{A_Yagasaki2000}. The same adjustment on the annulus uses angular rotation and
translation of its real coordinate. Its point-fixing subgroup is the
identity group of a hyperbolic three-ended sphere and is contractible,
so $A^0\simeq S^1$. On the plane first fix $0$ by translations and
then fix $1$ using multiplication by a nonzero complex number. The
remaining identity group fixes $0,1,\infty$ and is contractible; hence
$A^0\simeq S^1$ with generator the rotation circle. The component
groups are respectively $C_2$, $C_2\times C_2$, and
$\operatorname{GL}_2(\Z)$: orientation and end signs give the first
two, and the torus action on its lattice gives the third. Cutting the
plane or annulus along a ray or spanning arc verifies the stated
kernels; the torus case is the genus-one Dehn--Nielsen--Baer theorem.

By the simultaneous finite normal forms proved in B.1 and B.14,
conjugate the given finite image $H$ into $G$. We will prove
\[
                    C_G(H)\longrightarrow C_A(H)                 \tag{C.10.1}
\]
is a weak equivalence by examining its fibres. Conjugation by the chosen
homeomorphism transports both the fixed and derived-fixed comparisons,
so this reduction is an isomorphism of the actual diagrams.

\emph{Torus displacement.}
For $f\in C_A(H)$ let $M\in\operatorname{GL}_2(\Z)$ be its lattice
action and lift it as $F(x)=Mx+d(x)$, with $d$ periodic. Put
$b=\int_{T^2}d(x)\,dx$ and $g(x)=Mx+b$ modulo $\Z^2$.
Changing the lift changes $b$ by an integral vector and leaves $g$
unchanged. For $\gamma(x)=A_\gamma x+t_\gamma$ the commutation equation
is
\[
 d(A_\gamma x+t_\gamma)-A_\gamma d(x)
       =(I-M)t_\gamma+k_\gamma,\qquad k_\gamma\in\Z^2.
\]
Integration preserves its left-hand change of variables, because the
affine operator preserves Haar measure. It gives the same equation
with $d$ replaced by $b$. Thus $g$ commutes with $H$ and
$g^{-1}f\in A^0$ has a unique lift of zero mean displacement.
Consequently multiplication is a homeomorphism of spaces
\[
 C_G(H)\times\mathcal K\ \xrightarrow{\sim}\ C_A(H),
\]
where $\mathcal K$ consists of these zero-mean maps. The averaging
map need not be a group homomorphism; only this explicit space
factorization is used. A zero-mean lift commutes exactly with every
lift of $H$: the possible integral commutator above has zero average
and is therefore zero.

First divide by the finite translation subgroup $N=H\cap T^2$.
If $a$ represents a member of $N$ of order $m$, commutation implies
$F(x+a)=F(x)+a+k$. Iterating $m$ times and using $ma\in\Z^2$
forces $k=0$ for a zero-mean identity lift. Hence its displacement is
periodic for the enlarged covering lattice. Descent and lift identify
$\mathcal K$ with the same zero-mean space on $T'=T^2/N$ for $H/N$.
This quotient group has faithful linear part, a finite subgroup $P$
of $\operatorname{GL}_2(\Z)$.

When $P$ contains a nonidentity positive operator, its positive subgroup
$K$ is cyclic, generated by an affine rotation with $I-A$ invertible.
It has a nonempty finite fixed set $R$ on $T'$, stable under $H/N$.
A zero-mean lift commutes with each lift of that rotation and fixes its
unique Euclidean fixed point; thus it fixes every point of $R$.
Choose $p\in R$. The point-fixing argument above places our map in
$\operatorname{Homeo}(T',p)^0$. C.7, applied to the cyclic action on
$T'\setminus\{p\}$, supplies an equivariant identity isotopy. This
isotopy fixes all of $R$, since $R$ is finite. Our map therefore belongs
to $\operatorname{Homeo}(T',R)^0$. Apply C.7 again, now to the full
$H/N$ action on the hyperbolic $T'\setminus R$. Its identity centralizer
is exactly $\mathcal K$: conversely an equivariant identity isotopy
relative to $R$ has commuting identity lifts, whose mean is fixed by
$A$ and hence is zero. This proves contractibility of $\mathcal K$.

The remaining linear groups are $P=1$ and a group of order two generated
by a reflection. For $P=1$, the space factorization on the identity
component is $\operatorname{Homeo}(T')^0\cong T'\times\mathcal K$.
The first factor is a section of the mean map. Since the left side has
the homotopy type $T'$, the mean map is an isomorphism on $\pi_1$ and
$\mathcal K$ is weakly contractible. It is a retract of the ANR on the
left, so it has CW homotopy type and is contractible.

For a reflected or glide-reflected torus the underlying quotient is an
annulus, a M\"obius band, or a Klein bottle. This exhausts the cases:
the action reverses orientation, has no isolated rotation point, and
the quotient has Euler characteristic zero; its boundary consists of
mirrors. A zero-mean lift $F$ commutes exactly with the Euclidean orbifold
group $E$. Each lifted reflection line is preserved, and bounded
displacement prevents reversal of its infinite directions. In
particular the compact annulus quotient's two boundary circles cannot
be interchanged. The explicit straight homotopy
\[
                  F_s(x)=(1-s)F(x)+sx
\]
is $E$-equivariant, because every element of $E$ is affine. On each
reflection line it stays on that same line. It therefore descends to
a homotopy of the quotient as a surface with boundary, preserving its
boundary throughout. Its quotient domain is compact, so it is proper.
Epstein's Theorem 6.3, with the topological Appendix
\cite{C_Epstein1966}, applies with exactly this proper,
boundary-preserving homotopy and gives an identity isotopy. This also
handles the boundary-free Klein bottle; Theorem 5.8 of that source
separately verifies the M\"obius-band case. Thus no unverified inference
from an abstract outer automorphism to an exceptional-surface isotopy
is needed here.

Descent consequently puts $\mathcal K$ in the identity group of the
labelled quotient. Every identity isotopy there lifts; among the two
lifts exactly one preserves the orientation upstairs. This identifies
the identity component of $C_{A^0}(H/N)$ with the quotient identity
homeomorphism group. By Yagasaki's actual exception list its homotopy
type is $S^1$ in all three cases. The mean factorization on that component
is $T'^P_0\times\mathcal K$, with $T'^P_0\cong S^1$. Here
$P$ acts linearly on the translation torus; this fixed translation
circle is present even for an affine glide with no fixed point.
The section by translations makes the mean map surjective on $\pi_1$;
a surjection $\Z\to\Z$ is an isomorphism. Hence $\mathcal K$ is
weakly contractible, and is contractible by the same ANR-retract argument.
This completes (C.10.1) for the torus, including free glide reflections.

\emph{Annulus displacement.}
Lift $f\in C_A(H)$ to
$F(\theta,t)=(\epsilon_f\theta+d(\theta,t),F_2(\theta,t))$ and put
$b=\int_0^1d(\theta,0)\,d\theta$. Define
$g(\theta,t)=(\epsilon_f\theta+b,\eta_f t)$, where $\eta_f$ is its
end sign. The commutation equation, integrated on the invariant core
$t=0$, proves that $g$ commutes with $H$. It gives the same continuous
factorization $C_A(H)\cong C_G(H)\times\mathcal K$, with zero-mean
identity lifts in $\mathcal K$. Divide first by the finite rotation
subgroup. The resulting faithful sign group is a subgroup of
$C_2\times C_2$.

When it contains $(-1,-1)$, that positive involution has two fixed
points. Its subgroup is normal. A zero-mean lift fixes each of the
corresponding fixed points in $\R\times\R$. The same two applications
of C.7 as in the torus argument, first marking one fixed point and then
the whole invariant fixed set, identify $\mathcal K$ with a contractible
centralizer on a hyperbolic punctured annulus. An equivariant identity
lift in that punctured model has mean zero because its angular linear
sign is $-1$, proving equality in both directions.

A subgroup without $(-1,-1)$ is trivial, generated by $(-1,1)$, or
generated by $(1,-1)$. In the first case the factorization and
$A^0\simeq S^1$ give contractibility as above. An angular reflection
has a strip as quotient, with its two mirror boundary lines and two
ends. Exact commutation of a zero-mean lift preserves both lines and
both ends. The same straight homotopy in $\R\times\R$ is equivariant
for all the affine covering operators and preserves the boundary
lines. It is proper on the strip quotient: for $t\to+\infty$ both
$t$ and $F_2(\theta,t)$ tend uniformly to $+\infty$ on a compact
angular fundamental interval, and for $t\to-\infty$ both tend to
$-\infty$. Every convex combination has the same uniform property.
Epstein's proper boundary-homotopy theorem therefore places the
descended map in the identity group. That group is contractible by
Yagasaki, so $\mathcal K$ is contractible.
Finally a radial reflection or radial glide reflection has quotient a
half-open annulus or a punctured projective plane, respectively. The
straight homotopy is again equivariant; it preserves the fixed mirror
line in the reflection case and has the just-verified properness at
the quotient's end. Epstein's Theorem 6.3, or Theorem 5.8 for the open
M\"obius band, gives the required identity isotopy. Descent and isotopy
lifting identify the upstairs centralizer identity component with that
group. Yagasaki gives $S^1$ in both cases. The mean factorization has a
section by the fixed angular circle, hence again forces
$\mathcal K$ to be contractible. These cases include every affine
finite action and retain the end permutation throughout.

\emph{Plane centralizers.}
For a cyclic rotation group of order $n>1$, a commuting map fixes $0$.
Descent through $z\mapsto z^n$ is a covering of the orientation-preserving
centralizer onto $\operatorname{Homeo}^+(\R^2,0)$, with deck group
$C_n$. The base is connected with homotopy type $S^1$. A loop of one
rotation downstairs lifts from the identity to the deck rotation by
$2\pi/n$. Its monodromy is therefore onto $C_n$, so the covering
centralizer is connected and the upstairs rotation circle induces the
isomorphism on $\pi_1$; all higher homotopy groups vanish. A negative
map centralizes the rotation only for $n=2$, as its action on the
punctured-plane fundamental group inverts the deck generator. In that
case conjugation supplies the second component. These are precisely
the components and circle types of $C_{O(2)}(C_n)$.

For a reflection, descent has half-plane quotient. A commuting map has
two independent signs: its orientation on the fixed line and its
orientation on the plane. The four diagonal orthogonal maps realize
them. Each fibre is the identity homeomorphism group of the half-plane,
which is contractible by Yagasaki. For a dihedral group with rotation
order $n>2$, a centralizing map preserves plane orientation and each
reflection line. Preservation of cyclic order of their rays permits
only the identity permutation or simultaneous reversal of every ray;
these are realized by $I$ and $-I$. For $n=2$ the two perpendicular
reflection lines permit four independent ray-sign choices, realized by
the four diagonal orthogonal maps. After dividing by the matching map,
descent is to a sector homeomorphism fixing its vertex and both boundary
rays. This group is contractible: the sector is a half-plane with a
marked boundary point, and its identity group is covered by Yagasaki's
boundary-point case. The boundary rays also make the lift unique.
The trivial group case was verified at the start. This proves (C.10.1)
for every finite subgroup of $O(2)$.

It remains to identify the derived comparison, including all components.
The following elementary torus calculation does so for each of our
three $G$. Let $T=V/\Lambda$, with $V$ real and $\Lambda$ a lattice,
and let a finite group act on a $T$-torsor by affine maps. Its translation
part is a cocycle $b\in Z^1(\Gamma,T)$, where $T$ and $\Lambda$
have the specified linear $\Gamma$-module structures. The homotopy-fixed space is
the section space of its Borel fibration. Its obstruction is the class
$\delta[b]\in H^2(\Gamma,\Lambda)$ of the extension obtained by lifting
the affine maps to $V$. Averaging cochains gives
$H^j(\Gamma,V)=0$ for $j>0$, so
$\delta:H^1(\Gamma,T)\to H^2(\Gamma,\Lambda)$ is an isomorphism.
Thus a homotopy fixed point exists exactly when $b$ is a coboundary,
which is exactly existence of an actual fixed point. After choosing one,
\[
 \pi_0(T^{h\Gamma})=H^1(\Gamma,\Lambda),\qquad
 \pi_1(T^{h\Gamma})=\Lambda^\Gamma,
\]
and there are no higher groups. The exact sequence
\[
 0\to\Lambda^\Gamma\to V^\Gamma\to T^\Gamma
       \to H^1(\Gamma,\Lambda)\to0
\]
identifies the actual fixed components and their fundamental groups
with exactly these groups. The constant-map comparison induces these
identifications. Apply this to every invariant component of $G$, each
of which is just such an affine torus torsor. Noninvariant components
contain neither fixed nor homotopy-fixed points. It follows that
$G^\Gamma\to G^{h\Gamma}$ is a weak equivalence. The equivariant weak
equivalence $G\to A$, together with (C.10.1), now proves the claimed
comparison for $A$ by the commutative square of fixed and derived-fixed
spaces.

Finally $BG$ is the split $2$-type with fundamental group
$\pi_0G$, second group $\Lambda=\pi_1G$, and zero Postnikov invariant.
Coherent finite actions with a specified homomorphism to $\pi_0G$ are
therefore classified by $H^2(\Gamma,\Lambda)$. The same connecting
isomorphism realizes each by a $T$-valued cocycle, hence by a strict
homomorphism $\Gamma\to G$. It also identifies coherent equivalences
in the identity component with coboundaries, namely conjugacies by
translations. The fixed-space calculation supplies their full
transport-groupoid comparison; a prescribed nonidentity transporter
component is treated by composing with its representative in $G$.
Thus the corrected four fidelity levels all hold. We finish by checking
that this is the actual peripheral comparison, not just an abstract
replacement of $A$. For the plane the diagram is $B\Z\to *$.
Its equivalence $2$-group has component group $\{\pm1\}$,
$2$-automorphism group $\Z$, and sign action. For the annulus use
$B\Z\amalg B\Z\to B\Z$ with the two maps $+1,-1$.
An equivalence is specified on components by a generic sign $\epsilon$
and an end permutation of sign $\eta$. Its peripheral sign is forced
to be $\epsilon\eta$. The two square conjugators can be removed by
peripheral natural transformations; the surviving automorphism of
such a diagram equivalence is one diagonal $\Z$, on which the generic
sign $\epsilon$ acts. Thus its $2$-group has exactly the two-sign
component group and lattice action of the annulus $G$ displayed above.
For the torus, the empty peripheral diagram over $B\Z^2$ has
component group $\operatorname{GL}_2(\Z)$ and $2$-automorphism group
$\Z^2$ with its standard action. In all three cases strict diagram
functors with zero square-conjugators give a section, so the
Postnikov class is zero. The actual comparison maps a full plane or
annulus rotation to the generator traced by a peripheral or generic
basepoint, and torus translation loops to the two lattice generators.
It is therefore an isomorphism on these specified $\pi_1$ groups,
as well as on components. This proves that the models and comparisons
used throughout the proof are the ones induced by C.1's peripheral
functor. In particular the annulus end permutation has never been
forgotten.
\end{proof}

\begin{claim}{C.11}{Proved}{Infinite cyclic failures at levels one and two}
For every hyperbolic surface containing a disc, there is a strict
$\Z$-action for which $(A^0)^u$ is disconnected. There also are two
strict $\Z$-actions with the same mapping-class homomorphism that are
not conjugate. Thus finite-group hypotheses cannot be removed from
levels $1$ and $2$.
\end{claim}
\passage{\S{}23; folios 140--141; scans 279--281}

\begin{proof}
In a coordinate disc choose two disjoint small round discs centered at
$q_\pm=\pm1$. Define $f$ to rotate each small disc by angle
$\sqrt2(1-r/r_0)$ at radius $r<r_0$, and to be the identity elsewhere.
This is a homeomorphism isotopic to the identity through the same
formula with the angle multiplied by a parameter. The two centres are
isolated fixed points; all other fixed points lie outside the interiors
of these discs. At radius $r_0/2$ the rotation angle is
$\sqrt2/2$, whose ratio to $2\pi$ is irrational, so $f$ has infinite
order. Let $h$ rotate the larger disc by angle $\pi$ on a concentric disc
containing the two small discs, tapering continuously to angle zero
before the outer boundary. It is also isotopic to the identity.
It exchanges the two small discs by rigid rotation and therefore
commutes with $f$.

Set $u(1)=f$. Every homeomorphism commuting with $f$ preserves its fixed
set and sends isolated fixed points to isolated fixed points. Along a
path in the centralizer the induced permutation of $\{q_-,q_+\}$ is
constant. The identity fixes both and $h$ exchanges them. They therefore
belong to different components of $(A^0)^u$. This proves the level-$1$
failure by the exact test C.3. The action $u$ and the trivial $\Z$-action
both project to the trivial homomorphism into $K$, but a conjugate of
$f\ne\id$ cannot be the identity. This proves the level-$2$ failure.
The formulas are checked on an explicit grid in the supplied program;
the isolation and irrationality arguments prove the assertions beyond
that grid.
\end{proof}

\begin{claim}{C.12}{Open}{The requested infinite-group zero-level example}
For the admissible surfaces and the exact comparison of C.1--C.3,
we have not produced an infinite operator group whose level-$0$
comparison has a nontrivial kernel, nor proved universal injectivity.
For a hyperbolic surface the required kernel is
$\pi_1(C_{A^0}(u(\Gamma)),1)$, since the target is zero.
The disconnected centralizer in C.11 proves level-$1$ failure.
The enlarged sphere comparison in C.25 has its stated nontrivial
kernel, but the unpunctured sphere is excluded from C.1's object class.
The requested admissible infinite-group zero-level decision remains
unsettled here.
\end{claim}
\passage{\S{}23; folios 140--141; scans 279--281}

The attempted construction using a supported radial rotation produces
isolated fixed points and a component obstruction, not a nonzero loop.
Adding finitely many marked fixed points still gives aspherical
configuration spaces and cannot by that argument produce the required
$\pi_2$ of an evaluation base. No implication from the examples for
levels $1$ and $2$ to level $0$ is asserted. The independent infinite-group
non-realization theorem C.23 gives level-$3$ failure and does not close
this missing zero-level test. In particular, the possibility that level
$0$ holds for all operator groups has not been ruled out by this return.

\paragraph{Scope of the remaining infinite-group obligation.}
The source fixed in C.1 is the admissible-surface category of scan~264,
folio~133: every genus-zero compact component meets the finite puncture
set. The unpunctured sphere is therefore absent. The last sentence of
Part~C(b) asks for an infinite-group example at each fidelity level
within this specified comparison. An example on one of its admitted
abelian types would meet that existential request; a hyperbolic example
would settle the more specific question in C.12. The examples C.11 and
C.23 supply the independent obstructions at levels~$1$, $2$, and~$3$.
The level-$0$ example inside the admitted source remains unprovided.

Here is the exact kernel that is required. For an arbitrary discrete
operator group and a strict action $u:\Gamma\to A$, put
$H=u(\Gamma)$ and let $N_u$ be its weak automorphism groupoid in
$\Pi_1A$. At its identity object, a $2$-automorphism is precisely a
loop class in $A$ invariant under conjugation by every $u(\gamma)$.
The comparison is thus the inclusion-induced map
\[
 \pi_1(C_A(H),1)\longrightarrow
 \pi_1(N_u,1)=\pi_1(A,1)^\Gamma.                 \tag{C.12.1}
\]
Pointwise commutation proves that the image lies in these invariants.
Every based loop and every based homotopy in the source lies in $A^0$,
so its source can equally be written $\pi_1(C_{A^0}(H),1)$.
For hyperbolic types the target vanishes, giving exactly C.12's test.
For the plane, annulus, and torus the target is the invariant subgroup
of $\Z$, $\Z$, or $\Z^2$, respectively; the required obstruction is a
nonzero kernel of (C.12.1).

C.39 proves that every class in this kernel represented by a loop
whose values commute pairwise is zero. Its contraction is global on
the plane, including arbitrary changes and accumulation of orbit types,
and its explicitly lifted filling disk descends to every connected
admissible surface. Thus a remaining example must involve a class
without a commuting-valued representative. C.40 constructs such a representative for every loop whose generated
closure, in the actual compact-open topology, is locally compact.
No representative reduction for arbitrary loops is proved here. C.41 proves the exact
inclusion injective for every operator group with a dense orbit,
without a local compactness assumption on either group. C.36--C.37 separately
provide the stated contraction on the closed union of dimension-two
proper-image open orbits for general loops. C.42 identifies the actual
boundary quotient as an open map and proves that every component of
every real lifted fibre is a properly embedded line; it also supplies
an exact example with disconnected whole circle fibres. C.43 proves
that the real lifted fibre is connected and supplies actual product
coordinates. C.44 gives the five downstairs covering forms, with
source a plane or annulus. C.45 gives the exact relation between two
quotients and its three planar coordinate domains. C.46 constructs a
simultaneous intrinsic contraction when the full set of lifted
boundary fibre partitions has exactly two members, retaining every
operator and quotient relation. C.47 proves that three independent
actual quotients force the full compatible group on the source plane
to be locally compact and Lie, with its actual topology. C.48 gives
the simultaneous intrinsic contraction for the full one-foliation
family, including every component of the proper operator group and
the dense lattice case. C.49 identifies the actual closed lifted
generated group in the three-quotient case and contracts each whole
orbit closure, with exact overlap agreement and finite pasting.
C.50 supplies an exact uniform word bound and extends these
contractions through arbitrary accumulation, giving a joint
contraction on $X_*=\overline{U_2\cup U_3}$. C.51 extends
the intrinsic contractions through the entire frontier when
$O/Q_O$ is compact and the full boundary-quotient family has
at most two members, with the exact median boundary value
and finite pasting. For an empty family every frontier point
is then fixed. C.52 supplies exact agreement with the
prescribed exponential values on $X_*$ by a natural hybrid
target rule and diagonal frontier continuity. C.53 enlarges
that region to $X^\dagger$ using all global nonopen quotient
maps, and controls every accumulation of all $H$-translates
of finitely many remaining cocompact representatives.
The remaining requirements concern noncocompact quotients,
arbitrarily many $H$-inequivalent remaining components, and
extension through nonopen regions outside these closed sets.

C.25 proves such a kernel for the separately enlarged
$\Pi_1\operatorname{Homeo}(S^2)$ comparison. Its relationship with the
admitted punctured comparisons can be computed explicitly. Write
$p_\pm$ for the poles of its irrational rotation. In the identity
component its centralizer $C$ fixes both poles, and it is the same
subgroup in the unmarked, once-marked, and twice-marked sphere
homeomorphism groups. Forgetting marks induces
\[
\begin{gathered}
 \pi_1C=\Z
 \xrightarrow{\;1\;}\pi_1\operatorname{Homeo}(S^2,\{p_+,p_-\})=\Z
 \xrightarrow{\;1\;}\pi_1\operatorname{Homeo}(S^2,\{p_+\})=\Z,\\
 \pi_1\operatorname{Homeo}(S^2,\{p_+\})=\Z
 \xrightarrow{\bmod 2}\pi_1\operatorname{Homeo}(S^2)=\Z/2.
\end{gathered}                                           \tag{C.12.2}
\]
The same full axis-rotation loop generates the first three groups by
the plane and annulus comparisons of C.10; its inclusion into $SO(3)$
gives the last map. Conjugation by the irrational rotation is based
homotopic to the identity in each ambient group, so these are also the
weak-transporter target groups in (C.12.1). Thus puncturing one or both
poles removes C.25's kernel. Every nonpolar point has an infinite orbit,
so this action admits no other finite invariant puncture set.

Keeping the peripheral-diagram recipe while adding the sphere gives
yet another explicitly related target: its diagram is
$\varnothing\to *$, whose equivalence $2$-group is terminal. The
comparison factors through the functor
$\Pi_1\operatorname{Homeo}(S^2)\to\operatorname{Eq}(\varnothing\to*)$,
which sends the identity $2$-automorphism group $\Z/2$ to zero. For the
irrational action its fundamental-group map is the composite
$\Z\to\Z/2\to0$; C.25 computes the first arrow. Neither enlargement
changes the admitted object class fixed in C.1. Consequently C.25 is
retained with its exact stated comparison, while the admissible
infinite-group zero-level obligation remains recorded here.

\paragraph{Exact lift comparison for the remaining zero-level question.}
The compact-open topology agrees with the marked compactification
topology used in C.1. More precisely, restriction
$\operatorname{Homeo}(X,S)\to\operatorname{Homeo}(X\setminus S)$
is a topological-group isomorphism. A homeomorphism extends uniquely
to the finite end compactification, giving the algebraic inverse.
For its continuity at $f_0$, choose disjoint connected target end
neighborhoods $W_j$ of small diameter and complementary compact
core $C$. Choose a source core $K$ with
$f_0^{-1}(C)\subset\operatorname{int}K$, and connected source tails
$V_i$ whose $f_0$-images lie in the prescribed $W_j$.
The conditions $f^{-1}(C)\subset\operatorname{int}K$ and
$f(z_i)\in W_j$ for one point $z_i\in V_i$ are open:
inversion is continuous in the compact-open homeomorphism group
of a locally compact, locally connected surface.
Connectedness forces $f(V_i)\subset W_j$, since this image avoids
$C$ and meets $W_j$. This controls the extended maps uniformly on
all tails. Compact-open control on the remaining compact core
controls them on the rest of $X$, proving continuity of the inverse.

Let $U$ be the hyperbolic surface of C.12, let $A^0=\operatorname{Homeo}(U)^0$, let $H$ be the image of the specified operator group, and let $p:\widetilde U\to U$ be the universal covering with deck group $\pi$. Write $\widetilde H$ for the group of all lifts of elements of $H$, regarded as actual homeomorphisms of $\widetilde U$. There is a continuous injective homomorphism
\[
 j:A^0\longrightarrow\operatorname{Homeo}(\widetilde U),
 \qquad p\,j(f)=f\,p,
\]
characterized by $j(f)\delta=\delta j(f)$ for every $\delta\in\pi$. It identifies the precise centralizers:
\[
 j\bigl(C_{A^0}(H)\bigr)
   =j(A^0)\cap C_{\operatorname{Homeo}(\widetilde U)}(\widetilde H).
\]
In particular every based loop in the left-hand centralizer lifts to a closed loop of homeomorphisms commuting with the entire lift group.
\begin{proof}
Lift an isotopy from $\id$ to $f$ starting at $\id_{\widetilde U}$. Conjugation on the discrete deck group is constant during the lift, so its endpoint commutes with $\pi$. Two such endpoints differ by a deck transformation centralizing $\pi$; they are equal because $Z(\pi)=1$ by Part A. Uniqueness also makes $j$ a homomorphism and an injection. For continuity, fix $\widetilde u\in\widetilde U$. The evaluation map $A^0\to U$ lifts continuously to $\widetilde U$, since $A^0$ is contractible and locally path connected by the surface homeomorphism-group theorem. This lift gives $j(f)(\widetilde u)$, and the local lifting property of a covering gives continuous dependence of the lifted homeomorphism in the compact-open topology.

Suppose $f$ centralizes $H$ and $\widetilde h$ is any lift of $h\in H$. Then $[j(f),\widetilde h]$ is a deck transformation. Its conjugation on $\pi$ is the commutator of the identity automorphism with the automorphism induced by $\widetilde h$, hence is trivial. Centre-freeness makes this deck transformation the identity. This proves one inclusion in the displayed equality; descent proves the other. Finally $j(\id)=\id$, so a based loop has the same endpoint after lifting.
\end{proof}

This comparison does not supply a contraction inside the right-hand centralizer. Choosing a hyperbolic metric and interpolating from $j(f)(x)$ to $x$ along geodesics has two unproved properties: the arbitrary lift group $\widetilde H$ need not preserve these geodesics, and an interpolated map need not be a homeomorphism. The first prevents asserting equivariance; the second prevents asserting that the interpolation lies in $j(A^0)$. Neither property follows from the displayed centralizer equality. The universal vanishing of $\pi_1(C_{A^0}(H))$, and the requested counterexample to it, therefore remain unresolved by this argument.

\paragraph{Current orbitwise reduction and the remaining compatibility problem.}
C.26 settles the globally transitive centralizer case on a closed
hyperbolic surface and also gives an explicit closed subgroup with an
ambient-null but intrinsically nontrivial loop; its latter example is
not asserted to be a centralizer. For the canonical plane lift above,
C.28 identifies the orbits of the group generated by the loop:
points, intrinsic lines, circles, and open surface orbits.
C.27 gives an order-natural contraction on each intrinsic line.
C.29 proves that every circle restriction has winding zero by the
actual plane-loop evaluation map, and contracts that restriction
in its circle centralizer.
C.31 supplies inverse and track control, and C.34 proves joint
continuity of the contractions and their inverses on the entire
closed union $Y$ of the nonopen orbits. This is a based contraction
through homeomorphisms of $Y$, commuting with every homeomorphism
that commutes with the original loop. In particular C.34 settles
the hyperbolic centralizer loop whenever its lifted generated
group has no open orbit.

The open orbits have also been treated intrinsically. C.30 gives
the exact restricted operator group, including its restriction
kernel, and proves that its closed image acts freely and properly.
C.32 handles a discrete image, and C.33 treats every locally
compact closed image supplied by C.30. In each case the restricted
loop has a based null-homotopy inside the centralizer of the full
restricted operator group. These results give contractions on
all individual open orbits and on the whole nonopen locus.
The remaining step is to choose these contractions compatibly
across the frontiers of the open components, with joint
continuity of the resulting homeomorphisms and their inverses.

C.35 proves that this step cannot require extending the existing
median prescription unchanged. Its explicit plane action has
the boundary restriction map
\[
 \R^2\longrightarrow\R^3,\qquad
                  (A,B)\longmapsto(A,B,A-B).
\]
At the displayed parameter the prescribed three edge medians
are $(1,0,0)$, which lies outside its image because
$r_1-r_2-r_3=1$. Thus that prescribed homeomorphism of $Y$
has no equivariant extension over an adjacent open orbit.
The original loop in the same example has the explicit global
equivariant contraction given in (C.35.5).
C.36 constructs compatible choices on closures of open orbits
whose closed proper operator image has dimension two, using the
actual open quotient onto their simply transitive Lie
centralizers. Scaling the logarithm of each based universal-cover
lift respects the exact common boundary action. C.37 supplies a
uniform three-word estimate and proves joint continuity, including
inverses, on
\[
 X_2=\overline{\bigcup_{\dim Q_O=2}O},
\]
including every infinite accumulation of those open orbits.
Its bounded residual equations determine the unique logarithmic
limit on line orbits. The actual null residual word isotopies
and the two-point configuration map determine the real lifted
limit on circle orbits, including its winding.

C.49--C.50 subsequently enlarge this closed set to $X_*$,
including all open orbits with at least three actual boundary
quotient foliations. C.51 proves whole-closure extension for
the remaining open orbits with compact operator quotient,
using the median frontier prescription and finite pasting.
C.52 makes these extensions exactly relative to $X_*$.
C.53 enlarges the exponential region to $X^\dagger$ and
proves continuity across all accumulations of the operator
translates of finitely many remaining cocompact representatives.
The outstanding problem includes noncocompact operator
quotients, infinitely many $H$-inequivalent remaining
representatives, and extension through exterior regions
consisting of nonopen orbits. The hybrid target prescription
has been proved continuous on the treated closed sets;
its transverse continuity on all of the nonopen locus has
not been proved. C.35 requires allowing changes to the
original median prescription. These exact remaining
requirements, and the general kernel in (C.12.1), remain
unresolved here.


\subsection{Configuration invariants and drilling}

\begin{claim}{C.13}{Proved}{The invariant braid group and nonabelian cohomology}
Let a finite group act faithfully and orientation preservingly on a
hyperbolic $U$, and let $R=U^!$ be its finite set of points with nontrivial
stabilizer. Put $V=U\setminus R$,
$B=\{f\in A^0:f|_R=\id\}$, and $B^0=\operatorname{Homeo}(V)^0$.
With the action by conjugation, we have
\[
 (B/B^0)^\Gamma=1,\qquad H^1(\Gamma,B^0)=1,
 \qquad (A^0)^\Gamma=B^\Gamma\cong
 \operatorname{Homeo}(V/\Gamma)^0.
\]
Here the quotient group identifies canonically with
$\pi_1(\operatorname{Conf}_R(U),R)$ by evaluation, and its invariant
subgroup is the $H^0$ asked for on scan 285.
\end{claim}
\passage{\S{}24; folios 142--144; scans 282--286}

\begin{proof}
Evaluation $A^0\to\operatorname{Conf}_R(U)$ is a locally trivial
principal $B$-bundle, by the local extension theorem for embeddings of
a finite set. Its total space is contractible. Its homotopy exact
sequence gives $\pi_0B\cong\pi_1\operatorname{Conf}_R(U)$, with the
specified base configuration, and its identity component is the
identity component after retaining the new punctures. The construction
is equivariant for $\Gamma$, so this is the stated action on the braid
group.

An invariant component $b\in B/B^0$ is a mapping class of $V$ in the
kernel of filling $R$ that centralizes the action of $\Gamma$ on $V$.
By C.6 and its transporter version, choose a representative $f$ on $V$
commuting with $\Gamma$. It extends across the punctures and induces the
identity mapping class on $U$. Thus $f\in(A^0)^\Gamma$, and C.7 joins
it to the identity by an equivariant isotopy on $U$. Such an isotopy fixes
each point of $R$: it preserves each finite critical-point set and
starts with the identity permutation. Restriction to $V$ proves $b=1$.

For a nonabelian $1$-cocycle $a_\gamma\in B^0$, the maps
$v(\gamma)=a_\gamma u(\gamma)$ form another finite action on $V$.
The cocycle equation is exactly the homomorphism equation for $v$.
Both actions induce the same homomorphism into the mapping class group
of $V$. By C.6 they are conjugate by an element $b\in B^0$; equivalently
$a_\gamma=b u(\gamma)b^{-1}u(\gamma)^{-1}$. Hence the cocycle is a
coboundary and $H^1(\Gamma,B^0)=1$.

Finally C.7 gives connectedness. Descent to the quotient with its branch
points retained is precisely descent to $V/\Gamma$ with its punctures
retained; the identity lift construction in C.7 is the displayed
isomorphism. In particular the unprimed $V$ in the transcription of the
identity-component line on scan 286 must be compared with the quotient
$V'$ of scans 282--285. The two homeomorphism groups need not be
isomorphic; the descent map identifies the quotient one.
\end{proof}

\begin{claim}{C.14}{Refuted}{Finite critical configurations do not include mirrors}
The finite-set configuration formulation of C.13 does not extend by
replacing an orientation-preserving action by an arbitrary action and
keeping $R=U^!$ equal to all points with nontrivial stabilizer. A reflection
has a one-dimensional fixed locus, and deleting finitely many points
cannot make that action free.
\end{claim}
\passage{\S{}24; folios 142--148; scans 282--294}

\begin{proof}
Take complex conjugation on
$U=\mathbf P^1(\C)\setminus\{0,1,\infty\}$. It acts faithfully on this
hyperbolic surface. Its fixed set is
$\mathbf P^1(\R)\setminus\{0,1,\infty\}$, a union of three open arcs.
Every point of these arcs has nontrivial stabilizer. A finite deletion
leaves points in each arc. At any remaining fixed point, the lift of
conjugation fixing a selected universal-cover point is an element of
order two in the lift extension. Therefore the drilled extension still
has torsion. This explicit counterexample refutes both the finiteness
premise and a torsion-free finite-drilling assertion for such an action.
The corrected geometry retains mirror boundary in the quotient or cuts
along the mirror locus; it is not the operation of adding finitely many
puncture loop classes. The equations C.13 were stated on scans 282--286
after the explicit restriction $\Gamma=\Gamma^+$.
\end{proof}

\begin{claim}{C.15}{Proved}{The oriented drilling equivalence and the bijection}
Fix a finite group $\Gamma$ and faithful orientation-preserving outer
actions on discrete hyperbolic groups with loops. Let $\Phi(\pi,\Gamma)$
be the set of $\pi$-conjugacy classes of nontrivial maximal finite
subgroups of the canonical extension. Filling a specified
$\Gamma$-stable set $J$ of punctures with nontrivial stabilizers induces
an equivalence between the category of such drilled actions and the
category of filled actions with a specified subset
$J\subset\Phi(\pi,\Gamma)$. In particular
\[
 \Phi(\pi',\Gamma)\amalg J\xrightarrow{\ \sim\ }\Phi(\pi,\Gamma)
\]
is a canonical $\Gamma$-equivariant bijection. Drilling all of $\Phi$
produces a torsion-free extension, uniquely up to the unique exterior
isomorphism over the specified filled object.
\end{claim}
\passage{\S{}24; folios 144--147; scans 286--293}

\begin{proof}
Realize the action by C.5. A finite nontrivial subgroup of the lift group
acts on $\mathbb H^2$ by orientation-preserving isometries and has a
unique fixed point. Its stabilizer at that point is the unique maximal
finite subgroup containing it. Conversely every point of $U$ with
nontrivial stabilizer gives that subgroup after choosing a lift.
Changing the lift conjugates by $\pi$. This identifies $\Phi$ with the
finite critical-point set of the realized action, including its
$\Gamma$ action. It also proves the assertion about maximality, rather
than assuming it for arbitrary finite subgroups containing reflections.

Deleting the points of $J$ gives the drilled surface. Remaining critical
points and deleted points partition the original critical-point set,
which proves the displayed bijection. The group map is the epimorphism
$\pi'\to\pi$ obtained by killing the normal closure of exactly the new
peripheral subgroups. Local punctured discs show that its induced map
on lift groups sends a cusp stabilizer to the corresponding finite
point stabilizer. Thus the geometric bijection is the algebraic map
of scan 290, with the same conjugacy ambiguities.

To prove fullness, realize an equivariant outer isomorphism of the
filled peripheral groups by an equivariant homeomorphism, using C.6's
transporter conclusion. It preserves the selected subset of $\Phi$ by
the preceding construction. Its restriction to the complements gives
the desired drilled outer isomorphism. To prove faithfulness, represent
a drilled equivariant outer automorphism in the kernel by an equivariant
homeomorphism of the drilled surface and extend it across the new
punctures. It is isotopic to the identity on the filled surface. By C.7
there is an equivariant identity isotopy there. Every point of $J$ stays
fixed throughout, since it is a member of the finite critical set, so
the isotopy restricts to the drilled surface. The kernel is therefore
trivial. These two arguments also prove uniqueness over a fixed filled
isomorphism, not merely uniqueness of isomorphism classes.

After all critical points are removed, the finite action is free. Its
lift extension acts freely on $\mathbb H^2$ and is torsion-free, since
a nontrivial finite-order isometry fixes a point. The operations of
deleting and filling are inverse by the just-proved full faithfulness.
The only use of geometric realization in this proof is C.5 together
with the equivariant homeomorphism and isotopy results C.6--C.7; the
comparison map on groups and the kernel argument have been specified.
\end{proof}

\begin{claim}{C.16}{Refuted}{Nonfaithful outer actions obstruct torsion-free drilling}
The unrestricted existence claim on scans 286--288 fails for the
trivial outer action of $C_2$ on the group of type $(0,3)$. The faithful,
orientation-preserving version is C.15.
\end{claim}
\passage{\S{}24; folios 144--145; scans 286--288}

\begin{proof}
Suppose a finite-puncture drilling of that action had torsion-free
extension. Realize its finite outer action on the drilled hyperbolic
surface by C.5. Filling the specified punctures extends these
homeomorphisms to a finite action on the filled three-punctured sphere.
The induced mapping class of every element is trivial. The finite-order
rigidity argument of C.5 makes every filled homeomorphism the identity.
Its restriction to the dense drilled surface is therefore the identity.
The upstairs outer action is trivial as well. For a centre-free kernel
its canonical extension is $\pi'\times C_2$, which contains
$\{1\}\times C_2$. This contradicts torsion-freeness. The argument proves
the same obstruction for every nonfaithful outer action after restricting
to a nontrivial element of its finite kernel.
\end{proof}

\begin{claim}{C.17}{Proved}{Torsion-free finite extensions with peripheral data}
Let $\pi$ be a discrete hyperbolic surface group, possibly punctured,
and let $1\to\pi\to E\to\Gamma\to1$ be a torsion-free extension by a
finite group whose outer action preserves the loop structure. Then $E$
is the fundamental group of a finite-type surface, possibly
nonorientable. Its peripheral subgroups are the centralizers in $E$ of
the peripheral subgroups of $\pi$. With their conjugacy classes it is a
surface group with loops in this extended sense. The associated compact
surface and boundary groups form a Poincar\'e duality pair of dimension
two with its orientation local system.
\end{claim}
\passage{\S{}24; folios 147--148; scans 292--294; \S{}28; folio 178; scan 353}

\begin{proof}
The centre-free extension is the pullback of
$\Aut(\pi)\to\Out(\pi)$. The centralizer of $\pi$ in $E$ maps
isomorphically to the kernel of the outer action: a representative
inducing an inner automorphism can be multiplied by the inverse inner
element to centralize $\pi$, and this representative is unique because
$Z(\pi)=1$. This centralizer is finite, so torsion-freeness forces it
to be trivial. Hence the outer action is faithful.

Apply C.5 and identify $E$ with the group of lifts of the realized
action, by uniqueness of the centre-free extension. A nontrivial
stabilizer of a point gives a finite subgroup of $E$, so the action on
$U$ is free. Its quotient $W=U/\Gamma$ is an ordinary surface and
$\mathbb H^2\to W$ is its universal covering. This gives the required
isomorphism $E\cong\pi_1(W)$, respecting the specified subgroup
$\pi$ and quotient $\Gamma$.

A stabilizer of an end of $U$ cannot contain an orientation-reversing
operator. The local finite-group model at its compactifying point would
then contain a reflection and fixed points arbitrarily close to that
point, contrary to freeness on $U$. Its local action is therefore cyclic
rotation. Consequently an upstairs peripheral subgroup is a finite-index
subgroup of the downstairs peripheral infinite cyclic group. In the
hyperbolic universal cover its centralizer is exactly that full cusp
group. These are all the cusps, since every end downstairs has a
preimage upstairs. Thus the stated peripheral classes, including their
full subgroups rather than merely inclusions, are obtained.

Choose disjoint cusp neighbourhoods in $W$ and truncate them to obtain
a compact surface $M$ with boundary. The inclusions of boundary circles
induce precisely these cyclic groups. Lift a triangulation to the
contractible universal cover. The cellular cap-product map with the
relative orientation class gives the chain equivalence between
cochains and complementary relative chains; the local sign is the
orientation character of $W$. This is the Poincar\'e--Lefschetz duality
pair, and its group-pair chain model is the relative cellular complex
because both the universal cover and the lifted boundary components
are contractible. This proves the duality-pair assertion rather than
inferring it from virtual freeness alone. The exact geometric input is
the finite realization theorem; a separate proof from abstract duality
axioms is not substituted for this argument.
\end{proof}

\begin{claim}{C.18}{Proved}{The discrete oriented realization criterion}
For a faithful orientation-preserving finite outer action on a discrete
hyperbolic group with loops, the criterion of scan 294 is valid with
the word admissible meaning the full-faithfulness and uniqueness
conditions of scan 293. The drilling in question is at the canonical
set $\Phi(\pi,\Gamma)$, and its extension carries the quotient loop
structure described in C.17.
\end{claim}
\passage{\S{}24; folios 147--148; scans 293--294}

\begin{proof}
A realized action can be drilled at its critical points. C.15 proves
that the resulting action is admissible and its extension torsion-free;
C.17 identifies that extension with a group with loops. This verifies
every requirement of the stated criterion.

Conversely let $\pi'\to\pi$ be the stipulated filling map and let
$E'$ be its extension, with the group-with-loops structure required in
the criterion. Realize $E'$ by a compact surface with its boundary
circles, then take its covering corresponding to $\pi'$. The comparison
of peripheral subgroups in the definition ensures that this cover has
exactly the prescribed upstairs loop classes. Its deck group is
$\Gamma$, acting freely before filling. The cyclic end stabilizers
extend across each new puncture by the local map $z\mapsto z^m$.
Filling the selected ends thus gives a $\Gamma$-action on the surface
whose peripheral group is $\pi$. The canonical map of lift extensions
identifies its outer action with the original one. Admissibility makes
this construction independent, up to the uniquely specified exterior
isomorphism, of the chosen drilled object. This proves the converse
by the typed covering and filling maps. The mirror counterexample
C.14 explains the necessary restriction of this formulation.
\end{proof}

\subsection{Special homeomorphisms: the topology and the extension class}

\begin{claim}{C.19}{Refuted}{Escaping support invalidates the compact-open version}
Let $SA\subset A$ consist of homeomorphisms equal to the identity on
some neighbourhood of every puncture, and give it the subspace
compact-open topology. The canonical special-group comparison of
scan 304 does not identify $\pi_0(SA)$ with the boundary-framed mapping
class group. A peripheral boundary twist is connected to the identity
inside this $SA$. In particular, the annulus conclusion
$\pi_0(SA)=\Z$ on scan 304 is false in this topology.
\end{claim}
\passage{\S{}25; folios 151--153; scans 301--304}

\begin{proof}
Use a coordinate $z=re^{i\theta}$ on a disc about a puncture at $0$.
Choose a continuous function $\eta:[0,\infty)\to[0,1]$ equal to $0$ on
$[0,1/3]$ and to $1$ on $[2/3,\infty)$. For $0<s\leq1$ define
\[
 h_s(re^{i\theta})=r\exp\bigl(i\theta+2\pi i\eta(r/s)\bigr),
 \qquad h_0=\id,
\]
and extend by the identity outside the coordinate disc. Each $h_s$ is
identity near the puncture and outside the annulus
$s/3<r<2s/3$. It is a boundary twist when viewed on a fixed compact
surface obtained by truncating the puncture. Nevertheless
$\sup_z|h_s(z)-z|\leq4s/3$, and the same bound holds for its inverse.
Thus $s\mapsto h_s$ is a continuous path in the stated subspace topology
from the identity to that twist.

For the annulus every compactly supported orientation-preserving
homeomorphism fixes both ends and is represented on a compact annulus
by an integer twist; cutting along a spanning arc and applying the
Alexander isotopy of a disc proves this integer classification. The
isotopy to the corresponding integer twist can be chosen in that fixed
compact annulus. The path above then kills that twist in $SA$. Hence
$SA$ is path connected in the compact-open topology, proving the
explicit contradiction. The programme checks the formula and its
uniform bound; the bound itself proves convergence for the full space,
not just on a finite sample.

The fixed-boundary formulation below retains a single compact surface
throughout the isotopy. That is the required correction. No change of
topology is silently made in this refutation.
\end{proof}

\begin{claim}{C.20}{Proved}{The fixed-boundary repair and its class}
Let $M$ be an oriented compact surface obtained by replacing the
$\nu>0$ punctures of a hyperbolic type by boundary circles. Put
$A_\partial=\operatorname{Homeo}^+(M\ {\rm rel}\ \partial M)$.
Then $A_\partial^0$ is contractible,
$\pi_0(A_\partial)\cong S\Gamma^{!+}_{g,\nu}$ for the framed group
of E.1--E.2, and there is a central extension
\[
 1\longrightarrow\Z^\nu\longrightarrow\pi_0(A_\partial)
 \longrightarrow\Gamma^{!+}_{g,\nu}\longrightarrow1.
\]
The $i$th kernel generator is the positive boundary twist. Its
extension class is the vector of Euler classes of the oriented tangent
circle lines at the marked points, with the sign fixed by that generator.
For the annulus the kernel is the quotient of $\Z^2$ by the diagonal
rotation relation; for the once-punctured sphere it is zero.
\end{claim}
\passage{\S{}25; folios 152--154; scans 302--306}

\begin{proof}
Allow boundary rotations first, retaining each boundary component
setwise. Restriction to the boundary is a locally trivial fibration
onto $\prod_i\operatorname{Homeo}^+(S^1)$: extend a small boundary
isotopy across collars. Its fibre is $A_\partial$.
Each boundary homeomorphism group deformation retracts onto rotations
by lifting to increasing maps of $\R$ and linearly interpolating the
periodic displacement. Thus the base has homotopy type $(S^1)^\nu$.
Capping each boundary by a disc with its centre marked identifies the
free-boundary mapping classes with the pure oriented punctured mapping
classes. The same collar-extension construction gives the homotopy
comparison on identity components.

The identity component of the free-boundary group is contractible in
a hyperbolic type by Yagasaki's compact-pair theorem. The homotopy
sequence therefore injects $\pi_1((S^1)^\nu)=\Z^\nu$ into the
components of the fibre and identifies the quotient as stated. Lifting
one rotation in the $i$th collar gives exactly the boundary twist.
Its conjugation action is trivial because the boundary circles are
retained individually and their orientations are fixed. Choose a point
on each boundary circle. A boundary-fixing homeomorphism acts on the
framed peripheral diagram of E.2, and an isotopy relative to the
boundary gives its prescribed natural isomorphism. This defines
$\pi_0(A_\partial)\to S\Gamma^{!+}_{g,\nu}$. On the quotient it is
the peripheral Dehn--Nielsen--Baer isomorphism. A positive boundary
twist changes the path from the generic basepoint to that boundary
point by one peripheral turn; hence the induced map on kernels sends
the $i$th boundary generator to the $i$th generator of $T^I$.
It is an isomorphism on both kernel and quotient, so the two short
exact sequences give the asserted isomorphism of extensions.
Contractibility of $A_\partial^0$ follows directly from Lemma 2.1(ii),
with the nonempty fixed boundary. The connecting homomorphism is consequently the
fundamental-group extension of the bundle of tangent-circle framings
at the marked points. Obstruction to a section of that circle bundle
is its Euler class, proving the class identification with its sign
convention. In complex notation this is the tangent-line first Chern
class; the usual cotangent $\psi_i$ has its opposite sign.

For the annulus the identity component of the free-boundary group is
$S^1$. Its fundamental-group map to the two boundary-rotation groups
is the diagonal when both circles use the same angular coordinate,
or the antidiagonal when both use boundary orientations. The cokernel
is $\Z$. For the disc this map is an isomorphism $\Z\to\Z$, so the
cokernel is zero. These are the exact boundary models underlying the
two abelian computations on scans 303--304.
\end{proof}

\begin{claim}{C.21}{Proved}{The four small extension classes}
With the positive boundary-twist convention of C.20, the special
fixed-boundary extension classes for $(0,3)$ and $(0,4)$ are zero. For
$(1,1)$ the class is a generator $e$ of
$H^2(\operatorname{SL}_2(\Z);\Z)=\Z/12$. For $(1,2)$ the vector class
is $(q^*e,q^*e)$, where
$q:\Gamma^{!+}_{1,2}\to\operatorname{SL}_2(\Z)$ forgets the second
marked point; its order is $12$.
\end{claim}
\passage{\S{}25; folios 153--154; scans 304--306}

\begin{proof}
The pure oriented mapping class group of a three-punctured sphere is
trivial. Filling one puncture of a four-punctured sphere gives the
point-pushing group $F_2$, since the target pure group is trivial and
the three-punctured surface is hyperbolic. A graph is a classifying space
for $F_2$, so $H^2(F_2;\Z^4)=0$. This proves the first two class
computations, including the noncanonical splitting in type $(0,4)$.

The integral matrices
\[
 S=\begin{pmatrix}0&-1\\1&0\end{pmatrix},\qquad
 R=\begin{pmatrix}0&-1\\1&1\end{pmatrix}
\]
have orders $4$ and $6$ and satisfy $S^2=R^3=-I$.
Their action on the Farey tessellation gives
$\operatorname{PSL}_2(\Z)=C_2*C_3$: the dual bipartite tree has one
edge as fundamental domain, vertex stabilizers $C_2,C_3$, and trivial
edge stabilizer. Connectivity and the absence of circuits follow from
the Euclidean algorithm for adjacent reduced fractions. Lifting the
central element $-I$ gives
$\operatorname{SL}_2(\Z)=C_4*_{C_2}C_6$. The cohomology sequence of
this tree of groups identifies $H^2$ with
\[
 \{(r,s)\in\Z/4\oplus\Z/6:r\equiv s\pmod2\}:
\]
the preceding $H^1(C_2;\Z)$ and $H^1(C_4;\Z)\oplus
H^1(C_6;\Z)$ vanish, while the restrictions in degree two reduce
modulo $2$. This group has twelve elements and is cyclic.
The tangent-circle bundle restricts on the square torus fixed by $S$
to the rotation representation by $\pi/2$, and on the hexagonal torus
fixed by $R$ to the representation by $\pi/3$. Lifting these circle
rotations to $\R$ gives fourth and sixth powers equal to one positive
full turn. By the convention of C.20 the Euler class therefore
restricts to $(1,1)$, an element of order $12$. This proves its generator
assertion without an unproved presentation for the boundary mapping
class group. The supplied code verifies the matrix orders and the
finite abelian-group calculation.

In a complex torus translation canonically identifies the tangent
spaces at any two points. These identifications vary with the marked
complex torus and commute with all isomorphisms. Hence the two tangent
circle bundles on the two-pointed moduli orbifold are isomorphic, and
both are pulled back from the one-pointed tangent bundle. This proves
the vector formula. It has order dividing $12$. A square torus with
marked points $0$ and $(1+i)/2$ admits the order-four rotation preserving
both; a hexagonal torus admits an order-three rotation with at least
two marked fixed points. The restrictions of the class to these cyclic
subgroups are generators, as can be seen from their nontrivial local
rotations, or from C.22 below. Its order is divisible by $4$ and $3$,
and thus is exactly $12$. The existence of the displayed fixed points
is also verified in the integral torus matrices supplied with the code.
\end{proof}

\begin{claim}{C.22}{Proved}{Restriction over a finite puncture stabilizer}
For a hyperbolic type, a nontrivial finite subgroup
$G\subset\Gamma_i^+$ is cyclic. The extension by the $i$th copy of
$T=\Z$ obtained from the special group is nonsplit over $G$; its
preimage is infinite cyclic and the kernel is the subgroup of index
$|G|$.
\end{claim}
\passage{\S{}25; folios 153--154; scans 305--306}

\begin{proof}
Realize $G$ by C.5. It fixes the marked point $i$ and preserves
orientation. The finite local action is a faithful rotation action:
an element trivial on a small disc is the identity as a periodic
surface homeomorphism, by the local classification and connectedness.
Thus $G$ embeds in the circle rotations and is cyclic. Delete that
point and lift a small punctured quotient disc to its universal cover.
Its group is $\Z$, while the upstairs punctured disc has subgroup
$|G|\Z$. The exact normalizer interpretation proved in E.3 identifies
this with the extension pushed out by the $i$th projection. An infinite
cyclic group contains no nontrivial finite subgroup, so the extension
cannot split. This proves the remark on scan 306 with its stated
orientation and puncture-stabilizer hypotheses.
\end{proof}


\begin{claim}{C.23}{Proved}{An infinite-group failure of existence}
For every closed orientable surface $X_g$ of genus $g\geq2$, the
identity homomorphism from the discrete group
$\Gamma=\operatorname{Mod}^+(X_g)$ to its subgroup in $\pi_0A_g$
has no homomorphic lift to $A_g$. Thus level $3$ fails for an infinite
operator group already on each such closed hyperbolic surface.
\end{claim}
\passage{\S{}23; folios 140--141; scans 279--280}

\begin{proof}
Markovi\'c--\v Sari\'c's Theorem 1.1
\cite{C_MarkovicSaric2008} states that the projection
$\operatorname{Homeo}^+(X_g)\to\operatorname{Mod}^+(X_g)$ has no group
section for $g\geq2$. The theorem is about homeomorphisms, not merely
diffeomorphisms; its first-page statement and definitions have been
checked in the original paper. A lift of the displayed homomorphism
would take every element to an orientation-preserving homeomorphism,
since orientation is invariant under isotopy. Its composite with that
projection would be the identity, giving the forbidden section. This
is the exact reduction to the source theorem. The group is infinite:
a twist about a nonseparating simple curve acts on $H_1(X_g;\Z)$ by
an integral transvection, whose distinct powers are distinct.
\end{proof}

\subsection{The integral sphere comparison}

\begin{claim}{C.24}{Refuted}{The sphere comparison already fails for an order-three rotation}
Let $A=\operatorname{Homeo}(S^2)$ with its compact-open topology, and let
$\Gamma=C_3$ act on $S^2$ by rotation through $2\pi/3$ about a fixed axis.
Give $A$ the resulting conjugation action and base both fixed-point spaces
at the identity. The actual comparison has
\[
 \pi_1(A^\Gamma,1)=\Z,\qquad
 \pi_1(A^{h\Gamma},1)=\Z_3\times\Z/2,
 \qquad n\longmapsto(n,n\bmod2).
\]
Consequently the requested integral homotopy equivalence fails in type
$(0,0)$. This failure concerns the actual homeomorphism mapping space;
it is independent of the earlier loss of the sphere's $\pi_2$ in its
fundamental groupoid. For every odd prime $p$, the same construction with
$C_p$ gives $\Z_p\times\Z/2$ and the same comparison formula with $p$
in place of $3$.
\end{claim}
\passage{\S{}23; folio 141; scans 280--281}

\begin{proof}
We give the comparison maps and compute their effect, including the integral
rather than only completed homotopy group. Write $G=SO(3)$ and
$\rho:C_p\hookrightarrow G$ for the rotation subgroup, with $p$ odd.
The inclusion $G\hookrightarrow A^+=\operatorname{Homeo}^+(S^2)$ is
$C_p$-equivariant and a weak equivalence. Indeed evaluation on three
ordered distinct points is a locally trivial fibration onto
$\operatorname{Conf}_3(S^2)$, by isotopy extension. Its fibre is the
pure homeomorphism group of the three-punctured sphere. Its component
group is trivial, and its identity component is contractible by the
surface homeomorphism-group theorem used in C.7. The unique M\"obius
transformation carrying one ordered triple to another identifies this
configuration space with $PSL_2(\C)$; polar decomposition retracts this
group onto $SO(3)$. The resulting evaluation diagram proves the stated
weak equivalence for the actual inclusion. Homotopy fixed points preserve
this objectwise weak equivalence, after applying singular complexes. A
path through the identity in $A^{hC_p}$ stays in $A^+$, because $EC_p$
is connected and orientation takes values in the discrete group $C_2$.
It therefore suffices to calculate $(G^{hC_p},1)$.

The identity component of the actual centralizer fixes the two poles
individually. Quotient by $z\mapsto z^p$ identifies its descent map to
$\operatorname{Homeo}^+(S^2;0,\infty)^0$ as a covering of degree $p$.
Every map of the quotient fixing both poles lifts: on the complementary
annulus its action on the meridian is $+1$, so it preserves the subgroup
$p\Z$ of the covering group. Each lift commutes with the deck group,
since the induced automorphism of $\Z/p$ is $+1$; it extends uniquely
over the two poles. The usual local lifting neighbourhoods establish
that descent is a covering in the compact-open topology. Its total space
is connected: the quotient group is connected, and all deck lifts lie
on the connected circle of axis rotations. The annulus calculation of
C.10 identifies the base with a space of homotopy type $S^1$. The
rotation circle maps to its quotient rotation circle by $z\mapsto z^p$;
its inclusion in the total centralizer therefore induces an isomorphism
on every homotopy group. In particular the actual fundamental group is
$\Z$, with generator one complete axis rotation.

Set $Y=BG=BSO(3)$. The adjoint bundle
$EC_p\times_{C_p}G\to BC_p$, with action through $\rho$, has section
space $G^{hC_p}$. It is also the bundle of gauge automorphisms of the
principal $G$-bundle classified by $B\rho$. The universal bundle gives
the natural weak equivalence
\[
 G^{hC_p}\simeq
 \Omega_{B\rho}\operatorname{Map}(BC_p,Y).
\]
One way to verify the map is to send a gauge automorphism to the loop
of its classifying map obtained by gluing the bundle over the two ends
of $BC_p\times[0,1]$. Conversely a loop of classifying maps gives a
bundle over this cylinder with its specified endpoint identifications;
trivialization in the interval direction gives the gauge automorphism.
The spaces of such trivializations are contractible. This construction
is natural for the inclusion $C_G(\rho(C_p))\hookrightarrow G^{hC_p}$.

We now compute $\pi_2$ of this mapping space using the arithmetic square
of the simply connected finite-type space $Y$. Let
\[
 \widehat Y=\prod_qY_q^{\wedge},\qquad
 V=\Bigl(\prod_q\Z_q\Bigr)\otimes\Q.
\]
Here $q$ runs over primes and superscript $\wedge$ means homotopical
$q$-completion, throughout this proof. We have
$Y_{\Q}\simeq K(\Q,4)$ and
$(\widehat Y)_{\Q}\simeq K(V,4)$. This follows from
$SO(3)\simeq\mathbb {RP}^3$: its rational homotopy is concentrated in
degree three, and each of its higher homotopy groups is finite. The
square
\[
 \begin{matrix}
 Y&\longrightarrow&\widehat Y\\
 \downarrow&&\downarrow\\
 K(\Q,4)&\longrightarrow&K(V,4)
 \end{matrix}
\]
is a homotopy pullback. For completeness its homotopy-group verification
reduces on each free summand to
\[
 0\longrightarrow\Z\longrightarrow
 \Q\oplus\prod_q\Z_q\longrightarrow V\longrightarrow0,
 \quad (a,(b_q))\longmapsto(b_q-a)_q.
\]
Surjectivity follows by clearing a common denominator and choosing its
finitely many required residues by the Chinese remainder theorem; the
kernel consists of rational numbers integral at every prime and is
$\Z$. On finite summands completion gives the product of their primary
parts and rationalization is zero. The long exact sequence of the
homotopy pullback and Whitehead's theorem prove the square, also giving
the usual arithmetic-square construction in the exact form needed here.

Apply $\operatorname{Map}(BC_p,-)$ to this homotopy pullback and take the
component containing $B\rho$. For every rational vector space $W$,
positive-degree cohomology $H^k(C_p;W)$ vanishes by averaging. The two
rational mapping spaces consequently have vanishing $\pi_2$ and
$\pi_3$. The homotopy-pullback exact sequence therefore gives
\[
 \pi_2\operatorname{Map}(BC_p,Y)_{B\rho}
 \cong\prod_q
 \pi_2\operatorname{Map}(BC_p,Y_q^{\wedge})_{B\rho}.
\]
For $q\ne p$, evaluation is a weak equivalence from this mapping-space
component to $Y_q^{\wedge}$. To check it, use the Postnikov tower:
positive-degree $C_p$-cohomology of every $q$-complete homotopy group
vanishes since multiplication by $p$ is invertible; at each tower stage
only constant maps contribute, and homotopy inverse limits preserve
the comparison. Thus only $q=2$ contributes to $\pi_2$ among these
primes, giving $\pi_2(Y_2^{\wedge})=\Z/2$.

At $q=p$ we use the exact finite-$p$-group mapping-space theorem:
for a finite $p$-group $P$, a compact Lie group $K$, and
$\alpha:P\to K$, the map induced by the centralizer identifies
\[
 \pi_n\operatorname{Map}(BP,(BK)_p^{\wedge})_{B\alpha}
  =\bigl(\pi_n B C_K(\alpha(P))\bigr)_p^{\wedge}
  \quad(n\ge2).
\]
This is the Dwyer--Zabrodsky--Notbohm formula, stated and derived in
\cite[Theorem 1.9 and its proof, using Theorem 3.2(iii) of their earlier
paper]{C_JMO1995}. Its hypotheses apply to the displayed finite cyclic
group and $SO(3)$. Since a nontrivial odd-order rotation has centralizer
exactly the axis circle $SO(2)$, the contribution is $\Z_p$.
This proves $\pi_1(G^{hC_p},1)=\Z_p\times\Z/2$.

Naturality of the centralizer map makes the $p$-component of the map
from the rotation circle the completion $\Z\to\Z_p$. At the prime
$2$ it is the map $\pi_1SO(2)\to\pi_1SO(3)$, namely reduction modulo
$2$: lifting a full rotation to unit quaternions gives the endpoint
$-1$. This proves the stated formula for the exact comparison.
For an explicit missing element at $p=3$, put
$\beta=\sum_{j\ge1}3^{j!}\in\Z_3$. Its base-three digits have
infinitely many zeros and infinitely many ones, whereas an ordinary
integer has eventually only zeros or eventually only twos. Hence
$(\beta,0)$ is outside the image. This also proves nonsurjectivity
without a cardinality argument; for general odd $p$ the same factorial
series has digits zero and one, rather than eventually zero or $p-1$. The explicit action, its finite order, the axis
centralizer and the finite reductions of this comparison are checked in
\texttt{code/verify\_sphere\_comparison.py}. The inverse-limit argument
above, rather than a finite computation, proves the nonsurjectivity.
\end{proof}

\begin{claim}{C.25}{Refuted}{Zero-level fidelity for a faithful infinite cyclic sphere action}
Let $\Gamma=\Z$ act faithfully on $S^2$ by an irrational axis rotation,
$u(1)(\theta,t)=(\theta+2\pi\alpha,t)$ with
$\alpha\in\R\setminus\Q$ and the usual collapse of the latitude
circles at $t=\pm1$. Let $M=A^\Gamma$ be its actual automorphism
space, based at the identity, and let $N$ be its weak automorphism
space in the categorical group $\Pi_1A$, where
$A=\operatorname{Homeo}(S^2)$. The induced transporter map has
\[
 \pi_1(M,1)=\Z\longrightarrow\pi_1(N,1)=\Z/2,
 \qquad n\longmapsto n\bmod2.
\]
Its kernel is $2\Z$, so level $0$ fails for this faithful action of an
infinite group. This is the exact fundamental-groupoid comparison;
C.24 separately retains the higher homotopy of $A$ when computing the
untruncated homotopy-fixed space. The hyperbolic zero-level question
recorded in C.12 has its stated, stronger surface hypothesis.
\end{claim}
\passage{\S{}23; folios 140--141; scans 279--281}

\begin{proof}
Commutation with an irrational rotation implies commutation with every
axis rotation: its powers are dense in $S^1$, and composition and
inversion in the compact-open homeomorphism group are continuous.
An equivariant homeomorphism preserves the two poles as a set and
sends each other rotation orbit to another orbit. In the orientation-
preserving identity component it has exactly the form
\[
 f(\theta,t)=(\theta+a(t),b(t)),\quad
 a\in C((-1,1),\R/2\pi\Z),\quad
 b\in\operatorname{Homeo}^+([-1,1]).
\]
To prove that this parametrization is onto, equivariance forces the
angular function on each orbit to be translation by a unique $a(t)$;
the orbit quotient supplies $b$. Orientation forces $b$ to be increasing.
Conversely every such pair extends continuously at both poles, because
those circle orbits collapse there, and the same argument applies to
its inverse. Thus the parametrization is a homeomorphism of spaces
with the compact-open topologies. Continuity at a pole imposes no
limit on $a(t)$; on any compact subinterval of $(-1,1)$ it is exactly
ordinary compact-open convergence.

The homotopy $a_s(t)=a((1-s)t)$ contracts the angular-function factor
to its evaluation circle at $0$. The convex interpolation
$b_s(t)=(1-s)b(t)+st$ contracts the increasing-homeomorphism factor
to the identity: it remains strictly increasing and fixes both endpoints.
These homotopies are continuous in the stated topologies and extend
continuously to the sphere. The actual identity centralizer therefore
has homotopy type $S^1$, with its axis-rotation circle giving the
specified generator of $\pi_1M$.

In the weak transporter groupoid, an automorphism of its identity
object is a loop in $A$ fixed by the conjugation action of $\Gamma$.
This is exactly the groupoid calculation of C.9, now applied to
$\Pi_1A$ rather than asserting any higher truncation of $A$ itself.
By the evaluation and polar-decomposition comparison in C.24,
$\pi_1(A,1)=\pi_1SO(3)=\Z/2$. Conjugation by an axis rotation acts
trivially on this group: joining that rotation to the identity in
$SO(3)$ gives a based homotopy of its conjugation map to the identity.
Hence $\pi_1N=\Z/2$.

The transporter comparison sends a loop of equivariant homeomorphisms
to that same loop in $A$. Its effect is consequently the usual
$\pi_1SO(2)\to\pi_1SO(3)$, reduction modulo two. More explicitly,
a loop of two complete rotations has degree two under the evaluation
$a\mapsto a(0)$ in the actual centralizer, so it is nonzero there.
Its lift to the quaternion double cover $S^3\to SO(3)$ is the closed
loop $t\mapsto\cos(2\pi t)+\mathbf i\sin(2\pi t)$, which contracts
in the simply connected $S^3$. It is therefore killed by the weak
transporter map. This proves the exact nontrivial kernel and the
zero-level failure, with no inference from disconnectedness.
\end{proof}

\begin{claim}{C.26}{Proved}{Transitive centralizers on a closed hyperbolic surface}
Let $X$ be a closed connected orientable surface of genus $g\geq2$,
and give $A=\operatorname{Homeo}(X)$ the compact-open topology.
For every subgroup $L\leq A$ acting transitively on $X$, the centralizer
$C_A(L)$ is finite. Consequently, for an arbitrary strict action
$u:\Gamma\to A$ and a based loop $f_t$ in $C_{A^0}(u(\Gamma))$,
the loop is equivariantly null-homotopic whenever the subgroup generated
by its values acts transitively on $X$.
There also exists a closed torsion-free Polish subgroup
$G\leq\operatorname{Homeo}(D^2,\partial D^2)$ with $\pi_1G=\Z$;
extension by the identity embeds it closedly in every hyperbolic
$A^0$, where its displayed generator becomes null-homotopic. This
last subgroup is not asserted to be a centralizer.
\end{claim}
\passage{23; folios 140--141; scans 279--281; Part C(b) infinite-operator request; exact partial reduction for C.12}

\begin{proof}
Put $K=\overline L$ and $N=C_A(K)$. Continuity of composition gives
$C_A(L)=C_A(K)$, and $K$ is a closed Polish subgroup of $A$ that
still acts transitively. Fix $x\in X$ and write $K_x$ for its stabilizer.
The open mapping principle for Polish actions gives a homeomorphism
\[
 K/K_x\longrightarrow X,\qquad kK_x\longmapsto kx.
                                                        \tag{C.26.1}
\]
Precisely, we use the ordinary-action case of
\cite[Theorem~3.5]{C_GomezPinedoUzcategui2017}: the acting group is
Polish, $X$ is nonmeagre Hausdorff, the action is continuous and
transitive, and the domain $K^x$ of the partial-action notation is all
of $K$, hence open. The source proves this open mapping statement;
no local section of (C.26.1) is being assumed.

Every $n\in N$ is determined by $n(x)$, since
$n(kx)=k n(x)$. Its possible values form exactly the closed subset
\[
 N_K(K_x)/K_x\ \subseteq\ K/K_x=X.                      \tag{C.26.2}
\]
Indeed, $n(x)=kx$ forces $K_x=K_{kx}=kK_xk^{-1}$.
Conversely a normalizing $k$ defines the homeomorphism
$gK_x\mapsto gkK_x$, commuting with $K$. The normalizer of the closed
subgroup $K_x$ is closed: both inclusions in
$kK_xk^{-1}=K_x$ are intersections of closed conditions. Its image
under (C.26.1) is closed because its inverse image under the quotient
map is that same normalizer. Thus (C.26.2) is compact.

Evaluation is a homeomorphism of $N$ onto (C.26.2). Only inverse
continuity needs proof. Fix $\varepsilon>0$. For every $y=g_yx$ choose
an identity neighborhood $W_y$ in $K$ such that each
$g_ywg_y^{-1}$, $w\in W_y$, moves every point of $X$ by less than
$\varepsilon/3$. The open sets $g_yW_yx$ cover $X$; take a finite
subcover with elements $g_1,\ldots,g_m$. When $n(x)$ is sufficiently
close to $n_0(x)$, each $g_i n(x)$ is within $\varepsilon/3$ of
$g_i n_0(x)$. At $z=g_iwx$ in the corresponding set, commutation gives
\[
 n(z)=g_iw n(x),\qquad n_0(z)=g_iw n_0(x).
\]
The two uniformly small movements by $g_iwg_i^{-1}$ and the middle
$\varepsilon/3$ bound imply $d(n(z),n_0(z))<\varepsilon$. This is
uniform in $z$ and proves inverse continuity in the compact-open
topology. In particular $N$ is compact.

An element of $N$ fixing one point fixes every point, by transitivity
and commutation. Every element of $A^0$ has Lefschetz number
$\chi(X)=2-2g\ne0$, hence has a fixed point. Therefore
\[
                       N\cap A^0=1.
\]
Projection to the discrete group $\pi_0A$ is thus a continuous
injection on $N$. Its image is compact and discrete, hence finite,
as asserted.

For the stated loop, put $L=\langle f_t:0\leq t\leq1\rangle$.
Every $u(\gamma)$ commutes with $L$, so the transitivity hypothesis
and the first assertion place $u(\Gamma)$ inside a finite group.
Write $H=u(\Gamma)$. The equality
$C_{A^0}(u(\Gamma))=C_{A^0}(H)$ is equality of subspaces, irrespective
of the cardinality of $\Gamma$. C.7 applied to the finite group $H$
contracts the loop in this exact centralizer. The assertion therefore
does not replace the operator group by a different action without a
comparison map.
\end{proof}

\paragraph{A proved obstruction to a closed-subgroup-only argument.}
There is a closed torsion-free Polish subgroup of
$\operatorname{Homeo}(D^2,\partial D^2)$ with fundamental group
$\Z$. We give the embeddings, because this example explains exactly
why closedness and absence of compact circles do not by themselves
settle C.12. It is not asserted to be a centralizer.

Let
\[
 E=\{f\in L^1([0,1],\R):f\text{ is integer-valued almost everywhere}\}.
\]
An $L^1$-convergent sequence has an almost-everywhere convergent
subsequence; hence $E$ is closed and Polish. For $0<s\leq1$ define
$D_sf(t)=f(t/s)$ on $[0,s]$ and zero on $(s,1]$, and put $D_0f=0$.
These maps remain in $E$ and satisfy $\|D_sf\|_1=s\|f\|_1$.
Strong continuity of dilation follows first for step functions and
then by $L^1$ approximation. Thus $(s,f)\mapsto D_sf$ contracts $E$
and each of its norm balls at zero. In particular $E$ is contractible
and locally contractible.

The constant integer functions form a closed discrete subgroup $\Z$.
The quotient $G=E/\Z$ is Polish, and $E\to G$ is a covering: translates
of a norm ball of radius less than $1/3$ by distinct integers are
disjoint. It is the universal cover, and
\[
 \pi_1(G)=\Z,\qquad
 \gamma(s)=[\mathbf1_{[0,s]}]\quad(0\leq s\leq1)          \tag{C.26.3}
\]
is a generator, because its lift joins zero to the constant function
$1$. Nevertheless $G$ is torsion-free: $n[f]=0$ says that $nf$ is
a constant integer, so the integer-valued $f$ is constant and $[f]=0$.

There is a closed isometric group embedding
\[
 G\longrightarrow V:=L^1([0,1],\R)/\R,
       \qquad[f]_{\Z}\longmapsto[f]_{\R},                 \tag{C.26.4}
\]
where the second denominator consists of all constant real functions.
For integer-valued $h$, the convex function $c\mapsto\int|h-c|$ is
linear between consecutive integers, tends to infinity at both ends,
and attains a minimum at an integer. Hence
\[
 \inf_{c\in\R}\|h-c\|_1=\inf_{n\in\Z}\|h-n\|_1,
\]
proving the isometry. The quotient metric on $G$ is complete: choose
a Cauchy subsequence with successive quotient distances below $2^{-j}$
and adjust each representative by an integer so that the successive
$L^1$ distances have that bound. The representatives converge in $E$,
and the entire original Cauchy sequence then converges in $G$.
This proves closedness in the Banach space $V$.

We spell out a closed linear isometric embedding
$T:V\hookrightarrow C_0([0,1],\R)$, where the subscript here means
vanishing at both endpoints. The weak-star unit ball $B$ of $V^*$ is
compact metrizable, and $v\mapsto(\ell\mapsto\ell(v))$ embeds $V$
isometrically into $C(B)$. A continuous surjection from the Cantor
set to $B$ exists by nested finite covers of decreasing diameters.
Pullback followed by linear interpolation across each complementary
interval of the Cantor set embeds $C(B)$ linearly and isometrically
in $C([0,1])$. Place this interval in the middle third and interpolate
to zero at the two outer endpoints. The resulting map is the required
$T$; its image is closed because $V$ is complete.

For $F\in C_0([0,1])$, the formula
\[
 \Phi_F(x,y)=
   \bigl(\tanh(\operatorname{artanh}(x)+F(y)),y\bigr)
          \qquad(|x|<1)
                                                        \tag{C.26.5}
\]
on $[-1,1]\times[0,1]$, extended by the identity on the vertical
sides, fixes the entire boundary and satisfies
$\Phi_F\Phi_Q=\Phi_{F+Q}$. It is a topological-group embedding:
$\tanh$ is Lipschitz, giving continuity, and evaluation at $(0,y)$
recovers $\tanh(F(y))$, from which $F$ is recovered uniformly near
any given profile. Its image is closed. A homeomorphism limit fixes
the boundary and preserves the horizontal leaves; the image of the
compact segment $\{0\}\times[0,1]$ stays a positive distance from the
vertical sides. Applying $\operatorname{artanh}$ therefore recovers
uniformly bounded, uniformly convergent profiles.

Composing (C.26.4), $T$, and (C.26.5) embeds $G$ as a closed subgroup
of $\operatorname{Homeo}(D^2,\partial D^2)$. Extension by the identity
outside a coordinate disc is a closed topological-group embedding
into $A^0$ for any hyperbolic surface: restriction recovers the maps,
and being the identity off the disc is a closed condition. Its loop
(C.26.3) is ambient-null but nontrivial in $G$.
The precise enlargement killing it is already visible: if
$F_s=T([\mathbf1_{[0,s]}]_{\R})$, then $F_0=F_1=0$ and
$(r,s)\mapsto\Phi_{rF_s}$ is a based contraction in the larger real
shear group. No operator group cutting out a centralizer that excludes
this contraction has been constructed here. Thus this example refutes
the auxiliary assertion of fundamental-group injectivity for all
closed subgroups, while C.12 retains its stated centralizer question.

\begin{claim}{C.27}{Proved}{An order-natural contraction of line centralizer loops}
Give $\operatorname{Homeo}(\R)$ its compact-open topology, let
$H\leq\operatorname{Homeo}(\R)$ be an arbitrary subgroup, and put
\[
              C=C_{\operatorname{Homeo}^+(\R)}(H).
\]
Every based loop in $C$ has an explicit based null-homotopy in $C$.
Orientation-reversing elements of $H$ are allowed. The construction
commutes with conjugation by every homeomorphism of the line and
requires no finiteness, differentiability, or local compactness of $H$.
\end{claim}
\passage{23; folios 140--141; scans 279--281; Part C(b) infinite-operator request; exact partial reduction for C.12}

\begin{proof}
Let $t\mapsto f_t$ be a compact-open continuous loop in $C$, with
$f_0=f_1=\id$. For $0\leq r\leq1$ define
\[
 m_r(x)=\min_{0\leq u\leq r}f_u(x),\qquad
 M_r(x)=\max_{0\leq u\leq r}f_u(x).                    \tag{C.27.1}
\]
Both are increasing homeomorphisms. Indeed, for $x<y$, compactness
gives
\[
 \min_{0\leq u\leq1}\bigl(f_u(y)-f_u(x)\bigr)>0,
\]
which proves strict increase for each minimum and maximum. Their
properness is uniform in $r$: compactness of the inverse family
$\{f_u^{-1}:0\leq u\leq1\}$ gives, for every bounded target interval,
a common bound on all its inverse images. Continuity of (C.27.1),
jointly in $r$ and $x$, follows from uniform continuity on compact
parameter rectangles. Proper continuous strictly increasing maps of
$\R$ are homeomorphisms.

Set
\[
 F_{s,t}(x)=\operatorname{med}
      \bigl(m_{1-s}(x),f_t(x),M_{1-s}(x)\bigr)
   =\max\{m_{1-s}(x),\min\{f_t(x),M_{1-s}(x)\}\}.
                                                        \tag{C.27.2}
\]
Finite minima and maxima of increasing homeomorphisms are again
increasing homeomorphisms, by the same strict-increase and properness
arguments. Formula (C.27.2) is continuous in $(s,t,x)$, uniformly on
compact $x$-sets; it is therefore a continuous family in the
compact-open homeomorphism group. At $s=0$, the prefix interval is
the entire loop, so $m_1\leq f_t\leq M_1$ and $F_{0,t}=f_t$.
At $s=1$, $m_0=M_0=\id$, giving $F_{1,t}=\id$. Since $u=0$ occurs
in every prefix interval, $m_r\leq\id\leq M_r$, and hence
$F_{s,0}=F_{s,1}=\id$ for every $s$.

It remains to check the exact equivariance. An increasing
$h\in H$ fixes both envelopes under conjugation, because conjugation
commutes with their pointwise minimum and maximum and fixes every
$f_u$. A decreasing $h\in H$ interchanges $m_r$ and $M_r$ instead.
Every order-preserving or order-reversing bijection preserves the
median of three points. Thus both cases give
\[
                   h F_{s,t}h^{-1}=F_{s,t}
       \qquad(h\in H, 0\leq s,t\leq1).
\]
This proves that the entire based contraction lies in $C$. Replacing
the loop by its conjugate by an arbitrary line homeomorphism gives
the conjugate of the same formula, proving the asserted naturality.
\end{proof}

The surface problem in C.12 still requires a continuous construction
across different orbit strata. Applying (C.27.2) separately after
choosing coordinates on individual line orbits does not establish
continuity on the surface. No such gluing statement follows from
C.27, and none is used here.

\begin{claim}{C.28}{Proved}{Orbit types of a loop-generated plane group}
Let $t\mapsto f_t$ be a compact-open continuous based loop in
$\operatorname{Homeo}^+(\R^2)$, with $f_0=f_1=\id$, and put
\[
 L=\{f_t:0\leq t\leq1\},\qquad P=L\cup L^{-1},
 \qquad G=\langle P\rangle.
\]
Every $G$-orbit is one of the following: an open subset of the plane;
a point; a Jordan curve; or the image of a continuous injective map
$\eta:\R\to\R^2$ that restricts to a homeomorphism on every compact
interval. In the last case its intrinsic arc order identifies it with
the line, but no global homeomorphism from $\R$ to the orbit with its
subspace topology is asserted. The loop restricts continuously in the
compact-open topology to the intrinsically ordered line. Every
homeomorphism commuting with $G$ carries these orbit types to the same
types and induces a homeomorphism between the corresponding intrinsic
lines, allowing reversal of their orders.
\end{claim}
\passage{23; folios 140--141; scans 279--281; Part C(b) infinite-operator request; exact partial reduction for C.12}

\begin{proof}
Fix $x\in\R^2$ and put
\[
                  K_n=P^n x,\qquad O=Gx=\bigcup_{n\geq1}K_n.
                                                        \tag{C.28.1}
\]
The sets $K_n$ are nested compact Peano continua. Indeed, $P$ is a
continuous image of the wedge of two circles, using the given loop
and its inverse; it is compact, connected and locally connected.
A finite product of such spaces has the same properties, and the
word-evaluation map has image $K_n$. To recall why local connectedness
passes through this compact surjection, for an open neighborhood of a
point in the image cover its compact inverse fibre by finitely many
connected open subsets of its inverse image. Their images all meet
that point and lie in its component. The complement of the image of
the remaining compact set is a smaller open neighborhood in that
component. Thus components of open sets in the image are open.
In particular $K_n$ is arcwise connected. The identity lies in $P$,
which proves the nesting in (C.28.1). Every element of $G$ can be
joined to the identity by a path in $G$, by writing it as a finite word
and using the specified paths to its factors.

Suppose first that $O$ contains a Jordan curve $J$. For any $g\in G$
choose a path $g_s$ from $\id$ to $g$, so that $g_sJ\subseteq O$.
Let $D_s$ be the bounded Jordan domain of $g_sJ$. If $gJ\ne J$, then
$D_0\mathbin\triangle D_1$ contains a nonempty open set away from the
two boundary curves. Every point $q$ in this open set lies on some
$g_sJ$: otherwise the winding number of the continuously varying curve
$g_sJ$ around $q$ would be constant, whereas its endpoint values are
$0$ and $1$. Hence $O$ contains a nonempty open set. Transitivity of
$G$ on $O$ carries this open set to a neighborhood of each point of
$O$, so $O$ is open. If every $g$ preserves $J$, then for a point
$y\in J$ its orbit is contained in $J$; since $Gy=O$ and $J\subseteq O$,
we obtain $O=J$. Thus a nonopen orbit containing a Jordan curve is
itself that Jordan curve.

We may now consider a nonopen orbit containing no Jordan curve.
Each $K_n$ is a dendrite: a compact, connected, locally connected
metric space containing no Jordan curve. Arcs between two points are
unique, since two distinct arcs would contain a Jordan curve. This
also proves uniqueness of arcs in all of $O$: two proposed arcs in
$O$ would give the same forbidden curve. Existence follows by putting
their endpoints into one $K_n$.

Each $K_n$ has countably many branch points. Choose a countable dense
subset $D_n$. Three distinct components of the complement of a branch
point are open there, so choose a point of $D_n$ in each. The branch
point is the unique common point of the three pairwise connecting
arcs for this triple: two branches meet only at their point of
separation. Thus each branch point is determined by one of the
countably many triples from $D_n$.

A triod point of $O$ is a branch point of some $K_n$. Put the three
outer endpoints of its triod into one $K_n$; the unique connecting
arcs in $K_n$ must be the arcs of that triod. The triod points of $O$
are consequently countable. Their set is invariant under every element
of $G$. If it were nonempty, transitivity would make every point of
$O$ a triod point. A nontrivial arcwise connected orbit contains an
arc and is uncountable, giving a contradiction. Thus a nontrivial
$O$ has no triod points, and each nondegenerate $K_n$ is an arc.

The nested arcs give $O$ a consistent linear arc order, unique up to
reversal. To see the absence of additional branches directly, fix two
distinct points $a,b$. Every other point is on $[a,b]$, beyond $a$,
or beyond $b$, because any fourth alternative would give a triod.
There are at most two endpoints of this ordered union. The endpoint
set is invariant under $G$, so transitivity and uncountability exclude
endpoints. We can therefore pass to a cofinal sequence of arcs $A_j$
such that $A_j$ lies in the relative interior of $A_{j+1}$: extend
both endpoints in $O$ and then put the two extensions and the next
$K_n$ into a later $K_m$.

Choose compatible homeomorphisms $[-j,j]\to A_j$, extending at each
step along the two added arc intervals. Their union is a continuous
injective map
\[
                 \eta:\R\longrightarrow\R^2,
                        \qquad\eta(\R)=O,              \tag{C.28.2}
\]
whose restriction to every compact interval is a homeomorphism onto
its image. These assertions follow on each finite interval from one
of the compact arc parametrizations. This supplies the stated
intrinsic line structure. It does not prove that the inverse of
(C.28.2) is continuous for the global subspace topology; accumulations
from far along an orbit have not been excluded.

The restricted loop is continuous for that intrinsic line topology.
A compact arc in $O$ is contained in some $K_m$, and
$f_t(K_m)\subseteq K_{m+1}$ for all $t$. The latter is a compact arc,
on which the intrinsic and subspace topologies coincide. Joint
continuity of $(t,z)\mapsto f_t(z)$ on each such compact arc therefore
gives compact-open continuity of
$\eta^{-1}f_t\eta\in\operatorname{Homeo}(\R)$.
Its order sign is positive throughout the loop: images of two distinct
points vary continuously in the line and cannot cross. Finally a
commuting homeomorphism $h$ satisfies the exact equality
\[
                  h(K_n(x))=K_n(hx).
\]
It therefore maps compact arcs homeomorphically to compact arcs and
induces a homeomorphism of the intrinsic ordered lines; its inverse
has the same property. It also preserves the point, Jordan-curve and
open alternatives, completing the proof.
\end{proof}

\paragraph{Exact application and remaining step.}
The canonical lift in the paragraph after C.12 sends a hyperbolic
centralizer loop to a closed based loop in
$\operatorname{Homeo}^+(\widetilde U)$, where
$\widetilde U\cong\R^2$. The entire group $\widetilde H$ of lifts of
the operator image commutes with it and hence with the $G$ of C.28.
Thus C.28 applies to this actual lift group, and C.27 provides an
order-natural contraction on each intrinsic line orbit after
restriction to its stabilizer in $\widetilde H$.
These are typed restriction and lift maps, but they do not yet give a
homotopy through homeomorphisms of the whole plane. Continuity across
distinct orbit strata, compatibility with their frontiers, and a
relative contraction on proper open orbits remain unproved.
C.26 treats global transitivity on a closed hyperbolic surface; a
proper open orbit in C.28 does not satisfy that hypothesis. The
admissible infinite-group level-$0$ question therefore remains C.12.

\begin{claim}{C.29}{Proved}{Circle centralizers and the winding of a plane loop}
For every subgroup $H\leq\operatorname{Homeo}(S^1)$, put
$C=C_{\operatorname{Homeo}^+(S^1)}(H)$, with the compact-open topology.
The inclusion induces an injection
\[
 \pi_1(C,\id)\longrightarrow
 \pi_1(\operatorname{Homeo}^+(S^1),\id)=\Z,              \tag{C.29.1}
\]
where the integer is the winding of the image of one point. Its image
lies in $\Z^H$, with $H$ acting on $\Z$ by its orientation character.
A based loop of winding zero has an explicit contraction in $C$ that
is natural under all homeomorphisms between circles.
In particular, a based null-homotopic loop in
$\operatorname{Homeo}^+(\R^2)$ preserving a Jordan curve restricts to
a winding-zero loop on that curve. Its restriction therefore has this
natural contraction against every commuting circle homeomorphism.
\end{claim}
\passage{23; folios 140--141; scans 279--281; Part C(b) infinite-operator request; exact partial reduction for C.12}

\begin{proof}
Write the circle as $\R/\Z$, and set
\[
 \widetilde A=\{F\in\operatorname{Homeo}^+(\R):
                        F(x+1)=F(x)+1\}.
\]
The quotient map $\widetilde A\to\operatorname{Homeo}^+(S^1)$ is a
covering with deck transformations $F\mapsto F+n$, $n\in\Z$.
Indeed, requiring $F(0)$ to lie in a fixed interval of length less
than one gives continuous unique lifts over the corresponding
evaluation neighborhood. The homotopy
$(r,F)\mapsto(1-r)F+r\id$ contracts $\widetilde A$: strict
increase, the translation equation, and properness hold at every
parameter. Thus the target fundamental group in (C.29.1) is precisely
$\Z$. For a based loop $f_t$, its unique lift $F_t$ starting at
$\id$ ends at $x\mapsto x+n$, where $n$ is its point winding.

Suppose that $f_t\in C$ and $n=0$. Every lift $\widetilde h$ of an
element $h\in H$, including an orientation-reversing one, commutes
with every $F_t$. The commutator projects to the identity on the
circle, so it is an integer translation depending continuously on
$t$; it is the identity at $t=0$, and hence at all $t$.
Consequently $F_t$ is a based loop in the line centralizer of the
entire lift group of $H$. This group also contains the integer
translations. Apply C.27, explicitly
\[
 \begin{split}
 m_r(x)&=\min_{0\leq u\leq r}F_u(x),\qquad
 M_r(x)=\max_{0\leq u\leq r}F_u(x),\\
 Q_{s,t}(x)&=\operatorname{med}
       \bigl(m_{1-s}(x),F_t(x),M_{1-s}(x)\bigr).
 \end{split}                                                   \tag{C.29.2}
\]
Every $Q_{s,t}$ commutes with all lifts and with $x\mapsto x+1$,
so (C.29.2) descends to a contraction through elements of $C$.
This proves injectivity in (C.29.1).

The winding of a conjugated loop $h f_t h^{-1}$ is the winding of
$f_t$ multiplied by the orientation sign of $h$. Since these loops
agree for $h\in H$, the image is contained in $\Z^H$. For naturality,
let $q$ be a homeomorphism between two circles and choose any lift
$\widetilde q$ to their universal lines. The based lift of the
conjugated loop is $\widetilde q F_t\widetilde q^{-1}$.
An increasing $\widetilde q$ carries the two extrema in (C.29.2)
to the corresponding extrema; a decreasing one interchanges them.
Every line homeomorphism preserves the median of three points.
Thus (C.29.2) is carried to exactly the same formula for the
conjugated loop. Changing the lift of $q$ by an integer translation
does not change the descended contraction.

It remains to prove the stated plane winding comparison. Conjugate
by a Schoenflies homeomorphism so that the preserved Jordan curve is
the unit circle. Choose two distinct points $a,b$ on it. Evaluation
and the difference direction give the continuous maps
\[
 \operatorname{Homeo}^+(\R^2)
 \longrightarrow\operatorname{Conf}_2(\R^2)
 \longrightarrow S^1,
 \qquad
 f\longmapsto(f(a),f(b))
 \longmapsto\frac{f(b)-f(a)}{|f(b)-f(a)|}.                 \tag{C.29.3}
\]
For the given circle-preserving loop choose continuous angular lifts
$\alpha(t),\beta(t)$ with
$\alpha(t)<\beta(t)<\alpha(t)+2\pi$.
They both acquire the increment $2\pi n$: their difference remains
in $(0,2\pi)$ and returns to its initial value modulo $2\pi$.
The identity
\[
 e^{i\beta(t)}-e^{i\alpha(t)}
 =2i\,e^{i(\alpha(t)+\beta(t))/2}
            \sin\frac{\beta(t)-\alpha(t)}2
\]
shows that the winding of (C.29.3) is exactly $n$.
A null-homotopy of the ambient plane loop maps under (C.29.3) to a
null-homotopy in $S^1$, forcing $n=0$. This proves the exact
comparison and the last assertion.
\end{proof}

\paragraph{Application to C.12.}
The canonical lift of a hyperbolic centralizer loop is null-homotopic
as a loop of plane homeomorphisms: it is the image under the continuous
canonical lift homomorphism of a loop in the contractible $A^0$.
Every Jordan-curve orbit in C.28 is preserved throughout this loop.
Its restriction therefore has winding zero by (C.29.3), and (C.29.2)
contracts it naturally against its stabilizer in the full operator
lift group. Commuting operators carrying one circle orbit to another
carry these contractions to each other. This gives a circle-orbit
contraction with the exact required equivariance. C.34 subsequently
proves joint continuity, including inverses, on the full nonopen
locus. C.35 shows that this particular prescription need not admit
an equivariant extension over an adjacent open orbit, and C.36
constructs compatible alternatives on finite unions of closures
with two-dimensional proper operator image. The remaining global
compatibility requirements are stated in C.12, which remains open.

\begin{claim}{C.30}{Proved}{Properness of a transitive-action centralizer}
Let $V$ be a connected second-countable surface without boundary,
and let $K$ be a closed subgroup of $\operatorname{Homeo}(V)$ acting
transitively on $V$, with the compact-open topology. Put
$N=C_{\operatorname{Homeo}(V)}(K)$. Then $N$ is locally compact,
its action on $V$ is free and proper, and for every $x\in V$
evaluation is a homeomorphism onto the closed subset
\[
 N\xrightarrow{\ \mathrm{ev}_x\ }
 N_K(K_x)/K_x\ \subseteq\ K/K_x\cong V.                  \tag{C.30.1}
\]
Under the normalizer-quotient group structure, this is an
anti-isomorphism of topological groups. For an open orbit $O$ of the
loop-generated group in C.28, the image in $\operatorname{Homeo}(O)$
of its operator stabilizer lies in such an $N$; the closure of that
image acts freely and properly on $O$.
\end{claim}
\passage{23; folios 140--141; scans 279--281; Part C(b) infinite-operator request; exact partial reduction for C.12}

\begin{proof}
The group $K$ is Polish. The open mapping theorem
\cite[Theorem~3.5]{C_GomezPinedoUzcategui2017}, with the ordinary
continuous transitive action on the nonmeagre Hausdorff surface $V$,
identifies $K/K_x$ homeomorphically with $V$. In particular the
orbit map $k\mapsto kx$ is open. We will use the following exact
consequence, requiring no local continuous section: for every
sequence $x_j\to x$, there are $k_j\to1$ in $K$ with $k_jx=x_j$.
To prove it, choose a decreasing identity-neighborhood basis $W_m$.
The open sets $W_mx$ contain $x$, so successive tails of the sequence
lie in them. Choose $k_j$ from the corresponding $W_{m(j)}$, with
$m(j)\to\infty$.

An element of $N$ is determined by its value at $x$, since
$n(kx)=kn(x)$. If $n(x)=ax$, equality of stabilizers gives
$aK_xa^{-1}=K_x$. Conversely, for $a$ in this normalizer the map
$kK_x\mapsto kaK_x$ is a homeomorphism commuting with $K$.
The normalizer is closed: both inclusions defining
$aK_xa^{-1}=K_x$ are intersections of closed conditions.
It is saturated for the quotient by $K_x$, so its image is a closed
subset of $K/K_x=V$. This proves the claimed set-theoretic image
and closedness in (C.30.1). If $n(x)=ax$ and $m(x)=bx$, then
$(nm)(x)=bax$, proving the stated anti-isomorphism.

We prove inverse continuity of evaluation, including the issue of
noncompact $V$. Suppose that $n_j(x)\to n(x)$, with $n_j,n\in N$.
For any compact set $B\subset V$, a failure of uniform convergence
on $B$ would provide a subsequence and $y_j\in B$ with
$y_j\to y$ witnessing that failure. Choose $a\in K$ with $ax=y$,
and then $b_j\to1$ with $b_jy=y_j$, using open evaluation at $y$.
Commutation gives
\[
 n_j(y_j)=n_j(b_jax)=b_ja\,n_j(x)\longrightarrow a n(x)=n(y).
                                                               \tag{C.30.2}
\]
Also $n(y_j)\to n(y)$, a contradiction. Hence $n_j\to n$ uniformly
on every compact set, which is compact-open convergence.
Inversion can be checked directly as well. Put $c_j=n^{-1}n_j$,
so $c_jx\to x$, and choose $a_j\to1$ with $a_jx=c_jx$.
Because $c_j$ commutes with $a_j$,
$c_j^{-1}x=a_j^{-1}x\to x$; applying (C.30.2) to $c_j^{-1}$
gives convergence of the inverses. All the spaces are metrizable,
so this sequential argument proves inverse continuity.
Thus (C.30.1) is a homeomorphism, and $N$ is locally compact as a
closed subspace of the locally compact surface.

If $n\in N$ fixes one point, transitivity and commutation make it
fix every point, proving freeness. To prove properness, consider
sequences with
\[
                  x_j\longrightarrow x,\qquad
                  n_jx_j\longrightarrow y.
\]
Choose $k_j\to1$ with $k_jx=x_j$. Then
\[
                n_jx=k_j^{-1}n_jx_j\longrightarrow y.    \tag{C.30.3}
\]
The image of evaluation at $x$ is closed, so $y=nx$ for a unique
$n\in N$. Its inverse continuity gives $n_j\to n$.
For any compact subset of $V\times V$, every sequence in its
inverse image under $(n,z)\mapsto(z,nz)$ therefore has a convergent
subsequence, with limit in the same inverse image. Metrizability
turns this sequential compactness into compactness, proving that
the action map is proper. Every closed subgroup of $N$ inherits
the free proper action.

Finally let $G$ be the group in C.28 and let $O$ be one of its open
orbits. It is connected because every element of $G$ is joined to
the identity by a path in $G$. The closure
\[
        K=\overline{\{g|_O:g\in G\}}
                       ^{\operatorname{Homeo}(O)}
\]
is a closed Polish group still transitive on $O$. Let $\widetilde H$
be any group commuting with $G$, and let $\widetilde H_O$ be its
setwise stabilizer of $O$. The precise restriction homomorphism is
\[
 \widetilde H_O\longrightarrow\operatorname{Homeo}(O),\qquad
 h\longmapsto h|_O,
 \quad
 \ker=\{h\in\widetilde H_O:h|_O=\id\}.                 \tag{C.30.4}
\]
Its image commutes with the restriction of $G$ and hence with $K$.
The closure of this image is a closed subgroup of $N$, and therefore
acts freely and properly on $O$, as claimed.
\end{proof}

\paragraph{Precise remaining use.}
The closure in the last assertion has its subspace topology from
$\operatorname{Homeo}(O)$. Formula (C.30.4) does not assert that the
original discrete operator group acts freely, since its kernel may
be nontrivial; nor does it equip a possibly dense abstract image with
the discrete topology and claim properness for that topology.
The result supplies the closed locally compact transformation group
actually present on an open orbit. No Lie-group classification or
quotient-bundle contraction is invoked here. A contraction relative
to orbit frontiers, and its continuous assembly with C.27 and C.29
on the one-dimensional strata, remain unproved. Thus C.30 does not
change the unresolved status of C.12.

\begin{claim}{C.31}{Proved}{Track control and continuity at the fixed stratum}
Let $f_t$ be a based null-homotopic loop in
$\operatorname{Homeo}^+(\R^2)$, let $G$ be its generated group as in
C.28, and let $Y$ be the union of the nonopen $G$-orbits. This is a
closed subset of the plane. Define $Q_{s,t}$ on $Y$ by C.27 on each
intrinsic line orbit, by the descended formula of C.29 on each circle
orbit, and by the identity on point orbits. Then every $Q_{s,t}$ is
a bijection of $Y$ preserving each orbit, and
\[
 Q_{s,t}^{-1}=Q_{s,t}[f^{-1}],\qquad
 Q_{s,t}(x)\in\{f_u(x):0\leq u\leq1\},\qquad
 Q_{s,t}^{-1}(x)\in\{f_u^{-1}(x):0\leq u\leq1\}.
                                                        \tag{C.31.1}
\]
Here $Q[f^{-1}]$ means the same construction applied to the based
inverse loop $u\mapsto f_u^{-1}$, which generates the same group.
Both evaluation maps $(s,t,x)\mapsto Q_{s,t}^{\pm1}(x)$ are continuous
at every point of $[0,1]^2\times Z$, where
$Z=\{x:f_u(x)=x\text{ for all }u\}$ is the point-orbit stratum.
The construction commutes with every homeomorphism commuting with
all $f_u$. No continuity at a nonfixed orbit is asserted here.
\end{claim}
\passage{\S{}23; folios140--141; scans279--281; Part C(b) infinite-operator request; fixed-stratum gluing reduction for C.12}

\begin{proof}
The union of open orbits is open, proving that $Y$ is closed.
C.28 supplies exactly the remaining point, circle and intrinsic-line
types. The circle restrictions have zero winding by C.29, since
the ambient loop is null-homotopic. Thus the formulas defining $Q$
are available on every orbit in $Y$, and are independent of the
chosen line or circle coordinates by their proved naturality.
Each formula gives a homeomorphism of its single intrinsic line
or circle, so together they give a bijection of $Y$.

We prove the inverse identity before making any continuity claim.
For increasing homeomorphisms of a line, pointwise inversion
reverses order and exchanges finite minimum and maximum:
\[
 (a\wedge b)^{-1}=a^{-1}\vee b^{-1},\qquad
 (a\vee b)^{-1}=a^{-1}\wedge b^{-1}.
                                                        \tag{C.31.2}
\]
For example, at a point $x$ the first inverse is the larger of
$a^{-1}(x)$ and $b^{-1}(x)$: at this larger point one of $a,b$
equals $x$, and the other is at least $x$.
The same argument works for compact families, whose extrema
and inverse extrema are attained. Therefore, with the notation
of C.27,
\[
 m_r^{-1}(x)=\max_{0\leq u\leq r} f_u^{-1}(x),\qquad
 M_r^{-1}(x)=\min_{0\leq u\leq r} f_u^{-1}(x).
                                                        \tag{C.31.3}
\]
The median of three line homeomorphisms is the lattice expression
$(a\wedge b)\vee(a\wedge c)\vee(b\wedge c)$; its dual expression
$(a\vee b)\wedge(a\vee c)\wedge(b\vee c)$ has the same value.
Applying (C.31.2) thus proves
\[
 \operatorname{med}(a,b,c)^{-1}
       =\operatorname{med}(a^{-1},b^{-1},c^{-1}).
                                                        \tag{C.31.4}
\]
Combining (C.31.3)--(C.31.4), and using symmetry of the median
in its two envelope entries, gives the first identity in (C.31.1)
on every line orbit. On a circle, the inverse of the unique based
lift $F_u$ is exactly the unique based lift of $f_u^{-1}$.
The same calculation on the universal line descends to the circle.
It also holds on point orbits.

At any point the median chooses one of its three entries. Each
envelope entry is $f_u(x)$ for a parameter attaining its extremum,
and the middle entry is $f_t(x)$ itself. This proves the second
assertion of (C.31.1) on lines; projection of the same statement
for the based lifts proves it on circles. The inverse identity
then proves the third assertion.

Fix $p\in Z$ and an open neighborhood $V$ of $p$ in the plane.
Joint continuity of the action, compactness of $[0,1]$, and
$f_u(p)=p$ give a neighborhood $W$ of $p$ such that
$f_u(W)\subseteq V$ for every $u$. The compact inverse family
gives the same assertion for all $f_u^{-1}$ after shrinking $W$.
By (C.31.1),
\[
 Q_{s,t}^{\pm1}(W\cap Y)\subseteq V\cap Y
                   \quad\text{for all }s,t\in[0,1].       \tag{C.31.5}
\]
Since $Q_{s,t}^{\pm1}(p)=p$, this proves joint continuity at every
$(s,t,p)$, uniformly in the two homotopy parameters.

A homeomorphism commuting with all $f_u$ permutes the $G$-orbits,
preserves their four types, and maps intrinsic lines by line
homeomorphisms by C.28. Naturality in C.27 and C.29 gives
$hQ_{s,t}h^{-1}=Q_{s,t}$ on all of $Y$. This proves the asserted
equivariance without any continuity assertion at the remaining
points.
\end{proof}

\paragraph{Scope of the gluing step.}
Formula (C.31.5) resolves continuity at the fixed stratum for the
orbitwise construction on $Y$. It does not yet make the bijections
$Q_{s,t}:Y\to Y$ homeomorphisms, because continuity at the nonfixed
line and circle strata remains unproved. Nor does it control a
contraction on the omitted open orbits. These are the remaining
global assembly problems for this route to C.12.

\begin{claim}{C.32}{Proved}{Discrete proper operator groups on open orbits}
Let $O$ be a connected noncompact second-countable surface without
boundary, and let a discrete group $Q\leq\operatorname{Homeo}(O)$ act
freely and properly. With compact-open topology on the homeomorphism
groups, inclusion induces an injection
\[
 \pi_1\!\left(C_{\operatorname{Homeo}(O)}(Q),\id\right)
       \longrightarrow \pi_1(\operatorname{Homeo}(O),\id).
                                                        \tag{C.32.1}
\]
Moreover, a based null-homotopic loop in
$\operatorname{Homeo}^+(\R^2)$ preserving a connected open subset
$O\subseteq\R^2$ restricts to a null-homotopic loop in
$\operatorname{Homeo}(O)$.
Consequently, for a hyperbolic centralizer loop and an open orbit in
C.28 whose closed restricted operator image in C.30 is discrete, its
restriction has a based equivariant null-homotopy on that orbit.
The original restriction kernel is retained; no extension of this
null-homotopy to the orbit frontier is asserted.
\end{claim}
\passage{Section 23, folios 140-141, scans 279-281; Part C(b) infinite-operator request; discrete-closure open-orbit reduction for C.12}

\begin{proof}
First consider the free proper discrete action on $O$. It gives an
actual covering
\[
                       p:O\longrightarrow V=O/Q.          \tag{C.32.2}
\]
Here discreteness refers to the topology in the proper-action
hypothesis. For completeness, a compact neighborhood of a point
meets only finitely many of its translates, by properness. Shrink
an open neighborhood inside it to avoid the nonidentity translates
among this finite set, using freeness. Distinct translates of the
result are disjoint, which gives the evenly covered neighborhoods
for (C.32.2). The quotient $V$ is Hausdorff by properness, is a
surface without boundary by these local charts, and is second
countable because images of a countable open basis of $O$ form an
open basis. The connected covering is regular, with deck group
exactly $Q$.

We record the homeomorphism-group input precisely. Every
second-countable surface admits a PL triangulation. Applied with
the compact subpolyhedron consisting of one point $v$, the compact
case of \cite[Lemma~2.1]{A_Yagasaki2000} and the noncompact case of
\cite[Theorem~1.1]{A_Yagasaki2000} give
\[
 \operatorname{Homeo}(V,v)^0\simeq *
 \quad\text{for }V\not\cong S^2,\ \R P^2,\ \R^2.          \tag{C.32.3}
\]
These are all exceptions relevant to a surface without boundary
with that one-point subpolyhedron. Evaluation
$\operatorname{Homeo}(V)^0\to V$ is a locally trivial bundle:
in a coordinate disc a continuously chosen homeomorphism supported
in that disc moves its center to a nearby point, giving a local
section. Its identity fibre component is exactly the group in
(C.32.3). The homotopy exact sequence therefore gives an injection
\[
 \pi_1(\operatorname{Homeo}(V),\id)
          \longrightarrow\pi_1(V,v)
       \quad(V\not\cong S^2,\R P^2,\R^2).                \tag{C.32.4}
\]
The base $V$ of (C.32.2) cannot be $S^2$ or $\R P^2$:
their fundamental groups are finite, so every connected covering
has finitely many sheets and compact total space, contrary to
noncompactness of $O$. If $V=\R^2$, its connected covering is
one-sheeted, so $Q=1$ and (C.32.1) is the identity homomorphism.

Let $f_t$ be a based loop in the centralizer that is null-homotopic
in $\operatorname{Homeo}(O)$. It descends to a based loop $b_t$
on $V$, defined by the exact equation
\[
                         b_t p=p f_t.                    \tag{C.32.5}
\]
The family is continuous: joint continuity descends through the
open quotient map $[0,1]\times O\to[0,1]\times V$, and the
compact-open exponential law applies since $V$ is locally compact.
For $x\in O$, the evaluation loop $f_t(x)$ is null-homotopic in
$O$. Thus $b_t(p(x))=p(f_t(x))$ is null-homotopic in $V$.
By (C.32.4), $b_t$ is a null-homotopic based loop of homeomorphisms
of $V$. Choose a continuous filling
$B:D^2\to\operatorname{Homeo}(V)^0$ with the prescribed boundary
$b_t$, and a point $z_*\in\partial D^2$ corresponding to the based
endpoint, so that $B_{z_*}=\id$.

The lift is a filling through homeomorphisms and through the exact
centralizer; we verify both assertions. Form
\[
 F:O\times D^2\longrightarrow V,\qquad
                       F(x,z)=B_z(p(x)).                 \tag{C.32.6}
\]
At $(x_*,z_*)$, the induced fundamental-group image is
$p_*\pi_1(O,x_*)$, since $D^2$ is simply connected and
$F(x,z_*)=p(x)$. The covering lifting criterion gives a unique
continuous lift $\widetilde B:O\times D^2\to O$ with
$\widetilde B(x,z_*)=x$.
Apply the identical construction to
$F^-(x,z)=B_z^{-1}(p(x))$, giving a continuous lift
$\widetilde C$ that is the identity at $z_*$.
The maps
$(x,z)\mapsto\widetilde B_z\widetilde C_z(x)$ and
$(x,z)\mapsto x$ lift the same map $(x,z)\mapsto p(x)$
and agree at the basepoint. Uniqueness of lifts makes them equal;
the reverse composition is equal for the same reason.
Hence every $\widetilde B_z$ is a homeomorphism, with inverse
$\widetilde C_z$, and the family is compact-open continuous.

For $q\in Q$, the two maps
$\widetilde B_z q$ and $q\widetilde B_z$ lift the same map
$B_zp$ and agree at $z_*$, so they are equal for every $z$.
On $O\times\partial D^2$, both $\widetilde B_z$ and the specified
$f_t$ lift (C.32.5) and agree at the based endpoint. They are
therefore equal on that entire connected boundary product.
Thus $\widetilde B$ fills the original loop in its centralizer.
This proves (C.32.1).

Now let $O$ be a connected open subset of the plane preserved by
a based ambient-null loop. Its restriction is a continuous loop
in $\operatorname{Homeo}(O)^0$. By
\cite[Theorem~1.1]{A_Yagasaki2000}, with the empty subpolyhedron,
this identity group is homotopically trivial unless $O$ is the
plane or the open annulus; the nonorientable and boundary
exceptions in the source cannot occur for $O\subseteq\R^2$.
In either remaining case its fundamental group is $\Z$.
Choose model coordinates on $O$, a coordinate Jordan circle
$J\subset O$, and two distinct points $a,b$ of $J$. A Schoenflies
homeomorphism $h$ of the ambient plane sends $J$ to the unit circle.
The map
\[
 \Theta:\operatorname{Homeo}(O)\longrightarrow S^1,\qquad
 k\longmapsto
       \frac{h(k(b))-h(k(a))}
            {|h(k(b))-h(k(a))|}
                                                        \tag{C.32.7}
\]
has induced homomorphism $\pi_1\operatorname{Homeo}(O)\to\Z$
that is an isomorphism. Indeed, the full rotation loop in the
chosen plane or annulus model preserves $J$ and has restriction
winding one. The explicit computation in C.29 shows that its
image under (C.32.7) also has winding one. A homomorphism
$\Z\to\Z$ whose image contains $1$ is an isomorphism.
Formula (C.32.7) extends as a continuous map on all ambient
plane homeomorphisms by the same two evaluations. The
ambient-null loop therefore has zero image under it, so its
restricted fundamental-group class is zero. This proves the
second assertion in all cases.

Finally take an open orbit from C.28 and let $Q$ be the closed
restricted operator image supplied by C.30. When this group is
discrete, C.30 gives a free proper action with precisely that
discrete topology, so (C.32.2) applies. The restricted canonical
plane loop commutes with $Q$ and is null-homotopic in
$\operatorname{Homeo}(O)$ by the second assertion. Its
null-homotopy supplied by (C.32.1) commutes with $Q$, hence with
the image of the operator stabilizer. The kernel in (C.30.4)
acts identically on $O$ throughout. This proves the stated
equivariant conclusion with the actual restriction homomorphism.
\end{proof}

\paragraph{Transport and frontier control.}
One may choose the obtained contraction on one representative of
an operator orbit of open $G$-orbits and transport it to the others
by conjugation. Equivariance under the setwise stabilizer makes
this independent of the transporting operator. This establishes
well-defined orbitwise families. The covering lift above supplies
no uniform estimate at orbit frontiers. C.33 subsequently treats
all positive-dimensional proper closed operator images
intrinsically, and C.34 proves joint continuity on the full
nonopen locus. C.35 records the obstruction to extending its
particular median prescription, while C.36 constructs compatible
alternatives on finite unions of closures with two-dimensional
proper operator image. The remaining frontier and accumulation
requirements are stated in C.12, which remains open.

\begin{claim}{C.33}{Proved}{All proper operator groups on an open orbit}
Let $O$ be a connected noncompact second-countable surface without
boundary. Let $Q$ be a closed locally compact subgroup of
$\operatorname{Homeo}(O)$, with its compact-open topology, acting
freely and properly. Then inclusion induces an injection
\[
 \pi_1\!\left(C_{\operatorname{Homeo}(O)}(Q),\id\right)
       \longrightarrow \pi_1(\operatorname{Homeo}(O),\id).
                                                        \tag{C.33.1}
\]
For every open orbit of the canonically lifted hyperbolic
centralizer loop in C.28, its restriction therefore has a based
null-homotopy commuting with the entire restricted operator
stabilizer. This conclusion retains the restriction kernel in
C.30 and imposes no finiteness or discreteness hypothesis on the
original operator group.
It supplies an intrinsic null-homotopy on each open orbit;
continuity across orbit frontiers remains unproved.
\end{claim}
\passage{Section 23, folios 140-141, scans 279-281; Part C(b) infinite-operator request; arbitrary proper-closure open-orbit reduction for C.12}

\begin{proof}
Write $N=C_{\operatorname{Homeo}(O)}(Q)$.
The topological group $Q$ acts faithfully on the connected
three-manifold $O\times\R$ by
\[
                         q(x,s)=(qx,s).
                                                        \tag{C.33.2}
\]
This action is continuous: on any compact subset of $O\times\R$
the first coordinate is controlled on its compact projection to
$O$. Pardon proves that a locally compact group acting faithfully
on a connected three-manifold is a Lie group
\cite[Theorem~1.5 and the reduction preceding it]{C_Pardon2013}.
Thus $Q$ is a Lie group. Since it is a subspace of the
second-countable homeomorphism group, it is second countable.
Set $A=Q^0$ and $d=\dim A$.
Properness and freeness make $q\mapsto qx$ a proper embedding of
$A$ in $O$ for each $x$. Invariance of dimension gives $d\leq2$.
The case $d=0$ is precisely C.32.

If $d=2$, invariance of domain makes every $A$-orbit open.
Connectedness gives a single orbit, and evaluation $A\to O$
is a homeomorphism. For $q\in Q$, choose $a\in A$ with $qx=ax$.
Then $a^{-1}q$ fixes $x$, so freeness gives $q=a$ and $Q=A$.
Under $O\cong A$ with its left translation action, the elements of
$N$ are exactly the right translations. The correspondence
is an anti-isomorphism of topological groups, and evaluation
$N\to O$ is a homeomorphism. A loop in $N$ null-homotopic
in $\operatorname{Homeo}(O)$ has a null evaluation loop in $O$,
and hence is null-homotopic in $N$. This proves (C.33.1) in
dimension two.

It remains to treat $d=1$. The connected one-dimensional Lie
group $A$ is $\R$ or $\R/\Z$, written additively.
Its restricted action is still free and proper.
Palais' slice theorem gives a locally trivial principal bundle
\[
                 p_0:O\longrightarrow B=O/A.
                                                        \tag{C.33.3}
\]
Here all hypotheses of the theorem hold: the surface is completely
regular, $A$ is Lie, and a relatively compact neighborhood has
relatively compact transporter by properness, which is the
Cartan condition. A slice at a point has trivial stabilizer.
The resulting local trivialization is the map
$(a,s)\mapsto as$ from $A\times S$ onto the open saturated
neighborhood of that slice. This is the free-action specialization
of \cite[Theorem~2.3.3]{C_Palais1961}; its maps and inverse are
also given in \cite[Theorem~3.7]{C_KramerMuratGarcia2026}.

We verify the dimension and local topology of the base rather
than assuming differentiability of this action. The elementary
fact needed is that a metrizable space $S$ for which
$S\times(-1,1)$ is a surface is a one-manifold without boundary.
Indeed, projecting path-connected open neighborhoods in the
product shows that $S$ is locally path connected. It has no
isolated point, since an isolated point would give an open
interval in a surface. Every embedded open arc in $S$ is open:
its product with $(-1,1)$ maps injectively into a surface, and
invariance of domain applies.

Fix $p\in S$. Local path connectedness and the absence of isolated
points give an embedded arc $\alpha:[0,1]\to S$ with
$\alpha(0)=p$. Choose a path-connected open neighborhood $W$
of $p$ avoiding $\alpha(1)$. If $W$ contains a point outside
$\alpha([0,1])$, join $p$ to that point by an embedded arc in
$W$. This second arc cannot enter and subsequently leave
$\alpha((0,1))$: that open arc has boundary contained in
$\{p,\alpha(1)\}$, and an embedded arc in $W$ cannot revisit
$p$ or pass through $\alpha(1)$. Thus the two arcs meet only
at $p$, and together contain an arc having $p$ as an interior
point. Such an open arc is an open neighborhood of $p$ by the
preceding paragraph.
In the other case $W\subseteq\alpha([0,1))$, and a smaller
neighborhood of $p$ is a half interval. Its product with
$(-1,1)$ has zero local second homology at $(p,0)$,
whereas a surface without boundary has local second homology
$\Z$. This is impossible. The elementary fact follows.
Applied to the slices in (C.33.3), it proves that $B$ is a
one-manifold without boundary. Properness gives a Hausdorff
quotient, the open quotient map gives second countability, and
connectedness is inherited from $O$. Consequently $B$ is
$\R$ or $\R/\Z$.

Every principal $A$-bundle over either of these one-manifolds
is trivial. One can see the actual section as follows.
Over an interval, local trivializations are joined successively
across overlapping subintervals, adjusting one section by its
continuous $A$-valued transition function. A locally finite
interval exhaustion gives a section on $\R$.
Over the circle, first cut at one point and choose a section
on the resulting closed interval. Its endpoint discrepancy is
an element of the connected group $A$; multiplying the section
by a path from $0$ to the negative of that element makes the
endpoints agree, so that it descends to the circle.
The map $(a,b)\mapsto a\,s(b)$, with inverse given in each
bundle chart, therefore identifies
\[
                      O\cong A\times B,\qquad
             c(a,b)=(c+a,b)\quad(c\in A).                \tag{C.33.4}
\]
The pair $A=B=\R/\Z$ is excluded because $O$ is noncompact.

Let $\Delta=Q/A$. This discrete group acts freely and properly
on $B$. Freeness follows directly: an element of $Q$ fixing a
point of $B$ sends any point of its fibre to an $A$-translate;
after that translate is removed, it fixes a point of $O$ and
is the identity.
For properness, suppose $b_j\to b$ and
$\delta_j b_j\to b'$ in $B$. Local sections give $x_j\to x$
above $b_j$ and $y_j\to y$ above $\delta_j b_j$.
Choose a representative $q_j$ of $\delta_j$ and $a_j\in A$
such that $a_jq_jx_j=y_j$. Properness of the $Q$-action gives
a convergent subsequence of $a_jq_j$, so their components
$\delta_j$ are eventually constant on that subsequence.
This proves properness of the discrete action on the locally
compact metrizable space $B$.

We spell out this discrete action. On $\R$, an orientation
reverser has a fixed point and is excluded by freeness.
If $\Delta$ is nontrivial, its orbit of a point is closed and
discrete, with a nearest point to the right. The element
taking the chosen point to that next point generates
$\Delta$: its iterates are unbounded in both directions, and
any intervening orbit point would contradict the choice of
the nearest one. Choosing coordinates on one fundamental
interval and extending by its iterates identifies the action
with
\[
                   \Delta=\Z,\qquad \beta(b)=b+1.
                                                        \tag{C.33.5}
\]
On $\R/\Z$, properness makes $\Delta$ finite and freeness
again excludes orientation reversers. The action on one
finite orbit preserves its cyclic order, so it identifies
$\Delta$ with a cyclic group $C_n$. Coordinates on the
successive complementary intervals identify a generator
with
\[
                   \beta(b)=b+1/n\quad\text{in }\R/\Z.
                                                        \tag{C.33.6}
\]
The trivial group is also allowed on either base.

Suppose $\Delta$ is nontrivial and choose a representative
$q\in Q$ of the indicated generator. Conjugation by $q$ on
$A$ is multiplication by $\lambda$, where
$\lambda\in\R^\times$ for $A=\R$, and
$\lambda\in\{1,-1\}$ for $A=\R/\Z$.
Equation (C.33.4) therefore expresses its action exactly as
\[
                 q(a,b)=(\lambda a+\tau(b),\beta(b))
                                                        \tag{C.33.7}
\]
for a continuous function $\tau:B\to A$.
We next give a coordinate change removing this translation
term, with all component relations retained.

For $B=\R$ and $\Delta=\Z$, choose a continuous function
$c:[0,1]\to A$ whose endpoints satisfy
$c(1)=\lambda c(0)+\tau(0)$, and extend it in both directions
by
\[
                      c(b+1)=\lambda c(b)+\tau(b).
                                                        \tag{C.33.8}
\]
The endpoint equation makes the extension continuous at every
integer; invertibility of multiplication by $\lambda$ gives
the extension towards negative integers. In the new fibre
coordinate $a'=a-c(b)$, (C.33.7) becomes
$q(a',b)=(\lambda a',b+1)$.

For $B=\R/\Z$ and $\Delta=C_n$, necessarily $A=\R$.
Write $q^n=T_{a_0}$, where $T_{a_0}$ is translation by $a_0$
on the fibre. Conjugation gives $\lambda^n=1$ and
$\lambda a_0=a_0$.
If $\lambda=1$, replacing $q$ by $T_{-a_0/n}q$ makes
$q^n=1$ and does not change the group generated with $A$.
If $\lambda=-1$, then $n$ is even and $a_0=0$, so $q^n=1$
already. Thus the negative value of $\lambda$ is retained.
For this finite-order representative define the real-valued
function
\[
 c(b)=\frac1n\sum_{k=0}^{n-1}
      \operatorname{pr}_A\!\left(q^k(0,\beta^{-k}b)\right).
                                                        \tag{C.33.9}
\]
The affine nature of (C.33.7), followed by cyclic
reindexing using $q^n=1$, gives
$c(\beta b)=\lambda c(b)+\tau(b)$.
The same coordinate change $a'=a-c(b)$ now gives
$q(a',b)=(\lambda a',\beta b)$.
We have proved, by explicit changes of coordinates, that
the nonidentity components can be represented as
\[
                         q(a,b)=(\lambda a,\beta b).
                                                       \tag{C.33.10}
\]
When $\Delta$ is trivial no such additional equation is needed.

Commuting with all translations in (C.33.4) forces a
homeomorphism of $A\times B$ to have the form
\[
       h(a,b)=(a+\phi(b),\psi(b)),\qquad
       \phi\in C(B,A),\quad\psi\in\operatorname{Homeo}(B).
                                                       \tag{C.33.11}
\]
The inverse is
$(a,b)\mapsto(a-\phi(\psi^{-1}b),\psi^{-1}b)$.
These formulas and evaluation at the zero section give
mutually inverse continuous maps for the compact-open
topologies. In the nontrivial-component case, commuting
with (C.33.10) is exactly the pair of equations
\[
             \psi\beta=\beta\psi,\qquad
             \phi(\beta b)=\lambda\phi(b).
                                                       \tag{C.33.12}
\]
Thus the map (C.33.11) identifies $N$, as a space, with
the product of the function space defined by the second
equation and the homeomorphism centralizer defined by the
first. In particular the two constraints can be contracted
separately without changing the group $Q$ being centralized.

Let $h_t$ be a based loop in $N$ null-homotopic in
$\operatorname{Homeo}(O)$, and let $(\phi_t,\psi_t)$ be its
coordinates. The path $\psi_t$ preserves the orientation of
the one-manifold $B$.
When $A=\R$, multiplication of $\phi_t$ by $1-s$ contracts
the function-coordinate loop to zero and preserves its
equation in (C.33.12), including $\lambda=-1$.
For $B=\R$, C.27 contracts $\psi_t$ in the required
centralizer, with the trivial operator group when
$\Delta=1$. For $B=\R/\Z$, the evaluation loop
$h_t(0,b)=(\phi_t(b),\psi_t(b))$ is null-homotopic in
$\R\times(\R/\Z)$; hence $\psi_t(b)$ has winding zero.
C.29 contracts $\psi_t$ in the centralizer of the finite
rotation group in (C.33.6).

When $A=\R/\Z$, necessarily $B=\R$.
The based family $\phi_t(b)$ has a unique continuous lift
$\varphi:[0,1]\times\R\to\R$ with $\varphi_0=0$.
Evaluation of $h_t$ in the annulus $A\times\R$ shows that
$\varphi_1(b)=0$: this endpoint is the winding number of
the first coordinate of the null evaluation loop.
If $\Delta=\Z$, equation (C.33.12) says that
\[
           \varphi_t(b+1)-\lambda\varphi_t(b)\in\Z.
                                                       \tag{C.33.13}
\]
Its left side is a continuous function of $(t,b)$, equal
to zero at $t=0$, and is therefore identically zero.
Multiplying $\varphi$ by $1-s$ and descending to $\R/\Z$
contracts $\phi_t$ while preserving exact equivariance.
The same scaling applies when $\Delta=1$, without (C.33.13).
C.27 contracts the remaining $\psi_t$ on the line.
All these contractions keep $t=0,1$ fixed. Equations
(C.33.11)--(C.33.12), including the inverse formula, show
that their product is a continuous based contraction in
$N$. This completes the proof of (C.33.1).

Finally, for an open loop-group orbit from C.28, C.30
supplies the actual closed locally compact free proper
operator image $Q$. The restricted canonical plane loop
commutes with $Q$. It is null-homotopic in
$\operatorname{Homeo}(O)$ by the second assertion of
C.32, so (C.33.1) contracts it in the exact centralizer.
Every original operator in the setwise stabilizer acts
through its restriction to $Q$, and the restriction
kernel acts identically on $O$. The stated conclusion
follows with that kernel retained.
\end{proof}

\paragraph{What the open-orbit result provides.}
Equations (C.33.11)--(C.33.13) treat every component of the
proper operator group, including a circle fibre with a
nontrivial $\Z$-action on its line base. There is no
replacement of the base-translation centralizer by the
homeomorphism group of its quotient circle.
As in C.32, the obtained contraction can be transported to
other open orbits in the same operator orbit; exact
stabilizer equivariance makes the transport independent
of choices. This proof supplies no uniform control near
the boundary of the open orbit. C.34 subsequently proves
joint continuity, including inverses, on the entire closed
union of nonopen orbits. C.35 exhibits an exact obstruction
to extending that particular median prescription over an
adjacent open orbit. C.36 constructs compatible contractions
on closures with two-dimensional proper operator image and
proves continuity on every finite union of those closures.
C.37 supplies the uniform three-word estimate and gives a
continuous contraction on the closed union $X_2$ of all those
open orbits, including infinite accumulation. Extension across
$X_2$ to its exterior and frontier-compatible choices for the
dimension-zero and dimension-one operator images remain
additional requirements. C.12 retains its unresolved status.

\begin{claim}{C.34}{Proved}{The median contraction on all nonopen orbits}
Let $f_t$ be a based null-homotopic loop in
$\operatorname{Homeo}^+(\R^2)$, and use $P,G,Y$ from
C.28 and C.31. The orbitwise bijections $Q_{s,t}:Y\to Y$
defined in C.31 are homeomorphisms. Both maps
\[
 [0,1]^2\times Y\longrightarrow Y,\qquad
                  (s,t,x)\longmapsto Q_{s,t}^{\pm1}(x)
                                                        \tag{C.34.1}
\]
are continuous. Thus these formulas give a based
compact-open null-homotopy on the entire closed nonopen
locus, commuting with every homeomorphism commuting
with all $f_t$.
In particular, a hyperbolic centralizer loop whose
canonically lifted generated group has no open orbit
is null-homotopic in the original surface centralizer.
No extension over the omitted open orbits is asserted.
\end{claim}
\passage{Section 23, folios 140-141, scans 279-281; Part C(b) infinite-operator request; nonopen-stratum assembly in C.12}

\begin{proof}
We prove continuity at a nonfixed orbit. Continuity at
the common fixed set, bijectivity on $Y$, equivariance,
and the exact inverse identity have already been proved
in C.31.

\emph{Short paths moving in either direction.}
Let $q$ be a point of a nonopen nontrivial orbit. It lies
on an intrinsic line or a circle by C.28. There are
paths $g_s$ from the identity, formed from two factors
in $P$, which move $q$ a nonzero amount in either
chosen intrinsic direction while keeping its whole
track inside any prescribed sufficiently small
intrinsic arc about $q$.
To see this, $f_t(q)$ is not identically $q$; otherwise
every generator, and hence every element of $G$, would
fix $q$. Choose a component $(a,b)$ of
$\{t:f_t(q)\ne q\}$ and a sufficiently short initial
subinterval $[a,a+\eta]$. The path
\[
          h_s=f_{a+s\eta}f_a^{-1},\qquad 0\leq s\leq1,
                                                        \tag{C.34.2}
\]
starts at the identity, and $h_s(q)=f_{a+s\eta}(q)$.
C.28 gives compact-open continuity on an intrinsic
line; the circle restriction is continuous as well.
Consequently $h_s$ and $h_s^{-1}$ can be made uniformly
close to the identity on an intrinsic neighborhood of
$q$. Their endpoint displacements at $q$ have opposite
intrinsic signs. One of these two paths has the chosen
sign, and its track stays in the prescribed arc.
The paths also become uniformly small on each fixed
ambient compact set as $\eta\to0$.
Every such $g_s$ is a product of two elements of $P$.
For any $z\in K_n(x)=P^nx$,
\[
                        g_s(z)\in K_{n+2}(x).
                                                        \tag{C.34.3}
\]
In particular, the same compact-word argument as in
C.28 proves intrinsic continuity of its track on
every nonopen orbit, including when $z$ varies along
a compact intrinsic subarc.

\emph{A small cap cannot persist at a moving limit point.}
We give the local argument used below. Straighten a
compact arc of a limit orbit, extending it a little
at both ends, so that it is an interval of the
horizontal axis in a plane chart. Suppose compact
intrinsic arcs $A_j$ in nearby nonopen orbits lie in
strips of height tending to zero around this interval.
Let $q_j\in A_j$ attain the minimum $m_j$ of the
horizontal coordinate on $A_j$, suppose
$q_j\to q=(m,0)$, and suppose both endpoints of
$A_j$ have horizontal coordinates at least $m+\delta$
for a fixed $\delta>0$.
Take a rectangle
\[
          U=(m-\rho,m+\rho)\times(-\varepsilon,\varepsilon),
          \qquad 0<\rho<\delta/3,
                                                        \tag{C.34.4}
\]
whose closure stays in the chart. For large $j$,
$q_j$ lies in $U$, all of $A_j$ has vertical coordinate
of absolute value less than $\varepsilon/2$, and its
horizontal minimum is greater than $m-\rho/2$.
Starting at $q_j$ in each intrinsic direction towards
the two endpoints of $A_j$, take the first point
outside $U$. Both first-exit points lie on its right
edge. They are distinct points of the embedded arc.
Let $J_j$ be the closed subarc between these points.
Its interior lies in $U$, and every point of $J_j$
has horizontal coordinate at least $m_j$.

Use (C.34.2), taking the inverse when needed, to choose
a path $g_s$ whose $q$-track stays in $U$ and whose
endpoint is strictly to the left of $q$. This choice
is possible because $q$ is nonfixed. The track is a
compact subset of $U$, so joint continuity gives
$g_s(q_j)\in U$ for every $s$ and every sufficiently
large $j$.
Its intrinsic continuity from (C.34.3) now matters:
a continuous path in the intrinsic line or circle
starting at an interior point of $J_j$ cannot leave
$J_j$ without meeting one of its two endpoints.
Both endpoints lie outside the open rectangle $U$.
Thus the entire $g_s(q_j)$-track stays in $J_j$.
At $s=1$ its horizontal coordinate is at least $m_j$,
whereas
\[
        g_1(q_j)\longrightarrow g_1(q),\qquad
        \operatorname{pr}_1(g_1(q))<m=\lim_jm_j.
                                                        \tag{C.34.5}
\]
This contradiction excludes the cap.
The argument with horizontal maxima and a path moving
to the right is identical. The two first-exit points
are the explicit obstruction to escaping the cap
along another part of the same orbit.

The same argument excludes a whole Jordan-curve orbit
contained in such a shrinking strip and having a
nondegenerate interval as its limit: use a global
minimum of the horizontal coordinate on the circle,
and in its two directions take first exits towards
a point whose coordinate is bounded above that
minimum. The resulting subarc $J_j$ again has
two distinct endpoints outside $U$.

\emph{No excess in subarcs near an intrinsic line.}
Fix a point $p$ on an intrinsic line orbit and
take $x_j\to p$ in $Y$. For each fixed $n$, compactness
of $P^n$ and continuity of word evaluation give
\[
                         K_n(x_j)\longrightarrow K_n(p)
              \quad\hbox{in Hausdorff distance}.         \tag{C.34.6}
\]
The limit is a nondegenerate compact arc. Straighten
a slightly larger compact arc of the limit orbit.
Then (C.34.6) places the nearby $K_n(x_j)$ in shrinking
horizontal strips in that chart. Each $K_n(x_j)$ is
an arc, a point, or a whole circle, by C.28 and
connectedness. The whole-circle possibility is
excluded for all sufficiently large $j$ by the
preceding cap argument.

Suppose $u_j,v_j\in K_n(x_j)$ tend to $u,v\in K_n(p)$.
Write $[u_j,v_j]_j$ for the subarc contained in
$K_n(x_j)$. Then
\[
 \begin{gathered}
 z_j\in[u_j,v_j]_j,\quad z_j\longrightarrow z
             \quad\Longrightarrow\quad z\in[u,v].
                                                        \tag{C.34.7}
 \end{gathered}
\]
Indeed, an alleged limit point strictly to the left
of both $u,v$ produces a horizontal minimum of
$[u_j,v_j]_j$ whose two endpoints stay a fixed
distance to its right. A subsequence of these
minimum points converges, by (C.34.6), to a point
of the nonfixed limit orbit. This is the forbidden
cap. A point to the right is excluded by the
maximum version. The assertion includes $u=v$:
then all of the intervening subarc tends to $u$.
This proves the needed continuity of subarcs using
the path action; Hausdorff convergence alone has
not been used to infer it.

It follows that medians of three points in $K_n(x_j)$
converge to the median of their limits. Each median
lies in all three pairwise connecting subarcs.
Every subsequential limit therefore lies, by
(C.34.7), in all three corresponding limit intervals.
Their intersection is the single limit median,
including when two or three of the limiting points
coincide.

\emph{A circle limit: the precise local covering.}
Let the orbit of $p$ be a Jordan curve $C$. Choose a
collar annulus $W$ of $C$, with coordinates
$(\R/\Z)\times(-1,1)$ sending $C$ to the middle circle,
and write
\[
 \pi_W:\widetilde W=\R\times(-1,1)\longrightarrow W
                                                        \tag{C.34.8}
\]
for its universal covering. The compact family
$f_t$ preserves $C$. Therefore a sufficiently small
open disc $V$ about $p$ satisfies
$f_t(V)\subset W$ for all $t$.
The actual map
\[
 [0,1]\times V\longrightarrow W,\qquad
                         (t,x)\longmapsto f_t(x)
                                                        \tag{C.34.9}
\]
has a unique continuous lift $\widehat f_t(x)$ after
choosing a lift of the initial inclusion $V\to W$.
Its endpoint differs from its starting point by
an integer deck translation depending continuously
on $x\in V$. It is zero at $p$, by C.29's exact
plane-to-circle winding comparison, so it is zero
on all of the connected disc $V$.
Thus these are closed lifted point tracks.
For $x_j\to p$ in $V\cap Y$ put
\[
 T_j(r)=\{\widehat f_u(x_j):0\leq u\leq r\},\qquad
 T(r)=\{\widehat f_u(p):0\leq u\leq r\}.
                                                        \tag{C.34.10}
\]
Joint continuity on the compact parameter interval
shows $T_j(r_j)\to T(r)$ in Hausdorff distance
whenever $r_j\to r$. The set $T(1)$ is a
nondegenerate compact interval on the line
$\R\times\{0\}$.

We verify the intrinsic meaning of the lifted
nearby tracks. An intrinsic line orbit, or a
proper component of a circle orbit inside $W$,
has an intrinsic interval as its component through
$x_j$. Its continuous point track stays in that
component and lifts to an injectively parametrized
interval in $\widetilde W$, with embedded compact
subintervals. If the entire circle orbit lies in
$W$, it is a Jordan curve. Its winding in $W$
is $0$ or $\pm1$: the planar Jordan winding
around the complementary inner disc has exactly
these values. Its lifted components are therefore
Jordan curves in the first case and intrinsic
lines in the second.
Consequently $T_j(1)$ is a compact arc, a point,
or a whole lifted Jordan curve, with the
corresponding intrinsic topology.

The cap argument applies in this covering chart
with the original group paths. Here is the exact
local map needed to justify that use. At any
limiting extremum $q\in T(1)$ choose the small
rectangle $U$ of (C.34.4) with horizontal width
less than one, so that $\pi_W|_U$ is injective.
Choose a path (C.34.2) at $\pi_W(q)$ with the
required intrinsic sign and sufficiently small
that its images of a compact neighborhood of
$\pi_W(q)$ stay in $W$.
The map
\[
 (s,z)\longmapsto g_s(\pi_W(z))
                                                        \tag{C.34.11}
\]
on that small lifted neighborhood has a unique
continuous lift starting with $z\mapsto z$.
Its $q$-track lies in $U$ and has the required
horizontal displacement. For $q_j\to q$ its
whole $q_j$-track also stays in $U$, by compact
parameter continuity. Downstairs it is an
intrinsically continuous path in the original
orbit, by (C.34.3). Its lift cannot leave the
cap subarc without meeting one of its first-exit
endpoints on $\partial U$. Thus (C.34.5) applies
unchanged. No action of $G$ on the entire
annulus or its covering has been assumed.

This excludes a whole lifted Jordan curve for
large $j$ and proves (C.34.7), and hence median
continuity, for compact subarcs of $T_j(1)$.
These intrinsic interval operations are exactly
the ones used by C.27 and C.29. For a line or
a proper circle component the lift of its
interval is order preserving or reversing.
For an essential circle the annular covering
is its universal line covering. For a
contractible circle, the now proper arc
$T_j(1)$ lies in a single interval of its
universal line covering; a path leaving that
interval would have to cover the whole circle.
In all cases the based circle lift or intrinsic
line formula, followed by projection through
(C.34.8), gives the original $Q_{s,t}(x_j)$.
Median invariance under reversal removes any
choice of the interval orientation.

\emph{Continuity of the prefix-clamping formula.}
We can now handle a sequence
$(s_j,t_j,x_j)\to(s,t,p)$ with $p$ nonfixed.
For a line limit use the actual point tracks
$T_j(r)=\{f_u(x_j):0\leq u\leq r\}$ inside
$K_1(x_j)$. For a circle limit use (C.34.10).
In either case $T_j(1)$ is eventually an
intrinsic compact arc, its subarcs satisfy
(C.34.7), and its prefix tracks have their
stated Hausdorff limits.
Set $r_j=1-s_j$ and $r=1-s$.
Let $a_j,b_j$ be the two unordered endpoints
of the prefix arc $T_j(r_j)$, allowing equality.
Every subsequence has a further subsequence
on which both endpoints converge, say to
$a,b\in T(r)$. By (C.34.7), the Hausdorff
limit of their intervening arc is contained
in $[a,b]$. It is also exactly $T(r)$ by
the prefix-track continuity. Since $a,b$
belong to $T(r)$, their interval is contained
in $T(r)$. Therefore
\[
                         T(r)=[a,b].                    \tag{C.34.12}
\]
In particular the limiting unordered pair is
precisely the endpoint pair of the prefix
interval, including a degenerate prefix.

The formula of C.27 or its C.29 lift is
the median of $a_j,b_j$ and the current
track point at $t_j$. The median continuity
just proved and (C.34.12) give its limit as
the median of the limiting prefix endpoints
and the current point at $t$. This is
$Q_{s,t}(p)$, with projection through
(C.34.8) in the circle case.
Every subsequence has this same limit.
The domains are metrizable, so this proves
joint continuity on $[0,1]^2$ times the
nonfixed part of $Y$. C.31 supplies the
remaining fixed points.

Apply the same argument to the based inverse
loop $u\mapsto f_u^{-1}$. It is ambient-null,
generates the same group $G$, and has the
same closed locus $Y$. The exact identity
$Q[f]^{-1}=Q[f^{-1}]$ in C.31 proves joint
continuity of the inverse family in
(C.34.1). Thus all the orbitwise bijections
are homeomorphisms of $Y$. The compact-open
continuity follows from the evaluation
exponential law, since $Y$ is a closed,
locally compact subspace of the plane.
The based endpoints and equivariance are
those already proved orbitwise in C.31.

Finally, when the canonical lifted group of
a hyperbolic surface loop has no open orbit,
$Y=\widetilde U\cong\R^2$.
The obtained plane homeomorphisms commute
with the full operator lift group, including
all deck transformations. They descend by
the exact equations
\[
                  \overline Q_{s,t}\,p_U=p_UQ_{s,t}
                 \quad(p_U:\widetilde U\to U)
                                                        \tag{C.34.13}
\]
to homeomorphisms of $U$ commuting with every
original operator. Their joint continuity
and that of their inverses descend through
the open covering quotient, giving
compact-open continuity downstairs.
The endpoints in (C.34.13) are the original
loop and the identity, with both based
edges fixed. Hence this is its null-homotopy
inside the original centralizer.
At every parameter it lies in the identity
component by these same paths.
Equivalently, its lift is the canonical
one: two deck-commuting lifts of such a
map differ by a central deck transformation,
and the hyperbolic surface group has trivial
center. This proves the stated corollary
with the actual lift and descent maps.
\end{proof}

\paragraph{Remaining global extension.}
C.34 establishes the full nonopen-locus
homeomorphism family and its inverse.
C.33 provides intrinsic equivariant
contractions on each open orbit.
C.35 proves that this particular median
family need not extend equivariantly over
an adjacent open orbit. C.36 and C.37
construct compatible alternatives on the
closed union $X_2$ of open orbits with
two-dimensional proper operator image,
including every infinite accumulation.
Matching that boundary family to the
exterior of $X_2$, and treating the other
proper operator images at their frontiers,
remain the assembly requirements in
C.12, whose general status remains open.

\begin{claim}{C.35}{Proved}{An obstruction to extending the fixed median prescription}
There is a continuous $\R^2$-action on the plane and a based
ambient-null loop in its centralizer such that the C.34 median
homeomorphisms of the closed nonopen locus do not extend
equivariantly over an adjacent open orbit at one parameter.
The closed restricted operator image on that open orbit is
the free proper translation group $\R^2$.
The original loop nevertheless has an explicit equivariant
null-homotopy on the entire plane.
Thus this example refutes extension of the fixed C.34 median
prescription; it is not a failure of the level-$0$ comparison.
\end{claim}
\passage{Section 23, folios 140-141, scans 279-281; Part C(b) infinite-operator request; obstruction to extending a fixed contraction prescription in the C.12 proof route}

\begin{proof}
Use the closed simplex
\[
 \Delta=\{(x,y,z)\in\R_{\geq0}^3:x+y+z=1\}.
\]
For $(a,b)\in\R^2$ define
\[
 T_{a,b}(x,y,z)
   =\frac{(e^a x,e^b y,z)}{e^a x+e^b y+z}.
                                                        \tag{C.35.1}
\]
The denominator is positive on $\Delta$.
Direct substitution gives
$T_{a,b}T_{c,d}=T_{a+c,b+d}$ and
$T_{a,b}^{-1}=T_{-a,-b}$, so (C.35.1) is a continuous
action by homeomorphisms, fixing the three vertices and
preserving each open edge.
On the interior the coordinate homeomorphism
\[
 (x,y,z)\longmapsto(u,v)=\bigl(\log(x/z),\log(y/z)\bigr),
 \qquad
 (u,v)\longmapsto
       \frac{(e^u,e^v,1)}{1+e^u+e^v}
                                                        \tag{C.35.2}
\]
turns this action into $(u,v)\mapsto(u+a,v+b)$.
Its restricted image is therefore the closed free proper
translation group on $\R^2$.

Double $\Delta$ along its entire boundary. Since $\Delta$
is a closed disc, this double is a sphere. Apply (C.35.1)
on both copies; the formulas agree on their identified
boundary and hence give a continuous sphere action.
Remove the fixed vertex $(0,0,1)$. The remaining surface
$M$ is a plane, and the action restricts to it.
It is orientation preserving: each map preserves the
orientation on either triangle interior, and preserves
the two sheets of the oriented double.

Write $Q=\{T_{a,b}:(a,b)\in\R^2\}$ for this operator
image. Its orbits in $M$ are the two triangle interiors,
the three open edges, and the remaining two vertices.
Use the following coordinates on the edges:
\[
\begin{array}{c|c|c}
 \text{edge}&\text{coordinate in }\R&
                      \text{translation by }T_{a,b}\\ \hline
 y=0&\log(x/z)&a\\
 x=0&\log(y/z)&b\\
 z=0&\log(x/y)&a-b .
\end{array}                                               \tag{C.35.3}
\]
Thus the union $Y$ of the three edges and the two
vertices is exactly the closed nonopen orbit locus.

Let $w(t)=(a(t),b(t))$ be the piecewise linear based
loop with successive values
\[
       (0,0),\quad(1,1),\quad(1,0),\quad(0,0)
           \quad\text{at }t=0,\tfrac13,\tfrac23,1,
                                                        \tag{C.35.4}
\]
and put $f_t=T_{a(t),b(t)}$.
All its values commute with $Q$, and
\[
                  F_{\sigma,t}
                  =T_{(1-\sigma)a(t),(1-\sigma)b(t)}
                                                        \tag{C.35.5}
\]
is a based equivariant null-homotopy on the entire plane.
Continuity, including that of the inverse family, follows
from (C.35.1), which is already continuous on the compact
double and preserves the removed point.
The subgroup generated by the loop is exactly $Q$:
its first segment contains all diagonal translations
$(r,r)$ for $0\leq r\leq1$, and the differences of
points of the second segment contain all $(0,r)$
for $-1\leq r\leq1$. Integer sums and inverses give
both full coordinate axes and hence $\R^2$.
Consequently C.34 uses precisely the orbit locus $Y$
just described.

At $s=2/3$ the prefix in the C.27 formula ends at
$r=1-s=1/3$. By (C.35.3)--(C.35.4), the three
prefix translation intervals are, in their displayed
order,
\[
                         [0,1],\quad[0,1],\quad\{0\}.
                                                        \tag{C.35.6}
\]
At $t=2/3$, the current edge translations are
$(1,0,1)$. Taking the prescribed medians with
the prefix endpoints gives
\[
       \left.Q^{\mathrm{med}}_{2/3,\,2/3}\right|_Y:
           \quad\text{edge translations }(1,0,0).
                                                        \tag{C.35.7}
\]
The vertices stay fixed. These edge maps fit
continuously at the vertices, as also guaranteed
by C.34.

Every $Q$-equivariant homeomorphism from one triangle
interior to itself is a translation in (C.35.2).
Indeed, commutation with all translations gives
$h(u,v)=(u,v)+h(0,0)$. If it maps one interior to
the other, the same equation holds in the identical
coordinates on the target sheet. The sheet exchange
fixes the common boundary, so in either case any
continuous boundary extension has edge translation
parameters
\[
  R:\R^2\longrightarrow\R^3,\qquad
                    R(A,B)=(A,B,A-B).                  \tag{C.35.8}
\]
This is the exact restriction morphism.
Its image is the kernel of
\[
       \ell:\R^3\longrightarrow\R,\qquad
                         \ell(r_1,r_2,r_3)=r_1-r_2-r_3.
                                                        \tag{C.35.9}
\]
But $\ell(1,0,0)=1$, so (C.35.7) does not lie in
that image. It cannot extend equivariantly even
over a single adjacent interior. This proves
all assertions, while (C.35.5) gives the explicit
global contraction of the original loop.
\end{proof}

\paragraph{Executed endpoint certificate and scope.}
The standard-library rational calculation
\texttt{code/part\_b/c35\_median\_extension\_certificate.py}
computes the prefix extrema, the three median
values, the restriction matrix in (C.35.8),
and the nonzero obstruction (C.35.9).
Its JSON output records $(1,0,0)$ and residual $1$.
This certificate verifies the finite endpoint
calculation; the preceding simplex, doubling,
and restriction arguments prove the topological
statement.

C.34 is a homeomorphism construction on the whole
nonopen locus, and C.33 contracts the open-orbit
restrictions. C.35 proves that a global argument
must permit changing the nonopen-locus contraction
to satisfy the actual boundary restriction
morphisms of the open components. The remaining
problem is therefore the existence of compatible
equivariant contractions, not extension of the
particular median family without modification.
C.36 and C.37 meet the actual boundary equations on
the closed union $X_2$ of open orbits with
two-dimensional proper operator image, including
infinite accumulation. Extending this family
compatibly to the exterior of $X_2$ and the other
proper operator cases remains required.
The general C.12 comparison remains unresolved.

\begin{claim}{C.36}{Proved}{Compatible contractions on two-dimensional proper closures}
Let $f_t$ be a based ambient-null plane loop, let $G$ be
its generated group, and let $H$ be any group of plane
homeomorphisms commuting with every $f_t$.
For an open $G$-orbit $O$, let $Q_O$ be the closure of
the restricted setwise $H$-stabilizer, as in C.30.
For every such $O$ with $\dim Q_O=2$, there is a based
null-homotopy on the entire closure $\overline O$
through homeomorphisms of $\overline O$ that commute
with its operator stabilizer.
These closure contractions agree on every intersection
$\overline O_1\cap\overline O_2$ and are carried to one
another by $H$. They therefore give a jointly continuous
family, with jointly continuous inverse, on every
finite union of these closures.
No continuity across an infinite accumulating family
of closures is asserted.
\end{claim}
\passage{Section 23, folios 140-141, scans 279-281; Part C(b) infinite-operator request; frontier-compatible dimension-two operator closures in the C.12 proof route}

\begin{proof}
\emph{The actual quotient action on the closure.}
C.30 makes $Q_O$ a closed locally compact free proper
group of homeomorphisms of $O$, and C.33 makes it Lie.
Its dimension-two argument shows that $Q_O$ is
connected and simply transitive on $O$, and that
\[
 N_O=C_{\operatorname{Homeo}(O)}(Q_O)
                    \cong Q_O^{\mathrm{op}}             \tag{C.36.1}
\]
is also a connected, simply transitive two-dimensional
Lie group. Evaluation $N_O\to O$ is a homeomorphism.
The image of $G$ in $N_O$ is all of $N_O$: for any
$n\in N_O$, transitivity of $G$ supplies $g\in G$
with $gx=nx$ at one point $x\in O$, and the
injectivity of evaluation on $N_O$ then gives
$g|_O=n$.

Put $L=\overline G$ in $\operatorname{Homeo}^+(\R^2)$.
This is a Polish group. It preserves every open
$G$-orbit $O$. Indeed, if $g_j\to l$ with $g_j\in G$,
the identities $g_j(\partial O)=\partial O$ and
the convergent inverse family give
$l(\partial O)=\partial O$. For $x\in O$,
$lx\in\overline O$ by convergence, and
$lx\notin\partial O$ because $l$ preserves that
boundary and is injective. Thus $lx\in O$;
applying the same argument to $l^{-1}$ proves
$l(O)=O$.
The elements of $L$ still commute with $H$.
Restriction consequently gives a continuous
surjective homomorphism
\[
                   r_O:L\longrightarrow N_O.            \tag{C.36.2}
\]
It is open by the open mapping theorem for Polish
groups. Equivalently, apply the open orbit-map
principle from
\cite[Theorem~3.5]{C_GomezPinedoUzcategui2017}
to the transitive action of $L$ on the Baire
space $N_O$ through (C.36.2).

The kernel of (C.36.2) fixes $O$ pointwise, and
hence fixes $\overline O$ pointwise. Therefore
there is a well-defined action
\[
 N_O\times\overline O\longrightarrow\overline O,
       \qquad(n,z)\longmapsto lz
                 \quad\text{for any }l\text{ with }r_O(l)=n.
                                                        \tag{C.36.3}
\]
This action is jointly continuous: $r_O\times\id$
is an open quotient map, and the continuous
$L$-action factors through it. No continuous
section of (C.36.2) is required.
In fact each $n$ has a preimage in $G$, as proved
above. Hence (C.36.3) preserves every original
$G$-orbit contained in $\overline O$, and its
image on such an orbit is exactly the image of
$G$. This conclusion does not require a claim
that $L$ preserves every other $G$-orbit in the
plane.

\emph{The two-dimensional Lie contraction.}
We record the Lie groups and the maps being used.
A two-dimensional real Lie algebra is either
abelian or has a basis $X,Y$ with $[X,Y]=Y$:
in the nonabelian case its one-dimensional
derived algebra is an ideal, and scaling a
complementary basis vector gives that bracket.
The simply connected groups are respectively
$\R^2$ and the positive affine group of the
line. In coordinates $(a,b)$, the latter has
product
\[
           (a,b)(a',b')=(a+a',\,b+e^a b').
                                                        \tag{C.36.4}
\]
Its exponential map is
\[
 \exp(u,v)=\left(u,v\,\frac{e^u-1}{u}\right),
 \quad\text{with }\frac{e^u-1}{u}=1\text{ at }u=0.
                                                        \tag{C.36.5}
\]
The scalar factor is positive for every real
$u$, so this is a global diffeomorphism.
For $\R^2$ the exponential is the identity.
The affine group has trivial center, so its
connected groups with this Lie algebra have
no further covering quotients. In the abelian
case the connected group is a quotient of
$\R^2$ by a discrete central subgroup.
The compact torus is excluded here because
$N_O\cong O$ is noncompact.
Thus in every case the universal covering
group $\widetilde N_O$ has a globally defined
inverse $\log$ to its exponential.

Let $n_O(t)=f_t|_O\in N_O$.
By the restriction assertion of C.32 this loop
is null-homotopic in $\operatorname{Homeo}(O)$.
Evaluation on $O$ and the homeomorphism
$N_O\to O$ show that it is null-homotopic
in $N_O$. Its unique based lift
$\widetilde n_O(t)$ to $\widetilde N_O$
is therefore a closed loop.
With $p_O:\widetilde N_O\to N_O$ the covering
homomorphism, define
\[
 R^O_{s,t}
   =p_O\!\left(\exp\bigl((1-s)\log\widetilde n_O(t)\bigr)\right)
       \quad\text{acting on }\overline O
       \text{ by (C.36.3)}.                             \tag{C.36.6}
\]
This is a continuous family of homeomorphisms
with continuous inverse, starts with $f_t$ on
$\overline O$, ends with the identity, and
keeps both based edges fixed.
On $O$ it lies in the exact centralizer of
$Q_O$. Its commutation with the original
setwise operator stabilizer extends to
$\overline O$ by continuity and density.

\emph{The common boundary image and its topology.}
Take two of these closures and a $G$-orbit
$B$ in their intersection. It is a point,
an intrinsic line, or a circle: an open orbit
cannot lie in the boundary of another open
orbit. The point case is fixed throughout.
For a nontrivial $B$, each $N_i=N_{O_i}$
acts transitively on $B$ by (C.36.3), and
both action images are exactly $G|_B$.
We verify that this common image has a single
Lie topology, rather than identify abstract
quotients without their topology.

For $p\in B$, its stabilizer in $N_i$ is
closed. The homogeneous space
$N_i/(N_i)_p$ maps continuously and
bijectively to $B$. Its dimension is one:
dimension zero would make the connected
orbit a point, and dimension two would make
$B$ open in the plane by invariance of
domain. Its topology agrees with the
intrinsic topology of $B$.
Indeed, a compact interval in this
homogeneous one-manifold maps to an
embedded compact arc. In an intrinsic
line orbit that arc is the unique
connecting arc supplied by C.28; on a
circle the same conclusion holds locally.
The interiors of these arcs give
intrinsic interval neighborhoods, so the
bijective orbit map is a local, and
therefore global, intrinsic homeomorphism.
Thus the action is continuous also in
the intrinsic topology.

Divide $N_i$ by its closed kernel on $B$.
The resulting connected Lie group $M_i$
acts faithfully and transitively on that
one-manifold, and has dimension at most
two. In dimension one it is abelian, so
a point stabilizer is normal and is
therefore the kernel; faithfulness makes
the stabilizer trivial. Evaluation then
identifies its action with the translation
group of $\R$ or the rotation group of
$\R/\Z$.
An abelian group of dimension two would
give the same freeness conclusion and
an orbit of dimension two, so cannot
occur faithfully here.
In the remaining case $M_i$ is the
positive affine group (C.36.4).
The identity component of a point
stabilizer is a one-dimensional subgroup.
It cannot be the translation subgroup,
which is normal and would lie in the
action kernel. Every other such subgroup
is conjugate to the dilation subgroup.
Its normalizer is that dilation subgroup
itself, as is seen by conjugating its
unique fixed point on the line.
The full point stabilizer is consequently
the dilation subgroup, and the homogeneous
action is the usual positive affine
action on $\R$.

These descriptions prove topological
embedding of each actual image in
$\operatorname{Homeo}(B)$ with its
intrinsic compact-open topology.
For translations or rotations one
evaluation recovers the parameter.
For the affine group, the images of
two distinct ordered points recover
the positive slope and the intercept
continuously. They also prove that
the image groups are closed in that
homeomorphism group.
Since both images are the same set
$G|_B$, we have obtained a common
topological Lie group $M_B$, and
continuous quotient homomorphisms
\[
                       \rho_i:N_i\longrightarrow M_B.
                                                        \tag{C.36.7}
\]
The universal cover of $M_B$ is
$\R$ or the positive affine group,
and its exponential is again a
global diffeomorphism.

\emph{Equality on overlaps.}
The homomorphisms (C.36.7) lift
uniquely to based homomorphisms
\[
               \widetilde\rho_i:
                    \widetilde N_i\longrightarrow
                    \widetilde M_B.
                                                        \tag{C.36.8}
\]
The two paths
$\widetilde\rho_i(\widetilde n_i(t))$
are based lifts of the identical
loop $f_t|_B$, so covering uniqueness
makes them equal to one path
$\widetilde b(t)$.
Naturality of the Lie exponential
and its global injectivity give
\[
 \begin{split}
 d\widetilde\rho_i(\log\widetilde n_i(t))
      &=\log\widetilde b(t),\\
 \widetilde\rho_i
   \!\left(\exp\bigl((1-s)\log\widetilde n_i(t)\bigr)\right)
      &=\exp\bigl((1-s)\log\widetilde b(t)\bigr).
 \end{split}                                             \tag{C.36.9}
\]
After projection to $M_B$, this says
precisely that $R^{O_1}_{s,t}$ and
$R^{O_2}_{s,t}$ agree on $B$.
The intersection of the two closures
is a union of such $G$-orbits, so
they agree everywhere on the
intersection. This includes circle
orbits with their actual based
universal-cover lifts; no choice
of an unrecorded integer rotation
is made.

For $h\in H$, conjugation carries
$Q_O$ onto $Q_{hO}$ and $N_O$
onto $N_{hO}$, as follows from
the exact restriction of the
setwise stabilizers. This is a
topological Lie-group isomorphism.
It carries $n_O(t)$ to
$n_{hO}(t)$, since $h$ commutes
with $f_t$. Its based lift and
exponential naturality show
\[
                    R^{hO}_{s,t}
                         =hR^O_{s,t}h^{-1}.
                                                        \tag{C.36.10}
\]
Thus the compatible families have
the asserted operator equivariance.

Finally, on a finite union of
these closures, their continuous
maps agree on every overlap.
The finite closed-set pasting
lemma gives joint continuity.
Each intersection is preserved,
because its $G$-orbits are
preserved; hence the inverse
maps also agree there.
Applying the same pasting
argument to the inverse families
gives the stated family of
homeomorphisms on the finite
union. This proves the claim.
\end{proof}

\paragraph{Exact scope of the closure construction.}
The quotients (C.36.2)--(C.36.3)
provide continuity on each
individual closure, and
(C.36.7)--(C.36.9) enforce the
actual common boundary restriction
relations. In particular they
respect the three dependent edge
characters in C.35.
The present proof uses finite
closed-set pasting. C.37
subsequently proves a uniform
three-word estimate and extends
these contractions continuously,
with their inverses, to the
closed union $X_2$ of all such
open orbits, including infinite
accumulation. Extension across
$X_2$ to its exterior, including
frontier contractions for
$\dim Q_O=0,1$ and compatible
choices elsewhere on $Y$,
remains required for a global
conclusion in C.12.

\begin{claim}{C.37}{Proved}{Uniform control on all two-dimensional operator closures}
Let $f_t$, $0\leq t\leq1$, be an ambient-null based loop in
$\operatorname{Homeo}^+(\R^2)$, let $G=\langle f_t:0\leq t\leq1\rangle$,
and let $H$ be any subgroup commuting with every $f_t$.
For an open $G$-orbit $O$, let $Q_O$ be the actual closed restricted
operator image of C.30. Put
\[
 U_2=\bigcup_{\dim Q_O=2}O,\qquad X_2=\overline{U_2}.
                                                        \tag{C.37.1}
\]
The contractions of C.36 extend to a jointly continuous based
contraction through homeomorphisms of $X_2$, with jointly continuous
inverses, commuting with $H$. This includes arbitrary accumulation
of the open orbits in (C.37.1).

The uniform estimate used in the proof is the following exact
word statement. In either simply connected two-dimensional Lie
group $\R^2$ or $\operatorname{Aff}^+(\R)$, every continuous path
$k:[0,t]\to N$ from $1$ to $g$ satisfies
\[
 \exp(r\log g)
   =k(b)k(a)^{-1}
 \quad\hbox{or}\quad
 \exp(r\log g)=g\,k(b)k(a)^{-1}
 \quad(0\leq r\leq1)                                   \tag{C.37.2}
\]
for suitable $a,b\in[0,t]$. The same assertion is used in the
universal covering group when the actual group is $\R\times S^1$.
Consequently every contraction value on an open orbit in $U_2$
is represented by a word of length at most three in the original
loop and its inverse, with equality of the specified based
universal-cover lifts.
\end{claim}
\passage{Section 23, folios 140-141, scans 279-281; Part C(b) infinite-operator request; uniform dimension-two open-orbit closure step for C.12}


\begin{proof}
\emph{The three-word estimate.}
For a continuous path $\gamma:[0,1]\to\R^2$ with
$\gamma(0)=0$, $\gamma(1)=v$, one has
\[
 rv=\gamma(b)-\gamma(a)
 \quad\hbox{or}\quad
 rv=v+\gamma(b)-\gamma(a) \qquad(0\leq r\leq1).           \tag{C.37.3}
\]
Here is an exact measure-theoretic proof that applies to arbitrary
paths. For a polygonal path apply the signed-measure splitting
theorem of Grandoni--Ravi--Singh--Zenklusen
\cite[Theorem 4.7 and \S4.3.2]{C_GrandoniRaviSinghZenklusen2014}
to the two signed measures with densities $\gamma_1'$ and
$\gamma_2'$. A single circular interval has fraction $r$ of both
total measures. Its integral is either a path difference or the
total vector plus a path difference, giving (C.37.3).
Approximate an arbitrary continuous path uniformly by polygonal
paths with the same endpoints. Compactness of the two endpoints
and of the two alternatives passes the equation to the limit.
The planar case of the cited splitting proof uses at most two
cuts of an interval; the two equal-mass halves together have at
most three pieces, so one is a single interval. This is the
one-circular-interval count required here.

For $N=\R^2$, (C.37.2) is (C.37.3).
Write the affine group in the exact coordinates of C.36:
\[
 (a,b)(a',b')=(a+a',\,b+e^a b').
                                                        \tag{C.37.4}
\]
Every nonidentity $g=(a,b)$ lies on its unique exponential
one-parameter subgroup. If $a\ne0$, conjugation by
$(0,b/(e^a-1))$ sends $g$ to $(a,0)$.
The global coordinate map
\[
 (a',b')\longmapsto (a'/a,\,e^{-a'}b')
                                                        \tag{C.37.5}
\]
then sends left multiplication by $g^r$ to translation by
$(r,0)$, and sends $1,g$ to $(0,0),(1,0)$.
If $g=(0,b)$ with $b\ne0$, use $(a',b')\mapsto(b'/b,a')$
instead. Apply (C.37.3) in these coordinates.
A difference $(r,0)$ says $k(b)=g^r k(a)$; the other
alternative says $k(a)=g^{1-r}k(b)$. These are exactly
(C.37.2). The identity case uses the identity word.
This proves the estimate also for the simply connected
covering group of every group in C.36.

\emph{The compact family of actual word isotopies.}
Use $r=1-s$ for the contraction parameter, and set
\[
 R^O_{r,t}=p_O\exp\bigl(r\log\widetilde n_O(t)\bigr),
 \qquad n_O(t)=f_t|_O,                                  \tag{C.37.6}
\]
as in C.36. Consider the following compact family of actual
plane isotopies, with the compact-open topology on maps and
inverses:
\[
 W_{\epsilon,a,b,t}(u)
       = f_{ut}^{\,\epsilon} f_{ub}f_{ua}^{-1},
 \quad 0\leq a,b\leq t\leq1,\quad
       \epsilon\in\{0,1\},\quad 0\leq u\leq1 .
                                                        \tag{C.37.7}
\]
Compactness follows from continuity of the original loop and
of multiplication and inversion. Formula (C.37.2), applied to
the based lift of its prefix $[0,t]$, represents every
$R^O_{r,t}$ by the endpoint of one of (C.37.7).
More precisely, the based lift of that word isotopy in $N_O$
has endpoint
$\exp(r\log\widetilde n_O(t))$. Thus the representation retains
the lift, not only the projected homeomorphism.

\emph{The exact groups on nonopen limit orbits.}
Each $N_O$ in (C.37.6) is metabelian. Therefore the identity
\[
          [[g_1,g_2],[g_3,g_4]]=1
                                                        \tag{C.37.8}
\]
holds on $U_2$ for all $g_i\in G$, and hence on $X_2$ by
continuity. The set $X_2$ is $G$-invariant, so it is a union
of whole $G$-orbits. Its open $G$-orbits are precisely those
already in $U_2$: any other open orbit has an open
neighborhood disjoint from $U_2$.

Let $B\subset X_2$ be a nontrivial nonopen orbit, endowed
with the intrinsic line or circle topology of C.28.
The image $G_B$ in $\operatorname{Homeo}^+(B)$ is transitive,
metabelian, and path connected. For the last assertion,
every finite word has a path from the identity obtained from
the original paths and their inverses; C.28 gives intrinsic
compact-open continuity of all these finite-word paths.
We now prove the exact classification needed at $B$:
\[
 G_B\cong
 \begin{cases}
   (\R,+)\ \hbox{or }\operatorname{Aff}^+(\R),& B\text{ a line},\\
   S^1,& B\text{ a circle},
 \end{cases}                                            \tag{C.37.9}
\]
with the usual transitive actions after conjugacy.

First, a transitive abelian subgroup $A$ of
$\operatorname{Homeo}^+(B)$ has abelian closure $\overline A$,
and this closure acts freely: an element fixing one point
commutes with $A$, so fixes its whole transitive orbit.
The closure is Polish and transitive. The open orbit-map
theorem used in C.26 identifies it homeomorphically with $B$
by evaluation, since its point stabilizer is trivial.
In particular it is locally compact. Its faithful product
action on the three-manifold $B\times\R^2$ makes it a Lie
group by the same Hilbert--Smith input as C.33.
It is the one-dimensional connected group $\R$ or $S^1$,
and evaluation conjugates its action to translation.
Since $A$ was already transitive, injectivity of evaluation
also proves $A=\overline A$.

If $G_B$ is nonabelian, its derived subgroup $D$ is
nontrivial, abelian, normal, and path connected.
Path connectedness follows by taking commutator paths and
finite products of them. The common fixed set of $D$ is
$G_B$-invariant. Transitivity and nontriviality make it empty.
Every $D$-orbit is therefore a connected nontrivial interval
or circle. Such an orbit has no included endpoint, since
that endpoint would be preserved by its transitive
orientation-preserving action. All these orbits are open in
$B$; connectedness of $B$ forces $D$ to be transitive.
The preceding paragraph identifies $D$ with the full
translation group.

The orientation-preserving normalizer of translations on
the line consists exactly of $x\mapsto cx+d$, $c>0$.
Indeed normalization says $h(x+a)=h(x)+\chi(a)$;
continuity and additivity give $\chi(a)=ca$.
The multiplier image of the path-connected group $G_B$
is a connected subgroup of $\R_{>0}$. It is nontrivial in
the present case, so it is all $\R_{>0}$; because $G_B$
contains every translation, it is the full affine group.
On the circle, the orientation-preserving normalizer of
the full rotation group is just that rotation group:
the induced continuous automorphism of $S^1$ has degree
$+1$. This would contradict nonabelianness. These
arguments prove (C.37.9), including the topology of its
image in the intrinsic homeomorphism group.

For a line $B$, define $R^B_{r,t}$ by the logarithmic formula
in this simply connected group. For a circle, use the
unique based real lift of $f_t|_B$ to the cover of its
rotation group and scale that lift by $r$.
Its endpoint is zero by the exact winding comparison C.29.
On a point orbit define $R^B_{r,t}=1$.
These definitions are independent of the conjugacies in
(C.37.9), since Lie homomorphisms commute with exponential
maps and with based universal-cover lifts.
They agree with C.36 on every boundary orbit to which
that construction already applies.

\emph{Limits from the dense open union.}
Take
\[
 x_i\in O_i\subset U_2,\quad x_i\longrightarrow x\in X_2,
 \qquad (r_i,t_i)\longrightarrow(r,t).
                                                        \tag{C.37.10}
\]
Choose word isotopies $W_i$ from (C.37.7) representing
$R^{O_i}_{r_i,t_i}$ with their specified based lifts.
After a subsequence all word parameters converge.
Then $W_i$ and their inverses converge uniformly on compact
sets, uniformly in $u$, to a word isotopy $W$ from the
same family. Write $w=W(1)$ and $g=f_t$.
If $x$ is in an open orbit of $X_2$, eventually $x_i$ is
in that same open orbit, where (C.37.6) is continuous.
If $x$ is a point orbit, $w(x)=x$ for every possible
limiting word, which gives the required limit immediately.
It remains to identify $w$ on the intrinsic line or circle
$B=Gx$.

We will repeatedly use the following exact passage of
identities. For every fixed $h\in G$,
$hx_i\in O_i$ and $hx_i\to hx$. An identity of word
homeomorphisms on $O_i$, whose word isotopies converge,
therefore passes to every point $hx$ of $B$.
In particular
\[
                        [w|_B,g|_B]=1,                 \tag{C.37.11}
\]
because a fractional power commutes with its endpoint
on $O_i$.

Fix a positive integer $n$ and set $m_i=\lfloor nr_i\rfloor$.
In the universal covering group of $N_{O_i}$ the exact
residual is
\[
 \widetilde w_i^{\,n}\,
       \widetilde n_{O_i}(t_i)^{-m_i}
    =\exp\bigl((nr_i-m_i)
                  \log\widetilde n_{O_i}(t_i)\bigr).
                                                        \tag{C.37.12}
\]
Its exponent fraction lies in $[0,1)$.
By (C.37.2), the right side is again represented by one of
the compact family (C.37.7), with equality of based lifts.
Passing to a further subsequence for this fixed $n$,
the integer $m_i$ is constant, say $m_n$, and the residual
word isotopies converge to $Z_n$. Hence
\[
 (w|_B)^n(g|_B)^{-m_n}=z_n,\qquad
 |m_n-nr|\leq1,\qquad z_n=Z_n(1)|_B.                    \tag{C.37.13}
\]
All $z_n$ belong to the same compact set of restricted
three-word endpoints, independently of $n$.
No simultaneous subsequence for all $n$ is required.

For a line, (C.37.9) and (C.37.11) allow logarithms in
$\R$ or $\operatorname{Aff}^+(\R)$.
In the latter group the centralizer of every nonidentity
element is its exponential one-parameter subgroup, as is
seen directly from (C.37.4). Thus logarithms add for
commuting elements in both groups. The compact set in
(C.37.13) has bounded logarithms, giving
\[
 n\log(w|_B)-m_n\log(g|_B)=\log z_n,\qquad
 \sup_n\|\log z_n\|<\infty.
                                                        \tag{C.37.14}
\]
Divide by $n$ and let $n$ tend to infinity.
It follows that $\log(w|_B)=r\log(g|_B)$.
This includes $g|_B=1$, in which case boundedness forces
$w|_B=1$. The limiting word therefore has exactly the
required fractional action on the line.

\emph{The lifted residual bound on a circle.}
Identify the action in (C.37.9) on a circle with rotations,
and let $\theta(u)$ be the continuous real angle of
$f_u|_B$, normalized by $\theta(0)=0$.
The normalization uses rotations $\R/\Z$.
C.29 gives $\theta(1)=0$.
Every isotopy in (C.37.7), restricted to $B$, has a
canonical real endpoint angle
\[
 \epsilon\theta(t)+\theta(b)-\theta(a).
                                                        \tag{C.37.15}
\]
Let $A$ be this angle for $W$, let $T=\theta(t)$,
and let $C_n$ be the angle for $Z_n$.
Writing $M=\max_{0\leq u\leq1}|\theta(u)|$ gives
$|C_n|\leq3M$.
We must prove the real, rather than merely circle-valued,
identity
\[
                     nA-m_nT=C_n.                      \tag{C.37.16}
\]

The word isotopy
\[
 \Lambda_i(u)=W_i(u)^n f_{u t_i}^{-m_i}Z_i(u)^{-1}
                                                        \tag{C.37.17}
\]
restricts to a based loop in $N_{O_i}$.
Equation (C.37.12), with the specified lifts of $W_i$ and
$Z_i$, says that its universal-cover lift ends at the
identity. Hence that loop is null-homotopic in $N_{O_i}$.
Choose $h\in G$ with $hx\ne x$, and evaluate (C.37.17)
at the two points $x_i,hx_i$ of $O_i$.
For large $i$ they are distinct. The resulting loop in
$\operatorname{Conf}_2(\R^2)$ has difference-direction
winding zero. The word isotopies converge uniformly on
compact sets to
$\Lambda(u)=W(u)^n f_{ut}^{-m_n}Z_n(u)^{-1}$.
The two-point loops therefore converge uniformly to its
evaluation at $x,hx$. This limiting pair stays distinct,
so winding is preserved.
On the Jordan circle $B$ its winding, by C.29, is exactly
$nA-m_nT-C_n$ (up to a common orientation sign).
It is zero, proving (C.37.16).

Using (C.37.13) and $|C_n|\leq3M$ now gives
\[
        |A-rT|\leq\frac{3M+|T|}{n}\quad(n\geq1).
                                                        \tag{C.37.18}
\]
Consequently $A=rT$. This determines the based lifted
fractional action even when $f_t|_B=1$ but $\theta(t)$
is a nonzero integer. Thus no hidden winding survives
the limit from the open orbits.

\emph{Joint continuity, inverses, and equivariance.}
We have shown that every convergent subsequence of word
parameters in (C.37.10) gives the same endpoint
$R^B_{r,t}(x)$. Compactness of (C.37.7) proves
\[
          R^{O_i}_{r_i,t_i}(x_i)
                  \longrightarrow R^B_{r,t}(x).
                                                        \tag{C.37.19}
\]
This holds for every sequence from the dense set $U_2$
approaching every point of $X_2$, with parameters also
varying. It proves joint continuity on all of $X_2$.
Explicitly, for a sequence of points of $X_2$, approximate
each by a point of $U_2$ so closely that both its distance
and its image distance for that fixed parameter are
less than $1/i$. Such approximations exist by
(C.37.19) with parameters fixed. Applying (C.37.19) to
the chosen approximating sequence gives the full
continuity assertion.

The inverse loop $f_t^{-1}$ has the same group $G$,
the same sets $U_2,X_2$, and the same actual operator
images. Its construction gives precisely
$(R^B_{r,t})^{-1}$, by inversion of the based lifts.
Applying the preceding argument to it proves joint
continuity of the inverse family as well. Each member
is an orbitwise bijection, hence a homeomorphism of
$X_2$. The identities $R_{0,t}=1$, $R_{1,t}=f_t|_{X_2}$,
and $R_{r,0}=R_{r,1}=1$ give the based contraction.

Finally every $h\in H$ permutes the open $G$-orbits and
conjugates their actual closed operator images.
It therefore preserves $U_2$ and $X_2$.
C.36 gives exact transported equivariance on $U_2$.
The established joint continuity extends that equality
to its closure $X_2$. This proves all assertions.
\end{proof}

\paragraph{The remaining global boundary requirement.}
C.37 treats the whole closed set $X_2$, including infinite
accumulation of its open orbits. Its boundary contraction
is the based Lie-group fractional action forced by
(C.37.14) and (C.37.16). Extending this family to the
complement of $X_2$, compatibly with the dimension-zero
and dimension-one proper operator cases, remains a
separate requirement. Even where that complement
contains only nonopen orbits, its median contraction
from C.34 has not been identified with this boundary
family. C.35 proves that imposing the fixed median
prescription can violate the actual restriction map.
Thus the uniform estimate proved here closes the
dimension-two accumulation step, while the global
centralizer conclusion in C.12 remains open.

\begin{claim}{C.38}{Proved}{Discrete cyclic operators on nonabelian line and circle orbits}
Let $B$ be a line or circle and let
$T\leq\operatorname{Homeo}^+(B)$ be transitive and
nonabelian. Its full compact-open centralizer
\[
                  N_B=C_{\operatorname{Homeo}(B)}(T)
                                                        \tag{C.38.1}
\]
is discrete and acts freely and properly. On a line it is
trivial or infinite cyclic; on a circle it is finite cyclic.

In the plane-loop setting of C.28, let $B$ be a nonopen
orbit whose intrinsic image $G|_B$ is nonabelian, and let
$H_B=\{h\in H:hB=B\}$ be the operator stabilizer.
Its actual restriction has the exact sequence
\[
 1\longrightarrow K_B\longrightarrow H_B
       \xrightarrow{\ \rho_B\ } P_B\longrightarrow1,
 \quad
 K_B=\{h\in H_B:h|_B=1\},\quad P_B=\rho_B(H_B)\leq N_B.
                                                        \tag{C.38.2}
\]
In the intrinsic compact-open topology, $P_B$ is already
closed and discrete, acts freely and properly, and
\[
                    p_B:B\longrightarrow B/P_B
                                                        \tag{C.38.3}
\]
is an actual regular covering. Its deck transformations
are exactly $P_B$. The quotient is a line for trivial
$P_B$ on a line, and otherwise a circle. On a circle
$P_B$ is finite and the covering degree is $|P_B|$.
The kernel $K_B$ acts identically on $B$ throughout this
comparison.
\end{claim}
\passage{Section 23, folios 140-141, scans 279-281; Part C(b) infinite-operator request; exact one-dimensional operator-image reduction at the boundary remaining in C.12}


\begin{proof}
Let $L$ be the closure of $T$ in the compact-open
homeomorphism group of $B$, choose $b\in B$, and put
$J=L_b$. The group $L$ is Polish and transitive, and
$C(L)=C(T)$ by continuity of composition.
The open orbit-map theorem used in C.26 gives
$L/J\cong B$. The normalizer $N_L(J)$ is closed, and
the same exact evaluation construction as in C.30 gives
\[
 \operatorname{ev}_b:N_B
       \xrightarrow{\ \cong\ }N_L(J)/J\ \subset B.
                                                        \tag{C.38.4}
\]
The right side is closed. The inverse sends the coset
represented by $k$ to the map $lb\mapsto lkb$.
The normalizer equation makes this well defined.
Its compact-open continuity follows by the same
microtransitivity argument as C.30: convergent points
of $B$ lift locally through the open orbit map to
elements of $L$ tending to the identity.
No surface-dimensional step occurs in that argument.
The construction also proves freeness, local compactness,
and properness. Explicitly, if
$n_i b_i=c_i$, $b_i\to b$, and $c_i\to c$, choose
$l_i\to1$ with $l_i b=b_i$. Then
$n_i b=l_i^{-1}c_i\to c$, and the closed image and
homeomorphism (C.38.4) give a convergent subsequence
in $N_B$. This is the required properness test.

Every orientation-reversing homeomorphism of a line
or circle has a fixed point. Freeness therefore places
$N_B$ in $\operatorname{Homeo}^+(B)$.
Its faithful product action on the three-manifold
$B\times\R^2$ makes it a Lie group by the verified
Hilbert--Smith theorem used in C.33.
Equation (C.38.4) bounds its dimension by one.
If that dimension were one, every orbit of $N_B^0$
would be open in $B$: its one-dimensional orbit map
is injective and hence open by invariance of domain.
Connectedness of $B$ would make $N_B^0$ transitive.
Freeness would then give $N_B=N_B^0$, since an element
of $N_B^0$ with the same evaluation as an element of
$N_B$ equals it. This is the translation group $\R$
or the rotation group $S^1$ in evaluation coordinates.
Its centralizer is that same abelian group of right
translations. But $T\leq C(N_B)$, contradicting the
assumed nonabelianness of $T$.
Consequently $N_B$ is a zero-dimensional Lie group,
so is discrete.

A free proper discrete action supplies evenly covered
neighborhoods, and its orbit space is a connected
one-manifold without boundary. For a line this
quotient is a line or a circle. Since the line is
simply connected, the first case has trivial deck
group and the second has deck group $\Z$.
For a circle the quotient is a circle and the
covering has finitely many sheets; its deck group
is finite cyclic. This proves the first assertion.

In the stated plane-loop application, C.28 proves
that each operator stabilizing $B$ induces an actual
homeomorphism for its intrinsic line or circle
topology. Commutation with $G$ puts its image in
(C.38.1). Every subgroup of this discrete group is
closed; restriction of a proper free action to a
closed subgroup is proper and free. These facts
give (C.38.3). Its fibers are precisely the
$P_B$-orbits. A deck transformation is determined
by one point in the connected covering, and each
point of a fiber is obtained by exactly one
element of $P_B$, so the whole deck group is
exactly $P_B$. The same one-manifold covering
classification gives the stated quotient and
degree. Equation (C.38.2) retains the original
stabilizer and its restriction kernel.
\end{proof}

\paragraph{Topology and the remaining relative extension.}
The topology asserted for $P_B$ is its intrinsic
compact-open image topology. The operator group may
be regarded as discrete as in the original problem;
no assertion about continuity of the restriction
map from an inherited plane topology to a recurrent
intrinsic line topology is needed.
C.38 gives exact covering quotients on every
nonabelian one-dimensional orbit. It supplies no
transverse control between different orbits and no
relative extension of the C.37 family from $X_2$
to its exterior. Those requirements remain in C.12.

\begin{claim}{C.39}{Proved}{A global contraction for loops with commuting values}
Let $f_t$, $0\leq t\leq1$, be a based loop in
$\operatorname{Homeo}^+(\R^2)$ with
$f_s f_t=f_t f_s$ for every $s,t$.
Suppose the loop is null-homotopic in that ambient group.
For every subgroup $H$ commuting with all $f_t$, the loop
has a based null-homotopy in
$C_{\operatorname{Homeo}^+(\R^2)}(H)$.
The homotopy and its inverse family are jointly continuous
on the whole plane, including every change and accumulation
of orbit types.

Consequently, on every connected admissible surface $U$
of C.1, an element of the kernel of the exact inclusion map
\[
 \pi_1 C_{\operatorname{Homeo}(U)}(H)
      \longrightarrow\pi_1\operatorname{Homeo}(U)
                                                        \tag{C.39.1}
\]
that is represented by a loop with pairwise commuting values
is zero. On a hyperbolic surface every such based loop in
$C_{\operatorname{Homeo}(U)^0}(H)$ is therefore null-homotopic
in that centralizer.
\end{claim}
\passage{Section 23, folios 140-141, scans 279-281; Part C(b) infinite-operator request; complete global commuting-value subcase of the exact C.12 level-zero comparison}


\begin{proof}
\emph{The actual image on every orbit.}
Set $G=\langle f_t\rangle$. This is an abelian group.
The orbit classification C.28 applies.
On an open orbit $O$, the closure $K_O$ of $G|_O$ in
$\operatorname{Homeo}(O)$ is Polish, abelian, connected,
and transitive. It acts freely: an element fixing a point
commutes with the transitive subgroup $G|_O$ and hence fixes
every point. The Polish open orbit-map theorem identifies
evaluation $K_O\to O$ homeomorphically, with trivial
stabilizer. Thus $K_O$ is locally compact. Its faithful
action on $O\times\R$ and the Hilbert--Smith input of C.33
make it a two-dimensional abelian Lie group.
Since $O$ is a noncompact open subset of the plane, the
possibilities are $\R^2$ and $\R\times S^1$.
The original image is already the whole group:
for each $k\in K_O$, transitivity of $G|_O$ gives an
element with the same evaluation, and freeness makes
the two elements equal.

On an intrinsic line or circle orbit $B$, C.28 gives
compact-open continuity of the restricted original paths.
The same abelian closure and free evaluation argument,
now in the intrinsic topology, identifies $G|_B$ with
the full translation group $\R$ or rotation group $S^1$.
This is also the abelian case of (C.37.9).
On a point orbit the image is trivial.
Denote the resulting actual Lie image on any orbit $B$
by $K_B$. Its simply connected covering group is a vector
group of dimension at most two.

The original restricted loop has a closed based lift
$\widetilde n_B(t)$ to this vector group.
For an open orbit, C.32 gives its null-homotopy in
$\operatorname{Homeo}(O)$; evaluation and the
homeomorphism $K_O\cong O$ then make its class in
$\pi_1K_O$ zero. For a line there is no covering kernel.
For a circle the based lift closes by the exact winding
comparison C.29. Define on each orbit
\[
              R^B_{r,t}
                =p_B\bigl(r\,\widetilde n_B(t)\bigr),
                    \qquad 0\leq r,t\leq1.
                                                        \tag{C.39.2}
\]
Here $p_B$ is the Lie-group covering homomorphism,
with the trivial interpretation for a point orbit.
The action of the right side on $B$ is its actual
homeomorphism action. Its endpoints are the identity
at $r=0$ and the original loop at $r=1$; both based
edges $t=0,1$ are fixed.

\emph{One compact word family for all orbit types.}
In dimension two, C.37 represents every value of
(C.39.2) by a word isotopy
\[
 W_{\epsilon,a,b,t}(u)
        =f_{ut}^{\,\epsilon}f_{ub}f_{ua}^{-1},
 \quad \epsilon\in\{0,1\},\quad 0\leq a,b\leq t\leq1,
                                                        \tag{C.39.3}
\]
with equality of its specified based universal-cover
endpoint. In dimension one the intermediate value
theorem already finds $b\leq t$ with
$\widetilde n_B(b)=r\widetilde n_B(t)$.
This is a member of (C.39.3) with $\epsilon=0,a=0$.
The identity word handles point orbits.
Thus the same compact family of original plane isotopies
works for every orbit, with the based lifts retained.

Take arbitrary points $x_i\to x$ in the plane and
parameters $(r_i,t_i)\to(r,t)$, and write $B_i=Gx_i$,
$B=Gx$. Choose representing word isotopies $W_i$ for
$R^{B_i}_{r_i,t_i}$. Compactness of the parameter set
in (C.39.3) gives a subsequence converging, together
with inverses and uniformly in $u$, to a word isotopy
$W$. Put $w=W(1)$.
If $B$ is open, all sufficiently large $x_i$ are in
that same orbit and continuity follows from its
Lie action. If $B$ is a point, every possible word
limit fixes $x$, giving the required convergence.
It remains to identify the limiting word on a line
or circle. This argument explicitly allows the $B_i$
to have different types.

For each fixed positive integer $n$, set
$m_i=\lfloor nr_i\rfloor$. In the vector covering
group of the actual image $K_{B_i}$, one has
\[
 n\,\widetilde w_i
       -m_i\,\widetilde n_{B_i}(t_i)
       =(nr_i-m_i)\widetilde n_{B_i}(t_i).
                                                        \tag{C.39.4}
\]
The right side is represented by another word isotopy
$Z_i$ from (C.39.3), with equality of based lifted
endpoints. For this fixed $n$, pass to a further
subsequence so that $m_i=m_n$ is constant and
$Z_i\to Z_n$. Then $|m_n-nr|\leq1$.
Identities on $B_i$ pass to the whole orbit $B$ by
evaluation at $hx_i\to hx$, for any fixed $h\in G$.
In particular, with $g=f_t$ and $z_n=Z_n(1)$,
\[
             (w|_B)^n(g|_B)^{-m_n}=z_n|_B.
                                                        \tag{C.39.5}
\]

For a line, let $\theta(u)$ be the real translation
parameter of $f_u|_B$, and put
$M=\max_u|\theta(u)|$.
The real parameters $A$ of $w|_B$ and $C_n$ of
$z_n|_B$ come from their word paths:
each has the form
$\epsilon\theta(t)+\theta(b)-\theta(a)$.
Consequently $|C_n|\leq3M$, and (C.39.5) is the
real identity $nA-m_n\theta(t)=C_n$.
Division by $n$ gives $A=r\theta(t)$, as required.

For a circle, use the based real rotation lift
$\theta$ of the original loop, so $\theta(0)=\theta(1)=0$.
The word paths again specify real endpoint angles
$A,C_n$ with $|C_n|\leq3M$, but (C.39.5) initially
gives their relation only modulo integers.
The exact lift in (C.39.4) removes this ambiguity.
Indeed,
\[
       \Lambda_i(u)
         =W_i(u)^n f_{ut_i}^{-m_i}Z_i(u)^{-1}
                                                        \tag{C.39.6}
\]
is a based loop in the actual group $K_{B_i}$ and
has a closed based universal-cover lift.
It is therefore null-homotopic in that group.
Choose $h\in G$ with $hx\ne x$.
For large $i$ the two points $x_i,hx_i$ of $B_i$
are distinct. Evaluation of (C.39.6) at them has
difference-direction winding zero in
$\operatorname{Conf}_2(\R^2)$, regardless of whether
$B_i$ is an open, line, or circle orbit.
This follows from the actual null-homotopy in
$K_{B_i}$ and its continuous evaluation into the
plane; C.28 supplies that evaluation for intrinsic
line orbits.

The word isotopies and the two evaluation points
converge. Their limiting pair stays distinct.
The winding therefore remains zero. On the limiting
Jordan circle its value, by C.29, is
$nA-m_n\theta(t)-C_n$, up to a common orientation
sign. Thus the real identity
\[
 nA-m_n\theta(t)=C_n,\qquad
 |A-r\theta(t)|
       \leq\frac{3M+|\theta(t)|}{n}
                                                        \tag{C.39.7}
\]
holds for every $n$. It forces $A=r\theta(t)$.
This also treats an identity endpoint rotation
whose based angle $\theta(t)$ is a nonzero integer.
All possible subsequential word limits therefore
give the prescribed value (C.39.2).

Compactness of (C.39.3) now proves joint continuity
of (C.39.2) on the entire plane.
The inverse loop has the same orbits and image
groups, and its construction gives precisely the
inverse of (C.39.2); the same proof gives joint
inverse continuity. Each map is an orbitwise
bijection, so these are plane homeomorphisms.
Every commuting operator carries the actual
abelian Lie image on one orbit isomorphically to
that on the transported orbit. These are continuous
Lie-group isomorphisms in the specified intrinsic
or open-orbit topology. Their based covering lifts
are linear, and hence commute with multiplication
by $r$. Formula (C.39.2) is therefore exactly
$H$-equivariant. Taking $r=1-s$ proves the plane
assertion.

\emph{The precise surface lift and descent.}
Every connected admissible surface in C.1 has
universal covering $p:\widetilde U\to U$ with
$\widetilde U\cong\R^2$. Let $f_t$ represent an
element of the kernel in (C.39.1), with its
commuting-value property. Choose an ambient
filling $A:D^2\to\operatorname{Homeo}(U)$ whose
boundary is this loop and whose specified based
boundary point maps to the identity.
The actual map
\[
 D^2\times\widetilde U\longrightarrow U,\qquad
                  (z,x)\longmapsto A(z)(p(x))
                                                        \tag{C.39.8}
\]
has a unique lift equal to the identity on the
specified based fiber. Its domain is simply
connected. Lifting the inverse family as well,
and using uniqueness, proves that its slices are
homeomorphisms with jointly continuous inverses.
It is an ambient filling of a closed based
plane loop $\widetilde f_t$.

For every lift $\widetilde h$ of an operator,
$[\widetilde f_t,\widetilde h]$ is a deck
transformation. It depends continuously on $t$
and starts at the identity, so it is identically
the identity. Similarly
$[\widetilde f_s,\widetilde f_t]$ is a deck
transformation depending continuously on $(s,t)$;
it is the identity on the axes, and hence
everywhere. Thus the lifted loop has commuting
values and centralizes the full group of
operator lifts, including the deck group.
Apply the plane assertion. Its equivariant
contraction descends through $p$, with both
joint continuities, to the original centralizer
on $U$. Its slices lie in the identity component
by the displayed paths. This proves (C.39.1).
For a hyperbolic $U$ the ambient identity
component is contractible, so the ambient
kernel condition is automatic.
\end{proof}

\paragraph{Scope of the global subcase.}
This theorem treats loops with commuting values on
every orbit at once, including open orbits whose
operator image has dimension zero or one.
It does not identify every fundamental-group class
of an arbitrary centralizer with such a loop.
No deformation of a general centralizer loop to
this subcase has been constructed here.
C.12 remains open for general loops.

\paragraph{The earlier abelian closed subgroup cannot be a centralizer.}
Fix the specific embedding $j:B\hookrightarrow A^0$ of C.26, where
$B=L^1([0,1],\Z)/\Z$, $A=\operatorname{Homeo}(U)$ and $U$ is
hyperbolic. For every $H\leq A$ commuting with $j(B)$, the induced map
$\pi_1(B)\to\pi_1(C_A(H),1)$ is zero. Indeed the displayed generator
of $\pi_1(B)=\Z$ has pairwise commuting values and is ambient-null,
so C.39 contracts its image; the induced homomorphism therefore
vanishes on the whole group. The based contraction lies in $A^0$,
so the same conclusion holds with target $\pi_1(C_{A^0}(H),1)$.
Consequently $j(B)$ cannot equal either of these centralizers,
since its topological-group identification with $B$ induces an
isomorphism on the nonzero fundamental group. This rules out that
exact realization as a zero-level counterexample.

\begin{claim}{C.40}{Proved}{The locally compact loop-closure case}
Let $U$ be any connected admissible surface of C.1, let
$A=\operatorname{Homeo}(U)$ with its compact-open group
topology, and let $H\leq A$ be arbitrary. Put $C=C_A(H)$.
For a based loop $f:[0,1]\to C$, form the actual closed
subgroup
\[
 L_f=\overline{\langle f_t:0\leq t\leq1\rangle}^{\,A}
\]
with the topology induced from $A$. Every class in the
kernel of
\[
 i_*:\pi_1(C,1)\longrightarrow\pi_1(A,1)^H
                                                        \tag{C.40.1}
\]
that has a representative with locally compact $L_f$
is zero. Here $i:C\hookrightarrow A$ is the inclusion;
its induced map lands in the displayed fixed subgroup
because each $h\in H$ fixes the entire representative
loop under conjugation.

In particular, the same conclusion holds for loops
contained in any closed locally compact subgroup of
$C$. If $C$ itself is locally compact, (C.40.1) is
injective. On a hyperbolic surface every based loop
in $C_{\operatorname{Homeo}(U)^0}(H)$ with locally
compact $L_f$ is null-homotopic in that centralizer.
\end{claim}
\passage{Section 23, folios 140-141, scans 279-281; Part C(b) infinite-operator request; exact inclusion map from the centralizer fundamental group to the H-invariant ambient fundamental group}


\begin{proof}
\emph{The exact subgroup and its Lie topology.}
Centralizers are closed in the Hausdorff topological
group $A$. Consequently $L=L_f$ is a closed subgroup
of $C$. The subgroup generated by the loop is
path connected: each finite product of loop values
and their inverses can be joined to the identity by
the corresponding product of paths. Its closure $L$
is connected. Thus no identity-component replacement
or change of topology is being made.
The topology on $L$ is the subspace topology from
$A$, and its action on $U$ is continuous and faithful.

Suppose $L$ is locally compact. The product action
\[
 L\times(U\times\R)\longrightarrow U\times\R,
       \qquad (l,(x,a))\longmapsto(lx,a)
                                                        \tag{C.40.2}
\]
is a faithful continuous action on a connected
three-manifold. The verified Hilbert--Smith theorem
used in C.33, \cite[Theorem 1.5 and the introductory
reduction]{C_Pardon2013}, therefore makes $L$ a Lie
group with precisely this topology. It is a connected
Lie group. All homotopies constructed inside $L$ below
are consequently homotopies through actual elements
of $C$, continuous in its given topology.

\emph{A based deformation into a compact subgroup.}
Borel's principal theorem gives a connected maximal
compact subgroup $K$ and closed one-parameter
subgroups $H_1,\ldots,H_d$, each isomorphic to $\R$,
such that multiplication is a homeomorphism
\[
 K\times\R^d\longrightarrow L,\qquad
 (k,a_1,\ldots,a_d)\longmapsto
       k\,h_1(a_1)\cdots h_d(a_d).
                                                        \tag{C.40.3}
\]
This is the actual product assertion in
\cite[\S1, p.~271; \S\S3--4, pp.~277--278]{C_Borel1952}:
the product notation is explicitly defined there
by uniqueness and continuous dependence of all
factors. The maps $h_j:\R\to H_j$ are topological
group isomorphisms with $h_j(0)=1$.
The identity has the unique coordinates $(1,0,\ldots,0)$.
Thus
\[
 D_s\bigl(kh_1(a_1)\cdots h_d(a_d)\bigr)
       =k h_1((1-s)a_1)\cdots h_d((1-s)a_d)
                                                        \tag{C.40.4}
\]
is a continuous strong deformation retraction
$D:L\times[0,1]\to L$ onto $K$, fixing the identity.
In particular $D_s(f_t)$ is a based homotopy in
$C$ from $f$ to a based loop $\beta$ in $K$.

\emph{A commuting representative of the same based class.}
Average an inner product on the Lie algebra of $K$
over its compact adjoint action. The resulting
positive definite inner product is adjoint
invariant, so left translation gives a bi-invariant
Riemannian metric on $K$. Its pullback to the
universal covering group
$p:\widetilde K\to K$ is also bi-invariant.
For left-invariant vector fields its Levi-Civita
connection satisfies
\[
               \nabla_XY=\tfrac12[X,Y].
                                                        \tag{C.40.5}
\]
Indeed the Koszul formula and adjoint invariance
give this identity directly. Hence all geodesics
through the identity are the one-parameter
subgroups, and their left translates are defined
for all real time. This proves geodesic completeness
of the lifted metric. The Hopf--Rinow theorem
therefore joins the identity to every point of
$\widetilde K$ by a geodesic.

Lift $\beta$ from the identity to a path
$\widetilde\beta$ with endpoint $z\in\ker p$.
Choose $X$ in the Lie algebra of $\widetilde K$
with $\exp_{\widetilde K}(X)=z$, using the geodesic
just obtained. The paths $\widetilde\beta(t)$
and $\exp_{\widetilde K}(tX)$ have the same
endpoints in the simply connected space
$\widetilde K$, so they are homotopic relative
to both endpoints. Projection through $p$ is a
based homotopy in $K$ from $\beta$ to
\[
             \gamma(t)=p\bigl(\exp_{\widetilde K}(tX)\bigr).
                                                        \tag{C.40.6}
\]
This is a based loop, since $p(z)=1$.
It has pairwise commuting values because it is
the restriction of a one-parameter subgroup.
Every step from $f$ to $\gamma$ stays in $C$.

Now let $[f]$ be in the kernel of (C.40.1).
The inclusion of the fixed subgroup
$\pi_1(A,1)^H$ into $\pi_1(A,1)$ is injective,
so $f$ is ambient-null. The based centralizer
homotopy just constructed makes $\gamma$
ambient-null as well. C.39 applies on this
same admissible surface to this actual commuting
loop and the same operator subgroup $H$.
It supplies a based contraction of $\gamma$
inside $C$. Concatenating the two based
homotopies contracts $f$ inside $C$.

A loop contained in a closed locally compact
subgroup $M\leq C$ has $L_f\subseteq M$ closed;
it is therefore locally compact. Taking $M=C$
proves the assertion about a locally compact
centralizer. In the hyperbolic case
$\operatorname{Homeo}(U)^0$ is contractible by
the earlier surface theorem, so the ambient
kernel condition holds for every based loop.
The constructed contraction lies in the identity
path component because it is a based homotopy;
it consequently lies in the displayed
centralizer inside $\operatorname{Homeo}(U)^0$.
\end{proof}

\paragraph{Exact remaining scope.}
C.40 proves a reduction to commuting values for
the specified locally compact closed subgroup.
No such based homotopy has been constructed for
a general loop whose generated closure is not
locally compact. The unrestricted comparison
C.12 remains open. In particular, the proof
does not replace the induced topology on that
closure by a different locally compact topology.

\begin{claim}{C.41}{Proved}{A dense operator orbit settles the zero-level kernel}
Let $U$ be a connected admissible surface of C.1, let
$A=\operatorname{Homeo}(U)$ with its compact-open group
topology, and let $H\leq A$ have a dense orbit in $U$.
Then the exact inclusion-induced comparison
\[
 \pi_1(C_A(H),1)\longrightarrow\pi_1(A,1)^H
                                                        \tag{C.41.1}
\]
is injective. In particular, on a hyperbolic surface
$\pi_1(C_{A^0}(H),1)=0$. No local compactness assumption
is imposed on $H$ or on the subgroup generated by
the values of the loop.
\end{claim}
\passage{Section 23, folios 140-141, scans 279-281; Part C(b) infinite-operator request; exact centralizer inclusion on fundamental groups for a dense operator orbit}


\begin{proof}
\emph{The plane argument.}
Let $f_t$ be a based ambient-null loop of plane
homeomorphisms centralizing a subgroup $H$ with
a dense orbit $D$. Put $G=\langle f_t\rangle$.
Since the loop is based, its maps preserve
orientation, and the orbit results C.28--C.37
apply. Every $h\in H$ transports $G$-orbits
according to the exact identity
\[
                       h(Gx)=G(hx).
                                                        \tag{C.41.2}
\]
If there is no open $G$-orbit, C.34 already
gives the required global equivariant contraction.

Otherwise let $V$ be the nonempty union of the
open $G$-orbits. It is open and $H$-invariant.
Density makes $D$ meet $V$, and invariance then
gives $D\subseteq V$. Consequently $V$ is dense
in the plane.

Fix an open $G$-orbit $O$ and choose $y\in D\cap O$.
Let $H_O=\{h\in H:hO=O\}$ be its operator
stabilizer. The actual orbit equality is
\[
                     H y\cap O=H_O y.
                                                        \tag{C.41.3}
\]
Indeed, $hy\in O$ implies that the two $G$-orbits
$hO=G(hy)$ and $O$ meet, hence are equal. The
reverse inclusion is immediate. Since $Hy=D$
is dense, (C.41.3) is dense in $O$.

Retain the exact restriction homomorphism
$H_O\to\operatorname{Homeo}(O)$, with its
pointwise-fixing kernel. Let $Q_O$ be the
closure of its image in that homeomorphism
group. C.30 proves that $Q_O$ is locally
compact and acts freely and properly on $O$;
C.33 makes it a Lie group of dimension at
most two. Properness makes every $Q_O$-orbit
closed. In particular $Q_Oy$ is a closed
subset of $O$ containing the dense set
$H_Oy$, so
\[
                         Q_Oy=O.
                                                        \tag{C.41.4}
\]
Thus this free proper action is transitive.
Its orbit map $Q_O\to O$ is a homeomorphism,
and $\dim Q_O=2$.

This holds for every open $G$-orbit. In the
notation of C.37, therefore,
\[
 U_2=V,\qquad
 X_2=\overline{U_2}=\R^2.
                                                        \tag{C.41.5}
\]
The contraction proved there is defined on
the whole plane, with jointly continuous
inverses, and centralizes every element of
the original $H$. This settles both alternatives
of the plane argument.

\emph{The actual surface comparison.}
Let a based loop $f_t$ represent a kernel class
in (C.41.1). The inclusion of the fixed
subgroup in $\pi_1(A,1)$ is injective, so
choose an ambient filling
$F:D^2\to A$ of that loop.
For the universal covering
$p:\widetilde U\to U$, every admissible
$\widetilde U$ is a plane. Lift the actual map
\[
 D^2\times\widetilde U\longrightarrow U,\qquad
           (z,x)\longmapsto F(z)(p(x))
                                                        \tag{C.41.6}
\]
with the identity prescribed on the based
boundary fiber. Its domain is simply connected.
As in C.39, lifting the inverse family and
using uniqueness gives an ambient filling by
plane homeomorphisms with jointly continuous
inverses. Its boundary is a closed based
plane loop $\widetilde f_t$.

Let $\widetilde H$ be the group of all lifts
of elements of $H$, including the deck group.
For every $\widetilde h\in\widetilde H$,
the commutator
$[\widetilde f_t,\widetilde h]$ is a deck
transformation depending continuously on $t$
and equal to the identity at $t=0$.
It is therefore always the identity.
The lifted loop centralizes precisely the
same operator lifts required for descent.

Choose $x\in U$ with dense $Hx$ and
$\widetilde x\in p^{-1}(x)$. One has the
exact equality
\[
       \widetilde H\,\widetilde x=p^{-1}(Hx).
                                                        \tag{C.41.7}
\]
One inclusion follows by projection. For the
other, choose a lift of the specified
$h\in H$ and then a deck transformation
taking its value at $\widetilde x$ to the
specified point over $hx$.
The right side is dense because $p$ is open:
every nonempty open set upstairs projects
to a nonempty open set meeting $Hx$.

Apply the plane result to $\widetilde H$.
Its contraction commutes with all deck
transformations and all operator lifts.
The equation
$\overline R_{s,t}p=pR_{s,t}$ therefore
defines a based contraction downstairs
inside $C_A(H)$. The open covering quotient
and its local sections give joint continuity
of these maps and their inverses.
It contracts the original class in the
source of (C.41.1), proving injectivity.
For a hyperbolic $U$, the ambient identity
component is contractible, so its based
centralizer loops all satisfy the kernel
condition. The based contraction stays in
$A^0$.
\end{proof}

\paragraph{Remaining scope.}
C.41 settles every operator action with a
dense orbit. It does not assert that a
general operator group has such an orbit
or supply an equivariant decomposition into
subsets on which this argument applies.
The general operator problem without a dense orbit
in C.12 remains open.

\begin{claim}{C.42}{Proved}{The actual boundary quotient has proper line components upstairs}
Let $G$ be generated by a based plane loop as in C.28,
let $O$ be an open $G$-orbit, and put
$L=\overline G\leq\operatorname{Homeo}(\R^2)$.
Fix $x\in O$ and $J=L_x$. For
$z\in\operatorname{Fix}(J)$ whose $G$-orbit $B=Gz$
is a nontrivial intrinsic line or a circle, the formula
\[
 \Phi_z(lx)=lz
\]
defines the actual continuous open quotient
$O=L/J\to B=L/L_z$, with the intrinsic topology on $B$.
For the universal coverings $p:\widetilde O\to O$
and $q:\R\to B$, every lift
$F:\widetilde O\to\R$ of $\Phi_zp$ is open and onto,
and every connected component of every fiber of $F$
is a properly embedded line in the covering plane
$\widetilde O$.
Every nondegenerate compact connected component of
a fiber of $\Phi_z$ is a circle.
Whole fibers of $\Phi_z$ can be disconnected, including
for quotients obtained as limits of the original
commuting operators.
\end{claim}
\passage{Section 23, folios 140-141, scans 279-281; Part C(b) infinite-operator request; exact boundary quotients in the remaining centralizer extension problem}


\begin{proof}
\emph{The quotient and its topology.}
Every element of $L$ preserves $O$: forward and inverse
limits preserve its boundary, while a forward limit
sends an interior point into $\overline O$ and cannot
send it to that boundary. The Polish orbit-map theorem
used in C.26 and C.30 therefore makes $l\mapsto lx$
an open quotient $L\to O$.
If $l_1x=l_2x$, then $l_2^{-1}l_1\in J$ fixes $z$,
so $\Phi_z$ is well defined and continuous.
More precisely, the continuous action map
\[
 \operatorname{Fix}(J)\times L\longrightarrow\R^2,
 \qquad (z,l)\longmapsto lz
\]
descends through the open quotient
$\operatorname{id}\times(L\to O)$.
Consequently
\[
 \operatorname{Map}_L(O,\R^2)
       \xrightarrow[\cong]{\,\operatorname{ev}_x\,}
       \operatorname{Fix}(J)
                                                        \tag{C.42.1}
\]
is a homeomorphism in the compact-open topology:
equivariance proves bijectivity, and the displayed
joint action proves continuity of the inverse.
For $l\in L$, choose $g\in G$ with $gx=lx$.
Then $g^{-1}l\in J$, hence $lz=gz$.
Thus the range is exactly $Lz=Gz=B$.

The map into an intrinsic line is continuous in its
intrinsic topology. Indeed, the image of a compact
coordinate disk in $O$ is a compact Peano continuum
in $B$. By C.28 it contains neither a Jordan curve
nor a triod, and hence is a point or a compact arc;
on such an arc the intrinsic and ambient topologies
agree. These disks give neighborhoods at every point
of $O$. For a circle the two topologies already agree.
Each $l\in L$ restricts to an intrinsic homeomorphism
of $B$, because compact arcs are carried to compact
arcs. For $b=gz$ one has
\[
                  lb=\Phi_z(lgx).
                                                        \tag{C.42.2}
\]
This makes every point evaluation of the restriction
$L\to\operatorname{Homeo}(B)$ continuous.
Finite ordered partitions of compact intervals show
that pointwise convergence of monotone homeomorphisms
to a homeomorphism is uniform on compact intervals;
apply the same argument to inverses. Finite cyclic
partitions give the circle assertion. Thus restriction
is continuous in the compact-open group topology.
All these restrictions preserve orientation, since
the original loop does and the restriction map is
continuous.

Apply the Polish orbit-map theorem to this transitive
action on the Polish one-manifold $B$. It identifies
$B$ with $L/L_z$ in its stated topology. Since
$J\subseteq L_z$, the map $\Phi_z$ is the open quotient
$lJ\mapsto lL_z$.
In particular, if operators $h_n$ commute with $G$
and $h_nx\to z$, then they commute with $L$ and
(C.42.1) gives $h_n|_O\to\Phi_z$ uniformly on compact
subsets of $O$.

\emph{Both covering maps and the compatible group.}
The universal cover of the connected open planar
surface $O$ is a plane. Fix a lift
\[
 F:\widetilde O\longrightarrow\R,\qquad
                      qF=\Phi_zp.
                                                        \tag{C.42.3}
\]
For a line $B$, take $q$ to be an identification.
Local covering charts and openness of $\Phi_zp$
show that $F$ is open.
Consider the following subgroup, with the topology
induced from the indicated product:
\[
\begin{split}
\mathcal L=\{(l,a,b)\in
 L\times\operatorname{Homeo}(\widetilde O)
       \times\operatorname{Homeo}^+(\R):\quad
 pa&=(l|_O)p,\\
 qb&=(l|_B)q,\qquad Fa=bF\}.
\end{split}
                                                        \tag{C.42.4}
\]
The three equations are closed conditions, as may
be checked by evaluation on countable dense sets.
Hence $\mathcal L$ is a Polish group, and its
actions by $a$ and by $b$ are continuous.
Every lift $a$ of $l|_O$ has a unique compatible
lift $b$ of $l|_B$: the maps $Fa$ and $bF$ lift
the same map to $B$, and the target deck translation
is chosen to make them agree at one point.
Uniqueness of lifts then gives equality everywhere.
For a line target there is no target deck ambiguity.

The action of $\mathcal L$ on $\widetilde O$ is
transitive. Choose $l$ carrying the projections of
two specified points to one another and adjust a
source lift by a deck transformation to carry the
specified points themselves.
The image of $F$ is a nonempty open interval $I$.
It is invariant under every compatible $b$.
If $B$ is a circle and $I$ has a finite endpoint,
every increasing $b$ fixes that endpoint, so its
image under $q$ would be fixed by every element
of the transitive $L$-action on $B$. This is
impossible. Therefore $I=\R$; for a line target
surjectivity is already immediate from (C.42.3).

\emph{A component that separates the plane.}
Fix $t\in\R$ and write $E=F^{-1}(t)$.
Choose $u,v$ with $F(u)<t<F(v)$. The closed set
$E$ separates $u$ from $v$. Frankel's planar
separation lemma \cite[Lemma 2.7, p.~1258]{C_Frankel2015}
gives a connected component $C$ of $E$ separating
them. The lemma applies to unbounded closed sets.
For completeness, take the complementary domain
$U$ of $E$ containing $u$, and then the component
$V$ of $\widetilde O\setminus\overline U$ containing
$v$. The planar Brouwer property makes
$\partial V$ connected. It lies in $E$ and
separates $u,v$, so the component of $E$ containing
$\partial V$ has the required property.

Let $\mathcal K$ be the closed subgroup of
$\mathcal L$ whose $b$-coordinate fixes $t$.
It acts transitively on $E$: transitivity on the
covering plane finds an element carrying any
two points of $E$ to one another, and
$Fa=bF$ forces $b(t)=t$.
Thus all components of $E$ are ambiently
homeomorphic to $C$, and all separate the plane.
The stabilizer $\mathcal K_C$ is a closed Polish
subgroup and acts transitively on $C$.
Closedness follows by applying forward and inverse
limits to the closed set $C$.

\emph{Accessibility from each fixed complementary domain.}
Let $D$ be any component of
$\widetilde O\setminus C$.
Its stabilizer in $\mathcal K_C$ is open:
for a fixed $d\in D$, every sufficiently small
group element sends $d$ into $D$ and therefore
sends the whole complementary domain onto $D$.
The open orbit-map theorem for
$\mathcal K_C\curvearrowright C$ makes every
orbit of this open stabilizer open in $C$.
Since $C$ is connected, the stabilizer acts
transitively on $C$.

There is an accessible point of $\partial D$:
an interior point and a closest boundary point
give a straight end-cut in a coordinate copy of
the plane. Its boundary endpoint belongs to $C$.
Transporting this end-cut by the domain
stabilizer proves that every point of $C$
is accessible from this same $D$.
This holds for every complementary domain.

There are exactly two such domains.
There are at least two because $C$ separates.
For three domains, choose three distinct points
$c_1,c_2,c_3\in C$. Inside each domain, join
end-cuts to these points to an embedded tree
with these three boundary leaves and one
interior branch point. Such a tree is obtained
by joining two end-cuts to an arc and stopping
a third connecting arc at its first intersection.
The three trees have disjoint interiors and
share just their specified boundary leaves.
Their union embeds $K_{3,3}$ in the plane,
contradicting the bipartite planar Euler
inequality $9\leq2\cdot6-4$.

\emph{The proper line conclusion.}
If $C$ is unbounded, put
$C^*=C\cup\{\infty\}$ in the sphere; if $C$
is bounded, put $C^*=C$.
This is a compact connected proper subset
with exactly two complementary domains.
Every point except possibly the added
point $\infty$ is accessible from both.
Choose a stereographic pole in one complementary
domain. Deleting that pole leaves that domain
connected, and accessibility at its boundary
is unaffected, since end-cuts can be chosen
to avoid the pole. In the resulting plane
both inaccessible subsets of $C^*$ have
at most one point.
Polulyakh's converse Jordan theorem
\cite[Theorem 1.1 and Corollary 1.1]{C_Polulyakh2000}
therefore makes $C^*$ a Jordan curve.
The bounded alternative is impossible:
the real-valued open map $F$ is constant
on that Jordan curve, and its restriction
to the bounded disk either has an interior
extremum or is constant on a nonempty
open set. Both contradict openness.
Thus $C$ is a Jordan curve with the
compactification point removed, a properly
embedded line. Transitivity of $\mathcal K$
transfers this conclusion to every component
of every fiber.

Downstairs, a fiber of $\Phi_z$ is homogeneous
under $L_z$, and each of its component
stabilizers is transitive by the same
orbit argument. A nondegenerate compact
component is therefore homogeneous by
ambient neighborhood homeomorphisms.
Brechner--Erkama
\cite[Theorem 2]{C_BrechnerErkama1979}
makes it a Jordan curve.

\emph{An exact disconnected-fiber example.}
On $r>1$ in the plane put $s=\log(r-1)$.
Let $\beta_b$ be the time-$b$ flow of
$\sin\delta\,\partial_\delta$ on $\R$.
It fixes $\pi\mathbb Z$, is $2\pi$-periodic
as a lifted circle map, and
$|\beta_b(\delta)-\delta|\leq\pi$.
Define
\[
\begin{aligned}
 A_a(r,\theta)&=(r,\theta+a) &&(r\leq1),\\
 A_a(r,\theta)&=(1+e^a(r-1),\theta+a)&&(r>1),\\
 B_b|_{\{r\leq1\}}&=\operatorname{id},\\
 B_b(s,\theta)&=(\theta+\beta_b(s-\theta),\theta)&&(r>1).
\end{aligned}
                                                        \tag{C.42.5}
\]
The last first coordinate is $s$.
Its expression is independent of the chosen
representative of $\theta$.
The bound on $\beta_b-\operatorname{id}$
gives continuity at $r=1$ and properness at
infinity for $B_b$ and $B_{-b}$.
The $A$-formulas have these properties directly.
They are continuous orientation-preserving
one-parameter actions and commute, since
$A_a$ preserves $s-\theta$.

The open spiral strip
\[
 (t,\delta)\longmapsto
  (1+e^{t+\delta},\,t\bmod2\pi),
 \qquad t\in\R,\quad0<\delta<\pi,
                                                        \tag{C.42.6}
\]
is an open orbit $O$. This parametrization is
injective because an angular ambiguity is
$2\pi k$, while two $\delta$-values differ
by less than $\pi$.
In these coordinates $A_a$ translates $t$
and $B_b$ sends $\delta$ to $\beta_b(\delta)$.
The resulting group is $\R^2$, simply
transitive on $O$ and closed in the plane
homeomorphism group. Indeed, the open-orbit
preservation argument at the start of the proof
shows that any homeomorphic limit preserves
$O$, using both the forward and inverse limits.
Its evaluation at $(0,\pi/2)$ therefore remains
inside $O$ and recovers finite parameters
$a=t$ and $b=\log\tan(\delta/2)$ continuously.
The parameter limits recover the limiting
plane homeomorphism, proving closedness.
A triangular parameter loop
$0\to(1,0)\to(0,1)\to0$ generates this
group and contracts by scaling both parameters.

For the actual commuting operator subgroup
$H=\langle A_{-2\pi}\rangle$, its powers
$h_n=A_{-2\pi n}$ converge on compact
subsets of $O$ to
\[
 \Phi(t,\delta)=(1,t\bmod2\pi).
                                                        \tag{C.42.7}
\]
This is exactly the quotient associated to
the limit of $h_n(0,\pi/2)$.
Its fiber over $\theta_0$ is the disjoint
union of the lines $t=\theta_0+2\pi k$,
$k\in\mathbb Z$. Its real lift is $t$.
All maps and the ambient contraction have
been specified, so this example concerns
disconnected whole fibers rather than a
nonzero centralizer fundamental-group class.
The executed symbolic certificate
\path{code/part_b/c42_spiral_fibre_certificate.py}
checks the flow law and differential equation
in the exact half-angle coordinate,
the $A$--$B$ commutation identities, the
integer-labelled whole fibers and their real
lift, and the based triangular parameter loop.
It uses no numerical approximations; the
topological conclusions are proved above.
\end{proof}

\paragraph{Scope of the reduction.}
C.42 proves the component statement for the
specified real-valued lift. It does not
identify a lifted component with a whole
downstairs connected component without
checking the actual covering restriction.
It does not yet give local product charts
for the component decomposition, simultaneous
charts for several boundary quotients, or
the frontier-compatible contraction needed
in C.12. Those global extension questions
remain open.

\begin{claim}{C.43}{Proved}{The lifted boundary quotient is a product projection}
For the actual quotient and covering maps in C.42, every
fiber of the real lift
$F:\widetilde O\to\R$ is connected.
There is a homeomorphism
\[
 T:\R\times\R\longrightarrow\widetilde O,
 \qquad F(T(u,t))=t.                                    \tag{C.43.1}
\]
In particular the components in C.42 form a foliation
with genuine local product charts. The original map
$\Phi_z:O\to B$ also has local product charts with its
specified target coordinate. The construction treats
one quotient map at a time; its auxiliary flow and
the homeomorphism $T$ carry no asserted equivariance
under the original operator family.
\end{claim}
\passage{Section 23, folios 140-141, scans 279-281; Part C(b) infinite-operator request; actual boundary-quotient maps in the remaining relative centralizer extension}


\begin{proof}
Write $X=\widetilde O$ and let $\mathcal L$ be the
closed Polish group (C.42.4), acting transitively on
the covering plane $X$. For $v\in X$, let $C_v$ be
the connected component of $F^{-1}(F(v))$ containing
$v$. By C.42 these sets partition $X$ into properly
embedded lines. The equation $Fa=bF$ for an element
$(l,a,b)\in\mathcal L$ implies
\[
                         a(C_v)=C_{a(v)}.                \tag{C.43.2}
\]
Thus this is the actual component partition preserved
by the given compatible action.

\emph{Regularity and local products.}
Let $I\subset C_v$ be a compact arc containing $v$,
and let $\epsilon>0$. In a fixed metric on the
covering plane, the condition
\[
       \sup_{w\in I}d(a(w),w)<\epsilon
\]
defines an identity neighborhood in $\mathcal L$.
The Polish open orbit-map theorem makes its image
at $v$ a neighborhood of $v$. Consequently, for
every sufficiently near $v'$, some such $a$ sends
$v$ to $v'$. Restricting $a$ to $I$ gives an
$\epsilon$-close embedding of $I$ into $C_{v'}$
carrying $v$ to $v'$. This verifies precisely the
regularity condition of Aarts--Oversteegen
\cite[Definition 5.3]{C_AartsOversteegen1991}.
Their local-section and flowbox theorems
\cite[Theorems 5.6--5.7]{C_AartsOversteegen1991}
therefore provide open product neighborhoods
$S\times(-1,1)\to X$, with each vertical plaque
contained in one $C_v$.

The transverse factor is an interval locally.
Indeed it is metrizable, and its product with
an interval is a surface without boundary.
The elementary product-cancellation argument
proved in C.33 applies: local path connectedness
descends to $S$, embedded open arcs in $S$ are
open by invariance of domain, and an endpoint
would have zero local second homology in the
product. Thus $S$ is locally a one-manifold
without boundary. Shrink to an interval and
write one of these charts as
\[
 \chi:J\times(-1,1)\longrightarrow V\subset X.
\]
Because every plaque lies in a fiber,
$F\chi(s,u)=f(s)$ for a continuous function $f$.
For each open $W\subset J$, the set
$\chi(W\times(-1,1))$ is open in $X$, so the
openness of $F$ proves that $f(W)$ is open.
A continuous open map between intervals is
strictly monotone: failure of strict monotonicity
produces a local maximum, local minimum, or a
constant interval, each contradicting openness.
It follows that $f$ is a homeomorphism onto
an open interval. Replacing the transverse
coordinate by $f(s)$ gives actual charts
in which $F$ is the projection.

\emph{The component quotient and its topology.}
Let $Y$ be the set of components $C_v$, equipped
with the quotient topology, and let
$\pi:X\to Y$ be the quotient map.
This map is open. To verify this directly,
let $U\subset X$ be open and let $v$ belong
to the saturation of $U$. Choose
$w\in C_v\cap U$ and the compact arc in $C_v$
joining $v$ to $w$. Apply the regularity
argument above with an error small enough
that the image of $w$ remains in $U$.
For every $v'$ in a neighborhood of $v$
the resulting image arc joins $v'$ to a
point of $U$ in $C_{v'}$. The saturation
of $U$ is therefore open, which proves
the assertion about $\pi$.

In a chart where $F$ is projection, distinct
plaques have distinct $F$-values. They therefore
belong to distinct global components.
The transversal maps homeomorphically onto
$\pi(V)$: its inverse is the restriction
of the continuous factor map $\overline F$
in
\[
 X\xrightarrow{\ \pi\ }Y
   \xrightarrow{\ \overline F\ }\R,
 \qquad F=\overline F\pi.                               \tag{C.43.3}
\]
Here $\pi(V)$ is open by the preceding
paragraph, and the forward inverse in
the chosen target coordinate is the
continuous transversal section.
Thus $Y$ is a one-manifold without boundary
in the sense that it is locally an interval;
no Hausdorff property has been assumed.
The same charts show that $\overline F$
is a local homeomorphism.

The quotient $Y$ is connected and second
countable: an open quotient carries a
countable base to a countable base.
Every element of $\mathcal L$ induces a
homeomorphism of $Y$ by (C.43.2), and
transitivity on $X$ makes $Y$ homogeneous.
Baillif--Gabard's theorem
\cite[Theorem 1.1, proved in Section 4]
{C_BaillifGabard2008} now proves that
$Y$ is Hausdorff.
The topological hypotheses here are exact:
second countability gives Lindel\"ofness,
and the actual quotient charts give the
local Euclidean structure required in
that theorem.

One can also see the separation mechanism
directly from their proof. In a locally
Hausdorff space, Zorn's lemma supplies
a maximal Hausdorff open set through
each point; maximality makes it dense.
A countable subcover exists by
Lindel\"ofness. A manifold is locally
Baire and hence Baire, so those dense
open sets have a common point. That
point is separated from every other
point inside one member of the subcover.
Homogeneity transports this property
to every point.

A connected second-countable Hausdorff
one-manifold without boundary is a line
or a circle. The circle alternative is
impossible because $\overline F$ is open
and has image $\R$ by C.42. Thus $Y$
is a line. Its local homeomorphism
$\overline F:Y\to\R$ is globally
strictly monotone: the local direction
cannot change on a connected line.
It is onto, hence a homeomorphism.
Each fiber of $F$ is consequently one
of the proper lines $C_v$.

\emph{An exact global product for this map.}
Orient each leaf using the orientation
of the covering plane together with
the increasing transverse coordinate
$F$. On an overlap of two
orientation-preserving charts $(u,F)$,
the leaf-coordinate transition is
increasing. This supplies the coherent
orientation required for the
Aarts--Oversteegen global-flow theorem
\cite[Main Theorem 0.5, its proof in
Section 4, and Remark 4.4]
{C_AartsOversteegen1991}.
It gives a continuous flow
\[
 R:\R\times X\longrightarrow X
\]
whose orbits are exactly the specified
leaves. The flow is free, since those
leaves are lines rather than circles,
and a nontrivial period would make
the corresponding orbit a circle.

Take any small transversal section
$\sigma:I\to X$ with $F\sigma(t)=t$.
The map
\[
 \Theta_\sigma:\R\times I\longrightarrow F^{-1}(I),
 \qquad (u,t)\longmapsto R(u,\sigma(t))                 \tag{C.43.4}
\]
is continuous and bijective: every
fiber is exactly one flow orbit, and
freeness gives uniqueness of $u$.
Both spaces are surfaces, so invariance
of domain makes (C.43.4) a homeomorphism.
These are actual local trivializations
of a principal $\R$-bundle over $\R$.

For completeness, its global section
can be constructed with the group
coordinate retained. Choose a locally
finite interval cover $(I_i)$ with
sections $\sigma_i$ and a subordinate
partition of unity $(\rho_i)$ whose
supports lie in $I_i$. On overlaps,
let the continuous real transition
functions be defined by
\[
             \sigma_j(t)=R(a_{ij}(t),\sigma_i(t)).
\]
They obey $a_{ij}=a_{ik}+a_{kj}$.
Set $b_i(t)=\sum_j\rho_j(t)a_{ij}(t)$
on $I_i$, extending the individual
terms by zero outside their supports.
The sum is locally finite and
$b_i=a_{ik}+b_k$ on overlaps.
Consequently
$\sigma(t)=R(b_i(t),\sigma_i(t))$
is a well-defined continuous global
section with $F\sigma(t)=t$.
The same map as (C.43.4), now over
all $\R$, gives (C.43.1).

Finally, the identities
$qF=\Phi_zp$ from C.42, together with
local charts for the two covering
maps $p$ and $q$, carry the projection
charts to the original map $\Phi_z$.
This proves its stated local product
property with the actual target
coordinate.
\end{proof}

\paragraph{Scope of the reduction.}
The auxiliary flow in the proof is used
to trivialize one specified quotient.
It has not been shown to commute with
the original operators, preserve another
boundary quotient, or extend across the
frontier of $O$. C.43 therefore strengthens
the structure of the necessary boundary
maps without completing the simultaneous
relative extension in C.12.

\begin{claim}{C.44}{Proved}{Exact downstairs forms of a nontrivial boundary quotient}
For the actual quotient $\Phi_z:O\to B$ of C.42,
the open planar orbit $O$ is a plane or an annulus.
Choose a coordinate $B=\R$ when $B$ is a line,
or a coordinate $B=\R/\Z$ when it is a circle,
and use the corresponding target covering in
the definition of $F$. There are source
homeomorphisms giving one of the following
normal forms for the actual map:
\[
\begin{array}{c|c|c|c}
 O & B & \Phi_z & \text{whole fiber}\\ \hline
 \R\times\R & \R & (u,t)\mapsto t & \R\\
 \R\times\R & \R/\Z & (u,t)\mapsto t+\Z
       & \coprod_{\Z}\R\\
 (\R/\Z)\times\R & \R & ([u],t)\mapsto t
       & \R/\Z\\
 (\R/\Z)\times\R & \R/\Z & ([u],t)\mapsto t+\Z
       & \coprod_{\Z}(\R/\Z)\\
 \R\times(\R/m\Z) & \R/\Z & (u,[t]_m)\mapsto[t]_1
       & \coprod_{j=0}^{m-1}\R,\quad m\geq1 .
\end{array}                                             \tag{C.44.1}
\]
In the last row the actual induced image
$\Phi_{z*}(\pi_1O)$ is $m\Z\subset\pi_1B=\Z$;
in all other rows that image is zero.
Every downstairs fiber component is a
properly embedded line or a circle, and
the source covering restricts onto each
such component as the line covering
specified by these formulas.
\end{claim}
\passage{Section 23, folios 140-141, scans 279-281; Part C(b) infinite-operator request; exact quotient and covering constraints in the remaining frontier-extension problem}


\begin{proof}
Let $\Gamma$ be the deck group of
$p:\widetilde O\to O$.
Use C.43 to identify the covering plane
with $\R_u\times\R_t$ so that $F(u,t)=t$.
The equation $qF=\Phi_zp$ implies, for
each $\gamma\in\Gamma$,
\[
 F\gamma=F+\eta(\gamma),                                \tag{C.44.2}
\]
where $\eta(\gamma)=0$ for a line target
and $\eta(\gamma)\in\Z$ for a circle
target. For the circle, the difference
is a continuous integer-valued function
on the connected covering plane and
is therefore constant. Composition
shows that $\eta:\Gamma\to\Z$ is a
homomorphism. Its value is exactly the
winding number of the image under
$\Phi_z$ of the source loop whose lift
ends at $\gamma$ times its initial
point. Thus its image is the actual
induced fundamental-group image.

Write $K=\ker\eta$. Every deck map
preserves the orientation pulled back
from the planar surface $O$. It has
the form
\[
 \gamma(u,t)=
       (h_{\gamma,t}(u),\,t+\eta(\gamma)),                \tag{C.44.3}
\]
and each $h_{\gamma,t}$ is increasing:
the base coordinate is increasing,
and reversing the fiber direction
would reverse the surface orientation.
The functions in (C.44.3) and their
fiberwise inverses are jointly
continuous, by the deck map and its
inverse.

\emph{The exact kernel on a real fiber.}
The action of $K$ on any fiber
$\R\times\{t\}$ is free and proper.
Freeness follows from deck freeness.
For properness, compact subsets of
this fiber are compact in the
covering plane, and the deck action
has finite transporters of compact
pairs. The restriction is faithful,
since a deck transformation fixing
one point is the identity.
An increasing free proper group on
a line is trivial or infinite cyclic.
Here is the direct argument.
Its orbit of a point is closed and
discrete. In the nontrivial case
choose the element taking that point
to its nearest orbit point on the
right. Its iterates tend to the
two ends, because a finite limiting
point would be fixed. Every other
orbit point is one of these iterates:
an intervening point could be
translated back between the original
point and its nearest right neighbor.
Freeness then identifies every
element with the corresponding power.
Consequently
\[
                          K=1\quad\hbox{or}\quad K\cong\Z.
                                                               \tag{C.44.4}
\]

Suppose $K=\langle k\rangle\cong\Z$,
choosing the generator to move one
fiber point to the right.
Since the displacement
$h_{k,t}(u)-u$ is continuous and
nowhere zero on the connected plane,
it is positive everywhere.
Put $k_t=h_{k,t}$.
For $v=n+r$, where $n\in\Z$ and
$0\leq r\leq1$, define
\[
 H(v,t)=\bigl(k_t^n(r\,k_t(0)),\,t\bigr).
                                                               \tag{C.44.5}
\]
At integral $v$ the two formulas
agree. The intervals with endpoints
$k_t^n(0)$ exhaust the line in both
directions, so $H$ is a continuous
bijection of the plane. Invariance
of domain makes it a homeomorphism.
It preserves $F$ and satisfies
\[
                        H(v+1,t)=kH(v,t).
                                                               \tag{C.44.6}
\]
Thus the actual intermediate quotient
is
\[
 \widetilde O/K\cong(\R/\Z)\times\R,
 \qquad F_K([v],t)=t.                                  \tag{C.44.7}
\]

\emph{The compactness obstruction retaining the quotient map.}
The case $K\cong\Z$ cannot coexist
with a nonzero image of $\eta$.
Indeed that image would be $m\Z$
for some $m\geq1$. The quotient
$\Gamma/K$ is then infinite cyclic,
and its chosen generator acts on
(C.44.7) by a circle homeomorphism
in the first coordinate and by
$t\mapsto t+m$ in the second.
The compact strip
$(\R/\Z)\times[0,m]$ maps onto the
full quotient by $\Gamma/K$:
each point can be translated into
that strip using its second coordinate.
It follows that
\[
     O=\widetilde O/\Gamma
       =(\widetilde O/K)/(\Gamma/K)
\]
is compact, contrary to its being
a nonempty open subset of the plane.
Therefore $K\cong\Z$ forces
$\eta=0$ and $\Gamma=K$.
Equations (C.44.5)--(C.44.7) give
the third and fourth rows of
(C.44.1), according to the target.

If $K=1$ and $\eta=0$, then
$\Gamma=1$, so C.43 gives the
first two rows.

\emph{The remaining nonzero covering degree.}
It remains that $K=1$ and
$\eta(\Gamma)=m\Z$, with $m\geq1$.
Then $\Gamma$ itself is infinite
cyclic. Choose its generator
$\gamma$ with $\eta(\gamma)=m$,
and write
\[
                   \gamma(u,t)=(h_t(u),t+m).
                                                               \tag{C.44.8}
\]
We give a coordinate change
preserving the real target
coordinate and removing $h_t$.
For $0\leq t\leq m$ set
\[
 \rho_t(u)=
          \left(1-\frac{t}{m}\right)u
                   +\frac{t}{m}h_0(u).
                                                               \tag{C.44.9}
\]
Each $\rho_t$ is an increasing
homeomorphism of the line, with
$\rho_0=\id$ and $\rho_m=h_0$.
Extend to all real $t$ in both
directions by the exact recursion
\[
                     \rho_{t+m}=h_t\circ\rho_t.
                                                               \tag{C.44.10}
\]
The endpoint identity makes the
extension jointly continuous at
every multiple of $m$.
The map
$H_1(u,t)=(\rho_t(u),t)$ is a
continuous bijection of the plane
and hence a homeomorphism.
Equations (C.44.8)--(C.44.10) give
\[
              H_1^{-1}\gamma H_1(u,t)=(u,t+m).
                                                               \tag{C.44.11}
\]
The quotient is therefore
$\R\times(\R/m\Z)$, and the
original map, still satisfying
$qF=\Phi_zp$, is precisely
$(u,[t]_m)\mapsto[t]_1$.
This is the final row of (C.44.1).

All changes of source coordinates
preserved $F$; the target was
changed only by its initially
specified coordinate.
The whole-fiber descriptions now
follow directly. For a circle
target, levels $F=t_0+j$ are
indexed by $j\in\Z$. The deck
group identifies their labels
by $\eta(\Gamma)$ and acts
within each level by $K$.
Thus the component labels are
$\Z/\eta(\Gamma)$ and each
component is $\R/K$ in the
explicit coordinates above.
For a line target there is
only one real level and the
component is again $\R/K$.
These formulas also prove
surjectivity of the restricted
source covering onto each
downstairs component and give
its deck group. In particular,
they establish the component
assertion without assuming
component-surjectivity for an
arbitrary restricted covering.
\end{proof}

\paragraph{Scope of the reduction.}
C.44 classifies each of the actual
quotient maps separately, including
its whole fibers and covering
labels. Several boundary points
can impose several such maps on
the same open orbit. No common
product coordinates or homotopy
preserving all those maps have
been constructed here. That
simultaneous compatibility and
frontier extension remain part
of C.12.

\begin{claim}{C.45}{Proved}{Two actual boundary quotients give explicit planar coordinates}
Let $\Phi_1:O\to B_1$ and $\Phi_2:O\to B_2$
be two actual nontrivial boundary quotients
of the same open orbit as in C.42.
Use the same source universal covering
$p:X\to O$ and the real target coverings
$q_i:\R\to B_i$, and choose lifts
$q_iF_i=\Phi_i p$.
Exactly one of the following alternatives holds.
Either there is a homeomorphism
$f:\R\to\R$ with
\[
                         F_2=fF_1,                     \tag{C.45.1}
\]
or the pair $(F_1,F_2)$ is a homeomorphism
onto an open domain $\Omega\subset\R^2$
whose horizontal and vertical sections
are intervals and whose two projections
are onto $\R$.
In the second alternative, explicit
independent changes of the two real
coordinates, and a possible exchange
of them, carry $\Omega$ onto one of
\[
 \R^2,\qquad
 \{(x,y):x<y\},\qquad
 \{(x,y):x<y<x+1\}.                                    \tag{C.45.2}
\]
The full increasing coordinate-pair
groups preserving these three domains
are, respectively,
\[
 \operatorname{Homeo}^+(\R)^2,\qquad
 \{(g,g):g\in\operatorname{Homeo}^+(\R)\},\qquad
 \{(g,g):g\in\operatorname{Homeo}^+(\R),\
                         g(x+1)=g(x)+1\}.               \tag{C.45.3}
\]
The actual compatible lifted action
maps into the corresponding group.
These statements retain both original
covering equations and make no
identification of their downstairs
component labels.
\end{claim}
\passage{Section 23, folios 140-141, scans 279-281; Part C(b) infinite-operator request; simultaneous actual boundary-quotient constraints in the remaining extension problem}


\begin{proof}
\emph{A common compatible group.}
Use the closed subgroup
\[
\begin{split}
 \mathcal P=\{(l,a,b_1,b_2)\in
 L\times\operatorname{Homeo}(X)
       \times\operatorname{Homeo}^+(\R)^2:\qquad&\\
 pa&=(l|_O)p,\\
 q_i b_i&=(l|_{B_i})q_i,\qquad
 F_i a=b_iF_i\quad(i=1,2)\}.
\end{split}                                             \tag{C.45.4}
\]
The equations are closed in the
product of the indicated Polish
groups. For every lift $a$ of
$l|_O$, each target lift $b_i$
exists uniquely after its deck
translation is chosen at one point,
exactly as in C.42. Consequently
$\mathcal P$ acts transitively on
the source covering plane.
Both $F_i$ are open onto $\R$
with connected proper-line fibers,
by C.43.

Fix $t\in\R$ and set
$C=F_1^{-1}(t)$.
The closed subgroup of $\mathcal P$
with $b_1(t)=t$ acts transitively
on this whole line. Its action
on the interval $I=F_2(C)$ is
also transitive, by the equality
$F_2a=b_2F_2$. All these target
maps are increasing.
A nondegenerate interval admitting
such a transitive action has no
included endpoint, so it is an
open interval.
The Polish open orbit-map theorem,
applied to the actions on $C$
and $I$, makes
$F_2|_C:C\to I$ the open map
between their homogeneous-space
quotients: the point stabilizer
in $C$ is contained in the
stabilizer of its image in $I$.
It is therefore a strictly
monotone homeomorphism.

If instead $I$ is a point,
transitivity on the source
plane transports this constancy
to every $F_1$-fiber. There is
then a function $f$ satisfying
(C.45.1). Since $F_1$ is an
open quotient, $f$ is continuous.
For every open $W\subset\R$,
\[
              f(W)=F_2(F_1^{-1}(W))
\]
is open. Thus $f$ is open and
onto, hence a homeomorphism
of the line. This is the first
alternative.

In the nonconstant case,
$(F_1,F_2)$ is injective,
since $F_2$ is injective on
each $F_1$-fiber. Invariance
of domain makes it a
homeomorphism onto an open
domain $\Omega$.
Each vertical section is the
open interval $I$ just proved.
Interchanging the two indices
gives the horizontal statement;
the constant alternative in
that direction would contradict
the already proved two-dimensional
open image.
Surjectivity of the two $F_i$
gives both full coordinate
projections.

\emph{The endpoint maps are actual homogeneous quotients.}
The coordinate action
$(b_1,b_2)$ of $\mathcal P$
preserves $\Omega$ and is
transitive on it.
Write its vertical intervals as
\[
 \Omega_x=(\alpha(x),\beta(x)),
 \quad -\infty\leq\alpha(x)<\beta(x)\leq+\infty.
                                                               \tag{C.45.5}
\]
The two endpoints satisfy the
exact equations
\[
 \alpha(b_1x)=b_2(\alpha(x)),\qquad
 \beta(b_1x)=b_2(\beta(x)),                               \tag{C.45.6}
\]
where increasing homeomorphisms
fix the appropriate infinite
endpoints.
Transitivity of the first
coordinate action implies that
$\alpha$ is either everywhere
finite or identically $-\infty$,
and that $\beta$ is either
everywhere finite or identically
$+\infty$.

Each finite endpoint map is a
homeomorphism of the real line.
For example, choose $x_0$ with
$\alpha(x_0)\in\R$. The formula
\[
 \alpha(b_1x_0)=b_2(\alpha(x_0))
\]
descends the continuous second
evaluation through the open
first orbit map
$\mathcal P\to\R$.
It proves continuity of
$\alpha$ in its actual
topology. Its image is
nonempty and invariant under
the transitive second coordinate
action, hence is all $\R$.
The stabilizer of $x_0$ in
the first action fixes
$\alpha(x_0)$ in the second
by (C.45.6). The two open
orbit maps identify $\alpha$
with the associated open
coset-space map.
It follows that $\alpha$ is
continuous, open, and onto,
and therefore strictly monotone
and a homeomorphism.
The same proof applies to
finite $\beta$.
This verifies the endpoint
topology without assuming
regularity of a boundary graph.

\emph{The three domains and the coordinate changes.}
If both endpoints are infinite,
$\Omega=\R^2$.
If precisely one endpoint is
finite, use its inverse as
the second coordinate.
The defining inequality becomes
$x<y$ or $y<x$, according to
which endpoint and monotonicity
occur. In the latter case
exchange the two coordinates.

Suppose both endpoints are
finite. They have the same
monotonicity: an increasing
homeomorphism and a decreasing
homeomorphism cannot satisfy
$\alpha(x)<\beta(x)$ for all
real $x$, by their limits at
the two ends.
When both increase, replace
$y$ by $\alpha^{-1}(y)$ and put
$\tau=\alpha^{-1}\beta$.
When both decrease, replace
$y$ by $\beta^{-1}(y)$ and put
$\tau=\beta^{-1}\alpha$.
In either case the domain is
now exactly
\[
                        x<y<\tau(x),
 \qquad \tau\in\operatorname{Homeo}^+(\R),\
                          \tau(x)>x.                    \tag{C.45.7}
\]
Choose $x_0$ and any increasing
homeomorphism
$h:[x_0,\tau(x_0)]\to[0,1]$.
Extend it by
$h(\tau^n x)=h(x)+n$.
The iterates of $\tau$ exhaust
the two ends, since $\tau$
has no fixed point, and the
endpoint formulas agree.
Thus $h$ is an increasing
homeomorphism of $\R$ with
$h\tau h^{-1}(x)=x+1$.
Applying $h$ to both
coordinates in (C.45.7)
gives the final domain in
(C.45.2).
These are the asserted
explicit source coordinate
maps; a reversal in one
coordinate still conjugates
increasing coordinate actions
to increasing actions.

Finally, an increasing
coordinate pair preserving
$x<y$ must preserve its
boundary graph $y=x$,
which says $g=h$.
For $x<y<x+1$, its lower
and upper endpoint graphs
are separately preserved,
which says
$h(x)=g(x)$ and
$h(x+1)=g(x)+1$.
Conversely these identities
preserve the indicated
domains exactly.
The unrestricted plane case
has no additional equation.
This proves (C.45.3).

All induced actions used
above arise from (C.45.4).
An element acting trivially
on the covering plane may
still act elsewhere under
its $l$-coordinate; no
faithfulness of that global
coordinate has been imposed.
The original equations
$q_iF_i=\Phi_i p$ remain
in force, so the separate
deck shifts and component
labels of C.44 are retained.
\end{proof}

\paragraph{Scope of the reduction.}
C.45 gives the relation between
two specified actual boundary
maps, including all coordinate
changes. A third quotient can
impose an additional equation
on the group in (C.45.3).
Preserving all such equations
and obtaining a global
frontier-compatible contraction
are not consequences of this
pairwise classification.

\begin{claim}{C.46}{Proved}{Two boundary foliations admit a compatible intrinsic contraction}
Use the actual based null-homotopic plane loop $f_t$, its generated group $G$, the
open orbit $O$, and $L,x,J$ of C.42. For every
$z\in\operatorname{Fix}(J)\cap\partial O$ with nontrivial
intrinsic line or circle orbit, lift its quotient $\Phi_z$
to the common source universal covering $p:X\to O$ as in
C.42. Let $\mathscr F$ be the set of fiber partitions of
these real lifts; changing their target coordinates or
their lift constants does not change this set.
For $|\mathscr F|=2$, the loop $f_t|_O$ has a based
compact-open null-homotopy
\[
 K_{s,t}\in\operatorname{Homeo}(O),\qquad
 K_{0,t}=f_t|_O,\quad K_{1,t}=K_{s,0}=K_{s,1}=\id.
                                                        \tag{C.46.1}
\]
It commutes with every plane homeomorphism commuting
with all $f_t$ and preserving $O$. For every indicated
point $z$, it also satisfies the exact quotient relation
\[
                 \Phi_z K_{s,t}=Q^{B_z}_{s,t}\Phi_z,
                  \qquad B_z=Gz,                       \tag{C.46.2}
\]
where $Q^{B_z}$ is the natural median contraction of the
original restricted loop on that intrinsic line or
circle, as in C.27--C.29. Thus all the specified
boundary quotient relations hold simultaneously.
The asserted homotopy is intrinsic to $O$; continuity
at its ambient frontier is a further assertion.
\end{claim}
\passage{Section 23, folios 140-141, scans 279-281; Part C(b), infinite-operator zero-level comparison; simultaneous boundary relations in C12}


\begin{proof}
Choose representatives $F_1,F_2:X\to\R$ of the two
partitions. C.45 supplies a homeomorphism from $X$
onto one of its three domains $D$, with these two
coordinates changed only by independent line
homeomorphisms and a possible exchange. Lift the
isotopy $f_t|_O$ from the identity to $a_t$ on $X$.
For each original boundary target, lift $f_t|_{B_i}$
to its based real path $b_{i,t}$. Covering uniqueness
gives
\[
                       F_i a_t=b_{i,t}F_i.              \tag{C.46.3}
\]
Each $b_{i,t}$ is increasing. For a line target its
endpoint is the identity. For a circle target C.29
identifies its endpoint translation with the winding
of the original ambient-null plane loop and proves
that translation zero. Hence both $b_{i,1}$ are the
identity. Since the pair $(F_1,F_2)$ is injective,
$a_1=\id$. We have obtained an actual based loop in
the corresponding coordinate-pair group (C.45.3),
without discarding a source or target deck shift.

In these coordinates write $a_t=(g_{1,t},g_{2,t})$.
Apply C.27 to each coordinate loop and set
\[
\begin{split}
 m_{i,r}(u)&=\min_{0\leq v\leq r}g_{i,v}(u),\qquad
 M_{i,r}(u)=\max_{0\leq v\leq r}g_{i,v}(u),\\
 R_{i,s,t}(u)&=\operatorname{med}
       \bigl(m_{i,1-s}(u),g_{i,t}(u),M_{i,1-s}(u)\bigr),\\
 \widetilde K_{s,t}(u,v)&=(R_{1,s,t}(u),R_{2,s,t}(v)).
\end{split}                                             \tag{C.46.4}
\]
The coordinate maps and their inverses are jointly
continuous by C.27 and C.31. They preserve the actual
domain $D$. For $D=\R^2$ this follows directly from
their being line homeomorphisms. For $D=\{u<v\}$,
C.45 gives $g_{1,t}=g_{2,t}$, so the two new coordinate
maps are identical and preserve the strict inequality.
For $D=\{u<v<u+1\}$ the same equality holds and each
$g_{i,t}$ commutes with translation by one. Minima,
maxima and the median preserve that equation, so the
new equal coordinate maps preserve both inequalities.
The inverse family has the same property. Formula
(C.46.4) is therefore a based null-homotopy in
$\operatorname{Homeo}(X)$ with the asserted endpoint
identities.

We establish the operator comparison before descending.
Let $h$ commute with all $f_t$ and preserve $O$.
It commutes with $L$ by continuity of composition and
inversion. The stabilizers of $x$ and $hx$ in $L$
are equal: an element fixes one precisely when it
fixes the other, by commutation with $h$ and $h^{-1}$.
Choose $a_h\in L$ with $a_hx=hx$. The exact stabilizer
identity gives $a_hJa_h^{-1}=J$. For every specified
$z$ and every $l\in L$, evaluation gives
\[
 \Phi_z h(lx)=\Phi_z(l a_hx)=l a_hz
                         =\Phi_{a_hz}(lx).             \tag{C.46.5}
\]
The point $a_hz$ lies in $\operatorname{Fix}(J)\cap
\partial O$: normalizing $J$ proves the first property,
and $a_h$ preserves $O$ and its boundary. Its orbit
is the same $Lz=Gz$ as that of $z$, so it has the
same intrinsic type. Thus pullback by $h$ permutes
the full set $\mathscr F$. The inverse $h^{-1}$
gives the inverse permutation. After lifting to $X$,
the two sides of (C.46.5) differ at most by the
specified target deck translation, which leaves
their fiber partitions identical.

Every lift $\widetilde h$ of $h|_O$ consequently
permutes the two coordinate foliations of $D$.
It is a coordinate pair of line homeomorphisms,
possibly followed by exchange of the coordinates.
For clarity, preservation of a fiber partition
defines its map on the corresponding real quotient;
the open quotient projection proves continuity of
that map and of its inverse. Surjectivity of both
coordinate projections then makes it a line
homeomorphism. It may be increasing or decreasing.
These quotient maps give exactly the stated
coordinate form of $\widetilde h$.

Moreover $\widetilde h$ commutes with $a_t$.
Their commutator projects to the identity on $O$
and is a continuous deck-valued function of $t$,
equal to the identity at $t=0$. The discrete deck
group makes it the identity for every $t$.
C.27's naturality under every line homeomorphism,
together with the identical construction in both
coordinates, now gives
\[
      \widetilde h\widetilde K_{s,t}\widetilde h^{-1}
                           =\widetilde K_{s,t}.         \tag{C.46.6}
\]
In particular this holds for every source deck
transformation, which is a lift of $h=\id$.
Consequently the homotopy descends to $K_{s,t}$
on $O$. The evaluation maps and their inverse
maps descend continuously through the open covering
quotient. This proves compact-open continuity,
the endpoint identities, and the asserted
commutation with the original operators.

Finally, take any other specified $z$. Its lifted
fiber partition equals one of the two chosen ones.
The first alternative of C.45 identifies its real
map as a line homeomorphism composed with that
coordinate. Naturality of the median formula under
this exact reparametrization proves
\[
 F_z\widetilde K_{s,t}=\widetilde Q^{B_z}_{s,t}F_z.
\]
For a circle target the formula commutes with its
integer deck translations, so it descends; for a
line target it already is the stated target map.
Using $q_zF_z=\Phi_zp$ proves (C.46.2).
No ambient limit of points approaching $\partial O$
has been inferred from these quotient equations.
\end{proof}

\begin{claim}{C.47}{Proved}{Three independent boundary quotients force a Lie action}
Let $\Phi_i:O\to B_i$, $1\leq i\leq3$, be three actual boundary
quotients of the same open orbit as in C.42. Use their common source
universal covering and real target coverings, with
\[
 p:X\longrightarrow O,\qquad
 q_iF_i=\Phi_i p,\qquad F_i:X\longrightarrow\R .
                                                        \tag{C.47.1}
\]
Suppose the three lifted fiber partitions are pairwise distinct;
equivalently, none of the pairs satisfies the first alternative
of C.45. The full group
\[
 \mathcal W=
 \{a\in\operatorname{Homeo}(X):
       F_i a=b_iF_i\text{ for some }
       b_i\in\operatorname{Homeo}^+(\R),\ 1\leq i\leq3\}
                                                        \tag{C.47.2}
\]
is a closed, locally compact subgroup in the actual compact-open
topology. It is a Lie group of dimension two or three and acts
transitively on $X$. If any pair $(F_i,F_j)$ has image different
from $\R^2$, this action is simply transitive and has dimension two.
In particular, the closure in $\operatorname{Homeo}(X)$ of the
compatible lifted loop group is locally compact.
All three covering equations in (C.47.1) and the kernel of the
global-to-restricted action are retained. The conclusion concerns
the action on the source covering plane; it does not assert an
extension of an intrinsic homotopy across the frontier of $O$.
\end{claim}
\passage{Section 23, folios 140-141, scans 279-281; Part C(b) infinite-operator request; simultaneous actual boundary-quotient constraints in the remaining frontier extension problem}


\begin{proof}
\emph{The actual group topology.}
Include all three target lifts in the single closed Polish group
of tuples $(l,a,b_1,b_2,b_3)$ satisfying
\[
 pa=(l|_O)p,\qquad q_i b_i=(l|_{B_i})q_i,\qquad
 F_i a=b_iF_i \quad(1\leq i\leq3).
                                                        \tag{C.47.3}
\]
For each source lift, the three compatible target lifts are
obtained exactly as in C.42--C.45. This group acts transitively
on $X$, so its image in $\mathcal W$ does also.
Each $b_i$ in (C.47.2) is unique because $F_i$ is onto.
C.43 supplies a continuous section
$\sigma_i:\R\to X$ of each $F_i$. Consequently
\[
              b_i=F_i a\sigma_i,\qquad
              b_i^{-1}=F_i a^{-1}\sigma_i.               \tag{C.47.4}
\]
These formulas are continuous for uniform convergence on compact
sets, for maps and inverses.

They also prove closedness. If $a_n\in\mathcal W$ converges to a
homeomorphism $a$ together with its inverse, (C.47.4) gives locally
uniform limits $b_i$ and $d_i$ of the target maps and their
inverses. The identities $b_{i,n}d_{i,n}=d_{i,n}b_{i,n}=\id$
pass to the limit, since the images of each compact set remain
in a common compact set. Thus $d_i=b_i^{-1}$; each $b_i$ is an
increasing homeomorphism, and $F_i a=b_iF_i$. Hence $\mathcal W$
is a closed Polish subgroup of $\operatorname{Homeo}(X)$.

C.45 identifies $X$ by any two of the $F_i$ with an open domain
having interval coordinate sections. In coordinates $(x,y)=(F_1,F_2)$,
write $F_3=w(x,y)$. Every restriction of $w$ to a horizontal or
vertical section is strictly monotone. The signs are constant
over the connected domain: for two fixed nearby values in one
coordinate, the continuous difference of their $w$-values cannot
change sign. The coordinate representation of $\mathcal W$ is
therefore exactly the triples $(g_1,g_2,k)$ satisfying
\[
 (g_1,g_2)\Omega=\Omega,\qquad
 w(g_1x,g_2y)=k(w(x,y)).                                 \tag{C.47.5}
\]
The map from $\mathcal W$ to these triples and its inverse are
continuous, by (C.47.4) and the pair-coordinate homeomorphism.
All topological group assertions below use this topology.

\emph{A strictly interior mean makes an ordered-pair stabilizer trivial.}
We first prove an elementary fact used in the non-plane cases.
Let
\[
 D=\{(a,b):a<b\}
 \quad\hbox{or}\quad
 D=\{(a,b):a<b<a+1\}.
                                                        \tag{C.47.6}
\]
Suppose a group of increasing line homeomorphisms preserves $D$
diagonally and preserves a continuous function $M:D\to\R$ with
\[
 a<M(a,b)<b,
 \qquad M\text{ injective in each argument separately}.   \tag{C.47.7}
\]
A group element fixing one pair $(a,b)\in D$ is the identity.
Indeed, it fixes all the points obtained by repeatedly inserting
$M$ between consecutive fixed endpoints. These subdivisions
have mesh tending to zero. To verify this, fix $\varepsilon>0$.
On the compact set of pairs
\[
 a\leq u<v\leq b,\qquad v-u\geq\varepsilon,
 \]
both $M(u,v)-u$ and $v-M(u,v)$ have a positive minimum.
Every such pair is still in $D$, including in the band case
because $v\leq b<a+1\leq u+1$.
Every subdivision of an interval of length at least
$\varepsilon$ therefore shortens each child by a fixed positive
amount. No sufficiently deep subdivision can have that length.
Continuity now makes the homeomorphism fix $[a,b]$ pointwise.

Let $I$ be the maximal interval of fixed points containing
$[a,b]$. A finite right endpoint $c$ belongs to the fixed set.
Choose $x<c$ in the interior of $I$ with $(x,c)\in D$.
For $y>c$ sufficiently close to $c$, the pair $(x,y)$ remains in
$D$ and $M(x,y)$ remains in the interior of $I$. Equivariance
gives
\[
 M(x,gy)=M(gx,gy)=gM(x,y)=M(x,y).
 \]
Injectivity in the second argument forces $gy=y$, extending
$I$ past $c$. The analogous argument, with a fixed $y>c$ and
$x<c$ close to a finite left endpoint, uses injectivity in the
first argument. Both finite endpoints are impossible.
Thus $I=\R$ and $g=\id$.

\emph{A non-plane pair and two full remaining pairs.}
Suppose some pair image is not $\R^2$. Relabel and apply the
explicit coordinate changes of C.45 so that
$\Omega=D$ is one of (C.47.6).
The first two target actions then coincide:
$g_1=g_2=g$; in the band case $g(x+1)=g(x)+1$.
First suppose the images of $(F_1,F_3)$ and $(F_2,F_3)$ are both
the full plane. Each $w$-level is then the graph
\[
                   y=r_z(x),\qquad z\in\R,               \tag{C.47.8}
\]
where $(x,z)\mapsto r_z(x)$ is continuous and every $r_z$ is
a homeomorphism of $\R$ onto $\R$.
It is increasing, since $r_z(x)>x$ for all $x$, which a
decreasing real-line homeomorphism cannot satisfy.
Reverse the $z$-coordinate when necessary. Then
$z\mapsto r_z(x)$ is strictly increasing, and its range is
$(x,+\infty)$ in the half-plane case and $(x,x+1)$ in the
band case. The sign is the same for every $x$ by connectedness.
The action equation is exactly
\[
                      g r_z g^{-1}=r_{kz}.               \tag{C.47.9}
\]

For $(a,b)\in D$, there is a unique $s=s(a,b)$ with
\[
                   r_s^2(a)=b.
 \quad\hbox{Define}\quad M(a,b)=r_s(a).                  \tag{C.47.10}
\]
Here are the endpoint and continuity checks. The left side is
continuous and strictly increasing in $s$.
As $s\to-\infty$, $r_s(x)\to x$, and hence $r_s^2(a)\to a$.
This last limit follows by bounding $r_s(a)$ between $a$ and
$a+\delta$ and then using monotonicity of $r_s$ at these two
fixed points. As $s\to+\infty$ the square tends to $+\infty$
in the half-plane case. In the band case $r_s(x)\to x+1$,
and the same fixed-point bracketing gives
$r_s^2(a)\to a+2$. The required value $b$ lies strictly between
the endpoint limits. Strict monotonicity gives uniqueness,
and two nearby $s$-values bracketing a solution continue to
bracket it when $(a,b)$ varies slightly; this proves continuity.

The value in (C.47.10) lies strictly between $a$ and $b$.
It is strictly increasing in $b$ for fixed $a$.
For fixed $b$ and $a_1<a_2$, the corresponding parameters satisfy
$s_1>s_2$, because $r_s^2$ increases in its argument.
If $r_{s_1}(a_1)\geq r_{s_2}(a_2)$, applying these increasing
maps once more and using $s_1>s_2$ gives
\[
 b=r_{s_1}(r_{s_1}(a_1))
   >r_{s_2}(r_{s_2}(a_2))=b,
 \]
a contradiction. Thus $M$ is strictly increasing in its first
argument as well.
Equation (C.47.9) and uniqueness in (C.47.10) prove
$M(ga,gb)=gM(a,b)$.
The preceding mean argument therefore makes every source
point stabilizer trivial: once $g=\id$, surjectivity of $w$
in (C.47.5) also gives $k=\id$.

\emph{A further non-plane pair identifies all three line actions.}
It remains, with $\Omega=D$, to consider the case where at least
one of the other pair images is not the full plane. One of its
finite endpoint homeomorphisms from C.45 gives an exact equation
\[
                         f g=k f.                       \tag{C.47.11}
\]
Replace the third coordinate by $f^{-1}F_3$.
Now all three target actions are the same $g$; this remains true
when $f$ is decreasing. Denote the new third coordinate again
by $w(x,y)$. An equality $w=x$ or $w=y$ at one point would
hold everywhere by transitivity and contradict independence.
The relative order of the three values is therefore constant.

In the half-plane, the projected group of $g$'s acts transitively
on all ordered pairs. Its centralizer in
$\operatorname{Homeo}(\R)$ is trivial: the stabilizer of a point
$x$ acts transitively on each of $(-\infty,x)$ and $(x,+\infty)$,
so its only common fixed point is $x$.
A commuting map must send $x$ to that same fixed point.
It follows that every finite endpoint of either remaining
pair domain, now an equivariant self-homeomorphism of the
same line, is the identity.

When $x<w(x,y)<y$, take $M=w$.
When $w(x,y)>y$, the $(x,w)$ image is exactly
$\{(x,z):x<z\}$: its finite lower endpoint must be the identity,
and a finite upper endpoint would be the same map and leave
an empty interval. For each $a<b$, define $M(a,b)=c$ by
\[
                         w(a,c)=b.                      \tag{C.47.12}
\]
The restriction in $c$ is increasing and maps $(a,+\infty)$
onto $(a,+\infty)$; its increase also follows from
$w(a,c)>c$ and the limit as $c\to+\infty$.
The unique solution satisfies $a<c<b$.
When $w(x,y)<x$, the $(y,w)$ image is exactly
$\{(y,z):z<y\}$; define instead
\[
                         w(c,b)=a.                      \tag{C.47.13}
\]
The restriction in $c$ is increasing and maps
$(-\infty,b)$ onto $(-\infty,b)$, so again $a<c<b$.
Both inverse constructions are continuous by the pair-coordinate
homeomorphisms. They are injective in each argument, by the
strict injectivity of the corresponding restrictions of $w$.
Equation (C.47.5), with all three actions equal, proves their
equivariance.

In the band, $g$ commutes with unit translation.
The equalities $w=x+n$ and $w=y+n$, for any integer $n$, are
excluded by the same transitivity and independence argument.
There is consequently one integer $n$, constant on the connected
domain, such that
\[
                         x<w(x,y)-n<x+1.
 \]
Either $w(x,y)-n<y$ everywhere, or the reverse inequality holds
everywhere. In the first case take $M(x,y)=w(x,y)-n$.
In the second case use the exact domain homeomorphism
\[
          (u,v)=(y,x+1),\qquad
          M(u,v)=w(v-1,u)-n.                            \tag{C.47.14}
\]
It sends the band onto itself, and $u<M(u,v)<v$.
Both constructions are continuous and injective in each argument.
They are equivariant because every $g$ commutes with the
integer shifts. We have obtained a mean satisfying
(C.47.7) in all cases.

The mean argument proves simple transitivity of $\mathcal W$
whenever a pair image is not the full plane.
The Polish open orbit-map theorem makes evaluation
$\mathcal W\to X$ a homeomorphism in this case.
Thus $\mathcal W$ is locally compact and is a topological
two-manifold. Its faithful product action on $X\times\R$
and the Hilbert--Smith theorem in dimension three
\cite[Theorem~1.5 and the introductory reduction]{C_Pardon2013}
make it a two-dimensional Lie group.

\emph{Three full pair images give a global ordered loop.}
Suppose now that all three pair images are $\R^2$.
The function $w:\R^2\to\R$ is then continuous, and each partial
map is a homeomorphism onto $\R$. Reverse input coordinates
as necessary to make it increasing in both arguments.
At a chosen source point $(a,b)$ put $c=w(a,b)$ and use the
explicit coordinates
\[
 u=w(x,b)-c,\qquad v=w(a,y)-c,\qquad z=w(x,y)-c.
                                                        \tag{C.47.15}
\]
They give a continuous multiplication $m:\R^2\to\R$ with
\[
 m(u,0)=u,\qquad m(0,v)=v,                              \tag{C.47.16}
\]
strictly increasing in both variables, each partial map onto
$\R$. Its two division maps are continuous: solving a strictly
increasing continuous equation is continuous by fixed
bracketing on either side of the solution.

Let $S$ be the stabilizer of $(0,0)$ in $\mathcal W$.
Setting either argument equal to zero in its action equation
shows that the three target maps coincide. Thus, in its actual
subspace topology, $S$ is exactly
\[
 \operatorname{Aut}^+(m)=
 \{f\in\operatorname{Homeo}^+(\R):
       f(0)=0,\quad f(m(x,y))=m(fx,fy)\}.                 \tag{C.47.17}
\]
This is a closed subgroup. We next prove its local compactness
by uniform estimates for both maps and inverses.

Fix $a>0$ and $0<A<B$. Set
\[
 \mathcal F_{A,B}=\{f\in\operatorname{Aut}^+(m):
                            A\leq f(a)\leq B\}.
                                                        \tag{C.47.18}
\]
Define $D(x)=m(x,x)$ and $T_c(x)=m(c,x)$.
The map $D$ is an increasing homeomorphism,
$D(x)>x$ for $x>0$, and $D(x)<x$ for $x<0$.
The numbers
\[
 v_n=D^{-n}(a)>0,\qquad
 u_n\text{ defined by }m(u_n,v_n)=0                      \tag{C.47.19}
\]
therefore converge to zero, with $u_n<0$.
Every automorphism respects these formulas, and hence
\[
 D^{-n}(A)\leq f(v_n)\leq D^{-n}(B).
                                                        \tag{C.47.20}
\]
The two bounds tend to zero. If $\ell(t)$ is the unique
solution of $m(\ell(t),t)=0$, then
$f(u_n)=\ell(f(v_n))$; continuity of $\ell$ makes this tend
to zero uniformly over $\mathcal F_{A,B}$ as well.

For $c>0$, $T_c$ is increasing, onto, and satisfies $T_c(x)>x$.
Its integer iterates of zero tend to both ends of $\R$.
For every integer $j$,
\[
                       f(T_a^j(0))=T_{f(a)}^j(0).       \tag{C.47.21}
\]
To bound $f$ on a given compact interval, enclose it between
two such iterates with $j=-N,N$.
The right sides of (C.47.21) range over a compact set as
$f(a)$ ranges over $[A,B]$; the negative iterates depend
continuously on $c$ by the division-map continuity.
Monotonicity then bounds all $f$ uniformly on that interval.

For equicontinuity, let $I$ be a compact interval.
For fixed $n$, both gaps in
\[
                 m(u_n,x)<x<m(v_n,x)\qquad(x\in I)
 \]
have positive minima on $I$. Thus a single $\delta_n>0$
ensures
\[
 |y-x|<\delta_n,\ x\in I
 \quad\Longrightarrow\quad
                   m(u_n,x)<y<m(v_n,x).
 \]
Applying $f$ bounds $f(y)$ between
$m(f(u_n),f(x))$ and $m(f(v_n),f(x))$.
The already proved compact bound on $f(I)$, the uniform limits
$f(u_n),f(v_n)\to0$, and continuity of $m$ show that both
bounds approach $f(x)$ uniformly in $x$ and $f$.
This is equicontinuity on every compact interval.

The inverse family has the same estimates with new fixed
positive bounds. Choose $n$ so large that $D^{-n}(B)<a$.
Then $f(v_n)<a$ and $f^{-1}(a)>v_n>0$.
Choose $N$ with $T_A^N(0)>a$. Since $T_c^N(0)$ increases
with $c>0$, (C.47.21) gives
$f(T_a^N(0))\geq T_A^N(0)>a$, and hence
\[
                     v_n<f^{-1}(a)<T_a^N(0).            \tag{C.47.22}
\]
Applying the same compact bounds and equicontinuity argument
to these positive bounds controls all inverses uniformly.

Arzel\`a--Ascoli on a countable compact exhaustion now gives,
from any sequence in $\mathcal F_{A,B}$, locally uniform
limits $f$ and $h$ of a common subsequence of maps and inverses.
Their compact image bounds permit passage to the composition
identities, so $fh=hf=\id$.
The multiplication equation and the bounds in (C.47.18)
pass to the limit. Thus $\mathcal F_{A,B}$ is compact in the
homeomorphism-group topology.
Taking $A<a<B$ gives a compact neighborhood of the identity
in $S$, and proves its local compactness.

Evaluation $f\mapsto f(a)$ is also injective.
An automorphism fixing $a$ fixes every $v_n$ and therefore every
$T_{v_n}^j(0)$, for $j\in\Z$.
For fixed $x$, choose the first such iterate at or to the right
of $x$. It lies in $[x,T_{v_n}(x)]$, whose length tends to
zero since $v_n\to0$. These fixed points tend to $x$, so
the automorphism fixes every $x$.
The compactness of (C.47.18) proves that this continuous
evaluation is a proper embedding
\[
                              S\lhook\joinrel\longrightarrow(0,\infty).
                                                        \tag{C.47.23}
\]

\emph{From the stabilizer to the full group, without changing topology.}
The locally compact group $S$ acts faithfully and continuously
on $\R^3$ by its action on the first coordinate.
Pardon's theorem makes $S$ a Lie group.
Its embedding (C.47.23) gives $\dim S\leq1$.
The full group $\mathcal W$ is still the closed Polish group
already defined; it has not yet been assumed locally compact.

Consider the free right action of $S$ on $\mathcal W$.
It is a Cartan action in its actual topology.
Indeed, choose an identity neighborhood $U$ in $\mathcal W$
such that $U\cap S$ has compact closure in $S$, using the
proved local compactness and the subspace topology of $S$.
Choose an identity neighborhood $V$ with $V^{-1}V\subset U$.
Then
\[
       \{s\in S:Vs\cap V\ne\varnothing\}
                   \subseteq V^{-1}V\cap S
 \]
has compact closure in $S$, and translates provide the
required neighborhoods at every point.
The space $\mathcal W$ is metrizable and completely regular.
The free-action form of Palais' slice theorem therefore gives
local sections and local trivializations for
\[
                  \mathcal W\longrightarrow\mathcal W/S
 \quad\cite[Theorem~2.3.3]{C_Palais1961}.
 \]
For clarity, the slice map has the form
$\sigma(u)s\mapsto(u,s)$; its inverse is multiplication.
If $\theta$ is the equivariant slice function to $S$, the
inverse coordinates of $w$ are
$(w\theta(w)^{-1},\theta(w))$, as in the explicit free-action
statement \cite[Theorem~3.7]{C_KramerMuratGarcia2026}.
The open orbit-map theorem identifies
$\mathcal W/S$ homeomorphically with $X=\R^2$.
These actual local products of the locally compact spaces
$S$ and an open subset of $X$ prove that $\mathcal W$ is
locally compact. Its faithful action on $X\times\R$ then
makes it a Lie group by Pardon again.
The same local products give
$\dim\mathcal W=2+\dim S$, which is two or three.

Finally, every compatible lifted loop value belongs to
$\mathcal W$, and its closure in $\operatorname{Homeo}(X)$
is a closed subgroup of this locally compact group.
It is therefore locally compact with precisely the induced
homeomorphism-group topology.
None of the arguments replaces the $l$-coordinate of
(C.47.3) by its action on $X$: the kernel of that restriction
can still act elsewhere.
The coordinate changes act on the specified real target lifts
and leave the original equations (C.47.1) available throughout.
This proves the asserted intrinsic reduction and leaves the
frontier extension step explicit.
\end{proof}

\begin{claim}{C.48}{Proved}{One boundary foliation admits a compatible intrinsic contraction}
Use the based ambient-null plane loop and the full family
$\mathscr F$ of lifted boundary fiber partitions in C.46.
For an open orbit $O$ with $|\mathscr F|=1$, there is a based
compact-open null-homotopy $K_{s,t}$ of $f_t|_O$ which
commutes with its entire original operator stabilizer and
satisfies, simultaneously for every specified boundary point,
\[
                 \Phi_zK_{s,t}=Q^{B_z}_{s,t}\Phi_z.
                                                        \tag{C.48.1}
\]
Here $Q^{B_z}$ is precisely the natural median homotopy of
C.27--C.29. The original source and target covering maps are
retained. This is an intrinsic assertion on $O$; ambient
continuity at its frontier is an additional requirement.
\end{claim}
\passage{Section 23, folios 140-141, scans 279-281; Part C(b), infinite-operator zero-level comparison and actual boundary quotient compatibility}


\begin{proof}
\emph{The actual lifted operator group.}
Let $p:X\to O$ be the universal covering and choose one
real quotient $F:X\to\R$ representing the unique partition.
By C.43, $X$ is a plane, $F$ has connected proper-line fibers,
and a product chart identifies it with the second projection.
Let $Q$ be the closure, in $\operatorname{Homeo}(O)$, of the
restrictions of the original $O$-stabilizing operators.
C.30 proves that this exact closed group acts freely and
properly. Define its group of all lifts by
\[
 \widetilde Q=
 \{a\in\operatorname{Homeo}(X):pa=qp\text{ for some }q\in Q\}.
                                                        \tag{C.48.2}
\]
It includes the source deck group. Its topology is the
subspace homeomorphism-group topology, with the maps and
their inverses controlled on compact subsets.

We give the needed covering comparisons explicitly.
A lift is determined by its value at one point once its
downstairs map is specified. For a convergent sequence of
downstairs maps and a convergent compatible choice of those
values, the lifts converge on compact subsets: cover a
compact path from that point by finitely many covering
charts and continue the chosen sheet successively; the
same argument on finitely many such neighborhoods controls
a compact set. Apply it to the inverse maps as well.
Conversely a convergent sequence of lifts gives convergence
downstairs on each compact set by finitely many local
sections of $p$. These observations prove that (C.48.2) is
closed and that its projection to $Q$ has the stated
topology of lifts.

The action of $\widetilde Q$ on $X$ is free. A lift fixing
a point has downstairs map fixing a point, hence equal to
the identity; the resulting deck map fixing a point is the
identity too. It is proper in this topology. Indeed, from
$x_n\to x$ and $a_nx_n\to y$, properness downstairs gives
a subsequence $q_n\to q$. There is a unique lift $a$ of $q$
with $ax=y$, and the covering-chart argument gives
$a_n\to a$. The same argument makes evaluation at a point
a homeomorphism onto a closed subset of $X$: convergence
of its values supplies convergence of the corresponding
downstairs maps and then of the selected lifts.
Thus $\widetilde Q$ is locally compact. Its faithful product
action on $X\times\R$ makes it a Lie group by the already
checked Hilbert--Smith theorem of C.33. Its dimension is
at most two, since each orbit is properly embedded in $X$.

The original operators preserve the unique partition by
(C.46.5). This property passes to their restrictions'
closure and to all lifts. More explicitly, choose lifts
of an approximating sequence whose values at one point
converge to the chosen lift. Equality of $F$-values on a
fiber passes to the limit, and the inverse maps give
equality of the resulting fiber partitions. Since $F$ is
an open quotient, every $a\in\widetilde Q$ therefore has a
unique induced line homeomorphism $\bar a$ with
\[
                         Fa=\bar aF.                   \tag{C.48.3}
\]
A continuous section of $F$ gives $\bar a=Fa\sigma$, and
the inverse formula uses $a^{-1}$; hence this induced
homomorphism is continuous in the actual topologies.

Lift the original restricted isotopy from the identity
to $a_t$ on $X$. C.32 proves its null evaluation loop in
$O$, so $a_1=\id$. Every element of (C.48.2) commutes with
$a_t$: its commutator is a continuous deck-valued function
of $t$, initially the identity. Write
\[
                         Fa_t=b_tF.                    \tag{C.48.4}
\]
These are the based target lifts, and $b_0=b_1=\id$ by
C.29 for a circle target and directly for a line target.
The group generated by the paths $b_t$ acts transitively
on $\R$. Its orbits are intervals, because every word is
joined to the identity by a path of words. A finite
endpoint of a proper orbit interval would be fixed by
every word. Such a common fixed point would project to a
fixed point of the nontrivial transitive boundary orbit,
which is impossible. Thus the orbit interval is all of
$\R$.

\emph{The base operator image and its kernel.}
Take the closed transitive group generated on this real
line. The evaluation proof of C.30 applies to the line
without change and shows that its full centralizer is
locally compact, free and proper. Its product action on
$\R^3$ makes it Lie. No decreasing element is possible:
it would have a fixed point, contrary to freeness.
In dimension one its identity component acts transitively
on the line, and freeness excludes additional components;
it is the real translation group. In dimension zero the
nearest-orbit-point argument of C.44 gives either the
identity group or a cyclic translation group. Choose the
resulting line coordinate. We have thereby proved, with
the comparison (C.48.3), that every base operator is a
translation. Write
\[
 Fa=F+\alpha(a),\qquad
 \alpha:\widetilde Q\longrightarrow(\R,+),\qquad
 K=\ker\alpha,
                                                        \tag{C.48.5}
\]
after this specified change of target coordinate.

Restriction of $K$ to any fiber is faithful, free and
proper. Evaluation gives it the same topology as its
closed orbit in that line. The preceding one-dimensional
classification therefore gives exactly
\[
                         K=1,\quad K\cong\Z,
                         \quad\text{or}\quad K\cong\R.
                                                        \tag{C.48.6}
\]
All coordinate changes below are homeomorphisms of the
source and reparametrizations of this same target.
We construct them and keep the action of every component.

\emph{The real kernel.}
For $K\cong\R$, its orbits are the entire fibers. A global
section of its principal bundle over the line, constructed
as in C.33, gives coordinates $(u,r)$ with $F(u,r)=r$ and
\[
                         k_c(u,r)=(u+c,r).              \tag{C.48.7}
\]
The quotient $\widetilde Q/K$ acts freely and properly on
$X/K=\R$, by the exact local-section argument in C.33.
Its image under $\alpha$ is consequently closed, and is
$0$, $d\Z$ with $d>0$, or $\R$.
For image zero (C.48.7) is the complete operator action.
For image $d\Z$, choose a label-$d$ element. Its action is
\[
                 q(u,r)=(\lambda u+\tau(r),r+d),
                         \lambda\in\R^\times.          \tag{C.48.8}
\]
The scalar is constant because conjugation on the actual
group $K$ is a single continuous group automorphism.
Choose $c$ on $[0,d]$ with
$c(d)=\lambda c(0)+\tau(0)$ and extend it by
$c(r+d)=\lambda c(r)+\tau(r)$. In the coordinates
$u'=u-c(r)$ the complete additional operator is
\[
                         q(u',r)=(\lambda u',r+d).
                                                        \tag{C.48.9}
\]
This retains negative $\lambda$ as well.

For image $\R$, the group $\widetilde Q$ acts transitively
and freely on $X$. Both its kernel and quotient are
connected lines, so it is connected. A Lie-algebra element
mapping to $1$ integrates to a section $s_r$ with
$\alpha(s_r)=r$. Conjugation on $K$ has the form
$s_rk_us_r^{-1}=k_{e^{cr}u}$ for a real constant $c$.
Evaluation gives the exact group coordinates and product
\[
 X\cong\widetilde Q\cong\R^2,\qquad
 (u,r)(v,z)=(u+e^{cr}v,r+z),\qquad F(u,r)=r.
                                                        \tag{C.48.10}
\]
The formula proves directly that these coordinates are
global; $k_us_r$ is the unique element with the specified
base and fiber values. The operator action is left
multiplication, and its full centralizer consists of
right multiplications.

\emph{The cyclic kernel.}
For $K\cong\Z$, the construction (C.44.5) supplies
coordinates with $F(u,r)=r$ and generator
$k(u,r)=(u+1,r)$. The quotient is the actual cylinder
$X/K=(\R/\Z)\times\R$. The action of
$\widetilde Q/K$ on it is free and proper: lift convergent
point sequences by local sections and correct by elements
of $K$, then apply properness upstairs, exactly as in
C.33. The compact circle fibers show that its base image
is closed. Indeed, a sequence of base translations in a
compact interval sends one fixed point into a compact
cylinder strip; properness gives a convergent subsequence
of the quotient elements. Thus again the image is
$0$, $d\Z$, or $\R$.

For image $d\Z$, the additional generator has the form
\[
 q(u,r)=(g_r(u),r+d),\qquad
 g_r(u+1)=g_r(u)+\epsilon,\quad\epsilon\in\{1,-1\}.
                                                        \tag{C.48.11}
\]
Here $g_r$ is increasing for $\epsilon=1$ and decreasing
for $\epsilon=-1$. Choose an increasing unit-translation
equivariant path $\rho_r$ on $0\leq r\leq d$ from
$\rho_0=\id$ to $\rho_d=g_0\circ(\epsilon\id)$.
The convex interpolation of these two increasing maps
is such a path. Extend it in both directions by
\[
                 \rho_{r+d}=g_r\rho_r(\epsilon\id).
                                                        \tag{C.48.12}
\]
The endpoint equation makes this extension continuous.
Every $\rho_r$ remains increasing and commutes with unit
translation. The homeomorphism $(u,r)\mapsto(\rho_r(u),r)$
therefore changes the two generators into exactly
\[
                 k(u,r)=(u+1,r),\qquad
                 q(u,r)=(\epsilon u,r+d).              \tag{C.48.13}
\]

For base image $\R$, the projection with discrete kernel
is a covering homomorphism of Lie groups. Its identity
component maps isomorphically onto the simply connected
line. It centralizes $K$, since its conjugation map to
$\operatorname{Aut}(\Z)=\{1,-1\}$ is continuous and is
initially the identity. Thus
$\widetilde Q=\Z\times\R$. Use this real subgroup to
transport the fiber $F^{-1}(0)$ to every other fiber,
and use the nearest-orbit-point coordinate for $K$ on
that first fiber. These are global coordinates in which
\[
           k(u,r)=(u+1,r),\qquad s_z(u,r)=(u,r+z).
                                                        \tag{C.48.14}
\]
For base image zero only $k$ remains.

\emph{Trivial kernel and the possible dense base image.}
For $K=1$, the group $\widetilde Q$ embeds algebraically
in the translation group and is abelian. A positive
dimensional identity component maps isomorphically onto
$\R$; injectivity excludes other components. The same
transport from $F^{-1}(0)$ gives just $s_z$ in (C.48.14).
In dimension zero, $X\to X/\widetilde Q$ is a regular
surface covering with simply connected source. Its deck
group is the actual fundamental group of the quotient.
A noncompact surface has free fundamental group, and
an abelian free group is trivial or infinite cyclic.
For a compact quotient the surface classification gives
only the torus with abelian torsion-free fundamental
group: the sphere and projective plane do not have plane
universal cover, the Klein-bottle group is nonabelian,
and every remaining closed surface group is nonabelian.
Consequently the discrete possibilities are exactly
$1$, $\Z$, and $\Z^2$.

The cyclic case has a nonzero base translation; choose
its positive generator. The construction
(C.48.11)--(C.48.13), now without the unit-translation
condition, changes it to $(u,r)\mapsto(\epsilon u,r+d)$.
The trivial case has no equations.
For $\widetilde Q\cong\Z^2$ the quotient is a torus,
and its base image is dense: an injective homomorphism
from $\Z^2$ into $\R$ has two generators with irrational
ratio. The foliation descends to an oriented minimal
topological foliation of that torus. To verify minimality,
the saturation of any lifted fiber under the deck group
is $F^{-1}(r+\alpha(\widetilde Q))$, which is dense because
$F$ is open and the displayed set of levels is dense.
The quotient torus is orientable, and its deck maps
preserve both the surface and the transverse orientation,
so the leaf orientation descends as asserted.

We give the elementary foliation conjugacy needed here,
including its maps. It is the minimal-torus conclusion
recorded in \cite[Proposition~2.5]{C_MionMouton2025}; no
bi-foliation classification or asymptotic-cycle statement
is needed. The oriented local flowboxes give a nonsingular
continuous flow with these leaves by the construction
used in C.43. Choose a small closed transverse interval,
with endpoints off the orbit of an interior point, and a
first-return arc from that point to a different point of
the interval. A flowbox excludes sufficiently short
positive returns; minimality supplies a later return,
and compactness gives a first one. The endpoints chosen
off the orbit put that return in the interval's interior.
The arc's interior avoids the interval, and its endpoint
differs from its initial point because there is no
periodic orbit. Its union with the intervening
transverse subinterval is a simple closed curve.
In a thin flow strip along the arc, tilt the arc strictly
in the transverse direction pointing from the return
point back to the initial point. The tilt may be chosen
arbitrarily small, because the two endpoint transverse
coordinates are distinct; choose it so that the new
terminal point remains strictly between the old
endpoints in the return interval. Join the new endpoints
inside the transverse interval. Both pieces now have
the same strict transverse direction; straightening
their two corners in flowboxes gives an embedded closed
transversal $\gamma$.

Every leaf meets $\gamma$, since it is dense and a
flowbox around $\gamma$ is open. Forward and backward
returns exist: on a compact minimal flow each forward
and backward limit set is a nonempty closed invariant
set, hence the entire torus. Flowboxes and avoidance
of $\gamma$ on a compact interior return arc show that
the first return time is continuous and positive.
The return map $h:\gamma\to\gamma$ is a minimal
orientation-preserving circle homeomorphism. The map
which follows each point for a fraction of its return
time gives the leaf-preserving homeomorphism from its
mapping torus
\[
 [0,1]\times\gamma/((1,z)\sim(0,hz))\longrightarrow X/\widetilde Q.
                                                        \tag{C.48.15}
\]
Surjectivity uses the preceding returns; uniqueness of
the last intersection gives injectivity, and the compact
source and Hausdorff target give the inverse continuity.

For completeness, choose a probability measure invariant
under $h$ by taking a weak limit of finite orbit averages.
Minimality gives full support, and the absence of periodic
points gives no atoms. Its cumulative distribution is
therefore a circle homeomorphism. It conjugates $h$ to a
Lebesgue-measure-preserving increasing homeomorphism,
hence to a rotation: preservation of the measure of
every interval makes the displacement constant.
The angle $\theta$ is irrational by minimality.
After this conjugacy, the explicit map
\[
              [s,z]\longmapsto[(s,z+s\theta)]
                         \quad\text{in }\R^2/\Z^2
                                                        \tag{C.48.16}
\]
is well defined at the identified endpoints and sends
the suspended leaves to the lines of slope $\theta$.
This proves the required foliation conjugacy.

Lift (C.48.15)--(C.48.16) to the universal covers. It
intertwines the original deck action with translations
by $\Z^2$, through the induced isomorphism of their
fundamental groups. Each original $F$-fiber becomes a
whole affine line of slope $\theta$. A linear transverse
coordinate $r$ and a complementary linear coordinate
$u$ thus give
\[
       q(u,r)=(u+\beta(q),r+\alpha'(q))\quad(q\in\widetilde Q),
                                                        \tag{C.48.17}
\]
where these translation vectors form a lattice, and
$\alpha'(\widetilde Q)$ is dense. The original function
$F$ is a line homeomorphism composed with $r$: both
are open quotients with exactly the same connected
fibers. This gives the exact target reparametrization,
rather than identifying the two quotient functions
without a comparison map.

\emph{Contraction in the established normal forms.}
Except for the separately displayed transitive case
(C.48.10), write the actual lifted loop as
\[
                 a_t(u,r)=(h_t(u,r),b_t(r)).             \tag{C.48.18}
\]
Each $h_t(\,\cdot\,,r)$ is an increasing line
homeomorphism. Its endpoints are the identity family.
Let $B_{s,t}$ be the natural median contraction of
$b_t$ from C.27, and set
\[
 H_{s,t}(u,r)=
       \bigl((1-s)h_t(u,r)+su,\ B_{s,t}(r)\bigr).
                                                        \tag{C.48.19}
\]
We check every operator equation rather than changing
the operator group. For $k(u,r)=(u+1,r)$ the equation
is $h_t(u+1,r)=h_t(u,r)+1$, which (C.48.19) preserves.
For $q(u,r)=(\epsilon u,r+d)$ it is
\[
 h_t(\epsilon u,r+d)=\epsilon h_t(u,r),\qquad
 b_t(r+d)=b_t(r)+d.
                                                        \tag{C.48.20}
\]
Convex interpolation preserves the first identity and
median naturality preserves the second.
For the real kernel (C.48.7), commutation says
$h_t(u,r)=u+\phi_t(r)$; with (C.48.9) the further
equation is $\phi_t(r+d)=\lambda\phi_t(r)$.
Scaling $\phi_t$ by $1-s$ preserves it, including
negative $\lambda$.
For the real base translations in (C.48.14), the
equations are that $h_t$ is independent of $r$ and
$b_t(r)=r+c_t$. Both are preserved by (C.48.19).

For the lattice (C.48.17), the dense set of base
translations and continuity force $b_t(r)=r+c_t$.
The fiber equation is exactly
\[
 h_t(u+\beta(q),r+\alpha'(q))=h_t(u,r)+\beta(q).
                                                        \tag{C.48.21}
\]
It is preserved by convex interpolation with $u$.
The median of translations is again a translation,
so the base equation is preserved as well.
The trivial operator case has no extra equation.
These cases exhaust (C.48.6).

In (C.48.10), write the loop as right multiplication
by $(v_t,c_t)$. Its action is
$(u,r)\mapsto(u+e^{cr}v_t,r+c_t)$.
The median base contraction is translation by some
$C_{s,t}$. Right multiplication by
\[
                         ((1-s)v_t,C_{s,t})             \tag{C.48.22}
\]
is the required contraction and still commutes with
every left multiplication.

Every map in (C.48.19) is a homeomorphism: its base
map is a line homeomorphism, and its fiber map is
an increasing convex combination of two increasing
onto maps. The maps are jointly continuous. Their
inverses are jointly continuous as well: on compact
parameter sets, monotonicity and two fixed bracketing
fiber values bound the inverse, and continuity and
uniqueness identify every limit. This is the inverse
continuity argument of C.31 applied to the displayed
fiber family and the continuous inverse base family.
Formula (C.48.22) has the same properties directly
from the group law. All endpoint identities are
preserved. We have therefore obtained a based
null-homotopy commuting with all of $\widetilde Q$.

In particular it commutes with the original source
deck group, so it descends to a compact-open
null-homotopy $K_{s,t}$ on $O$, with continuous
inverse family. Its commutation with every original
operator follows from (C.48.2), retaining the
restriction kernel, which acts identically on $O$.
Finally each other boundary real quotient is an
exact reparametrization of $F$ by the first case
of C.45. Median naturality under that line
homeomorphism gives its required quotient identity.
For a circle target it commutes with the target
integer translations and descends. The equations
$q_zF_z=\Phi_zp$ now give (C.48.1) for every $z$.
No limit across the ambient frontier has been
deduced from this intrinsic construction.
\end{proof}

\begin{claim}{C.49}{Proved}{Lie contractions on closures with three boundary foliations}
Let $f_t$ be an ambient-null based loop in
$\operatorname{Homeo}^+(\R^2)$, let $G$ be its generated group,
and let $H$ commute with every $f_t$.
For every open $G$-orbit $O$ having at least three distinct actual
lifted boundary-quotient fiber partitions of C.42, the group
generated by the based lifts of $f_t|_O$ to the universal covering
$X\to O$ is closed in $\operatorname{Homeo}(X)$. With this actual
topology it is one of
\[
 \R^2,\qquad \operatorname{Aff}^+(\R),\qquad
 \mathcal A_2=\R^2\rtimes\R ,
                                                        \tag{C.49.1}
\]
where the last factor acts on $\R^2$ by positive scalar
dilations. Each group is simply connected and has a globally
invertible exponential map.

There is a jointly continuous based contraction on the whole
$\overline O$, with jointly continuous inverse, obtained from
these actual lifted groups by fractional powers.
These contractions preserve every actual boundary-quotient
equation, agree on intersections of such orbit closures and
with the contraction of C.37 on $X_2$, and are transported to
one another by $H$.
They consequently give a continuous homeomorphism family on
every finite union of these closures and $X_2$.
The conclusion does not assert continuity across an arbitrary
infinite accumulating family of the new three-dimensional
cases.
\end{claim}
\passage{Section 23, folios 140-141, scans 279-281; Part C(b) infinite-operator request; actual simultaneous boundary-quotient and frontier-extension obligations}


\begin{proof}
\emph{The actual based source lifts are closed loops and act transitively.}
Choose three partitions and lifts $F_i:X\to\R$ as in C.47.
Lift the isotopy $f_t|_O$ from the identity to
$\widetilde f_t\in\operatorname{Homeo}(X)$.
The covering homotopy property gives continuity for maps and
inverses in their compact-open topology.
Every $\widetilde f_t$ commutes with every deck transformation:
the commutator is a continuous deck-valued function of $t$
and is the identity at $t=0$.
Its target lifts satisfy
\[
                    F_i\widetilde f_t=b_{i,t}F_i .
                                                        \tag{C.49.2}
\]
For a real target, $b_{i,1}=\id$ directly. For a circle target,
the exact two-point winding comparison of C.29 makes the
based real lift return to the identity, since the original
plane loop is ambient-null. Thus $b_{i,1}=\id$ for all $i$.
The injectivity of $(F_1,F_2)$ from C.45 now gives
$\widetilde f_1=\id$.

Let $N_{\mathrm{alg}}=\langle\widetilde f_t:0\leq t\leq1\rangle$.
This actual subgroup is path connected: each finite word is
joined to the identity by simultaneously shortening its
factor paths. Put
\[
 P=\{\widetilde f_t:0\leq t\leq1\}
       \cup\{\widetilde f_t^{-1}:0\leq t\leq1\}.
 \]
It is compact and contains the identity.
The source deck group $\Delta$ is countable, because a fiber
of a covering of a second-countable surface is a discrete
countable subset of its second-countable covering plane.
For $x\in X$, transitivity of $G$ on $O$ gives
\[
                       X=\bigcup_{\delta\in\Delta}
                            \ \bigcup_{n\geq1}\delta P^n x .
                                                        \tag{C.49.3}
\]
All sets in this countable union are compact.
Baire's theorem gives interior to one of them, hence to
$N_{\mathrm{alg}}x$. Transitivity on that orbit then makes
the entire orbit open.
Every other orbit is a deck translate of it, since the
generated lifted group commutes with $\Delta$.
These disjoint open orbits partition the connected plane;
there is only one. Thus $N_{\mathrm{alg}}$ is transitive
on $X$, not just on its projection $O$.

\emph{The connected three-web groups and their exact actions.}
Let $\mathcal W$ be the full group in C.47 and write
$W=\mathcal W^0$. The group $W$ is open in $\mathcal W$.
The open orbit map makes every $W$-orbit open in $X$;
normality of $W$ and connectedness of $X$ imply transitivity.
For a point $x$, let $J=W_x$.
The Lie principal bundle $W\to W/J=X=\R^2$ shows that
$J$ is connected: the exact sequence on components has
$\pi_1(X)=0$ and $\pi_0(W)=0$.

If $\dim W=2$, then $J$ is discrete and connected, hence
trivial. Evaluation identifies $W$ homeomorphically with
the simply connected plane. Its Lie algebra and group are
the two explicit possibilities $\R^2$ and
$\operatorname{Aff}^+(\R)$ in C.36.
Since $N_{\mathrm{alg}}\subset W$ is already transitive
and this action of $W$ is simply transitive, one has
$N_{\mathrm{alg}}=W$.

Suppose $\dim W=3$. Here C.47 puts us in the case where all
three pair images are full planes.
The connected one-dimensional isotropy $J$ is a subgroup
of the ordered-loop automorphism group in that proof.
It cannot be a circle. Indeed, every compact group of
increasing line homeomorphisms is trivial: the compact orbit
of any point has a largest element, which all group elements
must fix, so transitivity on that orbit makes it a singleton.
Thus $J\cong\R$. The principal bundle over the contractible
plane now gives $\pi_1(W)=0$.

We identify its Lie algebra and homogeneous action explicitly.
Let $\mathfrak w=\operatorname{Lie}(W)$ and
$\mathfrak j=\operatorname{Lie}(J)$.
For each $i$, the stabilizer $K_i$ of $F_i(x)$ in the
induced real target action is a closed subgroup containing
$J$. The open orbit map identifies $W/K_i$ with the actual
real target. The same principal-bundle component argument
makes $K_i$ connected, and its Lie algebra $\mathfrak k_i$
has dimension two.
The three $\mathfrak k_i$ are distinct: otherwise the
corresponding connected subgroups, and hence the quotient
fibers through $x$, would coincide, contradicting C.45.

For $0\ne Z\in\mathfrak j$, the induced linear map
$\operatorname{ad}Z$ on $\mathfrak w/\mathfrak j$ preserves
the three distinct lines $\mathfrak k_i/\mathfrak j$.
A real linear map on a two-dimensional space preserving
three distinct lines is scalar; write that scalar $\lambda$.
If $\lambda=0$, then
$[\mathfrak j,\mathfrak w]\subset\mathfrak j$, so
$\mathfrak j$ is an ideal.
Its connected subgroup $J$ is then normal in $W$.
As a normal point stabilizer in a transitive action it fixes
all of $X$, contrary to the faithful inclusion
$W\subset\operatorname{Homeo}(X)$.
Therefore $\lambda\ne0$.

The adjoint map on $\mathfrak w$ has the eigenspace
decomposition
\[
 \mathfrak w=\mathfrak j\oplus V,\qquad
 \dim V=2,\qquad [Z,v]=\lambda v\quad(v\in V).
 \]
For $v,w\in V$, the derivation identity gives
$[Z,[v,w]]=2\lambda[v,w]$.
The value $2\lambda$ is absent from the spectrum
$\{0,\lambda\}$, so $[V,V]=0$.
Replacing $Z$ by $Z/\lambda$ identifies the Lie algebra
with scalar homotheties:
\[
              [Z,v]=v,\qquad [V,V]=0.                  \tag{C.49.4}
\]
Since $W$ is simply connected, the corresponding group
is exactly $\mathcal A_2$, with multiplication
\[
        (a,v)(b,w)=(a+b,\ v+e^a w),
        \qquad a,b\in\R,\quad v,w\in\R^2.               \tag{C.49.5}
\]
Under this isomorphism $J$ is the dilation subgroup
$\{(a,0)\}$. The orbit identification
$W/J\to X$ therefore gives an actual equivariant
homeomorphism of the source plane with the standard action
\[
                         z\longmapsto v+e^a z .
                                                        \tag{C.49.6}
\]

\emph{The generated subgroup is closed, rather than merely analytic.}
Yamabe's theorem makes the path-connected subgroup
$N_{\mathrm{alg}}\subset W$ an analytic Lie subgroup
\cite[pp.~13--14]{C_Yamabe1950}.
This statement initially gives its analytic subgroup
structure; it does not identify that structure with the
induced topology. We now establish the latter point by
classifying the transitive subgroups of (C.49.6).

A transitive analytic subgroup has dimension at least two.
In dimension two its Lie algebra is either the translation
plane $V$, or has a nonzero projection to the dilation
coordinate. In the latter case it is
\[
                    \R(Z+v_0)\oplus L,                 \tag{C.49.7}
\]
where $L\subset V$ is a one-dimensional translation subspace.
On the quotient line $V/L$, the generator $Z+v_0$ has flow
\[
                 \bar z\longmapsto
                       e^t\bar z+(e^t-1)\bar v_0 .
 \]
It fixes $-\bar v_0$, while $L$ acts trivially there.
Thus the analytic subgroup in (C.49.7) preserves the proper
affine line $-v_0+L$ in the source plane, contradicting
transitivity. The remaining two-dimensional subgroup is
exactly all translations $V$.
A three-dimensional analytic subgroup is open and equals
the connected group $\mathcal A_2$.
Hence, in this case,
$N_{\mathrm{alg}}=V$ or $N_{\mathrm{alg}}=\mathcal A_2$.

Both are closed with their actual homeomorphism-group
topology. One evaluation recovers a translation vector.
For a homothety, evaluation at two fixed distinct points
recovers its positive scale and translation continuously;
a limit that is still a homeomorphism has positive scale.
Together with the dimension-two case already treated,
this proves that $N=N_{\mathrm{alg}}$ is closed and is
one of (C.49.1), with its induced topology.
It also proves that choosing a different triple of
partitions does not change $N$: it is the same subgroup
generated by the same based lifted loop.

For the semidirect products in (C.49.1), with vector
dimension one or two, the exponential and its inverse
are given by
\[
 \exp(u,v)=
   \left(u,\frac{e^u-1}{u}\,v\right),
 \qquad \frac{e^u-1}{u}=1\quad(u=0).                    \tag{C.49.8}
\]
The displayed scalar factor is positive everywhere.
For $\R^2$ the exponential is the identity.
This proves the asserted global exponential coordinates
without an unspecified choice of logarithm.

\emph{The exact Polish quotient extending the action to the closure.}
Put $L=\overline G$ in the plane homeomorphism group.
It preserves $O$ by the forward-and-inverse boundary argument
of C.36. Define the closed Polish subgroup
\[
 E_O=\{(l,n)\in L\times N:
                  p n=(l|_O)p\}.                       \tag{C.49.9}
\]
Projection $E_O\to N$ is onto: each $n\in N$ is an actual
finite word in the based lifted loop, so the identical
word in the original plane loop supplies its $l\in G$.
This continuous surjective homomorphism of Polish groups
is open.
Its kernel consists of pairs $(l,1)$ with $l$ fixing $O$
pointwise, and therefore fixing $\overline O$ pointwise.
Consequently
\[
 N\times\overline O\longrightarrow\overline O,\qquad
     (n,z)\longmapsto lz\quad\bigl((l,n)\in E_O\bigr)     \tag{C.49.10}
\]
is well defined and jointly continuous.
Indeed, the product of the open projection with
$\id_{\overline O}$ is an open quotient, and the continuous
$E_O$-action factors through it.
This argument keeps the exact source covering and the
possibly nontrivial global restriction kernel.

Every $n$ has its chosen preimage in $G$, so (C.49.10)
preserves every original $G$-orbit in $\overline O$ and
has on that orbit exactly the image of $G$.
It commutes with the original setwise $H$-stabilizer of
$O$, since those preimages commute with that stabilizer.
All source deck transformations commute with $N$ because
they commute with its generating based lifts, so the
restriction of (C.49.10) to $O$ is also the exact descended
source action.

The required closure homotopy is now
\[
 R^O_{s,t}
   =\exp\bigl((1-s)\log\widetilde f_t\bigr)
           \quad\hbox{acting by (C.49.10)}.              \tag{C.49.11}
\]
The source loop is closed and based, as proved above.
Thus (C.49.11) starts with $f_t|_{\overline O}$, ends
with the identity, and fixes both based edges.
The action and group inversion give joint continuity
of the family and its inverse.

\emph{All actual boundary quotients and common boundary images.}
Consider any nontrivial nonopen $G$-orbit
$B\subset\overline O$, with its intrinsic line or circle
topology from C.28. The action (C.49.10) is transitive
on its underlying set. Its point stabilizer $N_z$ is closed.
The continuous bijective orbit map $N/N_z\to B$ has
one-dimensional domain.
Dimension zero is excluded by connectedness and nontriviality.
If its dimension were at least two, restricting to a small
two-dimensional coordinate disk and applying invariance
of domain in the ambient plane would give interior to
$B$, making this whole transitive orbit open.

The bijection agrees with the intrinsic topology.
A compact interval in the homogeneous one-manifold maps
to an embedded arc. On a non-circular orbit that arc is
the unique connecting arc of C.28; on a circle the same
statement holds in a sufficiently small interval.
Its interior is therefore an intrinsic interval
neighborhood. This proves that the orbit identification
is a homeomorphism, and that the $N$-action is continuous
in the intrinsic topology as well.

Every group in (C.49.1) is metabelian.
Its image on $B$ is path connected and equals the actual
$G|_B$. The exact metabelian image argument of C.37
therefore identifies this common image, with its induced
intrinsic compact-open topology, as
\[
 M_B=
 \begin{cases}
   \R\text{ acting by translations, or }
        \operatorname{Aff}^+(\R),&B\text{ a line},\\
   \R/\Z\text{ acting by rotations},&B\text{ a circle}.
 \end{cases}                                           \tag{C.49.12}
\]
These groups are closed in the intrinsic homeomorphism
group: one evaluation recovers the translation or
rotation, and two recover the affine parameters.
The action supplies a continuous surjective Lie
homomorphism
\[
                            \rho_B:N\to M_B.            \tag{C.49.13}
\]
In particular every actual boundary quotient retains its
original equivariance equation under the closure action.
To verify this directly for
$\Phi_z(lx)=lz$, take the preimage $g\in G$ of $n\in N$.
The equality $\Phi_z(gy)=g\Phi_z(y)$ follows from the
defining quotient formula; replacing $g$ by another
preimage changes neither side because its restriction
kernel fixes $\overline O$. Thus
\[
            \Phi_z(R^O_{s,t}y)
                   = R^B_{s,t}(\Phi_z(y))               \tag{C.49.14}
\]
with the same original line or circle target and its
covering equation.

Since $N$ is simply connected, (C.49.13) has a unique
based lift $\widetilde\rho_B:N\to\widetilde M_B$.
The latter group is $\R$ or $\operatorname{Aff}^+(\R)$,
whose exponential is globally invertible.
The path $\widetilde\rho_B(\widetilde f_t)$ is precisely
the based lift $\widetilde b_t$ of the original
$f_t|_B$. Naturality and injectivity of the exponential
give the exact identity
\[
 \widetilde\rho_B
       \left(\exp\bigl((1-s)\log\widetilde f_t\bigr)\right)
       =\exp\bigl((1-s)\log\widetilde b_t\bigr).          \tag{C.49.15}
\]
This defines $R^B_{s,t}$ in (C.49.14) after projection.
On a point orbit the action is the identity.

\emph{Intersections, operators, and the finite-union conclusion.}
For two orbit closures, their intersection consists of
nonopen $G$-orbits. On each nontrivial one, both images
(C.49.12) are the identical group $G|_B$ with the identical
intrinsic topology. The two based target-cover paths
are the same lift of $f_t|_B$. Formula (C.49.15)
therefore makes the closure contractions equal there.
It is also exactly the boundary formula in C.36--C.37,
so it agrees on every nonopen orbit in the intersection
with $X_2$.
If $O$ itself belongs to $U_2$, evaluation identifies the
simply transitive group of C.36 with $O$ and its universal
cover with $X$. The lifted translations are the unique
based source lifts used here; their transitive generated
group is exactly $N$. The two fractional-power formulas
are consequently identical on $O$ and its closure.
An open orbit not in $U_2$ is disjoint from $X_2$, because
it is an open neighborhood disjoint from $U_2$.
These observations prove agreement on the entire
intersection with $X_2$.

For $h\in H$ with $hO=O'$, lift $h|_O$ to a source-cover
homeomorphism $\widetilde h:X\to X'$.
Covering uniqueness for the two based isotopies gives
\[
              \widetilde h\,\widetilde f_t\,
                 \widetilde h^{-1}=\widetilde f'_t .
                                                        \tag{C.49.16}
\]
Conjugation is consequently a continuous group isomorphism
$N\to N'$. It intertwines their exponential maps and
the unique logarithms in (C.49.8).
It therefore carries (C.49.11) to the corresponding
homotopy on $O'$. Equality extends to the closures.
Changing the lift of $h$ by a deck map does not change
this identity, since deck maps commute with the generated
lifted groups.
The actual boundary-quotient families are transported
at the same time: with basepoint $hx$, their exact
formula is
$\Phi^{O'}_{hz}(lhx)=h\Phi^O_z(lx)$.
Thus the class of open orbits under consideration is
$H$-invariant; no arbitrary selected triple is assumed
to be preserved by the operators.

Every closure and $X_2$ is closed in the plane.
The matching families, and separately their matching
inverse families, paste continuously over a finite union.
Each preserves its member closures and is inverted by
the pasted inverse, so it is a homeomorphism of that union.
These assertions include all original covering and
operator labels. Uniform control for infinitely many
new homothety orbits, and extension to the remaining
one- and two-foliation regions, are not supplied by
this finite-pasting argument.
\end{proof}

\begin{claim}{C.50}{Proved}{Uniform contraction through all the new Lie orbit closures}
In the setting of C.49, let $U_3$ be the union of the open
$G$-orbits having at least three distinct actual lifted
boundary-quotient fiber partitions, and retain $U_2$ from C.37.
Set
\[
                  U_*=U_2\cup U_3,\qquad X_*=\overline{U_*}.
                                                        \tag{C.50.1}
\]
The contractions of C.36 and C.49 extend to a jointly
continuous based contraction through homeomorphisms of $X_*$,
with jointly continuous inverse, commuting with $H$.
This includes arbitrary accumulation of the three-dimensional
homothety orbits.

The uniform algebraic bound is exact. For every continuous
path $k:[0,t]\to\mathcal A_2$ with $k(0)=1$, put
$P=k([0,t])\cup k([0,t])^{-1}$ and $g=k(t)$.
Then
\[
        \{\exp(r\log g):0\leq r\leq1\}
                                      \subseteq P^{1913}.
                                                        \tag{C.50.2}
\]
The identity belongs to $P$, so shorter words can be padded.
For $\R^2$ and $\operatorname{Aff}^+(\R)$ the three-letter
bound of C.37 applies. All word equalities used on an orbit
hold in its specified simply connected source group.
The bound does not identify an additive vector sum with a
noncommutative product: the latter's complete error is
corrected explicitly in the proof.
\end{claim}
\passage{Section 23, folios 140-141, scans 279-281; Part C(b) infinite-operator request; uniform frontier control for the actual three-quotient Lie reductions}


\begin{proof}
\emph{Two additive path identities.}
The planar identity of C.37 represents any fraction of the
endpoint of a based path in $\R^2$ by at most three signed
path values. We also require its three-dimensional version:
for a continuous $\gamma:[0,t]\to\R^3$ from $0$ to $v$,
\[
              rv=\sum_{j=1}^{m}\varepsilon_j\gamma(s_j),
 \qquad m\leq5,\quad
       \varepsilon_j\in\{1,-1\},\quad 0\leq s_j\leq t .
                                                        \tag{C.50.3}
\]
For a polygonal path, apply the signed-measure theorem
\cite[Theorem~4.7 and \S4.3.2]{C_GrandoniRaviSinghZenklusen2014}
to its three coordinate derivatives.
The chosen set is a union of at most two circular intervals.
Its integral is a sum of two path differences, with at most
one additional total vector from an interval crossing the
identified endpoint. This gives five signed path values.
Zero terms $\gamma(0)$ permit padding to a nonempty word
when necessary.
The interval count in the halving proof is the same corrected
count used in C.37: two intervals cut at at most three points
give at most five pieces, so one half uses at most two.
For the whole interval, three cuts give at most four pieces,
so one of the two halves uses at most two intervals; these
are also circular intervals.
Complementation, halving, and compactness give every
fraction. Uniform polygonal approximation with the same
endpoints, followed by compactness of the five parameters
and the finite sign set, proves (C.50.3) for arbitrary
continuous paths.

\emph{A maximizing anchor and its exact translation commutators.}
Use the coordinates
\[
        (a,v)(b,w)=(a+b,v+e^a w),\qquad
              v,w\in\R^2,
 \]
and write $k(s)=(a_s,v_s)$.
Let $A=\max_{0\leq s\leq t}|a_s|$.
When $A=0$, the path lies in the translation plane and
the three-letter bound already proves (C.50.2).
For $A>0$, choose $u$ attaining the maximum and let
$h=k(u)$ or $k(u)^{-1}$, choosing its sign so that its
logarithmic scale is $A$.
Thus $h$ is one letter from $P$, not a power of unbounded
length. Conjugate all coordinates by the translation
which sends its unique fixed point to zero.
This is one group automorphism applied to both sides
of every forthcoming identity; it adds no letters to
any word. In the new coordinates $h=(A,0)$ and still
$|a_s|\leq A$. Put
\[
 D=e^A-1,\qquad
 C_s=[h,k(s)]
       =h k(s)h^{-1}k(s)^{-1}
       =(0,Dv_s).
                                                        \tag{C.50.4}
\]
The path $C_s$ is continuous in the translation plane
and begins at the identity.
Every one of its values and inverses uses at most four
letters of $P$.

For any fixed $s$ and $|\beta|\leq B$, where $B$ is a
positive integer, the translation $C_s^\beta$ uses at
most $12B$ letters of $P$.
Indeed, write it as $B$ copies of
$C_s^{\beta/B}$.
For nonnegative $\beta/B\leq1$, the planar identity
applied to the prefix $C|_{[0,s]}$ gives at most three
signed commutator values. For a negative fraction invert
the resulting word. Substituting (C.50.4) costs at most
four letters per commutator.
No continuous choice of the parameters in these words
is required.

\emph{The large-scale case.}
Suppose $A\geq\log2$, so $D\geq1$.
For $0\leq q\leq1$, the translation $T_{qv_s}$ equals
$C_s^{q/D}$ and therefore uses at most twelve letters.
Writing $D_a=(a,0)$ gives
\[
                    D_{a_s}=T_{v_s}^{-1}k(s),
                                                        \tag{C.50.5}
\]
which uses at most thirteen.
Let $a=a_t$ and $v=v_t$. Define
\[
 \theta(a,r)=
 \begin{cases}
  (e^{ra}-1)/(e^a-1),&a\ne0,\\
  r,&a=0.
 \end{cases}
 \qquad
       g^r=(ra,\theta(a,r)v).                           \tag{C.50.6}
\]
The function $\theta$ is continuous and lies in $[0,1]$.
The intermediate value theorem supplies $s$ with
$a_s=ra$, since the path starts at logarithmic scale zero
and ends at $a$.
Then
$g^r=T_{\theta(a,r)v}D_{a_s}$ uses at most
$12+13=25$ letters.

\emph{The small-scale case and its full noncommutative correction.}
Now suppose $0<A\leq\log2$.
Apply (C.50.3) to
$\gamma(s)=(a_s,v_s)$, obtaining
\[
       ra=\sum_{j=1}^{m}\varepsilon_j a_{s_j},\qquad
       rv=\sum_{j=1}^{m}\varepsilon_j v_{s_j},
                         \qquad m\leq5.                 \tag{C.50.7}
\]
Form the actual product
$w=k(s_1)^{\varepsilon_1}\cdots k(s_m)^{\varepsilon_m}$.
For an inverse letter, its translation is
$-e^{-a_{s_j}}v_{s_j}$.
Consequently, with
\[
 p_j=
 \begin{cases}
  \sum_{\ell<j}\varepsilon_\ell a_{s_\ell},
                                       &\varepsilon_j=1,\\
  \sum_{\ell\leq j}\varepsilon_\ell a_{s_\ell},
                                       &\varepsilon_j=-1,
 \end{cases}
                                                        \tag{C.50.8}
\]
the product is exactly
\[
 \begin{split}
 w&=\left(ra,\sum_{j=1}^{m}
                    \varepsilon_j e^{p_j}v_{s_j}\right),\\
 \operatorname{pr}_v(w)
   &=rv+\sum_{j=1}^{m}
             \varepsilon_j(e^{p_j}-1)v_{s_j}.
 \end{split}                                            \tag{C.50.9}
\]
Every $|p_j|$ is at most $mA$.
The pure translation which changes (C.50.9) to (C.50.6)
by left multiplication is
\[
 c=(\theta(a,r)-r)v
       -\sum_{j=1}^{m}
                  \varepsilon_j(e^{p_j}-1)v_{s_j}.
                                                        \tag{C.50.10}
\]
Using the translation commutators (C.50.4), this is the
exact group equation
\[
 g^r=C_t^{\beta_0}
          \left(\prod_{j=1}^{m}C_{s_j}^{\beta_j}\right)w,
 \quad
 \beta_0=\frac{\theta(a,r)-r}{D},\qquad
 \beta_j=-\varepsilon_j\frac{e^{p_j}-1}{D}.
                                                        \tag{C.50.11}
\]
The powers before $w$ commute because they are translations.

Here are uniform bounds, including $a=0$.
The integral formula
\[
        \theta(a,r)=
          \frac{\int_0^r e^{a q}\,dq}
               {\int_0^1 e^{a q}\,dq}                  \tag{C.50.12}
\]
holds at zero as well as away from it.
For $|a|\leq A$, the inequality
$|e^{aq}-1|\leq D$ holds on $0\leq q\leq1$, and the
denominator is at least $e^{-A}\geq1/2$.
Subtracting $r$ in (C.50.12), the numerator is the
difference of an integral over $[0,r]$ and $r$ times
an integral over $[0,1]$, each with integrand
$e^{aq}-1$. Its absolute value is at most $2D$.
Therefore $|\beta_0|\leq4$.
For every other coefficient,
\[
 \frac{|e^{p_j}-1|}{D}
 \leq\frac{e^{mA}-1}{e^A-1}
   =\sum_{\ell=0}^{m-1}e^{\ell A}
   \leq 2^m-1\leq31.                                  \tag{C.50.13}
\]
The same bound covers negative $p_j$ because
$1-e^{-mA}\leq e^{mA}-1$.
The word construction following (C.50.4) now applies with
$B=4$ and $B=31$. Equation (C.50.11) uses at most
\[
                      5+12(4+5\cdot31)=1913             \tag{C.50.14}
\]
original letters. This proves (C.50.2) in all regimes.
The executable certificate accompanying this claim checks
the inverse-letter exponents in (C.50.8), the complete
product correction, the commutator and conjugacy formulas,
and the integer count by exact symbolic arithmetic.
The measure selection and the analytic coefficient bounds
are proved above, independently of that certificate.

\emph{One compact family of actual word isotopies.}
Set $L_0=1913$.
For the original plane loop, all words just proved belong
to the compact parameter family
\[
 W_{\boldsymbol\varepsilon,\boldsymbol s}(u)
       =\prod_{j=1}^{L_0}f_{u s_j}^{\varepsilon_j},
 \qquad
 \varepsilon_j\in\{1,-1\},\quad
 0\leq s_j\leq t\leq1,\quad0\leq u\leq1.
                                                        \tag{C.50.15}
\]
Identity letters are obtained by $s_j=0$.
Continuity of finite multiplication and inversion in the
plane homeomorphism group gives compactness for maps and
inverses, uniformly in $u$.
On a C.49 orbit, the word equality is in its simply
connected actual group $N$.
On a C.36 orbit, use the based path in its specified
universal covering group and the three-letter C.37 bound,
then pad it. Thus the endpoint of the based lift of
(C.50.15) is exactly the required fractional power,
not just a homeomorphism with an unspecified deck shift.

Use $r=1-s$ for the contraction parameter and write
$R^O_{r,t}$ for these fractional actions on an open
orbit in $U_*$. C.36 and C.49 already prove their
individual closure continuity and finite compatibility.
Their groups are metabelian, so
$[[g_1,g_2],[g_3,g_4]]=1$ on $U_*$ for every
$g_i\in G$, and hence on $X_*$.
Every nontrivial nonopen orbit $B\subset X_*$ therefore
has the exact intrinsic image
$M_B=\R$, $\operatorname{Aff}^+(\R)$, or the circle
rotation group, by the path-connected metabelian
argument of C.37.
Define $R^B_{r,t}$ by its unique line-group logarithm
or its based real rotation lift, and use the identity
on point orbits. The based circle lift ends at zero
by C.29. These are the same boundary prescriptions
as in C.36 and C.49.

\emph{Limits, with the bounded residual in the specified cover.}
Take
\[
 x_i\in O_i\subset U_*,\quad x_i\to x\in X_*,
       \qquad (r_i,t_i)\to(r,t).
                                                        \tag{C.50.16}
\]
Choose representations (C.50.15) for
$R^{O_i}_{r_i,t_i}$ and pass to convergent word parameters.
Let $W$ be the limiting plane isotopy, $w=W(1)$,
and $g=f_t$.
If $x$ belongs to an open orbit of $X_*$, then that
orbit already lies in $U_*$ and eventually all $x_i$
lie in it. Its intrinsic continuity applies.
If $x$ is fixed by $G$, every limiting word fixes it.
For a remaining intrinsic line or circle orbit $B=Gx$,
identities on $O_i$ pass to every point $hx\in B$
by evaluating them at $hx_i$ for fixed $h\in G$.
In particular $w|_B$ commutes with $g|_B$.

For a fixed positive integer $n$, set
$m_i=\lfloor nr_i\rfloor$. In each specified simply
connected group, the residual is exactly
\[
       \widetilde w_i^{\,n}\widetilde g_i^{-m_i}
           =\exp\bigl((nr_i-m_i)\log\widetilde g_i\bigr).
                                                        \tag{C.50.17}
\]
Its exponent fraction is in $[0,1)$, so it has another
representation from the same compact word family.
For this fixed $n$, pass to a subsequence on which
$m_i=m_n$ and the residual word isotopies converge
to $Z_n$. Then
\[
       (w|_B)^n(g|_B)^{-m_n}=z_n,\qquad
       |m_n-nr|\leq1,\qquad z_n=Z_n(1)|_B,              \tag{C.50.18}
\]
and all $z_n$ belong to the same compact family of
intrinsic restricted word endpoints.
For a line, logarithms of this compact family are bounded.
In the affine group commuting nonidentity elements belong
to the same exponential one-parameter subgroup; the
identity case gives the same additive logarithm equation.
Thus
\[
          n\log(w|_B)-m_n\log(g|_B)=\log z_n .
 \]
Division by $n$ proves
$\log(w|_B)=r\log(g|_B)$.

For a circle, let $\theta(u)$ be the based real angle
of $f_u|_B$, and put $M=\max_u|\theta(u)|$.
Let $A_W$ be the real endpoint angle of $W|_B$,
$T=\theta(t)$, and $C_n$ that of $Z_n|_B$.
The fixed word length gives $|C_n|\leq L_0M$.
The real, rather than merely circle-valued, equality is
\[
                         nA_W-m_nT=C_n.                \tag{C.50.19}
\]
To prove it, use the plane word isotopies
$W_i(u)^n f_{u t_i}^{-m_i}Z_i(u)^{-1}$.
Their restrictions lift to loops ending at the identity
by (C.50.17). The specified source groups are simply
connected, so these loops are null-homotopic there.
Evaluate at the two points $x_i,hx_i$, where $hx\ne x$.
Their difference-direction winding is zero.
The resulting loops converge uniformly to the analogous
two-point loop at $x,hx$, which stays off the diagonal;
its winding is also zero.
The exact circle comparison of C.29 identifies that
winding, up to the common orientation sign, with
$nA_W-m_nT-C_n$. This proves (C.50.19).
Consequently
\[
                  |A_W-rT|
                       \leq\frac{L_0M+|T|}{n}
                       \quad(n\geq1),
 \]
so $A_W=rT$.
This also covers an identity endpoint on the circle
whose specified real angle $T$ is a nonzero integer.

\emph{Joint continuity and the remaining scope.}
Every convergent subsequence of the compact word
parameters in (C.50.16) has therefore the prescribed
endpoint action at $x$. This proves convergence to
$R^B_{r,t}(x)$ for the entire sequence.
It gives joint continuity from the dense set $U_*$.
For a general sequence in $X_*$, approximate its
$i$th point by a point of $U_*$ with both point distance
and image distance at that fixed parameter less than
$1/i$. The just-proved dense-set limit supplies those
approximations and proves full joint continuity.

Apply the same proof to $f_t^{-1}$.
It has the same generated group, the same open-orbit
families, and the same boundary partitions.
Its fractional actions are the orbitwise inverses,
so the inverse family is jointly continuous as well.
Every member is an orbitwise bijection inverted by
this family and is therefore a homeomorphism of $X_*$.
The equations
$R_{0,t}=1$, $R_{1,t}=f_t|_{X_*}$, and
$R_{r,0}=R_{r,1}=1$ give the based contraction.
The operator transport statements of C.36 and C.49
make $U_*$ and $X_*$ invariant under $H$ and give
commutation on $U_*$. Continuity extends that equality
to $X_*$.

This settles infinite accumulation for the orbit closures
treated by C.49. The remaining global step is a compatible
extension from $X_*$ through the regions with fewer
actual boundary-quotient foliations and their frontiers.
The particular median choice of C.34 cannot be imposed
there without proving compatibility, as C.35 demonstrates.
The general conclusion in C.12 consequently remains open.
\end{proof}

\begin{claim}{C.51}{Proved}{Cocompact proper operators control the entire frontier}
Let $f_t$ be the based ambient-null plane loop of C.49,
let $G=\langle f_t\rangle$, and let $H$ commute with all
its values. For an open $G$-orbit $O$, put
\[
 H_O=\{h\in H:hO=O\},\qquad
 Q_O=\overline{\{h|_O:h\in H_O\}}^{\operatorname{Homeo}(O)}.
                                                        \tag{C.51.1}
\]
Consider all such $O$ for which $O/Q_O$ is compact and
the full family $\mathscr F_O$ of actual lifted boundary
quotient fiber partitions in C.46 has at most two members.
For every one of these orbits the restricted loop has
a based contraction on its entire closure,
\[
 \overline K^O_{s,t}\in\operatorname{Homeo}(\overline O),
 \qquad
 \overline K^O_{0,t}=f_t|_{\overline O},\qquad
 \overline K^O_{1,t}=\overline K^O_{s,0}
                         =\overline K^O_{s,1}=\id,
                                                        \tag{C.51.2}
\]
jointly continuous with its inverse. It commutes with
$H_O$ and equals exactly the C.34 median contraction
on every nonopen orbit in $\partial O$.
When $\mathscr F_O$ is empty, every point of
$\partial O$ is fixed by $G$.

The contractions can be chosen with exact transport
under $H$. They agree on all common frontiers and with
C.34, and hence paste over the closed nonopen locus
together with any finite collection of these closures.
This assertion does not prescribe their values on the
closed set $X_*$ of C.50: the comparison there with
its exponential prescription remains to be proved.
\end{claim}
\passage{Section 23, folios 140-141, scans 279-281; Part C(b), infinite-operator zero-level comparison and ambient frontier extension}


\begin{proof}
Write $Q=Q_O$ and $L=\overline G$ in the plane
homeomorphism group. C.30 proves that the actual group
$Q$ in (C.51.1) is closed, locally compact, free and
proper on $O$, with the specified compact-open topology.
The restriction kernel of $H_O\to Q$ acts trivially
on $O$ and therefore on its closure.
We retain that map; we do not replace an element of
$Q$ by an unspecified global extension.

\emph{A compact set meeting every operator orbit.}
The quotient map $O\to O/Q$ is open. For each point
of $O$, choose a relatively compact open neighborhood
inside $O$. The images of these neighborhoods cover
the compact quotient. Choose finitely many whose
images still cover it, and let $C\subset O$ be the
union of their closures. Then $C$ is compact and
\[
                              Q C=O.                   \tag{C.51.3}
\]
This argument supplies the exact compact fundamental
set required below; it uses no fundamental-domain
boundary or continuous choice of orbit representatives.

\emph{Every frontier sequence produces an actual quotient.}
Take a sequence $x_n\in O$ converging in the plane to
$z\in\partial O$. By (C.51.3), write
$x_n=q_ny_n$ with $q_n\in Q$ and $y_n\in C$.
Pass to a subsequence with $y_n\to y\in C$.
For the moment choose $h_n\in H_O$ with
\[
                    |h_ny_n-q_ny_n|<1/n.               \tag{C.51.4}
\]
Such a choice follows from the definition of $Q$ as
the closure of the actual restrictions in (C.51.1).
Here and below $|\cdot|$ is the Euclidean distance
in the original ambient plane.

Every $h_n$ commutes with $L$: it commutes with all
generators of $G$, and commutation passes to its
closure in the plane homeomorphism group.
The orbit map
\[
                      L\longrightarrow O,\quad l\mapsto ly
                                                        \tag{C.51.5}
\]
is open by the same Polish orbit-map theorem used in
C.30 and C.42. Thus there are $l_n\to1$ in $L$ with
$l_ny=y_n$.
For clarity, take a decreasing countable identity
neighborhood basis $V_k$ in $L$. Each $V_k y$ is a
neighborhood of $y$, so $y_n\in V_k y$ eventually.
Choose an increasing sequence of such eventual
indices and then $l_n\in V_{k(n)}$ with
$k(n)\to\infty$. This proves the stated selection.
Commutation now gives the exact equality
\[
                       h_ny=l_n^{-1}h_ny_n\longrightarrow z.
                                                        \tag{C.51.6}
\]
Convergence follows from (C.51.4), $x_n\to z$, and
compact-open convergence of $l_n^{-1}$ to the identity
on a compact neighborhood of $z$.

Let $J_y=L_y$. For $j\in J_y$,
$j h_ny=h_nj y=h_ny$, so (C.51.6) gives $jz=z$.
Consequently $z\in\operatorname{Fix}(J_y)$, and
C.42's actual evaluation comparison gives
\[
 h_n|_O\longrightarrow\Phi_z^y
       \quad\hbox{uniformly on compact subsets of }O,
 \qquad
 \Phi_z^y:O\longrightarrow\R^2,\quad \Phi_z^y(ly)=lz.
                                                        \tag{C.51.7}
\]
This uses the homeomorphism
$\operatorname{Map}_L(O,\R^2)\to\operatorname{Fix}(J_y)$
of (C.42.1), including its compact-open topology.
In particular its assertion applies when $Gz$ is
a single point, in which case $\Phi_z^y$ is constant.
There is no convergence claim for arbitrary proper
operator maps outside their actual global approximants.

We compare (C.51.7) with the quotient family defined
at the fixed basepoint $x$ in C.46. Choose $a\in L$
with $ax=y$, and put $z_0=a^{-1}z$.
Then $z_0\in\operatorname{Fix}(L_x)\cap\partial O$.
For every $l\in L$ the equality
\[
             \Phi_{z_0}^x(ly)
                =\Phi_{z_0}^x(lax)
                =la z_0=lz=\Phi_z^y(ly)                \tag{C.51.8}
\]
identifies the actual maps, with the same target orbit.
The orbit is contained in $\partial O$, since $L$
preserves $O$ and its boundary. It is nonopen; an
open orbit meeting $\partial O$ would be an open
neighborhood disjoint from $O$, contradicting the
definition of its closure. Thus C.28 makes a
nontrivial such orbit an intrinsic line or circle.
For that case (C.51.8) is exactly one of the original
boundary quotients whose lifts define $\mathscr F_O$.
It follows in particular that an empty $\mathscr F_O$
forces every frontier orbit to be a point: apply
the construction to a sequence approaching any
prescribed frontier point.

\emph{The intrinsic family and every boundary equation.}
For $|\mathscr F_O|=1$ use the contraction $K^O_{s,t}$
of C.48, and for $|\mathscr F_O|=2$ use C.46.
Both commute with $H_O$, hence with its restricted
closure $Q$, and both have jointly continuous inverses.
They retain every quotient equation
\[
                    \Phi K^O_{s,t}=Q^B_{s,t}\Phi,
                                                        \tag{C.51.9}
\]
where $Q^B$ is the C.34 median family on the actual
boundary orbit $B$, with its based real covering
lift when that orbit is a circle.
For an empty family, C.32 makes $f_t|_O$ null in
$\operatorname{Homeo}(O)$ and C.33 supplies a based
null-homotopy $K^O$ in its exact $Q$-centralizer.
All frontier orbits are then points, as just proved,
so (C.51.9) for each constant $\Phi$ is automatic.

Define $\overline K^O$ to be $K^O$ on $O$ and
the C.34 family on $\partial O$. We prove its full
joint continuity, rather than continuity only for
fixed parameters. Let
\[
       x_n\in O,\quad x_n\to z\in\partial O,
       \qquad(s_n,t_n)\to(s,t).
                                                        \tag{C.51.10}
\]
Use $x_n=q_ny_n$ and a subsequence $y_n\to y$ as
above. The set
\[
 S=C\ \cup\!\bigcup_{(u,v)\in[0,1]^2}
                    \bigl(K^O_{u,v}(C)\cup(K^O_{u,v})^{-1}(C)\bigr)
                                                        \tag{C.51.11}
\]
is compact in $O$, by joint continuity on a compact
parameter space. Refine the choice of actual operators
in (C.51.4) to satisfy
\[
                       \sup_{w\in S}|h_nw-q_nw|<1/n.
                                                        \tag{C.51.12}
\]
For each $n$ this is a compact-open neighborhood
condition on $q_n\in Q$, so density of the specified
restrictions supplies $h_n\in H_O$.
The previous argument still proves (C.51.6)--(C.51.8).

Commutation with $Q$ and (C.51.12) give
\[
 \begin{split}
 K^O_{s_n,t_n}(x_n)&=q_nK^O_{s_n,t_n}(y_n),\\
 \bigl|K^O_{s_n,t_n}(x_n)
            -h_nK^O_{s_n,t_n}(y_n)\bigr|&<1/n.
 \end{split}                                            \tag{C.51.13}
\]
The points $K^O_{s_n,t_n}(y_n)$ converge in $O$ to
$K^O_{s,t}(y)$ and lie in the compact set $S$.
Uniform convergence (C.51.7), followed by the exact
quotient identity (C.51.9), therefore yields
\[
 \begin{split}
 h_nK^O_{s_n,t_n}(y_n)
      &\longrightarrow\Phi_z^y K^O_{s,t}(y)\\
      &=Q^{Gz}_{s,t}(\Phi_z^y(y))
        =Q^{Gz}_{s,t}(z).
 \end{split}                                            \tag{C.51.14}
\]
For a point orbit the same computation uses the
constant map and gives $z$.
Every subsequence in (C.51.10) has a further subsequence
of this form because $C$ is compact. The value on the
right is independent of all choices. Thus (C.51.14)
holds for the entire sequence and proves continuity
from $O$ to the prescribed frontier value.
Continuity for sequences contained in $\partial O$
is precisely C.34. Splitting a sequence according
to its two pieces proves joint continuity everywhere
on $[0,1]^2\times\overline O$.

The inverse maps have the same argument, using
the inverse images already included in $S$.
Equation (C.51.9) implies
$\Phi(K^O_{s,t})^{-1}=(Q^B_{s,t})^{-1}\Phi$
by composing with these actual inverse maps.
The inverse boundary family is jointly continuous
by C.34. This proves the inverse version of
(C.51.14), and hence joint continuity of the
piecewise inverse on the entire closure.
The maps preserve $O$ and each frontier orbit and
are bijections there; the piecewise inverses are
their actual inverses. They are homeomorphisms.
The endpoint and based-edge identities and
commutation with $H_O$ hold on both pieces.

\emph{Transport and finite pasting.}
The class of orbits in the claim is $H$-invariant:
conjugation by $h\in H$ identifies the two exact
groups $Q_O,Q_{hO}$, their quotient spaces, and
their full boundary-quotient families by C.46.
Choose one orbit from each $H$-orbit of these open
components, make the construction there, and set
\[
             \overline K^{hO}_{s,t}
                   =h\overline K^O_{s,t}h^{-1}.
                                                        \tag{C.51.15}
\]
The equality with every stabilizing operator makes
this independent of the choice of $h$ carrying the
selected representative to a given component.
Median naturality in C.27--C.34 makes its frontier
restriction the identical prescribed C.34 family.

Distinct open-orbit closures meet only in the
closed nonopen locus $Y$. Every one of the new
families and its inverse equals the corresponding
C.34 map there. The finite closed-set pasting
lemma thus gives the asserted maps and inverse
maps on $Y$ together with finitely many closures.
For every finite collection $\mathcal I$, the
homeomorphism $h$ carries this union to the union
indexed by $h\mathcal I$ and intertwines their
homotopies by (C.51.15). This is the exact
operator-transport comparison on those finite unions.

The proof uses one fixed compact fundamental set
inside each treated orbit. It supplies no uniform
control over infinitely many different such sets,
and it does not identify the median boundary value
with the exponential value forced on $X_*$ by C.50.
These are explicit remaining comparisons in C.12.
\end{proof}

\begin{claim}{C.52}{Proved}{Cocompact frontier extensions relative to the exponential region}
Use the actual based ambient-null plane loop $f_t$, the groups
$G=\langle f_t\rangle$ and $H$, and the closed set $X_*$ with its
contraction $E_{s,t}$ from C.50.
For every nonopen $G$-orbit $B$, define a based family
$R^B_{s,t}$ as follows:
\[
 R^B_{s,t}=
 \begin{cases}
 \id,&B\text{ a point},\\
 \text{the exponential fraction with parameter }1-s,
       &G|_B\text{ metabelian},\\
 \text{the C.27--C.29 median family},
       &G|_B\text{ nonmetabelian}.
 \end{cases}                                             \tag{C.52.1}
\]
The exponential fraction uses the actual intrinsic Lie image
classified in C.37. On a circle it uses the entire based lift
of the original loop to the covering group, including its
real angle at each intermediate time.

For every open orbit $O$ disjoint from $X_*$ whose full family
$\mathscr F_O$ of actual lifted boundary partitions has
cardinality at most two, there is an intrinsic based
null-homotopy $K^O_{s,t}$ commuting with its original operator
stabilizer and satisfying every actual boundary quotient
equation
\[
                 \Phi_zK^O_{s,t}=R^{Gz}_{s,t}\Phi_z.      \tag{C.52.2}
\]
For each such $O$ with compact proper-operator quotient
$O/Q_O$, this family extends to a based contraction through
homeomorphisms of $\overline O$, jointly continuous with
its inverse, and equal to (C.52.1) on its entire frontier.
It agrees with $E_{s,t}$ on $\overline O\cap X_*$.

These extensions can be chosen with exact transport under
$H$. They agree on common frontiers. Consequently they
paste with the prescribed C.50 contraction on $X_*$ over
every finite union
\[
                  X_*\ \cup\!
                  \bigcup_{O\in\mathcal I}\overline O,    \tag{C.52.3}
\]
where $\mathcal I$ is a finite collection of the
cocompact orbits just specified.
No continuity of (C.52.1) on the entire nonopen locus,
and no extension over an infinite collection or over the
other noncocompact open orbits, is asserted.
\end{claim}
\passage{Section 23, folios 140-141, scans 279-281; Part C(b) infinite-operator request; relative frontier extension preserving the already prescribed exponential region}


\begin{proof}
\emph{The precise intrinsic target family.}
For a nontrivial intrinsic line or circle orbit, $G|_B$ is
path connected and transitive. The proof of (C.37.9)
classifies every such metabelian image, with its actual
intrinsic homeomorphism-group topology: on a line it is
conjugate to all translations or to the full positive
affine group, and on a circle to all rotations.
That proof uses only these three stated properties of
the image, so it applies to every metabelian orbit in
(C.52.1), without requiring that orbit to lie in $X_*$.

For a real lift $F:X\to\R$ of an actual boundary quotient,
let $b_t$ be its based target path, characterized by
$Fa_t=b_tF$ for the based source lift $a_t$.
The exact source and target coverings are those of
C.42--C.44. The path $b_t$ is a based loop: for a circle
target this is the zero-winding conclusion of C.29,
and for a line target it follows from the identity
endpoint of $f_t$.
Its generated group
\[
                         T=\langle b_t:0\leq t\leq1\rangle
                                                               \tag{C.52.4}
\]
is transitive on $\R$. Indeed its orbits are intervals,
as each word is joined to the identity by a path of
words. A finite endpoint of such an orbit interval is
a common fixed point, which would project to a common
fixed point of the nontrivial transitive target action.
Thus no proper orbit interval is possible.

For a line target, (C.52.4) is the conjugate of the
actual image $G|_B$. For a circle, projection gives a
surjection $T\to G|_B$. When the latter image is
metabelian, its rotation conjugacy lifts to a real
coordinate in which every $b_t$ is a translation.
Transitivity then makes $T$ the full translation group.
Conversely a metabelian $T$ has metabelian quotient
$G|_B$. This proves the exact equivalence used in the
case selection below.

Define, in the specified real lift,
\[
 B_{s,t}=
 \begin{cases}
 \exp_T\bigl((1-s)\log_T b_t\bigr),&T\text{ metabelian},\\
 \text{the C.27 median of the path }b_t,&T\text{ nonmetabelian}.
 \end{cases}                                             \tag{C.52.5}
\]
In the first case $T$ is a simply connected translation
or positive affine Lie group with a global exponential
and logarithm. Formula (C.52.5) and its inverse are
jointly continuous, with all four required based-edge
identities. They commute with every target homeomorphism
commuting with the original path: in the exponential
case each new value lies in $T$, and in the median case
this is C.27.

This rule is natural under any line reparametrization,
including a decreasing one. Such conjugation identifies
the actual groups (C.52.4) topologically, preserves
metabelianness, and, in the Lie cases, intertwines
exponential maps and their unique logarithms. The
median case has the same naturality by C.27.
The source of a real lift of a circle may also be
changed by a covering conjugacy; the same argument
retains the specified based target path.
In particular the resulting map on $B$ is exactly
(C.52.1), independent of all these coordinates.

For clarity, at a time when $f_t|_B=\id$ on a circle,
its based real lift can still be a nonzero integer
translation. Formula (C.52.5) scales that real
translation. It does not replace it by zero.
The endpoint $t=1$ is zero by C.29. This is the same
covering convention as C.50.

\emph{Replacement in the two-foliation construction.}
For $|\mathscr F_O|=2$, retain exactly the source cover,
the two real quotient coordinates, and the domain $D$
of C.45--C.46. Replace each of the two coordinate
median families in (C.46.4) by (C.52.5).
On $D=\R^2$ this is a product of increasing line
homeomorphisms, with a jointly continuous inverse.
On $D=\{u<v\}$ the original two coordinate loops are
identical. Their generated groups and their selected
families are therefore identical, so the strict
inequality is preserved. On $D=\{u<v<u+1\}$ the same
argument applies, and the selected maps commute with
unit translation. In the exponential case this holds
because that translation centralizes $T$; in the
median case it is the original median naturality.
Thus both inequalities are preserved, as are they
by the inverse family.

Every original stabilizing operator, and every one
of its source lifts, permutes the full two-element
family of partitions by the exact comparison
(C.46.5). In these coordinates it is a pair of line
homeomorphisms, possibly followed by exchange.
Commutation with the based source loop identifies
the corresponding target paths by these same
conjugacies. Rule (C.52.5) is natural under them
and makes the same choice on conjugate paths.
Consequently the modified source family commutes
with every such lift, including all source deck maps.
It descends to $O$ with continuous inverse.

Every other actual real quotient is a line
homeomorphism composed with one of the two selected
coordinates, by C.45. The preceding naturality
therefore proves its exact lifted equation too.
Descending through $q_zF_z=\Phi_zp$ proves (C.52.2),
retaining the original source and target covers.

\emph{Replacement in the one-foliation construction.}
For $|\mathscr F_O|=1$, keep the actual closed proper
operator group, its group of all source lifts, and
every coordinate change of C.48.
In (C.48.19) replace only the median base map by
(C.52.5), leaving the fiber interpolation as displayed:
\[
 (u,r)\longmapsto
       \bigl((1-s)h_t(u,r)+su,\ B_{s,t}(r)\bigr).        \tag{C.52.6}
\]
Here the choice of target coordinate is the exact
reparametrization already established in C.48.
We check why all of its operator equations survive
this replacement.

For the unit fiber translation, the fiber equation
$h_t(u+1,r)=h_t(u,r)+1$ is unchanged by the base
replacement. For an operator $(\epsilon u,r+d)$,
the equations are
$h_t(\epsilon u,r+d)=\epsilon h_t(u,r)$ and
$b_t(r+d)=b_t(r)+d$. The first is preserved by
the same convex interpolation; the second says
that translation by $d$ centralizes $T$, so it
holds for (C.52.5).
For the real fiber kernel and a further
$(\lambda u,r+d)$, the fiber function is
$h_t(u,r)=u+\phi_t(r)$ with
$\phi_t(r+d)=\lambda\phi_t(r)$. Scaling $\phi_t$
preserves this equality for either sign of $\lambda$,
and the base comparison is again the same one.

In the normal forms with all real base translations
or a dense lattice base image, the original base
loop is $b_t(r)=r+c_t$. Its selected family is
$B_{s,t}(r)=r+(1-s)c_t$ and commutes with every base
operator. The fiber equations, respectively
independence of $r$ or the exact affine lattice
relation (C.48.21), are preserved by (C.52.6).
The trivial operator case imposes no extra equation.
These are the C.48 normal forms outside its
transitive two-dimensional case. That latter case
already lies in $U_2\subset X_*$ and is absent from
the open orbits of the present claim.

Thus the modified family commutes with the exact
group of all operator lifts, with its original
restriction kernel retained downstairs.
The fiber and base inverse-continuity argument
following (C.48.22) applies unchanged to (C.52.6):
the base family has continuous inverses by
(C.52.5), and fixed bracketing fiber values bound
the increasing convex-interpolation inverses.
The homotopy descends to $O$, with the asserted
endpoint identities. All other actual quotients
are reparametrizations of the one selected
partition; naturality proves (C.52.2) for them.

For an empty family, C.32--C.33 supply a based
null-homotopy in the exact proper-operator
centralizer. There is then no nontrivial actual
quotient equation to check. This finishes the
intrinsic assertion for all three cardinalities.

\emph{Cocompactness and the entire frontier.}
Now assume $O/Q_O$ compact and put $Q=Q_O$.
The first part of C.51 supplies a compact
$C\subset O$ with $QC=O$. It also proves that
every nontrivial frontier orbit gives an actual
quotient in the specified family, after the
explicit basepoint comparison (C.51.8).
In particular an empty family makes all frontier
points fixed.

Define $\overline K^O$ by the just-constructed
$K^O$ on $O$ and by (C.52.1) on $\partial O$.
We first prove its limit along points of $O$;
we do not assume transverse continuity of
(C.52.1). Let
\[
 x_n\in O,\quad x_n\to z\in\partial O,\qquad
                       (s_n,t_n)\to(s,t).              \tag{C.52.7}
\]
Write $x_n=q_ny_n$ with $y_n\in C$ and pass to
a subsequence $y_n\to y\in C$.
The compact set consisting of $C$ and all its
images and inverse images under $K^O$ is exactly
the type of set used in (C.51.11).
Approximate $q_n$ uniformly on this set to
accuracy $1/n$ by restrictions of actual
$h_n\in H_O$, as in (C.51.12).
The open evaluation map of $L=\overline G$ gives
$l_n\to1$ with $l_ny=y_n$. Exact commutation
gives $h_ny=l_n^{-1}h_ny_n\to z$.

The C.42 evaluation homeomorphism, with the
basepoint comparison (C.51.8), now gives
$h_n|_O\to\Phi_z^y$ uniformly on compact subsets.
For a nontrivial target this is an actual
quotient covered by (C.52.2); for a fixed target
it is constant. Commutation with $Q$ therefore
gives precisely
\[
 \begin{split}
 K^O_{s_n,t_n}(x_n)
   &=q_nK^O_{s_n,t_n}(y_n)\\
   &\longrightarrow
       \Phi_z^y K^O_{s,t}(y)
        =R^{Gz}_{s,t}(z). 
 \end{split}                                             \tag{C.52.8}
\]
The approximation error in the second line is
less than $1/n$ before passage to the limit.
Compactness of $C$ and uniqueness of the value
give (C.52.8) for the entire sequence.
This repeats C.51's exact operator comparison
with the new quotient equation; no element of
the restricted closure is assigned an
unproved global extension.

\emph{The diagonal argument at frontier points.}
Take now $z_n\in\partial O$ with $z_n\to z$
and $(s_n,t_n)\to(s,t)$.
For each fixed triple $(s_n,t_n,z_n)$,
(C.52.8) lets us choose $x_n\in O$ such that
\[
 |x_n-z_n|<1/n,\qquad
 \bigl|K^O_{s_n,t_n}(x_n)
            -R^{Gz_n}_{s_n,t_n}(z_n)\bigr|<1/n.         \tag{C.52.9}
\]
Then $x_n\to z$, so applying (C.52.8) to this
new sequence proves
$R^{Gz_n}_{s_n,t_n}(z_n)\to R^{Gz}_{s,t}(z)$.
This establishes the needed joint continuity
on this entire frontier, even across different
target types in (C.52.1).
Splitting a sequence into its interior and
frontier terms proves joint continuity on
$[0,1]^2\times\overline O$.

The same proof applies to the inverse family.
The compact approximation set already contains
all inverse images, and (C.52.2) implies the
exact inverse equation
$\Phi_z(K^O_{s,t})^{-1}=(R^{Gz}_{s,t})^{-1}\Phi_z$.
Use this equation in (C.52.8), then repeat the
diagonal choice (C.52.9) for inverses.
Thus the piecewise inverse is jointly continuous.
All maps are bijections on $O$ and on each
frontier orbit, with these actual inverses;
they are homeomorphisms of the closure.

\emph{Agreement with the fixed exponential region.}
Every group image on an open orbit of $U_2\cup U_3$
is metabelian, by C.36 and C.49. Hence, for all
$g_1,g_2,g_3,g_4\in G$, the plane homeomorphism
$[[g_1,g_2],[g_3,g_4]]$ fixes $U_2\cup U_3$
pointwise. Continuity makes it fix $X_*$ pointwise.
The set $X_*$ is $G$-invariant. Therefore every
nonopen orbit $B\subset X_*$ has metabelian
actual image $G|_B$.

An open orbit outside $U_2\cup U_3$ is disjoint
from $X_*$, because it is an open neighborhood
disjoint from that union. Thus the intersections
$\overline O\cap X_*$ occurring here contain only
nonopen orbits. On a point orbit both
contractions are the identity. On every other
such orbit, (C.52.1) and C.50 use the same
actual intrinsic Lie image and the same
based target-cover path. Exponential naturality
therefore gives equality, including the
intermediate real circle angles described above.

Distinct open-orbit closures meet only in the
nonopen locus. The present construction assigns
there the same rule (C.52.1), so their contractions
and inverses agree. Conjugation by $h\in H$
identifies the actual target images and their
based paths, preserves metabelianness, and
intertwines both branches of (C.52.1).
Choose the intrinsic construction on one
representative of each $H$-orbit of components
and transport it by $h$. Commutation with the
full stabilizer makes this choice independent
of the transporting element, exactly as in
(C.51.15).

Finite closed-set pasting now gives the
homeomorphism families and inverse families
on (C.52.3). A general $h$ intertwines this
union with the union indexed by $h\mathcal I$;
when it preserves the union, it commutes with
the displayed contraction. The family on
$X_*$ is exactly the prescribed one throughout.
The diagonal proof concerns a single fixed
closure. It supplies no uniform control over
infinitely many different compact fundamental
sets, and no compact fundamental set for the
remaining noncocompact quotients. These are
still retained parts of C.12.
\end{proof}

\begin{claim}{C.53}{Proved}{Global quotients and finitely many operator representatives}
Use the ambient-null based plane loop $f_t$, its generated group
$G$, the commuting operator group $H$, and $L=\overline G$ of
C.42. For every open $G$-orbit $O$, choose $x\in O$ and
put $J=L_x$. Define $\mathscr F_O^{\rm all}$ to be the family
of fiber partitions on the common universal covering of $O$
of all real lifts of
\[
 \Phi_z^x:O\longrightarrow Gz,\qquad \Phi_z^x(lx)=lz,
 \quad z\in\operatorname{Fix}(J),\quad
                  Gz\text{ nonopen and nontrivial}.
                                                        \tag{C.53.1}
\]
The target points need not belong to $\partial O$.
Let $U^\dagger$ be the union of all open orbits with either
$\dim Q_O=2$ or $|\mathscr F_O^{\rm all}|\geq3$, and set
$X^\dagger=\overline{U^\dagger}$.
There is a jointly continuous based contraction on
$X^\dagger$, with jointly continuous inverse and exact
$H$-commutation, extending the prescribed contraction on
$X_*$ in C.50. On every nonopen orbit it has the
corresponding exponential value of C.52.1.

Now choose finitely many open orbits
$O_1,\ldots,O_m$ disjoint from $X^\dagger$, each with compact
$O_j/Q_{O_j}$, and put
\[
 U=\bigcup_{j=1}^m\ \bigcup_{h\in H}hO_j,
 \qquad X=\overline U.
                                                        \tag{C.53.2}
\]
The original loop has a based contraction through
homeomorphisms of $X^\dagger\cup X$, jointly continuous
with its inverse, commuting with $H$ and equal to the
specified exponential contraction on $X^\dagger$.
On every nonopen orbit in $X$ its value is exactly the
hybrid prescription $R^B_{s,t}$ of C.52.1.
This assertion allows infinitely many different open
components in (C.53.2), with arbitrary accumulation.
It retains the stated finite number of $H$-representatives
and compactness of their actual operator quotients.
\end{claim}
\passage{Section23, folios140-141, scans279-281; Part C(b), infinite-operator zero-level comparison, global quotient maps and arbitrary accumulation under finitely many operator representatives}


\begin{proof}
\emph{The global quotient family and its exact comparisons.}
C.42 is already stated for every $z\in\operatorname{Fix}(J)$
with nontrivial line or circle orbit; it imposes no
frontier condition on $z$. It gives the actual continuous
open map (C.53.1), its intrinsic target topology, the
common covering construction, and its proper fiber lines.
C.43--C.45 use these maps and their transitive compatible
action to give the connected real fibers, product
coordinates and the exact three pairwise coordinate
domains. Thus each of those conclusions applies to
(C.53.1) with the same maps.

Changing the source basepoint also preserves the complete
family. If $y=ax$ for $a\in L$, then
$L_y=aJa^{-1}$ and, for $z\in\operatorname{Fix}(L_y)$,
\[
                \Phi_z^y(ly)=lz
                     =\Phi_{a^{-1}z}^x(lax).
                                                        \tag{C.53.3}
\]
The targets are the same actual $G$-orbit, because
$Lz=Gz$ for each such quotient, as proved in C.42.
This proves the asserted identification of partitions,
including the original target coverings.

For $h\in H$ stabilizing $O$, choose $a_h\in L$ with
$a_hx=hx$. Exact commutation gives
$L_{hx}=L_x$, hence $a_hJa_h^{-1}=J$, and
\[
                   \Phi_z^x h=\Phi_{a_hz}^x.             \tag{C.53.4}
\]
The point $a_hz$ is fixed by $J$ and lies on the same
nonopen orbit as $z$. Therefore (C.53.4) permutes the
entire family $\mathscr F_O^{\rm all}$. For $hO=O'$, use
basepoint $hx$ in $O'$; then
\[
               \Phi_{hz}^{hx}h=h\Phi_z^x.              \tag{C.53.5}
\]
These equalities identify the families under operator
transport and prove that their cardinalities are
well defined and $H$-invariant.

\emph{Three global quotients give the same actual Lie group.}
Choose three distinct partitions in
$\mathscr F_O^{\rm all}$. The full three-foliation group
of C.47 is constructed exclusively from the real maps
in C.43--C.45 and their common transitive action.
Its proper-pair or ordered-loop argument therefore
applies to these exact maps and gives the same locally
compact Lie group of dimension two or three.
The based target lifts in C.49.2 are loops: for circle
targets C.29 uses the ambient-null plane loop and two
points of that circle, independently of the position
of $O$; real target endpoints are the identity.
The Baire transitivity proof C.49.3, its isotropy
calculation C.49.4--C.49.7 and its closedness argument
now apply to the same based source lifts. Their actual
generated group $N$ is consequently one of
$\mathbb R^2$, $\operatorname{Aff}^+(\mathbb R)$, or
$\mathcal A_2$, with exactly the topology, global
exponential and unique logarithm proved in C.49.

We verify the action comparison for every global
quotient, rather than only for frontier targets.
Use the closed Polish group
\[
 E_O=\{(l,n)\in L\times N:pn=(l|_O)p\}.
                                                        \tag{C.53.6}
\]
Its projection to $N$ is onto, since every $n$ is an
actual word in the lifted loop; the identical global
word supplies $l\in G$. It is an open homomorphism
of Polish groups. If $(l,1)$ is in its kernel, $l$
fixes $O$ pointwise. Equivariance of (C.53.1) then
makes $l$ fix its entire image $B=Gz$ pointwise:
\[
                         l\Phi_z^x(u)=\Phi_z^x(lu)
                                      =\Phi_z^x(u).
                                                        \tag{C.53.7}
\]
C.42 proves continuity of $L\to\operatorname{Homeo}(B)$
in the intrinsic compact-open topology. Its restriction
to $E_O$ therefore factors continuously through the
open quotient to an actual action of $N$ on $B$.
That image is exactly $G|_B$: actual loop words prove
both inclusions. It is metabelian and hence is the
translation, affine, or rotation Lie image of C.37.
The resulting continuous homomorphism and its unique
based covering lift give exactly
\[
 \Phi_z^x K^O_{s,t}=R^B_{s,t}\Phi_z^x,
 \qquad K^O_{s,t}=
        \exp_N\bigl((1-s)\log_N\widetilde f_t\bigr)
                       \text{ descended to }O.         \tag{C.53.8}
\]
Here $R^B$ is the exponential branch of C.52.1,
with the full based real angle for a circle.
The original source covering, every original quotient,
and every restriction kernel are retained.

\emph{The enlarged closed exponential region.}
For a dimension-two proper-operator orbit use the
based lifted contraction of C.36--C.37. For an orbit
with three global partitions use (C.53.8).
Where both constructions apply, their source groups
are the same actual group generated by the same
based lifted loop, as in C.49, so the maps agree.
Operator conjugation transports these groups and
intertwines their exponential maps, by (C.53.5) and
the based covering identity C.49.16.

Every source group just specified has the uniform
fractional-word bound $1913$: C.37 gives the stronger
bound three for the first two groups, and C.50.1--C.50.14
give the stated bound for $\mathcal A_2$.
Each word is a word in the original based path and
its inverse. In particular the full argument
C.50.15--C.50.19 applies to this enlarged union with
precisely the same compact set of actual global word
isotopies. We spell out the transfer of its remaining
inputs. All the source images are metabelian, so
each global commutator-of-commutators fixes
$U^\dagger$, hence its closure. At a nonopen limit
orbit the actual transitive image is therefore the
same translation, affine or rotation group classified
in C.37. For every integer $n$, the residual after
removing $\lfloor n(1-s)\rfloor$ endpoint powers is a
fractional power in the specified simply connected
source group, and thus has the same uniform word
bound. On a line the bounded residual logarithms
force the exponential limit. On a circle the source
residual word loop is null in that same simply
connected group; its exact two-point winding passes
to the limiting orbit, giving the real equality
C.50.19, including its integer part. These are all
inputs of C.50's limit identification and joint
continuity proof. Applying it also to the inverse
loop gives a jointly continuous inverse on
$X^\dagger$.

The original $U_2\cup U_3$ lies in $U^\dagger$.
On that original union the maps are the same
source fractional powers. Continuity extends this
equality to $X_*$. The operator identities also
extend from the dense union. This proves the first
assertion and fixes its prescription for the rest
of the proof.

\emph{Compatible intrinsic maps on the remaining representatives.}
Each $O_j$ in (C.53.2) has
$|\mathscr F_{O_j}^{\rm all}|\leq2$ and
$\dim Q_{O_j}\leq1$, by its disjointness from
$X^\dagger$. For a family of size one or two,
apply the construction of C.52 to this full global
family. Here is the exact transfer. Its source
coordinates and operator normal forms use
C.42--C.45, already verified above for (C.53.1).
Its full-partition invariance is (C.53.4).
The base-family replacement C.52.5 is natural under
every target reparametrization and exchange, and
retains every affine operator equation C.52.6.
No step of those intrinsic constructions uses the
position of the target point on a frontier.
They therefore give a based intrinsic contraction
$K^j_{s,t}$, jointly continuous with its inverse,
commuting with $H_{O_j}$ and with its actual
restricted closure $Q_{O_j}$, and satisfying
\[
                    \Phi_z K^j_{s,t}
                             =R^{Gz}_{s,t}\Phi_z
                 \quad\text{for every map in (C.53.1)}.
                                                        \tag{C.53.9}
\]
For an empty family, C.32--C.33 give the based
intrinsic contraction in that exact proper-operator
centralizer. Every constant equivariant quotient
also satisfies (C.53.9), since its target is fixed.

Discard repetitions among the finitely many
$H$-orbits represented by $O_j$, and transport
$K^j$ to each $hO_j$ by $hK^jh^{-1}$.
Commutation with the full stabilizer makes this
well defined. Its union gives a jointly continuous
homeomorphism family $K$ on the open set $U$;
components are open, and each carries its intrinsic
family and inverse. It commutes with $H$ exactly.

\emph{One compact comparison set for every accumulating sequence.}
For each representative, compactness of
$O_j/Q_{O_j}$ supplies a compact $C_j\subset O_j$
with $Q_{O_j}C_j=O_j$, by C.51's finite-cover
argument. Suppose
\[
 x_n\in U,\quad x_n\longrightarrow z\in X\setminus U,
                  \qquad(s_n,t_n)\longrightarrow(s,t).
                                                        \tag{C.53.10}
\]
The orbit $Gz$ is nonopen. Indeed an open orbit
containing $z$ would contain every sufficiently
late $x_n$, and so would itself be one of the
components in $U$, contradicting $z\notin U$.
There are only finitely many representatives.
On a subsequence, write
\[
                  x_n=h_nq_ny_n,
 \quad h_n\in H,\quad q_n\in Q_{O_j},\quad
                  y_n\in C_j,\quad y_n\longrightarrow y\in C_j
                                                        \tag{C.53.11}
\]
for one fixed $j$. Fix the compact subset of $O_j$
\[
 S_j=C_j\ \cup\!
       \bigcup_{(u,v)\in[0,1]^2}
       \bigl(K^j_{u,v}(C_j)\cup(K^j_{u,v})^{-1}(C_j)\bigr).
                                                        \tag{C.53.12}
\]
Both joint continuity assertions prove compactness.

Choose actual $a_n\in H_{O_j}$ so that
\[
 \sup_{w\in S_j}
          |h_na_nw-h_nq_nw|<1/n,
 \qquad g_n=h_na_n\in H.
                                                        \tag{C.53.13}
\]
For each fixed $n$, the first condition is an
open neighborhood of $q_n$ in
$\operatorname{Homeo}(O_j)$: composition with the
fixed ambient homeomorphism $h_n$ and restriction
to the compact set $S_j$ are continuous.
The defining density of the actual restrictions
$H_{O_j}\to Q_{O_j}$ supplies $a_n$.
Thus the estimate is made \emph{after} the possibly
distorting operator $h_n$; no uniform modulus of
continuity for the sequence $h_n$ is assumed.

The open $L$-orbit map supplies $l_n\to1$ with
$l_ny=y_n$, exactly as in C.51. Since $g_n$
commutes with $L$, (C.53.10)--(C.53.13) give
\[
                       g_ny=l_n^{-1}g_ny_n\longrightarrow z.
                                                        \tag{C.53.14}
\]
Every element of $L_y$ fixes this limit.
The evaluation homeomorphism C.42.1 therefore gives
\[
                g_n|_{O_j}\longrightarrow\Phi_z^y
                     \quad\text{uniformly on compact sets}.
                                                        \tag{C.53.15}
\]
The exact basepoint change (C.53.3) identifies this
map with one of the global quotients in (C.53.1)
when its target is nontrivial; otherwise it is
constant. This is the reason the enlarged family
is required in this argument.

Transport and commutation with $Q_{O_j}$ now yield
\[
 \begin{split}
 K_{s_n,t_n}(x_n)&=h_nq_nK^j_{s_n,t_n}(y_n),\\
 \bigl|K_{s_n,t_n}(x_n)
              -g_nK^j_{s_n,t_n}(y_n)\bigr|&<1/n.
 \end{split}                                            \tag{C.53.16}
\]
The inner points converge in $S_j$ to $K^j_{s,t}(y)$.
Equations (C.53.15) and (C.53.9) consequently give
\[
                   K_{s_n,t_n}(x_n)\longrightarrow
                         \Phi_z^yK^j_{s,t}(y)
                           =R^{Gz}_{s,t}(z).            \tag{C.53.17}
\]
Every subsequence admits a further subsequence of
this form, by the finite number of $j$ and the
compactness of their $C_j$. The value in (C.53.17)
is independent of every such choice. Hence the
entire sequence has that limit. If a representative
has an empty global family, (C.53.15) also proves
that every limit reached through its translates
is a fixed point.

Define the extension on $X\setminus U$ to be
$R^{Gz}_{s,t}(z)$. The diagonal approximation
C.52.9 now applies with the dense set $U$ in $X$:
for each varying boundary point and fixed parameter
pair choose a point of $U$ close to it and with
close image, using (C.53.17). Applying (C.53.17)
to the resulting sequence proves joint continuity
on all of $X$. This argument does not require
prior transverse continuity of the hybrid rule.

Use the inverse images already present in $S_j$
and the exact inverse version of (C.53.9) to
repeat (C.53.13)--(C.53.17) for inverse maps.
The same diagonal approximation then proves
joint inverse continuity on $X$.
The maps preserve $U$ and each nonopen limit
orbit, with the displayed maps as mutual
inverses. Thus they are homeomorphisms of $X$.
Their based-edge identities and $H$-commutation
extend from $U$ by continuity; the restriction
to every nonopen orbit is the prescribed $R^B$.

Finally, $U$ is disjoint from $X^\dagger$.
Every point of $X\cap X^\dagger$ lies on a
nonopen orbit; otherwise its open orbit would
meet $U$ and be a component of $U$.
The metabelian identity on $X^\dagger$ and the
first part of the proof show that both
contractions there use precisely the same
exponential branch of C.52.1, with the same
based covering path. They and their inverses
agree on the full intersection. Pasting the two
closed sets $X$ and $X^\dagger$ gives the
asserted contraction on their union, preserving
the prescribed family on $X^\dagger$.

The finite-representative argument does not
supply compact comparison sets for arbitrary
noncocompact quotients or infinitely many
$H$-inequivalent representatives. Nor does it
extend the family through an exterior region
of nonopen orbits. These remain exact parts
of the general C.12 obligation.
\end{proof}

\section{D. Profinite extensions and the limits of the realization problem}

We write $P=\widehat\pi_0$, where $\pi_0=\pi_1(U,p)$ for a connected orientable surface of finite type with $\chi(U)<0$. All maps between profinite groups in this section are continuous. A lift of an outer action means a homomorphism inducing the prescribed outer homomorphism, with its labels in the operator group retained. For a finite group $G$ acting faithfully on $U$, $E$ denotes the group of lifts to the universal cover, not a mapping class group. The completion map $E\to\widehat E$ is injective: $E$ is a finitely generated linear hyperbolic orbifold group and is residually finite. We use $E^+$ for its orientation-preserving subgroup. The passages answered here are \S31, folios 198--199, scans 391--394; the inserted $G=G^+$ on scan 391 is material.

\begin{claim}{D.1}{Proved}{Exact completion of the lift extension}
The completion of $1\to\pi_0\to E\xrightarrow qG\to1$ is the exact sequence
\[
1\longrightarrow P\longrightarrow\widehat E\xrightarrow{\widehat q}G\longrightarrow1.
\]
It is the extension attached to the completed outer action. An invariant discretization therefore supplies a canonical comparison of extensions, with identity maps on $P$ and $G$.
\end{claim}
\passage{31; 198--199 verso; 391--394}

\begin{proof}
Every subgroup of finite index in $\pi_0$ contains a subgroup normal and of finite index in $E$: take its core in $E$, using $[E:\pi_0]<\infty$. Thus $\pi_0$ has its full profinite topology induced from $E$, and its closure is $P$. The kernel of $\widehat q$ is that closure, since $\pi_0=q^{-1}(1)$ is open. Conjugation in the completed extension is the continuous extension of conjugation in the dense subgroup $E$. The uniqueness of a center-free kernel extension, proved in Part A, identifies this extension with the pullback of $\Aut(P)\to\Out(P)$. The identification fixes the named kernel and quotient, so the assertion of canonicity has no unmentioned conjugator.
\end{proof}

\begin{claim}{D.2}{Proved}{Finite subgroups of completed surface orbifold groups, including reflections}
Every finite subgroup of $\widehat E$ is conjugate to a finite subgroup of $E$. Two finite subgroups of $E$ conjugate in $\widehat E$ are conjugate in $E$. These finite subgroups are cyclic rotation groups or dihedral groups, with a single reflection included as the dihedral case of order two.
\end{claim}
\passage{31; 198--199 verso; 391--394}

\begin{proof}
We specify the two inputs and the reduction. Bridson--Conder--Reid, Theorem 5.1 \cite{R-BCR}, applies to the finitely generated Fuchsian group $E^+$: every finite subgroup of $\widehat{E^+}$ is conjugate into $E^+$. Boggi--Zalesskii, Theorem 6.3 \cite{R-BZ}, gives conjugacy and normalizer comparison for finite $p$-subgroups of $E$:
\[
N_{\widehat E}(A)=\overline{N_E(A)}.
\]
Here its hypotheses hold because $E$ has a finite-index surface group (a free group in the noncompact case), hence is virtually compact special; its good-group and finiteness hypotheses are precisely those in Theorem 6.3. We use only this prime-power theorem, not an assertion obtained by commuting arbitrary intersections with closure.

Put $K^+=K\cap\widehat{E^+}$ for a finite $K\leq\widehat E$. The first input makes $K^+$ cyclic. When $K^+=1$, the orientation character embeds $K$ in $C_2$, and the prime-power input handles $K$. Otherwise choose the unique subgroup $A$ of order $p$ in $K^+$, for one prime $p\mid |K^+|$. It is characteristic in $K^+$ and normal in $K$. Conjugate $A$ into $E^+$. A nontrivial elliptic rotation of the hyperbolic plane has one fixed point $x$. Its normalizer in the discrete group $E$ stabilizes $x$, and conversely every element of $E_x$ normalizes the unique subgroup of order $p$ in its cyclic rotation subgroup. Thus $N_E(A)=E_x$ is finite cyclic or dihedral. The displayed comparison gives $N_{\widehat E}(A)=E_x$, so $K\leq E_x\leq E$.

For injectivity, suppose $K,L\leq E$ and $tKt^{-1}=L$, with $t\in\widehat E$. The orientation character gives $tK^+t^{-1}=L^+$. In the case $K^+=1$ apply prime-power conjugacy separation. Otherwise choose $A\leq K^+$ as above and $B=tAt^{-1}\leq L^+$. Prime-power conjugacy separation supplies $e\in E$ with $eAe^{-1}=B$. Now $e^{-1}t$ normalizes $A$, hence belongs to $N_E(A)\subset E$. Therefore $t\in E$, after this adjustment. This proves the required conjugacy assertion. Finally finite discrete hyperbolic-plane isometry groups stabilize a point and embed in $O(2)$; the listed cyclic and dihedral possibilities follow.
\end{proof}

\begin{claim}{D.3}{Proved}{The comparison of full and partial section classes}
For every subgroup $H\leq G$, completion induces a bijection
\[
\operatorname{Sec}(E_H\to H)/\pi_0
\ \longrightarrow\
\operatorname{Sec}(\widehat E_H\to H)/P,
\qquad E_H=q^{-1}(H).
\]
Both sets are finite. In particular the first conjecture on scan 393 holds for geometrically realized actions, with orientation-reversing operators allowed.
\end{claim}
\passage{31; 198--199 verso; 391--394; 18; 99--101; 197--203}

\begin{proof}
We prove the change from total-group conjugacy to kernel conjugacy, since omission of this step changes the question. A profinite section has a finite image $K$ mapping isomorphically to $H$. Apply D.2 to $E_H$, and choose $t\in\widehat E_H$ with $tKt^{-1}\leq E_H$. Choose $e\in E_H$ with $q(e)=\widehat q(t)$. Then $e^{-1}t\in P$, and conjugation by this element puts $K$ in $E_H$. The restriction of $q$ identifies the resulting subgroup with $H$, so it defines a discrete section with the correct labels.

Conversely, let discrete section images $K,L$ be conjugate by an element of $P$. D.2, applied in $E_H$, gives an $e\in E_H$ with $eKe^{-1}=L$. Let $l\in L$ be the unique element satisfying $q(l)=q(e)$. Then $l^{-1}e\in\pi_0$ still conjugates $K$ to $L$. Because the projection is injective on each section image, this equality of subgroups is equality of the labelled section homomorphisms after conjugation. This proves injectivity.

There are finitely many $E_H$-conjugacy classes of finite subgroups: they occur in the finite list of cone stabilizers, corner stabilizers, and mirror components of a finite-type quotient orbifold, and each finite stabilizer has finitely many subgroups. Each such class splits into at most $[E_H:\pi_0]=|H|$ kernel-conjugacy classes. The asserted finiteness follows.
\end{proof}

\begin{claim}{D.4}{Proved}{The exact profinite fixed subgroup of a section}
Let $s:H\to E_H$ be a nontrivial discrete section. Its fixed subgroup $P^{s(H)}=C_P(s(H))$ is trivial when $s(H)$ contains a nontrivial rotation. For a reflection $s(H)=\langle t\rangle$, it is $1$ for a noncompact real component and is the closure of the primitive oval subgroup, isomorphic to $\widehat\Z$, for a compact real component. The same assertions hold for every profinite section by D.3.
\end{claim}
\passage{31; 199; 393; 33; 222 verso--225; 440--445}

\begin{proof}
For a rotation choose $A$ of prime order inside $s(H)^+$. D.2's proof gives $N_{\widehat E}(A)=E_x$, a finite group. Thus $C_P(s(H))\leq E_x\cap P=1$, because $P$ is torsion-free. For a reflection the prime-power centralizer comparison gives
\[
C_{\widehat E}(t)=\overline{C_E(t)},\qquad
C_P(t)=\overline{C_{\pi_0}(t)}.
\]
The second equality uses that $P$ is open and $P\cap E=\pi_0$: the intersection of the closure of a subgroup with an open subgroup is the closure of their intersection. On the universal cover, $C_{\pi_0}(t)$ is the group of deck transformations preserving and commuting with reflection in its fixed line. It acts freely on that line. Its quotient is the corresponding real component of $U$. The deck group is trivial for a line and infinite cyclic for a circle; the circle is an essential primitive simple curve. Its induced profinite topology is full, so its closure is $\widehat\Z$.

This proves, rather than assumes, the correction to the expected answer in Part D(a): a dihedral section with a nontrivial rotation has fixed subgroup $1$. The value $\widehat\Z$ belongs to a reflection in an oval. It is not the fixed subgroup of a larger dihedral stabilizer containing that reflection.
\end{proof}

\begin{claim}{D.5}{Refuted}{Unique maximal sections require a rotation}
A finite subgroup of $\widehat E$ containing a nontrivial rotation is contained in a unique maximal finite subgroup. The assertion for every nontrivial finite subgroup is false, already for type $(0,3)$ with reflections.
\end{claim}
\passage{18; 99--101; 197--203; 31; 198--199; 391--393}

\begin{proof}
For the positive assertion conjugate the subgroup into $E$ using D.2. Its rotation part has a subgroup $A$ of prime order. Every finite overgroup has cyclic orientation-preserving part, so normalizes this same $A$. All finite overgroups therefore lie in the finite group $N_{\widehat E}(A)=E_x$, which is itself the unique maximal one.

For the counterexample use the $(2,3,\infty)$ reflection group
\[
W=\langle s_0,s_1,s_2\mid s_i^2=1,
(s_0s_1)^2=(s_1s_2)^3=1\rangle
       = V_4*_{C_2}S_3.
\]
Map it onto $S_3\times C_2$, with $r^3=a^2=\tau^2=1$, $ara=r^{-1}$ and $\tau$ central, by
$s_0\mapsto\tau$, $s_1\mapsto a\tau$, $s_2\mapsto ra\tau$.
The vertex groups inject. The kernel is an orientable torsion-free group of index $12$, Euler characteristic $-1$, with three cusps; hence it is $\pi_{0,3}$. The two maximal finite subgroups
$\langle s_0,s_1\rangle$ and $\langle s_1,s_2\rangle$ contain the same reflection $s_1$ and have distinct images in the quotient. They remain distinct and are not kernel-conjugate after completion. Their generated subgroup is $W$, so no finite overgroup contains both. Part B supplies the finite certificate and the geometric kernel computation.
\end{proof}

\begin{claim}{D.6}{Proved}{A well-defined transfer of the geometric rotation invariant}
For a realized $E$, its rotation invariant on finite-order elements transfers uniquely to the torsion of $\widehat E$ by conjugating each torsion element into $E$. This assertion uses the orientation module, so that conjugation by an orientation-reversing element acts by $-1$ on its values. Its values embed canonically in $\widehat T\otimes_{\Z}\Q/\Z$, where $\widehat T=T\otimes_{\Z}\widehat\Z$.
\end{claim}
\passage{17; 77--79 verso; 155--160; 31; 198--199; 391--393}

\begin{proof}
Use D.2 for existence. Explicitly, for torsion $h\in\widehat E$ choose $t\in\widehat E$ with $u=tht^{-1}\in E$ and set
\[
 \mu_{\widehat E}(h)=\varepsilon(t)\mu_E(u),
\]
where $\varepsilon:\widehat E\to\{\pm1\}$ is the orientation character, acting on $\widehat T$. Suppose $u,v\in E$ and $v=tut^{-1}$ in $\widehat E$. For a rotation apply the proof of D.2 to an order-$p$ subgroup of $\langle u\rangle$: first align it by discrete conjugation, then put the remaining conjugator in the finite discrete stabilizer. Thus the elements themselves, and not just their cyclic subgroups, are discretely conjugate, with the same orientation character of the conjugator. For reflections apply prime-power element conjugacy separation, or D.2 because an order-two subgroup has a unique nonidentity element. This establishes independence with the indicated action on $T$. A rotation of order $n$ has a value in $T/nT\subset T\otimes\Q/\Z$; the identification $T/nT=\widehat T/n\widehat T$ defines the claimed inclusion.
\end{proof}

% Replacement for D.7.  The bibliography entry bWilkes2017Seifert is supplied
% in d_mu_references.bib.  This fragment has no global preamble.
\begin{claim}{D.7}{Proved}{Equivariance of the profinite rotation invariant}
Let $E$ be the lift group of a finite action on an orientable hyperbolic
surface of finite type, and put $P=\widehat\pi\triangleleft\widehat E$.
Every automorphism $\alpha$ of $\widehat E$ preserving $P$ and its indexed
peripheral structure satisfies
\[
 \mu_{\widehat E}(\alpha(u))=\alpha_T\bigl(\mu_{\widehat E}(u)\bigr)
 \qquad (u\in\widehat E\text{ of finite order}).
\]
Here $\alpha_T$ is the induced automorphism of the orientation line
$\widehat T$ defined in A.2, and $\mu_{\widehat E}$ is the invariant of
D.6.  The assertion holds for compact and noncompact surfaces, and for
actions containing orientation-reversing elements.
\end{claim}
\passage{17; 77--79 verso; 155--160; 31; 198--199; 391--393}

\begin{proof}
Choose an orientation, so that $\alpha_T$ is multiplication by some
$\lambda\in\widehat\Z^\times$.  Conjugation by an element of
$\widehat E$ acts on $\widehat T$ by its geometric orientation sign.
Naturality of this action and commutativity of the automorphism group of a
rank-one module show that $\alpha$ preserves that sign.  Thus $\alpha$
restricts to an automorphism of $\widehat B$, where $B=E^+$ is the
orientable Fuchsian subgroup.  A nonidentity negative torsion element is a
reflection and has invariant $1/2$; multiplication by a profinite unit
fixes this value.  It remains to treat torsion in $\widehat B$.

We first make explicit the behavior of orientation under finite covers.
For a degree-$d$ connected finite cover of oriented surfaces, the map on
orientation homology takes the fundamental class upstairs to $d$ times
the fundamental class downstairs.  For punctured surfaces this follows
directly in A.2's kernel model: the peripheral factors above any one
puncture map by their positive local degrees, whose sum is $d$.
For closed surfaces it is the degree formula for fundamental classes.
The same formulas hold over $\widehat\Z$, either by this peripheral
matrix or by the natural finite-coefficient cohomology comparison.
Consequently, if $\alpha$ carries an open surface subgroup $V$ to $V'$,
the multipliers on their orientation lines and on $P$ satisfy
$d\lambda_V=d\lambda$.  Multiplication by a positive integer is
injective on $\widehat\Z=\prod_p\Z_p$, so $\lambda_V=\lambda$.
This proves the required compatibility with further finite covers; it
does not divide by $d$ in a finite coefficient ring.

\emph{The noncompact case.}
Let $L$ be a peripheral subgroup of $\pi$, and let $K$ be the primitive
cusp subgroup of $B$ containing it.  Write $K=\langle w\rangle$ and
$L=\langle w^n\rangle$, with compatible positive generators.
We have
\[
 N_{\widehat B}(\overline L)=\overline K\simeq\widehat\Z.
\]
Here is the reduction from Part A, rather than an assumption that arbitrary
normalizers commute with completion.  For an element of the normalizer,
choose a representative in $B$ of its coset modulo $P$.  That coset
stabilizes the profinite peripheral class.  Distinct discrete peripheral
classes remain distinct in $P$, so the same coset stabilizes the discrete
peripheral class.  Adjust its representative by an element of $\pi$ to
put it in $N_B(L)=K$.  The remaining factor belongs to
$N_P(\overline L)=\overline L$, as proved in Part A.  This proves the
displayed equality.

The loop condition on $\alpha|_P$ sends a positive generator of
$\overline L$ to a $P$-conjugate of the $\lambda$th power of the
corresponding generator of another peripheral subgroup $\overline L'$.
After making this conjugation, $\alpha$ sends $\overline K$ to
$\overline K'$ and has the form $w\mapsto(w')^a$.
The indices
$[\overline K:\overline L]=[\overline K':\overline L']=n$
are equal because $\alpha$ preserves $P$.  Thus $na=n\lambda$ in
$\widehat\Z$, and hence $a=\lambda$.  Every primitive orbifold cusp
therefore has precisely the same multiplier $\lambda$.

Write the signature of $B$ as $(h;m_1,\ldots,m_r;s)$ with $s\geq1$.
Choose cone generators $c_j$, of orders $m_j$, and cusp generators $w_i$
in a common orientation.  The abelianized presentation is
\[
 \widehat B^{\mathrm{ab}}
 =\left(\widehat\Z^{2h}\oplus
       \bigoplus_{j=1}^r\Z/m_j\,c_j\oplus
       \bigoplus_{i=1}^s\widehat\Z\,w_i\right)
       \big/\left\langle\sum_jc_j+\sum_iw_i\right\rangle.
\]
D.2 applied to $B$ makes $\alpha$ permute the maximal cone subgroups.
Hence, up to conjugacy in $\widehat B$,
$\alpha(c_j)=c_{\tau(j)}^{k_j}$, where
$m_{\tau(j)}=m_j$ and $k_j$ is a unit modulo $m_j$.
The preceding paragraph gives
$\alpha(w_i)\sim w_{\sigma(i)}^\lambda$.
Apply $\alpha$ to the displayed relation and subtract $\lambda$ times
that relation.  This yields
\[
 \sum_j(k_j-\lambda)c_{\tau(j)}=0.
\]
Eliminating $w_s$ exhibits the cone factors as an independent direct sum,
so $k_j\equiv\lambda\pmod{m_j}$ for every $j$.
The local formula of B.12 now proves the assertion on every cone
subgroup and, by conjugation in $\widehat B$, on all its torsion.

\emph{The compact case.}
First suppose that there are at least two cone points.  We use Wilkes's
cohomological calculation with its exact scope:
for a closed orientable hyperbolic orbifold group $B$ with cone
generators $c_j$ of orders $m_j$, an automorphism of $\widehat B$
sending each $c_j$ to a conjugate of $c_j^{k_j}$ acts on
$H^2(\widehat B;\Z/n)$, for every $n$, by one scalar
$\kappa\in\widehat\Z$ satisfying
$\kappa\equiv k_j\pmod{m_j}$.
This is \cite[Proposition 5.7]{bWilkes2017Seifert}, numbered Proposition
6.6 in arXiv:1512.05587v1.  Its proof concerns arbitrary automorphisms
of the indicated form; no lifting to a Seifert group is used in that
calculation.  In particular the scalar is independent of $n$.

For transparency, the part of the calculation that supplies both facts
is as follows.  After taking coinvariants, the relation cells $y_0,y_j$
have boundaries
$\partial y_0=\sum_jx_j$ and $\partial y_j=m_jx_j$.
The automorphism sends $x_j$ to $k_jx_j$, $y_j$ to $k_jy_j$, and
$y_0$ to $\kappa y_0+\sum_j\rho_jy_j$.
Commutation with the boundary gives
$\kappa+m_j\rho_j=k_j$ in $\widehat\Z$.
The dual coboundaries are $y_0^*+m_jy_j^*$, which shows that the
induced action is multiplication by $\kappa$ for every finite
coefficient group.  The completed resolution and the naturality of this
calculation are established in the cited proposition and its use of
Proposition 2.9 there.

The cone permutation of our $\alpha$ is realized by an
orientation-preserving homeomorphism of the underlying surface carrying
the marked cone points to the correspondingly ordered points.  It
induces an orbifold automorphism $\beta$ with the same permutation and
power $1$ on each cone.  Put $\psi=\beta^{-1}\alpha$ and apply the
preceding result to $\psi$.  Let
\[
 \eta_n\in H^2(\widehat B;\Z/n)
\]
be the image of the oriented fundamental cohomology class of the
underlying closed surface under the canonical map from the orbifold
classifying space to its underlying space.  Equivalently, in the
displayed resolution it is $[y_0^*]$, up to the common orientation sign.
The finite-coefficient class is transported to continuous cohomology by
goodness of Fuchsian groups, as in Wilkes's calculation.
Since $\beta$ preserves orientation, $\beta^*\eta_n=\eta_n$.
Thus $\alpha^*\eta_n=\psi^*\beta^*\eta_n=\kappa\eta_n$.
This argument does not require $\beta$ to preserve the chosen subgroup
$P$.

Set $N=[B:\pi]$.  Restriction to the surface subgroup gives
\[
 \operatorname{res}_{P}(\eta_n)=N\omega_n,
 \qquad\omega_n\in H^2(P;\Z/n)
\]
with $\omega_n$ the positive surface class: the underlying map from
the covering surface to the orbifold's underlying surface is an
orientation-preserving branched map of degree $N$.  Naturality of
restriction and A.2's duality yield
\[
 N\kappa\,\omega_n=N\lambda\,\omega_n
       \quad\text{for every }n.
\]
Indeed $\alpha_*[U]=\lambda[U]$ implies
$\langle\alpha^*\omega_n,[U]\rangle
=\langle\omega_n,\alpha_*[U]\rangle=\lambda$.
Taking all finite coefficients gives $N(\kappa-\lambda)=0$ in
$\widehat\Z$, hence $\kappa=\lambda$.  The local congruences are
therefore $k_j\equiv\lambda\pmod{m_j}$, exactly as required.

We remove the single-cone case from the resolution input explicitly.
When $r=1$, the underlying genus $h$ is positive and
$B^{\mathrm{ab}}=\Z^{2h}$, because the one cone generator is zero in
abelianization.  The characteristic subgroup
\[
 B_1=\ker\bigl(B\longrightarrow B^{\mathrm{ab}}/2B^{\mathrm{ab}}\bigr)
\]
has index $2^{2h}$ and contains every conjugate of the cone group.
Its orbifold has $2^{2h}$ cone points of the same order.  The subgroup
$\widehat{B_1}$ and $P\cap\widehat{B_1}$ are $\alpha$-invariant.
The argument for at least two cones applies there, and the finite-cover
calculation at the beginning identifies its surface multiplier with
$\lambda$.  For $r=0$ there is no nonidentity torsion in $\widehat B$,
by D.2.  Thus no exceptional compact signature remains.

Finally, choose positive local cone generators, for which
$\mu(c_j)=1/m_j$ in the chosen orientation line.  The proved congruence
says $\mu(\alpha(c_j))=\lambda/m_j$; its powers have the same
property.  Every positive torsion element is conjugate in $\widehat B$
to one of these powers by D.2.  Together with the reflection case this
proves the stated equation on all torsion of $\widehat E$.
\end{proof}

\noindent\emph{Why the earlier rank-one argument was insufficient.}
Integral relative orbifold cohomology can have a torsion summand:
for $C_3*C_3$ with cusp $\langle ab\rangle$ it has the presentation
$\langle u,v,w\mid3u=w=3v\rangle\cong\Z\oplus\Z/3$.
The proof above uses the cusp relation in the noncompact case and the
action on the full finite-coefficient cohomology in the compact case.
It therefore supplies the missing local-power congruences without
choosing an unproved canonical generator through that torsion summand.

\subsection{Completion of geometric mapping class groups and full loop automorphisms}

\begin{claim}{D.8}{Proved}{All finite subgroups in the geometric completion in the three requested types}
Let $M=\Gamma_{g,\nu}$ include puncture permutations and orientation reversal. For $(g,\nu)=(0,3),(0,4),(1,1)$, completion induces a bijection on conjugacy classes of finite subgroups of $M$ and $\widehat M$. Every finite subgroup of $\widehat M$ is solvable. For every finite $K\leq M$,
\[
 N_{\widehat M}(K)=\overline{N_M(K)},\qquad
 C_{\widehat M}(K)=\overline{C_M(K)}.
\]
These assertions concern $\widehat M$ with its given embedding in $\Out_{\rm lac}(P)$.
\end{claim}
\passage{30; 194 verso; 384; 30; 196 verso; 388}

\begin{proof}
For the three-punctured sphere $M=S_3\times C_2$ is finite. For the punctured torus $M=\mathrm{GL}_2(\Z)$; its principal congruence subgroup of level $3$ is a torsion-free free group of finite rank. For the four-punctured sphere the subgroup fixing each puncture and preserving orientation is a free group of rank two and has finite index. Thus all three $M$ are finitely generated virtually free groups.

Bridson--Piwek, Proposition 5.10 \cite{R-BP}, states for every finitely generated virtually free group that finite subgroup conjugacy classes are preserved bijectively by profinite completion. Its proof uses the finite graph of finite groups for $M$: completion gives the profinite fundamental group of the same graph, a finite subgroup fixes a vertex of its profinite Bass--Serre tree, and subgroup conjugacy separability distinguishes the resulting discrete subgroups. This applies without an orientation restriction and proves the first assertion. The normalizer equality is Ribes--Zalesskii, Theorem 2.6 \cite{R-RZ}, for finitely generated subgroups of finitely generated virtually free groups. To deduce the centralizer equality, the homomorphism $N_{\widehat M}(K)\to\Aut(K)$ has finite discrete target. Its kernel is open; the intersection of a dense subgroup with this open kernel is dense in the kernel. This is exactly $C_M(K)$.

For solvability, a finite subgroup of $M_{0,4}$ embeds in $S_4\times C_2$ by the torsion-free pure oriented kernel; one in $\mathrm{GL}_2(\Z)$ embeds in $\mathrm{GL}_2(\F_3)$ by the level-$3$ kernel. The former group is solvable. The latter has normal series with factors inside $C_2$, $A_4$, and the central $C_2$, since $\mathrm{PGL}_2(\F_3)=S_4$; it too is solvable. The finite case $S_3\times C_2$ is immediate. Transfer by the proved conjugacy bijection gives solvability in $\widehat M$.
\end{proof}

\begin{claim}{D.9}{Proved}{The nonsolvable discrete possibilities in genus at most two}
A nonsolvable finite subgroup of the extended mapping class group of a punctured orientable surface of genus at most two is realized on a sphere and has group $A_5$ or $A_5\times C_2$. Its puncture set is invariant under the corresponding icosahedral action. In particular it has at least twelve punctures. This classification does not assert the absence of additional nonsolvable finite subgroups in completions for all larger puncture numbers.
\end{claim}
\passage{30; 194 verso; 384}

\begin{proof}
Realize the finite subgroup by Part C. On the sphere its orientation-preserving subgroup is cyclic, dihedral, $A_4$, $S_4$, or $A_5$, by the finite rotation-group classification. The first four are solvable and an extension of one of them by an orientation group of order at most two is solvable. The orthogonal normalizer of the icosahedral rotation group is $A_5\times\{\pm I\}$: a normalizing orthogonal transformation preserves its rotation axes and its icosahedral cellulation, and composition with an icosahedral rotation leaves either $I$ or $-I$. A nonsolvable full group is consequently $A_5$ or $A_5\times C_2$. The smallest nontrivial orbit has twelve vertices; an invariant puncture set has at least twelve points.

On a torus the kernel of the action on $H_1$ consists of a finite translation group, which is abelian. The orientation-preserving linear quotient is cyclic of order $1,2,3,4$, or $6$, and the extended linear quotient is cyclic or dihedral. These groups and their extensions by the translation group are solvable.

In genus two the hyperelliptic involution is central in the full automorphism group of the invariant Riemann surface. Adjoin it when it is absent. The quotient is a sphere with six branch points. A nonsolvable orientation-preserving quotient would contain $A_5$, whose smallest orbit has twelve points, contradicting preservation of the six branch points. The quotient is therefore solvable; so is its central extension by $C_2$ and its extension by orientation reversal. Punctures can only reduce the finite group. This proves the asserted list and its restriction to the discrete setting.
\end{proof}

\begin{claim}{D.10}{Proved}{Finite torsion in the odd-prime normalized pro-$p$ centralizer}
Let $p$ be odd, let $F^{(p)}$ be the free pro-$p$ group on $x,y$, and define $GT_{0,p}$ by the normalized automorphisms
\[
 F_{\lambda,f}(x)=x^\lambda,\qquad
 F_{\lambda,f}(y)=f^{-1}y^\lambda f,
 \quad \lambda\in\Z_p^\times,\quad f\in(F^{(p)})',
\]
satisfying the two-cycle and three-cycle equations. Every finite subgroup of $GT_{0,p}$ is cyclic of order dividing $p-1$. There is a subgroup of order $p-1$, and every finite subgroup is conjugate in $GT_{0,p}$ to a subgroup of it. For $p\geq5$ this supplies finite loop-preserving outer automorphism groups which are not conjugate to geometric groups of the three-punctured sphere.
\end{claim}
\passage{25 bis; 157 verso--160; 312--317; 30; 194 verso; 384}

\begin{proof}
The normalized description is faithful: a conjugator between two normalized automorphisms first belongs to $C(x)=\langle x\rangle$; the normalization $f\in F'$ and the $y$ equation remove its exponent. The same observation identifies these automorphisms with the pure outer centralizer of the three puncture permutations. The equations are closed conditions, composition is the usual substitution law, and their inverses satisfy the same equations. Hence this is a profinite group. The Galois action on $\mathbf P^1_{\overline\Q}-\{0,1,\infty\}$ gives elements of this group with multiplier the $p$-adic cyclotomic character. This character surjects onto $\Z_p^\times$: the finite cyclotomic extensions have Galois groups $(\Z/p^n)^\times$, and their inverse limit is $\Z_p^\times$.

The kernel of $\Aut(F^{(p)})\to\Aut(F^{(p)}/\Phi(F^{(p)}))$ is pro-$p$; the exact statement and its application are also given in Herfort--Zalesskii, proof of Theorem 6.3 \cite{R-HZ}. For the normalized group the action on this Frattini quotient is the scalar $\lambda\bmod p$. Thus the kernel of
\[
 GT_{0,p}\longrightarrow\F_p^\times
\]
is pro-$p$, and the map is surjective. The profinite Schur--Zassenhaus theorem supplies a complement $C_{p-1}$ and conjugacy of complements in the inverse image of any subgroup of $\F_p^\times$. One can see the passage from the finite theorem directly: apply the finite theorem to every quotient by an open normal subgroup of the pro-$p$ kernel, then take the inverse limit of the nonempty finite sets of compatible complements and conjugators. The finite intersection property follows by passing to a common finite quotient.

We check the necessary absence of torsion in the kernel of the multiplier. A free pro-$p$ group has its injective Magnus map into
$\Z_p\langle\!\langle X,Y\rangle\!\rangle^\times$, with $x\mapsto1+X$, $y\mapsto1+Y$. The map is injective because reduction modulo $p$ is the usual faithful free pro-$p$ Magnus map; it extends continuously from the free group for the $(p,X,Y)$-adic topology on the power-series algebra. The ordinary homogeneous degree is used only in the following leading-term calculation, not as the topology for that extension. An automorphism acting identically on $F^{(p)}_{\rm ab}=\Z_p^2$ acts identically on the degree-one quotient of this power-series algebra. For a nonidentity such automorphism choose the least degree $d\geq2$ of a nonzero correction to $X$ or $Y$. The degree-$d$ correction to its $n$th iterate is $n$ times that correction: substitution of the higher terms can change only degrees above $d$. Each homogeneous coefficient module is a torsion-free $\Z_p$-module, so no positive iterate is the identity. Consequently the multiplier is injective on a finite subgroup.

The torsion subgroup of $\Z_p^\times$ is cyclic of order $p-1$. Therefore every finite subgroup maps injectively to $\F_p^\times$. It is a complement in the inverse image of its image, so the preceding conjugacy theorem proves the full assertion. Finally the geometric pure group of type $(0,3)$ consists of $1$ and complex conjugation, with multipliers $1$ and $-1$. A generator with multiplier a primitive $(p-1)$st root cannot be conjugated into it; the multiplier is unchanged by conjugation even in the full loop-preserving outer group. The first example has order four at $p=5$.
\end{proof}

\begin{claim}{D.11}{Proved}{Executed finite checks and their exact scope}
The program \texttt{code/root/finite\_checks.py} verifies the two-cycle and three-cycle equations in every Heisenberg group $H(\F_p)$ for $p=5,7,11$ and every $\lambda\in\F_p^\times$, using
\[
 f(a,b)=[a,b]^{(\lambda^2-1)/24},\qquad
 m=(\lambda-1)/2.
\]
In particular it verifies an order-four normalized automorphism at $(p,\lambda\bmod p)=(5,2)$ and an order-three one at $(7,2)$. The multipliers of the corresponding infinite finite-order subgroups are Teichm\"uller lifts of these residues; the integer $2\in\Z_5^\times$ itself has infinite order. D.10, not the finite checks alone, proves the existence of compatible infinite pro-$p$ lifts.
\end{claim}
\passage{25 bis; 159 verso--160; 316--317; 30; 194 verso; 384}

\begin{proof}
The program defines all $p^3$ elements as triples and uses
$(a,b,c)(a',b',c')=(a+a',b+b',c+c'+ab')$.
It enumerates inverses and powers and evaluates each side of both equations in that group. Its output records every multiplier and both truth values. Since all commutators are central, the associated automorphism is
$(a,b,c)\mapsto(\lambda a,\lambda b,\lambda^2c)$; the program checks multiplication preservation and computes its order. Division by $24$ is only used at the three displayed primes. No assertion of this formula at $2$ or $3$ is made. At $p=3$, D.10 gives only order-two prime-to-$p$ torsion. At $p=2$, the Magnus argument still proves that a finite subgroup injects into $\{\pm1\}$, but the complement argument does not settle conjugacy of involutions in $GT_{0,2}$.
\end{proof}

\begin{claim}{D.12}{Open}{Torsion in the full profinite $GT_0$ and the full low-type groups}
D.29 refutes the finite geometric-recognition assertion for $(g,\nu)=(0,4)$ by an actual involution with a complete preserved peripheral system. We do not settle whether every nonidentity torsion element of the full profinite $GT_0$ is conjugate to $(-1,1)$, or the full finite-subgroup recognition questions for $(g,\nu)=(0,3),(1,1)$. The geometric-completion assertion in all three low types is proved in D.8, and the odd-prime analogue is decided in D.10. D.32 proves that every finite subgroup of the full pure three-cusp group, including $GT_0$, has order at most two; it gives the faithful map $G\to S_3\times C_2$ for the full three-cusp group. D.34, D.36, D.41 and D.42 identify the global orientation character of every such finite group: every involution has its global $2$-adic sign, and every order-three element has multiplier $1$. D.35 identifies the exact remaining third-cusp conjugator equation, proves full generation from it, and identifies its conjugator sets in the pure group and in $GT_0$. D.38 gives a complete characteristic $A_5^{19}$ quotient test, including proved exclusions of finite candidates from full lifts. D.44 solves the coboundary and the first symmetry equation simultaneously in the full group; the third-cusp equation and hence existence of the conjugator for each actual full action remain unresolved. D.33 proves that full homology detects every finite subgroup in type $(1,1)$ and that its order divides $24$. D.43 proves that every involution in this full punctured-torus group has its global orientation sign. D.45 identifies the exact characteristic binary-icosahedral refinement and its full finite solution counts; the same cusp-triangle survivors pass. D.46 proves that every cyclic order-four outer action on the punctured-torus group has an actual order-four lift, using the three full section classes of its positive square. D.47 checks the exact conjugator coordinates in the binary refinement: precisely 10,125,000 of the 324,000,000 coordinate pairs pass, with the count unchanged by finite invertibility and both symmetry relations. D.48 proves that a full punctured-torus order-three action has exactly two section classes and a torsion-free cusp normalizer whose cube exponent is a 3-adic unit. D.49 computes the standard pure involution centralizer and the exact two-element fibres of its conjugacy orbit map. D.50 classifies the finite two-subgroups of the full punctured-torus group as exactly $1,C_2,C_4,C_2\times C_2,D_8$, exhibits all five, and proves the global multiplier character on every four-group or dihedral group of order eight. D.51 identifies the second torsion-free normal index-three kernel of every order-three torus extension, with its exact type-$(1,3)$ intersection cover and three labelled peripheral normalizers. D.52 gives the twelve possible abstract finite-subgroup types, realizes nine of them, and proves the reflection multiplier for every subgroup of type $S_3,D_{12},D_{24}$. D.53 proves the unique full labelled order-four section class and its exact normalizer. D.54 excludes every nonzero prime-to-three invariant character for a cubic torus action by an explicit peripheral norm comparison in a cyclic cover. D.55 supplies the corresponding square comparison, proves that every positive involution acts as $-I$ at every prime, and proves faithfulness of every primary homology map on every finite subgroup. The finite torus abstract types are consequently exactly the nine realized types $1,C_2,C_3,C_4,C_2\times C_2,C_6,S_3,D_8,D_{12}$; the three extra types in the earlier allowed list are excluded. D.56 computes the prime-section permutations of every cyclic-six action and proves that its full extension has no order-six subgroup or finite-order representative of the outer generator. D.57 proves actual labelled lifts for every four-group and order-eight dihedral action, with the exact three-or-one and one full section counts. D.58 determines the exact negative-involution section counts and their full fixed-subgroup mod-two cohomology, and proves that the two internal cusp-reflection classes fuse under full conjugacy. Pure-involution orbit exhaustion, the remaining scalar cubic and quartic torus multipliers, and geometric conjugacy remain unresolved.
\end{claim}
\passage{25 bis; 157 verso--160; 312--317; 30; 194 verso; 384}

Here is the exact obstruction reached, rather than an extrapolation from computation. The homomorphism from a free profinite group to the product of its maximal pro-$p$ quotients does not detect all its finite quotients. Our program constructs $A_5$ from a $5$-cycle and a $3$-cycle, enumerates its sixty elements, and verifies that its commutator subgroup has sixty elements. Thus $A_5$ is a nontrivial two-generated finite quotient invisible to every nilpotent quotient. The pro-$p$ complements constructed in D.10 do not give a compatible solution of the defining equations in such quotients.

We also checked the $f=1$ attack: the three-cycle relation becomes the power-product relation of Lochak--Schneps, Lemma 9 \cite{R-LS}; their finite modular quotients force the global alternatives $\lambda=1$ or $-1$. A choice of independent signs in the factors of $\widehat\Z$ is not a solution of this relation. The finite nilpotent search therefore supplies no full-profinite counterexample. D.32 excludes every element of odd prime order in the full pure three-cusp group and every element of finite order greater than two. The smallest unresolved full-$GT_0$ case is therefore an additional involution up to conjugacy; D.34 proves that its multiplier is globally $-1$, supplies its three section classes and a commutator coboundary, and leaves the conjugator inside the specified loop group to be determined. D.23 supplies the exact order-bounded compatibility criterion over all isolated finite shadows. D.24 executes nonnilpotent refinement tests and proves that specific order-three shadows cannot lift with order three; D.25 computes the entire dihedral inverse limit and the exact residual multiplier-one norm equation. D.28 proves the additional prime-order section and fixed-subgroup constraint, including its exact maximal pro-$p$ comparison. The finite-cover argument of D.32 now excludes all odd-order full lifts, beyond those bounded computations; the involution conjugacy problem retains its exact isolated-refinement formulation. D.29 uses an arbitrary Galois pants isomorphism and its inverse on two different pieces; this produces a finite involution in type $(0,4)$ without requiring that pants isomorphism itself have finite order. For genus at most two and arbitrary puncture number, the solvable comparison is supplied by the checked argument in F.7. D.31 now proves that every nonsolvable finite subgroup in these completions is abstractly $A_5$ or $A_5\times C_2$ on a punctured sphere, and gives exactly the geometric positive puncture-orbit sizes. Its proof also excludes every nonsolvable finite subgroup in the genus-one and genus-two completions. What remains is geometric conjugacy of the actual $A_5$ embeddings and the corresponding centralizer comparison of F.7.

The later tests also close two specific attempted constructions of a lower-type counterexample. D.39 constructs an actual rank-two Haran involution but proves, using its unique section class and the trivial centralizer of that section, that it preserves neither hyperbolic loop marking under any full profinite change of basis. D.37 gives the exact four-cusp descent condition for a negative transposition. D.40 excludes the required additional deck symmetry for the actual D.29 example by the unequal homomorphism counts $10200$ and $4920$ in $\operatorname{PSL}_2(\F_{11})$; the argument also excludes deck maps moving the original separator. These are proved obstructions for the specified constructions. They do not provide the conjugator still required in D.35.

\subsection{Descent along a finite covering}

Put $A(P)=\Aut_{\rm lac}(P)$ and let a discretization mean the image of a loop-preserving isomorphism $\widehat{\pi_0}\to P$ restricted to $\pi_0$. Thus it is an actual dense subgroup with its peripheral classes. Define $\operatorname{Discr}(P)$ using actual subgroups, and define $\operatorname{Discr}^{\rm out}(P)$ by quotienting this set by conjugation by $P$. This convention is necessary for an outer action: an outer automorphism acts canonically on the second set. Write $J=\Aut_{\rm lac}(\pi_0)$, $B=\Out_{\rm lac}(P)$ and $M=\Out_{\rm lac}(\pi_0)$ for their embedded images.

\begin{claim}{D.13}{Proved}{Restriction and normal descent with the peripheral step included}
For an open subgroup $V\leq P$, restriction of the corresponding finite covering defines an injective map
\[
 \operatorname{Discr}(P)\longrightarrow\operatorname{Discr}(V),
 \qquad\Delta\longmapsto\Delta\cap V.
\]
For normal $V$, its image consists exactly of the lattices $\Lambda\subset V$ whose outer deck action $P/V\to\Out_{\rm lac}(V)$ is induced by orientation-preserving loop automorphisms of $\Lambda$. The cyclic peripheral subgroups of the descended lattice have closures equal to the specified peripheral subgroups of $P$.
\end{claim}
\passage{28; 178--179; 352--354}

\begin{proof}
We prove the self-normalizer input, including the class-preserving step. Paris \cite[Theorems 1.7 and 3.2]{R-Paris} proves conjugacy $p$-separability for closed orientable surface groups and free groups. Thus two elements of $\Lambda$ conjugate in $\widehat\Lambda$ are already conjugate in $\Lambda$.

For a free group $F_r=\langle x_1,\ldots,x_r\rangle$, let $\varphi$ preserve every conjugacy class. After an inner adjustment we have $\varphi(x_1)=x_1$. Write $\varphi(x_i)=w_ix_iw_i^{-1}$. Choose a reduced representative $w$ for the double coset $\langle x_1\rangle w_i\langle x_i\rangle$, so that $w_i=x_1^awx_i^b$ and $w$ begins with neither $x_1$ nor $x_1^{-1}$ and ends with neither $x_i$ nor $x_i^{-1}$. The word $x_1wx_iw^{-1}$ is cyclically reduced with length $2+2|w|$. Since $\varphi(x_1x_i)$ is conjugate to $x_1x_i$, we obtain $w=1$. Hence
$\varphi(x_i)=x_1^{n_i}x_ix_1^{-n_i}$. For distinct $i,j\geq2$, conjugating $\varphi(x_ix_j)$ by $x_1^{-n_i}$ gives
$x_i x_1^{n_j-n_i}x_jx_1^{n_i-n_j}$; cyclic length is two exactly when $n_i=n_j$. Thus all $n_i$ agree, and $\varphi$ is inner. The rank-two case has only one exponent and gives the same conclusion.

For a closed orientable surface group of genus $g\geq2$, Dehn--Nielsen--Baer realizes the outer class of a class-preserving automorphism by a mapping class $f$. It fixes every oriented closed-curve class and acts identically on $H_1$, so preserves orientation. Choose a finite filling Alexander system of simple curves. The Alexander method (\cite[\S2.3]{A_FarbMargalit2012}) makes the subgroup fixing all those curve classes finite: its possible permutations of intersection points and edges form a finite group; the kernel fixes the resulting graph and is isotopic to the identity on its complementary discs. Therefore $f$ has finite order. Realize the cyclic group generated by $f$ using Part C. For an action of order $d$, rational cochains of the quotient are the invariant cochains upstairs, by an equivariant triangulation and averaging over the finite stabilizers. The quotient's first Betti number is consequently $2g$, because the action on $H_1$ is trivial. Its underlying genus is $g$. Riemann--Hurwitz now gives
\[
 2-2g=d\left(2-2g-\sum_j(1-1/m_j)\right).
\]
For $d>1$ the right side is strictly smaller than $2-2g<0$, a contradiction. Thus $f=1$ and the automorphism is inner.

Now let $h\in N_{\widehat\Lambda}(\Lambda)$. Its conjugation automorphism is class-preserving by the first paragraph, hence is conjugation by an element of $\Lambda$ by the two arguments just given. Dividing $h$ by that element centralizes the dense lattice and therefore all of $\widehat\Lambda$. Part A gives trivial centre, so $h\in\Lambda$. This proves $N_{\widehat\Lambda}(\Lambda)=\Lambda$ with the exact comparisons needed in LM3 B.8.

For normal $V$, a descended lattice $\Delta$ normalizes $\Lambda=\Delta\cap V$. Conversely, an element of $N_P(\Lambda)$ has a coset in the finite group $P/V$. The intersection $N_P(\Lambda)\cap V$ is $\Lambda$ by the just-proved input. Given that the deck action is induced on $\Lambda$, form the pullback discrete extension $E_\Lambda$ of $P/V$ by $\Lambda$ inside $\Aut_{\rm lac}(\Lambda)$. We construct its map into $P$ explicitly. An element of the pullback is a pair $(g,\varphi)$, with $g\in P/V$, $\varphi\in\Aut_{\rm lac}(\Lambda)$, and prescribed outer class. Extend $\varphi$ continuously to $\widehat\varphi\in\Aut(V)$ and choose a lift $p_g\in P$ of $g$. Equality of outer classes gives a $v\in V$ with $\widehat\varphi=\Inn(vp_g)|_V$. This $v$ is unique, since $C_P(V)=1$ by Part A. Send $(g,\varphi)$ to $vp_g$. Changing $p_g$ changes $v$ inversely and leaves the product unchanged; multiplication is respected because both products induce the same conjugation on $V$ and have the same quotient label. This defines an injective homomorphism: an element in its kernel has $g=1$ and induces the identity on the dense subgroup $\Lambda$. Its image is exactly $N_P(\Lambda)$, by restricting a normalizing conjugation to $\Lambda$ and applying the same uniqueness. The discrete group is torsion-free, as a subgroup of $P$. The finite-extension realization theorem of Part C gives a surface lattice $\Delta=E_\Lambda$; its orientation is positive because the deck action on $V$ has multiplier $+1$. Its completion is $P$ by D.1.

We supply the peripheral argument left incomplete in LM3 B.10. Let $C\leq\Lambda$ be a primitive peripheral subgroup, and let $Q=C_P(\overline C)$ be the maximal procyclic peripheral subgroup in $P$ containing $\overline C$. Put $D=C_\Delta(C)$. The quotient surface supplied by Part C makes $D$ its primitive cusp subgroup. We have $\overline D\leq Q$ and
\[
 \overline D\cap V=\overline C=Q\cap V.
\]
The image of $D$ in $P/V$ is the stabilizer of the $\Lambda$-conjugacy class of $C$. To see this, adjust a representative of a stabilizing coset by an element of $\Lambda$ to normalize $C$; orientation preservation changes normalization into centralization. Conversely a centralizing element fixes the class. The identical argument in the profinite group identifies the image of $Q$ with the stabilizer of the $V$-conjugacy class of $\overline C$. The map from discrete cusp classes to their profinite classes is bijective and equivariant by peripheral noncommensurability and the finite covering construction of Part A. These two stabilizers are therefore equal. The two closed subgroups $\overline D\leq Q$ have the same intersection with the open subgroup $V$ and the same image in $P/V$, so they are equal. This proves the promised loop identification.

In particular a descended lattice is forced to equal $N_P(\Lambda)$, proving injectivity for normal $V$. For arbitrary $V$, take $W=\bigcap_{p\in P}pVp^{-1}$, an open normal subgroup, and use the actual intersection map $\Lambda\mapsto\Lambda\cap W$. Equality after restriction to $V$ implies equality after restriction to $W$, hence equality of the original lattices by the normal case. This proves injectivity without assuming $V$ normal.
\end{proof}

\begin{claim}{D.14}{Proved}{The image for a nonnormal finite covering}
Let $V\leq P$ be open, $W=\operatorname{core}_P(V)$ and $\Lambda\in\operatorname{Discr}(V)$. Set $\Lambda_W=\Lambda\cap W$. The image condition for $\Lambda$ is that the full deck action of $P/W$ on $W$ is induced by orientation-preserving loop automorphisms of $\Lambda_W$. This condition is independent of a choice of coset representatives. The unique extension is $N_P(\Lambda_W)$, and its intersection with $V$ is $\Lambda$.
\end{claim}
\passage{28; 178--179; 352--354}

\begin{proof}
Necessity follows by restriction of an actual lattice in $P$. For sufficiency apply the normal descent statement D.13 to $W\triangleleft P$, obtaining $\Delta=N_P(\Lambda_W)$. Both $\Delta\cap V$ and $\Lambda$ are discretizations of $V$ with intersection $\Lambda_W$ with the normal open subgroup $W$ of $V$. The injectivity already proved inside $V$ gives $\Delta\cap V=\Lambda$. The deck action is an outer action, so changing a representative multiplies its automorphism by an inner automorphism of $W$; its membership in the embedded geometric outer subgroup is unaffected. This proves the asserted independence and the exact nonnormal image description.
\end{proof}

\begin{claim}{D.15}{Proved}{The explicit nonextending discretization in scan 354}
Let $U=\mathbf P^1-\{0,1,\infty\}$ and
$U'=\mathbf P^1-\{0,\infty,1,-1\}$, with the double cover
$z\mapsto z^2$. Put $P=\widehat{\pi_1(U)}$, let $V\triangleleft P$ be its index-two subgroup, and let $\Lambda=\pi_1(U')\subset V$ be the geometric discretization. There is an explicit continuous loop automorphism $\gamma$ of $V$ for which $\gamma(\Lambda)$ is a discretization of $V$ that does not extend to $P$, although it has the same prediscretization as $\Lambda$.
\end{claim}
\passage{28; 178--179; 352--354}

\begin{proof}
The deck involution $s:z\mapsto-z$ fixes $0,\infty$ and exchanges $1,-1$. Let $a$ be a simple curve separating $\{0,1\}$ from $\{\infty,-1\}$ and let $b=s(a)$; write $T_a,T_b$ for positive Dehn twists. They give a free basis of the pure oriented mapping class group $M^{!+}_{0,4}\simeq F_2$: in the ordered moduli space $\mathcal M_{0,4}\simeq\mathbf P^1-\{0,1,\infty\}$ they are monodromies around two distinct boundary points. One can verify the completed embedding directly from A.12: filling the fourth puncture identifies this pure group with the point-pushing group $\pi_{0,3}$, and the peripheral-normalization map sends its completion into $\Aut(\widehat\pi_{0,3})$ as $\Inn(\widehat\pi_{0,3})$. The latter is isomorphic to $\widehat\pi_{0,3}$ by triviality of its centre. The composite is injective, so the completion embeds already in $\Out_{\rm lac}(V)$. Thus the pure group's closure is its free profinite completion, with abelianization $\widehat\Z^2$ on the indicated twists.

Choose the fully specified profinite integer
\[
 \alpha_2=0\in\Z_2,\qquad
 \alpha_\ell=1\in\Z_\ell\quad(\ell\ {\rm odd}),
 \qquad \alpha\in\widehat\Z.
\]
It is not an integer: its $2$-adic coordinate would force that integer to be zero, contrary to its $3$-adic coordinate. Let $u=T_a^\alpha$ in the closure of $M^{!+}_{0,4}$. Since $s$ preserves orientation,
$sT_as^{-1}=T_b$, and the continuous power maps give
\[
 u^{-1}su=s\,T_b^{-\alpha}T_a^\alpha .
\]
The pure factor on the right has abelianization $(\alpha,-\alpha)$. Every element of the discrete pure group has integral abelianization. Moreover
$M_{0,4}\cap\overline{M^{!+}_{0,4}}=M^{!+}_{0,4}$: the puncture permutation and orientation sign are detected by finite quotients. Consequently $u^{-1}su$ does not belong to the discrete extended mapping class group.

Choose any continuous automorphism $\gamma\in\Aut_{\rm lac}(V)$ representing $u$. Its ambiguity is conjugation by an element of $V$, which does not change membership in the outer geometric subgroup. The geometric outer subgroup associated to $\gamma(\Lambda)$ is $uM_{0,4}u^{-1}$. The calculation says that the deck involution is absent from this subgroup. The necessary direction of D.13 therefore proves that $\gamma(\Lambda)$ does not extend to $P$. Since $u$ lies in the closure of the pure geometric group, its coset modulo the closed geometric stabilizer is the identity, proving equality of the prediscretizations.

The code in \texttt{code/root/finite\_checks.py} verifies compatible Chinese-remainder residues of $\alpha$ and the puncture partitions interchanged by $s$. Nonintegrality and the infinite completed-group conclusion follow from the coordinate argument above. This spells out and verifies the proposed element in LM3's review, \S4.2.
\end{proof}
\noindent\emph{Manuscript:} \S28, folios 178--179, scans 352--354, especially \emph{Un tel $\gamma$ existe toujours}. We have proved the indicated index-two example; the word \emph{toujours} is not used to assert the construction for every finite covering.


\begin{claim}{D.16}{Proved}{A precise fixed-set formula for every actual action}
Let $A$ be a group, $J\leq A$, and $G\leq A$ a finite subgroup. Under the left action of $G$ on $A/J$,
\[
 (A/J)^G=\{aJ:a^{-1}Ga\leq J\}.
\]
It is the disjoint union of the $C_A(G)$-orbits
\[
 \coprod_{[a]} C_A(G)\big/\bigl(aC_J(a^{-1}Ga)a^{-1}\bigr),
\]
where $[a]$ ranges over $J$-conjugacy classes of the labelled embeddings $g\mapsto a^{-1}ga$ occurring in $J$. This formula applies to actual discretizations with $(A,J)=(A(P),\Aut_{\rm lac}(\pi_0))$, to outer discretizations with $(A,J)=(B,M)$, and to prediscretizations with $(A,J)=(B,\overline M)$.
\end{claim}
\passage{28; 178--181; 352--358}

\begin{proof}
The fixed-coset condition is $gaJ=aJ$ for every $g$, exactly the displayed inclusion. Two fixed cosets $aJ,bJ$ are in the same centralizer orbit exactly when there are $c\in C_A(G)$ and $j\in J$ with $b=caj$. Substitution gives $b^{-1}gb=j^{-1}a^{-1}gaj$ for every labelled $g$. Conversely this equality makes $b j^{-1}a^{-1}$ centralize $G$. The stabilizer of $aJ$ is $C_A(G)\cap aJa^{-1}=aC_J(a^{-1}Ga)a^{-1}$, proving the orbit formula.

The homogeneous-space identifications use the extension of each discrete loop automorphism to completion and its injectivity. The stabilizer of an actual lattice in $A(P)$ is exactly $J$. Its inner-conjugacy quotient is $B/M$. Since the closure of $\Inn(\pi_0)$ is $\Inn(P)$, the closed based stabilizer contains all inner automorphisms and its quotient is $\overline M$; hence the prediscretization set is $B/\overline M$. These identifications also show explicitly why a bare outer element is not assigned an action on the set of actual subgroups without choosing a lift.
\end{proof}

\subsection{Torsion-free extensions and recognition}

\begin{claim}{D.17}{Proved}{Torsion-free extensions with an invariant discretization}
Let $1\to P\to\mathcal E\to G\to1$ be a profinite extension with $G$ finite, $P=\widehat\pi_0$ hyperbolic, and the outer action preserving the loop structure and a discretization. A torsion-free $\mathcal E$ is the profinite completion of a torsion-free discrete surface group, with the descended peripheral system. The quotient surface may be nonorientable.
\end{claim}
\passage{28; 178 verso; 353; 24; 147 verso; 293}

\begin{proof}
First the outer action is faithful. Its kernel is the image of $C_{\mathcal E}(P)$: an element acting inner can be divided by its unique inner element of $P$. The centralizer meets $P$ trivially, so embeds in the finite quotient $G$ and is finite. Torsion-freeness makes it trivial. Construct the discrete extension $E$ of the stabilized lattice with this faithful outer action, using centre-freeness. By D.1, its completion is the given $\mathcal E$ with the identity on kernel and quotient. The discrete extension embeds in its completion because it is the finite extension realized by Part C and hence residually finite. Torsion-freeness of $\mathcal E$ therefore implies torsion-freeness of $E$. In a realization, a nontrivial point stabilizer gives a nontrivial finite-order element of the lift group. Thus $G$ acts freely on the surface associated to $\pi_0$, and its quotient is a surface. Its fundamental group is $E$ by lifting paths. The primitive cusp group in $E$ is the centralizer of the covering cusp group; D.13 proves equality of their closures with the specified procyclic peripherals. Orientation-reversing free deck transformations yield the nonorientable quotient and its orientation character. This is the exact scope of the descent argument.
\end{proof}

\begin{claim}{D.18}{Proved}{The unmarked pro-$p$ extension statement for punctured surfaces}
A finitely generated torsion-free pro-$p$ group $\mathcal E$ containing an open free pro-$p$ group $P$ of rank $r$ is free pro-$p$. Its rank is
\[
 1+\frac{r-1}{[\mathcal E:P]}.
\]
In particular the displayed quantity is an integer. This proves the underlying-group assertion for punctured pro-$p$ surfaces, without identifying the peripheral system. More generally, every torsion-free profinite group with an open normal pro-$p$ subgroup is itself pro-$p$. Thus D.19--D.20 apply to arbitrary finite-quotient extensions of their pro-$p$ kernels; the quotient is forced to be a $p$-group.
\end{claim}
\passage{28; 178 verso; 353; 24; 147 verso; 293}

\begin{proof}
Herfort--Zalesskii, Theorem 1.1 \cite{R-HZ}, writes such a finitely generated virtually free pro-$p$ group as the fundamental pro-$p$ group of a finite graph of finite $p$-groups, with vertex groups embedded in $\mathcal E$. Torsion-freeness makes all vertex and edge groups trivial. The fundamental pro-$p$ group is therefore the free pro-$p$ group on the cycles of the graph. The pro-$p$ Schreier rank formula gives $r-1=[\mathcal E:P](\operatorname{rank}(\mathcal E)-1)$, proving the formula. The theorem retains no designated peripheral subgroups in its graph-of-groups conclusion; their identification is a further assertion, not supplied by this calculation.

We prove the assertion about a general finite quotient. Write $G=\mathcal E/P$ and fix a prime $q\ne p$. For every open normal $U\triangleleft\mathcal E$ contained in $P$, the kernel of $\mathcal E/U\to G$ is a finite $p$-group. A Sylow $q$-subgroup of $\mathcal E/U$ therefore has order $|G|_q$ and maps isomorphically to a Sylow $q$-subgroup of $G$. Fix one such subgroup $Q\leq G$. The finite nonempty sets of Sylow $q$-subgroups above $Q$ admit compatible choices under quotient maps: an image is again a Sylow subgroup, and any specified one lifts by taking a Sylow subgroup of its inverse image. Compactness gives an inverse-limit subgroup of $\mathcal E$ isomorphic to $Q$, since all these transition maps on the chosen subgroups are isomorphisms. Torsion-freeness forces $Q=1$. Hence every prime divisor of $|G|$ is $p$, and every finite quotient of $\mathcal E$ is a $p$-group. This is exactly the assertion that $\mathcal E$ is pro-$p$.
\end{proof}

\begin{claim}{D.19}{Proved}{Recognition of the peripheral pro-$p$ extension}
Let $p$ be any prime, let $P$ be the pro-$p$ completion of an orientable hyperbolic surface group of type $(g,\nu)$ with $\nu>0$, and let
\[
 1\longrightarrow P\longrightarrow\mathcal E\longrightarrow G\longrightarrow1
\]
be an exact sequence of pro-$p$ groups with $G$ finite. Assume that $\mathcal E$ is torsion-free and its outer action preserves the $P$-conjugacy classes of the peripheral subgroups. Choose one $L_j$ from each $\mathcal E$-orbit of these classes. Then
\[
 \mathcal N_j=N_{\mathcal E}(L_j)=C_{\mathcal E}(L_j)\simeq\Z_p,
\]
and the pair $(\mathcal E;\{\mathcal N_j\})$ is continuously isomorphic, preserving the indexed peripheral conjugacy classes, to the pro-$p$ completion of a compact surface with its boundary circles. Its interior is orientable for odd $p$. At $p=2$ the interior may be nonorientable; its orientation character is the induced action on the orientation line of $P$.
\end{claim}
\passage{28; 178 verso; 353; 24; 147 verso; 293}

\begin{proof}
Each $\mathcal N_j$ is closed, and
$\mathcal N_j\cap P=N_P(L_j)=L_j$ by Part A. Thus $\mathcal N_j$ is a torsion-free pro-$p$ group with an open subgroup $\Z_p$. It is finitely generated, by lifting generators of its finite quotient and adjoining a generator of $L_j$. Apply D.18 to $\mathcal N_j$: the free pro-$p$ rank formula gives rank one. Consequently $\mathcal N_j\simeq\Z_p$, and it centralizes $L_j$.

We identify the restricted peripheral family, including its multiplicities. The restriction of $(\mathcal E;\{\mathcal N_j\})$ to $P$ contains the subgroups
\[
 P\cap e\mathcal N_je^{-1}=eL_je^{-1},
 \qquad e\in P\backslash\mathcal E/\mathcal N_j .
\]
The map from this double-coset set to the orbit of the $P$-conjugacy class of $L_j$ is surjective by definition. It is injective: an equality
$e'L_je'^{-1}=a eL_je^{-1}a^{-1}$ with $a\in P$ is equivalent to
$e'^{-1}ae\in N_{\mathcal E}(L_j)=\mathcal N_j$, exactly equality of the double cosets. Thus the restricted family consists of all original peripheral classes, each once.

Wilkes, \cite[Proposition 5.11]{R-WilkesRelative}, proves finite-index ascent and descent of Poincar\'e duality pairs at $p$ for a $p$-torsion-free profinite ambient group, with compatible dualizing modules. The pair $(P;\{L_i\})$ is orientable of dimension two: this follows from its standard free pro-$p$ surface presentation and the converse direction of \cite[Theorem 3.3]{R-Wilkes}. We have checked the exact restricted family and torsion hypothesis. Hence $(\mathcal E;\{\mathcal N_j\})$ is a Poincar\'e duality pair of dimension two.

Its dualizing module has underlying group $\Q_p/\Z_p$, with action through a continuous character $\chi:\mathcal E\to\Z_p^\times$ whose restriction to $P$ is trivial. Naturality of duality under conjugation identifies its character with the action on the surface orientation line, using the inverse convention on the dual cohomology line. The image of $\chi$ is a finite $p$-subgroup because it factors through $G$: it is trivial at odd $p$ and is trivial or $\{\pm1\}$ at $p=2$. Thus inversion of the character does not affect the answer. Its restriction to each $\mathcal N_j$ is trivial, by the boundary-duality restriction in \cite[Proposition 5.9]{R-WilkesRelative} and the trivial duality character of $\Z_p$.

D.18 makes $\mathcal E$ free pro-$p$. Write the peripheral count as $b+1$ and the rank as $n+b$. The condition $n\geq0$ in the classification follows from the relative cohomology sequence: the cokernel of
\[
 H^1(\mathcal E;\F_p)\longrightarrow
       \bigoplus_{j=0}^{b}H^1(\mathcal N_j;\F_p)
\]
has dimension one, since $H^2(\mathcal E;\F_p)=0$ and the relative top group has dimension one. Hence $b\leq\operatorname{rank}(\mathcal E)$.
The classification in \cite[Theorem 3.3]{R-Wilkes} gives a free basis
$x_1,\ldots,x_n,s_1,\ldots,s_b$, with each $s_i$ generating a conjugate of $\mathcal N_i$, and a generator $s_0$ of $\mathcal N_0$ in one of these forms:
\[
\begin{array}{ll}
 \chi=1:&
 s_0=[x_1,x_2]_{-}\cdots[x_{n-1},x_n]_{-}\,s_1\cdots s_b,
 \\[2pt]
 \chi\ne1,\ n\ {\rm odd}:&
 s_0=x_1^2[x_2,x_3]_{-}\cdots[x_{n-1},x_n]_{-}\,s_1\cdots s_b,
 \\[2pt]
 \chi\ne1,\ n\ {\rm even}:&
 s_0=x_1^2[x_1,x_2]_{-}[x_3,x_4]_{-}\cdots
 [x_{n-1},x_n]_{-}\,s_1\cdots s_b .
\end{array}
\]
Here $[x,y]_{-}=x^{-1}y^{-1}xy$ matches that source's convention. The first line has $n$ even. The other two have $p=2$; their parameter $2^f$ is zero, since $\chi$ has no infinite principal-unit part.

These are surface words, with explicit changes of basis. Inverting each handle generator changes the commutator convention in the first line. The second line gives one crosscap and orientable handles. For the third, put
\[
 a=x_1x_2^{-1},\qquad b'=x_2,\qquad
 x_1=ab',\quad x_2=b'.
\]
Then $x_1^2[x_1,x_2]_{-}=a^2(b')^2$, which gives two crosscaps and orientable handles. Replacing $s_0$ by its inverse gives the usual boundary relation and preserves its subgroup. Sending the discrete surface generators to these free pro-$p$ generators induces the asserted isomorphism of completed pairs, with exactly the indexed peripheral conjugacies already specified. The free rank formula gives
$\chi(\mathcal E)=\chi(P)/[\mathcal E:P]<0$, so its interior is hyperbolic. This proves the peripheral statement at every prime, including $2$ and $3$.
\end{proof}

\begin{claim}{D.20}{Proved}{Recognition of the closed pro-$p$ extension}
Let $P$ be the pro-$p$ completion of a closed orientable surface group of genus $g\geq2$. A torsion-free pro-$p$ group $\mathcal E$ containing $P$ as an open normal subgroup is the pro-$p$ completion of a closed surface group. The surface is orientable at odd $p$. At $p=2$ a nonorientable surface is allowed, and its orientation character is the action on the orientation line of $P$.
\end{claim}
\passage{28; 178 verso; 353; 24; 147 verso; 293}

\begin{proof}
We give the absolute finite-index duality calculation. The orientable surface presentation makes $P$ a Poincar\'e duality pro-$p$ group of dimension two, with
\[
 H^i(P,\Z_p[[P]])=
 \begin{cases}\Z_p,&i=2,\\0,&i\ne2,\end{cases}
\]
and trivial right $P$-action on the top group. This is the trivial-character case of the Demushkin classification in \cite[Theorem 1.9]{R-Wilkes}.

Serre's theorem \cite[\S1, Theorem]{R-SerreCD} equates the $p$-cohomological dimensions of an open subgroup and a profinite group without elements of order $p$. Therefore $\operatorname{cd}_p(\mathcal E)=2$. Wilkes \cite[Proposition 1.26]{R-WilkesRelative} gives finite-index ascent and descent of type $p$-$FP_\infty$ for profinite groups. The surface resolution gives this property for $P$, and hence for $\mathcal E$.

At finite index, induction and coinduction of compact modules agree, and continuous Shapiro gives
\[
\begin{split}
 H^i(\mathcal E,\Z_p[[\mathcal E]])
 &=
 H^i\!\left(\mathcal E,
        \operatorname{coind}_{P}^{\mathcal E}\Z_p[[P]]\right)\\
 &\simeq H^i(P,\Z_p[[P]]).
\end{split}
\]
The first coefficient isomorphism and compact Shapiro are \cite[Propositions 1.16 and 1.21, and the remark after 1.21]{R-WilkesRelative}. We make their module compatibility explicit. Put $\Lambda_H=\Z_p[[H]]$ and let $\operatorname{pr}_P:\Lambda_{\mathcal E}\to\Lambda_P$ retain the $P$-coset component. The coefficient isomorphism is
\[
 \Lambda_{\mathcal E}\longrightarrow\operatorname{Coind}_P^{\mathcal E}\Lambda_P,
 \qquad \lambda\longmapsto \bigl(x\mapsto\operatorname{pr}_P(x\lambda)\bigr).
\]
Here coinduction consists of left-$P$-linear functions and its left $\mathcal E$-action is $(e f)(x)=f(xe)$. For representatives $t$ of the finite set $P\backslash\mathcal E$, the inverse is $f\mapsto\sum_t t^{-1}f(t)$. These continuous maps are left-$\mathcal E$-linear and right-$P$-linear. Evaluation at $1$ in Shapiro is natural for the right-$P$ coefficient endomorphisms. Consequently the top right duality character restricts to that of $P$. The cohomology criterion for Poincar\'e duality now makes $\mathcal E$ a duality pro-$p$ group of dimension two with dualizing module $\Q_p/\Z_p$. Its orientation character $\chi:\mathcal E\to\Z_p^\times$ has finite $p$-group image, trivial at odd $p$ and contained in $\{\pm1\}$ at $p=2$.

Its one-dimensional degree-two cohomology and nondegenerate degree-one cup product make it an infinite Demushkin group. The classification \cite[Theorem 1.9]{R-Wilkes} gives, for these characters, exactly the three words of D.19 with the boundary letters omitted. The same invertible free-basis change converts the even-rank nonorientable word to the surface form. Hence an actual closed surface presentation induces a continuous isomorphism of its pro-$p$ completion with $\mathcal E$. Finite-index multiplicativity for the Euler characteristic of the finite two-dimensional resolution gives
\[
 \chi(\mathcal E)=\frac{2-2g}{[\mathcal E:P]}<0.
\]
The resulting surface is hyperbolic, and conjugation-naturality of the dualizing module gives the stated orientation.
\end{proof}

\begin{claim}{D.21}{Refuted}{A full profinite normalizer shortcut}
The assertion that every torsion-free profinite group with an open subgroup $\widehat\Z$ is itself $\widehat\Z$ is false. The pro-$p$ version used in D.19 is proved in D.18.
\end{claim}
\passage{28; 178 verso; 353}

\begin{proof}
Define
\[
 N=(\Z_3\rtimes\Z_2)\times\prod_{\ell\ne2,3}\Z_\ell,
 \qquad (a,b)(a',b')=(a+(-1)^{\,b\bmod2}a',\,b+b')
\]
on the first two factors, and use addition on the remaining factors.
This is a profinite group with the open normal subgroup
\[
 L=(\Z_3\times2\Z_2)\times\prod_{\ell\ne2,3}\Z_\ell
       \simeq\widehat\Z
\]
of index two. A finite-order element projects trivially to the torsion-free quotient
$\Z_2\times\prod_{\ell\ne2,3}\Z_\ell$, and then lies in the torsion-free kernel $\Z_3$; it is therefore the identity. Yet conjugation by $(0,1)$ sends $(1,0)$ to $(-1,0)$, so $N$ is nonabelian. The code in \texttt{code/root/finite\_checks.py} checks its quotient
$C_3\rtimes C_4$, the index-two cyclic subgroup of order six, and this noncommutation. The proof in the inverse limit establishes torsion-freeness; no finite quotient is claimed to be torsion-free.

The multiplier on $L$ is $-1$ in its $3$-adic coordinate and $+1$ in every other coordinate. Thus passing separately to pro-$p$ groups loses precisely this interaction between two different primes. D.29 uses this $N$ as an explicit retract in its full peripheral counterexample. Its cusp normalizers are all procyclic, so that example does not depend on making $N$ itself a cusp normalizer.
\end{proof}

\begin{claim}{D.22}{Refuted}{Full profinite recognition without an invariant discretization}
The assertion that every torsion-free finite extension of a full
profinite hyperbolic surface group with preserved loop structure
is a completed surface pair is false.  D.29 constructs a
counterexample with kernel of type $(0,4)$ and quotient $C_2$,
including the exact peripheral comparisons and a finite quotient
excluding every free profinite presentation. D.30 constructs a
counterexample with closed genus-two kernel and quotient $C_2$;
its degree-two cohomology at two odd primes excludes every
surface completion. Recognition with an
invariant discretization holds by D.17.  The separate pro-$p$
recognition theorems, with their precise group and peripheral
hypotheses, are proved in D.18--D.20.
\end{claim}
\passage{28; 178 verso; 353; 24; 147 verso; 293}

\begin{proof}
The basis-derived four-cusp presentation, faithful outer involution,
continuous twisted extension and torsion-free retraction are given
in D.29.  Its order-eighteen quotient needs three generators, whereas
an index-two extension of a free rank-three group would have rank two
if it were free.  The exact finite-cover comparison excludes closed
surface completions.  Thus the counterexample satisfies the full
peripheral hypothesis and fails the claimed conclusion.

For the closed case, D.30 gives the extension by its continuous
twisted multiplication and by an equivalent profinite presentation.
Its torsion-free maximal pro-$2$ quotient detects the parity of
every element, proving that every torsion element lies in its
torsion-free surface kernel. The two restriction isomorphisms give
$H^2(\mathcal E;\mathbf F_3)=\mathbf F_3$ and
$H^2(\mathcal E;\mathbf F_5)=0$. The cohomology of every orientable,
nonorientable, or punctured surface completion is computed in
D.30 and excludes this pair. This settles the closed hyperbolic
assertion in its smallest genus as well.

This also identifies the failed step in the earlier recognition
attempt.  Separate pro-$p$ comparisons do not produce a single
full profinite surface presentation.  The mixed-prime multiplier
in D.29 is an actual finite-order loop automorphism, and all its
cusp normalizers are procyclic.  The preserved global surface
relation therefore does not restore the proposed conclusion.
The first possible punctured nonabelian rank is three: D.27 excludes
proper torsion-free finite extensions of $\widehat F_2$.
The explicit free-basis comparison in D.29 settles the formerly
undecided peripheral marking on D.26's rank-three extension.

We checked Collas--Philip--Yamaguchi \cite{R-CPY}, Definition 2.1,
Theorems 1.3 and 2.4, against the supplied source.  Their profinite
F-groups already have the stated orbifold presentation.
Theorem 2.4 constructs a torsion-free open normal subgroup, and
Theorem 1.3 bounds the derived length of a finite quotient in its
nonperfect case.  The maps go from a prescribed presentation to
an open subgroup; they do not assert the reverse reconstruction
from an arbitrary finite extension.  D.29 refutes that reverse
assertion under the present full profinite hypotheses.  The
positive conclusions of D.17--D.20 retain their proved maps and
hypotheses throughout.
\end{proof}


\subsection{A full profinite lifting criterion and nonnilpotent finite tests}

\begin{claim}{D.23}{Proved}{Order-preserving lifting and conjugacy in the full shadow system}
Let $\mathcal I$ be the cofiltered poset of isolated finite-index normal
subgroups $N\unlhd B_3$ contained in the pure braid group $PB_3$, and put
$Q_N=GT(N)$.  The finite groups and their reduction homomorphisms
$R_{N,H}:Q_N\to Q_H$ for $N\subseteq H$ are those of
\cite[Proposition 3.14, Theorem 5.2 and Proposition 5.5]{A-DGShadows}.
We use the continuous isomorphism
\[
 \Psi:GT_0\xrightarrow{\ \sim\ }\varprojlim_{N\in\mathcal I}Q_N,
 \qquad (\lambda,f)\longmapsto
 \bigl([ (\lambda-1)/2,f]_N\bigr)_N.
\]
For a positive integer $d$, write $T_d(Q)=\{q\in Q:q^d=1\}$.
For each $H\in\mathcal I$, the image of $T_d(GT_0)$ in $Q_H$ is exactly
\[
 \bigcap_{\substack{N\in\mathcal I\\N\subseteq H}}
 R_{N,H}\bigl(T_d(Q_N)\bigr).                         \tag{D.23a}
\]
An element of order exactly $d$ in this intersection has a lift of
order exactly $d$. Every element of order exactly $d$ in $GT_0$ has
an image of order exactly $d$ at some isolated level.
For $a,b\in GT_0$, conjugacy in $GT_0$ is equivalent to conjugacy
of $\Psi(a)_N$ and $\Psi(b)_N$ in every $Q_N$.
\end{claim}
\passage{\S{}31; folios 198--199 verso; scans 391--394; \S{}30; folio 194 verso; scan 384; Part D(b) of the task: full GT\_0 torsion}

\begin{proof}
The cited theorem identifies the two profinite groups, including their
topologies. Its Proposition 5.1 makes the coordinates cofinal in the
finite quotients of $\widehat\Z$ and $\widehat F_2$. The map between
the normalization in that paper and ours is exactly
$(\widehat m,\widehat f)\mapsto(2\widehat m+1,\widehat f)$;
Proposition 5.5 checks the equations and the multiplication under this
map. Thus its inverse limit concerns the full $GT_0$ of Part E.

Fix $q\in Q_H$ in the right side of (D.23a). For every isolated
$N\subseteq H$ define the finite nonempty set
\[
 S_N=\{v\in Q_N:v^d=1,\ R_{N,H}(v)=q\}.
\]
Every reduction $R_{N,M}$ restricts to a function $S_N\to S_M$.
The inverse limit of these sets is nonempty. To verify the exact
compactness step, any finite list of coordinates and compatibility
conditions is dominated by one isolated subgroup
$L\subseteq H\cap N_1\cap\cdots\cap N_t$; Proposition 3.14 supplies
an isolated refinement of the intersection. Choose an element of
$S_L$ and project it to the specified coordinates. This satisfies
all the finitely many conditions. The closed conditions in the
compact product $\prod_NS_N$ consequently have the finite
intersection property. A compatible element maps under
$\Psi^{-1}$ to $g\in GT_0$ with $g^d=1$ and $g_H=q$.
The reverse inclusion follows by taking the coordinates of such a
$g$. The levels below $H$ are cofinal, so the constructed family is
a full inverse-limit point.

When $q$ has order $d$, the order of $g$ both divides $d$ and is a
multiple of the order of its image $q$, hence equals $d$.
Conversely, for $g$ of order $d$, choose, for each prime $p\mid d$,
a finite coordinate detecting $g^{d/p}\ne1$. A common isolated
refinement detects all these powers. The order of its image divides
$d$ and divides none of the proper divisors $d/p$, so it is $d$.

For the conjugacy assertion, the sets
\[
 C_N(a,b)=\{v\in Q_N:v\Psi(a)_Nv^{-1}=\Psi(b)_N\}
\]
are finite nonempty sets with restriction maps. The same finite
intersection argument gives a compatible conjugator, which
$\Psi^{-1}$ sends to a conjugator in $GT_0$. Necessity follows by
projection. The construction does not assume surjectivity of a
reduction map or that every element of $Q_N$ comes from $GT_0$.
\end{proof}

\begin{claim}{D.24}{Proved}{Dihedral tests with an order bound and the full dyadic quotient}
The program \texttt{code/part\_a/gt\_shadow\_refinement.py} executes
the finite order-two and order-three tests below. Each entry counts
nonidentity elements with order dividing $d$ at the lower level,
and then those admitting a lift with order dividing $d$ at the
upper level.
\begin{center}
\begin{tabular}{rrrrr}
\toprule
$n$&$q$&$d$&Lower elements&Elements admitting such a lift\\
\midrule
6&18&2&7&7\\
6&18&3&2&0\\
18&54&3&8&0\\
54&162&3&8&0\\
14&98&3&14&14\\
4&8&2&3&1\\
8&16&2&7&2\\
16&32&2&11&4\\
32&64&2&19&8\\
\bottomrule
\end{tabular}
\end{center}
The levels are isolated dihedral shadows $Q_{K^{(n)}}$. Their
subsystem with $n=2^\alpha$, $\alpha\geq2$, has inverse limit
\[
 H_2=\{(k,(-1)^a5^b)\in
 \Z_2\rtimes\Z_2^\times:
 a\in\{0,1\},\ b\in\Z_2,\ k\equiv b\pmod2\}.
                                                               \tag{D.24a}
\]
Every finite subgroup of $H_2$ is trivial or cyclic of order two,
and all its order-two subgroups are conjugate to
$\langle(0,-1)\rangle$. There is a continuous surjection
$GT_0\to H_2$. The family of all dihedral levels is not cofinal
in $\mathcal I$; an explicit $A_5$ quotient below proves this.
\end{claim}
\passage{\S{}31; folios 198--199 verso; scans 391--394; \S{}30; folio 194 verso; scan 384; Part D(b) of the task: finite-quotient torsion program}

\begin{proof}
Write $D_n=\langle r,s\mid r^n=s^2=1,\ srs^{-1}=r^{-1}\rangle$.
The homomorphism
\[
 \psi_n:PB_3=\langle x,y\rangle\times\langle c\rangle
       \longrightarrow D_n^3,
 \quad x\longmapsto(r,s,s),\quad
 y\longmapsto(rs,r,rs),\quad c\longmapsto1
\]
defines $K^{(n)}=\ker\psi_n$. For even $n\geq4$, the explicit
finite model of \cite[Theorem 4.3 and Proposition 4.5]{A-BDHPShadows}
is
\[
 Q_{K^{(n)}}\simeq
 \left\{(k,u)\in(\Z/(n/2))\rtimes(\Z/(2n))^\times:
 \begin{array}{l}
 k\equiv\nu(u)\pmod2\text{ when }4\mid n
 \end{array}\right\},                                  \tag{D.24b}
\]
where $\nu(u)=(u-1)/4$ for $u\equiv1\pmod4$ and
$\nu(u)=(u+1)/4$ for $u\equiv3\pmod4$.
Multiplication is $(k,u)(l,v)=(k+ul,uv)$, with the indicated
moduli. For $n\mid q$, reduction is
$(k\bmod(q/2),u\bmod2q)\mapsto(k\bmod(n/2),u\bmod2n)$.

The code constructs the actual subgroup
$G_n=\langle(r,s,s),(rs,r,rs)\rangle\leq D_n^3$ for
$n=4,6,8,12,14$. It computes its commutator subgroup and exhausts
all pairs $(m,f)$ with $0\leq m<n$, $\gcd(2m+1,n)=1$ and
$f\in G_n'$. It tests the two cycle equations
\[
 \theta(f)f=1,\qquad
 \tau^2(y^mf)\tau(y^mf)y^mf=1,
\]
where $\theta(x)=y$, $\theta(y)=x$, $\tau(x)=y$ and
$\tau(y)=(xy)^{-1}$. It checks descent of these maps, and of
$x\mapsto x^{2m+1}$, $y\mapsto f^{-1}y^{2m+1}f$, on every
labelled edge of the finite Cayley graph, and checks that the
resulting maps are bijective. The accepted pairs are exactly
those in (D.24b), of respective orders $4,12,16,24,84$.
This checks the reduction to the quoted finite model on the
tested groups, with no reliance on floating point arithmetic.
Projection of $G_6,G_{12},G_{14}$ onto $D_3,D_3,D_7$
exhibits finite quotients in which the odd-order rotation and
order-two reflection do not commute; their nilpotence would
force these coprime Sylow subgroups to commute. These tests
therefore include nonnilpotent finite quotients of $\widehat F_2$.

For the refinement counts, the code enumerates every element of
(D.24b), computes its $d$th power using the affine multiplication,
reduces every upper root, and counts the exact fibres. For example,
at $n=6$ the group is
$\operatorname{Aff}(\F_3)\times C_2$, and its two order-three
elements are the nonzero translations. At $q=18$, an element of
order dividing three has multiplier $u\equiv1,4,7\pmod9$ and
translation divisible by $3$. Both its affine coordinates reduce
to the identity modulo $3$. Thus neither order-three element
at $n=6$ has an order-dividing-three lift, although the unbounded
reduction map is surjective. At $n=14$ the multipliers $2,4$
modulo $7$ have order three. Each lifts to its order-three
root modulo $49$, and each of the seven lifts of a translation
coordinate is allowed because $1+u+u^2=0$ there. This gives
seven order-three lifts of each of the fourteen lower elements,
as the exhaustive output records. The same program gives the
remaining rows and the full table in its accompanying CSV.

The compatible affine coordinate reductions identify the dyadic
inverse limit with (D.24a): the residue of $b$ modulo $2$ is the
condition in (D.24b), and the inverse limits of the two coordinate
rings are $\Z_2$ and $\Z_2^\times$. The surjection from $GT_0$
is \cite[Theorem 5.7 and Proposition 5.8]{A-BDHPShadows}; more
precisely, that theorem proves surjectivity already after
restriction along $G_{\Q}\to GT_0$.

The translation kernel $\Z_2$ is torsion-free, so a finite
subgroup of $H_2$ injects into the finite-order subgroup of
$\Z_2^\times$. This subgroup is $\{1,-1\}$: the parametrization
$\Z_2^\times=\{\pm1\}\times5^{\Z_2}$ has torsion-free second
factor, as follows from
$v_2(5^{2^j}-1)=j+2$ and passage to the inverse limit.
Every nonidentity finite-order element of $H_2$ is consequently
$(k,-1)$ with $k\in2\Z_2$. Put $t=k/2$ and choose
$v=5^{t\bmod2}$. Then $(t,v)\in H_2$, and direct multiplication
gives
\[
 (t,v)(0,-1)(t,v)^{-1}=(2t,-1)=(k,-1).
\]
This proves the full finite-subgroup and conjugacy assertions
for this infinite quotient, rather than extrapolating the
finite dyadic counts.

Finally every $G_n$ is metabelian: the commutator subgroup of
$D_n^3$ lies in its abelian rotation subgroup, and $G_n'$ is
a subgroup of that rotation subgroup. Define the epimorphism
$\eta:F_2\to A_5$ by
$x\mapsto(1\ 2\ 3\ 4\ 5)$ and $y\mapsto(1\ 2\ 3)$.
The code enumerates its image and derived subgroup; both have
order $60$, proving directly that this image is a nontrivial
perfect group. Proposition 5.1 of \cite{A-DGShadows} gives an
isolated $N$ with $N\cap F_2\subseteq\ker\eta$.
An inclusion $K^{(n)}\subseteq N$ would give the surjection
\[
 G_n=F_2/(K^{(n)}\cap F_2)
   \longrightarrow F_2/\ker\eta=A_5.
\]
Its image would be metabelian, contradicting the computed
nontrivial perfect derived subgroup. Hence no such inclusion
exists. This proves the asserted failure of cofinality by an
explicit quotient and the exact factoring map.
\end{proof}

These tests give explicit finite obstructions for D.12.
D.32 proves, by its full fixed-subgroup and finite-cover argument,
that $GT_0$ has no nonidentity element of odd order and that every
finite subgroup has order at most two, detected by the dyadic
multiplier. D.34 further proves that a nonidentity involution
has the global multiplier $-1$.
The fourteen order-three shadows at $K^{(14)}$ survive the
executed test through $K^{(98)}$. Combining D.32 with D.23
shows that each fails the order-three lifting condition at some
isolated refinement; a common refinement excludes all fourteen.
This existence conclusion does not identify that refinement or
alter the recorded finite counts. D.25 explains why the entire
dihedral subsystem alone cannot supply the exclusion.
The remaining full $GT_0$ question is conjugacy of its
involutions. D.35 gives the exact third-cusp equation and the
finite conjugator criterion, and D.44 supplies a coboundary
satisfying its first symmetry. D.42 also proves the global
multiplier restriction for every finite subgroup of the full
three-cusp loop-outer group.

\begin{claim}{D.25}{Proved}{All-prime dihedral compatibility and the multiplier-one obstruction}
The inverse limit $D$ over all the dihedral levels of D.24 has the
following explicit isomorphism of profinite groups:
\[
 D\xrightarrow{\ \sim\ }
 \left(\prod_{p\text{ odd}}\operatorname{Aff}(\Z_p)\right)
 \times H_2.                                                   \tag{D.25a}
\]
There is a canonical continuous homomorphism $GT_0\to D$.
The group $D$ contains an element of every prime order. In
particular, imposing order three throughout this entire
dihedral system still permits a multiplier whose $7$-adic
coordinate is a primitive cube root and whose other prime
coordinates are $1$.

Put $F=\widehat F_2=\langle x,y\rangle$ and define the continuous
epimorphism and its closed kernel by
\[
 \kappa:F\longrightarrow\prod_pF_2^{(p)},\qquad
 R=\ker\kappa.
\]
Write $\lambda:GT_0\to\widehat\Z^\times$ for the multiplier.
For every integer $d\geq2$, the elements $a\in\ker\lambda$
with $a^d=1$ correspond exactly to the following elements
$f\in R$. The continuous endomorphism
\[
 \alpha_f:F\to F,\qquad x\mapsto x,\quad y\mapsto f^{-1}yf
\]
must be an automorphism and satisfy
\[
 \theta(f)f=1,\qquad
 \tau^2(f)\tau(f)f=1,\qquad
 \prod_{j=0}^{d-1}\alpha_f^{j}(f)=1,                         \tag{D.25b}
\]
where the product is in increasing order of $j$,
$\theta(x)=y$, $\theta(y)=x$, $\tau(x)=y$, and
$\tau(y)=(xy)^{-1}$. Every such finite-order element has identity
image in $D$. The kernel $R$ itself maps onto the explicit
perfect group $A_5$ used in D.24.
The finite-order solutions of (D.25b) in the stated
automorphism domain are exactly $f=1$, by D.32's subsequent
full fixed-subgroup argument.  Thus the nontrivial prime-order
points of the dihedral inverse limit exhibited here do not
lift with that order to $GT_0$ when the prime is odd.
\end{claim}
\passage{\S{}25 bis; folios 157 verso--160; scans 312--317; \S{}30; folio 194 verso; scan 384; \S{}31; folios 198--199 verso; scans 391--394}

\begin{proof}
It suffices to use the even levels $n=2^\alpha n_0$ with
$\alpha\geq1$ and $n_0$ odd; every odd level agrees with its
double. Apply the Chinese remainder isomorphism separately to
the translation coordinate modulo $n/2$ and to the multiplier
coordinate modulo $2n$ in (D.24b). The odd factors give
$\operatorname{Aff}(\Z/n_0)$, and the powers of two give
the finite groups defining $H_2$. These coordinate maps commute
with every reduction $n\mid q$. Conversely, specified odd
coordinates and a specified point of $H_2$ reconstruct every
finite pair $(k,u)$ by the Chinese remainder theorem. They
satisfy the parity condition because that condition involves
only the power-of-two coordinates. The resulting inverse
maps are continuous, being continuous at every finite
coordinate. This proves (D.25a), including its topology.
Restriction of the full inverse-limit point in D.23 defines
the indicated map from $GT_0$; surjectivity onto $D$ is not
used or asserted.

Let $q$ be a prime, and choose a prime divisor $p$ of $2^q-1$.
The residue of $2$ in $\F_p^\times$ has order $q$: its order
divides $q$, and order one would give $p\mid2-1$. Hence
$q\mid p-1$, in particular $p$ is odd and $p\ne q$.
The element $2\bmod p$ lifts to a unique $q$th root
$\zeta\in\Z_p^\times$. Explicitly, after a root $a$ modulo
$p^j$ is chosen, substitute $a+p^jt$, $0\leq t<p$,
in $X^q-1$ modulo $p^{j+1}$. Its coefficient of $t$ is
$p^jqa^{q-1}$, and $qa^{q-1}$ is a unit modulo $p$;
there is exactly one permitted value of $t$.
The compatible limit has order exactly $q$ because its
residue already does. Put $(0,\zeta)$ in the $p$-factor
of (D.25a), and put the identity in all the other factors.
This gives the asserted element of order $q$ in $D$.

For $q=3$, take $p=7$. The executed program
\texttt{code/part\_a/gt\_shadow\_all\_prime.py} constructs ten
successive roots of $X^2+X+1$ starting with $2\bmod7$, and
uses the exact Chinese remainder formulas to construct
compatible shadow coordinates at sixteen levels through
$n=2744$. At $n=14$ its affine pair is $(0,9\bmod28)$;
at $n=98$ it is $(0,177\bmod196)$. Their orders are three
and the latter reduces to the former. The preceding
root-lifting proof constructs the infinite point, while
the program verifies the displayed finite coordinates.
Thus the equations tested over all dihedral levels admit
prime-dependent torsion multipliers. They do not force
the global alternatives $1$ and $-1$.

We next prove the statement about $R$. Each projection
$F\to F_2^{(p)}$ is the canonical maximal pro-$p$ quotient
map. Their product is onto: in any finite collection of
finite $p$-group quotients at distinct primes, its image
surjects to each factor. Its order is consequently divisible
by every factor order, hence by their product, since these
orders are pairwise coprime. The image is the full product.
These finite products form the finite quotients of
$\prod_pF_2^{(p)}$, so the compact image of $F$ is the whole
product. Also $R\subset F'$, since the abelianization
$F^{\mathrm{ab}}=\widehat\Z^2$ is detected by the
abelianizations of all these pro-$p$ quotients.

Let $a=(1,f)\in GT_0$ have finite order dividing $d$.
It induces a finite-order automorphism of each free pro-$p$
quotient with trivial action on its abelianization. The
Magnus leading-term argument in D.10, which also applies
at $p=2$, makes this automorphism the identity. Its normalized
formula then gives
$f_p\in C_{F_2^{(p)}}(y)\cap(F_2^{(p)})'=1$ by A.7 and
the primitive abelianization coordinate of $y$.
Thus $f\in R$.

The two first equations of (D.25b) are precisely the defining
cycle equations at multiplier one; the required invertibility
is stated explicitly. To obtain the last equation, put
$w_d=f\alpha_f(f)\cdots\alpha_f^{d-1}(f)$.
Induction on $d$ gives
\[
 \alpha_f^d(x)=x,\qquad
 \alpha_f^d(y)=w_d^{-1}yw_d.
\]
Therefore $\alpha_f^d=1$ is equivalent to
$w_d\in C_F(y)$. But $w_d\in F'$ and
$C_F(y)\cap F'=\overline{\langle y\rangle}\cap F'=1$,
again by A.7 and abelianization. This proves the necessity
of the last equation and its sufficiency. The normalized
identification of $GT_0$ with these automorphisms gives
the asserted bijection. The kernel of the multiplier on
$D$ is a subgroup of the additive translation group
$\prod_p\Z_p$, which is torsion-free. Consequently every
finite-order multiplier-one element maps to the identity
of $D$.

Finally let $\eta:F\to A_5$ be the epimorphism from D.24.
The quotient $A_5/\eta(R)$ is a finite quotient of
$F/R=\prod_pF_2^{(p)}$ and is therefore nilpotent.
It is also perfect, since the executed permutation
calculation in D.24 proves $A_5'=A_5$. A finite nilpotent
perfect group is trivial: its lower central series would
otherwise be constantly equal to the nontrivial group
and could not terminate. Hence $\eta(R)=A_5$.
This identifies a nontrivial part of the residual group in the
exact equation (D.25b). The full three-cusp argument in D.32
proves that this equation, together with its stated cycle and
invertibility requirements, has only the identity solution.
The surjection $R\to A_5$ and the compatible points of $D$
remain the exact quotient calculations proved above.
\end{proof}

\begin{claim}{D.26}{Refuted}{Unmarked recognition at index two and the peripheral obstruction}
The assertion that a torsion-free profinite group with an open
free profinite subgroup of rank three is free is false, already
at index two. We give the rank-three instance of
\cite[Proposition, pp.~350--352]{A-Haran1985} as an explicit
extension
\[
 1\longrightarrow H=\widehat F(l,a,b)
   \longrightarrow E\longrightarrow C_2\longrightarrow1.
                                                               \tag{D.26a}
\]
It has no section. Its outer involution preserves no loop
structure of type $(1,2)$ on $H$. Every preserved type-$(0,4)$
structure would have a double-transposition cusp permutation,
multiplier $-\epsilon$, and canonical character
$\chi(l)=1$, $\chi(a)=\chi(b)=-1$ sending all four cusp
generators to $-1$. Two explicit type-$(0,4)$ candidates fail,
with finite $S_3$ and $S_4$ certificates. D.29 supplies a
different complete type-$(0,4)$ marking, satisfying exactly
these necessary conditions, and thereby turns this extension
into a counterexample to the loop-preserving assertion D(d).
\end{claim}
\passage{\S{}28; folios 178--179; scans 352--354; Part D(d) of the task: full profinite recognition}

\begin{proof}
Choose $\epsilon,\kappa\in\widehat\Z$ by their prime
coordinates:
\[
 \epsilon_3=-1,\quad\epsilon_p=1\ (p\ne3),\qquad
 \kappa_2=1,\quad\kappa_p=0\ (p\ne2).
\]
Thus $\epsilon$ is a unit, $\epsilon^2=1$ and
$\epsilon\kappa=\kappa$. Define the continuous endomorphism
of the free profinite group $H$ by
\[
 \alpha(l)=l^\epsilon,\qquad
 \alpha(a)=b,\qquad
 \alpha(b)=l^\kappa a l^{-\kappa},\qquad c=l^\kappa.
                                                               \tag{D.26b}
\]
Substitution on the three free generators gives
$\alpha^2=\operatorname{Inn}(c)$ and $\alpha(c)=c$.
In particular $\operatorname{Inn}(c^{-1})\alpha$ is an
inverse, so $\alpha$ is an automorphism.

Define $E$ on the compact space $H\times\{0,1\}$ by
\[
 (h,i)(h',j)=
 \bigl(h\alpha^i(h')c^{\lfloor(i+j)/2\rfloor},
             i+j\bmod2\bigr).
                                                               \tag{D.26c}
\]
The two identities just checked prove associativity by
the eight triples of quotient labels. When the first two
labels are $1$, the required interchange is
$\alpha^2(h'')c=ch''$; the triple $(1,1,1)$ additionally
uses $\alpha(c)=c$. The other triples follow directly.
The identity and inverses follow by the same formulas.
All operations are continuous, so this is a profinite group.
Writing $t=(1,1)$ gives $tht^{-1}=\alpha(h)$ and $t^2=c$.
The projection to the second coordinate gives (D.26a).

The closed subgroup $N=\overline{\langle l,t\rangle}$
has the explicit isomorphism
\[
 N\simeq (\Z_3\rtimes\Z_2)
          \times\prod_{p\ne2,3}\Z_p,
\]
where the parity of the $\Z_2$ coordinate acts by inversion
on $\Z_3$. Namely, the normal form $l^u t^i$ maps to
$(u_3,2u_2+i,(u_p)_{p\ne2,3})$.
This is a continuous bijection preserving the displayed
relations and multiplication; its inverse reads the parity
$i$ and then all coordinates of $u$.
The projection of $N$ to
$\Z_2\times\prod_{p\ne2,3}\Z_p$ has torsion-free kernel
$\Z_3$ and torsion-free target, hence $N$ is torsion-free.

There is a continuous retraction $E\to N$ given on the
generators by $l\mapsto l$, $t\mapsto t$, $a,b\mapsto1$;
(D.26b) verifies its compatibility with the group law.
An element of finite order maps trivially under this
retraction and therefore has trivial image in $C_2$.
It lies in $H$, which is torsion-free by A.7, so it is
the identity. Thus $E$ is torsion-free and the extension
has no section. The induced outer action is faithful:
its nonidentity element acts nontrivially on $H^{\mathrm{ab}}$.

For comparison with Haran's construction, the homomorphism
$N\amalg\widehat\Z\to E$ from the free profinite product,
identity on $N$ and sending the other generator to $a$,
is onto since $b=tat^{-1}$. In the opposite direction send
$l,a,b,t$ to $l,a,tat^{-1},t$ in that free product.
Equations (D.26b) hold there, using the two relations in
$N$, so (D.26c) gives a continuous homomorphism in the
opposite direction. The composites fix every generator.
These are mutually inverse isomorphisms.

To prove that $E$ is not free, form the finite group
\[
 Q=(\F_3^2)\rtimes C_2,
 \qquad (v,i)(w,j)=(v+(-1)^iw,i+j).
\]
The formulas $l\mapsto(e_1,0)$, $a\mapsto(e_2,0)$,
$t\mapsto(0,1)$, and $b\mapsto(-e_2,0)$ give an
epimorphism $E\to Q$: $\epsilon\equiv-1\pmod3$ and
$\kappa\equiv0\pmod3$ verify all relations.
This group requires three generators. Two elements of
$\F_3^2$ do not generate its $C_2$ quotient. When one
of two proposed generators lies outside $\F_3^2$, it is
an involution. Multiplying the other by it when necessary
puts that other generator in $\F_3^2$. Conjugation then
only inverts this one vector, so the generated subgroup
has order at most six. Three generators, as displayed,
do generate $Q$.
Were $E$ a free profinite group of rank $r$, its open
index-two subgroup would have rank $2r-1$ by the finite
coset Schreier rewriting. Since $H$ has rank three,
this would give $r=2$, contradicting its quotient $Q$.
The program \texttt{code/part\_a/haran\_rank3.py}
checks all $18^3$ associativity instances and all $18^2$
ordered pairs in $Q$. Their generated subgroup orders
are $1,2,3,6,9$, with multiplicities $1,27,32,216,48$;
none is $18$.

We now check the peripheral condition rather than
discarding it. On the basis $[l],[a],[b]$ of
$H^{\mathrm{ab}}=\widehat\Z^3$, the outer involution has
matrix and determinant
\[
 M=\begin{pmatrix}\epsilon&0&0\\0&0&1\\0&1&0\end{pmatrix},
 \qquad\det M=-\epsilon.
                                                               \tag{D.26d}
\]
This determinant has value $+1$ at the prime $3$ and
$-1$ at every other prime, and is invariant under a
change of basis. For a preserved type-$(1,2)$ peripheral
structure, let $W$ be the primitive rank-one cusp
submodule in $H^{\mathrm{ab}}$. Write $\lambda$ for the
action on the orientation line and $\delta\in\{1,-1\}$
for the sign of the permutation of the two punctures.
The action on $W$ is multiplication by $\lambda\delta$.
The quotient $H^{\mathrm{ab}}/W$ is the first homology
of the filled torus. The orientation comparison of A.2
identifies its determinant with $\lambda$.
The square of the outer action is the identity, so
$\lambda^2=1$. Multiplicativity of determinant in the
exact sequence of these free modules gives
\[
 \det M=(\lambda\delta)\lambda=\delta.
\]
This contradicts (D.26d). Thus no such type-$(1,2)$
peripheral structure is preserved.

For type $(0,4)$ this determinant argument permits a
prime-dependent orientation multiplier. We test the
specific free basis $l,a,b^\epsilon$, which is a basis
because the power substitution on $b$ has inverse
$b\mapsto b^\epsilon$. Set
\[
 P=lab^\epsilon,
 \quad L_1=\overline{\langle l\rangle},\quad
 L_2=\overline{\langle a\rangle},\quad
 L_3=\overline{\langle b^\epsilon\rangle},\quad
 L_4=\overline{\langle P^{-1}\rangle}.
\]
This is a full type-$(0,4)$ peripheral system. The first
three classes transform as required with multiplier
$\epsilon$ and permutation $(2\ 3)$. The fourth would
require $\alpha(P)$ to be conjugate to $P^\epsilon$.
The actual words are
\[
 \alpha(P)=l^\epsilon b l^\kappa a^\epsilon l^{-\kappa},
 \qquad P^\epsilon=(lab^\epsilon)^\epsilon.
\]
Map the free group $H$ continuously onto $S_3$ by
\[
 l\longmapsto(1\ 2),\qquad
 a\longmapsto(2\ 3),\qquad
 b\longmapsto(1\ 2\ 3).
\]
At its exponent six, $\epsilon\equiv5$ and
$\kappa\equiv3$. Direct multiplication gives
$P\mapsto1$, $\alpha(P)\mapsto(1\ 2\ 3)$,
and $P^\epsilon\mapsto1$. These images are not conjugate,
so the fourth loop class fails. The program checks these
values and all $6^3$ possible assignments, finding
twenty-four nonconjugacy certificates.

There is a basis-independent restriction on any remaining
type-$(0,4)$ structure. Write $\sigma\in S_4$ for its cusp
permutation and $\lambda\in\widehat\Z^\times$ for its
multiplier. The square of the outer action is the identity,
so $\sigma^2=1$ and $\lambda^2=1$. Its action on
$H^{\mathrm{ab}}$ is $\lambda$ times the action of
$\sigma$ on $\widehat\Z^4/\widehat\Z(1,1,1,1)$.
Comparing traces with (D.26d) gives
\[
 \epsilon=\lambda\bigl(|\operatorname{Fix}(\sigma)|-1\bigr).
                                                               \tag{D.26e}
\]
The identity permutation would give $3\lambda=\epsilon$,
which is impossible modulo $3$. A single transposition
forces $\lambda=\epsilon$; a double transposition forces
$\lambda=-\epsilon$.

For each such peripheral structure there is a canonical
character $\chi:H\to\{1,-1\}$ which is $-1$ on all
four cusp generators. The product relation makes it
well-defined, and a free peripheral basis makes it
nonzero. It is preserved by $\alpha$, since $\lambda$
is odd and the cusp classes are permuted. In the basis
$l,a,b$ the three nonzero invariant characters are exactly
\[
 \chi_1=(-1,1,1),\qquad
 \chi_2=(1,-1,-1),\qquad
 \chi_3=(-1,-1,-1).                               \tag{D.26f}
\]
We compute the action on continuous cohomology with
coefficients $\Q_p(\chi)$ for an odd prime $p$; equivalently
one may use the inverse limit of the finite coefficient
modules and then tensor with $\Q_p$. A crossed
homomorphism $u$ is determined freely by
$(L,A,B)=(u(l),u(a),u(b))$. Put
$\eta=\chi(l)$ and $\zeta=\chi(a)=\chi(b)$, and define
\[
 e=\begin{cases}\epsilon_p&\eta=1,\\1&\eta=-1,\end{cases}
 \qquad
 k=\begin{cases}0&\eta=1,\\1&\eta=-1.\end{cases}
\]
Here $u(l^\epsilon)=eL$ and $u(l^\kappa)=kL$.
For $\eta=1$ this is multiplication by the $p$-coordinate
of the exponent; for $\eta=-1$ the geometric sum is
$(1-(-1)^v)/2=1$ for either odd exponent
$v=\epsilon,\kappa$. Substitution in (D.26b) gives
\[
 \alpha^*u\quad\longleftrightarrow\quad
 A_{\chi,p}\begin{pmatrix}L\\A\\B\end{pmatrix},
 \qquad
 A_{\chi,p}=
 \begin{pmatrix}
 e&0&0\\0&0&1\\(1-\zeta)k&\eta&0
 \end{pmatrix}.                                 \tag{D.26g}
\]
The coboundary line is spanned by
$(\eta-1,\zeta-1,\zeta-1)$, and this nonzero vector
is fixed by the displayed matrix. Thus
\[
 \det\bigl(\alpha^*|H^1(H,\Q_p(\chi_i))\bigr)
 =\begin{cases}
 1&i=1,3,\\-\epsilon_p&i=2.
 \end{cases}                                    \tag{D.26h}
\]

We identify the determinant required by the peripheral
structure through actual restriction and filling maps.
Let $K=\ker\chi$. Under the specified marking it is the
completion of the double cover of the four-punctured
sphere in which every cusp has nontrivial monodromy.
The cover has four punctures and genus one, since its
Euler characteristic is $-4$. The deck involution fixes
each of the four peripheral homology classes. Their
span is the three-dimensional invariant subspace of
$H_1(K,\Q_p)$: restriction and averaging identify the
invariant cohomology with $H^1(H,\Q_p)$, of dimension
three. Filling the four punctures therefore identifies
the anti-invariant part with the first homology of the
filled torus. The restriction of crossed homomorphisms
induces the precise comparison
\[
 H^1(H,\Q_p(\chi))
   \xrightarrow{\ \sim\ }H^1(K,\Q_p)^-
   \xleftarrow{\ \sim\ }
          H^1\bigl(\widehat{\pi_1(T^2)},\Q_p\bigr).
                                                               \tag{D.26i}
\]
The first isomorphism follows by averaging over the
order-two deck group, which is invertible in $\Q_p$;
the second is the dual of the filling quotient on
the anti-invariant part. Both commute with the induced
action of $\alpha$. Upstairs cusp generators are squares
of the downstairs generators. A conjugator in $H$ can
be moved into $K$ by multiplying it by the corresponding
downstairs cusp generator when necessary; that generator
centralizes its own square. Thus $\alpha$ has the same
multiplier $\lambda$ on this covering pair. The transfer
of the relative orientation class, followed by filling,
is the comparison of A.2; it identifies the determinant
on the filled torus with $\lambda_p$. Consequently the
determinant in (D.26h) must be $\lambda_p$.

At $p=3$ all three determinants in (D.26h) are $1$.
This excludes the single-transposition case, which by
(D.26e) would have $\lambda_3=\epsilon_3=-1$.
For the remaining double transposition,
$\lambda=-\epsilon$; at $p=5$ its value is $-1$.
Formula (D.26h) then excludes $\chi_1$ and $\chi_3$,
leaving exactly $\chi_2$. This proves the necessary
conditions stated at the beginning for every marking.

One can satisfy the first two cusp conditions in this
remaining case by the explicit free basis
\[
 x=a,\qquad y=b^{-\epsilon},\qquad z=b^\epsilon l,
 \qquad w=(xyz)^{-1}=(al)^{-1}.                    \tag{D.26j}
\]
The inverse substitution is $a=x$, $b=y^{-\epsilon}$,
$l=yz$, proving that this is a free basis.
It has canonical character $\chi_2$, and
$\alpha(x)$ is $y^{-\epsilon}$ while $\alpha(y)$
is conjugate to $x^{-\epsilon}$. For permutation
$(x\ y)(z\ w)$ and multiplier $-\epsilon$, the next
condition would require
\[
 \alpha(z)\sim w^{-\epsilon},\quad\hbox{equivalently}\quad
                 a^\epsilon l^\epsilon\sim(al)^\epsilon.
                                                               \tag{D.26k}
\]
Indeed $\alpha(z)=c a^\epsilon c^{-1}l^\epsilon$
is conjugate to $a^\epsilon l^\epsilon$ because
$c=l^\kappa$ commutes with $l$. In $S_4$ take
$a=(2\ 3\ 4)$, $l=(1\ 2\ 3\ 4)$ and $b=1$.
Its exponent is twelve, with $\epsilon\equiv5$.
The left member of (D.26k) is the transposition $(1\ 4)$,
whereas the right member is the four-cycle
$(1\ 3\ 2\ 4)$. They cannot be conjugate.
The supplied program checks all $24^2$ pairs in $S_4$,
finding $96$ such nonconjugacy witnesses. It also checks
all six matrices (D.26g) over $\F_3$ and $\F_5$, their
fixed coboundary lines and the determinants on their
two-dimensional cohomology quotients.

These calculations reduce the permissible type-$(0,4)$
markings to double transpositions with the specified
multiplier and character. The second candidate (D.26j)
is excluded, but D.29 constructs another marking by
gluing a genuine three-cusp automorphism to its inverse.
The exact basis change $b'=\tau(a)$ in that construction
identifies its extension with (D.26a). Thus an invariant
marking does exist; the finite failures above concern
the two displayed candidates and the proved exclusions
of the other permutation and character types.
\end{proof}

\begin{claim}{D.27}{Proved}{The full-profinite rank divisibility and prime-order lifting test}
For a torsion-free profinite extension
\[
 1\longrightarrow\widehat F_r\longrightarrow E
   \longrightarrow G\longrightarrow1,
 \qquad 2\le r<\infty,\quad |G|<\infty,
\]
the integer $|G|$ divides $r-1$. In particular a torsion-free
finite extension of $\widehat F_2$ equals its prescribed kernel.
Every homomorphism $C_p\to\Out(\widehat F_2)$, for any prime $p$,
has a homomorphic lift to $\Aut(\widehat F_2)$; a faithful one has
a lift of order exactly $p$. This last assertion holds in the full
profinite category, without an invariant discretization and without
first passing to a pro-$p$ quotient.
\end{claim}
\passage{\S{}28; folios 178--179; scans 352--354; \S{}31; folios 198--199 verso; scans 391--394}

\begin{proof}
We give the normal-subgroup version of Jarden's index argument
\cite[Theorems 1.2 and 4.1, pp.~496--498]{R-Jarden1982}, with the
quotient comparisons written out. Put $F=\widehat F_r$. For every
prime $q$, $\operatorname{cd}_q(F)=1$. The torsion-free finite-index
dimension theorem of Serre \cite[\S1]{R-SerreCD}, whose torsion
hypothesis applies to $E$ itself, gives $\operatorname{cd}_q(E)=1$.
The same is true of every open subgroup of $E$. Thus these groups
are projective profinite groups. Here we use the precise
projectivity criterion recalled in \cite[\S2]{R-Jarden1982}:
a profinite group is projective exactly when its cohomological
dimension at every prime is at most one; equivalently, all its
finite embedding problems have weak solutions. This is the
criterion for the full profinite category, including finite
embedding problems whose kernels are nonabelian.

Fix a prime divisor $p$ of $|G|$, and let $p^a$ be its exact
$p$-part. Choose a Sylow $p$-subgroup $G_p\le G$ and write
$E_p$ for its inverse image in $E$. Define $O^p(J)$, for a
profinite group $J$, as the kernel of its maximal pro-$p$ quotient.
The characteristic subgroup $O^p(F)$ of $F$ is normal in $E_p$.
There are exact continuous quotient maps
\[
 1\longrightarrow F/O^p(F)\longrightarrow E_p/O^p(F)
 \longrightarrow G_p\longrightarrow1.
\]
The middle group is pro-$p$, since both its kernel and quotient
are pro-$p$. Therefore $O^p(E_p)\le O^p(F)$. Conversely the
restriction to $F$ of the maximal pro-$p$ quotient of $E_p$ is a
pro-$p$ quotient of $F$, so $O^p(F)\le O^p(E_p)$. This proves
\[
 O^p(E_p)=O^p(F),\qquad
 [E_p^{(p)}:F^{(p)}]=p^a.                         \tag{D.27a}
\]

The quotient $E_p^{(p)}$ is projective as a pro-$p$ group. Indeed a
finite embedding problem with $p$-group target and covering group
pulls back to one for $E_p$; its projective solution factors through
$E_p^{(p)}$ because its target is a $p$-group. A projective pro-$p$
group is free pro-$p$. It has finite rank $d_p$ because $E_p$ is
finitely generated. The subgroup $F^{(p)}$ in (D.27a) is free
pro-$p$ of rank $r$, with its full induced topology, as the exact
quotient construction proves. Finite-coset Schreier rewriting
therefore gives
\[
                         r-1=p^a(d_p-1).          \tag{D.27b}
\]
Consequently each exact prime-power divisor $p^a$ of $|G|$ divides
$r-1$. These divisors are pairwise coprime, proving the first
assertion. At $r=2$, this gives $|G|=1$.

Now let $\rho:C_p\to\Out(\widehat F_2)$ be a prescribed continuous
outer action, with the finite group given its discrete topology.
Form the pullback
\[
 E_\rho=\{(a,g)\in\Aut(\widehat F_2)\times C_p:
                         [a]=\rho(g)\}.
\]
This is a profinite group. Centre-freeness identifies its kernel
over $C_p$ with $\Inn(\widehat F_2)=\widehat F_2$ by the injective
map $x\mapsto(\Inn(x),1)$. If $E_\rho$ were torsion-free, the
already proved rank-two assertion would force its quotient to be
trivial, contradicting $|C_p|=p$. Hence $E_\rho$ contains a
nonidentity torsion element $v$. Its finite cyclic subgroup
intersects the torsion-free kernel trivially, so projection embeds
it in $C_p$. Its order is therefore exactly $p$, and it projects
isomorphically onto $C_p$. The inverse of that projection is a
section with the prescribed labels. Composing with the first
projection $E_\rho\to\Aut(\widehat F_2)$ gives the asserted lift.
For faithful $\rho$, the lift is faithful because its outer
projection is. No choice of conjugacy class of this section is
claimed; the argument proves existence.

Applied to a prime-order element of $GT_0$, this gives a genuine
finite-order loop-preserving automorphism representing its outer
class. It does not impose an unproved restriction on its multiplier,
or identify its fixed subgroup. In particular it neither solves
the residual norm equation D.25 nor proves conjugacy to the
geometric involution. It does prove that the nontrivial
torsion-free recognition test of D.22 first reaches free rank
three, where (D.27b) permits index two and D.26 supplies an
explicit unmarked example.
\end{proof}

\begin{claim}{D.28}{Proved}{Prime-order sections and the full fixed subgroup}
Let $p$ be a prime and let $\rho:C_p\to GT_0$ be a continuous
homomorphism whose multiplier has $p$-adic component $1$.
This condition is automatic for odd $p$. Let $\alpha$ be the
normalized automorphism of $F=\widehat F(x,y)$ associated to
the chosen generator of $C_p$, and put
\[
 E=F\rtimes_\alpha C_p,\qquad S=\{1\}\times C_p,
 \qquad C=C_F(S).
\]
There is exactly one $F$-conjugacy class of sections of
$E\to C_p$. Each of the three peripheral normalizers contains
a section. The inclusion $C\hookrightarrow F$ induces an
isomorphism of maximal pro-$p$ quotients
\[
                         C^{(p)}\xrightarrow{\sim}F^{(p)}.
                                                               \tag{D.28a}
\]
This assertion concerns the maximal pro-$p$ quotient of the
actual closed fixed subgroup. It does not identify that
subgroup with $F$. An explicit order-three automorphism below
has a proper fixed subgroup and exactly the isomorphism
(D.28a), while its third peripheral class fails in a finite
dihedral quotient.
\end{claim}
\passage{\S{}31; folios 198--199verso; scans 391--394; \S{}30; folio 194verso; scan 384; Part D(a) of the task: actual full GT\_0 torsion}

\begin{proof}
The normalized map $GT_0\to\Aut(F)$ is a homomorphism:
for the pair $(\lambda,f)$ it sends
$x\mapsto x^\lambda$, $y\mapsto f^{-1}y^\lambda f$, and
composition gives precisely the product $f\alpha(g)$ in
the second coordinate. Consequently $\alpha^p=1$.
On the maximal pro-$p$ quotient its multiplier is $1$;
the Magnus argument of D.10 and D.25 makes this finite-order
normalized automorphism the identity. Thus the quotient of
the semidirect product is the explicit map
\[
 \pi:E\longrightarrow F^{(p)}\times C_p,
                  \qquad(h,s^j)\longmapsto(\bar h,s^j).
                                                               \tag{D.28b}
\]
For odd $p$, a finite-order unit in $\Z_p^\times$ of order
dividing $p$ is $1$: reduction modulo $p$ eliminates the
prime-to-$p$ part, and the logarithm on $1+p\Z_p$ eliminates
the remaining torsion. This explains the automatic condition.

Write $c_1=x$, $c_2=y$, $c_3=(xy)^{-1}$, and
$L_i=\overline{\langle c_i\rangle}$. Preservation of the
three peripheral classes gives an exact sequence
\[
 1\longrightarrow L_i\longrightarrow N_E(L_i)
                \longrightarrow C_p\longrightarrow1;
                                                               \tag{D.28c}
\]
its kernel is $L_i$ by A.7. Choose $a_i$ over the specified
generator of $C_p$. Then
\[
 a_i c_i a_i^{-1}=c_i^\lambda,\qquad
 a_i^p=c_i^{m_i},\qquad (\lambda-1)m_i=0.
                                                               \tag{D.28d}
\]
The first component of $\pi(a_i)$ centralizes $\bar c_i$,
since $\lambda_p=1$. The centralizer of this primitive
element in the free pro-$p$ group is
$\overline{\langle\bar c_i\rangle}$. Hence
$\pi(a_i)=(\bar c_i^{v_i},s)$ for $v_i\in\Z_p$, and
comparison of $p$th powers yields
\[
                            (m_i)_p=pv_i.
                                                               \tag{D.28e}
\]
Let $N_\lambda=1+\lambda+\cdots+\lambda^{p-1}$.
In the $q$-adic coordinate where $\lambda_q\ne1$,
the equation $\lambda_q^p=1$ and the integral-domain
property give $(N_\lambda)_q=0$ and $(m_i)_q=0$.
Where $\lambda_q=1$, the norm is $p$; division by $p$
is possible for $q\ne p$, and (D.28e) supplies it at
$q=p$. Define $u_i\in\widehat\Z$ coordinatewise by
$(u_i)_q=-(m_i)_q/p$ in the latter coordinates and zero
in the former. This gives
$N_\lambda u_i=-m_i$. The element
$b_i=a_i c_i^{u_i}$ has order $p$, by (D.28d), and
projects to the prescribed generator. Thus every
normalizer in (D.28c) contains a section $S_i$.

Let $e_p\in\widehat\Z$ be the idempotent with
$p$-coordinate $1$ and all other coordinates $0$.
Each $S_i$ centralizes $c_i^{e_p}$. Its image in
$F^{\rm ab}\otimes\F_p$ is the nonzero primitive vector
$[c_i]$. It follows that
\[
       \dim_{\F_p}H^1(C_F(S_i),\F_p)\ge1.
                                                               \tag{D.28f}
\]
This inference uses restriction of a linear functional
on $F^{\rm ab}/p$ that is nonzero on the displayed
element; it does not infer first homology from the mere
presence of a pro-$p$ subgroup.

We now use the finite-cohomology form of the profinite
Brown theorem, quoted as Theorem~2 and applied in
\cite[\S1.1, equations (1)--(6)]{A-KochloukovaZalesskii2007}.
There is no subgroup $C_p\times C_p$ in $E$, because
every finite subgroup injects into $E/F=C_p$.
The Lyndon--Hochschild--Serre spectral sequence and
$\operatorname{cd}_p(F)=1$ give, for every $n>1$,
\[
 \dim_{\F_p}H^n(E,\F_p)\le
             1+\dim_{\F_p}H^1(F,\F_p)=3.
\]
For each subgroup $T$ of order $p$ in $E$, projection
to $C_p$ shows that its normalizer equals its centralizer:
a normalizing element must induce the identity on the
generator of the quotient. Dividing by the unique element
of $T$ with the same quotient gives
\[
                 N_E(T)=C_F(T)\times T.
\]
The Brown restriction map is a dense embedding into the
product of the cohomology of these normalizers. Its
finite-dimensional source makes it an isomorphism and
makes the number of factors finite. Each closed subgroup
$C_F(T)$ has $p$-cohomological dimension at most one.
The direct-product cohomology calculation therefore gives
the exact finite sum and bound
\[
 \sum_{[T]}
    \left(1+\dim_{\F_p}H^1(C_F(T),\F_p)\right)
       =\dim_{\F_p}H^n(E,\F_p)\le3,
                                                               \tag{D.28g}
\]
where $[T]$ runs through $E$-conjugacy classes of order-$p$
subgroups. These are also their $F$-conjugacy classes:
for any such $T$, the decomposition $E=FT$ gives the
same conjugacy orbit under $E$ and $F$.

By (D.28f), two distinct cusp-section classes would
contribute at least four to (D.28g). Thus all three
$S_i$ are conjugate. Conjugating $S_1$ and $S_2$ to a
single representative places conjugates of
$x^{e_p}$ and $y^{e_p}$ in its fixed subgroup.
Their images in $F^{\rm ab}/p$ are still the independent
vectors $[x],[y]$. This gives two independent restricted
linear functionals, so that common factor contributes
at least three. Bound (D.28g) is exhausted. There are
no other classes, and its fixed subgroup has
$\dim H^1(-,\F_p)=2$. Every order-$p$ subgroup projects
isomorphically onto $C_p$ and determines one labelled
section; hence there is exactly one conjugacy class of
sections. In particular these conclusions apply to $S$.

The image of $C$ in $F^{(p)}$ is closed and surjects
on its Frattini quotient by the two independent cusp
vectors, so it is all of $F^{(p)}$. Since $C$ is a
closed subgroup of a free profinite group, it is
projective. Pulling a finite pro-$p$ embedding problem
back to $C$ shows that $C^{(p)}$ is projective pro-$p$,
and hence free pro-$p$. Its rank equals
$\dim H^1(C,\F_p)=2$. The resulting surjection from
a rank-two free pro-$p$ group onto $F^{(p)}$ is an
isomorphism: after identifying free bases it is a
surjective endomorphism of a finitely generated profinite
group, which is injective. This proves (D.28a) in the
claimed direction.

Here is an explicit limitation of that last comparison.
Choose $u\in\widehat\Z^\times$ with $u_q=1$ for
$q\ne7$, and $u_7$ the order-three root with
$u_7\equiv2\pmod7$. The Hensel construction is given
in D.25. The substitutions
\[
                    \beta(x)=x^u,\qquad\beta(y)=y^u
                                                               \tag{D.28h}
\]
have inverse given by exponent $u^{-1}$ and have order
exactly three. The fixed subgroup $C_\beta$ is proper,
since $\beta(x)\ne x$ in the cyclic quotient of order
seven. But it contains $x^{e_3}$ and $y^{e_3}$, so its
map to $F^{(3)}$ is onto. The same Brown bound, now
using this one section and these two independent vectors,
gives $\dim H^1(C_\beta,\F_3)=2$. The preceding
projectivity argument proves
$C_\beta^{(3)}\cong F^{(3)}$, despite $C_\beta\ne F$.

The precise missing cusp condition is detected by
\[
 D_7=\langle r,s\mid r^7=s^2=1, srs=r^{-1}\rangle,
                  \qquad x\mapsto s,\quad y\mapsto sr.
\]
Here $xy\mapsto r$, $u\equiv9\pmod{14}$, and the
two involutions are unchanged under exponent $u$.
Thus $x^u y^u\mapsto r$ whereas $(xy)^u\mapsto r^2$.
The conjugacy class of $r$ consists of $r,r^{-1}$,
so these are not conjugate. Consequently (D.28h) is
not an element of the three-cusp loop-outer group, and
in particular not a $GT_0$ torsion construction. The
program \texttt{code/part\_a/gt\_torsion\_sections.py}
executes this finite certificate and the coordinate
norm-solvability tests used above. D.32 uses the third cusp
in a finite-cover rank argument to strengthen (D.28a) to
$C=F$ under the three-cusp hypotheses, and thereby excludes
the nontrivial prime-order actions considered here.
\end{proof}

\begin{claim}{D.29}{Refuted}{A preserved four-puncture structure on the torsion-free Haran extension}
The full profinite recognition assertion in D(d) is false, including
its peripheral hypothesis.  There is an exact sequence of profinite
groups
\[
 1\longrightarrow H=\widehat\pi_{0,4}
 \longrightarrow E\longrightarrow C_2\longrightarrow1
                                                        \tag{D.29.1}
\]
in which $E$ is torsion-free and its faithful outer action preserves
a complete type-$(0,4)$ loop structure, while $E$ is the completion
of no finite-type surface group, orientable or nonorientable.
The extension is explicitly isomorphic to D.26's group.
Its outer involution has a mixed-prime orientation multiplier and
is conjugate to no subgroup of the geometric completion inside
$\Out_{\rm lac}(\widehat\pi_{0,4})$.
For this involution both the discretization and prediscretization
fixed sets of F.1 are empty.
\end{claim}
\passage{28; 178--179; 352--354; 30; 194--196; 383--387; 31; 198--199 verso; 391--394}

\begin{proof}
Define units and an exponent by their prime coordinates:
\[
 \epsilon_3=-1,\qquad\epsilon_q=1\ (q\ne3),\qquad
 \kappa_2=1,\quad\kappa_q=0\ (q\ne2),\qquad
 \lambda=-\epsilon.
                                                        \tag{D.29.2}
\]
Thus $\epsilon^2=\lambda^2=1$ and
$\epsilon\kappa=\kappa$.  In particular $\lambda_3=1$
and $\lambda_q=-1$ for every other prime.

We first obtain the required pants isomorphism, without any
finite-order requirement on it.  The compatible assignments
$\zeta_n\mapsto\zeta_n^{\lambda\bmod n}$ define an automorphism
of the cyclotomic extension $\Q^{\rm cyc}/\Q$.
An automorphism of an algebraic extension extends to its algebraic
closure: extend its embedding along algebraic generators, or take
a maximal extension of the embedding and adjoin a root of a missing
minimal polynomial to prove maximality gives all of
$\overline\Q$.  We therefore choose an element
$\sigma\in G_{\Q}$ with cyclotomic character $\lambda$.
The Galois action on the three-punctured projective line commutes
outerly with its rational $S_3$ action and preserves each cusp;
on its local inertia it has exponent $\lambda$.
Harbater--Schneps's normalized description
\cite[\S0.2, Main Theorem (a), Theorem A]{R-HS}
therefore supplies an element $(\lambda,f)\in GT_0$, with
$f\in\widehat F(x,y)'$, whose automorphism is
\[
 x\longmapsto x^\lambda,\qquad
 y\longmapsto f(x,y)^{-1}y^\lambda f(x,y).
                                                        \tag{D.29.3}
\]
Its third cusp $(xy)^{-1}$ goes to a conjugate of its
$\lambda$th power.  This is the only property of $f$ needed
below, besides invertibility of the displayed automorphism.
Every choice of the extension $\sigma$ has these properties,
and every one gives the construction and conclusions that follow.

Let $H=\widehat F(l,a,b)$ and define its closed free factors
\[
 P_1=\widehat F(l,a),\qquad P_2=\widehat F(l,b),\qquad
 L=\overline{\langle l\rangle}.
\]
Their shared-factor universal property is
$H=P_1\mathbin{*_{L}}P_2$: homomorphisms out of the two
free factors agreeing on $l$ are exactly arbitrary continuous
assignments of $l,a,b$.  Give $P_1$ the ordered pants basis
$(l^{-1},a)$ and $P_2$ the ordered pants basis $(l,b)$.
Formula (D.29.3), between these bases, defines a continuous
isomorphism $\tau:P_1\to P_2$ with
\[
 \begin{split}
 \tau(l)&=l^\epsilon,\\
 \tau(a)&=f(l,b)^{-1}b^\lambda f(l,b),\\
 \tau(a^{-1}l)&\sim_{P_2}(b^{-1}l^{-1})^\lambda.
 \end{split}                                           \tag{D.29.4}
\]
The symbol $\sim_{P_2}$ means conjugacy by an element of
$P_2$; in particular it specifies a conjugacy in $H$.
Because $\lambda^{-1}=\lambda$, the inverse isomorphism has
the same peripheral multiplier.

Put $c=l^\kappa$ and prescribe
\[
 \alpha|_{P_1}=\tau,\qquad
 \alpha|_{P_2}=\operatorname{Inn}(c)\tau^{-1}.
                                                        \tag{D.29.5}
\]
Both maps send the shared generator $l$ to $l^\epsilon$:
$\tau^{-1}(l)=l^\epsilon$ and $c$ commutes with $l$.
They consequently extend to a continuous endomorphism of
$H$.  Direct substitution on the two free factors gives
\[
                  \alpha(c)=c,\qquad
                  \alpha^2=\operatorname{Inn}(c).      \tag{D.29.6}
\]
For example, on $P_1$ the second iterate is
$\operatorname{Inn}(c)\tau^{-1}\tau$; on $P_2$ it is
$\operatorname{Inn}(\alpha(c))\tau\tau^{-1}$.
These identities also prove invertibility, with inverse
$\operatorname{Inn}(c^{-1})\alpha$.

Here is the full peripheral structure, with its relation:
\[
 \begin{split}
 (u_1,u_2,u_3,u_4)&=(a,\ a^{-1}l,\ b,\ b^{-1}l^{-1}),\\
                         u_1u_2u_3u_4&=1.
 \end{split}                                           \tag{D.29.7}
\]
The first three entries form a free profinite basis, since
$a=u_1$, $l=u_1u_2$, and $b=u_3$ are its continuous
inverse substitutions.  Thus (D.29.7) is precisely a
type-$(0,4)$ presentation, not just four procyclic subgroups
whose generators happen to satisfy a relation.
Equations (D.29.4)--(D.29.5) give
\[
 \alpha(u_i)\sim_H u_{s(i)}^\lambda,
                  \qquad s=(1\ 3)(2\ 4).              \tag{D.29.8}
\]
For the first two entries this is (D.29.4), and for the last
two it is the inverse comparison, followed by conjugation
by $c$.  Hence the action preserves every required loop
source, its conjugacy ambiguity, and the common orientation
multiplier.  It permutes distinct cusp classes nontrivially,
so its outer class is nonidentity.  Equation (D.29.6) gives
that class exact order two.

On the profinite space $H\times\{0,1\}$ define
\[
 (h,i)(k,j)=
   \bigl(h\alpha^i(k)c^{ij},\ i+j\bmod2\bigr).
                                                        \tag{D.29.9}
\]
Associativity follows by the two carry identities in
(D.29.6); checking the binary indices gives precisely
$\alpha^2(k)c=ck$ and $\alpha(c)=c$.
The element $t=(1,1)$ satisfies $t^2=c$ and
$tht^{-1}=\alpha(h)$.  The inclusion $h\mapsto(h,0)$
and parity projection give (D.29.1), with the profinite
product topology and continuous operations.

We make the comparison with D.26 explicit.  Set
$b'=\tau(a)$.  The pair $(l^\epsilon,b')$ is a free basis
of $P_2$, and replacing its first member by its
$\epsilon$th power gives the free basis $(l,b')$.
Thus $(l,a,b')$ is a free basis of $H$, and
\[
 \alpha(l)=l^\epsilon,\qquad
 \alpha(a)=b',\qquad
 \alpha(b')=c a c^{-1}.                               \tag{D.29.10}
\]
The map from D.26's named free basis to $(l,a,b')$, with
$t\mapsto t$, is a continuous isomorphism of its entire
twisted extension onto $E$.  It is an isomorphism of
extensions with identity map on $C_2$.  Formula (D.29.7)
therefore supplies the previously missing preserved
peripheral system on that exact group.  In particular
D.26's necessary double-transposition and double-cover
character conditions are attained.

For clarity, we give the torsion and finite-quotient
arguments also in the present basis.  The retraction
$r:H\to\widehat\Z$ sends $l$ to $1$ and $a,b$ to $0$.
Since both $\tau(a)$ and $\tau^{-1}(b)$ lie in the
normal closures of the respective nonshared generators,
$r\alpha=\epsilon r$.  It extends to a retraction of
$E$ onto the group $N=\widehat\Z\times\{0,1\}$ with law
(D.29.9), action $\epsilon$, and square parameter
$\kappa$.  An element $(z,1)$ in $N$ has square
$((1+\epsilon)z+\kappa,0)$, whose $2$-coordinate is
$2z_2+1$ and is never zero.  Since its open subgroup
$\widehat\Z$ is torsion-free, every torsion element
outside it would have order two; hence $N$ is torsion-free.
A torsion element of $E$ has trivial image in $N$, so
has even parity and lies in the torsion-free $H$.
It is therefore the identity.

Let $V=\F_3^2$, with basis $v_1,v_2$, and send
\[
         l\longmapsto v_1,\qquad
         a\longmapsto v_2,\qquad
         b\longmapsto -v_2.
                                                        \tag{D.29.11}
\]
Here $\lambda_3=1$ and $\epsilon_3=-1$; conjugators
vanish in the abelian target.  The induced action of
$\alpha$ is $-1$ on $V$, and $c$ has image zero because
$\kappa_3=0$.  Sending $t$ to the involution consequently
extends (D.29.11) to a continuous epimorphism
\[
              E\longrightarrow V\rtimes_{-1}C_2.      \tag{D.29.12}
\]
The target has order eighteen and requires three
generators.  Two positive elements do not generate the
quotient $C_2$; one positive and one negative element
generate at most one line of $V$; two negative elements
reduce to the preceding case by taking their product.
The three displayed images do generate it.
The independent program
\path{haran_quotient_independent_audit.py} checks all
$18^3$ associativity cases, all $18^2$ ordered pairs,
all ordered triples, and the actual matrix identity in
(D.29.11).  The ordered-pair subgroup-size histogram is
\[
 \begin{array}{c|rrrrr}
 \text{generated subgroup order}&1&2&3&6&9\\\hline
 \text{ordered pairs}&1&27&32&216&48.
 \end{array}                                          \tag{D.29.13}
\]
These calculations agree with the independently executed
D.26 certificate under (D.29.10).

A free profinite $E$ of finite rank $d$ would have an
index-two free subgroup of rank $1+2(d-1)$.  The actual
rank of $H$ is three, forcing $d=2$, contrary to
(D.29.12).  Thus $E$ is not free.  This excludes all
punctured finite-type surface completions.
For a putative isomorphism from a closed surface
completion to $E$, the inverse image of $H$ is an open
index-two subgroup.  Its intersection with the dense
surface group has index two, and completion of this
intersection is exactly that open subgroup by the
finite-index topology argument of D.1.  It is the
completion of the fundamental group of the corresponding
closed double cover.  Every aspherical closed surface
completion has cohomological $p$-dimension two, whereas
$\widehat F_3$ has dimension one.  The spherical and
projective-plane groups and their completions are finite
and cannot give $H$.  This excludes every closed case
as well.

Finally A.2 supplies the continuous orientation
homomorphism
\[
 \operatorname{or}:\Out_{\rm lac}(H)
                       \longrightarrow\widehat\Z^\times.
\]
It is unchanged by conjugation.  Every discrete geometric
class has value $1$ or the global unit $-1$, and the
closure of the geometric group still has values in that
closed two-element subgroup.  The value of $[\alpha]$
is the mixed unit $\lambda$ in (D.29.2), which belongs
to neither value.  No conjugate of $\langle[\alpha]\rangle$
is therefore contained in the geometric completion.
The transporter formulas of D.16 and F.1 now give both
empty fixed sets claimed above.

The cusp normalizers themselves cause no exception to
this argument.  Since $s$ fixes none of the four cusp
classes, their normalizers in $E$ map trivially to $C_2$;
each equals its normalizer in $H$, namely its own
procyclic loop group by A.7.  The failure is present
even with these procyclic normalizers retained.
\end{proof}

\begin{claim}{D.30}{Refuted}{A closed genus-two counterexample to full profinite recognition}
The assertion that every torsion-free finite extension of a full
profinite closed hyperbolic surface group is a completed surface group
is false.  There is an exact extension
\[
 1\longrightarrow H=\widehat\pi_1(\Sigma_2)
 \longrightarrow E\longrightarrow C_2\longrightarrow1
                                                        \tag{D.30.1}
\]
with $E$ torsion-free and with a faithful outer action, such that
\[
 H^2(E;\mathbb F_3)=\mathbb F_3,\qquad
 H^2(E;\mathbb F_5)=0.                                  \tag{D.30.2}
\]
No profinite completion of a discrete surface group has these two
cohomology groups.  The exact maximal pro-$p$ quotients are
\[
 E^{(2)}\cong\widehat{\pi_1(N_3)}^{(2)},\qquad
 E^{(3)}\cong\mathbb Z_3^2,\qquad
 E^{(p)}\cong\widehat F_2^{(p)}\quad(p\ne2,3),           \tag{D.30.3}
\]
where $N_3$ is the closed nonorientable surface of genus three.
Thus the counterexample also gives the precise maps through which
its separate pro-$p$ surface presentations arise.
Its outer involution is conjugate to no element of the extended
geometric completion. Its fixed discretization and
prediscretization sets are both empty.
\end{claim}
\passage{28; 178 verso; 353; 24; 147 verso; 293; Part D(d) of the task}

\begin{proof}
We use the commutator convention $[x,y]=xyx^{-1}y^{-1}$.
Set
\[
 \epsilon_3=-1,\quad\epsilon_p=1\ (p\ne3),\qquad
 \kappa_2=1,\quad\kappa_p=0\ (p\ne2),\qquad
 \lambda=-\epsilon
                                                        \tag{D.30.4}
\]
in $\widehat{\mathbb Z}=\prod_p\mathbb Z_p$.  These elements
satisfy $\epsilon^2=\lambda^2=1$ and
$\epsilon\kappa=\kappa$.

\emph{An actual punctured-torus isomorphism.}
Take the elliptic curve $Y^2=4X^3-4X$ over $\mathbb Q$,
with its origin removed and its rational Weierstrass tangential
base point.  Nakamura's full profinite construction
\cite[\S1, equation (1.2), and \S2.5, p.~424]{C_Nakamura2013}
gives an actual homomorphism
\[
 G_{\mathbb Q}\longrightarrow\operatorname{Aut}
 \langle x,y,z\mid[x,y]z=1\rangle^{\widehat{\ }},
 \qquad \sigma(z)=z^{\chi(\sigma)}.
                                                        \tag{D.30.5}
\]
Here $\chi$ is the full cyclotomic character.  The local
coordinate at the rational puncture and its Kummer covers identify
its inertia with $\widehat{\mathbb Z}(1)$, which explains the exact
power equation.  Every finite cyclotomic extension has Galois group
$(\mathbb Z/n)^\times$; the resulting surjections and compactness
show that $\chi:G_{\mathbb Q}\to\widehat{\mathbb Z}^{\times}$
is onto.  Choose $\sigma$ with $\chi(\sigma)=\lambda$.
Renaming the target generators gives a continuous isomorphism
\[
 \tau:F_1=\widehat F(x,y)\xrightarrow{\sim}
       F_2=\widehat F(u,v),
 \qquad \tau([x,y])=[u,v]^\lambda.                     \tag{D.30.6}
\]
This uses a single Galois element of prescribed multiplier.
The construction below proves its outer involution directly from
$\tau$ and $\tau^{-1}$.

\emph{The automorphism of the closed surface group.}
Use its defining profinite presentation
\[
 H=\langle x,y,u,v\mid[x,y][u,v]=1\rangle^{\widehat{\ }},
 \qquad l=[x,y]^{-1}=[u,v],\qquad c=l^\kappa.           \tag{D.30.7}
\]
Define $\alpha$ on the first pair by the two words of $\tau$,
and on the second pair by the two words of
$\operatorname{Inn}(c)\tau^{-1}$.
Equation (D.30.6) and $\lambda^{-1}=\lambda$ give
\[
 \alpha([x,y])=l^\lambda,\qquad
 \alpha([u,v])=c\,[x,y]^\lambda c^{-1}=l^{-\lambda}.
\]
Their product is one, so this substitution defines a continuous
endomorphism of the presented group $H$.  It gives
$\alpha(l)=l^\epsilon$ and $\alpha(c)=c$.
Substitution on each of the four generators, using the actual inverse
$\tau^{-1}$, then gives
\[
                     \alpha^2=\operatorname{Inn}(c).
                                                        \tag{D.30.8}
\]
Consequently $\operatorname{Inn}(c^{-1})\alpha$ is its inverse.
These identities use the universal property of the displayed
presentation, so no assertion about efficiency of a profinite
amalgam is required.

Define $E$ on the compact space $H\times\{0,1\}$ by
\[
 (h,i)(k,j)=
 \bigl(h\alpha^i(k)c^{\lfloor(i+j)/2\rfloor},
                           i+j\bmod2\bigr).            \tag{D.30.9}
\]
For quotient labels $(1,1,0)$ associativity is precisely
$\alpha^2(k)c=ck$; for $(1,1,1)$ it also uses
$\alpha(c)=c$.  The other six label triples follow by direct
substitution.  The identity is $(1,0)$; the inverse of $(h,0)$ is
$(h^{-1},0)$, and that of $(h,1)$ is
$(\alpha^{-1}(h^{-1}c^{-1}),1)$.
All operations are continuous.  Thus $E$ is a profinite group
and its projection gives (D.30.1).  Write $t=(1,1)$, so that
$tht^{-1}=\alpha(h)$ and $t^2=c$.

\emph{An exact presentation and torsion-freeness.}
Set $u'=\tau(x)$ and $v'=\tau(y)$.  This is a free basis
of the formal second punctured-torus group, with
\[
 [u',v']=[u,v]^\lambda=l^\lambda,\qquad
       u'=txt^{-1},\quad v'=tyt^{-1}.
\]
The relations for $t$ on $u',v'$ follow from $t^2=c$.
Eliminating these two generators therefore gives the exact
profinite presentation
\[
 E\cong\langle x,y,t\mid
          tlt^{-1}=l^\epsilon,\quad t^2=l^\kappa
                 \rangle^{\widehat{\ }},
                 \qquad l=[x,y]^{-1}.                 \tag{D.30.10}
\]
For completeness, the homomorphism to (D.30.9) sends $x,y,t$
to the named elements and is onto because their conjugates
$u',v'$ generate the second factor.  In the reverse direction,
recover $u,v$ by applying the words of $\tau^{-1}$ to
$txt^{-1},tyt^{-1}$.  Their commutator is
$[txt^{-1},tyt^{-1}]^\lambda=l$, by (D.30.10).
Thus the surface relator holds, and the conjugation relations on
both pairs follow from $t^2=c$ and the defining words of
$\alpha$.  The two maps fix all named generators and are inverse.

The universal property of the maximal pro-$2$ quotient substitutes
$\epsilon_2=\kappa_2=1$ into (D.30.10).  Its first relation follows
from the second, giving
\[
 E^{(2)}\cong
       \langle x,y,t\mid[x,y]t^2=1\rangle_{\mathrm{pro}\text{-}2}.
                                                        \tag{D.30.11}
\]
This is the rank-three Demushkin presentation with $q=2$ and
$f=\infty$: cyclically move $t^2$ to the beginning and replace
$x,y$ by their inverses to match the commutator convention of
\cite[\S1.4 and Theorem~1.9]{R-Wilkes}.  The converse in that
theorem proves the Demushkin assertion.  Its abelianization is
$\mathbb Z_2^2\times C_2$, so it is infinite.  The same source's
preceding cohomological-dimension statement gives
$\operatorname{cd}_2(E^{(2)})=2$.  A subgroup of order two has
infinite $2$-cohomological dimension, so $E^{(2)}$ is torsion-free.

There is also an explicit free-basis identification with the
three-crosscap word.  The substitution
\[
 x\mapsto a d^{-1}b^{-1}a^{-1},\qquad
 y\mapsto a b d a b a^{-1},\qquad t\mapsto abd
                                                        \tag{D.30.12}
\]
sends $[x,y]t^2$ to $a^2b^2d^2$.  Its inverse is
\[
 a\mapsto xt,\qquad
 b\mapsto t^{-1}x^{-1}t^{-1}yxt,\qquad
 d\mapsto t^{-1}x^{-1}y^{-1}t^2.
\]
Free reduction checks the relator and both composites.  The maps
extend to mutually inverse pro-$2$ isomorphisms, proving the first
comparison in (D.30.3).

The parity quotient $E\to C_2$ factors through $E^{(2)}$ and
sends $t$ to its nonidentity element and $x,y$ to the identity.
Every finite-order element of $E$ maps trivially to the torsion-free
$E^{(2)}$, hence has trivial parity and lies in $H$.
The latter is torsion-free by A.7.  This proves torsion-freeness of
$E$ itself and, in particular, the absence of a section of (D.30.1).

For an odd prime $p$, (D.30.10) gives $t^2=1$ in $E^{(p)}$,
so $t=1$.  At $p=3$ the remaining relation is $l=l^{-1}$,
which forces $l=1$, and the quotient is
$\langle x,y\mid[x,y]=1\rangle_{\mathrm{pro}\text{-}3}
=\mathbb Z_3^2$.  At every other odd prime the remaining relation
is the identity $l=l$.  The universal properties in both directions
give the free pro-$p$ group on $x,y$.  This proves all of (D.30.3).

\emph{The orientation line and the obstruction to every surface.}
Let $A\in\operatorname{GL}_2(\widehat{\mathbb Z})$ be the
abelianization matrix of $\tau$.  In the finite Heisenberg group
modulo $n$, the commutator of two elements is the central generator
raised to the determinant of their two abelianization columns.
Applying (D.30.6) there gives $\det A=\lambda\pmod n$ for every
$n$, hence $\det A=\lambda$ in $\widehat{\mathbb Z}$.
On the basis $([x],[y],[u],[v])$ of $H^{\mathrm{ab}}$, inner
conjugation disappears and the matrix of $\alpha$ is
\[
 M=\begin{pmatrix}0&A^{-1}\\A&0\end{pmatrix}.
\]
For $J=\left(\begin{smallmatrix}0&1\\-1&0\end{smallmatrix}\right)$,
the identity $A^tJA=(\det A)J$ gives
\[
 M^t\operatorname{diag}(J,J)M
       =\lambda\operatorname{diag}(J,J),                \tag{D.30.13}
\]
because $\lambda^{-1}=\lambda$.  The companion identity
$M\operatorname{diag}(J,J)M^t=\lambda\operatorname{diag}(J,J)$
also holds by $AJA^t=(\det A)J$.  Since the matrix of $\alpha^*$
on degree-one cohomology is $M^t$, this last identity gives the
multiplier of its cup pairing with all dual maps specified.
The completed surface cup-product comparison of A.2 is natural
for automorphisms.  Its nondegenerate degree-one pairing therefore
shows that $\alpha^*$ acts by $\lambda_p$ on the one-dimensional
$H^2(H;\mathbb F_p)$.  The matrix $M$ exchanges two nonzero direct
summands, so $\alpha$ is not inner.  Equation (D.30.8) thus gives
a faithful outer $C_2$ action.

The orientation-line character is a continuous homomorphism
\[
 \operatorname{Out}(H)\longrightarrow
                    \widehat{\mathbb Z}^{\times}.
\]
Its restriction to the extended geometric mapping class group
has image the two diagonal elements $+1,-1$, and continuity gives
the same bound for its closure. Conjugation in $\operatorname{Out}(H)$
preserves this character because its target is abelian. Our element
has value $\lambda$, with $\lambda_3=1$ and $\lambda_5=-1$, so no
conjugate lies in that closure. D.16's exact transporter formulas
therefore give empty fixed discretization and prediscretization
sets for this actual involution.

For odd $p$, averaging over $C_2$ makes all its higher cohomology
with $\mathbb F_p$-coefficients vanish.  The continuous
Hochschild--Serre spectral sequence of (D.30.1) consequently gives
\[
 H^2(E;\mathbb F_p)\xrightarrow[\mathrm{restriction}]{\sim}
                      H^2(H;\mathbb F_p)^{C_2}.
                                                        \tag{D.30.14}
\]
At $p=3$ the scalar is $\lambda_3=1$, while at $p=5$ it is
$\lambda_5=-1$.  This proves (D.30.2).

Finally, the completion of a closed orientable surface group of
positive genus has one-dimensional $H^2(-;\mathbb F_p)$ at every
prime, by the surface comparison of A.2 (including its torus
calculation).  A closed nonorientable surface of genus at least
two has an orientable double cover.  Exact completion at index
two and the same restriction map (D.30.14) compute its degree-two
cohomology at odd primes as the invariants of the orientation
reversal, hence zero.  The sphere group is trivial and the
projective-plane group is $C_2$, so both have zero degree-two
cohomology at odd primes.  A surface with boundary or a noncompact
surface has a free fundamental group, whose profinite completion
has cohomological dimension at most one.  These cases exhaust the
surface groups and exclude the pair of values in (D.30.2).

The independently executed program
\path{code/part_b/closed_extension_d30_audit.py}
checks the free-word isomorphism in (D.30.12), the intersection
matrices and the finite cochain computation of the two invariant
lines.  The profinite construction, presentation comparisons and
torsion-freeness are proved above; the finite computations verify
the stated identities.
\end{proof}

\begin{claim}{D.31}{Proved}{Abstract finite groups in the geometric completions of genus at most two}
Let $M_{g,n}^+=\Gamma^+_{g,n}$ allow permutations of the $n$
punctures, and let $M_{g,n}^{\pm}$ also allow orientation reversal.
Every finite nonsolvable subgroup of $\widehat{M_{0,n}^+}$ is
isomorphic to $A_5$. Its permutation action on the punctures has
\[
       n=12a+20b+30c+60m,\qquad
       a,b,c\in\{0,1\},\quad m\geq0,                 \tag{D.31.1}
\]
where the summands are the actual orbit sizes. In particular, it
has at least twelve punctures. Every finite nonsolvable subgroup
of $\widehat{M_{0,n}^{\pm}}$ is isomorphic to $A_5$ or
$A_5\times C_2$. Every finite subgroup of
$\widehat{M_{g,n}^{\pm}}$ for $g=1,2$ is solvable.

These statements determine the abstract groups and, in the positive
sphere case, their puncture-permutation representations. The
surjectivity of the map from geometric $A_5$ conjugacy classes to
completed $A_5$ conjugacy classes, and the completed $A_5$
centralizer comparison of F.7, remain to be proved.
\end{claim}
\passage{\S{}30; folios 194--196; scans 383--387; \S{}31; folios 198--199 verso; scans 391--394; Part D(b): finite subgroups of the geometric completion and full completion comparison}

\begin{proof}
For $n=0,1,2$ the positive sphere mapping class groups are
respectively $1,1,C_2$; including orientation reversal gives
groups of order at most four. Their completions are the same
finite solvable groups. The sphere arguments below therefore
take $n\geq3$.

We give the completed exact sequences first, so that every use of
a subgroup or quotient has its specified completion map.

\emph{A completion and configuration-group preparation.}
For the forgetting sequences used below, the kernel $N$ is finitely
generated and the quotient $Q$ is good. These properties imply
that $\widehat N\to\widehat E$ is injective. Here is the required
argument. Take a characteristic finite-index subgroup $N_0$ of $N$;
these are cofinal because $N$ is finitely generated. In $E/N_0$ the
kernel $B=N/N_0$ is finite. Its centralizer has finite index and fits
into the central extension
\[
 1\longrightarrow Z(B)\longrightarrow C_{E/N_0}(B)
       \longrightarrow Q_0\longrightarrow1,
\]
where $Q_0$ has finite index in $Q$. Goodness passes to $Q_0$.
The class in $H^2(Q_0,Z(B))$ is therefore represented by a continuous
cohomology class of $\widehat Q_0$, hence is inflated from a finite
quotient. Restriction to the kernel of that finite quotient kills
it. The central extension splits over this finite-index subgroup;
a chosen complement meets $B$ trivially and has finite index in
$E/N_0$. Its core is a finite-index normal subgroup still meeting
$B$ trivially. Pulling this core back to $E$ gives a finite quotient
whose restriction to $N$ has kernel exactly $N_0$. This proves the
full induced topology and thus
\[
 1\longrightarrow\widehat N\longrightarrow\widehat E
       \longrightarrow\widehat Q\longrightarrow1.     \tag{D.31.2}
\]
The last equality follows directly from the finite quotients of
$E/N$.

The braid groups on the plane, on finitely punctured surfaces,
and on the closed surfaces of genus one or two used here are
good and have finite cohomological dimension.
We recall how the ingredients apply here. The ordered
configuration-space fibrations of Fadell--Neuwirth
\cite[Theorems~1 and~3]{A-FadellNeuwirth1962} forget one point and
have fibre the surface with finitely many points removed.
The base surface and the fibres are aspherical; hence their
configuration spaces are aspherical by the homotopy exact sequence.
The fibres' groups are finitely generated free groups once a point
has been removed. Starting with the good surface group, the
preceding finite-kernel argument and the Lyndon--Hochschild--Serre
comparison spectral sequences prove goodness inductively.
The required cohomology of each fibre with finite coefficients is
finite. For the plane the initial group is trivial; for the closed
surfaces its goodness is the usual surface-group goodness also
used in \cite{R-BZ}. The unordered group has the ordered one as a
finite-index subgroup, so goodness follows by the finite-index
comparison and Shapiro's lemma. The configuration spaces are
finite-dimensional aspherical manifolds, which give finite
cohomological dimension. Goodness therefore makes each completed
braid group torsion-free: its cohomology with every finite
$p$-primary coefficient module vanishes above that dimension,
whereas a subgroup $C_p$ has infinite $p$-cohomological dimension.
This also follows directly from the prime-subgroup comparison in
\cite[Theorem~A]{R-BZ}.

\emph{Finite subgroups of a completed puncture stabilizer are cyclic.}
For $n\geq4$, capping the boundary of a disc by a once-punctured
disc gives the exact sequence
\[
 1\longrightarrow\langle\Delta^2\rangle\cong\Z
 \longrightarrow B_{n-1}\longrightarrow
 (M_{0,n}^+)_{p}\longrightarrow1.                      \tag{D.31.3}
\]
Here the distinguished puncture is $p$, and the kernel is the
central full twist; this is
\cite[Lemma~2.2]{A-BigelowBudney2001}, whose proof specifies the
capping map. Braid abelianization sends each Artin generator to
$1$ and $\Delta^2$ to $(n-1)(n-2)$. Thus the composite
\[
 \widehat\Z\longrightarrow\widehat B_{n-1}
       \longrightarrow\widehat\Z
\]
is multiplication by this nonzero integer and is injective.
Consequently completion of (D.31.3) is the exact central extension
\[
 1\longrightarrow\widehat\Z\longrightarrow
 \widehat B_{n-1}\longrightarrow
 \widehat{(M_{0,n}^+)_{p}}\longrightarrow1.             \tag{D.31.4}
\]
The last group identifies with the stabilizer of $p$ in
$\widehat{M_{0,n}^+}$: its discrete source has finite index, and
stabilizing $p$ is the inverse image of a subgroup in the finite
puncture-permutation quotient.

Let $K$ be a finite subgroup of that completed stabilizer and let
$E_K$ be its inverse image in $\widehat B_{n-1}$. The subgroup
$A=\widehat\Z$ is central and open in $E_K$, of index $d=|K|$.
Continuous transfer over its finitely many cosets gives
\[
             \operatorname{Ver}:E_K\longrightarrow A,
             \qquad \operatorname{Ver}(a)=d a
             \quad(a\in A).                           \tag{D.31.5}
\]
Its kernel intersects $A$ trivially, so injects into the finite
group $E_K/A$. It is finite, and is trivial because $E_K$ is a
subgroup of the torsion-free completed braid group. Transfer is
therefore a continuous injection into $\widehat\Z$. Its compact
image is a closed procyclic subgroup. Hence $E_K$, and then its
finite quotient $K$, are procyclic and cyclic respectively.
For $n=3$ the same capping map is
$1\to2\Z\to B_2=\Z\to C_2\to1$, and the conclusion is immediate.
For $n=1,2$ the positive point stabilizer is trivial.

It follows in particular that a finite subgroup of a completed
positive sphere mapping class group preserving a specified set
of $1\leq b\leq3$ punctures is solvable. The kernel of its permutation
map to $S_b$ fixes a puncture and is cyclic, while $S_b$ is
solvable. This assertion also applies to every finite-index
decorated subgroup through its actual inclusion into the completed
sphere group.

\emph{The prime-power inputs and elimination of a solvable radical.}
The good-group prime-power theorem
\cite[Theorem~A]{R-BZ} conjugates each finite $p$-subgroup of
$\widehat{M_{0,n}^+}$ to a geometric one and gives its completed
normalizer. The checked prime-power topology argument in F.7
identifies that normalizer with the full completion of its
discrete normalizer. Thus, for a nontrivial normal $p$-subgroup
$A$ of a finite $G\leq\widehat{M_{0,n}^+}$, conjugate $A$ into
$M_{0,n}^+$ and use the exact covering-normalizer sequence
\[
 1\longrightarrow A\longrightarrow N_{M_{0,n}^+}(A)
       \longrightarrow Q_A\longrightarrow1.           \tag{D.31.6}
\]
The group $Q_A$ is a finite-index decorated subgroup of the
mapping class group of the quotient sphere. Its specified branch
set has two points for a cyclic rotation group, and three for a
noncyclic rotation group. This branch set includes the images of
ramification points that were already punctures upstairs; the
ramification indices are retained in the covering decoration.
Completion of (D.31.6) and the normalizer comparison give an
actual injection
\[
                 G/A\lhook\joinrel\longrightarrow\widehat Q_A.
                                                               \tag{D.31.7}
\]
It preserves that set of two or three branch points. The puncture
stabilizer argument makes $G/A$ solvable, and then $G$ is solvable.

For a nontrivial finite solvable normal subgroup $R$ of $G$, the
last nonzero derived term of $R$ is abelian. Its nontrivial
$p$-primary component $A$ is characteristic in $R$, hence normal
in $G$. The preceding argument again makes $G$ solvable.
Consequently a nonsolvable finite $G$ has trivial solvable radical.
Only the prime-power comparison was used in this step.

The same prime-power theorem and the finite rotation-group list
give cyclic odd Sylow subgroups of $G$, and cyclic or dihedral
Sylow $2$-subgroups, including $C_2^2$ in the latter case.
There is also the precise normalizer constraint
\[
 \operatorname{im}\bigl(N_G(C_p)\longrightarrow\Aut(C_p)\bigr)
                      \subset\{1,-1\}.               \tag{D.31.8}
\]
Indeed normalize a geometric rotation. A representative of a
discrete normalizer either fixes its two poles or exchanges them.
The oriented local rotation numbers are opposite, so the induced
automorphism is respectively $1$ or $-1$. The normalizer comparison
preserves this image because $\Aut(C_p)$ is finite. The same
argument with the local orientation sign proves (D.31.8) in the
extended geometric completion as well. Equivalently one can apply
the prime-power normalizer theorem to the good index-two extended
mapping class group; the discrete representatives still induce
only these two exponents.

We record the elementary finite-group reduction needed below.
A finite group all of whose Sylow subgroups are cyclic is solvable.
For the least prime $p$ dividing its order, its cyclic Sylow group
$P$ has $N(P)/C(P)$ of order both prime to $p$ and dividing
$p^{a-1}(p-1)$, where $|P|=p^a$. Minimality of $p$ forces this
quotient to be trivial. Transfer to $P$ has restriction
$x\mapsto x^{[G:P]}$ and is onto. To check this formula without
a fusion assumption, any two elements of the abelian Sylow group
$P$ conjugate in $G$ are conjugate by $N(P)$: the two Sylow
subgroups containing the conjugate element are conjugate inside
its centralizer. Thus they are equal when $P\leq Z(N(P))$.
The coset-cycle formula for transfer now gives the stated exponent.
Its kernel is a normal $p$-complement. Induction on the group order
proves solvability. It follows that every odd-order subgroup here
is solvable, and that a nonsolvable $G$ has a noncyclic dihedral
Sylow $2$-subgroup.

\emph{The verified Gorenstein--Walter application.}
Let $S$ be such a Sylow $2$-subgroup and $z\in Z(S)$ an involution.
The quotient of $C_G(z)$ by $\langle z\rangle$ in (D.31.7)
preserves the two branch points of the rotation quotient.
Its pointwise branch stabilizer is cyclic. Its inverse image
$C_0\leq C_G(z)$ has index at most two and is a central extension
of a finite cyclic group by $\langle z\rangle$, hence is abelian:
a lift of a quotient generator and the central kernel generate it.
The odd-order part of $C_0$ is characteristic in $C_0$, normal in
$C_G(z)$, and cyclic because all odd Sylow subgroups are cyclic.
It is a normal abelian $2$-complement of $C_G(z)$.

We can therefore apply the original, precisely matching theorem
\cite[Theorem~I, p.~553]{A-GorensteinWalter1962}. It states that
$G$ has a normal odd-order subgroup $K$ such that $G/K$ is one
of
\[
       \operatorname{PSL}_2(q),\quad
       \operatorname{PGL}_2(q)\quad(q\text{ odd}),
       \qquad A_7,\qquad\text{a Sylow }2\text{-subgroup}.
                                                               \tag{D.31.9}
\]
The source explicitly includes $C_2^2$ among dihedral groups.
All its hypotheses have just been verified. The odd group $K$
is solvable by the preceding transfer argument, so the triviality
of the solvable radical gives $K=1$.

The group $A_7$ contains the commuting independent cycles
$(1\ 2\ 3)$ and $(4\ 5\ 6)$, contradicting cyclic odd Sylow
subgroups. A $2$-group is solvable. Write $q=p^f$ in either
matrix case. The unipotent subgroup
$u(t)=\left(\begin{smallmatrix}1&t\\0&1\end{smallmatrix}\right)$
for $t\in\F_q$ is elementary abelian of order $p^f$, so $f=1$.
In $\operatorname{PSL}_2(p)$ the diagonal matrices
$\operatorname{diag}(a,a^{-1})$ act on this $C_p$ by
$u(t)\mapsto u(a^2t)$, supplying $(p-1)/2$ distinct automorphisms.
In $\operatorname{PGL}_2(p)$ the matrices
$\operatorname{diag}(a,1)$ supply all $p-1$ automorphisms.
Constraint (D.31.8) leaves $p\leq5$ in the first case and
$p\leq3$ in the second. At $p=3$, the faithful projective-line
action embeds these groups into $S_4$, so they are solvable.
The remaining group is $\operatorname{PSL}_2(5)\cong A_5$.

For this last isomorphism the exact map is conjugation on the
five Sylow $2$-subgroups. The supplied matrix enumeration finds
fifteen involutions, partitions them into the five four-groups,
and verifies that the conjugation image has order sixty and
consists of even permutations. With the recorded ordering of
the five subgroups, the matrices
\[
 \begin{pmatrix}1&1\\0&1\end{pmatrix},\qquad
 \begin{pmatrix}0&-1\\1&0\end{pmatrix}
\]
map respectively to $(1\ 2\ 4\ 3\ 5)$ and
$(2\ 4)(3\ 5)$. Thus the map is injective and its image is exactly
$A_5$, with its sources and targets fixed.

\emph{The puncture representation.}
The pure positive sphere group is the ordered configuration
group of $n-3$ points on the three-punctured sphere: fix the
first three marked points and use their contractible identity
homeomorphism component in the evaluation fibration. It is good
of finite cohomological dimension by the preceding configuration
argument, so its completion is torsion-free. This is exactly
the kernel of the completed map to $S_n$, by the finite-index
completion sequence. Thus every finite subgroup acts faithfully
on the puncture set.

For a finite $A_5$ subgroup, each puncture stabilizer is cyclic by
(D.31.4)--(D.31.5). The cyclic subgroups of $A_5$ have orders
$1,2,3,5$, so its transitive puncture orbits are precisely
$A_5/1$, $A_5/C_2$, $A_5/C_3$, and $A_5/C_5$.
The last three have sizes $30,20,12$. For each $d=2,3,5$,
an element generating $C_d$ fixes
$|N_{A_5}(C_d)/C_d|=2$ points in the corresponding orbit and
none in the other three types. Prime-power conjugacy makes its
puncture permutation that of a rotation of the sphere, which
has at most two fixed punctures. Each exceptional orbit occurs
at most once. This proves (D.31.1), including its actual
permutation representation, and gives the lower bound twelve.

\emph{The extended sphere group.}
Let $G\leq\widehat{M_{0,n}^{\pm}}$ be finite and nonsolvable.
Its positive subgroup $G^+$ has index at most two and is
nonsolvable, so is $A_5$. The centralizer $C_G(A_5)$ intersects
$A_5$ trivially and injects into $G/G^+$, so has order at most
two. With a nontrivial centralizer, its generator is negative
and $G=A_5\times C_2$. With a trivial centralizer, conjugation
embeds $G$ into $\Aut(A_5)\cong S_5$ with image containing the
inner $A_5$, leaving $A_5$ or $S_5$. The automorphism isomorphism
is induced by the action on the five subgroups $A_4\leq A_5$.
An automorphism fixing all five makes $g$ and its image induce
the same permutation of them for every $g\in A_5$; the faithful
natural action of $A_5$ then makes that automorphism the identity.
Conjugation by $S_5$ is faithful as well and supplies one hundred
twenty automorphisms, proving the isomorphism. In $S_5$, the normalizer of
a $5$-cycle acts on it by all four nonzero exponents modulo five.
This violates the extended version of (D.31.8). Therefore only
$A_5$ and $A_5\times C_2$ occur as abstract nonsolvable groups.

\emph{Forgetting punctures in genus two.}
Evaluation on the unordered puncture set, with isotopy extension,
gives the discrete forgetting sequence
\[
 1\longrightarrow B_n(S_2)\longrightarrow M_{2,n}^+
        \longrightarrow M_{2,0}^+\longrightarrow1.      \tag{D.31.10}
\]
The identity homeomorphism component of $S_2$ is contractible,
as used in Part C, so there is no preceding evaluation-loop
kernel. The quotient is good, the braid kernel is finitely
generated, and (D.31.2) completes this sequence exactly.
The completed kernel is torsion-free by the configuration
argument. Every finite subgroup of $\widehat{M_{2,n}^+}$
therefore injects under forgetting into $\widehat{M_{2,0}^+}$.
The hyperelliptic quotient is the exact sequence
\[
 1\longrightarrow C_2\longrightarrow M_{2,0}^+
        \longrightarrow M_{0,6}^+\longrightarrow1;
                                                               \tag{D.31.11}
\]
see \cite[Propositions~3.2--3.3]{A-BigelowBudney2001}.
The central involution is detected on homology modulo three,
so it injects into the completion. Finite-kernel completion
then makes (D.31.11) exact with hats. Every finite subgroup
has solvable image in $\widehat{M_{0,6}^+}$ by (D.31.1),
and kernel contained in $C_2$. It is solvable. The same holds
in the extended group because the positive subgroup has
index at most two. The case $n=0$ uses (D.31.11) directly.

\emph{Forgetting punctures in genus one.}
For $n\geq1$, evaluation gives
\[
 1\longrightarrow N_n\longrightarrow M_{1,n}^+
       \longrightarrow\operatorname{SL}_2(\Z)
       \longrightarrow1,\qquad
 N_n=B_n(T^2)/\Z^2.                                    \tag{D.31.12}
\]
Here $\Z^2$ is the image of the translation loops in the
identity homeomorphism component, identified by evaluation
as in C.9. It is central in the braid group: applying a loop
of translations to every configuration of a braid gives the
square homotopy interchanging the two loop products. Its injection
can be checked on the ordered configuration space by evaluating
the first strand. To verify its full topology in the unordered
braid group, sum the homology classes of all strand paths:
\[
 B_n(T^2)\longrightarrow H_1(T^2,\Z)=\Z^2.
\]
Endpoints are permuted, so the sum is a cycle and gives a
homomorphism. On diagonal translation loops it is multiplication
by $n$. The resulting completed injection proves the exact
central extension
\[
 1\longrightarrow\widehat\Z^2\longrightarrow
 \widehat{B_n(T^2)}\longrightarrow\widehat N_n
       \longrightarrow1.                              \tag{D.31.13}
\]
The completed braid group is torsion-free by the configuration
argument. For completeness the ordered torus configuration
has the explicit product coordinates
\[
 (x_1,\ldots,x_n)\longmapsto
 (x_1,x_2-x_1,\ldots,x_n-x_1)
 \in T^2\times\operatorname{Conf}_{n-1}(T^2\setminus\{0\}),
\]
so its goodness also follows directly from that of $\Z^2$
and the iterated punctured-surface configuration groups.

For a finite $K\leq\widehat N_n$, pull it back in (D.31.13).
Transfer to the central open $\widehat\Z^2$ has restriction
$a\mapsto |K|a$. Its kernel again injects into $K$ and is
finite, so torsion-freeness makes transfer injective.
The pullback, and hence $K$, is abelian. Since
$\operatorname{SL}_2(\Z)$ is good and $N_n$ is finitely
generated, (D.31.2) also completes (D.31.12) exactly.
The image of every finite subgroup lies in a finite subgroup
of $\widehat{\operatorname{SL}_2(\Z)}$, which is cyclic
of order $1,2,3,4$, or $6$ by D.8's virtually free comparison.
The finite kernel is abelian by (D.31.13), proving solvability.
For orientation reversal the positive subgroup has index at
most two, which preserves this conclusion. For $n=0$ use
D.8 directly with $\operatorname{GL}_2(\Z)$.

The executed program
\path{code/part_a/sphere_completion_restriction.py}
checks the projective-matrix automizer calculations at
$p=3,5,7,11,13$, constructs the displayed
$\operatorname{PSL}_2(5)\to A_5$ isomorphism, compares
the $C_5$ normalizer actions in $A_5$ and $S_5$, and
checks the exceptional coset fixed-point counts and the
$C_3^2$ subgroup of $A_7$. The general reductions are
the completed exact sequences and transfer arguments above.
\end{proof}

\begin{claim}{D.32}{Proved}{The three-cusp rank test and finite pure loop groups}
Let $F=\widehat F(x,y)$, with the three named peripheral generators
$c_1=x$, $c_2=y$, and $c_3=(xy)^{-1}$.  Let $P$ be the subgroup of
$\Out(F)$ preserving each of their three procyclic conjugacy classes.
Its common peripheral multiplier is denoted by $\lambda$.
Every finite subgroup of $P$ maps injectively, by $\lambda_2$, to
$\{1,-1\}\leq\Z_2^\times$.  In particular it has order at most two,
and a nonidentity element of finite order has order two and
$\lambda_2=-1$.  These statements hold for $GT_0\leq P$ as well.

For the full three-cusp loop-outer group, allowing puncture permutations,
every finite subgroup $G$ has the explicit faithful representation
\[
 G\longrightarrow S_3\times C_2,\qquad
       g\longmapsto\bigl(\sigma_g,\lambda_2(g)\bigr).
                                                        \tag{D.32a}
\]
This is an assertion about the displayed homomorphism and abstract
finite groups.  Conjugacy of its original action to the geometric
action with this image is not asserted here. D.34 determines the
global multiplier of a surviving involution; its geometric conjugacy
remains to be proved.
\end{claim}
\passage{30; folios 194--196; scans 383--387; 31; folios 198--199 verso; scans 391--394; Part D(b) of the task}

\begin{proof}
We first prove the closed-subgroup fact which uses all three cusps.
For a prime $p$, let $e_p\in\widehat\Z$ have $p$-coordinate one and
all other coordinates zero.  Suppose a closed subgroup $C\leq F$
has $\dim_{\F_p}H^1(C,\F_p)=2$ and contains, for each $i=1,2,3$,
a conjugate
\[
                    g_i c_i^{e_p\Z_p}g_i^{-1}.             \tag{D.32b}
\]
We prove $C=F$.

The comparison with open subgroups is the actual restriction map
\[
 \underset{U\text{ open},\ C\leq U\leq F}{\operatorname{colim}}
       H^1(U,\F_p)\ \xrightarrow{\ \sim\ } H^1(C,\F_p),
                                                        \tag{D.32c}
\]
where the transition maps are restrictions as $U$ decreases.
Indeed a continuous character of $C$ kills $C\cap V$ for some open
normal $V\triangleleft F$.  It extends to $CV$ by the quotient
$CV/V=C/(C\cap V)$.  Conversely a character on $U$ which vanishes
on $C$ vanishes after restriction to its open kernel, which contains
$C$.  This proves both surjectivity and injectivity of (D.32c).

If $C$ is open, finite-coset Schreier rewriting identifies it with
a free profinite group of rank $1+[F:C]$.  Its given first-cohomology
dimension therefore forces $[F:C]=1$.
Suppose instead that $C$ has infinite index.  The $p$-valuations
of the finite indices $[F:U]$, for open $U$ containing $C$, must
be unbounded.  To prove this, assume they were bounded and choose
$U_0$ attaining their maximum.  For $C\leq V\leq U_0$, every
further restriction between such open subgroups has prime-to-$p$
index and is injective on $H^1(-,\F_p)$: a nonzero character has
an index-$p$ kernel and cannot vanish on a subgroup of prime-to-$p$
index.  Thus (D.32c) contains the first cohomology of every such
$V$.  The indices $[F:V]$ are unbounded, since their intersection
is $C$ and $[F:C]=\infty$.  Their Schreier dimensions
$1+[F:V]$ are therefore unbounded, a contradiction.

Choose an open $U\geq C$ for which $p\mid [F:U]$.  Put
$\Delta=F(x,y)$, the named dense discrete subgroup of $F$, and
$U_0=U\cap\Delta$.  The finite covering of the pair of pants
defined by $U_0\leq\Delta$ has degree $n=[F:U]$, and its
completed fundamental group is exactly $U$: finite-index subgroups
of a discrete group carry the full induced profinite topology.
For each $g_i$ choose $h_i\in\Delta$ in its left $U$-coset and
write $g_i=u_i h_i$, with $u_i\in U$.  The subgroup
\[
 U\cap g_i\overline{\langle c_i\rangle}g_i^{-1}
  =u_i\overline{\langle h_i c_i^{d_i}h_i^{-1}\rangle}u_i^{-1}
                                                        \tag{D.32d}
\]
is consequently the completion, up to the displayed inner
conjugation, of a particular boundary subgroup in this finite
cover.  Here $d_i$ is its boundary covering degree.  Inclusion
(D.32b) proves $p\nmid d_i$.

The three selected boundaries lie over the three different original
boundaries, so they are distinct.  There is at least one more
boundary.  Otherwise each original boundary would have just one
lift, and its degree would be $n$, contrary to $p\nmid d_i$ and
$p\mid n$.  If the cover has genus $h$ and $b\geq4$ boundary
components, its presentation is
\[
 \left\langle a_1,b_1,\ldots,a_h,b_h,d_1',\ldots,d_b'
       \ \middle|\ \prod_{j=1}^h[a_j,b_j]\,
                         d_1'\cdots d_b'=1\right\rangle.
\]
Thus the only relation between its boundary homology classes over
$\F_p$ is their sum.  The three selected classes are independent,
and evaluation on them gives a surjection
$H^1(U,\F_p)\to\F_p^3$.

This remains a surjection after restriction of these characters to
$C$ and evaluation on the three elements in (D.32b).
To check the comparison explicitly, the element
$g_i c_i^{e_p}g_i^{-1}$ is the $e_p/d_i$ power of the boundary
generator in (D.32d).  Inner conjugation does not change a
character, and its evaluation is therefore multiplied by the
nonzero scalar $d_i^{-1}\in\F_p$.  We obtain three independent
characters of $C$, contradicting its dimension two.  This proves
the closed-subgroup assertion.

We now apply it to finite loop actions, retaining the precise
normalization map.  If an outer class fixes the three peripheral
classes, its exponents on $c_1,c_2,c_3$ are equal: in
$F^{\rm ab}=\widehat\Z^2$, the relation
$[c_1]+[c_2]+[c_3]=0$ gives equality of the two coordinates.
Conjugate a representative $\beta$ to have
$\beta(x)=x^\lambda$, and write
$\beta(y)=k^{-1}y^\lambda k$.  Write $[k]=(a,b)$ in the
named abelian basis.  Conjugating this representative by $x^a$
gives the normalized pair
\[
 \alpha(x)=x^\lambda,\qquad
 \alpha(y)=f^{-1}y^\lambda f,
       \qquad f=y^{-b}kx^{-a}\in F'.                     \tag{D.32e}
\]
Multiplication on the left of the conjugator by $y^{-b}$ changes
neither image of $y$ nor the automorphism.  Two normalized
representatives differ by conjugation by an element centralizing
$x^\lambda$, hence by $x^v$ by A.7.  Their two conjugators then
differ by left multiplication by a power of $y$ and right
multiplication by $x^{-v}$.  Since both are in $F'$, their two
abelian coordinates force both powers to be zero.  Normalization
is therefore unique.  A product of normalized pairs has second
coordinate $f\alpha(g)\in F'$, so (D.32e) is a homomorphic
section into $\Aut(F)$.  Its image is precisely the Belyi group
defined in \cite[p.~3120, equation (2)]{R-HS}; the argument has
not identified that entire group with its subgroup $GT_0$.

Let an element of $P$ have prime order $p$ with $\lambda_p=1$.
Its normalized automorphism has order $p$.  The proof of D.28
uses exactly the properties just established: finite order of the
normalized automorphism, the common multiplier, and preservation
of the three cusps.  Neither of the extra $GT_0$ symmetry equations
is used in its norm computation or Brown restriction argument.
In particular that proof gives a single conjugacy class of
sections in $F\rtimes C_p$ and the actual fixed subgroup
$C=C_F(C_p)$ with $\dim H^1(C,\F_p)=2$.
Each of the three peripheral sections centralizes the whole
$c_i^{e_p\Z_p}$; conjugating those sections to the specified
$C_p$ places the three groups (D.32b) in this same $C$.
The closed-subgroup assertion gives $C=F$, so the normalized
automorphism is the identity.  This contradicts its prime order.

For odd $p$, an element of order dividing $p$ in
$\Z_p^\times$ equals one, by reduction modulo $p$ and the
torsion-free group $1+p\Z_p$.  Thus no odd prime occurs in
the order of a finite subgroup of $P$.  The finite image of
$\lambda_2$ lies in $\{1,-1\}$, since $1+4\Z_2$ is
torsion-free.  Its kernel in any finite subgroup has no element
of order two by the preceding argument, and no element of odd
prime order either.  Cauchy's theorem makes that kernel trivial.
This proves the first assertion.

For a class permuting cusps, compose it with the inverse of the
corresponding positively oriented geometric permutation.  The
preceding abelian calculation gives its common multiplier too,
and composition on cusp generators proves that both $\sigma$
and $\lambda$ are homomorphisms.  On a finite subgroup, the
kernel of (D.32a) is a finite subgroup of $P$ on which
$\lambda_2=1$, hence is trivial.  All subgroups of the target
already have their geometric actions on the three-punctured
sphere: use an equilateral triple on the equator, the positive
rotational $S_3$, and reflection in that equator for the
commuting $C_2$.  This proves the claimed abstract group bound
and its attainment, while keeping the action-conjugacy question
separate through the actual map (D.32a).

The program \texttt{code/part\_c/three\_cusp\_rank\_d32.py}
enumerates the finite covers used for a validation through degree
six.  It computes the Schreier graph, every boundary cycle,
the unique boundary-homology relation, and explicit dual
characters for the selected boundaries, including the
$d_i^{-1}$ restriction scalars.  The proof above applies to
arbitrary degrees and closed subgroups; no bounded enumeration
is used in place of it.
\end{proof}

\paragraph{Source check for the fixed-second-coordinate assertion.}
Lochak--Schneps, Lemma~9 on p.~583 of
\cite{R-LS}, asserts that a fixed $f$ occurs only with the two
multipliers $\lambda$ and $-\lambda$.  The displayed class-two
identity $\lambda^2=24r+1$ determines the squares of the prime
coordinates.  The second proof in that source adds the required
full-group equation, whose elementary verification is as follows.

Comparing $(\lambda,f)$ and $(\mu,f)$ by the exact group law
gives $(\lambda^{-1}\mu,1)$.  Set
$n=(\lambda^{-1}\mu-1)/2\in\widehat\Z$.
The three-cycle relation is
\[
                       x^n(xy)^{-n}y^n=1.                 \tag{LS.1}
\]
In the class-two Heisenberg quotients its central coordinate
gives $n(n+1)=0$.  Thus in each integral domain $\Z_p$ one
has $n_p=0$ or $n_p=-1$.  These choices are linked by actual
finite quotients: for every odd prime $q$, evaluate (LS.1) in
\[
 D_q=\langle r,s\mid r^q=s^2=1,\ srs=r^{-1}\rangle,
                         \qquad x=s,\quad y=r.
\]
Both $x$ and $xy$ have order two.  If $n_2=0$, (LS.1) is
$r^{n_q}=1$, forcing $n_q=0$.  If $n_2=-1$, it is
$r^{1+n_q}=1$, forcing $n_q=-1$.  Hence $n$ is globally
zero or globally minus one, proving $\mu=\lambda$ or
$\mu=-\lambda$.  Both occur, because multiplication on the
right by $(-1,1)$ leaves the second coordinate unchanged.

This verifies the source's fixed-$f$ conclusion by exact
quotient morphisms.  Powers of a general torsion element
have second coordinates $f\alpha(f)\cdots\alpha^{j-1}(f)$;
the argument does not identify those coordinates with each
other.  D.32 supplies the separate argument restricting
finite subgroups, and leaves the stated involution questions
unchanged.

\begin{claim}{D.33}{Proved}{Finite punctured-torus actions are detected by their primary homology}
Put $F=\widehat F(x,y)$ and $c=[x,y]$, and let
$T=\Out_{\rm lac}(F;\overline{\langle c\rangle})$ be the full
once-punctured-torus loop-outer group.  For every prime $p$, the
homomorphism induced by the named abelianization,
\[
              A_p:T\longrightarrow\GL_2(\Z_p),            \tag{D.33a}
\]
is injective on every finite $p$-subgroup of $T$.
The full abelianization homomorphism
$T\to\GL_2(\widehat\Z)$ is therefore injective on every finite
subgroup.  Every such subgroup has order dividing $24$.

The fixed-subgroup calculation underlying these assertions is also
explicit.  For an automorphism $\alpha$ of $\widehat F_r$ with
$\alpha^p=1$ whose induced automorphism on the maximal pro-$p$
quotient is the identity, put $E=\widehat F_r\rtimes_\alpha C_p$.
There is one kernel-conjugacy class of sections of $E\to C_p$,
and for its fixed subgroup $C$ the actual restriction map is an
isomorphism
\[
 H^1(\widehat F_r,\F_p)\ \xrightarrow{\ \sim\ }H^1(C,\F_p).
                                                        \tag{D.33b}
\]
For the punctured-torus loop structure and $r=2$, its cusp forces
$C=F$ and hence $\alpha=1$.
\end{claim}
\passage{30; folios 194--196; scans 383--387; 31; folios 198--199 verso; scans 391--394; Part D(b) of the task}

\begin{proof}
We first prove (D.33b), keeping the multiplicative information in
the cohomology comparison.  Write $P=\widehat F_r^{(p)}$ and
$S=\{1\}\times C_p\leq E$.  The quotient homomorphism is
\[
 \pi:E\longrightarrow P\times C_p,\qquad
              (h,s^i)\longmapsto(\bar h,s^i).              \tag{D.33c}
\]
Restriction along $\widehat F_r\to P$ is an isomorphism on
mod-$p$ cohomology: in degree one both sides are the same
$r$-dimensional character group, and higher cohomology vanishes
because both groups have $p$-cohomological dimension one.
The Hochschild--Serre comparison for (D.33c) is consequently an
isomorphism on its second pages and gives a ring isomorphism
\[
 \pi^*:H^*(P\times C_p,\F_p)\xrightarrow{\sim}H^*(E,\F_p).
                                                        \tag{D.33d}
\]
All maps here are induced by the specified quotient, not by an
identification of the two groups.  In particular
$\dim H^n(E,\F_p)=r+1$ for every $n>1$.

Let $u\in H^2(C_p,\F_p)$ be the nonzero periodic class, and
inflate it to either side of (D.33d).  The K\"unneth calculation
for the direct product gives, for every $m\geq1$, an injective map
\[
 H^1(P,\F_p)\longrightarrow H^{2m+1}(E,\F_p),
                         \qquad\chi\longmapsto\pi^*(\chi u^m).
                                                        \tag{D.33e}
\]
For a subgroup $Q\leq E$ of order $p$, its projection to
$C_p$ is an isomorphism, and
$N_E(Q)=C_{\widehat F_r}(Q)\times Q$, exactly as in D.28.
The first coordinate of $\pi(Q)$ is trivial because $P$ is
torsion-free.  Thus every character $\pi^*\chi$ in (D.33e)
vanishes on $Q$.  Its restriction to this normalizer lies
entirely in $H^1(C_{\widehat F_r}(Q),\F_p)$, and multiplication
by $u^m$ places it in
\[
 H^1(C_{\widehat F_r}(Q),\F_p)\otimes H^{2m}(Q,\F_p).
                                                        \tag{D.33f}
\]

Apply the profinite Brown restriction theorem in the precise
form verified in \cite[Theorem 2 and \S1.1]{A-KochloukovaZalesskii2007}.
There is no $C_p\times C_p$ in $E$, since a finite subgroup
injects into $E/\widehat F_r=C_p$.  Its open subgroup
$\widehat F_r$ has $p$-cohomological dimension one.  The
finite-dimensional source makes the theorem's dense restriction
embedding an isomorphism onto a finite product.  If $k$ is
the number of conjugacy classes and
$d_Q=\dim H^1(C_{\widehat F_r}(Q),\F_p)$, its dimension formula is
\[
                   k+\sum_{[Q]}d_Q=r+1.                  \tag{D.33g}
\]
The injection (D.33e), followed by this same restriction map,
lands in the summands (D.33f).  Therefore
$\sum d_Q\geq r$.  Since $k\geq1$, (D.33g) forces
$k=1$ and $d_Q=r$.  For the specified $S$, (D.33e)--(D.33f)
then say exactly that restriction
$H^1(\widehat F_r,\F_p)\to H^1(C_{\widehat F_r}(S),\F_p)$
is injective; both spaces have dimension $r$, proving (D.33b).
As in D.28, conjugacy under $E$ equals conjugacy under the kernel,
because $E=\widehat F_r Q$ for any section $Q$.

Now take $r=2$ and assume that $\alpha$ preserves the conjugacy
class of $L=\overline{\langle c\rangle}$, where $c=[x,y]$.
Write its peripheral multiplier as $\lambda$.  In a class-two
Heisenberg quotient, the image of the commutator under an
abelianization matrix $A$ is the $\det(A)$ power of the
commutator.  The finite Heisenberg quotients for all moduli
therefore give the exact identity
\[
                        \lambda=\det(\alpha^{\rm ab}).   \tag{D.33h}
\]
In the present calculation $\alpha$ is the identity on $P$,
so $\lambda_p=1$.

The cusp normalizer has the exact sequence
\[
 1\longrightarrow L\longrightarrow N_E(L)
                        \longrightarrow C_p\longrightarrow1.
                                                        \tag{D.33i}
\]
The kernel is $L$ by A.7.  For a lift $a$ of the chosen
generator, write
$aca^{-1}=c^\lambda$ and $a^p=c^v$, so $(\lambda-1)v=0$.
The first coordinate of $\pi(a)$ centralizes the image $\bar c$
in $P$ and belongs to $\overline{\langle\bar c\rangle}$,
by the pro-$p$ peripheral centralizer assertion of A.7.
Thus it is $\bar c^w$ and $v_p=pw$.
This uses the whole cusp subgroup.  Its primitive pro-$p$
parameter can also be checked directly: in the mod-$p$ Magnus
expansion, $[x,y]-1$ starts with $XY-YX$.  A $p$th power
starts in degree divisible by $p$.  At $p=2$, the possible
degree-two square is $(aX+bY)^2$; its $X^2$ and $Y^2$
coefficients prevent it from equaling $XY+YX$.  Hence the
commutator is not a $p$th power.

Since $\alpha^p=1$, its multiplier satisfies $\lambda^p=1$.
Set $N_\lambda=1+\lambda+\cdots+\lambda^{p-1}$.
At a prime $q$ with $\lambda_q\ne1$, both
$(N_\lambda)_q$ and $v_q$ are zero.  Where $\lambda_q=1$
the norm is $p$; $v_q$ is divisible by $p$, using $v_p=pw$
at $q=p$ and invertibility of $p$ at all other primes.
Consequently there is $z\in\widehat\Z$ with
$N_\lambda z=-v$.  The element $ac^z$ has order $p$
and centralizes $c^{e_p\Z_p}$.
The unique section-conjugacy class just proved places a
conjugate $g c^{e_p\Z_p}g^{-1}$ inside
$C=C_F(S)$.

Every character of $C$ vanishes on $g c^{e_p}g^{-1}$:
by (D.33b) it is the restriction of a character of $F$,
and $c$ is a commutator.  We prove that this forces $C=F$.
Its mod-$p$ first-cohomology dimension is two.  If $C$ is
open, Schreier gives $1+[F:C]=2$.  If $C$ is proper and
not open, the restriction-colimit proof of D.32 shows that
its index has infinite $p$-part.  Choose an open $U\geq C$
with $p\mid n=[F:U]$, and take the degree-$n$ covering of
the punctured torus corresponding to $U\cap F(x,y)$.
The chosen cusp has a boundary lift of degree $d$ prime
to $p$, since its subgroup contains the full chosen
$c^{e_p\Z_p}$.  The comparison with a discrete boundary
generator is exactly (D.32d), with a dense discrete
representative for the profinite conjugator $g$.
There must be at least two boundary lifts: a unique lift
would have degree $n$, divisible by $p$.
The selected boundary has nonzero homology in this cover,
since the sole relation among the at least two boundary
classes is their sum.  A character of $U$ detecting this
boundary restricts to a character of $C$ detecting
$g c^{e_p}g^{-1}$; the value is multiplied by
$d^{-1}\bmod p$, as computed in D.32.  This is the required
contradiction.  Hence $C=F$ and $\alpha=1$.

To apply this to (D.33a), let an outer element have order
$p$ and $A_p=I$.  D.27 supplies an actual order-$p$
automorphism representing it.  Its induced pro-$p$
automorphism has identity abelianization.  The Magnus
argument in D.10 and D.25 makes this finite-order
automorphism the identity: a first nonzero homogeneous
deviation in the $\Z_p$ Magnus algebra would be multiplied
by $p$ under the $p$th iterate, and the homogeneous module
has no $p$-torsion.  Thus (D.33c) applies to this lift
directly; there is no subsequent conjugation which could
change its order.  The cusp argument proves that the lift,
and hence its outer class, is the identity.
The kernel of (D.33a) on a finite $p$-subgroup cannot
contain an order-$p$ element, so it is trivial.  Likewise,
the kernel of full abelianization on any finite subgroup
would contain an element of some prime order and is trivial.

We finish with the numerical bound on these finite groups.
For $p>3$, a matrix of order $p$ over $\Q_p$ has a primitive
$p$th root among its eigenvalues.  The polynomial
$\Phi_p(t)$ is irreducible of degree $p-1$ over $\Q_p$,
as seen by Eisenstein after $t\mapsto t+1$.  It cannot be
the factor of a degree-two characteristic polynomial.
Thus no such prime divides the order of a finite subgroup
of $T$.  A finite $3$-subgroup injects by $A_3$ into
$\GL_2(\Z_3)$ and then injects modulo $3$: the kernel
$1+3M_2(\Z_3)$ is torsion-free by the matrix logarithm.
Since $|\GL_2(\F_3)|=48$, its order divides $3$.

For completeness, a finite $2$-subgroup $K$ of
$\GL_2(\Q_2)$ has order at most eight.  The same
Eisenstein argument for $\Phi_{2^a}$ excludes order
$2^a$ for $a\geq3$.  Every matrix of order four has
minimal polynomial $t^2+1$, so its square is $-I$.
Consequently the projective image of $K$ has exponent
two.  Extend the matrices to the algebraic closure
$\overline{\Q}_2$; the scalar kernel is unchanged.
An exponent-two subgroup of $\PGL_2(\overline{\Q}_2)$
has order at most four.  To prove this
last bound, conjugate a nonidentity element to
$z\mapsto-z$.  Every commuting transformation preserves
its pair of fixed points, so is $z\mapsto az$ or
$z\mapsto a/z$.  The scaling subgroup has at most two
elements under the exponent-two condition and has index
at most two.  The scalar kernel of $K$ is contained in
$\{I,-I\}$, the finite $2$-torsion of $\Q_2^\times$.
Thus $|K|\leq8$.  Apply this to a Sylow $2$-subgroup
through $A_2$.  The Sylow bounds and the exclusion of
other primes give $|G|\mid 8\cdot3=24$.

The program \texttt{code/part\_c/torus\_homology\_d33.py}
checks the Magnus degree-two coefficients, the periodic
cohomology multiplication used in (D.33e), and the
one-cusp restriction argument on every connected finite
cover through degree six.  The matrices and boundary
characters are obtained from actual Schreier walks.
The bound and comparison maps above are proved for the
full profinite groups, independently of the bounded checks.
\end{proof}

\begin{claim}{D.34}{Proved}{Global signs and the cusp cycle of a pure involution}
Every nonidentity finite-order element of the full pure
three-cusp loop-outer group $P$ of D.32 is an involution
whose common multiplier is the global element
$-1\in\widehat\Z^\times$.  In particular this holds in
$GT_0$, and the multiplier-one kernel in either group is
torsion-free.

For its normalized lift
\[
 \alpha(x)=x^{-1},\qquad\alpha(y)=f^{-1}y^{-1}f,
                      \qquad f\in F',
                                                        \tag{D.34a}
\]
there are precisely three $F$-conjugacy classes of sections
of $F\rtimes_\alpha C_2\to C_2$.  Each of their fixed
subgroups has zero first cohomology with $\F_2$ coefficients.
The two section classes in each cusp normalizer are the
endpoints of one edge of a triangle.  Moreover there is
$h\in F'$ with
\[
                         f=h^{-1}\alpha(h).               \tag{D.34b}
\]
Thus $\alpha$ inverts both $x$ and $hyh^{-1}$.
The endomorphism $x\mapsto x$, $y\mapsto hyh^{-1}$ induces
an automorphism of every maximal pro-$p$ quotient.  Its
surjectivity on the full $F$, and conjugacy of the original
involution to complex conjugation in the pure loop group,
are not conclusions of this statement.
\end{claim}
\passage{30; folios 194--196; scans 383--387; 31; folios 198--199 verso; scans 391--394; Part D(b) of the task}

\begin{proof}
We first prove a genus-zero incidence calculation for an
arbitrary number $n\geq3$ of cusps.  Let
\[
 F_n=\widehat{\langle c_1,\ldots,c_n\mid
                              c_1\cdots c_n=1\rangle}
\]
and let $\beta$ be an actual involution preserving each
peripheral conjugacy class, with common multiplier $\mu$
and $\mu_2=-1$.  Put $E_n=F_n\rtimes_\beta C_2$.
The equality $\beta^2=1$ gives $\mu^2=1$; hence each
$\mu_q$ is $1$ or $-1$.

For each cusp $L_i=\overline{\langle c_i\rangle}$,
A.7 gives an exact sequence
\[
 1\longrightarrow L_i\longrightarrow N_{E_n}(L_i)
                          \longrightarrow C_2\longrightarrow1.
                                                        \tag{D.34c}
\]
Choose a lift $a_i$ of the nonidentity element and write
$a_i c_i a_i^{-1}=c_i^\mu$ and $a_i^2=c_i^{v_i}$.
Then $(\mu-1)v_i=0$.  In every coordinate with
$\mu_q=-1$ this forces $(v_i)_q=0$, including $q=2$.
In a coordinate with $\mu_q=1$ the prime $q$ is odd,
so division by two is possible.  Thus
$(1+\mu)u_i=-v_i$ has a solution in $\widehat\Z$,
and $b_i=a_i c_i^{u_i}$ is an involution.

The set of its further order-two translates by $L_i$,
modulo conjugation by $L_i$, is exactly
\[
 \frac{\ker(1+\mu:\widehat\Z\to\widehat\Z)}
      {(1-\mu)\widehat\Z}\ \simeq\ \F_2.               \tag{D.34d}
\]
At an odd prime with $\mu_q=-1$, multiplication by two
is onto, and a coordinate with $\mu_q=1$ contributes
zero to the kernel.  The remaining coordinate is
$\Z_2/2\Z_2$.  Representatives for the two classes are
$b_i$ and $c_i^{e_2}b_i$.

Since $\beta$ is pure and the cusp classes span homology,
its action on $V=H_1(F_n,\F_2)$ is the identity.  Therefore
the actual mod-two abelianization of the semidirect product is
\[
                  H_1(E_n,\F_2)=V\oplus\F_2.             \tag{D.34e}
\]
The two involutions just exhibited have images differing by
the nonzero vector $[c_i]$ in $V$.  They cannot be conjugate
in $E_n$ or in $F_n$.  Thus the two local classes in
(D.34d) give two distinct global section classes.

Build a finite graph with one edge for each of the $n$
cusps and with these two global classes as its endpoints.
This graph uses some of the global section classes; we
now bound all of them.  The Brown calculation of D.28,
with free rank $n-1$ and prime two, gives
\[
 \sum_{[Q]}
   \bigl(1+\dim_{\F_2}H^1(C_{F_n}(Q),\F_2)\bigr)
       =\dim_{\F_2}H^k(E_n,\F_2)\leq n\qquad(k>1).
                                                        \tag{D.34f}
\]
Every order-two subgroup is a section, and its total-group
and kernel conjugacy classes agree, because $E_n=F_nQ$.
Hence there are at most $n$ vertices available for the graph.
An edge has no loop by (D.34e).

Assign to every vertex its image $(v,1)$ in (D.34e).
The sum of the edge vectors around a graph cycle is zero,
since its successive vertex differences cancel.  But the
$n$ vectors $[c_1],\ldots,[c_n]$ have just one nonzero
linear relation over $\F_2$, namely
$[c_1]+\cdots+[c_n]=0$.  Consequently any simple cycle
in the graph uses all $n$ edges.  The graph cannot be a
forest: a forest with $n$ edges needs at least $n+1$
vertices, more than (D.34f) permits.  Thus it is exactly
a simple $n$-cycle.  It has $n$ vertices, so there are
no additional section classes outside this graph, and
(D.34f) gives
\[
            H^1(C_{F_n}(Q),\F_2)=0
                   \quad\text{for every section }Q.       \tag{D.34g}
\]
This construction is invariant under any automorphism
of the extension preserving its kernel and peripheral
classes, with any permutation of those classes: such
an automorphism sends the normalizer and its two classes
in (D.34c)--(D.34d) to those of the image cusp.

Return to a nontrivial finite-order element of $P$.
D.32 gives its order two, its normalized actual lift
$\alpha$, and $\lambda_2=-1$.  We show that
$\lambda_q=-1$ for every odd prime $q$.  The only other
possibility is $\lambda_q=1$.  Suppose this occurs and
use the exact cyclic quotient
\[
 \chi:F=\widehat F(x,y)\longrightarrow C_q,
          \qquad\chi(x)=1,\quad\chi(y)=0,
          \qquad N=\ker\chi.                            \tag{D.34h}
\]
The common multiplier gives
$\chi\alpha=\lambda\chi=\chi$, since its reduction
modulo $q$ is one.  In particular $N$ is invariant.
The corresponding degree-$q$ covering of the pair of
pants has one boundary over $x$, one over
$z=(xy)^{-1}$, and $q$ over $y$, with representative
generators
\[
 x^q,\qquad z^q,\qquad x^i yx^{-i}\quad(0\leq i<q).
                                                        \tag{D.34i}
\]
The two first degrees are $q$ and all others are one.
Its Euler characteristic is $-q$ and it has $q+2$
boundaries, so its genus is zero.  Finite-index
completion identifies its full profinite group with
the specified subgroup $N$.

The restriction of $\alpha$ is pure on this entire
peripheral structure, with the same multiplier
$\lambda$.  The first two boundaries are the unique
lifts of their original boundaries.  For the other
$q$ boundaries use the actual normalized formula
\[
 \alpha(x^i yx^{-i})
        =x^{\lambda i}f^{-1}y^\lambda f x^{-\lambda i}.
                                                        \tag{D.34j}
\]
Here $\chi(f)=0$, because $f\in F'$, and
$\lambda i\equiv i\pmod q$.  The conjugator in
(D.34j) therefore belongs to the same $N$-coset
as $x^i$: explicitly
$k_i=x^{\lambda i}f^{-1}x^{-i}\in N$ satisfies
$\alpha(x^i yx^{-i})=k_i(x^i yx^{-i})^\lambda k_i^{-1}$.
This is the required conjugacy inside $N$.  The action on
$H_1(N,\Z_2)$ is consequently the scalar $-1$,
so the restriction remains a nontrivial involution.

Put $E=F\rtimes_\alpha C_2$ and
$E_N=N\rtimes_\alpha C_2$.  The quotient maps give
\[
                     E/N=C_q\times C_2.                 \tag{D.34k}
\]
In particular $E_N$ is normal in $E$.  Conjugation by
$x$ defines an automorphism of $E_N$ preserving its
map to $C_2$, and thus acts on its $N$-conjugacy
classes of sections.  Its $q$th power is conjugation
by $x^q\in N$, so this action has order dividing
$q$ on those classes and on the cusp incidence graph.
On the cusp edges (D.34i), it fixes the $x^q$ and
$z^q$ classes and cycles the $q$ classes over $y$.
The second fixed class is the unique boundary over
$z$, or follows directly from its cyclic coset
decomposition.  Thus its graph action has order
exactly $q$.

By the preceding calculation the graph is a cycle
with $q+2$ edges.  Its automorphism group has order
$2(q+2)$: an image of a vertex and a choice of which
neighbor is its successor determine the automorphism.
The odd prime $q$ does not divide $2(q+2)$.
This contradicts the displayed order-$q$ graph action.
Therefore $\lambda_q=-1$ for every odd $q$, proving
the global multiplier assertion.  This argument used
the whole pure group $P$, without an additional
$GT_0$ symmetry assumption.

At $n=3$ the incidence graph is the stated triangle.
Write $s$ for the chosen section in $E=F\rtimes C_2$.
The normalized group law gives $f\alpha(f)=1$,
so $fs$ is also an involution and it normalizes $y$
by inversion.  The three involutions
\[
                     s,\quad x^{e_2}s,\quad y^{e_2}fs
                                                        \tag{D.34l}
\]
have the distinct images $(0,0,1)$, $(1,0,1)$,
and $(0,1,1)$ in (D.34e).  They represent all three
classes.  Since $f\in F'$, the image of $fs$ is
$(0,0,1)$, so it is conjugate to $s$ by an element
of $F$.  Choose $h$ with $fs=h^{-1}sh$.  This is
exactly (D.34b).  In abelianization it says
$0=(-2)[h]$ in the torsion-free group
$\widehat\Z^2$, hence $h\in F'$.
Substitution gives
\[
 \alpha(hyh^{-1})
   =(hf)(f^{-1}y^{-1}f)(f^{-1}h^{-1})
   =hy^{-1}h^{-1}.
\]
The two displayed inverted generators have the same
abelian images as $x,y$.  On each maximal pro-$p$
quotient their images therefore generate the Frattini
quotient and hence the whole group.  The induced
endomorphism of that finitely generated free pro-$p$
group is surjective and thus an automorphism.
No full-profinite surjectivity is inferred from those
separate comparisons.

The program \texttt{code/part\_c/pure\_involution\_d34.py}
checks the normalizer orbit calculation, the unique
binary cusp relation, and the exact cyclic-cover
boundary and deck permutations for several odd
primes.  The exclusion of all odd primes is the
unrestricted divisibility argument above.
\end{proof}

\begin{claim}{D.35}{Proved}{One exact conjugator equation and its finite obstructions}
Let $P$ be the full pure three-cusp group of D.32, let
$F=\widehat F(x,y)$, and put
$\delta(x)=x^{-1}$, $\delta(y)=y^{-1}$.
For a nonidentity involution $\alpha\in P$, use its normalized
lift from D.34,
\[
 \alpha(x)=x^{-1},\qquad
 \alpha(y)=f^{-1}y^{-1}f,\qquad f\in F'.
\]
Define the compact solution space
\[
\begin{split}
 \mathcal S_\alpha=\{(\mu,h,u)\in
       \widehat\Z^\times\times F'\times F:\quad&
       h^{-1}\alpha(h)=f,\\
       &x^\mu h y^\mu h^{-1}
                    =u(xy)^\mu u^{-1}\}.
\end{split}                                                   \tag{D.35a}
\]
The assignment
\[
 (\mu,h,u)\longmapsto
 \left[\beta_{\mu,h}\right],\qquad
 \beta_{\mu,h}(x)=x^\mu,\quad
 \beta_{\mu,h}(y)=h y^\mu h^{-1},                              \tag{D.35b}
\]
has image exactly
\[
 \operatorname{Conj}_{P}(\delta,\alpha)
       =\{\beta\in P:\ \beta\delta\beta^{-1}=\alpha\}.
                                                                  \tag{D.35c}
\]
Every map (D.35b) is an automorphism of the full profinite
group.  The fibers of the assignment are exactly
$(\mu,h,u)\mapsto(\mu,h,u(xy)^a)$ for $a\in\widehat\Z$.
It therefore induces a homeomorphism from the corresponding
quotient of $\mathcal S_\alpha$ to (D.35c).

For $\alpha\in GT_0$, the inclusion $GT_0\hookrightarrow P$
gives the equality of actual conjugator sets
\[
       \operatorname{Conj}_{GT_0}(\delta,\alpha)
             =\operatorname{Conj}_{P}(\delta,\alpha).             \tag{D.35d}
\]
The defining Harbater--Schneps relations for the conjugating
pair follow from the proof below; they are not discarded.
Finally, the existence of a solution of (D.35a) is equivalent
to the nonemptiness of its explicit finite solution sets
(D.35h) at every characteristic finite quotient of $F$.
\end{claim}
\passage{\S{}31; folios 198-199verso; scans 391-394; \S{}30; folio 194verso; scan 384; Part D(b): actual full GT\_0 torsion and conjugacy}

\begin{proof}
For $(\mu,h,u)\in\mathcal S_\alpha$, the free universal
property first gives a continuous endomorphism
$\beta_{\mu,h}:F\to F$ with the displayed generator images.
Write $H=\beta_{\mu,h}(F)$, a closed subgroup because $F$
is compact.  It is generated by two elements, so
$\dim H^1(H,\F_2)\leq2$.  The restrictions of the two
standard characters of $F$ to these generators are the
two coordinate characters: $\mu\bmod2=1$ and conjugation
does not change abelianization.  These restrictions are
independent, proving
\[
                         \dim H^1(H,\F_2)=2.                   \tag{D.35e}
\]
The unit exponent $\mu$ shows that $H$ contains
$\overline{\langle x\rangle}$ and
$h\overline{\langle y\rangle}h^{-1}$.
The second equation in (D.35a) puts
$u\overline{\langle xy\rangle}u^{-1}$ in $H$ as well.
In particular $H$ contains the three full cusp
$\Z_2$ subgroups required by D.32.  Its closed-subgroup
rank test, applied to (D.35e), gives $H=F$.
A surjective endomorphism of a finitely generated
profinite group is injective, so $\beta_{\mu,h}$ is
an automorphism.  Its images of the three cusps,
including the exact displayed equation for the third,
place its outer class in $P$.

The first equation in (D.35a) is $\alpha(h)=hf$.
Substitution on the two free generators gives
\[
 \alpha\beta_{\mu,h}(x)=x^{-\mu}
                          =\beta_{\mu,h}\delta(x),\qquad
 \alpha\beta_{\mu,h}(y)=h y^{-\mu}h^{-1}
                          =\beta_{\mu,h}\delta(y).
                                                                  \tag{D.35f}
\]
These are equalities of actual automorphisms, proving
that the image of (D.35b) lies in (D.35c).

Conversely, let $\beta\in P$ satisfy (D.35c), and take its
unique normalized representative from D.32.
Write it as $\beta_{\mu,h}$ with $h\in F'$.
The normalization is a homomorphic section
$P\to\Aut(F)$, so the outer conjugacy equality lifts
to the actual equality $\alpha\beta_{\mu,h}
=\beta_{\mu,h}\delta$.
Its value on $y$ implies that
$h^{-1}\alpha(h)f^{-1}$ centralizes $y^{-\mu}$.
By A.7 it is a power of $y$.  It also lies in $F'$,
so its abelianization forces this power to be zero.
Thus $h^{-1}\alpha(h)=f$.
The third-cusp condition on $\beta$ gives a conjugator
$u$ in the second equation of (D.35a).
This proves that the image is exactly (D.35c).

The normalized pair $(\mu,h^{-1})$ is unique.
For this fixed pair the two third-cusp conjugators
$u,u'$ differ by right multiplication by the
centralizer of $(xy)^\mu$, which is
$\overline{\langle xy\rangle}$ by A.7.
These are precisely the asserted fibers.
The assignment is continuous in the profinite
automorphism topology, and its domain is compact.
Its quotient is therefore homeomorphic to its
Hausdorff image.

For an alternative form using D.34, fix
$h_0\in F'$ with $h_0^{-1}\alpha(h_0)=f$, and put
$C=C_F(\alpha)$.
All solutions of the first equation of (D.35a)
are exactly $h=ch_0$ with $c\in C$:
$\alpha(h)=hf$ and $\alpha(h_0)=h_0f$ give
$\alpha(hh_0^{-1})=hh_0^{-1}$, and the reverse
substitution proves the converse.
Since $\alpha$ acts by $-1$ on $F^{\rm ab}$,
every $c\in C$ has $2[c]=0$, hence $[c]=0$.
Thus $C\subset F'$ and this substitution respects
the stated domain of (D.35a).  The remaining exact
equation is
\[
 x^\mu c h_0 y^\mu h_0^{-1}c^{-1}
                         =u(xy)^\mu u^{-1}.
                                                                  \tag{D.35g}
\]

We prove (D.35d) by identifying the relevant centralizer,
rather than assuming that a conjugator in $P$ lies in
$GT_0$.  Let $P^1=\ker(\lambda:P\to\widehat\Z^\times)$.
First,
$C_{P^1}(\delta)=1$.
Indeed a normalized element $(1,k)\in P^1$ commuting
with $\delta$ has $k=\delta(k)$, by uniqueness of the
normalized representative and its group law.
The following explicit mutually inverse homomorphisms
identify the semidirect product:
\[
 \begin{aligned}
 C_2\amalg C_2\amalg C_2
    &\xrightarrow{\ \sim\ }F\rtimes_\delta\langle s\rangle,
       & (s_0,s_1,s_2)&\longmapsto(s,xs,ys),\\
 (x,y,s)&\longmapsto(s_1s_0,s_2s_0,s_0).
 \end{aligned}
\]
The orders and conjugation relations follow by direct
substitution.  A free profinite factor has trivial
intersection with a conjugate by an element outside
that factor; the original full-profinite statement is
recalled in \cite[Theorem 5 of the inspected author
version]{A-HerfortZalesskii2008}.
Consequently the centralizer of $s_0$ in this free
product is $\langle s_0\rangle$.
Its intersection with $F$ is trivial.
Hence $C_F(\delta)=1$, so $k=1$ and
$C_{P^1}(\delta)=1$ as claimed.

Let $\Sigma\cong S_3$ be the geometric cusp-permutation
subgroup of the full loop-outer group, generated by
the outer classes
$\theta(x)=y,\ \theta(y)=x$ and
$\omega(x)=y,\ \omega(y)=(xy)^{-1}$.
The exact identification
$GT_0=C_P(\Sigma)$ follows from
\cite[Propositions 3--4 and Theorem 2]{R-HS}:
commutation with $\theta$ is relation (I), and,
after (I), commutation with $\omega$ is relation (II).
Both $\alpha$ and $\delta$ centralize $\Sigma$.
For any $\beta\in\operatorname{Conj}_P(\delta,\alpha)$
and $\sigma\in\Sigma$, form the actual element of
the outer group
$r_\sigma=\beta^{-1}\sigma\beta\sigma^{-1}$.
The pure subgroup $P$ is normal in the full
loop-outer group, and the common multiplier is a
homomorphism there.  Therefore $r_\sigma\in P^1$.
Moreover
\[
 \begin{split}
 r_\sigma\delta r_\sigma^{-1}
  &=\beta^{-1}\sigma\beta
        (\sigma^{-1}\delta\sigma)
        \beta^{-1}\sigma^{-1}\beta\\
  &=\beta^{-1}\sigma\alpha\sigma^{-1}\beta
    =\beta^{-1}\alpha\beta=\delta.
 \end{split}
\]
The proved equality $C_{P^1}(\delta)=1$ gives
$r_\sigma=1$.  Thus $\beta$ centralizes $\Sigma$,
so its normalized pair satisfies both (I) and (II)
and lies in $GT_0$.  This proves (D.35d).

Here is the exact finite form of the remaining
existence problem.  For each open characteristic
subgroup $N\triangleleft F$, put $Q_N=F/N$ and
$m_N=\exp(Q_N)$.
The automorphism $\alpha$ induces $\alpha_N$.
Use bars for the images of $x,y,f$ in $Q_N$ and define
\[
\begin{split}
 \mathcal S_{\alpha,N}=\{(a,b,v)\in
       (\Z/m_N)^\times\times Q_N'\times Q_N:\quad&
       b^{-1}\alpha_N(b)=\bar f,\\
       &\bar x^a b\bar y^a b^{-1}
              =v(\bar x\bar y)^a v^{-1}\}.
\end{split}                                                       \tag{D.35h}
\]
These are finite sets, obtained entirely from the
specified finite quotient and its induced
automorphism.  Reduction sends
$\mathcal S_\alpha$ to every
$\mathcal S_{\alpha,N}$.
Conversely, a point of a nonempty
$\mathcal S_{\alpha,N}$ lifts to
$\widehat\Z^\times\times F'\times F$:
units lift by the Chinese remainder theorem,
and the image of $F'$ in $Q_N$ is exactly $Q_N'$.
The lifts satisfying the two equations modulo $N$
form a nonempty closed subset of this compact space.
Every finite collection of such conditions is
implied by the one at a common characteristic
refinement.  Their finite intersection property
therefore gives a point satisfying both equations
in $F$ itself.  Open characteristic subgroups are
cofinal for a finitely generated profinite group,
so equality at all these quotients is equality in
$F$.  This proves the asserted finite criterion.

The need for the full third-cusp equation can be
checked by a short executed finite certificate.
Put $d=[x,y]=xyx^{-1}y^{-1}$ and take the actual
commutator word
\[
                 w=y d^{-1}y^{-1}x^{-1}d^{-1}x\in F'.
\]
Consider the endomorphism
$\rho(x)=x$, $\rho(y)=wyw^{-1}$.
Its abelianization is the identity, so the Frattini
argument makes its induced endomorphism on every
maximal pro-$p$ quotient an automorphism.
In the finite quotient
\[
 F\twoheadrightarrow A_5,\qquad
             x\longmapsto(1\,2\,3),\quad
             y\longmapsto(3\,4\,5),
\]
the program
\path{code/part_a/pure_involution_conjugator.py}
enumerates the generated group of order $60$ and
computes
\[
 \bar w=(1\,4\,3\,2\,5),\qquad
 \overline{wyw^{-1}}=(1\,2\,3).
\]
The image of $\rho(F)$ in this quotient therefore
has order $3$, proving $\rho(F)\ne F$.
Moreover $\overline{xy}$ has order $5$, whereas
$\overline{xwyw^{-1}}$ has order $3$; the
third-cusp equation fails in this same quotient.
This word is a certificate concerning the
generation test.  It is not asserted to solve
the coboundary equation for a nonstandard
involution.  No nonemptiness or emptiness
assertion for the spaces (D.35a) of arbitrary
involutions is inferred from this finite example.
\end{proof}

\begin{claim}{D.36}{Proved}{The global multiplier of a positive transposition}
Let $B$ be the full three-cusp loop-outer group, and let
$g\in B$ be an involution whose $2$-adic multiplier is
$\lambda_2(g)=1$.  Its puncture permutation is a
transposition, and its common multiplier is the global
element $1\in\widehat\Z^\times$.

Consequently a finite subgroup of $B$ whose map (D.32a)
has image $S_3\times C_2$ has the usual global orientation
character.  The same is true when that image is a
$V_4$ containing the pure involution, or is the positively
oriented $S_3\times\{1\}$.  These statements identify
the actual multiplier homomorphism; they do not assert
conjugacy of those finite actions.
\end{claim}
\passage{\S{}30; folios194--196; scans383--387; \S{}31; folios198--199verso; scans391--394; Part D(b) of the task}

\begin{proof}
We first record the genus-zero extension of the actual
automorphism calculation already proved in D.32--D.33.
Let $F_n=\widehat\pi_{0,n}$ with $n\geq3$, and let
$\beta$ be a pure automorphism with $\beta^p=1$ and
$p$-adic peripheral multiplier one.  Then $\beta=1$.
Indeed $F_n$ is free of rank $r=n-1$, and its cusp
classes span its abelianization.  Its induced pro-$p$
automorphism has identity abelianization and finite
order, so the $\Z_p$ Magnus argument of D.33 makes
it the identity.  The ring calculation (D.33c)--(D.33g)
gives one conjugacy class of sections of
$F_n\rtimes_\beta C_p\to C_p$, and its actual fixed
subgroup $C$ has
\[
                       \dim H^1(C,\F_p)=n-1.               \tag{D.36a}
\]
The norm calculation of D.28 applies to every cusp:
its image is a primitive element of the free pro-$p$
group, the first coordinate of its normalizer lift
centralizes that element, and the lift's $p$th power
has $p$-coordinate divisible by $p$.  Coordinatewise
division of the norm therefore supplies a section in
each peripheral normalizer.  Uniqueness of the
section class places conjugates of all $n$ full
peripheral $\Z_p$ subgroups inside the same $C$.

The finite-cover rank proof of D.32 now applies with
all $n$ cusps.  If $C$ is open, its first-cohomology
dimension is $1+[F_n:C](n-2)$, which equals $n-1$
only at index one.  If it has infinite index and
finite $p$-part of that index, the same exact
restriction colimit and eventual prime-to-$p$
injections give unbounded first-cohomology dimension.
Otherwise choose an open $U\geq C$ of $p$-divisible
index.  The $n$ selected boundary lifts all have
degrees prime to $p$ and lie over different original
boundaries.  They cannot be all the boundaries:
a single lift over any original boundary would
have degree $[F_n:U]$.  With at least $n+1$ boundaries,
the $n$ selected boundary homology classes are
independent.  Their dual characters restrict to
$C$ with the nonzero multipliers $d_i^{-1}\bmod p$
computed in (D.32d).  Thus $\dim H^1(C,\F_p)\geq n$,
contradicting (D.36a).  Hence $C=F_n$ and $\beta=1$.
This proof concerns an actual finite-order automorphism;
it does not presume a lift for arbitrary finite outer
actions in larger rank.

Return to $g$.  D.32 excludes a nonidentity pure
involution with $\lambda_2=1$, so its permutation
must be a transposition.  Relabel the original cusp
generators as $x,y,z=(xy)^{-1}$ so that $x$ is fixed
and $y,z$ are exchanged.  D.27 provides an actual
order-two lift $\alpha$ of this outer action.
Its peripheral multiplier satisfies $\lambda^2=1$,
so all its prime coordinates are $1$ or $-1$.

Suppose $\lambda_q=-1$ for an odd prime $q$, and use
the exact cyclic quotient
\[
 \chi:F\longrightarrow C_q,\qquad
       \chi(x)=0,\quad\chi(y)=1,\quad\chi(z)=-1,
                   \qquad N=\ker\chi.                    \tag{D.36b}
\]
The puncture permutation and the multiplier give
$\chi\alpha=\chi$: on $y$ the exponent $-1$ at
$q$ is multiplied by $\chi(z)=-1$, and similarly
for $z$.  Thus $N$ is invariant and the induced
action on its cyclic deck quotient is the identity.
The degree-$q$ cover is a sphere with $q+2$ cusps.
There are $q$ degree-one cusps over $x$, and one
degree-$q$ cusp over each of $y,z$.

Here are the exact covering and filling generators.
Finite-coset Schreier rewriting with coset representatives
$1,y,\ldots,y^{q-1}$ gives the free basis of $N$
\[
 a_i=y^i x y^{-i}\quad(0\leq i<q),\qquad b=y^q.
\]
The remaining boundary generator is
\[
 d=z^q=b^{-1}a_{q-1}^{-1}\cdots a_0^{-1},
                     \qquad a_0\cdots a_{q-1}bd=1.        \tag{D.36c}
\]
For example the latter relation is the actual free-word
identity
$(xy)^q=a_0\cdots a_{q-1}y^q$.
Completing the discrete Schreier comparison identifies
this presented free profinite group with the specified
open subgroup $N$.

The lift $\alpha$ fixes the conjugacy class of every
$a_i$ inside $N$.  To see this without choosing an
unproved normalization, write
$\alpha(x)=h^{-1}x^\lambda h$ for some $h\in F$.
Since $\chi\alpha(y)=1$, its action on the indices
of these cusp classes is the translation
$i\mapsto i-\chi(h)$ in $C_q$.  Its square is the
identity, so $2\chi(h)=0$ and $\chi(h)=0$.
Explicitly
\[
 k_i=\alpha(y)^i h^{-1}y^{-i}\in N,\qquad
                \alpha(a_i)=k_i a_i^\lambda k_i^{-1}.     \tag{D.36d}
\]
The two unique boundary classes over $y,z$ are
exchanged, with multiplier $\lambda$.

Fill exactly these two cusps.  The map is the quotient
by the closed normal subgroup of $N$ generated by $b,d$:
\[
 N\longrightarrow
 K=N/\overline{\langle\!\langle b,d\rangle\!\rangle}_N
   =\widehat{\langle a_0,\ldots,a_{q-1}\mid
                             a_0\cdots a_{q-1}=1\rangle}
   \cong\widehat\pi_{0,q}.                               \tag{D.36e}
\]
Both presentation maps are explicit: kill $b$ and
send the remaining $a_i$ to the named peripheral
generators; the inverse identifies $a_{q-1}$ with
$(a_0\cdots a_{q-2})^{-1}$.  These substitutions
verify both universal properties and compositions.
Since $\alpha$ exchanges the normal generators'
peripheral classes, the closed normal subgroup
in (D.36e) is invariant.  Consequently $\alpha$
induces an actual automorphism $\bar\alpha$ of
$K$ with $\bar\alpha^2=1$.  By (D.36d) it is
pure, with the same multiplier $\lambda$.

The proved genus-zero calculation at $p=2$ applies:
$q\geq3$ and $\lambda_2=1$, so $\bar\alpha=1$.
But its $q$-adic multiplier on each surviving
nontrivial cusp is $-1$, which is not the
identity on the cusp's $\Z_q$ subgroup.  Equivalently,
it acts as $-I$ on the nonzero $q$-adic abelianization
of the genus-zero group $K$.  This contradiction
excludes every such $q$.  Together with
$\lambda_2=1$ it proves $\lambda=1$ globally.

For the finite-group consequences, use the faithful
map (D.32a).  In a $V_4$ with pure kernel, the
pure involution has multiplier $-1$ by D.34,
and the positive transposition has multiplier
$1$ by the present proof.  These generate it.
The positively oriented $S_3$ is generated by
its positive transpositions.  Finally
$S_3\times C_2$ is generated by those
transpositions and its pure involution.
Multiplicativity of the common multiplier proves
each asserted character identity.

The program \texttt{code/part\_c/positive\_transposition\_d36.py}
checks (D.36c) as an exact free-word identity,
the finite cyclic cusp permutations and their
involutive translations, and the two presentation
substitutions in (D.36e) for several odd primes.
The proof above supplies the all-prime assertion.
\end{proof}

\begin{claim}{D.37}{Proved}{Negative transpositions and the specified four-cusp deck action}
Every involution of the full three-cusp loop-outer group whose
puncture permutation is a transposition and whose multiplier has
$\lambda_2=-1$ gives a faithful $V_4$ action on a specified
four-cusp double cover.  One generator is the standard free positive
deck involution; another has double-transposition permutation and
multiplier $\lambda$, and its center-free extension is torsion-free.
Conversely every such $V_4$ action containing the specified standard
deck action descends to a negative transposition on the three-cusp
group.  The constructions are inverse as marked extensions.

For the double-transposition permutation of D.29, neither of the
commuting single transpositions preserves any separating two-pants
partition whose pants the double transposition exchanges.  Thus an
additional deck action in this construction necessarily moves the
separator.  This is an obstruction to preserving that decomposition,
and asserts no obstruction to an action moving it.
\end{claim}
\passage{Part D(b) and D(d) of the task; \S{}24, folio 147 verso, scan 293; \S{}30, folios 194--196, scans 383--387; \S{}31, folios 198--199 verso, scans 391--394.}

\begin{proof}
We specify both the markings and all the extension maps.  Put
\[
 F=\widehat F(x,y),\qquad z=(xy)^{-1},\qquad
 \chi:F\longrightarrow C_2,\quad (\chi(x),\chi(y))=(0,1),
 \qquad N=\ker\chi.
                                                        \tag{D.37a}
\]
Schreier rewriting identifies $N$ with the full free profinite
group on
\[
 A_0=x,\qquad A_1=yxy^{-1},\qquad B=y^2,
 \qquad D=z^2=B^{-1}A_1^{-1}A_0^{-1}.
                                                        \tag{D.37b}
\]
Its four peripheral generators satisfy $A_0A_1BD=1$.
These are the actual degree-one lifts of $x$ and the degree-two
lifts of $y,z$; completion of the discrete index-two Schreier
isomorphism gives precisely the specified inclusion $N\subset F$.
Let $\delta\in\Out(N)$ be conjugation by $y$.  Its named
representative is
\[
 A_0\longmapsto A_1,\qquad
 A_1\longmapsto B A_0 B^{-1},\qquad B\longmapsto B,
 \qquad \delta^2=\operatorname{Inn}(B).
                                                        \tag{D.37c}
\]
Here the final equality concerns the named automorphism;
its outer class has order two, multiplier $1$, and permutation
$(A_0\ A_1)$.  Its center-free extension is the explicitly
identified group $F$, hence is torsion-free.  In the converse
below, ``the specified deck action'' means this action transported
by a given peripheral isomorphism.  It does not mean an arbitrary
positive outer involution with a torsion-free extension.

Let $g$ be a negative transposition, labeled to fix $x$ and
exchange $y,z$.  Its center-free extension $E$ contains $F$
as an open subgroup of index two.  The peripheral normalizer
surjects onto $\langle g\rangle$, and A.7 gives its kernel
$\overline{\langle x\rangle}=\widehat\Z$.
For a lift $v$ in this normalizer write
\[
 v x v^{-1}=x^\lambda,\qquad v^2=x^a,
 \qquad (\lambda-1)a=0,\qquad \lambda^2=1.
                                                        \tag{D.37d}
\]
At every prime where $\lambda_p=-1$ the last commutation
identity gives $a_p=0$.  At every prime where $\lambda_p=1$
the prime is odd, since $\lambda_2=-1$, so $2$ is invertible.
Choose $b_p=-a_p/2$ in these latter coordinates and $b_p=0$
in the former.  Then $s=x^b v$ has order two.  Thus
\[
 E=F\rtimes_\alpha\langle s\rangle,\qquad
 \alpha^2=1,\qquad \alpha(x)=x^\lambda.
                                                        \tag{D.37e}
\]
This constructs the required finite-order lift inside the
peripheral normalizer, without a further choice of inner
normalization.  The multiplier and the cusp permutation give
$\chi\alpha=\chi$.  Therefore $N$ is normal in $E$, and
$E/N=C_2\times C_2$, with basis the images of $y,s$.

There is exactly one $F$-conjugacy class of order-two subgroups
of $E$.  Indeed the action on $H^1(F,\F_2)$ is the two-dimensional
nontrivial unipotent module $J_2$, isomorphic to the regular
$\F_2[C_2]$ module.  Consequently
$H^i(C_2,H^1(F,\F_2))=0$ for $i>0$.
The Hochschild--Serre spectral sequence, with the split
inflation on its bottom row, gives
\[
                 \dim H^n(E,\F_2)=1\qquad(n\geq3).
                                                        \tag{D.37f}
\]
The full profinite Brown restriction theorem used in D.28
identifies this high-degree cohomology with the product over
the order-two conjugacy classes.  Each class contributes
$1+\dim H^1(C_F(T),\F_2)$, since its normalizer is
$C_F(T)\times C_2$; see
\cite[Theorem 2 and \S1.1]{A-KochloukovaZalesskii2007}.
The section in (D.37e) exists, so (D.37f) forces exactly one
class and zero first mod-two cohomology for its centralizer.
Conjugacy by $E$ agrees here with conjugacy by $F$: for any
section $T$, one has $E=FT$, and the $T$ factor centralizes
that subgroup itself.
Every finite-order element of $E$ injects into $E/F$,
because $F$ is torsion-free, and consequently has order two.
All nonidentity torsion therefore maps to the one line
$\langle\bar s\rangle$ in the abelian group $E/N$.

Conjugation gives a faithful action of $E/N$ on $N$ in
$\Out(N)$.  To check the kernel, first
$C_F(N)=1$: an element centralizing both $x$ and $y^2$
lies in $\langle x\rangle\cap\langle y\rangle=1$ by A.7.
For $e\in C_E(N)$ and $f\in F$, the commutator $[e,f]$
lies in $F\cap C_E(N)$ and hence is trivial.  Thus $e$
centralizes $F$; faithfulness of the original nonidentity
outer action and $Z(F)=1$ give $e=1$.

Write $a$ for the outer restriction of $\alpha$ to $N$.
It fixes the classes of $A_0,A_1$ and exchanges $B,D$.
For the second fixed class use the explicit conjugator
$\alpha(y)y^{-1}\in N$:
\[
 \alpha(A_1)
  =(\alpha(y)y^{-1}) A_1^\lambda
                         (\alpha(y)y^{-1})^{-1}.
                                                        \tag{D.37g}
\]
The unique lifts of the other two cusps are exchanged inside
$N$.  Thus $\delta$ and $a$ commute in $\Out(N)$, and
\[
 \tau=\delta a,\qquad
 \sigma(\tau)=(A_0\ A_1)(B\ D),\qquad
 \operatorname{or}(\tau)=\lambda.
                                                        \tag{D.37h}
\]
The preimage $E_\tau$ of the diagonal line
$\langle\bar y+\bar s\rangle$ in $E/N$ is the center-free
extension of this outer involution.  It is torsion-free:
all nonidentity torsion has image $\bar s$, and its
intersection with $N$ is already torsion-free.
For an explicit representative set
\[
 \beta=\operatorname{Inn}(y)\alpha|_N,\qquad
 c=y\alpha(y)\in N;\qquad
 \beta^2=\operatorname{Inn}(c),\quad \beta(c)=c.
                                                        \tag{D.37i}
\]
Then $E_\tau$ is exactly the twisted extension generated by
$N$ and $ys$, whose square is $c$; the multiplication rule
is the rule of (D.29.9) with these named $\beta,c$.

Conversely let a marked four-cusp group carry the specified
$\delta$ and an involution $\tau$ with the permutation and
multiplier in (D.37h), satisfying $[\delta,\tau]=1$.
The distinct permutations make
$\Gamma=\langle\delta,\tau\rangle\cong V_4$.
Take its center-free extension
\[
 K=\{(\varphi,\gamma)\in\Aut(N)\times\Gamma:
                         [\varphi]=\gamma\}.
                                                        \tag{D.37j}
\]
Here $N$ embeds by its inner automorphisms.  The preimage
$K_\delta$ of $\langle\delta\rangle$ is identified with
the exact group $F$ in (D.37a) by
\[
 f\longmapsto
       (\operatorname{Inn}(f)|_N,\delta^{\chi(f)}).
                                                        \tag{D.37k}
\]
Injectivity follows from $C_F(N)=1$; surjectivity follows
because an automorphism in either of these two outer
classes differs from $1$ or $\operatorname{Inn}(y)|_N$
by an inner automorphism from $N$.
The quotient $K/F=C_2$ therefore gives an outer action on
this very $F$.  It is faithful, since any element of $K$
centralizing $F$ centralizes its inner copy of $N$, and
an automorphism of the center-free $N$ with that property
is the identity.

We check descent of the peripheral powers, not only of the
four unordered procyclic subgroups.  The two classes
$A_0,A_1$ become the one $F$-class of $x$.
The other classes are $B=y^2$ and $D=z^2$.
For a representative of $\tau$, the induced automorphism
$\psi$ of $F$ satisfies
$\psi(y)^2=h z^{2\lambda}h^{-1}$ for some $h\in N$.
A.7 implies
$C_F(hz^{2\lambda}h^{-1})=h\langle z\rangle h^{-1}$,
so $\psi(y)=hz^\mu h^{-1}$ and $2\mu=2\lambda$.
Multiplication by two is injective on $\widehat\Z$;
hence $\mu=\lambda$.  The same proof applies to $z$,
and the $x$ exponent is already given by $A_0$.
Thus the descended action fixes $x$, exchanges $y,z$,
and has exactly the original $\lambda$, in particular
$\lambda_2=-1$.

For completeness its normalizer construction (D.37d)--(D.37e)
and (D.37f) also recover the torsion-free assertion upstairs,
without making that assertion an extra converse hypothesis.
An order-two lift in the normalizer of $\langle x\rangle$
fixes $A_0$, so its image in $\Gamma$ is $\delta\tau$,
not $\tau$.  All torsion has this same image by (D.37f).
The preimage of $\langle\tau\rangle$ is therefore torsion-free.
The conjugation maps (D.37j)--(D.37k) identify the recovered
extension with the original one on every element; conversely
restriction to $N$ recovers the given $\Gamma$ and both
named generators.  This proves the asserted inverse
correspondence of marked extensions.

We finish by making the obstacle for the D.29 gluing explicit.
Write its ordered cusps as $1,2,3,4$, so that
$\sigma(\tau)=(1\ 3)(2\ 4)$.  The unordered two-pants
partitions and their images are
\[
\begin{array}{c|c|c}
 \text{partition}&\tau\text{ on its two blocks}
                     &(1\ 3)\text{ or }(2\ 4)\\\hline
 12\mid34&\text{exchanges}&14\mid23\\
 13\mid24&\text{preserves each}&13\mid24\\
 14\mid23&\text{exchanges}&12\mid34.
\end{array}                                            \tag{D.37l}
\]
The only single transpositions commuting with
$(1\ 3)(2\ 4)$ are $(1\ 3)$ and $(2\ 4)$:
a commuting permutation must send its two two-cycles to
two-cycles, and a transposition can do this only by
exchanging the members of one of them.  Hence a deck
involution of the required permutation cannot act on the
same two-vertex pants graph as the inverse-pair gluing,
under any relabeling or choice of its exchanged partition.

One can nevertheless write the required geometric deck
map in the D.29 free basis.  For
$(u_1,u_2,u_3,u_4)=(a,a^{-1}l,b,b^{-1}l^{-1})$, the
change of four-cusp marking
\[
 A_0=a,\quad A_1=(a^{-1}l)b(a^{-1}l)^{-1},\quad
 B=a^{-1}l,\quad D=b^{-1}l^{-1}
                                                        \tag{D.37m}
\]
satisfies $A_0A_1BD=1$.  Its inverse substitutions are
$a=A_0$, $l=A_0B$, $b=B^{-1}A_1B$.
Transporting (D.37c) through these maps gives
\[
 \Delta(l)=a^{-1}lb,\qquad
 \Delta(a)=a^{-1}lbl^{-1}a,\qquad
 \Delta(b)=a,\qquad
 \Delta^2=\operatorname{Inn}(a^{-1}l).
                                                        \tag{D.37n}
\]
It has multiplier $1$ and permutation $(1\ 3)$.
The existence question for D.29 is the existence of a
deck action of this marked geometric kind commuting
outerly with its actual $\alpha$; (D.37l) alone does
not decide it.

In particular failure of commutation for one selected
deck representative cannot settle the question.  In the
geometric specialization $\lambda=-1$, $\epsilon=1$,
$\kappa=1$, $f=1$, D.29 gives
\[
 \alpha(l)=l,\quad\alpha(a)=b^{-1},\quad
                         \alpha(b)=la^{-1}l^{-1}.
\]
For (D.37n), the words
$\alpha\Delta(l)=bl^2a^{-1}l^{-1}$ and
$\Delta\alpha(l)=a^{-1}lb$ are not conjugate even
in the full profinite group: the finite quotient
$l\mapsto(1\ 2\ 3)$, $a,b\mapsto(2\ 3)$ in $S_3$
sends the first to $1$ and the second to a nonidentity
three-cycle.  These two automorphisms therefore do not
commute outerly.  Yet another geometric deck does commute:
\[
 \begin{aligned}
 \Delta_0(l)&=bla^{-1},\qquad \Delta_0(a)=b,\qquad
                  \Delta_0(b)=blal^{-1}b^{-1},\\
 \Delta_0^2&=\operatorname{Inn}(bl),\qquad
       \alpha\Delta_0=\operatorname{Inn}(b^{-1})\Delta_0\alpha.
 \end{aligned}
                                                        \tag{D.37o}
\]
Its origin is also explicit.  In the geometric extension
$E=\widehat F(a,t)$, put $l=t^2$ and
$b=ta^{-1}t^{-1}$.  The involution
$\eta(a)=a^{-1}$, $\eta(t)=a^{-1}t$ preserves the
index-two kernel.  Restricting
$\operatorname{Inn}(t)\eta$ to that kernel gives
exactly (D.37o).  The formulas prove its peripheral
permutation and multiplier as well as the two displayed
composition identities.  Thus this geometric
specialization supplies an explicit example of the
separator being moved.
To verify directly that $\Delta_0$ is a deck of the
specified geometric kind, adjoin its lift $r$ with
$r^2=bl$.  The extension presentation eliminates to
$\widehat F(a,r)$ by
$b=rar^{-1}$ and $l=b^{-1}r^2=ra^{-1}r$.
The peripheral generators downstairs are
$a,a^{-1}r,r^{-1}$, whose product is one.
Their lifts are the two classes of $a,b$, the square
$(a^{-1}r)^2=a^{-1}l$, and the square $r^{-2}$,
conjugate to $b^{-1}l^{-1}$.  These substitutions identify
its four-cusp inclusion with (D.37a)--(D.37b).

The executable \texttt{code/part\_c/negative\_transposition\_d37.py}
verifies the free-word covering and inverse substitutions,
the deck squares, all the equations in (D.37o), the
two-dimensional mod-two cohomology calculation, and the
complete permutation calculation (D.37l) under all
relabelings.  The extension and torsion arguments above
prove the statements for the full profinite groups.
\end{proof}

\begin{claim}{D.38}{Proved}{A nonnilpotent finite obstruction on the universal \(A_5\) quotient}
Let $R$ be the intersection of the kernels of all continuous
epimorphisms $\widehat F(x,y)\twoheadrightarrow A_5$.
This is an open characteristic subgroup, and the quotient
specified by the nineteen generating-pair representatives in
the accompanying program is
\[
                 Q=F/R\xrightarrow{\ \sim\ } A_5^{19}.
                                                               \tag{D.38a}
\]
Write $X,Y,Z=(XY)^{-1}$ for the images of the three cusps,
and $\theta,\omega$ for the automorphisms of $Q$ induced by
$x\leftrightarrow y$ and $x\mapsto y,\ y\mapsto(xy)^{-1}$.
The image of every full $GT_0$ involution belongs to the
following explicitly computed finite set:
\[
\begin{split}
 \mathcal T=\{t\in Q:\quad&
 t^2=1,\quad tXt^{-1}=X^{-1},\\
 &\omega(t)=Y^{-1}\theta(t)\}.
\end{split}                                                     \tag{D.38b}
\]
Its normalized coordinates are
$\alpha_Q=\operatorname{Inn}(t)$ and
$f_Q=\theta(t)t^{-1}$.
Every tuple in $\mathcal T$ satisfies the finite
normalized relations (I), (II), and
$f_Q\alpha_Q(f_Q)=1$, and defines an actual automorphism
of $Q$ of order two.

The executed exact enumeration gives
\[
\begin{array}{l|r}
\text{finite set or test}&\text{cardinality}\\ \hline
\mathcal T&10\,125\,000\\
\text{tuples surviving the full cusp-triangle obstruction}
       &5\,062\,500\\
\text{tuples conjugate to standard inversion in the finite
pure model (D.38f)}&2\,531\,250\\
\text{triangle survivors outside that finite conjugacy set}
       &2\,531\,250 .
\end{array}                                                     \tag{D.38c}
\]
Thus $5\,062\,500$ explicit finite tuples satisfying the
finite involution equations have no full pure-involution
lift.  The last row is disjoint from the image of the
standard full conjugacy class.  No tuple in that last
row is asserted to have a full lift.
\end{claim}
\passage{\S{}31; folios 198-199verso; scans 391-394; Part D(b): actual full GT\_0 torsion and conjugacy}

\begin{proof}
The program enumerates all $60^2$ ordered pairs in $A_5$,
tests generation by closure under multiplication, and
obtains $2\,280$ generating pairs.
Simultaneous conjugation by all $120$ elements of $S_5$
has nineteen free orbits on this set.
These are all $\Aut(A_5)$-orbits: an automorphism
permutes the five Sylow-$2$ subgroups, and equivariance
with the faithful inner $A_5$ action makes that
permutation representation of $\Aut(A_5)$ injective;
conjugation by $S_5$ already supplies its $120$ elements,
as checked in D.31.
Two epimorphisms to $A_5$ have the same kernel exactly
when they differ by an automorphism of $A_5$.
The nineteen kernels are therefore distinct.
The joint quotient map is a subdirect product of
nineteen copies of the nonabelian simple group $A_5$.
The elementary subdirect-product argument shows that
a proper subdirect product of copies of one simple
group identifies two coordinates by an automorphism:
induct on the number of factors and apply the
normal-subgroup description of a direct product of
nonabelian simple groups to the last projection.
Such an identification would equate two of the
nineteen kernels.  This proves surjectivity of the
specified map to $A_5^{19}$ and hence (D.38a).
The kernel is characteristic because every automorphism
of $F$ permutes all epimorphisms to $A_5$.

For each representative pair $(X_i,Y_i)$ the program
also records the ordered triple of conjugacy classes
of $X_i,Y_i,Z_i$ in $A_5$, allowing the simultaneous
exchange of the two classes of elements of order five
by odd conjugation in $S_5$.
All nineteen resulting triples are different.
Every triple contains an element of order five,
and has at least two elements of odd order.
These assertions are checked by complete finite
enumeration, with representatives and component
permutations retained in the output.

By D.34 a full pure involution has multiplier $-1$.
Inversion preserves each conjugacy class of $A_5$:
this is immediate for involutions, and reversal of
a $3$-cycle or $5$-cycle is achieved by an even
permutation.
An automorphism of $A_5^{19}$ permutes its minimal
normal simple factors and acts on each factor by
an element of $\Aut(A_5)=S_5$.
Preservation of the three cusp classes therefore
forces its factor permutation to preserve the
recorded class triple.  The nineteen distinct
triples force the permutation to be the identity.
Since each triple contains an order-five entry,
the component automorphisms must have even parity;
odd conjugation exchanges its two order-five
classes.  Thus the induced automorphism of $Q$
is $\operatorname{Inn}(t)$ for a unique $t\in Q$.
The group $Q$ has trivial center, so its
involutivity gives $t^2=1$.

The exact Harbater--Schneps automorphism identities
from \cite[Propositions 3--4]{R-HS}, with
$m=(-1-1)/2=-1$, descend to $Q$:
\[
 \theta\alpha_Q\theta^{-1}\alpha_Q^{-1}
      =\operatorname{Inn}(f_Q),\qquad
 \omega\alpha_Q\omega^{-1}\alpha_Q^{-1}
      =\operatorname{Inn}(Y^{-1}f_Q).
\]
Center-freeness yields
$f_Q=\theta(t)t^{-1}$ and
$\omega(t)t^{-1}=Y^{-1}f_Q$.
Together with the normalized equation on $X$ these
give precisely (D.38b).

Conversely take $t$ in (D.38b) and define
$f=\theta(t)t^{-1}$.
Then $\theta(t)$ inverts $Y$, so
\[
 tYt^{-1}=f^{-1}Y^{-1}f,\qquad
 \theta(f)f=1,\qquad f\operatorname{Inn}(t)(f)=1.
\]
The second relation in (D.38b) involving $\omega$
can be rewritten $f=Y\omega(t)t^{-1}$.
Apply $\omega$ and $\omega^2$ and use $\omega^3=1$:
\[
 \omega(f)=Z\omega^2(t)\omega(t)^{-1},\qquad
 \omega^2(f)=Xt\omega^2(t)^{-1}.
\]
Substitution now proves the complete finite
three-cycle relation by cancellation,
\[
       \omega^2(f)Z^{-1}\omega(f)Y^{-1}fX^{-1}=1.
                                                               \tag{D.38d}
\]
The tuple $X$ has components of odd order, so
$tXt^{-1}=X^{-1}$ excludes $t=1$ in $Q$.
Thus $\operatorname{Inn}(t)$ has exact order two.
This verifies the assertions for every member
of the finite set, independently of any lifting.

We next impose an obstruction from the full
group.  For each component the extension quotient
is explicitly
\[
 A_5\rtimes_{\operatorname{Inn}(t_i)}C_2
       \xrightarrow{\ \sim\ }A_5\times C_2,\qquad
                (g,s^j)\longmapsto(g t_i^j,s^j).
                                                               \tag{D.38e}
\]
The three vertices of the full cusp triangle in
D.34 are represented by $s$, $xs$, and $yf s$.
Their first coordinates under (D.38e) are
$t_i$, $X_i t_i$, and $Y_i f_i t_i$.
None can be the identity in $A_5$.
Indeed that would make the image of the
corresponding section central in $A_5\times C_2$.
This section class is incident to two different
cusp edges of the full triangle.  At either
edge choose its representative in the cusp
normalizer.  Its image is still central, because
centrality is unchanged by conjugation.
It must therefore both act trivially on the
cusp image and act by inversion there.
Both incident cusp images would have order
at most two.  No cusp image is trivial for
a generating $A_5$ pair.  Two involutions
generate a dihedral group, while any two
entries of the generating triple generate
$A_5$.  This is a contradiction.
Possible fusion of distinct section classes
in the finite quotient does not change this
argument.

For all components together, this obstruction
is exactly
$t_i\notin\{1,X_i\}$ for every $i$.
The third vertex condition follows from these:
$f_i t_i=\theta(t)_i$ and $\theta(X)=Y$,
so $Y_i\theta(t)_i=1$ is another occurrence
of $t_j=X_j$ under the component action of
$\theta$.

Here are the complete finite counts, with the
factorization used to execute them.
The relation $\omega(t)=Y^{-1}\theta(t)$
joins components in twelve independent
sets, using indices $0,\ldots,18$:
\[
 \{0,2\},\{1\},\{3,5\},\{4\},\{6,7\},\{8\},
 \{9,17\},\{10\},\{11,14\},\{12,18\},\{13,15\},\{16\}.
\]
Exhaustive propagation of each initial component
value through its specified $S_5$ conjugation
maps gives respectively
\[
                 3,3,3,3,4,2,5,5,5,5,5,5
\]
assignments.  The equations on a component set
involve no other set, so their Cartesian product
has exactly $10\,125\,000$ elements.
The triangle obstruction changes only the
count for $\{6,7\}$, from four to two.
Its surviving cardinality is therefore
$5\,062\,500$.
All forty-eight component assignments are
executed and checked, including the original
equations (D.38d) and the finite norm; the
Cartesian-product count is exact.

For precision, the finite conjugacy model in
(D.38c) consists of the following actual
automorphisms of $Q$ with a common unit exponent:
\[
\begin{split}
 \mathcal P_Q=\{(\mu,\gamma)\in
      (\Z/30)^\times\times\Aut(Q):\quad&
 \gamma(X)=X^\mu,\\
 &\gamma(Y)\sim_QY^\mu,\quad
 \gamma(Z)\sim_QZ^\mu\}.
\end{split}                                                     \tag{D.38f}
\]
Every normalized full pure automorphism
maps to this set, since $R$ is characteristic.
The recorded class triples again imply that
an element of (D.38f) acts on the $i$th
factor by conjugation by some $g_i\in S_5$,
without permuting factors.
Its parity is even for $\mu\bmod5\in\{1,4\}$
and odd for $\mu\bmod5\in\{2,3\}$.
The remaining condition is exactly
$g_iX_i g_i^{-1}=X_i^\mu$; with that parity,
the $Y_i$ and $Z_i$ class conditions follow.

The standard inversion induces
$\operatorname{Inn}(d)$ on $Q$, where the
program supplies each simultaneous inverter
$d_i$ of $(X_i,Y_i)$.
The conjugacy condition is the finite equation
$t_i=g_i d_i g_i^{-1}$ for every $i$.
The program tests all $120$ choices of $g_i$
for each of the eight units modulo thirty.
There are exactly two distinct compatible
patterns, each with component counts
\[
                  3,3,3,3,1,1,5,5,5,5,5,5.
\]
They are disjoint, so their total cardinality
is $2\cdot1\,265\,625=2\,531\,250$.
This proves the last two rows of (D.38c).
It proves an upper constraint on the image
of the known full conjugacy class, not
surjectivity from that class onto the
finite model.

The complete implementation is
\path{code/part_a/a5_involution_quotient.py};
its output retains the nineteen pairs,
the $\theta,\omega$ component maps, the
standard inversion tuple, every component
count, and an explicit rejected tuple.
No family of lifts through all further
characteristic quotients has been supplied
for a surviving tuple.
\end{proof}

\begin{claim}{D.39}{Proved}{A rank-two Haran involution admitting neither hyperbolic loop type}
Let $\epsilon,\kappa\in\widehat\Z$ have the coordinates
\[
 \epsilon_3=-1,\quad\epsilon_q=1\ (q\ne3),\qquad
 \kappa_2=1,\quad\kappa_q=0\ (q\ne2).
\]
There is an actual involution of the rank-two free
profinite group $H=\widehat F(a,b)$ given by
\[
 \alpha(a)=b^{-1}a^\epsilon b,\qquad
 \alpha(b)=b^{-1}a^\kappa.                                  \tag{D.39a}
\]
Its split extension has the explicit identification
\[
 H\rtimes_\alpha C_2\ \xrightarrow{\ \sim\ }\ N\amalg C_2,
 \qquad
 N=(\Z_3\rtimes\Z_2)\times\prod_{q\ne2,3}\Z_q,               \tag{D.39b}
\]
where parity in $\Z_2$ acts by inversion on $\Z_3$.
Its actual fixed subgroup $C_H(\alpha)$ is trivial.
The outer involution preserves no three-cusp
marking and no once-punctured-torus marking of $H$,
under any full profinite change of free basis.

For the natural once-punctured-torus marking
$\overline{\langle[a,b]\rangle}$, the peripheral
condition also fails: a quotient of order $21$
separates $\alpha([a,b])$ from every conjugate
of $[a,b]^{-\epsilon}$.
\end{claim}
\passage{\S{}31; folios 198-199verso; scans 391-394; Part D: full rank-two loop actions and precise peripheral conditions}

\begin{proof}
The relations $\epsilon^2=1$ and
$\epsilon\kappa=\kappa$ give, by substitution
on the specified free generators,
\[
 \begin{split}
 \alpha^2(a)
  &=(a^{-\kappa}b)(b^{-1}ab)(b^{-1}a^\kappa)=a,\\
 \alpha^2(b)
  &=(a^{-\kappa}b)(b^{-1}a^{\epsilon\kappa}b)=b.
 \end{split}
\]
Thus the universal continuous endomorphism
defined by (D.39a) is an automorphism with
square one.
On full abelianization its matrix is
\[
       A=\begin{pmatrix}\epsilon&\kappa\\0&-1\end{pmatrix},
       \qquad \det A=-\epsilon,\qquad
       \operatorname{tr}A=\epsilon-1.                       \tag{D.39c}
\]
It is nonidentity on abelianization, so both
the automorphism and its outer class have
exact order two.

For the exact identification (D.39b), use
the normal form for $N$ verified in D.26:
\[
 N=\overline{\langle l,t\rangle},\qquad
 tlt^{-1}=l^\epsilon,\qquad t^2=l^\kappa.
\]
The map $l^u t^i\mapsto
(u_3,2u_2+i,(u_q)_{q\ne2,3})$, $i\in\{0,1\}$,
is its displayed isomorphism with the group
in (D.39b).  The parity character sends
$l$ to zero and $t$ to one.
In $N\amalg\langle r\mid r^2=1\rangle$
extend this character by $r\mapsto1$.
The substitutions
\[
             a\longmapsto l,\qquad
             b\longmapsto tr,\qquad
             s\longmapsto r
\]
define a homomorphism from
$H\rtimes_\alpha\langle s\rangle$ to that
free profinite product.  To check its
relations, note that $r=(tr)^{-1}t$,
so
\[
 rlr=(tr)^{-1}l^\epsilon(tr),\qquad
 r(tr)r=(tr)^{-1}l^\kappa.
\]
In the reverse direction send
$l\mapsto a$, $t\mapsto bs$, and $r\mapsto s$.
Here
$(bs)^2=b\alpha(b)=a^\kappa$ and
$(bs)a(bs)^{-1}=b\alpha(a)b^{-1}=a^\epsilon$.
The normal-form multiplication of $N$
therefore gives its continuous homomorphism
to the semidirect product, and the free
product property extends it over $r$.
The two composites fix every displayed
generator, proving the isomorphism.
In particular the parity kernel is exactly
the free profinite group on the stated
generators $a,b$, and the extension has
the actual section $\langle r\rangle$.

The free factor $\langle r\rangle$ is
malnormal by the full-profinite
free-factor intersection property
\cite[Theorem 5 of the inspected author
version]{A-HerfortZalesskii2008}.
Its centralizer in $N\amalg C_2$ is
therefore precisely $\langle r\rangle$.
Intersecting with the parity kernel
proves $C_H(\alpha)=1$.

In fact every involution $u\in H\rtimes_\alpha C_2$
has $C_H(u)=1$.  To prove this with the specified
index-two kernel, reduce (D.39c) modulo two.
The resulting matrix is
$\left(\begin{smallmatrix}1&1\\0&1\end{smallmatrix}\right)$.
Thus its dual action on $H^1(H,\F_2)$ is the
regular two-dimensional $\F_2[C_2]$ module $J_2$.
Its positive-degree $C_2$ cohomology vanishes:
both differentials of the periodic resolution are
$1+\alpha$, and their image equals their kernel,
both of dimension one.
Since $H$ has cohomological dimension one and
the extension is split, the Hochschild--Serre
spectral sequence gives
\[
       \dim H^n(H\rtimes_\alpha C_2,\F_2)=1
                        \qquad(n\geq3).                 \tag{D.39g}
\]
The full profinite Brown restriction theorem
\cite[Theorem 2 and \S1.1]{A-KochloukovaZalesskii2007},
applied exactly as in (D.37f), expresses this
cohomology as the product over involution classes.
Every class contributes
$1+\dim H^1(C_H(u),\F_2)$, because its normalizer
is $C_H(u)\times C_2$.
The existing section $\langle s\rangle$ therefore
accounts for the unique conjugacy class.
Conjugacy by the extension agrees with conjugacy
by $H$: for any involution $u$ the decomposition
$E=H\langle u\rangle$ absorbs its second factor,
which centralizes $u$.
Every involution is consequently $H$-conjugate
to $s$, and the already proved equality
$C_H(s)=1$ gives
\[
                  C_H(u)=1
           \quad\text{for every involution }u\in E.     \tag{D.39h}
\]

There is a basis-independent obstruction
to a three-cusp structure.  An involution
on three cusps has either the identity
permutation or a transposition.
In the former case its homology is a
scalar matrix $\lambda I$ and has trace
$2\lambda$, whose $2$-coordinate is
nonzero.  This contradicts the zero
$2$-coordinate of the trace in (D.39c).
For a transposition the reduced
permutation module has trace zero, and
multiplication by a common unit keeps
that trace zero.  This contradicts the
$3$-coordinate $-2$ of (D.39c).
Trace is unchanged by conjugating the
homology matrix under a change of free
basis.  Hence neither peripheral
permutation can occur.

There is also a basis-independent obstruction
to a once-punctured-torus marking.
Take any full profinite free basis $a',b'$ of
$H$ and put $c'=[a',b']$,
$L=\overline{\langle c'\rangle}$.
Preservation of the conjugacy class of $L$
would give the multiplier
$\nu=\det A=-\epsilon$.
The exact comparison is the one in
(D.33h): in every finite class-two
Heisenberg quotient the commutator
exponent is the determinant of the
abelianization matrix, so equality for
all moduli identifies these elements
of $\widehat\Z^\times$.
Changing the free basis conjugates the
matrix and leaves its determinant unchanged.

The normalizer $N_E(L)$ surjects onto
$E/H=C_2$, and its kernel is $L$ by A.7
transported through this free-basis isomorphism.
For a lift $v$ of the nonidentity element,
write
\[
  vc'v^{-1}=(c')^\nu,\qquad
  v^2=(c')^d,
  \qquad(\nu-1)d=0.                                    \tag{D.39i}
\]
Here $\nu_3=1$ and $\nu_q=-1$ for every
$q\ne3$.  Therefore $d_q=0$ for $q\ne3$,
including $q=2$, since multiplication by two
is injective on $\Z_q$.
Choose $j\in\widehat\Z$ with
$j_3=-d_3/2$ and $j_q=0$ otherwise.
The lift $u=(c')^jv$ then satisfies
\[
 u^2=(c')^{(1+\nu)j+d}=1,
 \qquad
 u(c')^{e_3w}u^{-1}=(c')^{e_3w}
                       \quad(w\in\Z_3),                 \tag{D.39j}
\]
where $e_3$ is the idempotent supported at
the prime three.
The lift $u$ has exact order two because
its image in $E/H$ is nonidentity.
The peripheral injection
$\widehat\Z\xrightarrow{\sim}L$ makes
$(c')^{e_3\Z_3}$ a nontrivial subgroup of
$C_H(u)$, contradicting (D.39h).
This excludes every once-punctured-torus
marking, including those obtained from
arbitrary full profinite free bases.

For an additional finite certificate in
the original basis, put
$c=[a,b]=aba^{-1}b^{-1}$.

Use the finite group
\[
 Q=\F_7\rtimes C_3,\qquad
 (v,i)(w,j)=(v+2^i w,\ i+j),
                                                               \tag{D.39d}
\]
where the second coordinate is modulo
three and $2$ has order three in
$\F_7^\times$.
The assignment
\[
                    a\mapsto(0,1),\qquad b\mapsto(1,1)
                                                               \tag{D.39e}
\]
is an epimorphism: $a,b$ have order
three, their commutator is $(1,0)$,
and these elements generate the
normal $\F_7$ and the $C_3$ quotient.
Since $\epsilon\equiv-1\pmod3$
and $\kappa\equiv0\pmod3$, (D.39a)
has the following exact images:
\[
 \begin{array}{c|ccccc}
 \text{word}&c&\alpha(a)&\alpha(b)&\alpha(c)&c^{-\epsilon}\\ \hline
 \text{image}&(1,0)&(5,2)&(3,2)&(1,0)&(6,0).
 \end{array}                                                   \tag{D.39f}
\]
The last entry uses
$(-\epsilon)_7=-1$.
Conjugacy on the normal $\F_7$ in
(D.39d) is multiplication by a
power of two.  The two relevant
orbits are exactly
\[
             \{(1,0),(2,0),(4,0)\},\qquad
             \{(3,0),(5,0),(6,0)\}.
\]
They are disjoint.  Thus the images
of $\alpha(c)$ and $c^{-\epsilon}$
are not conjugate, proving failure
of the natural torus cusp.

The program
\path{code/part_a/haran_torus_test.py}
executes this order-$21$ quotient
and an independent matrix check
inside $\operatorname{PSL}_2(\F_7)$.
For the latter it enumerates all
$168$ elements and uses
\[
 a=\begin{pmatrix}0&1\\6&1\end{pmatrix},
 \qquad b=\begin{pmatrix}0&2\\3&1\end{pmatrix}
 \quad\bmod\{\pm I\}.
\]
Both have order three,
$c=\left(\begin{smallmatrix}1&1\\3&4\end{smallmatrix}\right)$
has order seven, and
$\alpha(c)=
\left(\begin{smallmatrix}0&3\\2&2\end{smallmatrix}\right)$.
There are seven conjugators from
$c$ to this image and none from
$c^{-1}$.  The matrix calculation
agrees with (D.39f).
The finite calculation detects the original
cusp explicitly.  The argument
(D.39g)--(D.39j) supplies the exclusion
for every alternative marking.
The program also verifies the regular
mod-two module used in (D.39g).
\end{proof}

\begin{claim}{D.40}{Proved}{A finite homomorphism-count obstruction to the extra deck symmetry}
The mixed-multiplier four-cusp involution constructed in D.29
commutes with no single-transposition involution of multiplier $1$ in its
full loop-outer group. In particular it admits none of the
specified geometric deck actions required for the descent of D.37,
including actions moving the two-pants separator.

More precisely, for the exact profinite group
\[
 E=\overline{\langle a,l,t\mid
       tlt^{-1}=l^\epsilon,\ t^2=l^\kappa\rangle},
 \qquad p=a^{-1}l,
                                                        \tag{D.40a}
\]
with $\epsilon,\kappa$ as in D.29 and $\lambda=-\epsilon$,
there is no continuous automorphism sending the conjugacy
classes of $a,p$ to those of $a^\lambda,p$, respectively.
There is also none sending those classes to $a,p^\lambda$.
The obstruction is the exact inequality
\[
 \begin{aligned}
 \#\{\phi:E\to Q:\phi(a),\phi(p)\in C\}&=10200,\\
 \#\{\phi:E\to Q:\phi(a)\in C^{-1},\phi(p)\in C\}&=4920,
 \end{aligned}
                                                        \tag{D.40b}
\]
where all homomorphisms are continuous,
$Q=\operatorname{PSL}_2(\F_{11})$, and $C$ is the
conjugacy class of
$\left(\begin{smallmatrix}0&1\\10&2\end{smallmatrix}\right)$
modulo scalar matrices $\{\pm I\}$.
\end{claim}
\passage{Section 24; folio 147 verso; scan 293. Section 30; folios 194-196; scans 383-387. Section 31; folios 198-199 verso; scans 391-394. Part D(b) and D(d).}

\begin{proof}
We first identify the group and its peripheral marking.
In the basis $(l,a,b')$ of (D.29.10), adjoining the outer
lift $t$ gives the relations
\[
 tlt^{-1}=l^\epsilon,\quad tat^{-1}=b',\quad
 tb't^{-1}=l^\kappa a l^{-\kappa},\quad t^2=l^\kappa.
\]
Eliminating $b'$ gives (D.40a): the third relation follows
from the second and fourth. Conversely $b'=tat^{-1}$
defines the inverse homomorphism from (D.40a) to this
presentation. Both composites fix the named generators.
Thus this is an isomorphism of profinite groups, and
\[
                         E=N_{\rm bad}\amalg\widehat\Z(a)
                                                        \tag{D.40c}
\]
with the first factor generated by $l,t$.
The index-two subgroup $H$ is the kernel of the character
sending $t$ to $1$ and $a,l$ to $0$ in $C_2$.

The four $H$-peripheral classes of D.29 are represented by
\[
 a,\quad p=a^{-1}l,\quad b,\quad b^{-1}l^{-1}.
\]
Since $\alpha$ interchanges the first and third classes,
and the second and fourth, with multiplier $\lambda$,
and conjugation by $t$ is $\alpha$, these are equally
the classes of
\[
 a,\quad p,\quad (tat^{-1})^\lambda,
                         \quad(tpt^{-1})^\lambda.       \tag{D.40d}
\]
Here $\lambda^2=1$, and the original cusp conjugators
belong to $H$; thus (D.40d) records actual conjugacy
classes under the specified inclusion into $E$.
This description is independent of the Galois conjugator
$f$ used to construct the two pants isomorphisms.

For a finite group $Q$ define the central function
\[
 F_Q(L)=\#\{T\in Q:T L T^{-1}=L^\epsilon,
                                      \ T^2=L^\kappa\}.
                                                        \tag{D.40e}
\]
Every assignment $a\mapsto A$, $p\mapsto P$ forces
$l\mapsto AP$. The presentation (D.40a) then gives
exactly $F_Q(AP)$ continuous homomorphisms to $Q$.
For two conjugacy classes $C,D$ put
\[
             K_Q(C,D)=\sum_{A\in C,\ P\in D}F_Q(AP).
                                                        \tag{D.40f}
\]
Conjugating $L$ conjugates all its possible $T$, so
$F_Q$ is central. Since $PA=A^{-1}(AP)A$, this also proves
$K_Q(C,D)=K_Q(D,C)$.

Precomposition by a continuous automorphism of $E$ is
a bijection on its continuous homomorphisms to $Q$.
An automorphism with the first pair of class images in
the claim would therefore give
$K_Q(C,D)=K_Q(C^\lambda,D)$ for every pair $C,D$.
The inverse map has the same power effect because
$\lambda^2=1$, so this is equality of the two full sets,
without a choice of conjugators in the target.
The second proposed class effect would similarly give
$K_Q(C,D)=K_Q(C,D^\lambda)$.

We compute these integers in one completely specified
finite group. Enumerate the determinant-one matrices
over $\F_{11}$, and identify each matrix with its
negative, using their lexicographically smaller
representative. Matrix multiplication followed by this
identification constructs all $660$ elements of $Q$.
The class $C$ has size $60$ and order $11$.
Its inverse class is distinct and has the same size.
Since $\lambda_{11}=-1$, we have $C^\lambda=C^{-1}$.
The following table gives the full product counts needed
in (D.40f). The representative of each nonidentity
column is $\left(\begin{smallmatrix}0&1\\10&r\end{smallmatrix}\right)$;
the last column is the identity. We write
$N_{C,D}(L)=\#\{(A,P)\in C\times D:AP\in[L]\}$.
\[
\begin{array}{c|rrrrrrrr}
 r&0&1&2&3&4&5&9&I\\\hline
 |L|&2&3&11&5&5&6&11&1\\
 \text{class size}&55&110&60&132&132&110&60&1\\
 F_Q(L)&0&6&1&1&1&0&1&56\\
 N_{C,C}(L)
   &0&1320&120&660&660&0&840&0\\
 N_{C^{-1},C}(L)
   &660&0&120&660&660&1320&120&60
\end{array}                                               \tag{D.40g}
\]
Multiplying either of the last two rows by the
$F_Q$ row and summing gives respectively
\[
 7920+120+660+660+840=10200,
 \qquad 120+660+660+120+3360=4920.
\]
These unequal integers prove the first nonexistence
assertion. Symmetry of $K_Q$ proves the second.

All entries of (D.40g) are verified by the finite
enumeration in \path{code/root/haran_pair_counts.py}.
It constructs the multiplication, inverse, power and
conjugacy tables directly from the displayed matrices,
computes the residues of $\epsilon,\kappa$ modulo each
element order, and tests both equations in (D.40e) for
every pair $(L,T)$. It then counts each product in the
two displayed class pairs. Two independent computations,
\path{code/part_a/haran_pair_count_audit.py} and
\path{code/part_c/deck_obstruction_psl11_independent.py},
instead generate the faithful twelve-point permutation
action by $z\mapsto z+1$ and $z\mapsto-1/z$ on
$\mathbf P^1(\F_{11})$. The first enumerates allowed
$(L,T)$ before setting $P=A^{-1}L$; the second also
enumerates the triples directly. Both reproduce (D.40b).
The attached outputs include the full integer data.

Finally suppose a positive single-transposition
involution $\delta$ commuted with the D.29 involution
$\tau$. Their distinct peripheral permutations make
$\langle\delta,\tau\rangle\simeq V_4$.
Its center-free pullback extension inside $\Aut(H)$
contains $E$ as the normal preimage of $\langle\tau\rangle$,
using the canonical conjugation identification of
center-free extensions. Conjugation by any lift of
$\delta$ therefore defines a continuous automorphism
$\eta$ of this exact $E$.
By D.37's permutation computation, $\delta$ is either
$(1\ 3)$ or $(2\ 4)$ in the labeling (D.40d).
Its positive multiplier is $1$.
In the first case $\eta$ sends $a$ to an $E$-conjugate
of $a^\lambda$ and $p$ to an $E$-conjugate of $p$;
in the second case the roles are reversed.
Both have just been excluded. The argument imposes
no preservation of the original pants separator and
no order condition on the chosen lift of $\delta$.
It proves the stated obstruction for this actual
four-cusp construction; the general negative
three-cusp conjugacy question remains in D.12.
\end{proof}

\begin{claim}{D.41}{Proved}{The global multiplier of a negative transposition}
An involution of the full three-cusp loop-outer group whose
puncture permutation is a transposition and whose $2$-adic
multiplier is $-1$ has the global multiplier
$-1\in\widehat\Z^\times$.
Together with D.34 and D.36, this identifies the multiplier
of every involution in that full group with its global
$2$-adic sign.  Consequently every finite subgroup generated
by involutions has the usual global orientation character.
\end{claim}
\passage{Part D(b) of the task; \S{}24, folio 147 verso, scan 293; \S{}30, folios 194--196, scans 383--387; \S{}31, folios 198--199 verso, scans 391--394.}

\begin{proof}
Label the original peripheral generators as
$x,y,z=(xy)^{-1}$ so that the involution fixes $x$ and
exchanges $y,z$.  Write its common multiplier as $\lambda$.
It satisfies $\lambda^2=1$, so every coordinate $\lambda_p$
is $1$ or $-1$.  The peripheral-normalizer construction
(D.37d)--(D.37e) supplies an actual order-two automorphism
$\alpha$ of $F=\widehat F(x,y)$ representing it and satisfying
$\alpha(x)=x^\lambda$.

Suppose a positive odd-prime coordinate exists, and let
$q$ be the least odd prime with $\lambda_q=1$.  Use the
exact cyclic quotient
\[
 \chi:F\longrightarrow C_q,\qquad
       \chi(x)=-2,\quad\chi(y)=\chi(z)=1,\qquad N=\ker\chi .
                                                        \tag{D.41a}
\]
The relation $xyz=1$ is respected.  Since $\lambda=1$
modulo $q$ and $\alpha$ exchanges the last two peripheral
classes, $\chi\alpha=\chi$.  The cover corresponding to
$N$ is the completion of the specified connected cyclic
cover of the pair of pants.  Each of its three original
cusps has one lift of degree $q$: all three displayed
monodromies are nonzero in $C_q$.  Its Euler characteristic
is $-q$, so its genus is
\[
                             g=(q-1)/2.
                                                        \tag{D.41b}
\]

Fill the two cusp lifts represented by $y^q,z^q$.
The quotient and its deck extension are explicit:
\[
\begin{split}
 R&=F/\overline{\langle\!\langle y^q,z^q\rangle\!\rangle}_F
       =\langle y,z\mid y^q=z^q=1\rangle_{\rm prof}
       =C_q\amalg C_q,\\
 K&=\ker(R\longrightarrow C_q,\ y,z\longmapsto1)\\
  &=N/\overline{\langle\!\langle y^q,z^q
                                      \rangle\!\rangle}_N
       \cong\widehat\pi_{g,1}.
\end{split}                                             \tag{D.41c}
\]
Here $x=z^{-1}y^{-1}$, and the one remaining cusp is
$c=x^q$.  To verify the equality of normal subgroups
used in (D.41c), note $F=N\langle y\rangle=N\langle z\rangle$.
Conjugating $y^q$ or $z^q$ by any element of $F$ therefore
gives an $N$-conjugate of the same element.  Thus its
closed $F$-normal closure is its closed $N$-normal closure.
The completion comparison is also exact: start with the
discrete cyclic surface cover, fill those two boundaries,
and complete the displayed quotient presentation.
Finite-index subgroup topology identifies its completed
kernel with the specified kernel in $C_q\amalg C_q$.
The inverse presentation substitutions are the original
$y,z$ and $x=z^{-1}y^{-1}$.

The automorphism $\alpha$ preserves this normal subgroup,
because it exchanges the two filled peripheral classes.
It therefore induces an actual automorphism of $R$,
preserves $K$, and induces an involution $\beta$ of $K$
with
\[
                           \beta(c)=c^\lambda.
                                                        \tag{D.41d}
\]
Let $T$ be conjugation by $y$ on $K$.  In $R$, $y^q=1$,
so $T^q=1$ as an actual automorphism.  The equality
$\chi\alpha=\chi$ gives $\alpha(y)y^{-1}\in K$; hence
$\beta T\beta^{-1}$ and $T$ differ by an inner
automorphism of $K$.  In particular their maps on
abelianization commute.

We compute that abelianization as a deck module, without
requiring a comparison of traces.  Reidemeister--Schreier
rewriting with representatives $1,y,\ldots,y^{q-1}$ gives
\[
 u_i=y^i z y^{-(i+1)}\quad(i\in\Z/q),\qquad
 u_0u_1\cdots u_{q-1}=1,\qquad T(u_i)=u_{i+1}.
                                                        \tag{D.41e}
\]
Exponents of $y$ here are interpreted in $R$, where
$y^q=1$.  The first $q-1$ elements form a free profinite
basis; eliminate the last by
$u_{q-1}=(u_0\cdots u_{q-2})^{-1}$.
These substitutions give both maps in the Schreier
presentation.  The remaining boundary also has the exact
word
\[
 c=x^q=\prod_{j=0}^{q-1}u_{\,q-1-2j}^{-1},
                       \qquad\text{indices modulo }q.
                                                        \tag{D.41f}
\]
Indeed $z=u_0y$ gives
$x=z^{-1}y^{-1}=y^{-1}u_0^{-1}y^{-1}
 =u_{q-1}^{-1}y^{-2}$.
Multiplying its $q$ successive translates by $y^{-2}$
gives exactly (D.41f), because $y^{2q}=1$.

For any prime $p$ put $V=H_1(K,\Q_p)$ and let $A$ be
the map induced by $\beta$.  From (D.41e) we obtain
the exact module isomorphism
\[
 V\cong
 \Q_p[C_q]\big/\Q_p(1+T+\cdots+T^{q-1})
 \cong\Q_p[X]/\Phi_q(X),\qquad
              A^2=1,\quad AT=TA.
                                                        \tag{D.41g}
\]
Every irreducible cyclotomic factor occurs with
multiplicity one.

The one-boundary surface structure in (D.41c) supplies
a nondegenerate alternating pairing, and its relation
to (D.41d) is algebraic.  In a surface free basis the
boundary word is the product of the $g$ commutators,
up to inversion and conjugation.  Its degree-two image
in the class-two free pro-$p$ quotient is a nondegenerate
bivector $\omega\in\bigwedge^2 V$.
Equation (D.41d) gives
$(\bigwedge^2 A)\omega=\lambda_p\omega$.
The deck map preserves its peripheral exponent, so
$(\bigwedge^2 T)\omega=\omega$.
If $\Omega$ is the nonsingular skew matrix of $\omega$,
these are precisely
\[
 A\Omega A^{\mathsf t}=\lambda_p\Omega,\qquad
 T\Omega T^{\mathsf t}=\Omega.
\]
Consequently the alternating form with matrix
$\Omega^{-1}$ satisfies
\[
 \langle Av,Aw\rangle=\lambda_p\langle v,w\rangle,
 \qquad
 \langle Tv,Tw\rangle=\langle v,w\rangle.
                                                        \tag{D.41h}
\]
This proves the multiplier comparison also at $p=2$:
one may first take the integral commutator calculation
in each finite class-two quotient and then tensor with
$\Q_2$.  It uses no geometric realization of $\beta$.

Now let $p\ne q$ be a prime with
$d=\operatorname{ord}_q(p)$ even.  The irreducible
factors of $\Phi_q$ over $\Q_p$ have root sets
the Frobenius orbits
\[
                  \{\zeta^{j p^k}:0\leq k<d\},
                     \qquad j\in(\Z/q)^\times.
                                                        \tag{D.41i}
\]
For completeness these roots lie in the unramified
extension with residue field $\F_{p^d}$: the primitive
$q$th roots in that field lift uniquely since $p\ne q$,
and Frobenius lifts their map $z\mapsto z^p$.
The least such degree is $d$, proving the factor
description.  Since $d$ is even, $p^{d/2}=-1$ modulo
$q$, and each of these root sets is closed under
inversion.

Write $V=\bigoplus V_i$ for those irreducible summands.
By (D.41g), the commutant of $T$ is the product of their
fields $\Q_p(\zeta^j)$.  Therefore $A|_{V_i}$ is
multiplication by an element of a field whose square
is one, and so is the scalar $1$ or $-1$ on that entire
summand.  Each $V_i$ is nondegenerate for the form in
(D.41h).  To see this explicitly after a splitting-field
extension, a $T$-eigenvector with eigenvalue $\zeta^j$
can pair only with the eigenvalue $\zeta^{-j}$, since
the pairing is $T$-invariant.  Both belong to $V_i$ by
(D.41i); different summands are orthogonal.  Global
nondegeneracy thus restricts to nondegeneracy on $V_i$.
Multiplication by either scalar $1$ or $-1$ preserves
that restricted form.  Comparing with (D.41h) gives
\[
 \lambda_p=1
       \quad\text{whenever }p\ne q
                    \text{ and }\operatorname{ord}_q(p)
                    \text{ is even}.
                                                        \tag{D.41j}
\]

There is a prime divisor $p$ of $q-1$ having this even
order.  Otherwise every prime factor of $q-1$ would
lie in the odd-order subgroup of the cyclic group
$(\Z/q)^\times$.  Their product, with multiplicities,
would also have odd order.  But that product is
$q-1=-1$ modulo $q$, of order two.  This contradiction
proves the assertion.  Such a prime satisfies $p<q$.
For $p=2$ the hypothesis gives $\lambda_p=-1$; for an
odd $p<q$ the minimal choice of $q$ gives the same
value.  Both contradict (D.41j).

Thus no positive odd-prime coordinate exists, and
$\lambda=-1$ globally.  D.34 handles a pure
involution, while D.36 handles an involution with
positive $2$-adic multiplier.  All involutions are
therefore accounted for.  Multiplicativity proves
the last assertion for every subgroup generated by
them.  This conclusion determines these actual
multiplier homomorphisms; it does not assert geometric
conjugacy of the finite actions or settle the
multiplier of an isolated order-three subgroup.

The program \texttt{code/part\_c/negative\_multiplier\_d41.py}
checks the exact Schreier boundary word, the induced
cyclotomic deck matrices and the nondegenerate
degree-two boundary form in examples, and checks
the prime-factor and Frobenius-orbit calculation.
The displayed argument proves the assertion for
every odd prime, rather than only for the executed
range.
\end{proof}

\begin{claim}{D.42}{Proved}{Order-three pants actions have global multiplier one}
Every element of order three in the full three-cusp
loop-outer group has orientation multiplier
$1\in\widehat\Z^\times$.
Consequently every finite subgroup of that full
group has multiplier image contained in the two
global units $\{1,-1\}$, and its multiplier is
determined by its $2$-adic coordinate.

The closed-surface calculation used in the proof is
also unconditional.  For $g\geq2$ and every prime $p$,
the homomorphism
\[
 \Aut(\widehat\pi_{g,0})
       \longrightarrow\GL(H_1(\widehat\pi_{g,0},\Z_p))
                                                        \tag{D.42a}
\]
is injective on every finite $p$-subgroup.
Here the source consists of actual continuous
automorphisms.  The order-three outer action in the
first assertion is supplied with an actual lift by
the exact rank-two lifting result D.27.
\end{claim}
\passage{\S{}30; folios 194--196; scans 383--387; \S{}31; folios 198--199 verso; scans 391--394; Part D: full loop-outer actions of type (0,3)}

\begin{proof}
We first prove (D.42a), keeping the two completion
maps and the fixed subgroup explicit.

\emph{The pro-$p$ automorphism calculation.}
Let $P=\pi_{g,0}^{(p)}$ be the pro-$p$ completion.
The closed lower central terms $\gamma_nP$
satisfy
\[
 \bigcap_n\gamma_nP=1,
 \qquad P/\gamma_nP\text{ is torsion-free for every }n.
                                                        \tag{D.42b}
\]
These are precisely
\cite[Propositions 2.14--2.15, pp.~1515--1516]{A-Morales2024},
applied to the orientable surface presentation.
The second proposition checks the primitive degree-two
initial form of its single defining relation; the
first passes the torsion-free nilpotent quotients to
the pro-$p$ completion by the stated canonical quotient
maps.  In particular every
$\gamma_nP/\gamma_{n+1}P$, as a subgroup of
$P/\gamma_{n+1}P$, is torsion-free.

Put
$K_n=\ker(\Aut(P)\to\Aut(P/\gamma_{n+1}P))$ for $n\geq1$.
For $b\in K_n$, the formula
\[
 d_n(b)(\bar z)=b(z)z^{-1}\pmod{\gamma_{n+2}P}
                                                        \tag{D.42c}
\]
defines a continuous homomorphism from $P^{\rm ab}$
to $\gamma_{n+1}P/\gamma_{n+2}P$.
Indeed this last quotient is central modulo
$\gamma_{n+2}P$, which makes the formula additive
in $z$ and trivial on commutators.
An automorphism acting trivially on $P^{\rm ab}$
acts trivially on every lower central graded quotient:
these quotients are generated by the iterated
commutators of its abelianization generators.
Consequently $d_n$ is a group homomorphism on $K_n$,
and its kernel is exactly $K_{n+1}$.
Its target is torsion-free by (D.42b).
A finite-order element of $K_1$ therefore lies
successively in every $K_n$ and is the identity.
Thus the kernel of
$\Aut(P)\to\GL(H_1(P,\Z_p))$ is torsion-free.

\emph{The full fixed subgroup.}
Write $\Pi=\widehat\pi_{g,0}$ and let
$a\in\Aut(\Pi)$ satisfy $a^p=1$ and act trivially
on $H_1(\Pi,\Z_p)$.
Its induced automorphism on $P=\Pi^{(p)}$ is the
identity by the preceding calculation.
There is therefore a specified quotient homomorphism
\[
 E=\Pi\rtimes_a C_p\longrightarrow P\times C_p,
             \qquad(h,s^i)\longmapsto(\bar h,s^i).
                                                        \tag{D.42d}
\]
The natural map $\Pi\to P$ induces an isomorphism
on mod-$p$ cohomology, including products.
Both groups have cohomological $p$-dimension two,
first cohomology $\F_p^{2g}$, and second cohomology
$\F_p$ with the standard nondegenerate surface
cup pairing.  The map is the identity on the
generator characters, so the pairing shows that
it is also an isomorphism in degree two; higher
degrees vanish.  These are the surface duality
comparisons of A.2 and A.7, with the full and
pro-$p$ formations respectively.
Hochschild--Serre comparison for the named map
(D.42d) now gives a ring isomorphism
\[
 H^*(P\times C_p,\F_p)\xrightarrow{\sim}H^*(E,\F_p).
                                                        \tag{D.42e}
\]
In particular $\dim H^n(E,\F_p)=2g+2$ for $n>2$.

Choose a nonzero surface class
$\eta\in H^2(P,\F_p)$ and a nonzero periodic
class $u\in H^2(C_p,\F_p)$.
For every $m\geq1$, the class
$\eta u^m\in H^{2m+2}(E,\F_p)$ is nonzero
under (D.42e).
Every subgroup $S\leq E$ of order $p$ projects
isomorphically to $C_p$, because $\Pi$ is
torsion-free.  Its first coordinate in $P\times C_p$
is trivial because $P$ is torsion-free.
Its normalizer is exactly
\[
               N_E(S)=C_\Pi(S)\times S.                 \tag{D.42f}
\]
An element normalizing $S$ fixes each of its elements
by projection to the abelian quotient $C_p$;
splitting off the uniquely determined $S$ component
proves (D.42f).

Apply the full profinite Brown restriction theorem
\cite[Theorem 2 and \S1.1]{A-KochloukovaZalesskii2007}
with virtual cohomological dimension two.
Its hypotheses hold: $\Pi$ is an open subgroup
of cohomological $p$-dimension two, and there is
no $C_p\times C_p$ in $E$, since every finite
subgroup injects into $E/\Pi=C_p$.
The finite-dimensional source makes its dense
restriction embedding an isomorphism onto a
finite product of the normalizer cohomologies.
Restriction of the nonzero class $\eta u^m$
is nonzero in at least one factor.
For that factor put $C=C_\Pi(S)$.
Because the first coordinate of $S$ in (D.42d)
is trivial, its restriction is exactly
\[
        (\eta|_C)\otimes(u^m|_S)
       \in H^2(C,\F_p)\otimes H^{2m}(S,\F_p).
                                                        \tag{D.42g}
\]
Thus $\eta|_C\ne0$.
The restriction
$H^1(\Pi,\F_p)\to H^1(C,\F_p)$ is injective:
for any nonzero $v$, nondegeneracy supplies $w$
with $v\smile w$ a nonzero multiple of $\eta$,
and its restriction cannot vanish.
Hence $\dim H^1(C,\F_p)\geq2g$ and
$\dim H^2(C,\F_p)\geq1$.
For every normalizer (D.42f), the K\"unneth formula
in degree $n>2$ contributes
\[
  1+\dim H^1(C_\Pi(S),\F_p)
                +\dim H^2(C_\Pi(S),\F_p).              \tag{D.42h}
\]
This holds also for odd $p$: the exterior generator
of $H^*(C_p,\F_p)$ and its degree-two polynomial
generator give one dimension in each nonnegative
degree.  At $p=2$ the degree-one polynomial
generator has the same dimension consequence.
The one factor in (D.42g) already accounts for
all $2g+2$ dimensions.  There is exactly one
conjugacy class of sections, and for its centralizer
\[
 \dim H^1(C,\F_p)=2g,\qquad H^2(C,\F_p)=\F_p.           \tag{D.42i}
\]
Conjugacy under $E$ equals conjugacy under $\Pi$,
since $E=\Pi S$ for each section $S$ and the
$S$ factor centralizes that section.
Consequently these conclusions apply to
$C=C_\Pi(a)$ for the specified section in (D.42d).

We now prove directly that this closed subgroup
equals $\Pi$.
Continuous cohomology has the restriction comparison
\[
 H^j(C,\F_p)=
  \mathop{\rm colim}_{C\leq U\leq_o\Pi}H^j(U,\F_p).
                                                        \tag{D.42j}
\]
For an explicit cochain comparison, choose an
open normal subgroup $V$ of $\Pi$ so small that
the given continuous finite-valued cochain on $C$
factors through $C/(C\cap V)$ in every variable.
The isomorphism $CV/V=C/(C\cap V)$ inflates it to
the open subgroup $CV$.  Cocycle equalities descend
to that finite quotient; for a coboundary choose
a common $V$ also for its primitive cochain.
This proves (D.42j) with the displayed restriction maps.
Each $U$ is the completion of its finite covering
surface, of genus
$g_U=1+[\Pi:U](g-1)$.
For $V\leq U$ the restriction on $H^2$ multiplies
the fundamental class by $[U:V]$.
A nonzero class surviving in (D.42j) in degree two
therefore has a representative at some $U$ for
which every smaller open $V\geq C$ has
$p\nmid[U:V]$.
For these subgroups the restriction on $H^1$
is injective: corestriction composed with restriction
is multiplication by this prime-to-$p$ index.
If $C$ had infinite index, the indices
$[\Pi:V]$ would be unbounded, making
\[
 \dim H^1(V,\F_p)=2+2[\Pi:V](g-1)
\]
unbounded inside the fixed finite-dimensional
$H^1(C,\F_p)$, contrary to (D.42i).
Thus $C$ is open.  The same genus formula and
(D.42i) give $[\Pi:C]=1$.
It follows that $a=1$.
The kernel of (D.42a) on a finite $p$-subgroup
cannot contain an element of order $p$, proving
its asserted injectivity.

\emph{The order-three cyclic cover.}
Let $g$ have order three in the full pants group,
with multiplier $\lambda$.
By D.32 its cusp permutation is a three-cycle.
Label $x,y,z=(xy)^{-1}$ so that it sends their
classes successively to $y,z,x$, with the common
power $\lambda$.
D.27 gives an actual order-three representative
$\alpha\in\Aut(\widehat F(x,y))$.
The equality $\lambda^3=1$ implies $\lambda_3=1$:
an order-three unit of $\Z_3$ reduces to one,
and for $v\in\Z_3^\times$, $r\geq1$, the
valuation of $(1+3^rv)^3-1$ is $r+1$.

Suppose a coordinate $\lambda_q$ were nonidentity.
Then $q\ne3$, and its reduction
$\zeta\in\F_q^\times$ has exact order three;
in particular $q\equiv1\pmod3$ and $q\geq7$.
Define the surjective character
\[
 \chi:F\longrightarrow C_q=\F_q,\qquad
 \chi(x)=\zeta,\quad\chi(y)=1,\quad\chi(z)=\zeta^2.
                                                        \tag{D.42k}
\]
The relation is respected because
$1+\zeta+\zeta^2=0$.
The three scalar identities
$\zeta\cdot1=\zeta$,
$\zeta\cdot\zeta^2=1$, and
$\zeta\cdot\zeta=\zeta^2$
show exactly that $\chi\alpha=\chi$;
all inner conjugators disappear under $\chi$.
Let $N=\ker\chi$.
The specified degree-$q$ cover has one cusp
above each of the original three, since all
three values in (D.42k) are nonzero.
Its type is $(h,3)$ with
$h=(q-1)/2\geq3$, by
$2-2h-3=-q$.
Fill these three cusps and put
\[
 \Pi_h=N/\overline{\langle\!\langle x^q,y^q,z^q
                         \rangle\!\rangle}_N
          =\widehat\pi_{h,0}.                            \tag{D.42l}
\]
Completion of the finite-cover presentation
gives this exact quotient map.
The peripheral powers under $\alpha$ preserve
the indicated closed normal subgroup, so it
induces an actual automorphism $b$ of $\Pi_h$
with $b^3=1$.
Conjugation by $y$ gives the deck automorphism
$T$ of this quotient; $T^q=1$, because $y^q$
has been killed.
The relation $\chi\alpha=\chi$ says that
$bT$ and $Tb$ differ by an inner automorphism.
Their actions on homology therefore commute.

Here is an explicit calculation of that deck
module.  First fill only $y^q,z^q$ upstairs.
The group downstairs is
$D=C_q(y)\amalg C_q(z)$, with its map to
$C_q$ given by $y\mapsto1$, $z\mapsto r=\zeta^2$.
Its kernel $K$ has Schreier generators
\[
 v_i=y^iz y^{-i-r}\quad(i\in\Z/q),\qquad
           v_0v_rv_{2r}\cdots v_{(q-1)r}=1.              \tag{D.42m}
\]
The displayed relation contains every generator
once and eliminates one; these are the completed
discrete Schreier presentation and maps.
Conjugation by $y$ sends $v_i$ to $v_{i+1}$.
The remaining cusp $x^q=(z^{-1}y^{-1})^q$
rewrites as
\[
             \prod_{j=0}^{q-1}
                    v_{-r-j(r+1)}^{-1}.                \tag{D.42n}
\]
Since $r+1\ne0$, it also uses every index once.
Its abelianization is the negative of the
already imposed sum relation.  Filling it changes
no first homology.  It follows, for the actual
map $K\to\Pi_h$, that
\[
 H_1(\Pi_h,\Q_3)
   =\Q_3[C_q]/\Q_3(1+T+\cdots+T^{q-1})
   \simeq\Q_3[X]/\Phi_q(X).                             \tag{D.42o}
\]
This is the regular augmentation module, with
each irreducible factor occurring once.

Because $q\ne3$, reduction of $\Phi_q$ modulo
three is square-free.  Hensel lifting its
irreducible factors expresses the last algebra
in (D.42o) as a product of finite unramified
extensions of $\Q_3$.
None contains a nonidentity cube root of one:
reduction to the residue field, whose multiplicative
group has order $3^d-1$, first gives residue one;
the same valuation calculation used for $\Z_3$
then applies, since three remains a uniformizer.
The centralizer of $T$ on its cyclic module
(D.42o) is exactly that product algebra.
Thus the commuting operator $b_*$, with
$b_*^3=1$, equals one in every factor.
The torsion-free lattice
$H_1(\Pi_h,\Z_3)$ embeds in (D.42o), so $b_*$
is also the identity on that lattice.
The first part of this proof, applied with $p=3$,
now gives $b=1$ on the full $\Pi_h$.

Finally the quotient filling comparison in A.2
identifies the relative orientation line of
$(N;\langle x^q\rangle,\langle y^q\rangle,
\langle z^q\rangle)$ with the closed orientation
line of $\Pi_h$.
The three cusp maps have the common power
$\lambda$, so $b$ acts on the latter line by
$\lambda$ as well.
Its being the identity forces
$\lambda_q=1$, contradicting the chosen coordinate.
There is no such $q$, proving $\lambda=1$.

For the finite-subgroup assertion, D.32 embeds
every finite subgroup in $S_3\times C_2$.
Its elements therefore have orders $1,2,3$, or $6$.
D.34, D.36, and D.41 give the global sign for
every involution; the preceding proof handles
order three.  An element $v$ of order six is
the product of its commuting powers $v^3$ and
$v^4$, of orders two and three respectively.
Multiplicativity gives the assertion for every
element, with its prescribed $2$-adic sign.

The program \path{code/part_a/pants_order_three.py}
executes the character invariance, both Schreier
index cycles, the cyclic deck matrix and its
cyclotomic polynomial, and the characteristic-three
factorization checks for both primitive cube roots
at every prime $q\equiv1\pmod3$ through $97$.
The arguments above prove the result for every
prime and for the full profinite groups.
\end{proof}

\begin{claim}{D.43}{Proved}{Global orientation of punctured-torus involutions}
Every involution in the full once-punctured-torus loop-outer
group has the global orientation multiplier $1$ or $-1$,
with the same sign as its $2$-adic determinant.
Consequently the orientation character of every finite
subgroup generated by involutions has values in the global
two-element subgroup $\{1,-1\}\subset\widehat\Z^\times$.
\end{claim}
\passage{Part D(b) of the task; \S{}24, folio 147 verso, scan 293; \S{}30, folios 194--196, scans 383--387; \S{}31, folios 198--199 verso, scans 391--394.}

\begin{proof}
Write $F=\widehat F(x,y)$ and $c=[x,y]$ for the torus
cusp.  D.27 supplies an actual order-two representative
$\alpha$ of the given outer involution.  Its peripheral
multiplier is
\[
 \lambda=\det(\alpha|_{F^{\rm ab}})\in\widehat\Z^\times,
                    \qquad \lambda^2=1,
                                                        \tag{D.43a}
\]
by the exact Heisenberg comparison of D.33.
Thus every $\lambda_p$ is $1$ or $-1$.

\emph{The characteristic Heisenberg cover.}
For an odd prime $q$, take the characteristic quotient
\[
 F\longrightarrow G_q,\qquad
 G_q=\langle X,Y,Z\mid X^q=Y^q=Z^q=1,\
              [X,Y]=Z,\ [X,Z]=[Y,Z]=1\rangle,
                                                        \tag{D.43b}
\]
with $x\mapsto X$, $y\mapsto Y$, $c\mapsto Z$.
Collection of these words identifies $G_q$ with the
order-$q^3$ Heisenberg group.  Explicitly it is
$\F_q^2\times\F_q$, with multiplication
\[
 (v,z)(w,t)=(v+w,z+t+\tfrac12\omega(v,w)),
 \qquad \omega((a,b),(a',b'))=ab'-ba'.
                                                        \tag{D.43c}
\]
Take $X=((1,0),0)$, $Y=((0,1),0)$ and $Z=((0,0),1)$.
The kernel $N_q$ is the closed subgroup generated by
the laws $\gamma_3F$ and all $q$th powers, so it is
characteristic.  More precisely it is their closed
normal closure: the exponent-$q$, class-two quotient
has the normal forms $X^aY^bZ^d$ in (D.43b), and
(D.43c) verifies that all $q^3$ forms are distinct.
This proves the characteristic-quotient identification.

The corresponding discrete cover of a once-punctured
torus has $q^2$ cusps, one for each coset of
$\langle Z\rangle$ in $G_q$, each of degree $q$.
Fill all these cusps.  Completing the finite-cover
presentation gives the exact quotient
\[
 N_q\longrightarrow
 \Pi_q=N_q/
   \overline{\langle\!\langle h c^q h^{-1}:
                       h\in F\rangle\!\rangle}_{N_q}
       \cong\widehat\pi_{g_q,0},
 \qquad 2g_q=2+q^3-q^2 .
                                                        \tag{D.43d}
\]
Representatives for the $q^2$ cosets suffice in the
normal closure.  The genus formula follows by multiplying
the punctured torus's Euler characteristic by $q^3$
and then adding its $q^2$ filled discs.
The original deck group induces its actual homology
representation on this closed surface.
Since $\alpha$ preserves $N_q$ and permutes all the
filled peripheral classes, it induces an actual
involution $\beta_q$ of $\Pi_q$.

For every prime $p$, put $V_p=H_1(\Pi_q,\Q_p)$,
let $\rho$ be its $G_q$ representation, and let
$A$ be the map induced by $\beta_q$.  If $\theta$
is the automorphism induced by $\alpha$ on $G_q$,
the precise deck comparison is
\[
 A^2=1,\qquad
 A\rho(h)A^{-1}=\rho(\theta(h)),\qquad
                    \theta(Z)=Z^{\lambda\bmod q}.
                                                        \tag{D.43e}
\]
Indeed representatives in $F$ for the two sides differ
only by an element of $N_q$, whose conjugation is inner
on $\Pi_q$ and hence trivial on its first homology.

The nondegenerate alternating surface pairing satisfies
\[
 \langle Av,Aw\rangle=\lambda_p\langle v,w\rangle,
       \qquad
 \langle\rho(h)v,\rho(h)w\rangle=\langle v,w\rangle .
                                                        \tag{D.43f}
\]
Here is an exact verification of the first multiplier.
The involution permutes the odd number $q^2$ of lifted
cusps, so one cusp class is fixed.  Fill every other
cusp first.  This is an invariant closed normal
quotient, and gives a surface group of genus $g_q$
with one cusp, which $\alpha$ preserves with power
$\lambda$.  Filling its last cusp does not change
first homology.  The degree-two boundary bivector
calculation in (D.41h) therefore proves the first
identity in (D.43f) for precisely $A$ on $V_p$.
The second identity comes from the orientation-preserving
geometric deck maps of the specified finite surface cover.
Thus no realization of $\alpha$ as a homeomorphism
has been used.

We compute the whole deck module by finite cover chains.
The cover before filling retracts onto the regular
cover of a two-loop rose, so its first homology has
representation class $1+\Q_p[G_q]$.
The relative second homology of the attached discs
is $\Q_p[G_q/\langle Z\rangle]$.
The exact homology sequence of this filling, including
the invariant orientation class in degree two, gives
\[
 [V_p]=
       2[1]+[\Q_p[G_q]]
                 -[\Q_p[G_q/\langle Z\rangle]].
                                                        \tag{D.43g}
\]
All these maps are deck-equivariant.  They compute
the homology of the completed surface groups through
the finite-index and surface-homology comparisons
already used in A.2 and D.42.

Over a characteristic-zero field containing a primitive
$q$th root $\zeta$, let $V_j$ be the subspace where $Z$
acts as $\zeta^j$, for $j\in\F_q^\times$.
The last permutation module in (D.43g) is trivial on
$Z$, and the regular representation has each of its
$q$ central eigenvalues with multiplicity $q^2$.
Semisimplicity therefore gives
\[
                         \dim V_j=q^2 .
                                                        \tag{D.43h}
\]
The $G_q$-invariant alternating form pairs $V_j$
only with $V_{-j}$, and gives a nondegenerate
pairing between them.

\emph{One negative odd-prime sign determines every sign.}
Suppose $\lambda_q=-1$ at an odd prime $q$.
For any prime $p$, extend scalars to $\Q_p(\zeta)$.
By (D.43e), $A$ exchanges $V_j$ and $V_{-j}$.
Choose bases of these two spaces and write $P$ for
their nonsingular pairing matrix, $B$ for
$A:V_j\to V_{-j}$ and $C$ for $A:V_{-j}\to V_j$.
Then $CB=BC=1$, and alternation in (D.43f) gives
\[
                  -B^{\mathsf t}P^{\mathsf t}C
                                      =\lambda_p P.
                                                        \tag{D.43i}
\]
The dimension $q^2$ is odd.  Taking determinants
in (D.43i) yields $-1=\lambda_p^{q^2}$, hence
$\lambda_p=-1$ by (D.43a).
This works also at $p=q$ and $p=2$, since the
calculation is over characteristic-zero fields.
Thus
\[
       \lambda_q=-1\text{ for an odd }q
                  \quad\Longrightarrow\quad\lambda=-1.
                                                        \tag{D.43j}
\]

\emph{The positive $3$-adic case.}
It remains to prove that $\lambda_3=1$ forces
$\lambda_2=1$.  The matrix of $\alpha$ on
$H_1(F,\Z_3)$ has square one and determinant one.
Its two eigenvalues over $\Q_3$ are consequently
both $1$ or both $-1$, and it is diagonalizable,
so the matrix is $I$ or $-I$.

In the group $G_3$ of (D.43c), the resulting $\theta$
has the form
\[
                \theta(v,z)=(Mv,z+\ell(v)),
                      \qquad M=I\text{ or }-I,
                                                        \tag{D.43k}
\]
with $\ell$ a linear functional.
For $M=I$, $\theta^2=1$ implies $2\ell=0$ and
$\ell=0$.
For $M=-I$, conjugation of $\theta$ by the inner
automorphism of $(u,0)$ replaces $\ell(v)$ by
$\ell(v)-2\omega(u,v)$.  Choose $u$ with
$\omega(u,v)=\ell(v)/2$.  Lifting $(u,0)$ to $F$
and conjugating the actual automorphism $\alpha$
by that inner automorphism preserves its order two
and its peripheral multiplier.  It gives exactly
\[
             \theta=1\quad\hbox{or}\quad
             \theta=\iota,\qquad\iota(v,z)=(-v,z).
                                                        \tag{D.43l}
\]
This is conjugation of an actual automorphism,
not multiplication by an arbitrary inner lift.

Work now over the unramified quadratic field
$L=\Q_2(\zeta_3)$, with
$\sigma(\zeta_3)=\zeta_3^{-1}$.
For $j=1,2$ the irreducible $G_3$ module $W_j$
with central character $\zeta_3^j$ has basis
$e_r$, $r\in\F_3$, and matrices
\[
 X e_r=e_{r+1},\qquad
 Y e_r=\zeta_3^{-jr}e_r,\qquad
 Z e_r=\zeta_3^j e_r .
                                                        \tag{D.43m}
\]
These matrices satisfy (D.43b).
The distinct eigenvalues of $Y$ and the cyclic
action of $X$ show irreducibility.  The diagonal
projectors for $Y$, followed by powers of $X$,
span the full matrix algebra.  Hence the central
block of $L[G_3]$ is $M_3(L)$.
Together with (D.43g)--(D.43h), this gives
\[
                       V_j=W_j\otimes_L U_j,
                             \qquad\dim U_j=3.
                                                        \tag{D.43n}
\]

The matrix $J=I$ implements $\theta=1$ on $W_j$.
For $\theta=\iota$, take $J e_r=e_{-r}$.
This rational permutation matrix implements
$X\mapsto X^{-1}$, $Y\mapsto Y^{-1}$, fixes $Z$,
and has $J^2=I$ and $J^{\mathsf t}J=I$.
Thus (D.43e) and the full matrix-algebra commutant
give
\[
                A|_{V_j}=J\otimes B_j,\qquad B_j^2=I.
                                                        \tag{D.43o}
\]
Choose a multiplicity basis for $U_1$ and its
$\sigma$-image for $U_2$.
Since the original representation and $A$ are
defined over $\Q_2$, we then have
$B_2=\sigma(B_1)$.
The canonical invariant pairing between $W_1$
and $W_2$ is the ordinary dot product in (D.43m).
Consequently the nondegenerate cross-pairing in
(D.43f) has the form $I_3\otimes S$ for an
invertible $3$-by-$3$ matrix $S$ on the multiplicity
spaces.  Equation (D.43f) becomes
\[
                 B_1^{\mathsf t}S\,\sigma(B_1)
                                      =\lambda_2 S.
                                                        \tag{D.43p}
\]
Because $B_1^2=I$, its determinant is $1$ or $-1$,
both fixed by $\sigma$.
Taking determinants in (D.43p) therefore gives
$1=\lambda_2^3$.  Together with $\lambda_2^2=1$,
this proves $\lambda_2=1$.

If $\lambda_2=-1$, the last conclusion forces
$\lambda_3=-1$, so (D.43j) gives $\lambda=-1$.
If $\lambda_2=1$, (D.43j) excludes every negative
odd-prime coordinate, and $\lambda=1$.
This proves the claim.  Multiplicativity proves
the stated consequence for subgroups generated
by involutions; no conclusion about their
conjugacy is inferred.

The program \texttt{code/part\_c/torus\_involution\_d43.py}
checks the characteristic Heisenberg presentations,
the finite-cover chain ranks and central-character
dimensions, and every $G_3$ involution with determinant
one in (D.43k), including the exact inner-conjugacy
formula in (D.43l).  It also verifies the
Schr\"odinger matrices and their rational implementer.
The field and determinant arguments above prove
the assertion for the full profinite groups.
\end{proof}

\begin{claim}{D.44}{Proved}{A first-symmetry choice for the involution coboundary}
Let $\alpha=(-1,f)\in GT_0$ be an involution in its
normalized actual action on $F=\widehat F(x,y)$:
\[
 \alpha(x)=x^{-1},\qquad
 \alpha(y)=f^{-1}y^{-1}f,\qquad f\in F'.
\]
Write $\theta(x)=y$, $\theta(y)=x$.
There exists $h\in F'$ satisfying both exact equations
\[
               f=h^{-1}\alpha(h),\qquad h\theta(h)=1.
                                                        \tag{D.44a}
\]
Thus the coboundary in D.35 can be chosen with the
first Harbater--Schneps symmetry.  For every unit
$\mu\in\widehat\Z^\times$, the associated continuous
endomorphism
$\rho_{\mu,h}(x)=x^\mu$,
$\rho_{\mu,h}(y)=hy^\mu h^{-1}$ satisfies
$\alpha\rho_{\mu,h}=\rho_{\mu,h}\delta$, where
$\delta$ inverts $x,y$.
Its full peripheral-conjugator condition is precisely
D.35's remaining equation
$x^\mu hy^\mu h^{-1}\sim_F(xy)^\mu$.

At the characteristic finite quotient $A_5^{19}$
of D.38, this first-symmetry condition is soluble
for every one of the $5\,062\,500$ tuples that
pass the triangle test.  This finite equation
therefore gives no further exclusion among them.
\end{claim}
\passage{\S{}30; folios 194--196; scans 383--387; \S{}31; folios 198--199 verso; scans 391--394}

\begin{proof}
The normalized involution and the first
Harbater--Schneps equation give
\[
                  f\alpha(f)=1,\qquad\theta(f)f=1.
                                                        \tag{D.44b}
\]
Checking the two free generators gives the
identity of actual automorphisms
\[
                 \theta\alpha\theta
                        =\operatorname{Inn}(f)\alpha.
                                                        \tag{D.44c}
\]
Indeed both sides send $x$ to $fx^{-1}f^{-1}$
and $y$ to $y^{-1}$.

Form $E=F\rtimes_\alpha\langle s\mid s^2=1\rangle$.
Define an automorphism $\Theta$ of this exact
extension by
\[
                      \Theta|_F=\theta,\qquad
                      \Theta(s)=fs.                    \tag{D.44d}
\]
The first equation of (D.44b) preserves $s^2=1$;
(D.44c) preserves the conjugation relation;
and the second equation of (D.44b) gives
$\Theta^2(s)=\theta(f)fs=s$.
Thus $\Theta^2=1$ on all of $E$.
Adjoin its named involution $\tau$ and put
\[
 G=E\rtimes_\Theta\langle\tau\rangle,\qquad
                         G/F=C_2\times C_2.
                                                        \tag{D.44e}
\]
The two quotient coordinates are represented
by $s,\tau$.

The preimage of the $\tau$ line is the normal
subgroup
$G_\theta=F\rtimes_\theta\langle\tau\rangle$.
It has the explicit free-product identification
\[
 G_\theta\xrightarrow{\sim}
                  C_2(\tau)\amalg\widehat\Z(v),
       \qquad x\mapsto v,\quad y\mapsto\tau v\tau.
                                                        \tag{D.44f}
\]
The inverse sends $v$ to $x$ and $\tau$ to
the named involution; both composites fix
the indicated generators.
The free-factor intersection theorem
\cite[Theorem 5 of the inspected author version]{A-HerfortZalesskii2008}
therefore gives $C_F(\tau)=1$.

There is exactly one $F$-conjugacy class of
involutions in $G_\theta$.
For a direct verification, its mod-two
$H^1(F,\F_2)$ module is the regular
two-dimensional $C_2$ module, since $\theta$
exchanges its two basis characters.
Its positive-degree $C_2$ cohomology vanishes.
The split Hochschild--Serre calculation
therefore gives
$\dim H^n(G_\theta,\F_2)=1$ for $n\geq3$.
Brown restriction
\cite[Theorem 2 and \S1.1]{A-KochloukovaZalesskii2007}
gives one involution class, by exactly the
dimension argument of (D.37f).
Conjugacy by $G_\theta$ equals conjugacy by
$F$, since $G_\theta=F\langle u\rangle$
for every involution $u$ and its second
factor centralizes that element.

Consider the closed normalizer
$V=N_G(\langle\tau\rangle)$.
For every $g\in G$, the conjugate $g\tau g^{-1}$
belongs to $G_\theta$, and its image in
$G/F$ is still the $\tau$ coordinate.
The preceding uniqueness gives an $F$-conjugator
taking it back to $\tau$.
It follows that $V\to G/F$ is onto.
Its kernel is $C_F(\tau)=1$, so
\[
                   V\xrightarrow{\sim}C_2\times C_2.
                                                        \tag{D.44g}
\]
In particular $V$ is a finite subgroup,
contains $\tau$, and has a unique involution
$u$ mapping to the $s$ coordinate.
That involution belongs to $E$ and commutes
with $\tau$.

The three $F$-conjugacy classes of involutions
in $E$ were determined in D.34.  They are
represented by
\[
                            s,\quad xs,\quad yfs.
                                                        \tag{D.44h}
\]
Their distinct mod-two abelianization vectors
also identify the classes.
Formula (D.44d) sends $s$ to $fs$, which
belongs to the class of $s$, and gives the
exact transposition
\[
                      \Theta(xs)=yfs,\qquad
                      \Theta(yfs)=xs.
                                                        \tag{D.44i}
\]
Thus the class of $s$ is the unique fixed
class under $\Theta$.
Since $u$ commutes with $\tau$, its class
is fixed, so $u=ksk^{-1}$ for some $k\in F$.
Conjugate the finite subgroup (D.44g)
by $k^{-1}$.
It now contains $s$ and
\[
                k^{-1}\tau k=h\tau,\qquad
                              h=k^{-1}\theta(k).
                                                        \tag{D.44j}
\]
The square of this second involution gives
$h\theta(h)=1$.
Its commutation with $s$ gives
$hf=\alpha(h)$, which is precisely the
first equation of (D.44a).
On abelianization, this equation reads
$-2h^{\rm ab}=f^{\rm ab}=0$.
Multiplication by two is injective on
$\widehat\Z^2$, hence $h\in F'$.
Finally substitution in $\alpha(hy^\mu h^{-1})$
using $\alpha(h)=hf$ proves
$\alpha\rho_{\mu,h}(y)=\rho_{\mu,h}(y^{-1})$; the
$x$ identity is immediate.

We check the finite scope explicitly.
Write $Q=A_5^{19}$, and use D.38's component
description of $\theta$ and its involution tuple
$t$.  The triangle test makes every component
$t_i$ a nonidentity involution of $A_5$.
The equations to solve in this quotient are
\[
 h^{-1}t h=\theta(t),\qquad h\theta(h)=1,
       \qquad f=\theta(t)t^{-1}.
                                                        \tag{D.44k}
\]
They separate over the one-element and
two-element orbits of $\theta$ on the factors.

For a two-element orbit, write
$\theta(v)_i=g v_jg^{-1}$ and
$\theta(v)_j=g^{-1}v_i g$, with $g\in S_5$.
All nonidentity involutions of $A_5$ are
conjugate within $A_5$.
Choose $h_i\in A_5$ satisfying
$h_i^{-1}t_i h_i=g t_jg^{-1}$, and set
$h_j=g^{-1}h_i^{-1}g$.
These substitutions give both equations
of (D.44k) on that orbit.

For a fixed factor, $\theta$ is conjugation
by an involution $g\in S_5$, allowing $g=1$.
Indeed its square is the identity and
$C_{S_5}(A_5)=1$.
Choose an element $v$ with $v^2=1$, commuting
with $t_i$ and having the same parity as $g$.
When $g$ is even choose $v=1$.
When $g$ is odd, take either transposition
in the disjoint-transposition expression
for $t_i$.
Then $h_i=vg^{-1}\in A_5$ satisfies both
equations in (D.44k).
Thus every triangle survivor has a finite
first-symmetry solution, including the
$2\,531\,250$ survivors failing D.38's
finite peripheral-conjugacy test.
No lifting of these finite solutions is
asserted.

The program
\path{code/part_a/pure_involution_first_symmetry.py}
executes the local construction for all
$390$ fixed-factor cases and all
$27\,000$ two-factor cases.
It checks both equations in (D.44k),
which cover every factor orbit in D.38.
\end{proof}

\begin{claim}{D.45}{Proved}{A characteristic binary-icosahedral refinement}
Let $F=\widehat F(x,y)$ and let $R$ be D.38's characteristic
kernel with $F/R=A_5^{19}$.  Intersecting the kernels of all
epimorphisms $F\twoheadrightarrow S=\mathrm{SL}_2(\F_5)$ gives
an open characteristic subgroup $\widetilde R\subset R$ and
the following exact quotient and reduction map:
\[
 \widetilde Q=F/\widetilde R
       \xrightarrow{\ \sim\ }S^{19}\times C_2^2,
 \qquad
 \pi:\widetilde Q\longrightarrow A_5^{19},\quad
       ((g_i),a)\longmapsto(g_i\bmod\{\pm I\})_i.
                                                        \tag{D.45a}
\]
For the $10\,125\,000$ finite tuples $t$ of D.38, impose at
this specified refinement the normalized multiplier $-1$
action, both Harbater--Schneps cycle equations, and the
exact norm $f\alpha(f)=1$.
Exactly $5\,062\,500$ base tuples admit such a refined
$f$-coordinate.  They are precisely D.38's cusp-triangle
survivors, and each has $64$ admissible $f$-coordinates.
Thus the total number of refined coordinate pairs is
$324\,000\,000$.
For a fixed $t$, these $64$ coordinates define the same
actual automorphism of $\widetilde Q$, of order two.
In particular this named characteristic refinement
detects exactly D.38's first excluded half.  All
$2\,531\,250$ triangle survivors outside D.38's finite
pure conjugacy set still pass this refinement.
\end{claim}
\passage{\S{}30; folios 194--196; scans 383--387; \S{}31; folios 198--199 verso; scans 391--394; Part D of the task: full GT\_0 torsion and involution conjugacy}

\begin{proof}
All matrix calculations below are in $\F_5$.
The program enumerates the $120$ determinant-one matrices
and verifies that their commutator subgroup has order $120$.
Their center is $\{I,-I\}$, and their projective quotient is
identified explicitly with D.38's $A_5$ by a generating pair
of orders two and three with product of order five.
The resulting bijection is checked against all $60^2$
products.  Conjugation by the $120$ projective invertible
matrices gives all automorphisms of $S$.
Indeed $\Aut(S)\to\Aut(A_5)$ is injective: an element of
its kernel has the form $g\mapsto c(g)g$, where
$c:S\to\{\pm I\}$ is a homomorphism, and perfection makes
$c=1$.  D.38 gives $|\Aut(A_5)|=120$, and the enumerated
projective conjugations already give $120$ distinct
automorphisms of $S$.

Choose lifts $X_i,Y_i\in S$ of the nineteen generating
pairs in D.38, and put $Z_i=(X_iY_i)^{-1}$.
Every pair of lifts generates $S$: its image generates
$A_5$, so a proper image would have index two in $S$,
contrary to perfection.  The program also checks all
$4\cdot19=76$ choices directly.
Two maps to $A_5$ with a given kernel differ by an
automorphism, and every automorphism lifts uniquely to
$S$.  Over each of the nineteen $A_5$ kernels there
are exactly four $S$ kernels, from the independent
sign choices for the lifts of $x$ and $y$.
If two such lifts had the same kernel, the induced
automorphism of $S$ would act trivially on $A_5$;
the proved injectivity would identify the lifts.
Consequently these are exactly the $76$ epimorphism
kernels defining $\widetilde R$.

The joint map
\[
 x\longmapsto((X_i),e_x),\qquad
 y\longmapsto((Y_i),e_y),
        \qquad e_x,e_y\text{ the basis of }C_2^2,
                                                        \tag{D.45b}
\]
is onto.  To see this without a rank inference, its
projection $D$ to $S^{19}$ maps onto $A_5^{19}$ by
D.38.  Thus $D Z(S^{19})=S^{19}$; centrality gives
$[S^{19},S^{19}]=[D,D]\leq D$, so perfection gives
$D=S^{19}$.  The commutator subgroup of the image
of (D.45b) therefore contains $S^{19}\times\{0\}$,
and its projection to $C_2^2$ is onto.  This proves
surjectivity of (D.45b).
Every other sign lift of the $i$th pair is obtained by
composing (D.45b) with
\[
 ((g_j),a)\longmapsto g_i(-I)^{\ell(a)},
                         \qquad \ell\in\Hom(C_2^2,C_2).
\]
The intersection of these four kernels, for all $i$,
is trivial: the zero character first makes each
$g_i=I$, and all characters then make $a=0$.
This proves (D.45a), including the specified kernel.
It is characteristic because every automorphism of
$F$ permutes its epimorphisms to $S$.

Here are the complete finite equations used in the
refinement.  The named word substitutions $\theta,
\omega$ of D.38 descend to $\widetilde Q$.
On its characteristic perfect subgroup $S^{19}$ they
permute the factors by D.38's recorded maps and apply
the unique lifted $S$ automorphisms just described.
Write $\widetilde t_i\in S$ for any lift of $t_i$ and
put
\[
 f_{0,i}=\theta(\widetilde t)_i\widetilde t_i^{-1},
 \qquad
 f_i=(-I)^{b_i}f_{0,i},\qquad b_i\in\F_2.
                                                        \tag{D.45c}
\]
These are all the $2^{19}$ lifts of
$f_Q=\theta(t)t^{-1}$ into
$\widetilde Q'=S^{19}$.
For every such lift, set
\[
 \begin{split}
 c_{x,i}&=X_i^{-1}
       (\widetilde t_i X_i\widetilde t_i^{-1})^{-1},\\
 c_{y,i}&=(f_i^{-1}Y_i^{-1}f_i)
       (\widetilde t_i Y_i\widetilde t_i^{-1})^{-1}.
 \end{split}                                             \tag{D.45d}
\]
D.38's normalization implies $c_{x,i},c_{y,i}\in
\{I,-I\}$.  The actual automorphism induced by the
normalized pair is exactly
\[
 \alpha((g_i),(a_x,a_y))=
 \bigl((\widetilde t_i g_i\widetilde t_i^{-1}
           c_{x,i}^{a_x}c_{y,i}^{a_y})_i,
                      (a_x,a_y)\bigr).
                                                        \tag{D.45e}
\]
This formula is a continuous automorphism; its
restriction to the first factor is an automorphism,
and its two correction characters are central.
It sends the two generators (D.45b) to the required
$X^{-1}$ and $f^{-1}Y^{-1}f$.
Since $\widetilde t_i^2$ is central and both correction
characters have exponent two, its square is the
identity.  Its reduction is the nonidentity order-two
automorphism of D.38, so its order is exactly two.
Multiplying $f_i$ by a central sign changes neither
(D.45d) nor (D.45e).  Conversely these images of the
two generators determine the automorphism uniquely.

The program tests the following three exact equations
in $S^{19}$ for all choices (D.45c):
\[
 \theta(f)f=1,\qquad
 \omega^2(f)Z^{-1}\omega(f)Y^{-1}fX^{-1}=1,
 \qquad f\alpha(f)=1.
                                                        \tag{D.45f}
\]
On $f$, which has zero $C_2^2$ coordinate, formula
(D.45e) is conjugation by $\widetilde t$.
Thus every operation in (D.45f) is a specified
matrix multiplication or a specified factor map;
the norm is tested as an equation, not inferred
from the order of the finite automorphism.

The factor maps for $\theta,\omega$ have six joint
orbits.  Exhaustion within these orbits gives the
following complete counts, with D.38's zero-based
factor indices:
\[
\begin{array}{c|r|l}
 \text{factor orbit}&\text{base tuples}&
       \text{numbers of admissible central lifts}\\ \hline
 \{0,2,10\}&15&2\text{ for each of }15\\
 \{1,9,17\}&15&2\text{ for each of }15\\
 \{3,5,6,7,11,14\}&60&0\text{ for }30;\quad4\text{ for }30\\
 \{4,12,18\}&15&2\text{ for each of }15\\
 \{8,13,15\}&10&2\text{ for each of }10\\
 \{16\}&5&1\text{ for each of }5.
\end{array}                                               \tag{D.45g}
\]
The rejected thirty assignments in the six-factor
orbit are exactly its triangle failures; the other
five orbits have no triangle failure.
The equations on one joint orbit involve no other
orbit, so their solution sets have the displayed
Cartesian-product decomposition.
The surviving count is
$15\cdot15\cdot30\cdot15\cdot10\cdot5=5\,062\,500$,
and the number of central lifts per survivor is
$2\cdot2\cdot4\cdot2\cdot2\cdot1=64$.
This proves every asserted count.
The executed implementation is
\path{code/part_a/a5_schur_refinement.py}; its output
records each orbit, the lift histogram, and sample
coordinates.  These computations impose a necessary
finite refinement on full involutions.  They assert
no full lift of a surviving finite coordinate pair.
\end{proof}

\begin{claim}{D.46}{Proved}{Three positive involution sections and order-four lifting}
Let $T$ be the full once-punctured-torus loop-outer group
of $F=\widehat F(x,y)$.
For every involution $b\in T$ whose $2$-adic orientation
is positive, its centre-free extension
\[
             1\longrightarrow F\longrightarrow E_b
                  \longrightarrow\langle b\rangle
                  \longrightarrow1
\]
has exactly three $F$-conjugacy classes of order-two
sections.  For every such section $S$,
\[
                         H^1(C_F(S),\F_2)=0.             \tag{D.46a}
\]
Every cyclic order-four subgroup of $T$ has a homomorphic
lift to $\operatorname{Aut}(F)$, with its prescribed
outer action.  Thus every element of order four in
$T$ has an actual order-four representative.
\end{claim}
\passage{Part D(b) of the task; \S{}24, folio 147 verso, scan 293; \S{}30, folios 194--196, scans 383--387; \S{}31, folios 198--199 verso, scans 391--394.}

\begin{proof}
\emph{The homology of a positive involution.}
D.27 gives an actual order-two representative $\beta$
of $b$, and identifies $E_b$ with
$F\rtimes_\beta\langle s\rangle$.
The map of D.33 on the finite $2$-group $\langle b\rangle$
is faithful.  The matrix of $\beta$ on
$H_1(F,\Z_2)$ is therefore nonidentity.
Its square is the identity and its determinant is $1$.
Over $\Q_2$ it is diagonalizable; the determinant
condition makes its two eigenvalues equal.
It follows that
\[
                    \beta^{\rm ab}_2=-I.               \tag{D.46b}
\]
This uses the punctured-torus statement D.33 to exclude
the identity matrix; no analogous exclusion for an
arbitrary unmarked rank-two action is assumed.

\emph{The full cohomology algebra and its Bockstein.}
Let $u:E_b\to\F_2$ be the quotient character, with
$u(s)=1$.  Choose a basis $\chi_1,\chi_2$ of
$H^1(F,\F_2)$ and extend these characters to $E_b$
by declaring $\chi_i(s)=0$.
These extensions exist because (D.46b) is the identity
modulo $2$.
The Hochschild--Serre spectral sequence has precisely
the two rows
\[
 E_2^{a,0}=\F_2 u^a,\qquad
 E_2^{a,1}=H^1(F,\F_2)\,u^a.                            \tag{D.46c}
\]
The transgression on $E_2^{0,1}$ is zero since the two
characters have just been extended to $E_b$.
Its other possible differentials are the products of
this transgression with powers of $u$, and hence are
zero as well.  There are no further nonzero rows.
Consequently, for every $n\geq1$,
\[
 H^n(E_b,\F_2)
   =\langle u^n,u^{n-1}\chi_1,u^{n-1}\chi_2\rangle_{\F_2},
 \qquad \dim H^n(E_b,\F_2)=3.                           \tag{D.46d}
\]
Independence follows from the two associated graded
pieces in (D.46c).
Multiplication by $u$ is an isomorphism in every
degree at least one: it is an isomorphism on those
associated graded pieces.

For each $i$, lift $\chi_i|_F$ to a continuous
homomorphism $\varphi_i:F\to\Z/4$.
Equation (D.46b) gives
$\varphi_i(\beta(h))=-\varphi_i(h)$.
The actual continuous one-cochain on $E_b$
\[
             \widetilde\chi_i(hs^\epsilon)=\varphi_i(h)
                      \quad(\epsilon=0,1)
\]
therefore satisfies
\[
 \begin{split}
 \delta\widetilde\chi_i(hs^\epsilon,ks^\eta)
  &=\varphi_i(h)+\varphi_i(k)
             -\varphi_i(h\beta^\epsilon(k))\\
  &=2\epsilon\,\varphi_i(k)
       \quad\text{in }\Z/4.
 \end{split}                                           \tag{D.46e}
\]
For the coefficient sequence
$0\to\F_2\xrightarrow{\;2\;}\Z/4\to\F_2\to0$,
this says that the Bockstein of $\chi_i$ is
$u\chi_i$.
It is also $\chi_i^2$.
For completeness, the lift of any $\F_2$ character
which takes its values in $\{0,1\}\subset\Z/4$
has coboundary, divided by $2$, equal to the cup
square of that character; comparison of two lifts
changes this cocycle only by a coboundary.
Thus
\[
                            \chi_i^2=u\chi_i.           \tag{D.46f}
\]

Invert the central class $u$ and take degree zero:
\[
             R=\bigl(H^*(E_b,\F_2)[u^{-1}]\bigr)^0.
\]
It has dimension three, with basis
$1,z_1=\chi_1/u,z_2=\chi_2/u$ by (D.46d).
Equation (D.46f) makes both $z_i$ idempotent.
In particular $R$ is a quotient of the algebra
\[
 \F_2[Z_1,Z_2]/(Z_1^2-Z_1,Z_2^2-Z_2)
           \cong\prod_{(a_1,a_2)\in\F_2^2}\F_2,
 \quad P\longmapsto(P(a_1,a_2))_{(a_1,a_2)}.             \tag{D.46g}
\]
The displayed evaluation map is an isomorphism:
the four products of $Z_i$ and $1-Z_i$ are its four
coordinate idempotents.
Every ideal of a finite product of fields is a product
of ideals in those fields.  Every quotient is therefore
reduced, so $R$ has no nonzero nilpotent element.

\emph{Brown restriction and the three full sections.}
Every finite subgroup of $E_b$ injects into
$E_b/F=C_2$, since $F$ is torsion-free.
Thus the nontrivial elementary abelian $2$-subgroups
are exactly its order-two sections, with no inclusions
between distinct ones.
For such a section $S$ put $C_S=C_F(S)$.
Its normalizer is the exact direct product
\[
                         N_{E_b}(S)=C_S\times S.        \tag{D.46h}
\]
Indeed every normalizer element centralizes $S$, since
$\operatorname{Aut}(C_2)=1$; multiplying it by its
unique image in $S$ leaves an element of $C_S$.

The full-profinite Brown restriction theorem used in
D.28 and D.33 applies because $F$ is open in $E_b$
and has $2$-cohomological dimension one
\cite[Theorem~2 and \S1.1]{A-KochloukovaZalesskii2007}.
In sufficiently high degrees it identifies cohomology
with the product of the cohomologies of (D.46h);
the identification is the actual restriction map and
preserves cup products.
Each factor contributes a copy of $\F_2$ in every
positive degree, so (D.46d) in particular makes the
number of factors finite.
Since $C_S$ is a closed subgroup of $F$, it has
$\operatorname{cd}_2(C_S)\leq1$.
After the same localization, restriction identifies
$R$ as a ring with
\[
       \prod_{[S]}
         \bigl(\F_2\oplus H^1(C_S,\F_2)\,u^{-1}\bigr).
                                                        \tag{D.46i}
\]
The second summand in each factor is an ideal with
square zero, since $H^2(C_S,\F_2)=0$.
Here $u|_{C_S}=0$ and $u|_S$ is the nonzero
character, so the localization in (D.46i) is exactly
the localization already used in (D.46g).
Reducedness of $R$ forces every second summand to
vanish.  Its dimension three then proves that there
are exactly three section classes, and proves
(D.46a).

These are also precisely the $F$-conjugacy classes.
To see this directly, let $e\in E_b$ and choose
$s'\in S$ with the same image as $e$ in $E_b/F$.
Then $eS e^{-1}=(es'^{-1})S(es'^{-1})^{-1}$,
and $es'^{-1}\in F$.

\emph{An invariant section in the order-four extension.}
Now let $a\in T$ have order four, and form its
centre-free extension
\[
 E=\{(\phi,j)\in\operatorname{Aut}(F)\times C_4:
                                      [\phi]=a^j\}.
                                                        \tag{D.46j}
\]
Its kernel is identified with $F$ by
$h\mapsto(\operatorname{Inn}(h),0)$.
Let $E_2$ be the inverse image of $\langle a^2\rangle$.
The orientation of $a^2$ at $2$ is positive:
the only roots of unity of order dividing four
in $\Z_2^\times$ are $1,-1$, whose squares are $1$.
Equivalently, D.33 embeds $\langle a\rangle$ in
$\operatorname{GL}_2(\Z_2)$, where its square is $-I$.
The first part of the proof thus applies to $E_2$.

Conjugation by $E$ permutes the three $F$-conjugacy
classes of sections of $E_2$.
Every element of $E_2$ acts trivially on this class
set, by the preceding paragraph, so the permutation
action factors through $E/E_2=C_2$.
An action of $C_2$ on a set of three elements has a
fixed element.
Choose a section $S$ in that class.
If $e\in E$ projects to the specified generator of
$C_4$, there is $f\in F$ with
$eSe^{-1}=fSf^{-1}$.
Hence $f^{-1}e\in N_E(S)$ still projects to that
generator.
Put $N=N_E(S)$ and $C=C_F(S)$.
We have proved the exact sequences
\[
 \begin{gathered}
 1\longrightarrow C\longrightarrow N
              \longrightarrow C_4\longrightarrow1,
       \qquad S\subseteq Z(N),\\
 1\longrightarrow C\longrightarrow J:=N/S
              \longrightarrow C_2\longrightarrow1.
 \end{gathered}                                         \tag{D.46k}
\]
The embedding of $C$ in the second line is injective
because $C\cap S=1$.

\emph{Projectivity supplies the order-four section.}
Equation (D.46a) implies
\[
                              C^{(2)}=1:                \tag{D.46l}
\]
every nontrivial pro-$2$ group has a nontrivial finite
$2$-group quotient, which in turn has a quotient
$C_2$.
In particular every homomorphism from $J$ to a finite
$2$-group kills $C$ and factors through the last
$C_2$ in (D.46k).  That quotient is nontrivial, so
\[
                              J^{(2)}=C_2.              \tag{D.46m}
\]

We show that $J$ has a nonidentity torsion element.
If it were torsion-free, its open subgroup $C$ has
$\operatorname{cd}_p(C)\leq1$ at every prime $p$,
as a closed subgroup of the free profinite group $F$.
Serre's torsion-free finite-index dimension theorem,
in exactly the form used in D.27, would give
$\operatorname{cd}_p(J)\leq1$ at every prime.
Thus $J$ would be projective as a full profinite
group.
Its maximal pro-$2$ quotient would then be a
projective, hence free, pro-$2$ group: pull a finite
$2$-group embedding problem back to $J$, solve it
by projectivity, and factor the solution through
$J^{(2)}$.
This contradicts (D.46m), since $C_2$ is not a free
pro-$2$ group.

Let $j\in J$ be a nonidentity torsion element.
Its cyclic subgroup injects into $J/C=C_2$, because
$C$ is torsion-free.
Consequently $j$ has order two and lies outside $C$.
Choose $n\in N$ over $j$.
Then $n^2\in S$, whereas its image in $C_4$ is an
odd element and has order four.
It follows that $n^2$ is the nonidentity element
of $S$, and $n$ has order exactly four.
After replacing $n$ by $n^{-1}$ if necessary, its
projection is the specified generator of $C_4$.
Thus $C_4\to E$, sending that generator to $n$,
is a section of (D.46j).
Projection to $\operatorname{Aut}(F)$ gives the
claimed homomorphic lift with its original labels.

The program \texttt{code/part\_c/torus\_order4\_lift\_d46.py}
checks the finite cochain formula (D.46e), its
comparison with the cup square, and the explicit
Boolean-algebra evaluation in (D.46g).
The dimension and normalizer arguments above prove
the statement for the full profinite extensions.
No claim of conjugacy to a geometric order-four
action is made.  D.43 applied to $a^2$ does also
give $\lambda(a)^2=1$ globally, so the remaining
orientation issue for $a$ concerns mixed signs,
rather than fourth roots of unity.
\end{proof}

\begin{claim}{D.47}{Proved}{Conjugator coordinates at the binary refinement}
Use the specified quotient
$\widetilde Q=S^{19}\times C_2^2$, $S=\mathrm{SL}_2(\F_5)$,
of D.45, and its $324\,000\,000$ refined pairs $(t,f)$.
Every one of these pairs satisfies D.44's finite
first-symmetry coboundary condition
\[
       h^{-1}\alpha(h)=f,\qquad h\theta(h)=1,
                         \qquad h\in S^{19}.             \tag{D.47a}
\]
Exactly $10\,125\,000$ of the pairs satisfy the two
finite conjugator equations of D.35 at this quotient.
The count stays the same after requiring both actual
invertibility of the finite substitution and both
Harbater--Schneps symmetries for the conjugating pair.
There are $5\,062\,500$ pairs with admissible unit
parameters congruent to a square modulo five, and
$5\,062\,500$ with parameters congruent to a nonsquare.
These two sets are disjoint.

Thus $313\,875\,000$ refined coordinate pairs lie
outside the image of the standard full $GT_0$
involution conjugacy class.  For each of D.38's
$2\,531\,250$ conjugator-compatible base tuples,
exactly four of its $64$ refined $f$-coordinates
pass this test; for each of the other
$2\,531\,250$ base tuples, none does.
These are statements about the exact normalized
coordinates: D.45's $64$ coordinates for a fixed
$t$ define the same finite automorphism.
No surviving pair is asserted to have a full
involution or conjugator lift.
\end{claim}
\passage{\S{}30; folios 194--196; scans 383--387; \S{}31; folios 198--199 verso; scans 391--394; Part D of the task: full GT\_0 involution conjugacy}

\begin{proof}
Write $X_i,Y_i,Z_i,\widetilde t_i$ and $\alpha$ for
the matrices and actual automorphism fixed in D.45.
Since $S$ has exponent dividing $60$, the power
parameter in D.35 factors through
$(\Z/60)^\times$.  For the second Harbater--Schneps
symmetry we retain the finer, explicit parameter map
\[
 a\in(\Z/120)^\times\longmapsto
       \left(a\bmod60,\ m=(a-1)/2\bmod60\right).
                                                        \tag{D.47b}
\]
It is well-defined, and every such $a$ is the
reduction of a unit in $\widehat\Z$.
Changing the representative of $a$ by $120$ changes
$m$ by $60$, so all subsequent equations are
independent of the choice of integer representative.

The image of $F'$ is exactly $S^{19}$ by D.45.
For fixed $(t,f)$ the first conjugator equation
therefore separates into the nineteen exact equations
\[
 h_i^{-1}\widetilde t_i h_i\widetilde t_i^{-1}=f_i,
                                  \qquad h_i\in S.       \tag{D.47c}
\]
Each has exactly four solutions.  The program
verifies this by enumeration of all $120$ matrices.
It also follows from the order-four conjugacy class
in $S$: $\widetilde t_i$ has square $-I$, and the
norm equation in D.45 gives
$(f_i\widetilde t_i)^2=\widetilde t_i^2=-I$.
The thirty matrices with square $-I$ form one
conjugacy class with centralizer of order four,
as the same matrix enumeration verifies.

For each parameter $a\bmod60$ and each of these
four matrices, the second equation of D.35 is
tested by the exact matrix conjugacy condition
\[
      X_i^a h_iY_i^a h_i^{-1}
                       \sim_S (X_iY_i)^a.               \tag{D.47d}
\]
Indeed the $C_2^2$ coordinates of both sides are
$e_x+e_y$, since $a$ is odd, and that direct factor
is central.  Conjugacy in $\widetilde Q$ is therefore
exactly these nineteen independent matrix conjugacies.
The program computes every conjugacy class in $S$
by conjugation by all $120$ matrices; (D.47d) is
tested against those complete classes.

The finite invertibility test is explicit as well.
Reduce the ordered pair
$(X_i^a,h_iY_i^a h_i^{-1})$ to $A_5$ and test whether
it belongs to the $i$th recorded $\Aut(A_5)$ orbit
of generating pairs in D.38.
This is sufficient: its unique component
automorphism lifts to an automorphism $\gamma_i$
of $S$, and the two central differences from
$\gamma_i(X_i),\gamma_i(Y_i)$ are corrected by
characters of $C_2^2$, exactly as in (D.45d)--(D.45e).
The resulting automorphism of $\widetilde Q$ has
the prescribed two generator images.
It is necessary for a finite pure automorphism
with these images: its reduction preserves the
three specified power-conjugacy classes, and
D.38's distinct passports force it to preserve
each simple factor.  Its pair must therefore be
in that recorded orbit.  This tests the actual
finite substitution, with its characteristic
quotient map retained.

For completeness the program also imposes both
symmetries on the conjugating pair, not just its
three peripheral images.
Put $k=h^{-1}$ and use $m$ from (D.47b).
The equations tested are
\[
       h\theta(h)=1,\qquad
       \omega^2(k)Z^m\omega(k)Y^m kX^m=1.
                                                        \tag{D.47e}
\]
All $h,k,f$ coordinates lie in the perfect factor
$S^{19}$; the actions of $\theta,\omega$ there
are the specified lifted factor maps of D.45.
The $C_2^2$ coordinate of the second equation is
$m(e_x+e_y+e_z)=0$.
Thus these are exact equations at the indicated
finite quotient, including the half-exponent
required by relation (II).

The first equation of (D.47e) is exhausted without
assuming independent choices where factors are
exchanged.  If
$\theta(v)_i=G_i v_jG_i^{-1}$, a choice of $h_i$
forces
\[
                       h_j=G_i^{-1}h_i^{-1}G_i.         \tag{D.47f}
\]
For a fixed factor it must satisfy the same equation
with $j=i$.  The program enumerates one of the
four candidates from (D.47c) for each orbit of
$\theta$, uses (D.47f), and checks membership in
the corresponding candidate sets.
It verifies that (D.47a) is soluble for all the
refined pairs.  For the full conjugator test it
then checks (D.47d), finite invertibility, and the
second equation of (D.47e) for every remaining
choice and all $32$ units modulo $120$.

Here is the entire factorization of the result.
Let $U_+$ and $U_-$ denote the subsets of
$(\Z/120)^\times$ whose residues modulo five are
respectively squares and nonsquares; each has
$16$ elements.
For each of D.45's six joint factor orbits, the
table gives the number of refined pairs having
the specified set of admissible parameters.
\[
\begin{array}{c|r|r|r|r}
 \text{factor orbit}&U_+&U_-&U_+\cup U_-&\varnothing\\ \hline
 \{0,2,10\}&15&15&0&0\\
 \{1,9,17\}&15&15&0&0\\
 \{3,5,6,7,11,14\}&30&30&0&60\\
 \{4,12,18\}&15&15&0&0\\
 \{8,13,15\}&10&10&0&0\\
 \{16\}&0&0&5&0.
\end{array}                                               \tag{D.47g}
\]
The admissible global parameter set is the
intersection of the six local sets, because
the parameter is common while the factor
equations have the exact orbit decomposition
of D.45.
It is nonempty precisely when the first five
nonempty choices all select the same one of
$U_+,U_-$.
Each choice gives
$15\cdot15\cdot30\cdot15\cdot10\cdot5
 =5\,062\,500$ pairs.
Their sum is $10\,125\,000$, and subtraction
from D.45's $324\,000\,000$ gives the stated
$313\,875\,000$ failures.
The program records the same table when only
(D.47c)--(D.47d) are imposed, and when finite
invertibility and the first symmetry are added.
Thus the asserted equality of these counts is
itself checked at every factor orbit.
It also records the counts for each base tuple
before multiplying over the six orbits.  Their
Cartesian product gives the exact histogram
\[
 \begin{array}{c|r}
 \text{admissible }f\text{-coordinates per base tuple}
          &\text{number of base tuples}\\ \hline
 0&2\,531\,250\\
 4&2\,531\,250.
 \end{array}                                             \tag{D.47h}
\]

Every full conjugator in D.35 reduces to solutions
of (D.47c)--(D.47d), and D.35's equality of the
$P$ and $GT_0$ conjugator sets also supplies
(D.47e).
Consequently failure of the displayed finite
test excludes that exact coordinate pair from
the image of the standard full conjugacy class.
The executed implementation is
\path{code/part_a/pure_involution_binary_conjugacy.py}.
Its output retains all unit sets, the six local
tables, the per-base histogram, and the separate
first-symmetry counts.  The proof uses only this
finite refinement and makes no assertion that
one of its surviving pairs lifts through all
characteristic refinements of $F$.
\end{proof}

\begin{claim}{D.48}{Proved}{The two sections of a punctured-torus order-three action}
Let $a$ have order three in the full once-punctured-torus
loop-outer group of $F=\widehat F(x,y)$, and let $E_a$
be its centre-free extension.
There are exactly two $F$-conjugacy classes of order-three
sections in $E_a$.  Each such section $S$ satisfies
\[
                         H^1(C_F(S),\F_3)=0.             \tag{D.48a}
\]
The normalizer in $E_a$ of the torus cusp
$L=\overline{\langle[x,y]\rangle}$ is torsion-free.
If $v$ is any element of this normalizer projecting
to the prescribed generator $a$, and
\[
             v[x,y]v^{-1}=[x,y]^\lambda,\qquad
                            v^3=[x,y]^\kappa ,
\]
then its $3$-adic coordinate satisfies
$\kappa_3\in\Z_3^\times$.
\end{claim}
\passage{Part D(b) of the task; \S{}24, folio 147 verso, scan 293; \S{}30, folios 194--196, scans 383--387; \S{}31, folios 198--199 verso, scans 391--394.}

\begin{proof}
\emph{The actual action and its abelian quotient.}
D.27 gives an actual order-three representative
$\alpha$ of $a$, so
$E_a=F\rtimes_\alpha\langle s\rangle$.
Let $A$ be its matrix on $H_1(F,\Z_3)$.
It has order exactly three by D.33.
Over $\Q_3$ the polynomial $X^2+X+1$ is irreducible:
a primitive cube root would give a square root of
$-3$, whose $3$-adic valuation is odd.
Since $A^3=I$ and $A\ne I$, it follows that
\[
             A^2+A+I=0,\qquad
             \det(A-I)=3,\qquad
             \operatorname{coker}(A-I)\cong C_3.
                                                        \tag{D.48b}
\]
The last equality follows by Smith reduction over
$\Z_3$, since the determinant has valuation one.
Abelianizing the actual semidirect product gives
the exact maximal pro-$3$ abelian quotient
\[
                     E_a^{\rm ab,(3)}\cong C_3^2.        \tag{D.48c}
\]
One factor is (D.48b), the other is the specified
section quotient.

The reduction of $A$ modulo $3$ is a nontrivial
unipotent Jordan block.  Indeed, if $A=I+3B$,
then $A^3=I$ implies
$B(I+3B+3B^2)=0$ after division by $9$.
The second factor is invertible, giving $B=0$,
contrary to $A\ne I$.
The induced action on
$M=H^1(F,\F_3)$ is therefore also one nontrivial
Jordan block of dimension two.
Writing $D$ for the action matrix on this dual
space, the exact cyclic resolution has successive
operators
\[
              D-I,\quad 1+D+D^2=0,\quad D-I,\quad0,\ldots,
 \qquad \operatorname{rank}(D-I)=1.
                                                        \tag{D.48d}
\]
Consequently $H^j(C_3,M)$ has dimension one for
every $j\geq0$.

\emph{Bockstein injectivity and periodicity.}
Let $u\in H^1(E_a,\F_3)$ be the quotient character,
and let $\chi$ be a nonzero invariant character of
$F$, extended to $E_a$ by $\chi(s)=0$.
They form a basis of $H^1(E_a,\F_3)$ by (D.48c).
For the coefficient sequence
\[
              0\longrightarrow\F_3
                 \xrightarrow{\;3\;}\Z/9
                 \longrightarrow\F_3\longrightarrow0
\]
write $\mathfrak b$ for the Bockstein.
Every homomorphism $E_a\to\Z/9$ factors through
(D.48c), so its values are in $3(\Z/9)$ and its
reduction modulo $3$ is zero.
The coefficient long exact sequence therefore makes
\[
     \mathfrak b:H^1(E_a,\F_3)\longrightarrow
                              H^2(E_a,\F_3)
                  \quad\hbox{injective}.                \tag{D.48e}
\]
Put $v=\mathfrak b(u)$ and $w=\mathfrak b(\chi)$.
The notation $v$ in this cohomology calculation
denotes a class, not the normalizer element in the
statement.
The class $v$ is the inflation of the usual degree-two
period class of $C_3$; it is nonzero on the section
$\langle s\rangle$.
The classes $v,w$ are linearly independent by
(D.48e).

The Hochschild--Serre spectral sequence has only
rows zero and one.  By (D.48d), at each positive
total degree its $E_2$ page has two
one-dimensional terms.  Thus its degree-two
abutment has dimension at most two.
The two independent Bockstein classes make its
dimension exactly two.
The only possible differential is
$d_2:E_2^{j,1}\to E_2^{j+2,0}$.
At $j=0$ it vanishes because $\chi$ extends to
$E_a$, and at $j=1$ it vanishes because the
one-dimensional term $E_2^{1,1}$ must survive
to obtain dimension two in degree two.

Multiplication by the period class $v$ identifies
each cyclic-resolution term with the term two
places later.
On the row with coefficient $M$ this also holds
starting at $j=0$, since the norm in (D.48d) is
zero, so $H^0(C_3,M)\to H^2(C_3,M)$ is the identity
on $\ker(D-I)$.
On the row with trivial coefficients it is the
usual period isomorphism.
The multiplicative spectral sequence therefore
makes every remaining $d_2$ a multiple of one of
the two already vanishing differentials.
It collapses, and multiplication by $v$ is an
isomorphism on its two associated graded pieces
at every positive total degree.  Hence
\[
 \begin{gathered}
        \dim H^n(E_a,\F_3)=2\quad(n\geq1),\\
        v:H^n(E_a,\F_3)\xrightarrow{\ \sim\ }
                             H^{n+2}(E_a,\F_3)
                       \quad(n\geq1),\\
                         wv^m\ne0\quad(m\geq0).
 \end{gathered}                                         \tag{D.48f}
\]

\emph{Exactly two full section classes.}
Every finite subgroup of $E_a$ injects into its
quotient $C_3$, because $F$ is torsion-free.
For an order-three section $S$ the normalizer is
the exact direct product
\[
                         N_{E_a}(S)=C_F(S)\times S.      \tag{D.48g}
\]
Conjugation by a normalizer element induces the
identity on $S$: projection to the abelian quotient
$C_3$ identifies every element of $S$ with its
label, so that label cannot change under
conjugation.
Multiplying by the section element with the same
quotient label then leaves precisely $C_F(S)$.

The full-profinite Brown restriction theorem
\cite[Theorem~2 and \S1.1]{A-KochloukovaZalesskii2007},
applied with the open subgroup $F$ of
$3$-cohomological dimension one, identifies
sufficiently high cohomology with the product
over (D.48g).
Since $C_F(S)$ is a closed subgroup of $F$,
its $3$-cohomological dimension is at most one.
Writing $k$ for the number of section classes,
the dimensions in (D.48f) therefore give
\[
                  k+\sum_{[S]}\dim H^1(C_F(S),\F_3)=2.
                                                        \tag{D.48h}
\]
In particular the class set is finite.

It cannot consist of a single class.
Otherwise take that class to be the chosen
$S=\langle s\rangle$.
On its normalizer (D.48g), the character $\chi$
is pulled back from $C_F(S)$, since $\chi(s)=0$.
Its Bockstein therefore restricts to zero:
\[
       w|_{N_{E_a}(S)}
        =\operatorname{pr}_{C_F(S)}^*
               \mathfrak b(\chi|_{C_F(S)})=0,
                                                        \tag{D.48i}
\]
because $H^2(C_F(S),\F_3)=0$.
Every nonzero class $wv^m$ from (D.48f) would
then have zero Brown restriction, a contradiction
for large $m$.
Thus $k\geq2$; (D.48h) makes $k=2$ and makes
both summands zero, proving (D.48a).
As in D.46, $E_a$-conjugacy equals $F$-conjugacy:
for any section $S$ choose its unique element
with the same quotient label as a given
conjugator and multiply the conjugator by its
inverse.

\emph{The peripheral normalizer.}
We spell out the boundary consequence of the
vanishing in (D.48a).
A closed subgroup $C\leq F$ with
$H^1(C,\F_3)=0$ cannot contain a conjugate of
the full $3$-part
$\langle[x,y]^{e_3}\rangle\cong\Z_3$
of the torus cusp.
First its supernatural index in $F$ must have
infinite $3$-part.  Otherwise choose an open
$U_0\geq C$ whose index has maximal $3$-valuation.
For every open $V$ with $C\leq V\leq U_0$,
the index $[U_0:V]$ is prime to $3$.
Transfer makes
$H^1(U_0,\F_3)\to H^1(V,\F_3)$ injective.
The restriction-colimit comparison of D.32
would then inject the nonzero group
$H^1(U_0,\F_3)$ into $H^1(C,\F_3)$,
a contradiction.

Choose an open $U\geq C$ with $3\mid[F:U]$.
Its corresponding finite cover of the punctured
torus has a lift of the selected cusp of degree
prime to $3$, because its completed boundary
subgroup contains the selected full $\Z_3$.
The exact discrete boundary comparison uses
the dense discrete representative for the
profinite conjugator, as in D.32 and D.33.
There are at least two boundary lifts: if there
were only one, its degree would be $[F:U]$.
The selected boundary class is nonzero in
$H_1(U,\F_3)$, since the only relation among
the boundary classes is their sum.
A boundary-dual character restricts nontrivially
to the selected $\Z_3$ in $C$; the scalar
comparison is the inverse of its covering
degree modulo $3$.
This again contradicts $H^1(C,\F_3)=0$ and proves
the assertion.

For $L=\overline{\langle[x,y]\rangle}$, the
peripheral normalizer has the exact sequence
\[
           1\longrightarrow L\longrightarrow N_{E_a}(L)
                    \longrightarrow C_3\longrightarrow1.
                                                        \tag{D.48j}
\]
Surjectivity follows from preservation of the cusp
class, and its kernel is $L$ by A.7.
The multiplier satisfies $\lambda^3=1$, and
$\lambda_3=1$, since $\Z_3^\times$ has no
nontrivial root of unity of order three.
An order-three section in (D.48j) would therefore
centralize the full cusp $\Z_3$.
Its fixed subgroup in $F$ would contradict the
boundary consequence of (D.48a).
Every torsion subgroup of (D.48j) injects into
$C_3$, so this proves that the normalizer is
torsion-free.

Finally take the element called $v$ in the
statement, with $v^3=[x,y]^\kappa$.
Commuting $v$ with its own cube gives
$(\lambda-1)\kappa=0$.
Where $\lambda_p\ne1$, the coordinate $\kappa_p$
and the norm $1+\lambda_p+\lambda_p^2$ are zero.
Where $\lambda_p=1$, that norm is $3$.
If $\kappa_3$ were divisible by $3$, the exact
norm equation
\[
                (1+\lambda+\lambda^2)z=-\kappa
                               \quad\text{in }\widehat\Z
                                                        \tag{D.48k}
\]
would consequently have a solution: divide by
$3$ at the positive coordinates, and use zero
at the other coordinates.
The element $v[x,y]^z$ would then have order
three and project to $a$, contradicting the
torsion-free conclusion.
Thus $\kappa_3$ is a unit, as asserted.

The program
\texttt{code/part\_c/torus\_order3\_sections\_d48.py}
checks the integral and mod-$3$ matrix calculations,
the abelian character and Bockstein input in its
finite semidirect quotient, and the cyclic
resolution operators.
The full restriction and finite-cover arguments
above establish the section and normalizer
assertions.  They do not assert global triviality
of the order-three multiplier or geometric
conjugacy.
\end{proof}

\begin{claim}{D.49}{Proved}{The standard centralizer and the two-element conjugator fibres}
Let $P$ be the full pure three-cusp group of D.32 and
$\delta(x)=x^{-1}$, $\delta(y)=y^{-1}$.
Its full centralizer and normalizer are
\[
             C_P(\delta)=N_P(\langle\delta\rangle)
                         =\langle\delta\rangle.         \tag{D.49a}
\]
Define the actual conjugacy orbits
\[
 \mathcal O_P=\{\beta\delta\beta^{-1}:\beta\in P\},
 \qquad
 \mathcal O_{GT_0}=\{\beta\delta\beta^{-1}:\beta\in GT_0\}.
\]
Every fibre of $P\to\mathcal O_P$ is exactly the
right coset $\{\beta,\beta\delta\}$, and
\[
                    \mathcal O_{GT_0}
                            =\mathcal O_P\cap GT_0.     \tag{D.49b}
\]
For $\alpha\in\mathcal O_P$, its two conjugators
have the exact normalized formulas
\[
 \begin{array}{ll}
  x\mapsto x^\mu,& y\mapsto h y^\mu h^{-1},\\
  x\mapsto x^{-\mu},& y\mapsto h y^{-\mu}h^{-1},
 \end{array}
       \qquad \mu\in\widehat\Z^\times,\quad h\in F'.   \tag{D.49c}
\]
Here $h$ is unique, and the unordered pair
$\{\mu,-\mu\}$ is unique.
In D.35's solution space, a choice $(\mu,h,u)$
therefore gives all the solutions over $\alpha$ as
\[
 (\mu,h,u(xy)^a),\qquad
 (-\mu,h,x^{-\mu}u(xy)^a),\qquad a\in\widehat\Z.
                                                        \tag{D.49d}
\]
These assertions concern the specified orbit map;
they do not assert that it contains every full
pure or $GT_0$ involution.
\end{claim}
\passage{\S{}30; folios 194--196; scans 383--387; \S{}31; folios 198--199 verso; scans 391--394; Part D of the task: full GT\_0 involution conjugacy}

\begin{proof}
Lochak--Schneps prove that complex conjugation is
self-centralizing in the usual $GT$
\cite[Proposition 4(ii), with Lemma 9]{R-LS}.
Their argument uses the first two defining relations
and applies to $GT_0$ as well.  We give the passage
from the full group $P$ explicitly, together with
the normalized fibre formulas.

Let $\beta\in C_P(\delta)$ have normalized pair
$(\mu,k)$, so
$\beta(x)=x^\mu$ and $\beta(y)=k^{-1}y^\mu k$,
with $k\in F'$.
D.32's normalization is a homomorphic section
into $\Aut(F)$, so the commutation is an actual
automorphism equality.
Its value on $y$ makes $k\delta(k)^{-1}$
centralize $y^\mu$.
The peripheral centralizer is $\langle y\rangle$
by A.7, while $k\delta(k)^{-1}\in F'$.
Abelianization therefore gives $k=\delta(k)$.
The exact free-product identification of D.35,
$F\rtimes_\delta C_2\cong C_2\amalg C_2\amalg C_2$,
gives $C_F(\delta)=1$.  Hence $k=1$.

We also need the precise inclusion in $GT_0$.
For each geometric cusp permutation $\sigma\in\Sigma$,
form $r_\sigma=\beta^{-1}\sigma\beta\sigma^{-1}$.
It lies in $P^1=\ker\lambda$, and commutes with
$\delta$, because both $\beta$ and $\sigma$ do.
D.35 proved $C_{P^1}(\delta)=1$, so $r_\sigma=1$.
The exact identity $GT_0=C_P(\Sigma)$ of D.35
therefore puts this $\beta$ in $GT_0$.

Its second Harbater--Schneps relation, with
$k=1$ and $m=(\mu-1)/2\in\widehat\Z$, is
\[
                (xy)^{-m}y^m x^m=1,
       \qquad\text{equivalently}\quad
                         (xy)^m=y^m x^m.                \tag{D.49e}
\]
For every positive integer $n$ use the explicit
finite Heisenberg quotient with coordinates
$(a,b,c)\in(\Z/n)^3$ and multiplication
\[
 (a,b,c)(a',b',c')=(a+a',b+b',c+c'+ab'),
       \quad x=(1,0,0),\quad y=(0,1,0).
\]
The two sides of (D.49e) have coordinates
\[
 (m,m,m(m+1)/2),\qquad(m,m,0),
                                                        \tag{D.49f}
\]
where the binomial expression is the continuous
integer-valued function on $\widehat\Z$.
These identities first follow for nonnegative
integers by multiplication, and then for all
$m\in\widehat\Z$ by continuity; modulo $n$ they
depend only on $m\bmod2n$.
Equality in every such quotient gives
$m(m+1)/2=0$ in $\widehat\Z$, and hence
\[
                           \mu^2=1.                     \tag{D.49g}
\]
Since $k=1$, the actual square of $\beta$ fixes
both $x$ and $y$ by (D.49g).  Thus $\beta^2=1$.
D.34 gives the global multiplier $-1$ for a
nonidentity pure involution; the identity has
multiplier $1$.
The normalized pair is therefore $(1,1)$ or
$(-1,1)$, proving $C_P(\delta)=\langle\delta\rangle$.
A normalizer of its order-two subgroup must
centralize its sole nonidentity element, proving
the normalizer equality in (D.49a).

Two elements $\beta,\beta'$ have the same image
under the orbit map exactly when
$\beta^{-1}\beta'\in C_P(\delta)$.
This proves its two-element fibres.
For an orbit element lying in $GT_0$, D.35's
proved equality of conjugator sets places both
conjugators in $GT_0$; the reverse inclusion is
immediate.  This proves (D.49b).

Write one conjugator in D.35's convention as
$\beta_{\mu,h}$.  Right composition with $\delta$
gives exactly $\beta_{-\mu,h}$, so (D.49c) follows
with the same $h$.
They are distinct because $\mu\ne-\mu$ in
$\widehat\Z^\times$.
Moreover actual normalized conjugation identifies
$C_F(\alpha)$ with $\beta(C_F(\delta))=1$.
The description of all coboundaries in (D.35g)
then makes $h$ the unique solution of
$h^{-1}\alpha(h)=f$ in $F'$.
The two-element fibre proves uniqueness of the
unordered multiplier pair.

Finally put $A=x^\mu$, $B=h y^\mu h^{-1}$ and
suppose $AB=u(xy)^\mu u^{-1}$.
The exact word identity
\[
           A^{-1}B^{-1}=A^{-1}(AB)^{-1}A
\]
shows that $x^{-\mu}u$ is a third-cusp conjugator
for the opposite multiplier.
D.35's centralizer calculation gives all other
third-cusp conjugators by right multiplication
by $(xy)^a$.  This proves (D.49d).
The program
\path{code/part_a/standard_involution_centralizer.py}
checks (D.49f) over sixteen finite moduli and
checks the displayed word identity.
The full conclusion (D.49a) uses the proved
global-sign theorem D.34, beyond those finite
checks.
\end{proof}

\begin{claim}{D.50}{Proved}{The finite two-subgroups of the full punctured-torus group}
Let
$T=\Out_{\rm lac}(\widehat F(x,y);
\overline{\langle[x,y]\rangle})$ be the full
once-punctured-torus loop-outer group.
The possible abstract isomorphism types of its finite
$2$-subgroups are exactly
\[
                  1,\quad C_2,\quad C_4,\quad
                         C_2\times C_2,\quad D_8,
                                                        \tag{D.50a}
\]
where $D_8$ has order eight. All five types occur.
For a subgroup of type $C_2\times C_2$ or $D_8$,
the full multiplier takes its values in the globally
constant signs $\{1,-1\}\subset\widehat\Z^\times$.
In $C_2\times C_2$ there is one positive involution
and two negative involutions. In $D_8$ the rotation
subgroup has multiplier $1$ and the four reflections
have multiplier $-1$.

For an arbitrary cyclic subgroup of order four,
a generator $a$ satisfies the already proved
restrictions
\[
           \lambda(a)^2=1,\qquad \lambda(a)_2=1,
           \qquad \lambda(a^2)=1.                       \tag{D.50b}
\]
The abstract classification (D.50a) does not establish
conjugacy of arbitrary such subgroups, or replace
(D.50b) by a global multiplier determination.
\end{claim}
\passage{Section 30, folios 194-196, scans 383-387; Section 31, folios 198-199 verso, scans 391-394; Part D(b), full loop-outer finite actions in type (1,1)}

\begin{proof}
D.33 makes the actual homology homomorphism
\[
                     A_2:T\longrightarrow\GL_2(\Z_2)
                                                        \tag{D.50c}
\]
injective on every finite $2$-subgroup $K$ and proves
$|K|\leq8$. We classify its image using that same
specified injection.

No matrix in $\GL_2(\Q_2)$ has order eight:
an eigenvalue of order eight has minimal polynomial
$X^4+1$ over $\Q_2$, which is irreducible by
Eisenstein after $X\mapsto X+1$.
The same argument excludes higher powers of two.
Every matrix $R$ of order four has minimal
polynomial $X^2+1$, since this polynomial is
irreducible over $\Q_2$ and $R^2\ne I$.
Consequently
\[
                         R^2=-I.                       \tag{D.50d}
\]
An exponent-two subgroup is abelian, because
$(ab)^2=a^2=b^2=1$ implies $ab=ba$.
Its matrices are commuting diagonalizable
involutions over $\Q_2$, and simultaneous
diagonalization embeds it in
$\{\operatorname{diag}(\epsilon_1,\epsilon_2):
\epsilon_i=\pm1\}$. Its order is at most four.
Thus an image of order eight contains an $R$
of order four.

The matrix centralizer of $R$ is the multiplicative
group of the field $\Q_2[R]\cong\Q_2(i)$.
For completeness, choose a nonzero vector $v$:
the vectors $v,Rv$ are independent, since an
eigenvalue in $\Q_2$ would contradict the
irreducibility of $X^2+1$.
A commuting endomorphism is determined by its
value on $v$ and is therefore $aI+bR$.
Every element of the finite $2$-subgroup in this
field has order at most four by the preceding
matrix argument. It is consequently one of
the at most four roots of $X^4-1$.
Since $\langle R\rangle$ already has four elements,
\[
                       C_K(R)=\langle R\rangle
                                                        \tag{D.50e}
\]
for a subgroup of order eight.

Choose $S\in K\setminus\langle R\rangle$.
The index-two subgroup $\langle R\rangle$ is normal;
by (D.50e) conjugation by $S$ induces its nonidentity
automorphism. Hence
\[
          SRS^{-1}=R^{-1},\qquad
                    S^2\in\{I,R^2\}.                  \tag{D.50f}
\]
The second inclusion follows because $S^2\in
\langle R\rangle$ and $S$ fixes its own square,
whereas it inverts $R$.
The alternative $S^2=I$ gives the dihedral
presentation of order eight. We exclude the
alternative $S^2=R^2=-I$ explicitly.

In the basis $(v,Rv)$ used above one has
\[
 R=\begin{pmatrix}0&-1\\1&0\end{pmatrix}.
\]
The equation $SR=-RS$ forces
\[
             S=\begin{pmatrix}a&b\\b&-a\end{pmatrix},
             \qquad S^2=(a^2+b^2)I.
                                                        \tag{D.50g}
\]
Thus the second alternative would require
$a^2+b^2=-1$ in $\Q_2$.
Here is the exact valuation obstruction.
Put $m=\min(v_2(a),v_2(b))$ and write
$a=2^m u$, $b=2^m v$, with $u,v\in\Z_2$
and at least one a unit. If both are units,
$u^2+v^2\equiv2\pmod8$, so this sum has
valuation one. If exactly one is a unit,
the sum has valuation zero.
On the other hand the equation says
\[
                      u^2+v^2=-2^{-2m};
                                                        \tag{D.50h}
\]
the valuation of its right side is even.
Both-unit values are therefore impossible.
The other case forces $m=0$; reduction modulo
four then gives $1=-1$, also impossible.
Zero entries cause no exception to this
argument. This excludes the quaternion
alternative and leaves only $D_8$ at order
eight. Orders four, two, and one give the
remaining groups in (D.50a).

We specify actual full-group realizations.
On $\widehat F(x,y)$ define
\[
 \begin{array}{c|cc}
          &x&y\\ \hline
 r        &y&x^{-1}\\
 s        &y&x .
 \end{array}                                            \tag{D.50i}
\]
The free universal property gives continuous
endomorphisms, and the displayed relations
$r^4=s^2=1$, $srs=r^{-1}$, verified on
$x,y$, give their inverses and make them
automorphisms.
For $c=[x,y]$ direct multiplication gives
\[
                   r(c)=x^{-1}cx,\qquad
                   s(c)=c^{-1}.                        \tag{D.50j}
\]
Thus their outer classes lie in $T$.
Their homology matrices are
\[
             R=\begin{pmatrix}0&-1\\1&0\end{pmatrix},
       \qquad S=\begin{pmatrix}0&1\\1&0\end{pmatrix},
                                                        \tag{D.50k}
\]
and the eight matrices $R^iS^j$
($0\leq i<4$, $0\leq j<2$) are distinct.
These maps therefore exhibit an embedded
$D_8$ in $T$. The subgroups
$\langle r\rangle$, $\langle r^2,s\rangle$,
$\langle s\rangle$, and the identity realize
$C_4$, $C_2\times C_2$, $C_2$, and $1$.

We finish with the asserted multiplier data.
D.43 makes every involution's multiplier a
global sign. A positive nonidentity involution
has determinant $1$ in (D.50c). Its order-two
matrix is diagonalizable in characteristic
zero, and faithfulness excludes $I$; its matrix
is consequently $-I$.
There is at most one such involution in any
finite $2$-subgroup.
For $C_2\times C_2$, simultaneous diagonalization
identifies its full image with the four sign
matrices above, giving exactly one positive
and two negative involutions.
For $D_8$, the square of a rotation is
positive by (D.50d), so every other involution
is negative. The reflections generate $D_8$,
and a rotation is a product of two reflections.
The homomorphic multiplier is therefore
globally $1$ on the rotations and globally
$-1$ on the reflections.

Finally let $a$ have order four in an arbitrary
cyclic subgroup. Its matrix in (D.50c) has
minimal polynomial $X^2+1$ and determinant one.
Thus $\lambda(a)_2=1$ and $\lambda(a^2)_2=1$.
D.43 applied to the involution $a^2$ gives
$\lambda(a^2)=1$ globally. The multiplier
homomorphism then gives (D.50b).
No assertion of conjugacy follows from this
matrix classification.

The finite certificate
\path{code/part_c/torus_two_subgroups_d50.py}
verifies the full-word maps (D.50i)--(D.50j),
their eight distinct matrices, and the unit
square computations modulo eight used in
(D.50h). The passage from that finite
calculation to all of $\Q_2$ is the valuation
argument just given.
\end{proof}

\begin{claim}{D.51}{Proved}{The second index-three kernel and its three peripheral normalizers}
Let $a$ have order three in the full punctured-torus
loop-outer group, put $F=\widehat F(x,y)$,
$c=[x,y]$, and let $E$ be its centre-free extension.
Besides $F$, there is exactly one other torsion-free
normal subgroup $U$ of index three in $E$.
Write $V=F\cap U$. Then
\[
 \begin{gathered}
 [F:V]=[U:V]=3,\qquad E/V\cong C_3^2,\\
 V=\widehat{\pi_1(S_{1,3})},\qquad
 U\ \hbox{is projective},\qquad
 U^{(3)}\cong\widehat F^{(3)}_2 .
 \end{gathered}                                         \tag{D.51a}
\]
The cusp normalizer $N=N_E(\langle c\rangle)$
is contained in $U$. The three cusp subgroups
$L_0,L_1,L_2$ of the displayed degree-three
cover have normalizers $N_i=N_U(L_i)$ with
the exact sequences
\[
              1\longrightarrow L_i
                 \longrightarrow N_i
                 \longrightarrow U/V\longrightarrow1.
                                                        \tag{D.51b}
\]
They are three distinct $U$-conjugacy classes,
cyclically permuted by $E/U$. All have the
given multiplier $\lambda$ under the same
generator label in $U/V$.

For a prime $q\ne3$ at which $\lambda_q$ is
a nonidentity cube root of unity, the
$U/V$-character multiplicities of
$H_1(V,\Q_q)$, ordered as $(1,\lambda_q,\lambda_q^2)$,
are exactly one of
\[
                         (1,3,0),\qquad (0,2,2).
                                                        \tag{D.51c}
\]
Correspondingly $U^{(q)}$ is free pro-$q$ of rank
one or zero. These are constraints on the actual
subgroup $U$; they do not identify it with a full
free profinite group.
\end{claim}
\passage{Section 30, folios 194-196, scans 383-387; Section 31, folios 198-199 verso, scans 391-394; Part D(b), full loop-outer finite actions in type (1,1)}

\begin{proof}
\emph{The section and cusp lines in the actual abelian quotient.}
Use D.27 to choose an actual order-three representative
$\alpha$ and write $E=F\rtimes_\alpha\langle s\rangle$.
Let $u:E\to\F_3$ be its quotient character, with $u(s)=1$.
D.48 gives exactly two full section classes,
$S_0=\langle s\rangle$ and $S_1$, and gives
$E^{\rm ab,(3)}\cong C_3^2$.

Let $M$ be the reduction of the homology matrix of
$\alpha$ modulo three. By D.48 it is a nonidentity
unipotent matrix with $(M-I)^2=0$.
On $\F_3^2$ write
$\omega((a,b),(a',b'))=ab'-ba'$.
Choose $w$ with $e=(M-I)w\ne0$.
Then $(e,w)$ is a basis, $Me=e$, $Mw=w+e$,
and $d=\omega(e,w)\ne0$.
The linear functional
\[
                  \chi(r)=\omega(e,r)/d
                                                        \tag{D.51d}
\]
is $M$-invariant. It gives a character of $F$,
which we extend to $E$ by $\chi(s)=0$.
Thus $(u,\chi):E\to\F_3^2$ is the actual
abelian quotient above.

The generator of $S_1$ labelled by $u=1$
has $\chi$-value $b_1\ne0$.
Here is the required separation, using the maps
already proved in D.48 rather than an assumption
about section labels.
The Bockstein of $\chi$ has nonzero high periodic
multiples. On a section normalizer its restriction
is the Bockstein of $\chi$ restricted to the section:
the fixed-group first cohomology vanishes by D.48.
If $\chi$ vanished on both section classes,
all those multiples would have zero full Brown
restriction. This contradicts D.48f.
It follows that $b_1\ne0$.

To locate the cusp line, use the characteristic
class-two exponent-three quotient
\[
 Q=\F_3^2\times\F_3,\qquad
 (r,z)(r',z')=(r+r',z+z'+\tfrac12\omega(r,r')) .
                                                        \tag{D.51e}
\]
The images of $x,y$ are the two standard vectors
with central coordinate zero, and $c$ maps to
$Z=(0,1)$.
Since $\lambda_3=1$, the induced automorphism is
\[
               \theta(r,z)=(Mr,z+\ell(r))
                                                        \tag{D.51f}
\]
for a linear functional $\ell$.
For an element $(r,z)s$ of $Q\rtimes_\theta C_3$
with quotient label one, direct multiplication gives
\[
       ((r,z)s)^3
          =Z^{\,d\chi(r)^2+\ell(e)\chi(r)} .
                                                        \tag{D.51g}
\]
To check the formula, write $r=be+tw$, so
$\chi(r)=t$. The successive vector parts are
$r,r+te,r+2te$, whose sum is zero.
The sum of the three central correction terms
is $\ell((2I+M)r)=t\ell(e)$.
The three alternating cross terms, multiplied
by $1/2$, sum to $dt^2$ in $\F_3$.
This proves (D.51g) for every element, including
the arbitrary central coordinate $z$.

The image of the actual generator of $S_1$
has cube one. Equation (D.51g) and $b_1\ne0$
therefore give $\ell(e)=-db_1$.
Thus its cube polynomial is
\[
                           db(b-b_1).
                                                        \tag{D.51h}
\]
Choose $v\in N$ with $u(v)=1$, as in D.48.
It satisfies $v^3=c^\kappa$ with
$\kappa_3$ a unit. Its image in (D.51e)
has nonzero cube exponent. The two roots of
(D.51h) are $0$ and $b_1$; the only other
element of $\F_3$ is $-b_1$. Consequently
\[
             (u,\chi)(S_0)=\langle(1,0)\rangle,\quad
             (u,\chi)(S_1)=\langle(1,b_1)\rangle,\quad
             (u,\chi)(N)=\langle(1,-b_1)\rangle .
                                                        \tag{D.51i}
\]
The last equality uses that $N$ is generated
topologically by $v$ and $\langle c\rangle$,
and both characters kill $c$.

\emph{The full kernels and peripheral labels.}
Define the continuous characters and subgroups
\[
           \psi=\chi+b_1u,\qquad
           U=\ker\psi,\qquad V=\ker u\cap\ker\chi .
                                                        \tag{D.51j}
\]
These are actual subgroups of $E$.
The character $\psi$ is nonzero on both
section classes in (D.51i), and it kills $N$.
Every nonidentity finite subgroup of $E$
is one of the order-three section groups,
up to conjugacy. Hence $U$ is torsion-free.
The four lines in $E^{\rm ab,(3)}=\F_3^2$
give all normal index-three subgroups of $E$.
Two contain respectively $S_0$ and $S_1$.
The other two are exactly $\ker u=F$ and
$\ker\psi=U$, proving the uniqueness assertion.
Also $E=U\rtimes S_0$, since $\psi|_{S_0}$
is an isomorphism onto $\F_3$.

The restriction of $\chi$ to $F$ is nonzero
and kills its commutator cusp.
Choose a word $f$ in the dense discrete
group $F(x,y)$ with $\chi(f)=1$.
The degree-three covering of the punctured
torus corresponding to $\ker(\chi|_{F(x,y)})$
has three boundary lifts of degree one.
Its Euler characteristic is $-3$; thus its
type is $(1,3)$, and its fundamental group
has rank four.
The induced profinite subgroup is exactly $V$,
with cusp representatives
\[
          L_i=f^i\langle c\rangle f^{-i},
          \qquad
          N_i=f^i N f^{-i}\quad(0\leq i<3).
                                                        \tag{D.51k}
\]
The full peripheral normalizer assertion A.7
gives $N_E(L_i)=N_i$.
Since $N\subset U$ and $U$ is normal in $E$,
all $N_i$ lie in $U$. Their kernels under
$U\to U/V$ are exactly $L_i$, because
$N_i\cap F=L_i$; their images are all of
$U/V$, because $u(v)=1$.
This proves (D.51b), with the common generator
label supplied by $u$.

For the class assertions, the actual quotient
$(u,\chi)$ identifies $E/V$ with $\F_3^2$,
and the image of every $N_i$ is the same
line $U/V$.
The three $V$-conjugacy classes of the cusp
lifts are therefore the three cosets of this
line in $E/V$. The subgroup $U/V$ fixes each
class, so they remain distinct under $U$.
The quotient $E/U$ acts regularly on these
three cosets, proving its asserted cyclic
permutation. This also proves that conjugation
by a generator of $U/V$ fixes each cusp
class with multiplier $\lambda$.

The open subgroup $V$ of $U$ is free profinite
of rank four and $U$ is torsion-free.
Serre's finite-index cohomological-dimension
theorem, used in D.27, makes $U$ projective.
Its maximal pro-three quotient is projective
as a pro-three group, by the finite embedding
problem characterization, and is therefore
free pro-three.
The pro-three topology on $V$ is exactly the
one induced from $U$: taking the intersection
of the three conjugates of a normal
three-power-index subgroup of $V$ produces
one normal in $U$, still of three-power index.
Thus $V^{(3)}$ embeds as an open subgroup
of index three in $U^{(3)}$.
Schreier's formula now reads
\[
                 4=1+3(\operatorname{rank}U^{(3)}-1),
\]
and gives rank two. This proves (D.51a).

\emph{The exact character restriction at another prime.}
The three oriented boundary classes of $V$
have their sole relation given by their sum.
They form a two-dimensional submodule of
$H_1(V,\Q_q)$ on which $U/V$ acts by
$\lambda_q$.
Filling the three boundaries gives the compact
degree-three cover of the original torus.
Its map on homology is the inclusion of a
rank-two sublattice of index three, hence an
isomorphism over $\Q_q$, and over $\Z_q$
when $q\ne3$.
Conjugation by $v$ projects to the original
outer action $a$ on the base; its inner
factor acts trivially on that base homology.
Therefore, if $A_q$ denotes the original
matrix, the exact boundary sequence and
averaging over $C_3$ give
\[
       H_1(V,\Q_q)\cong
                  \lambda_q^{\oplus2}
                    \oplus(\Q_q^2,A_q)
                 \quad\hbox{as }U/V\hbox{-modules}.
                                                        \tag{D.51l}
\]
This comparison is induced by the stated
boundary inclusion and covering homology map.

At the specified prime, $A_q^3=I$ and
$\det A_q=\lambda_q$.
The three cube roots $1,\lambda_q,\lambda_q^2$
are distinct in $\Q_q$. Diagonalization gives
exactly the two eigenvalue multisets
\[
                 \{1,\lambda_q\},\qquad
                 \{\lambda_q^2,\lambda_q^2\}.
                                                        \tag{D.51m}
\]
Adding the two boundary characters in (D.51l)
proves (D.51c).
For $q\ne3$, the same decomposition holds
over $\Z_q$: the averaging idempotents have
unit denominator three, and the cube-root
differences are units.
Hochschild--Serre for $V\triangleleft U$ then
identifies $H^1(U,\F_q)$ with the invariant
subspace of $H^1(V,\F_q)$, of dimension one
or zero respectively.
The projectivity of $U$ makes its maximal
pro-$q$ quotient free pro-$q$, with this
number of generators. This proves the last
assertion.

The program
\path{code/part_c/torus_order3_sibling_d51.py}
checks (D.51g) for every central-linear
automorphism above every nonidentity unipotent
matrix over $\F_3$, the resulting section and
cusp lines, and the finite eigenvalue patterns
in (D.51m). These finite calculations validate
the displayed formulas. The full subgroup,
topology, and projectivity conclusions are
proved by the maps above.
\end{proof}

\paragraph{Exact remaining recognition problem.}
The construction supplies a specified torsion-free
projective group $U$ with an open full free subgroup
of rank four and three specified normalizers.
It supplies no identification of $U$ with
$\widehat F_2$.
At the primes in (D.51c), its maximal pro-$q$
quotient has rank zero or one, whereas the
maximal pro-three quotient has rank two.
The existence or exclusion of such a full
$U$ with these three normalizers remains the
step needed for this approach to determine
every torus order-three multiplier.

\begin{claim}{D.52}{Proved}{An abstract allowed list for every finite full torus subgroup}
Let $T$ be the full once-punctured-torus loop-outer
group of D.50. Every finite subgroup of $T$ has
one of the following twelve abstract types:
\[
\begin{gathered}
1,\ C_2,\ C_3,\ C_4,\ C_2\times C_2,\ C_6,\ S_3,\ D_8,\\
C_{12},\ D_{12},\ \operatorname{Dic}_{12},\ D_{24}.
\end{gathered}                                             \tag{D.52a}
\]
Here $D_{2n}$ has order $2n$, and
$\operatorname{Dic}_{12}=C_3\rtimes C_4$, with
a generator of $C_4$ inverting $C_3$.
The nine types obtained by omitting
$C_{12},\operatorname{Dic}_{12},D_{24}$ all occur.
Occurrence of those three remaining types is
not determined by this allowed list.
The subsequent full covering arguments in D.54--D.55
exclude all three.

Every subgroup of type $S_3,D_{12}$ or $D_{24}$
has the global usual reflection multiplier:
rotations have multiplier $1$ and reflections
have multiplier $-1$.
In a subgroup of type $\operatorname{Dic}_{12}$,
the subgroup of order three has multiplier $1$,
and every generator $b$ of a cyclic subgroup
of order four satisfies
\[
       \lambda(b)^2=1,\qquad
       \lambda(b)_2=1,\qquad\lambda(b)_3=-1.              \tag{D.52b}
\]
For a subgroup of type $C_{12}$, its order-four
generator instead has $\lambda(b)_2=\lambda(b)_3=1$,
and still has $\lambda(b)^2=1$.
These restrictions retain the unresolved cyclic
order-three and order-four multiplier questions.
\end{claim}
\passage{Section 30, folios 194-196, scans 383-387; Section 31, folios 198-199 verso, scans 391-394; Part D(b), full loop-outer finite actions in type (1,1)}

\begin{proof}
\emph{The two primary representations and their actual kernels.}
For a finite subgroup $G\leq T$, D.33 gives
\[
 |G|\mid24,\qquad
 A_p:G\longrightarrow\GL_2(\Z_p),\qquad
 \det A_p(g)=\lambda(g)_p.                              \tag{D.52c}
\]
It makes $A_p$ injective on every finite
$p$-subgroup. Put $K=\ker(A_2|_G)$.
An element of order two in $K$ would contradict
that injection. Cauchy's theorem and (D.52c)
therefore give
\[
                      K=1\quad\hbox{or}\quad K\cong C_3.
                                                               \tag{D.52d}
\]
We classify the actual image $H=A_2(G)$ and
then the extension by this specified kernel.
We do not identify $G$ with a representation
at one prime when the kernel is nontrivial.

\emph{Finite images over $\Q_2$.}
The two-subgroups of $H$ have the types in D.50.
Indeed a Sylow two-subgroup of the preimage
of any such subgroup maps isomorphically onto
it, since the kernel has odd order.
The matrix arguments in D.50 also explicitly
exclude the quaternion group of order eight.

Suppose $3$ divides $|H|$.
Its number of Sylow three-subgroups is either
one or four. We exclude four.
Conjugation gives a transitive homomorphism
$H\to S_4$. Its kernel has order at most six.
If that kernel had order divisible by three,
its unique Sylow three-subgroup would be
characteristic in the kernel and normal in $H$,
contrary to the assumed four Sylow subgroups.
Thus the transitive image has order divisible
by both four and three, and is $A_4$ or $S_4$.
The kernel consequently has order one or two.

There is no faithful two-dimensional
representation of $A_4$ over $\Q_2$.
Its normal Klein four subgroup would be
simultaneously diagonalizable and, by
faithfulness, would be the full group
of diagonal sign matrices. One of its
three involutions would therefore be $-I$.
Conjugation fixes this scalar, whereas the
order-three elements of $A_4$ cyclically
permute all three involutions.
This is a contradiction.
It excludes the cases with trivial kernel,
whose images are $A_4$ or $S_4$.
In the remaining case $|H|=24$ and the image
is $A_4$ with a kernel of order two.
The preimage $P$ of its Klein four subgroup
is a normal subgroup of order eight, hence
is $D_8$ by D.50. But
$|\operatorname{Aut}(D_8)|=8$: a generator
of its rotation subgroup has two possible
images, and a chosen reflection has four.
Every element of order three therefore
centralizes $P$. Multiplication by a Sylow
three-subgroup gives $H=P\times C_3$, which
has only one Sylow three-subgroup.
This contradiction proves that $H$ has a
normal subgroup $P_3$ of order three.

Let $R$ generate $P_3$.
Its minimal polynomial is
$X^2+X+1$, irreducible over $\Q_2$ by
reduction modulo two.
The basis $(v,Rv)$ for any nonzero vector $v$
identifies the matrix centralizer with
\[
 L_2^\times,\qquad L_2=\Q_2(\zeta_3).
                                                               \tag{D.52e}
\]
Under this identification, a commuting
endomorphism is multiplication by $a+b\zeta_3$.
The full normalizer consists exactly of these
multiplication maps and their compositions
with $\sigma(\zeta_3)=\zeta_3^{-1}$:
conjugation must send $R$ to $R$ or $R^{-1}$,
and division by $\sigma$ in the second case
gives a commuting map.

The field $L_2$ contains no element of order
four. To check this directly, the equation
$(a+b\zeta_3)^2=-1$ has coefficient equations
\[
                   a^2-b^2=-1,\qquad b(2a-b)=0.
                                                               \tag{D.52f}
\]
They require either $a^2=-1$ or $a^2=1/3$
in $\Q_2$. Both are impossible: these are
units congruent respectively to seven and
three modulo eight, whereas an odd square
is congruent to one.
Since $|H|\mid24$, it follows that
$C=H\cap L_2^\times$ is $C_3$ or $C_6$.
Here its two-part is contained in
$\{1,-1\}$, and its three-part is $P_3$,
so no general classification of local
roots of unity is required.

If $H=C$, these are the two possible groups.
Otherwise choose $B=a\sigma\in H\setminus C$.
Then
\[
 B^2=N_{L_2/\Q_2}(a)\,I
        \in C\cap\Q_2^\times=\{I,-I\}\cap C.             \tag{D.52g}
\]
The last scalar restriction follows because
a finite-order scalar of order dividing
twenty-four in $\Q_2$ is a sign: there are
no primitive cube or fourth roots in $\Q_2$.
For $C=C_3$, equation (D.52g) gives $S_3$.
For $C=C_6$, its two values give respectively
$D_{12}$ and $\operatorname{Dic}_{12}$.
Thus the images with order divisible by
three have the complete allowed list
\[
                  C_3,\ C_6,\ S_3,\ D_{12},
                             \operatorname{Dic}_{12}.  \tag{D.52h}
\]

\emph{The three-adic normalizer and its determinant.}
We will need a statement about any normal
subgroup $P_3\cong C_3$ of the actual $G$.
The map $A_3$ is injective on $P_3$.
Its generator has minimal polynomial
$X^2+X+1$ over $\Q_3$, irreducible since
its discriminant $-3$ has odd valuation.
The same centralizer and normalizer
calculation identifies
\[
 C_{\GL_2(\Q_3)}(A_3(P_3))=L_3^\times,\qquad
 N_{\GL_2(\Q_3)}(A_3(P_3))=L_3^\times\rtimes\langle\sigma\rangle,
 \quad L_3=\Q_3(\zeta_3).                              \tag{D.52i}
\]
This field also has no element of order four:
(D.52f) would require $a^2=-1$, impossible
modulo three, or $a^2=1/3$, impossible by
valuation. Thus every finite-order element
in its centralizer with order dividing
twenty-four lies in $\mu_6$ and has norm one.

There is also no element of norm $-1$ in $L_3$.
Here is an explicit calculation, including
possible nonintegral coordinates.
The norm of $a+b\zeta_3$ is
$a^2-ab+b^2$. Put
$m=\min(v_3(a),v_3(b))$ and write
$a=3^m u,b=3^m v$, where $u,v\in\Z_3$
and at least one is a unit.
When $u+v\not\equiv0\pmod3$, the norm of
$u+v\zeta_3$ has valuation zero and residue
$(u+v)^2=1$ modulo three.
When $u+v\equiv0\pmod3$, both are units;
writing $u=-v+3k$ gives
\[
                  u^2-uv+v^2
                       =3(v^2-3vk+3k^2),               \tag{D.52j}
\]
of valuation one.
The norm $-1$ would have valuation zero.
The second case cannot give that valuation,
since $2m+1\ne0$.
The first forces $m=0$ and residue one,
whereas $-1$ has residue two.
This proves the asserted norm obstruction.

An element of the finite image $A_3(G)$
outside the centralizer has the form
$B=a\sigma$, and $B^2=N(a)I$ is a
finite-order scalar, hence a sign.
The norm obstruction forces $B^2=I$.
In the basis $(1,\zeta_3)$, the determinant
of $\sigma$ is $-1$, so
$\det B=-N(a)=-1$.
Together with the centralizer calculation,
this proves the exact character identity
\[
 \lambda(g)_3=
 \begin{cases}
  1,&g\text{ centralizes }P_3,\\
 -1,&g\text{ inverts }P_3.
 \end{cases}                                            \tag{D.52k}
\]
The argument uses faithfulness only on
$P_3$, not on all of $G$.

\emph{Reconstructing the actual finite subgroup.}
If $K=1$ in (D.52d), the injection into
$H$ gives precisely D.50 and (D.52h).
Now let $K=P_3\cong C_3$.
A Sylow two-subgroup $Q$ of $G$ has order
$|G|/3$ and maps isomorphically onto $H$.
The actual multiplication map
\[
                  P_3\rtimes Q\longrightarrow G,
                         \qquad (a,q)\longmapsto aq,    \tag{D.52l}
\]
is an isomorphism, with action given by
conjugation. D.50 lists the five possibilities
for $Q$.
For $Q=1,C_2,C_4$, their homomorphisms
to $\operatorname{Aut}(C_3)=C_2$ give
$C_3$; $C_6$ or $S_3$; and $C_{12}$ or
$\operatorname{Dic}_{12}$, respectively.

For $Q=C_2\times C_2$ or $D_8$, D.50
already determines its global multiplier.
Equation (D.52k) therefore makes its
action on $P_3$ exactly the reflection
character of D.50.
Write $a$ for a generator of $P_3$.
For $Q=C_2\times C_2$, write $z$ for its
positive involution and $s$ for a negative
one. The element $az$ has order six,
is inverted by $s$, and generates $a,z$
through $(az)^4=a$ and $(az)^3=z$.
This gives an explicit $D_{12}$ presentation.
For $Q=D_8=\langle r,s\rangle$, where
$r$ is its order-four rotation, the
element $ar$ has order twelve, is inverted
by $s$, and satisfies
\[
                       (ar)^4=a,\qquad (ar)^9=r.
                                                               \tag{D.52m}
\]
It gives the explicit $D_{24}$ presentation.
These actual maps finish the allowed list.

\emph{Multipliers and realized cases.}
In every group just obtained that inverts
its normal $C_3$, the multiplier on that
$C_3$ is one: a conjugation-invariant
homomorphism into the abelian group
$\widehat\Z^\times$ takes an inverted
element to its own inverse, and a value
of order dividing three with square one
is one.
In $S_3,D_{12},D_{24}$ each reflection
inverts $C_3$ and hence has three-adic
multiplier $-1$ by (D.52k).
D.43 makes its full multiplier $-1$.
Every rotation is a product of two
reflections, proving the stated global
reflection character.
For $\operatorname{Dic}_{12}$, an
order-four generator inverts $C_3$,
so (D.52k) gives its negative
three-coordinate. D.50 gives its positive
two-coordinate and global square one.
For $C_{12}$ the same generator centralizes
$C_3$, giving instead the positive
three-coordinate. This proves all the
claimed multiplier restrictions.

D.50 realizes its five types by explicit
full profinite free-group automorphisms.
The other four realized types have the
following simultaneous exact surface model.
On $\R^2/\Z^2$ the integral matrices
\[
       R=\begin{pmatrix}0&-1\\1&1\end{pmatrix},
       \qquad
       S=\begin{pmatrix}0&1\\1&0\end{pmatrix}            \tag{D.52n}
\]
fix the origin and satisfy
$R^6=S^2=I$, $SRS=R^{-1}$.
Their twelve matrices $R^iS^j$ are distinct.
Removing the origin therefore gives an
actual $D_{12}$ action on a once-punctured
torus. Passing to its fundamental group,
then to its full profinite completion and
outer automorphisms, gives a homomorphism
to $T$. A small peripheral circle is
carried to an ellipse of local degree
$\det R=1$ or $\det S=-1$, so this
homomorphism preserves the stated full
cusp with the indicated multiplier.
The induced homology matrices are exactly
$R,S$ in $\GL_2(\widehat\Z)$; their
distinctness proves that the homomorphism
is injective.
The subgroups
$\langle R^2\rangle,\langle R\rangle,
\langle R^2,S\rangle$, and
$\langle R,S\rangle$ realize
$C_3,C_6,S_3,D_{12}$.

The program
\path{code/part_c/torus_finite_types_d52.py}
checks the matrix relations, the local
norm and square residues, the eight
automorphisms of $D_8$, and the explicit
finite semidirect-product identifications
in (D.52l)--(D.52m).
The prime-wise kernel and valuation
arguments above prove the full allowed
list. No assertion of occurrence or
nonoccurrence of its three remaining
types is inferred from those finite checks.
\end{proof}

\begin{claim}{D.53}{Proved}{The unique labelled order-four section class}
Let $a$ have order four in the full punctured-torus
loop-outer group, and let
\[
          1\longrightarrow F\longrightarrow E
             \xrightarrow{\;\pi\;}C_4=\langle a\rangle
             \longrightarrow1,\qquad F=\widehat F_2,
\]
be its centre-free extension. Its homomorphic
sections form exactly one $F$-conjugacy class,
with the generator label $a$ retained.
For an actual lift $b\in E$ with
$b^4=1$ and $\pi(b)=a$, put $B=\langle b\rangle$.
Then
\[
       N_E(B)=C_F(b)\times B,\qquad
                         H^1(C_F(b),\F_2)=0.            \tag{D.53a}
\]
Conjugation by a lift of $a$ fixes exactly one
of the three full section classes for $a^2$
and interchanges the other two.
Consequently $E$ has two conjugacy classes
of subgroups of order two; exactly one of
them is contained in its class of cyclic
subgroups of order four.
\end{claim}
\passage{Section 30, folios 194-196, scans 383-387; Section 31, folios 198-199 verso, scans 391-394; Part D(b), full loop-outer finite actions in type (1,1)}

\begin{proof}
D.46 supplies the actual lift $b$ with its
prescribed label. Put $T=\langle b^2\rangle$
and
\[
               E_2=\pi^{-1}(\langle a^2\rangle)
                      =F\rtimes T.
                                                               \tag{D.53b}
\]
The action of $b^2$ on $H_1(F,\Z_2)$ is
$-I$ by D.46. Thus
\[
                      E_2^{\rm ab,(2)}\cong C_2^3.
                                                               \tag{D.53c}
\]
Use coordinates $(v,u)\in\F_2^2\times\F_2$,
where $u$ is the quotient character and
the first two coordinates are the extended
characters $\chi_1,\chi_2$ of D.46.
The three full section classes of $E_2$
give three distinct points of the affine
plane $u=1$.
Indeed, the degree-zero localized algebra
in D.46g--i has dimension three, is generated
by $\chi_1/u,\chi_2/u$, and restriction to
the three section normalizers identifies
it with $\F_2^3$.
Each restriction is evaluation at
$(\chi_1(t),\chi_2(t))$ for the nonidentity
element $t$ of that section.
If two such points agreed, the two
restriction coordinates would agree on
both generators and hence on the entire
algebra, contradicting the stated product
isomorphism.

Let $R=A_2(a)\in\GL_2(\Z_2)$.
D.50 gives $R^2=-I$.
Moreover $R\not\equiv I\pmod2$.
For otherwise $R=I+2D$ with an integral
matrix $D$, so $R^2=-I$ would give
$I+2D+2D^2=0$, an impossibility modulo two.
Conjugation by $b$ on (D.53c) fixes $u$
and induces the actual affine map
\[
               v\longmapsto\overline R v+d
                        \quad\text{on }u=1             \tag{D.53d}
\]
for a vector $d\in\F_2^2$.
The linear part follows from the
restriction of conjugation to the
specified subgroup $F$; the translation
is the image of a chosen point of $u=1$.
Conjugation by an element of $E_2$
acts trivially on the three
$F$-conjugacy classes, as shown in D.46.
Thus (D.53d) permutes their three points
with order dividing two.
Three distinct points of $\F_2^2$ affinely
span it. Fixing all three would therefore
make this affine map the identity,
contrary to $\overline R\ne I$.
Its permutation is consequently a
transposition, with exactly one fixed class.
The class of $T=\langle b^2\rangle$ is
that fixed class.

Put $C=C_F(T)$ and $N=N_E(T)$.
Since $b$ normalizes $T$, the image of $N$
under $\pi$ is all of $C_4$.
Its kernel is exactly $C$, because an
element of $F$ normalizing $T$ centralizes
its nonidentity element.
The chosen section gives the actual
semidirect-product identification
\[
          N=C\rtimes\langle b\rangle,\qquad
          T\subseteq Z(N),\qquad
          J:=N/T=C\rtimes\langle\overline b\rangle,
                    \quad|\overline b|=2.              \tag{D.53e}
\]
In the last expression the embedding of
$C$ is injective since $C\cap T=1$.
Conjugation by $b$ on $C$ has square the
identity because $b^2$ centralizes $C$.

As a closed subgroup of the full free
profinite group $F$, the group $C$ is
projective and has $\operatorname{cd}_2(C)\leq1$.
D.46 gives $H^1(C,\F_2)=0$.
All its positive-degree mod-two cohomology
therefore vanishes.
Hochschild--Serre for the last extension
in (D.53e) identifies inflation as
\[
               H^*(C_2,\F_2)\xrightarrow{\;\cong\;}
                                  H^*(J,\F_2).         \tag{D.53f}
\]
This uses the displayed full group $C$,
not an identification of $C$ with its
maximal pro-two quotient.

We spell out the full Brown hypotheses
for this new normalizer quotient.
The subgroup $C$ is open in $J$ and has
finite $2$-cohomological dimension.
It is torsion-free. Every finite subgroup
of $J$ consequently injects into $J/C=C_2$;
in particular $J$ has no subgroup $C_2^2$.
Theorem~2 of
\cite[\S1.1]{A-KochloukovaZalesskii2007}
therefore supplies, in degrees $n>1$,
the dense restriction embedding into the
product over the conjugacy classes of
order-two subgroups of $J$.
By (D.53f) its source has dimension one.
Each product factor is nonzero: for a
section $S$ of $J\to C_2$, its normalizer is
\[
                        C_C(S)\times S,
                                                               \tag{D.53g}
\]
and the restriction from the $S$ factor
already supplies one dimension.
Thus the product has finitely many
factors and the finite dense image is
the entire product.
Since $\operatorname{cd}_2(C_C(S))\leq1$,
the dimensions give the exact equality
\[
             1=\sum_{[S]}
                       \bigl(1+\dim H^1(C_C(S),\F_2)\bigr).
                                                               \tag{D.53h}
\]
There is a section $\langle\overline b\rangle$.
It follows that it is the unique
conjugacy class and that
$H^1(C_C(\overline b),\F_2)=0$.
These are also $C$-conjugacy classes:
an element of $J$ can be multiplied by
the appropriate element of a given
section to leave a conjugator in $C$.

Now let $B'$ be any subgroup of order
four in $E$. Its intersection with $F$
is trivial, so it is cyclic and $\pi|_{B'}$
is an isomorphism onto $C_4$.
Choose its unique generator $b'$ with
$\pi(b')=a$. Its square generates an
order-two section $T'$ of $E_2$.
Conjugation by $b'$ has the same action
on the section classes of $E_2$ as
conjugation by $b$, because
$b'b^{-1}\in F$.
The class of $T'$ is fixed, and is
therefore the unique class of $T$.
An $F$-conjugation puts $T'$ equal to $T$,
and then $B'\subseteq N_E(T)=N$.
Its image in $J=N/T$ has order two.
By (D.53h), a further conjugation by
an element of $C$ puts that image equal
to $\langle\overline b\rangle$.

The exact inverse image of this subgroup
under $N\to N/T$ has order four,
because its kernel is the specified
subgroup $T$ of order two.
It contains $\langle b\rangle$ and
hence equals $\langle b\rangle$.
Thus the conjugated $B'$ equals $B$.
All conjugations used lie in $F$ and
preserve $\pi$, so its labelled generator
$b'$ is sent to $b$, proving uniqueness
of the homomorphic section with labels.

The fixed group $C_C(\overline b)$ in
(D.53h) is exactly $C_F(b)$.
For an element of $C$ fixed by
$\overline b$, its commutator with $b$
lies in $T$ and also in $C$, and
$T\cap C=1$.
The converse is immediate, and an
element commuting with $b$ already
centralizes $T$.
This proves the cohomology statement
in (D.53a).
Finally, an element normalizing $B$
must fix its generator $b$: conjugation
preserves its label in the abelian
quotient $C_4$, and $\pi|_B$ is injective.
Multiplication by the matching power
of $b$ then leaves an element of
$C_F(b)$, proving the direct-product
normalizer in (D.53a).

The two orbits of the transposition
in (D.53d) are exactly the two
$E$-conjugacy classes of order-two
subgroups. Only its fixed class can
be the square of a subgroup of order
four, and $B$ realizes that containment.
This proves the final assertion.

The finite program
\path{code/part_c/torus_order4_sections_d53.py}
checks the mod-four matrix obstruction
and every invariant three-point affine
configuration above a nonidentity
order-two matrix in $\GL_2(\F_2)$.
The normalizer and restriction arguments
prove the full section result.
No conjugacy of the outer action $a$
to a prescribed geometric action is
asserted here.
\end{proof}

\begin{claim}{D.54}{Proved}{No fixed character for a cubic torus action}
Let $a$ have order three in the full
once-punctured-torus loop-outer group.
For every prime $q\ne3$, it fixes no nonzero
continuous character
\[
                         \chi:F\longrightarrow C_q,
                      \qquad F=\widehat F(x,y).         \tag{D.54a}
\]
Consequently the reduction of $A_q(a)$ modulo
$q$ has no eigenvalue one.
For every prime $p$, the endomorphism
$A_p(a)-I$ of $\Z_p^2$ is injective.
The full fixed group of every actual
order-three representative is therefore
contained in the full commutator subgroup
$F'$.

In particular, $A_2$ is injective on every
finite subgroup of the full torus group.
The allowed list of D.52 loses
$C_{12}$ and $D_{24}$.
At a nontrivial cubic multiplier
coordinate $q\ne3$, the remaining
homology matrix is exactly
\[
                         A_q(a)=\lambda(a)_q^2 I.
                                                               \tag{D.54b}
\]
This statement does not exclude the
remaining scalar cubic multiplier case.
\end{claim}
\passage{Section 30, folios 194-196, scans 383-387; Section 31, folios 198-199 verso, scans 391-394; Part D(b), full loop-outer finite actions in type (1,1)}

\begin{proof}
Use D.27 to choose an actual representative
$\alpha$ of order three and put
$E=F\rtimes_\alpha\langle s\rangle$,
where $s$ has quotient label $a$.
Write $c=[x,y]$ and $L=\langle c\rangle$.
D.48 supplies a labelled cusp-normalizer
element $v$ with
\[
       v\in N_E(L),\qquad v^3=c^\kappa,\qquad
       vcv^{-1}=c^\lambda,\qquad
                         \kappa_3\in\Z_3^\times.        \tag{D.54c}
\]
It also gives $\lambda_3=1$ and
\[
                  I+A_3(a)+A_3(a)^2=0
                         \quad\text{on }\Z_3^2.         \tag{D.54d}
\]

Suppose $\chi$ in (D.54a) were nonzero
and invariant. Its extension
\[
       \widetilde\chi:E\longrightarrow C_q,\qquad
              \widetilde\chi(hs^j)=\chi(h)
                                                               \tag{D.54e}
\]
is an actual continuous homomorphism.
Put
\[
               N=\ker\chi,\qquad
               E_N=\ker\widetilde\chi=N\rtimes\langle s\rangle .
                                                               \tag{D.54f}
\]
Both have their indicated index $q$
in $F$ and $E$. Since $\chi(c)=0$,
(D.54c) gives
$3\widetilde\chi(v)=0$. The prime $q$
differs from three, so $v\in E_N$.
Thus $v=ns$ for a specified $n\in N$.

The restriction of $\chi$ to the dense
discrete free group is onto $C_q$.
Its covering surface is the degree-$q$
cover of the punctured torus with
boundary monodromy zero.
It has $q$ boundary lifts of degree one,
Euler characteristic $-q$, and hence
type $(1,q)$. Its full fundamental-group
completion is exactly $N$, by the
finite-index completion map.
Choose a discrete word $f$ with
$\chi(f)=1$. Its boundary representatives
are
\[
                c_i=f^i c f^{-i}\quad(0\leq i<q).
                                                               \tag{D.54g}
\]
Their homology classes have exactly the
single relation $\sum_i[c_i]=0$.
In particular each $[c_i]$ is nonzero
in $H_1(N,\F_3)$, because $q\geq2$.
Equivalently, omitting any other boundary
from the usual genus-one presentation
puts $c_i$ into a free basis.

The action on every boundary class has
multiplier $\lambda$ and fixes its label.
This follows from the actual $v\in E_N$:
the element
$n_i=(vf^iv^{-1})f^{-i}$ belongs to $N$,
and
$vc_iv^{-1}=n_i c_i^\lambda n_i^{-1}$.
Conjugation by $s=n^{-1}v$ differs from
this by an inner automorphism of $N$.
Thus $\alpha$ fixes all $q$ boundary
homology classes over $\Z_3$.

We give the required splitting of
the entire homology module, including
its comparison with the base torus.
Let $D=F/N=C_q$ be the deck group and set
\[
          M=H_1(N,\Z_3),\qquad
                       e=\frac1q\sum_{d\in D}d .
                                                               \tag{D.54h}
\]
The action of $\alpha$ on $M$ commutes
with the deck action: it induces the
identity on $D$, and two lifts of a
deck transformation differ by an inner
automorphism of $N$ on this homology.
Consequently $e$ commutes with $\alpha$.

Filling the $q$ boundaries gives the
compact degree-$q$ torus cover.
Explicitly its lattice is the kernel
$\Lambda\leq\Z^2$ of the character
$\Z^2\to C_q$ induced by $\chi$,
and its map to the base torus is
$\R^2/\Lambda\to\R^2/\Z^2$.
The deck maps are translations and
act trivially on the filled homology.
The boundary submodule
\[
              B_\partial
                =\Z_3[D]/\Z_3\Bigl(\sum_{d\in D}d\Bigr)
                                                               \tag{D.54i}
\]
has no deck-invariant vector, because
$q$ is a unit in $\Z_3$.
The exact boundary sequence and (D.54h)
therefore give
\[
             (1-e)M=B_\partial,\qquad
             eM\cong H_1(\R^2/\Lambda,\Z_3)
                         \cong H_1(\R^2/\Z^2,\Z_3).
                                                               \tag{D.54j}
\]
The last isomorphism is induced by the
displayed covering: its lattice inclusion
has index $q$, a $3$-adic unit.
These maps intertwine the induced
automorphisms, so $\alpha|_{eM}$ is
exactly $A_3(a)$.
On $(1-e)M$ it is the identity by
(D.54g). Hence
\[
                  I+\alpha+\alpha^2=0
                       \quad\text{on }H_1(N,\F_3):
                                                               \tag{D.54k}
\]
on the first summand this is (D.54d),
and on the boundary summand the norm
is multiplication by three.

The cube of the actual element $ns$
now gives the contradiction:
\[
 \begin{split}
  [v^3]
      &=[n\,\alpha(n)\,\alpha^2(n)]\\
      &=(I+\alpha+\alpha^2)[n]=0
                       \quad\text{in }H_1(N,\F_3),\\
  [c^\kappa]&=(\kappa_3\bmod3)[c]\ne0 .
 \end{split}                                             \tag{D.54l}
\]
Equations (D.54c) and (D.54l) concern
the same two elements of $N$.
Thus no invariant character (D.54a)
exists.

For $q\ne3$, a matrix and its transpose
have the same rank, so a one-eigenvector
would supply a nonzero invariant
character. It has just been excluded.
In particular $A_q(a)-I$ is invertible
over $\Z_q$.
At $q=3$ its determinant is three,
by D.48, so it is still injective.
If $h$ is fixed by an actual representative,
its image in every $\Z_p^2$ is therefore
zero. The common kernel of the actual
abelianization map
$F\to\widehat\Z^2=\prod_p\Z_p^2$
is $F'$, proving the fixed-group assertion.

At a nontrivial cubic multiplier
coordinate the two eigenvalue multisets
were exactly
$\{1,\lambda_q\}$ and
$\{\lambda_q^2,\lambda_q^2\}$ in D.51.
The first is now excluded.
A finite-order matrix in characteristic
zero is diagonalizable, so the second
gives the actual scalar matrix
(D.54b).

Finally let $G$ be any finite torus
subgroup. D.33 gives $|G|\mid24$.
A nontrivial kernel of $A_2|_G$ would
contain a prime-order element.
Its prime can only be two or three.
D.33 excludes an element of order two,
and (D.54a), with $q=2$, excludes
an element of order three acting
trivially on two-adic homology.
Thus $A_2|_G$ is injective.
The faithful-image part (D.52h), together
with D.50, now excludes $C_{12}$ and
$D_{24}$ from that allowed list.

The program
\path{code/part_c/torus_cover_norm_d54_d55.py}
builds the cyclic covering homology,
its individual boundary vectors and
the deck-average splitting for several
primes $q\ne3$, and verifies the forced
norm contradiction over $\F_3$.
The displayed maps prove it for every
such prime.
\end{proof}

\begin{claim}{D.55}{Proved}{Full primary faithfulness and the nine finite torus types}
Every positive involution $b$ in the full
once-punctured-torus loop-outer group has
\[
                      A_p(b)=-I
                           \quad\text{for every prime }p.
                                                               \tag{D.55a}
\]
Consequently, for each prime $p$, the actual
homology map $A_p$ is injective on every
finite subgroup of that group.
The possible finite abstract subgroup
types are exactly the following nine,
all realized:
\[
       1,\ C_2,\ C_3,\ C_4,\ C_2\times C_2,\ C_6,\
                                      S_3,\ D_8,\ D_{12}.
                                                               \tag{D.55b}
\]
The groups $D_8,D_{12}$ have orders eight
and twelve respectively.

For an arbitrary order-four element $a$,
\[
                A_p(a)^2=-I\quad\text{for every }p.
                                                               \tag{D.55c}
\]
At a negative multiplier coordinate $q$
it is necessarily a scalar matrix
$i_q I$ or $-i_q I$, where $i_q^2=-1$
in $\Q_q$; thus $q\equiv1\pmod4$.
Together with (D.54b), this retains
the precise scalar cubic and quartic
multiplier cases. The abstract
classification (D.55b) does not
settle those scalar cases or the
conjugacy of arbitrary finite subgroups.
\end{claim}
\passage{Section 30, folios 194-196, scans 383-387; Section 31, folios 198-199 verso, scans 391-394; Part D(b), full loop-outer finite actions in type (1,1)}

\begin{proof}
\emph{The positive involution's full cusp normalizer.}
Let $b$ be any positive involution.
D.43 gives its global multiplier $1$.
D.27 gives an actual representative
$\beta$ of order two, and put
$E_b=F\rtimes_\beta\langle s\rangle$.
The first half of D.46 applies to
every such positive involution; it
does not require $b$ to have first
been exhibited as an order-four square.
It gives
\[
        A_2(b)=-I,\qquad
        H^1(C_F(S),\F_2)=0
          \quad\text{for every order-two section }S.
                                                               \tag{D.55d}
\]

We recall the prime-independent part
of the closed-subgroup argument in
D.48. If $C\leq F$ is closed and
$H^1(C,\F_p)=0$, it cannot contain
the full subgroup $c^{e_p\Z_p}$ of
the torus cusp, where $c=[x,y]$.
If the $p$-part of $[F:C]$ were
finite, choose an open neighborhood
$W\geq C$ with maximal such part.
All subsequent indices are prime
to $p$, so transfer makes the
restriction maps on first cohomology
injective. Their direct limit is
$H^1(C,\F_p)$, contradicting the
nonzero first cohomology of the open
full free group $W$.
If the $p$-part is infinite, choose
an open $W\geq C$ of index divisible
by $p$. Its selected cusp lift has
degree prime to $p$, because it
contains the full cusp $\Z_p$.
There must be at least two boundary
lifts: a unique lift would have
degree $[F:W]$.
A boundary character of this finite
cover then restricts nontrivially
to that cusp $\Z_p$ and hence to
$C$, again a contradiction.
These are the actual open-subgroup
restriction and finite-cover maps
used in D.48, with an arbitrary
prime $p$ in place of three.

Apply this with $p=2$ to (D.55d).
An involution in $N_{E_b}(\langle c\rangle)$
would be a section and would centralize
the entire cusp, since the global
multiplier is $1$.
Its full fixed group would therefore
contain the cusp $\Z_2$, which has
just been excluded.
Thus this cusp normalizer is
torsion-free.
Choose its labelled element $v$:
\[
              vcv^{-1}=c,\qquad v^2=c^\kappa.
                                                               \tag{D.55e}
\]
If $\kappa_2$ were even, the equation
$2z=-\kappa$ would be soluble in
$\widehat\Z$: at two by that
divisibility, and at every odd prime
because two is a unit.
Then $(vc^z)^2=1$, and its nontrivial
quotient label would make it an
involution in the cusp normalizer.
Consequently
\[
                             \kappa_2\in\Z_2^\times.
                                                               \tag{D.55f}
\]
The use of the global multiplier
$1$ in (D.55e) is essential to
this exact norm equation.

\emph{The primitive boundary rules out every odd fixed character.}
Let $q$ be odd and suppose a nonzero
$\beta$-invariant character
$\chi:F\to C_q$ existed.
Extend it by $\widetilde\chi(s)=0$
to $E_b$, and put
$N=\ker\chi$ and
$E_N=\ker\widetilde\chi=N\rtimes\langle s\rangle$.
Equation (D.55e) gives
$2\widetilde\chi(v)=0$, so $v\in E_N$.
Write $v=ns$ with $n\in N$.

The exact cover and its maps are
those of (D.54g)--(D.54j), now over
$\Z_2$. It has type $(1,q)$.
The deck-average idempotent
$e=q^{-1}\sum_{C_q}d$ commutes
with $\beta$, identifies $eH_1(N,\Z_2)$
with the base torus homology by the
index-$q$ lattice inclusion, and
has boundary augmentation as its
complement.
On the first summand $\beta=-I$
by (D.55d). On the second it is
the identity: the actual $v\in E_N$
fixes each cusp label and centralizes
its cusp generator up to $N$-conjugacy,
exactly as in (D.54g).
Hence
\[
                  I+\beta=0
                       \quad\text{on }H_1(N,\F_2).
                                                               \tag{D.55g}
\]
The same two elements in (D.55e)
then have incompatible homology:
\[
             [v^2]=[n\,\beta(n)]=0,\qquad
             [c^\kappa]=(\kappa_2\bmod2)[c]\ne0.
                                                               \tag{D.55h}
\]
The second class is nonzero because
$c$ is a primitive boundary in
the cover with $q\geq3$ cusps.
This excludes the invariant character.

At every odd prime $q$ the matrix
$A_q(b)$ squares to $I$ and has
determinant $1$.
Its characteristic-zero eigenvalues
are both $1$ or both $-1$.
The first possibility makes it $I$
and supplies a nonzero invariant
character modulo $q$, contrary
to (D.55h).
The second makes it $-I$.
Together with (D.55d), this proves
(D.55a).

\emph{Faithfulness at every prime and the exact abstract list.}
Fix a prime $p$ and a finite
subgroup $G$ of the full torus group.
By D.33 its order divides twenty-four.
A nontrivial kernel of $A_p|_G$
would contain an element of order
two or three.
A positive involution has image
$-I$ by (D.55a).
A negative involution has determinant
$-1$ at every prime by D.43, so
also cannot lie in the kernel.
An order-three element has nonidentity
image at $p=3$ by D.33, and at
every other prime by D.54's
fixed-character obstruction.
The kernel is therefore trivial,
as claimed.

D.54 already removed $C_{12}$
and $D_{24}$ from D.52.
The only additional type in that
allowed list was $\operatorname{Dic}_{12}$.
Its normal subgroup of order three
would be faithfully represented over
$\Q_3$, and an order-four element
inverting it would also have order
four by the just-proved faithfulness.
But the exact field-normalizer
calculation (D.52i)--(D.52k) proves
that every finite-order matrix
inverting this $C_3$ has square $I$:
the alternative would require norm
$-1$ in $\Q_3(\zeta_3)$.
This contradiction excludes
$\operatorname{Dic}_{12}$.
The nine remaining types all have
the explicit realizations of
D.50 and D.52, proving (D.55b).

For an arbitrary order-four $a$,
its square is a positive involution,
so (D.55a) gives (D.55c).
The full multiplier still has square
one by D.50.
At a negative coordinate the two
eigenvalues of $A_q(a)$ are roots
of $X^2+1$ whose product is $-1$,
so they are equal.
The matrix is diagonalizable in
characteristic zero, hence scalar
$i_q I$ or $-i_q I$.
The two-coordinate is positive by
D.50. For an odd $q$ a square
root of $-1$ in $\Q_q$ reduces to
one in $\F_q$, which requires
$q\equiv1\pmod4$.
This proves the stated quartic
restriction without excluding its
remaining scalar case.

The program
\path{code/part_c/torus_cover_norm_d54_d55.py}
also checks the square norm and
primitive-boundary comparison over
$\F_2$ in explicit cyclic odd-degree
covering homology modules.
The full subgroup, covering and
normalizer arguments above give
the all-prime conclusion.
\end{proof}

\begin{claim}{D.56}{Proved}{The prime sections of a cyclic-six torus action}
Let $a$ have order six in the full once-punctured-torus
loop-outer group of $F=\widehat F(x,y)$, and let
\[
 E=\{(\phi,j)\in\operatorname{Aut}(F)\times C_6:
                                      [\phi]=a^j\}
                                                               \tag{D.56a}
\]
be its centre-free extension, with kernel identified
with $F$ by the inner-automorphism map.
There are exactly three $F$-conjugacy classes of
order-two sections over $\langle a^3\rangle$ and
exactly two $F$-conjugacy classes of order-three
sections over $\langle a^2\rangle$.
Conjugation by an element with quotient label $a^2$
cycles the three former classes; conjugation by an
element with quotient label $a^3$ exchanges the two
latter classes.

Consequently $E$ has exactly one conjugacy class
of subgroups of order two and one of order three,
and these are its only nontrivial finite subgroup
types.  In particular the full extension (D.56a)
does not split, and $a$ has no finite-order
representative in $\operatorname{Aut}(F)$.
For each prime section $S$ of order $p=2$ or $3$,
\[
                 N_E(S)=C_F(S)\times S,
                 \qquad H^1(C_F(S),\F_p)=0.             \tag{D.56b}
\]
The assertions concern these specified extension
and section maps; they do not identify $E$ with
a free profinite product.
\end{claim}
\passage{Section 30, folios 194-196, scans 383-387; Section 31, folios 198-199 verso, scans 391-394; Part D(b), finite actions and actual lifts in type (1,1)}

\begin{proof}
Put
\[
 E_2=\pi^{-1}(\langle a^3\rangle),\qquad
 E_3=\pi^{-1}(\langle a^2\rangle),
 \quad \pi:E\longrightarrow\langle a\rangle .          \tag{D.56c}
\]
The involution $a^3$ is positive.
Indeed D.55 embeds $\langle a\rangle$ faithfully
in $\operatorname{GL}_2(\Q_2)$.
An order-six matrix there has its order-three
part in the irreducible quadratic field
$\Q_2(\zeta_3)$; its commuting involution is
the nonidentity scalar $-I$.
For completeness, the centralizer is that field
because the cubic part has irreducible polynomial
$X^2+X+1$, and its only elements squaring to one
are $1,-1$. Faithfulness excludes the former
for $a^3$.
Thus $a^3$ has determinant one at two, and D.43
gives its global positive sign.
D.27 supplies the actual prime-order sections
of both groups in (D.56c).

\emph{The three involution classes and their affine coordinates.}
Choose an order-two section generator $t$ in
$E_2$ and identify $E_2=F\rtimes\langle t\rangle$.
By D.46 and D.55, its action on two-adic
homology is $-I$, and its maximal elementary
abelian two-quotient has coordinates
\[
 (u,\chi_1,\chi_2):E_2\longrightarrow\F_2^3,
 \quad u(t)=1,\quad\chi_1(t)=\chi_2(t)=0,                \tag{D.56d}
\]
where $\chi_1|_F,\chi_2|_F$ form a basis of
$H^1(F,\F_2)$.
Every section generator has $u=1$, so its
coordinate is a point of the affine plane
$\F_2^2$.
The three $F$-conjugacy classes of D.46 give
three distinct such points.
One can read the distinctness directly from
that proof: Brown restriction identifies the
localized degree-zero ring with $\F_2^3$,
and this ring is generated by $\chi_1/u$ and
$\chi_2/u$. Equal values on two section
classes would prevent these two elements
from generating the corresponding two
coordinate factors.

Conjugation by an element $e\in E$ with
$\pi(e)=a^2$ preserves $u$ and acts on the
plane (D.56d) by
\[
                           r\longmapsto Rr+d,           \tag{D.56e}
\]
where the linear part is the reduction of
$A_2(a^2)$ in the chosen homology coordinates.
D.54, at $q=2$, says that $R-I$ is invertible.
Also $R^3=I$, so $R$ is a nonidentity cubic
matrix over $\F_2$ and $I+R+R^2=0$.
The affine map (D.56e) has one fixed point,
and its other three points form a single
orbit of length three: translate its unique
fixed point to zero and use the corresponding
linear statement.
Its invariant three-point set of section
classes must be that orbit.
This proves the asserted cyclic permutation.
The action on the class set factors through
$E/E_2=C_3$, because conjugation by $E_2$
is already $F$-conjugation on its sections,
as proved in D.46.

\emph{The two cubic classes and their affine coordinates.}
Choose an order-three section generator $s$
in $E_3$ with label $a^2$.
Use the actual abelian quotient of D.48 and
the separation proved in D.51:
\[
 (u_3,\chi):E_3\longrightarrow\F_3^2,
 \qquad u_3(s)=1,\quad\chi(s)=0.                        \tag{D.56f}
\]
The two $F$-conjugacy classes of section
generators with prescribed label $u_3=1$
have $\chi$-values $0$ and $b$, with
$b\ne0$.
For an element $e\in E$ with label $a^3$,
conjugation preserves $u_3$ because $E/F$
is abelian.  On the restriction of $\chi$
to $F$ it acts by $-1$, since D.55 gives
$A_3(a^3)=-I$.
The difference between the conjugated
character and $-\chi$ vanishes on $F$,
so is a scalar multiple of $u_3$.
On the line $u_3=1$ the action consequently
has the exact form
\[
                              r\longmapsto-r+d'.        \tag{D.56g}
\]
Preserving the two-point set $\{0,b\}$
forces $d'=b$: the alternative of fixing
both would give $2b=0$ in $\F_3$.
Thus this action exchanges the two classes.
Conjugation by $E_3$ is $F$-conjugation on
its sections, by D.48, so the action factors
through $E/E_3=C_2$.

\emph{The full finite subgroups and the failure of lifting.}
Every finite subgroup of $E$ injects into
$E/F=C_6$, because $F$ is torsion-free.
It is therefore cyclic of order one, two,
three or six.  A subgroup of order six
would normalize its order-two subgroup.
Its element with quotient label $a^2$
would then fix that subgroup's
$F$-conjugacy class, contradicting
(D.56e).  Thus no subgroup of order six
exists.
The two affine actions show respectively
that all order-two subgroups are conjugate
in $E$ and all order-three subgroups are
conjugate in $E$.

For a chosen section $S$ of order two,
the stabilizer of its $F$-conjugacy class
in $E/F$ is exactly $\langle a^3\rangle$;
for order three it is exactly
$\langle a^2\rangle$.
The image of $N_E(S)$ equals this
stabilizer: a conjugator whose class is
fixed can be multiplied by an element
of $F$ to normalize $S$ itself.
Thus the full normalizer is already
its normalizer in $E_p$.
Equations (D.46h) and (D.48g) identify
it as $C_F(S)\times S$ and give the
cohomology assertion in (D.56b).

The first projection of (D.56a) is
injective and identifies $E$ with the
inverse image of $\langle a\rangle$
in $\operatorname{Aut}(F)$.
Indeed the outer class of the first
coordinate determines its label $j$.
A homomorphic lift of $\langle a\rangle$
would give a subgroup of order six in
this exact inverse image, which has
been excluded.
More generally, a finite-order
representative $\phi$ of $a$ would
satisfy $\phi^6\in\operatorname{Inn}(F)$.
The latter group is isomorphic to the
torsion-free group $F$, so
$\phi^6=1$, making $\phi$ a forbidden
element of order six.
This proves the labelled nonlifting
assertion while preserving the actual
prime-order lifts in (D.56c).

The program
\path{code/part_a/torus_order6_sections_d56.py}
checks every affine cubic action on
$\F_2^2$ with nonidentity cubic linear
part, both nonzero cubic section
separations in (D.56f), and the
corresponding quotient stabilizers.
The extension and cohomological
arguments above establish that these
are the actions on the full section
class sets.
\end{proof}

\begin{claim}{D.57}{Proved}{Full lifts of the noncyclic torus two-groups}
Let $Q$ be a subgroup of the full
once-punctured-torus loop-outer group,
isomorphic to $C_2\times C_2$ or to
$D_8$ of order eight.  Put
\[
 E_Q=\{(\phi,q)\in\operatorname{Aut}(F)\times Q:
                        [\phi]=q\},\qquad F=\widehat F(x,y).
                                                               \tag{D.57a}
\]
The projection $E_Q\to Q$ has a homomorphic
section with exactly its specified outer labels.
For $Q=D_8$ there is exactly one
$F$-conjugacy class of such sections.

For $Q=C_2\times C_2$, let $z$ be its
positive involution and let $r$ be either
of its two negative involutions.  There
are exactly three section classes when
\[
                         A_2(r)\bmod2=I,                \tag{D.57b}
\]
and exactly one otherwise.  This choice
does not depend on which negative
involution is called $r$.
For the image $H$ of every full section,
\[
                  N_{E_Q}(H)=C_F(H)\times H,
                  \qquad H^1(C_F(H),\F_2)=0.            \tag{D.57c}
\]
Together with D.27 and D.46, this supplies
actual lifts for every finite two-subgroup
in the list of D.50.  It is a statement
about labelled automorphism lifts, not
a geometric conjugacy classification.
\end{claim}
\passage{Section 30, folios 194-196, scans 383-387; Section 31, folios 198-199 verso, scans 391-394; Part D(b), finite two-groups and actual labelled lifts in type (1,1)}

\begin{proof}
\emph{The invariant prime-section classes for a four-group.}
Suppose first that $Q=C_2\times C_2$.
By D.50 and D.55 its faithful dyadic
representation has exactly one nonidentity
positive element $z$, with $A_2(z)=-I$;
the two other involutions are negative.
In the inverse image $E_z$ of
$\langle z\rangle$, D.46 supplies exactly
three $F$-conjugacy classes of sections
$T$ of order two, each with
\[
                         H^1(C_F(T),\F_2)=0.             \tag{D.57d}
\]
As in (D.56d), their elementary abelian
coordinates are three distinct points
of the affine plane $\F_2^2$.

Conjugation by an element of $E_Q$
with label $r$ induces an affine map
\[
                           v\longmapsto Mv+d,
                    \qquad M=A_2(r)\bmod2.              \tag{D.57e}
\]
Its square is the identity on the plane:
the square of the conjugating element
lies in $F$, and inner automorphisms
act trivially on the elementary abelian
quotient of $E_z$.
Its three-point set of section classes
is invariant.

If $M=I$, a nonzero translation would
have only two-element orbits, so could
not preserve a three-point set.
Thus $d=0$, and all three classes are
fixed. If $M\ne I$, it is the unique
two-dimensional unipotent type over
$\F_2$, with
\[
          \ker(I+M)=\operatorname{im}(I+M),
                         \quad\dim\ker(I+M)=1.
\]
The square identity gives $(I+M)d=0$,
so (D.57e) has two fixed points and
one orbit of size two.
An invariant three-point subset must
contain that orbit and one fixed
point. Thus exactly one section
class is fixed in this case.
Since $A_2(zr)=-A_2(r)$, the reductions
of the two negative involutions are
equal, proving the independence in
(D.57b).

Choose any fixed class and one of
its section groups $T$.
Its normalizer $N=N_{E_Q}(T)$ maps
onto $Q$: adjust a conjugator with
label $r$ by an element of $F$ to
normalize $T$, and use $T$ itself
for the label $z$.
Its kernel is $C=C_F(T)$, since an
element of $F$ normalizing the
order-two group $T$ centralizes it.
Moreover $T$ is central in $N$.
The quotient $J=N/T$ therefore has
the exact sequence
\[
                  1\longrightarrow C\longrightarrow J
                     \longrightarrow Q/\langle z\rangle
                     \longrightarrow1,
                  \qquad Q/\langle z\rangle=C_2.        \tag{D.57f}
\]
The inclusion of $C$ in this quotient
is injective, because $C\cap T=1$.

\emph{The invariant order-four class for a dihedral group.}
Now suppose that $Q=D_8$ and let
$R\cong C_4$ be its rotation subgroup.
By D.53 the inverse image $E_R$ has
exactly one $F$-conjugacy class of
labelled order-four sections $B$,
and
\[
                         H^1(C_F(B),\F_2)=0.             \tag{D.57g}
\]
Every element of $E_Q$ preserves this
class, since $R$ is normal in $Q$.
The normalizer $N=N_{E_Q}(B)$ thus
maps onto $Q$, by the same adjustment
of a conjugator by an element of $F$.
Its kernel is $C=C_F(B)$: an element
of $F$ normalizing $B$ fixes the label
of every element of $B$, and therefore
centralizes $B$.
The subgroup $B$ is normal in $N$,
and the quotient $J=N/B$ has
\[
                  1\longrightarrow C\longrightarrow J
                        \longrightarrow Q/R\longrightarrow1,
                              \qquad Q/R=C_2.           \tag{D.57h}
\]
Again $C$ embeds, because $C\cap B=1$.
Conjugation on $B$ in this construction
is the actual conjugation on its
specified rotation labels in $D_8$.

\emph{Torsion and exact finite preimages in the two constructed quotients.}
In each of (D.57f) and (D.57h), the
group $C$ is a closed subgroup of
$F$, and hence is torsion-free with
$\operatorname{cd}_p(C)\leq1$ at every
prime. Equations (D.57d) and (D.57g)
give $C^{(2)}=1$.
Every map from $J$ to a finite two-group
therefore kills $C$, and the displayed
nontrivial quotient gives
\[
                              J^{(2)}=C_2.              \tag{D.57i}
\]
The same full cohomological-dimension
argument used in D.27 and D.46 proves
that $J$ has torsion. Indeed, a
torsion-free $J$ would be projective,
since $C$ is open and has the stated
cohomological dimensions. Its maximal
pro-two quotient would then be a free
pro-two group, contrary to (D.57i).
Every finite subgroup of $J$ injects
into $J/C=C_2$. Consequently $J$ has
an involution $j$ outside $C$.

Take the full inverse image $H$ of
$\langle j\rangle$ in $N$.
In the four-group construction its
kernel is $T$, so $|H|=4$; in the
dihedral construction its kernel is
$B$, so $|H|=8$.
In either case $H\cap F=1$.
To check this in the actual maps,
an element of $H\cap F$ has image
in $\langle j\rangle\cap C=1$ in
$J$, and then lies in $T\cap F=1$
or in $B\cap F=1$ respectively.
Thus the projection $H\to Q$ is
injective. Its order equals $|Q|$,
so it is an isomorphism with the
prescribed labels.
Its inverse is the required section.
This exact injection also fixes the
isomorphism type of the finite lift;
no alternative central extension
is substituted for $Q$.

\emph{Uniqueness above each invariant class.}
The section just constructed identifies
$N=C\rtimes H$ and
\[
                   J=C\rtimes\langle j\rangle .        \tag{D.57j}
\]
The Hochschild--Serre sequence, with
$H^1(C,\F_2)=0$ and
$\operatorname{cd}_2(C)\leq1$, gives
the actual inflation isomorphism
\[
                   H^*(J,\F_2)\cong H^*(C_2,\F_2).
                                                               \tag{D.57k}
\]
The full Brown restriction theorem
\cite[Theorem~2 and \S1.1]{A-KochloukovaZalesskii2007}
applies: $C$ is open of finite
two-cohomological dimension and no
finite subgroup of $J$ can be a
four-group, since it injects into
$C_2$.
In high degree the source of its
dense restriction map is the finite
one-dimensional space in (D.57k).
Restriction to the displayed section
in (D.57j) is already nonzero in
every degree, so this map is injective.
Its image is therefore closed as
well as dense, so the restriction
is an isomorphism onto the product
of the section-normalizer cohomologies.
For a section $S\leq J$ its normalizer
is $C_C(S)\times S$, by the quotient
labels. Each such factor contributes
\[
                          1+\dim H^1(C_C(S),\F_2)
\]
in every sufficiently high degree.
There is consequently exactly one
conjugacy class of sections in $J$,
and its fixed group has zero first
cohomology over $\F_2$.
These are also exactly the
$C$-conjugacy classes: multiply a
given conjugator by the section
element with its quotient label.

Full sections of $E_Q$ containing
the chosen $T$ or $B$ lie in $N$,
and are exactly the full inverse
images of the order-two sections
of $J$. This follows from their
orders and the same kernel calculation
used above.
They therefore form one $C$-conjugacy
class. Every full four-group section
has a positive subgroup $T$ in one
of the invariant classes of
(D.57e); distinct such classes
cannot be identified by $F$-conjugacy.
The numbers of full section classes
are thus exactly three and one
as asserted. Every full dihedral
section has a rotation subgroup
in the unique class of D.53, so
the dihedral count is exactly one.

The fixed group $C_C(j)$ in the
Brown calculation is precisely
$C_F(H)$ in the original extension.
For example, commuting with $j$
modulo $T$ or $B$ means that a
commutator lies in that finite
subgroup. The commutator also lies
in $F$, so is trivial. The elements
of $C$ already commute with $T$
or $B$. This proves the cohomology
part of (D.57c).
Finally an element normalizing $H$
can be multiplied by the inverse
of the unique element of $H$ with
its outer label. The result lies
in $F$ and normalizes $H$, and
therefore centralizes each element
of $H$ by preservation of labels.
This gives the exact direct product
in (D.57c).

The program
\path{code/part_a/torus_two_group_lifts_d57.py}
checks all affine involutions of
$\F_2^2$ together with their invariant
three-point sets and verifies the
three-versus-one fixed-class count.
The full projectivity, Brown and
finite-preimage arguments above
provide the labelled lifts and
their uniqueness statements.
\end{proof}

\begin{claim}{D.58}{Proved}{Negative torus involution sections and cusp fusion}
Let $b$ be a negative involution of the full
once-punctured-torus loop-outer group of
$F=\widehat F(x,y)$, put $c=[x,y]$, and let
$E_b$ be its centre-free extension.
Write $A_2$ for its action on $H_1(F,\Z_2)$.
The $F$-conjugacy classes of order-two sections and
their full fixed-subgroup cohomology are exactly
\[
\begin{array}{c|c|c}
 A_2\bmod2&
 \text{number of section classes}&
 \{\dim_{\F_2}H^1(C_F(S),\F_2)\}_{[S]}\\ \hline
 I&2&\{1,0\}\\
 \ne I&1&\{0\}.
\end{array}                                             \tag{D.58a}
\]
For every such section,
$N_{E_b}(S)=C_F(S)\times S$.
The cusp normalizer is the full profinite dihedral
group
\[
 N_{E_b}(\overline{\langle c\rangle})
       =\overline{\langle c\rangle}
                       \rtimes\langle v\rangle,\qquad
 v^2=1,\quad vcv^{-1}=c^{-1}.                            \tag{D.58b}
\]
Its two internal conjugacy classes of order-two
sections become the same $F$-conjugacy class in $E_b$.
In particular, for every choice of the cusp reflection
$v$, there exists $g\in F$ with
\[
                  gvg^{-1}=cv,\qquad
                  c=[g,v]=g\,(vgv^{-1})^{-1}.            \tag{D.58c}
\]
The dimensions in (D.58a) describe continuous cohomology
of the full closed fixed subgroups.
\end{claim}
\passage{Section 30, folios 194-196, scans 383-387; Section 31, folios 198-199 verso, scans 391-394; Part D(b), full loop-outer finite actions in type (1,1)}

\begin{proof}
\emph{The actual extension and the two integral module cases.}
D.27 supplies an actual order-two representative
$\alpha$ and the exact identification
$E_b=F\rtimes_\alpha\langle s\rangle$.
D.43 gives its global multiplier $\lambda=-1$,
so $A_2^2=I$ and $\det A_2=-1$.
Let $u:E_b\to\F_2$ be the quotient character,
$u(s)=1$.

First suppose $A_2\equiv I\pmod2$.
On the free $\Z_2$-module $H^1(F,\Z_2)$, the induced
involution $D$ also satisfies $D\equiv I\pmod2$.
The endomorphism
\[
                         P=(I-D)/2                     \tag{D.58d}
\]
is integral and satisfies $P^2=P$.
Its image and kernel are direct summands of rank one,
because $D$ has the two eigenvalues $1,-1$ over $\Q_2$.
Thus choose continuous characters
$\varphi_+,\varphi_-:F\to\Z_2$ which form a basis and
satisfy
\[
        \varphi_+\alpha=\varphi_+,\qquad
        \varphi_-\alpha=-\varphi_- .                    \tag{D.58e}
\]
This is a decomposition of the specified cohomology
module. All subsequent maps use these actual
characters; no free generating pair is changed.

Their mod-$2$ reductions extend to characters
$a,b:E_b\to\F_2$ by $a(s)=b(s)=0$.
The Hochschild--Serre spectral sequence has the two
rows $\F_2u^i$ and $\F_2^2u^i$.
The degree-one characters extend, so their
transgression vanishes.
All other possible differentials are its products
with powers of $u$ and also vanish.
The associated graded pieces give, for every $n\geq1$,
\[
 H^n(E_b,\F_2)
     =\langle u^n,u^{n-1}a,u^{n-1}b\rangle_{\F_2}.
                                                        \tag{D.58f}
\]
Multiplication by $u$ is an isomorphism in degrees
at least one, since it is an isomorphism on both
associated graded pieces.

\emph{The two Bocksteins and the actual localized ring.}
Reduce (D.58e) modulo four and use the continuous
one-cochains
\[
 \widetilde a(hs^\epsilon)=\varphi_+(h)\bmod4,\qquad
 \widetilde b(hs^\epsilon)=\varphi_-(h)\bmod4.
\]
The first is a homomorphism, while the second has
coboundary
\[
 \delta\widetilde b(hs^\epsilon,ks^\eta)
        =2\epsilon\varphi_-(k)\pmod4.                   \tag{D.58g}
\]
For the coefficient sequence
$0\to\F_2\xrightarrow{2}\Z/4\to\F_2\to0$,
the Bockstein of a degree-one character is its
cup square. This follows directly by taking its
$\{0,1\}$-valued lift; changing the lift changes
the divided coboundary by a coboundary.
Thus (D.58g) gives
\[
                         a^2=0,\qquad b^2=ub.           \tag{D.58h}
\]

Write $ab=\epsilon ua+\delta ub+\gamma u^2$ using
(D.58f), with coefficients in $\F_2$.
Restriction to the chosen section $\langle s\rangle$
gives $\gamma=0$.
Multiplication by $a$ and (D.58h) then give
\[
           0=a^2b
             =\delta\epsilon\,u^2a+\delta^2u^2b .
\]
The two displayed degree-three classes are independent
by (D.58f), so $\delta=0$.
Consequently the degree-zero localization
$R=(H^*(E_b,\F_2)[u^{-1}])^0$ has the exact presentation
\[
 R=\F_2[A,B]\big/
       (A^2,\ B^2-B,\ AB-\epsilon A),\qquad
 A=a/u,\quad B=b/u,\quad \dim_{\F_2}R=3.                 \tag{D.58i}
\]
For either value of $\epsilon$ the map
\[
 \begin{split}
 R&\longrightarrow\F_2[\nu]/(\nu^2)\times\F_2,\\
 1&\longmapsto(1,1),\\
 A&\longmapsto(\nu,0),\\
 B&\longmapsto(\epsilon,1-\epsilon)
 \end{split}                                           \tag{D.58j}
\]
is a ring isomorphism.
It preserves the stated relations and sends the
three-dimensional basis $1,A,B$ to a basis.
In particular $R$ has exactly two primitive
idempotents, and its nilpotent ideal has dimension one.

\emph{The nonidentity action modulo two.}
Now suppose $A_2\bmod2\ne I$.
The dual two-dimensional module
$M=H^1(F,\F_2)$ is the regular $\F_2[C_2]$-module:
$D-I=D+I$ has rank one, with image equal to kernel.
The periodic cyclic resolution therefore gives
\[
       \dim M^{C_2}=1,\qquad H^i(C_2,M)=0\quad(i>0).
\]
Choose the nonzero invariant character on $F$ and
extend it to $\chi:E_b\to\F_2$ with $\chi(s)=0$.
This again makes the only possible transgression zero.
The two-row spectral sequence now gives
\[
 H^1(E_b,\F_2)=\langle u,\chi\rangle,\qquad
 H^n(E_b,\F_2)=\F_2u^n\quad(n\geq2).                    \tag{D.58k}
\]
Both $u\chi$ and $\chi^2$ restrict to zero on
$\langle s\rangle$, whereas $u^2$ restricts nontrivially.
Thus $u\chi=\chi^2=0$.
In this case the degree-zero localization is
\[
             (H^*(E_b,\F_2)[u^{-1}])^0=\F_2.             \tag{D.58l}
\]

\emph{Full Brown restriction and the section labels.}
Every finite subgroup of $E_b$ injects into $E_b/F=C_2$,
because $F$ is torsion-free.
In particular $E_b$ has no subgroup $C_2\times C_2$.
For a section $S$, conjugation in its normalizer
induces the identity on $S$.
Multiplying a normalizer element by the section
element with its quotient label leaves an element
of $C_S=C_F(S)$. This proves the exact direct product
\[
                         N_{E_b}(S)=C_S\times S.        \tag{D.58m}
\]
The closed subgroup $C_S\leq F$ has
$\operatorname{cd}_2(C_S)\leq1$.

The full-profinite Brown restriction theorem applies
with the open subgroup $F$ of cohomological
$2$-dimension one
\cite[Theorem~2 and \S1.1]{A-KochloukovaZalesskii2007}.
For $n>1$ its restriction map to the product of the
cohomologies of (D.58m) is a dense embedding.
The source is finite by (D.58f) or (D.58k);
therefore that map is an isomorphism onto a finite
product, as in the cited proof.
These are the actual restriction maps and preserve
cup products. Since $u$ restricts trivially on
$C_S$ and nontrivially on $S$, localization gives
the ring isomorphism
\[
 R\ \xrightarrow{\ \operatorname{res}\ }\
   \prod_{[S]}
      \left(\F_2\oplus H^1(C_S,\F_2)u^{-1}\right).
                                                        \tag{D.58n}
\]
The second summand in each factor has square zero,
because $H^2(C_S,\F_2)=0$.
Each factor is local, with that summand as its
nilpotent ideal and with residue field $\F_2$.

If $A_2\bmod2\ne I$, (D.58l) and (D.58n) give
one factor and a zero second summand.
If $A_2\bmod2=I$, (D.58j) gives two factors,
one with a one-dimensional second summand and one
with a zero second summand.
This proves all the numbers in (D.58a).
The classes in this comparison are also precisely
the $F$-conjugacy classes:
for a conjugator $e\in E_b$, multiply it by the
inverse of the unique element of $S$ having the
same quotient label. This replaces $e$ by an
element of $F$ without changing $eSe^{-1}$.

In the identity-modulo-two case the two classes
are distinguished by the value of $b$ on their
nonidentity section element.
Indeed the two residue-field evaluations of
(D.58j) give $B=0$ and $B=1$.
Restriction to a section is exactly that
residue-field evaluation in (D.58n).
Also $a$ vanishes on every section element, since
$a^2=0$ whereas the square of the nonzero
character of $C_2$ is nonzero.

\emph{The exact cusp reflections and their fusion.}
Put $L=\overline{\langle c\rangle}\cong\widehat\Z$.
Preservation of the cusp class and A.7 give
\[
       1\longrightarrow L\longrightarrow N_{E_b}(L)
                     \longrightarrow C_2\longrightarrow1.
                                                        \tag{D.58o}
\]
Choose $v$ over its nonidentity label.
The global multiplier $-1$ gives $vcv^{-1}=c^{-1}$.
Write $v^2=c^\kappa$.
Commutation of $v$ with its own square implies
$c^\kappa=c^{-\kappa}$, hence $2\kappa=0$ in
$\widehat\Z$ and $\kappa=0$.
This proves (D.58b), with the multiplication and
specified label of its finite factor retained.

Every section of this dihedral group is
$\langle c^t v\rangle$, $t\in\widehat\Z$.
Conjugation by $c^z$ changes $t$ to $t+2z$,
and conjugation by $v$ changes it to $-t$.
Its two internal classes are therefore represented
by $v$ and $cv$.
These elements have the same image under every
abelian character of $E_b$, since $c=[x,y]\in[F,F]$.
When there is only one full section class they
are already $F$-conjugate.
When there are two, their equal $b$-values and
the preceding residue-field distinction again
put them in the same $F$-class.
An $F$-conjugator sends the unique nonidentity
element to the unique nonidentity element, giving
the first equation of (D.58c).
Multiplication by $v^{-1}$ gives its other
two exact equations.

The executable
\texttt{code/part\_c/torus\_negative\_involution\_d58.py}
checks the mod-four one-cochains, the cyclic
resolution modules, the two explicit localized
algebra isomorphisms, and geometric word examples
of the cusp equation.
The full-group conclusions above follow from
restriction and the actual normalizer maps.
They assert no full triviality or procyclicity
of the fixed subgroups from their first
cohomology dimensions.
\end{proof}

\section{E. Special groups and two operator groups}

\subsection{E(a). Special groups from conjugators and from framed diagrams}
\label{sec:ea}
We give the constructions of \S25, folios 149--154, scans 296--306,
with their precise hypotheses. Let $\pi$ be a discrete hyperbolic
surface group or its full profinite completion, and let $I$ be its finite
peripheral index set. Choose representatives $L_i$ and coherent
parametrizations $c_i:T\to L_i$, where $T=\Z$ or $\widehat\Z$.
The self-centralizing and self-normalizing properties needed below
are the results of Part A. Put
$\mathcal E=\Aut_{\rm lac}(\pi)$ and $\Gamma=\Out_{\rm lac}(\pi)$.
Every $\alpha\in\mathcal E$ has a permutation $\sigma_\alpha$ of $I$
and an automorphism $\chi_\alpha$ of $T$.

\begin{claim}{E.1}{Proved}{The conjugator construction and the product kernel}
Define $S\mathcal E$ to consist of pairs $(\alpha,(g_i)_{i\in I})$
with $\alpha\in\mathcal E$, $g_i\in\pi$, and
\[
 \alpha(c_i(t))=g_i^{-1}c_{\sigma_\alpha(i)}(\chi_\alpha t)g_i
 \quad(t\in T).
\]
The multiplication is
\[
 (\beta,(h_i))(\alpha,(g_i))
 =\bigl(\beta\alpha,
       (h_{\sigma_\alpha(i)}\,\beta(g_i))_{i\in I}\bigr).
\]
There are exact sequences and a canonically embedded normal subgroup
\[
 1\to T^I\to S\mathcal E\to\mathcal E\to1,
 \qquad
 \iota:\pi\hookrightarrow S\mathcal E,
 \quad a\mapsto(\operatorname{Int}(a),(a^{-1})_i).
\]
The groups $T^I$ and $\iota(\pi)$ commute and intersect trivially.
Consequently
\[
 S\Gamma=S\mathcal E/\iota(\pi),\qquad
 1\to T^I\times\pi\to S\mathcal E\to\Gamma\to1,
 \qquad 1\to T^I\to S\Gamma\to\Gamma\to1.
\]
The superscript $!$ takes the inverse image of the kernel of the
permutation action on $I$. All maps in the profinite case are continuous
maps of profinite groups, and all kernels and quotients are closed.
\end{claim}
\passage{\S{}25; folios 149--151; scans 296--300}

\begin{proof}
Apply $\beta$ to the defining equation for $\alpha$. The two
conjugating factors multiply as
$\beta(g_i)^{-1}h_{\sigma_\alpha(i)}^{-1}$, which is the inverse of the
factor displayed in the product. Associativity follows either by
substitution in that formula or by associativity of composition of the
specified conjugations. The identity is $(\id,(1)_i)$, and the inverse
has automorphism $\alpha^{-1}$ and $j$th conjugator
$\alpha^{-1}(g_{\sigma_\alpha^{-1}(j)}^{-1})$.

Surjectivity to $\mathcal E$ is the definition of preserving the loop
classes. In its kernel every $g_i$ centralizes $L_i$, hence belongs to
$L_i$; this identifies the kernel with $T^I$. Substituting
$(\operatorname{Int}(a),(a^{-1})_i)$ proves that $\iota$ is a
homomorphism and that conjugation by $(\alpha,(g_i))$ takes $\iota(a)$
to $\iota(\alpha(a))$. Its injectivity follows from $Z(\pi)=1$, including
when $I$ is empty. To check commutation with the kernel, the two products
with $(\id,(\ell_i))$, $\ell_i\in L_i$, both have conjugators
$\ell_i a^{-1}$. Their intersection is trivial because projection to
$\mathcal E$ sends $\iota(a)$ to $\operatorname{Int}(a)$. Finally an
element lying over the identity in $\Gamma$ has an inner automorphism
component and can be multiplied by a unique $\iota(a)^{-1}$ to put that
component equal to the identity. This proves that the whole kernel is
the displayed direct product, not merely an abstract extension of its
two factors.

For a finitely generated profinite $\pi$, its automorphism group has the
usual profinite topology obtained from characteristic open quotients.
The defining equations of $S\mathcal E$ are closed equations in
$\mathcal E\times\pi^I$. Thus it is profinite, the multiplication and
all maps above are continuous, and the compact subgroups giving the
quotients are closed. The same algebraic proof gives the exact
continuous sequences.
\end{proof}

\begin{claim}{E.2}{Proved}{Agreement with autoequivalences of the framed diagram}
The group $S\Gamma$ of E.1 is the group of isomorphism classes of
autoequivalences of
\[
 B_I\longrightarrow\coprod_{i\in I}BL_i\longrightarrow B\pi
\]
with the common orientation action retained. Adding a fixed base object
in $B\pi$ gives $S\mathcal E$. Thus the two constructions on scans
296--300 agree, with the comparison maps and inner ambiguities exactly
as in E.1.
\end{claim}
\passage{\S{}25; folios 150--151; scans 298--300}

\begin{proof}
An autoequivalence gives a permutation of the selected objects, the
specified automorphism $\chi$ on their cyclic automorphism groups,
and an automorphism $\alpha$ at the generic base object. The natural
isomorphism making the second square commute is, at its $i$th object,
an element $g_i\in\pi$. Its naturality equation is exactly the defining
equation of E.1. Pasting the squares gives the product there, including
the index $\sigma_\alpha(i)$.

Initially the first square also contains an isomorphism at each selected
peripheral object. A natural transformation at that object changes this
isomorphism by the corresponding element of $L_i$. Use it to make the
first-square isomorphism the identity. After this gauge choice there
is no remaining peripheral-object natural-transformation ambiguity.
A natural transformation at the generic object is an element $a\in\pi$
and changes the representative by $\iota(a)$, with the convention of
E.1. These are precisely the equivalences being divided out in
$S\Gamma$. Fixing the generic base object and its comparison
isomorphism removes that last ambiguity and gives $S\mathcal E$.
Conversely every pair in E.1 defines the three functors and their two
commutativity squares. The two constructions undo each other on the
chosen representatives and on the specified natural transformations,
so the group isomorphisms are canonical relative to the initial
peripheral representatives. Changing those representatives transports
the diagrams by their explicit conjugating isomorphisms.
\end{proof}

\begin{claim}{E.3}{Proved}{The puncture-stabilizer pushout is the normalizer group}
Let $\Gamma_i\subset\Gamma$ stabilize $i$. Restrict the extension
$S\Gamma\to\Gamma$ to $\Gamma_i$ and push its kernel out by
$\operatorname{pr}_i:T^I\to T$. The resulting extension is canonically
\[
 1\to L_i\xrightarrow{\operatorname{Int}}
 \{\alpha\in\Aut_{\rm lac}(\pi):\alpha(L_i)=L_i\}
 \longrightarrow\Gamma_i\to1.
\]
\end{claim}
\passage{\S{}25; folios 153--154; scans 305--306}

\begin{proof}
For a representative $(\alpha,(g_j))$ with $\sigma_\alpha(i)=i$, put
\[
 F_i[(\alpha,(g_j))]=\operatorname{Int}(g_i)\alpha.
\]
The defining equation shows that this automorphism normalizes the
actual subgroup $L_i$. Replacing the pair by its generic-basepoint
change replaces $\alpha$ by $\operatorname{Int}(a)\alpha$ and $g_i$ by
$g_i a^{-1}$; the displayed value is unchanged. Hence the map descends
to $S\Gamma_i$. The multiplication formula gives
\[
 \operatorname{Int}(h_i\beta(g_i))\beta\alpha
 =\operatorname{Int}(h_i)\beta\operatorname{Int}(g_i)\alpha,
\]
so it is a homomorphism. An automorphism normalizing $L_i$ has a
representative with $g_i=1$, proving surjectivity. On the kernel $T^I$
the map is $\ell_i\mapsto\operatorname{Int}(\ell_i)$; its kernel is
exactly the other factors because $Z(\pi)=1$. This is the universal
pushout claimed. The normalizer equality in Part A gives the last
kernel $L_i$ in the displayed sequence. This is the compatibility of
scans 305--306, with its actual group of automorphisms rather than an
unspecified outer class.
\end{proof}

\begin{claim}{E.4}{Refuted}{The product-kernel claim requires the rigidity hypotheses}
The formula $\ker(S\mathcal E\to\Gamma)=T^I\times\pi$ does not hold
unchanged for all abelian types. For the unpunctured torus,
$I=\varnothing$, $S\mathcal E=\operatorname{GL}_2(\Z)=\Gamma$ and the
kernel is trivial, whereas $\pi=\Z^2$. For the annulus the inner
basepoint subgroup intersects the peripheral kernel, and the special
outer kernel is a quotient of $T^2$ by one diagonal copy of $T$.
\end{claim}
\passage{\S{}25; folios 149--153; scans 296--304}

\begin{proof}
For the torus, all inner automorphisms of $\Z^2$ are the identity and
there are no conjugators indexed by punctures. Thus the pair
construction contains only the automorphism itself, giving the stated
identity map and its trivial kernel.

For the annulus, both peripheral images are the whole group
$\pi=\Z$. The conjugator kernel is $\pi^2$, but the generic-basepoint
change $a$ is $(\id,(a^{-1},a^{-1}))$ and lies in that same kernel.
Dividing by these changes gives $\Z^2/\Delta\Z\cong\Z$, with the
antidiagonal convention after orienting both boundary loops. Keeping
the two end labels and the independent orientation data supplies the
appropriate extended quotient; their omission would create an
additional error. For the plane the peripheral homomorphism
$\Z\to1$ is not injective, so the normalizer calculation used in E.1
is unavailable. Its correct fixed-boundary positive special group is
trivial, as C.20 computes. These explicit maps identify the failed
hypotheses in each case.
\end{proof}

\subsection{E(b). Two commuting operator groups}

We retain the names $\theta$ and $\omega$ for the standard order-two and order-three lifts on $\widehat F_2$; Grothendieck's $\sigma_0$ and $\rho$ denote the corresponding choices on scans 316--317, \S25 bis, folios 159--160. A change of base path conjugates a lift by $\widehat F_2$. We keep this ambiguity in the statements below.

\begin{claim}{E.20}{Proved}{The two lifting conditions cut out exactly the multiplier-$1$ subgroup modulo $3$}
In $B=\Out_{\rm lac}(\widehat\pi_{0,3})$ let $S_3$ be the geometric positive subgroup. Its centralizer is $GT_0$, with the normalized description $(\lambda,f)$ of Harbater--Schneps. Every element of $GT_0$ has an automorphism representative commuting with the chosen order-two lift $\theta$. The elements with an automorphism representative commuting with the chosen order-three lift $\omega$ are exactly
\[
 GT_0^{(3)}=\{(\lambda,f)\in GT_0:\lambda\equiv1\pmod3\}.
\]
This subgroup has index two. The two representatives need not be the same representative; this is the two-condition formulation on scans 316--317.
\end{claim}
\passage{25 bis; 157 verso--160; 312--317}

\begin{proof}
Harbater--Schneps, Main Theorem (a) and Theorem A(i) \cite{R-HS}, identify the centralizer and state precisely
\[
 [\theta,F]=\Inn(f),\qquad
 [\omega,F]=\Inn(y^mf),\qquad m=(\lambda-1)/2,
\]
for $F(x)=x^\lambda$, $F(y)=f^{-1}y^\lambda f$. Here $[a,b]=aba^{-1}b^{-1}$. A representative $\Inn(h)F$ commutes with $\theta$ precisely when
\[
 \theta(h)f=h,
\]
since $F_2$ has trivial centre. It commutes with $\omega$ precisely when $\omega(h)y^mf=h$. These are the two exact coboundary equations. We determine their solvability without imposing the pentagon equation.

The positive lift extension of the $S_3$ action on the three-punctured sphere is
\[
 E=C_2*C_3=\langle s,r\mid s^2=r^3=1\rangle,
 \qquad E\longrightarrow S_3,
\]
whose kernel is the peripheral free group $F_2$. Its cusp generator can be chosen as $e=sr$; $e^2$ is a primitive cusp in that kernel. Complete this exact sequence by D.1. An outer centralizer element, represented by $F$, acts on $\widehat E$ by conjugating its image in $\Aut(\widehat F_2)$ and induces the identity on $S_3$. It sends $e^2$ to a $\widehat F_2$-conjugate of $e^{2\lambda}$. After that conjugation it sends $e$ into $C_{\widehat E}(e^{2\lambda})=\overline{\langle e\rangle}$. To prove this equality, put $L=\overline{\langle e^2\rangle}$. A centralizing element preserves its $\widehat F_2$-conjugacy class, so its image in $S_3$ stabilizes the associated puncture. That stabilizer is $\langle q(e)\rangle$ of order two. Multiplication by $e$ when necessary puts the element in $\widehat F_2$, where Part A gives $C_{\widehat F_2}(e^{2\lambda})=L$. This proves the equality and identifies the whole centralizer with $\widehat\Z$. Doubling is injective in that group, so the image of $e$ is $e^\lambda$.

Consequently this automorphism acts on $\widehat E_{\rm ab}=C_6$ by multiplication by $\lambda$. Every order-two element is conjugate to $s$, and every order-three subgroup is conjugate to $\langle r\rangle$, by D.2 or the free profinite product theorem. A lift of the given transposition is $F_2$-conjugate to $s$: a total-group conjugator projects to $C_{S_3}(s)=\langle s\rangle$, and multiplication by $s$ adjusts it into the kernel. Thus there is one kernel-conjugacy class for the order-two lift. It is fixed by every centralizer element, proving solvability of the first equation.

There are exactly two kernel-conjugacy classes of lifts of the specified three-cycle. They are represented by $r$ and $s r^{-1}s^{-1}$. Indeed a total-group conjugator has either an even or an odd image in $S_3$; adjusting by $\langle r\rangle$ or by $s\langle r\rangle$ gives these two representatives. They are distinct because their images in the order-three quotient of $E_{\rm ab}$ are inverse nonidentity elements. The action of $(\lambda,f)$ preserves the class of $r$ exactly when $\lambda\equiv1\pmod3$, by the computed action on $C_6$. When it does, the same adjustment by $r$ puts its conjugator in the kernel. This is precisely solvability of the second coboundary equation. The argument treats elements, so it retains the label of the three-cycle.

For comparison, Lochak--Schneps, Proposition 3 and Theorem 2 \cite{R-LS}, compute the same two nonabelian cohomology sets: one class for $\theta$, two for $\omega$, with the second class selected by $\lambda\equiv-1\pmod3$. The first two parts of their proof use only the two-cycle and three-cycle equations, so their computation also applies to $GT_0$. Finally $(-1,1)\in GT_0$ exchanges the two classes. The homomorphism to $(\Z/3)^\times$ is therefore surjective and the asserted index is exactly two. The residue calculation is independently recorded in the executed program.
\end{proof}

\begin{claim}{E.21}{Proved}{The arithmetic subgroup and the two fixed points}
The intersection of the usual Galois image $G_{\Q}\to GT_0$ with $GT_0^{(3)}$ is the image of $G_{\Q(\sqrt{-3})}$. The order-three fixed points in $U=\mathbf P^1-\{0,1,\infty\}$ are $e^{\pm i\pi/3}$, and the order-two lift has one fixed point in $U$. Thus the requirement to fix the chosen order-three section is a proper finite-index restriction, even on the Galois subgroup.
\end{claim}
\passage{25 bis; 157 verso--160; 312--317}

\begin{proof}
The multiplier of the Galois action is the cyclotomic character, by its action on the tame inertia at $0$. Its reduction modulo $3$ is the action on $\mu_3$ and has fixed field $\Q(\zeta_3)=\Q(\sqrt{-3})$. E.20 identifies precisely its kernel. The three-cycle $z\mapsto1/(1-z)$ has fixed-point equation $z^2-z+1=0$, and the transposition $z\mapsto1-z$ has fixed-point equation $z=1/2$. These points lie in $U$. Complex conjugation exchanges the former pair and fixes the latter point, which also proves properness directly.

On scan 312 the relevant subgroup is obtained by fixing a finite set of critical data; it is therefore of finite index. An assertion that every element of $G_\Q$ lifts commuting with a specified one of the order-three lifts would contradict the displayed computation. We record the scan's wording as editorial evidence, not as a settled reading.
\end{proof}

\begin{claim}{E.22}{Proved}{Normalizer, stabilizer, and the forced lift}
Let $1\to\pi\to E\xrightarrow q\mathcal G\to1$ have centre-free kernel, and let $K\leq E$ be a finite subgroup on which $q$ is injective. Assume $C_\pi(K)=1$, as holds for a nontrivial rotation stabilizer in the anabelian geometric situation and its completion. Put $G_Q=q(K)$ and let $Q$ denote the $\pi$-conjugacy class of $K$. The restriction
\[
 q:N_E(K)\longrightarrow\mathcal G
\]
is injective and has image the stabilizer $\mathcal G_Q$ of $Q$. Its inverse on that image is a canonical homomorphism. In a two-group situation $\mathcal G=\Gamma G$ with $\Gamma$ centralizing $G$, its restriction gives the forced lift of $\Gamma_Q=\Gamma\cap\mathcal G_Q$, and this lift centralizes $K$.
\end{claim}
\passage{25 bis; 155--157 verso; 307--312}

\begin{proof}
An element of $N_E(K)\cap\pi$ acts trivially on $q(K)$, and the restriction of $q$ to $K$ is injective. It therefore centralizes $K$, so the kernel is $C_\pi(K)=1$. A quotient element stabilizes $Q$ precisely when one, and hence every, lift carries $K$ to a $\pi$-conjugate. Multiplying that lift by the appropriate element of $\pi$ produces an element of $N_E(K)$ above it. This proves the image statement and uniqueness of the lift. The inverse is a homomorphism because it is the inverse of an injective homomorphism onto its image. For $\gamma\in\Gamma_Q$, conjugation of its lift on $K$ projects to conjugation by $\gamma$ on $G_Q$, which is trivial. Injectivity on $K$ then makes this conjugation trivial. In the profinite case all these maps are continuous, and the inverse is continuous by compactness of the closed normalizer and the Hausdorff property of the quotient.

For the centralizer, the precise image is $\mathcal G_Q\cap C_{\mathcal G}(G_Q)$, rather than all of $\mathcal G_Q$. When $G_Q$ is nonabelian dihedral, its own elements already distinguish these images. For a reflection in an oval the kernel is $\Z$ or $\widehat\Z$ by Part B and D.4, so the stated injectivity is replaced by an exact sequence with that kernel. These computations identify exactly the rigidity used in the manuscript's argument.
\end{proof}

\begin{claim}{E.23}{Proved}{Explicit discrete centralizers and pointwise kernels}
For the three actions in the question, the full extended geometric groups, their centralizers, their normalizers, and their pointwise kernels on the exceptional points are as follows. Here $D_{2n}$ has order $2n$.
\[
\begin{array}{c|c|c|c|c}
 (g,\nu),G&M&C_M(G)&N_M(G)&C_M(G)^{\natural}\\ \hline
 (0,3),S_3&S_3\times C_2&C_2&S_3\times C_2&1\\
 (1,1),C_4&\mathrm{GL}_2(\Z)&C_4&D_8&C_2\\
 (1,1),C_6&\mathrm{GL}_2(\Z)&C_6&D_{12}&1
\end{array}
\]
The last column fixes every exceptional point individually, not merely its orbit under $G$.
\end{claim}
\passage{25 bis; 155--160; 307--317}

\begin{proof}
For type $(0,3)$, the geometric group is $S_3$ times complex conjugation. The exceptional points are the three transposition fixed points $-1,1/2,2$ and the two order-three fixed points of E.21. Complex conjugation fixes the first three and interchanges the last two. The centralizer of $S_3$ is therefore $C_2$ and its pointwise kernel is trivial; its normalizer is the whole product.

Use the matrices
\[
 R_4=\begin{pmatrix}0&-1\\1&0\end{pmatrix},\qquad
 R_6=\begin{pmatrix}0&-1\\1&1\end{pmatrix}
\]
on $\R^2/\Z^2$, punctured at $0$. A rational matrix commuting with either one is $aI+bR_n$. For an integral commuting matrix $a,b\in\Z$. Its determinant is $a^2+b^2$ for $R_4$ and $a^2+ab+b^2$ for $R_6$. Equating this nonnegative norm to $1$ gives four and six units respectively, proving the centralizer assertions. Conjugation in the normalizer embeds its quotient by the centralizer in $\Aut(C_n)=C_2$ for $n=4,6$. Reflection in the square or hexagonal lattice realizes inversion, proving the two normalizer assertions.

For $C_4$ the exceptional set outside the puncture is exactly the three nonzero two-torsion points. The matrix $-I$ fixes all three, while $R_4$ interchanges $(1/2,0)$ and $(0,1/2)$. Thus the kernel is $\langle-I\rangle$. For $C_6$ the exceptional set consists of the same three two-torsion points and $(1/3,1/3),(2/3,2/3)$. The three two-torsion points force a power of $R_6$ to be $I$ or $-I$, and the latter exchanges the two order-three points. Hence the kernel is trivial. These lists are exhaustive: each nonidentity power has its fixed points given by $(R_n^j-I)v\in\Z^2$; their union is the displayed finite set. The program enumerates these solutions, the kernels, and the matrix centralizers and normalizers. The norm calculation proves exhaustiveness beyond its bounded integer search.
\end{proof}

\begin{claim}{E.24}{Refuted}{Recovery of an arbitrary commuting group from its pointwise subgroup}
For an arbitrary specified $\Gamma\leq C_{\Gamma'}(G)$, its pointwise subgroup is
\[
 \Gamma^{\natural}=\Gamma\cap\Gamma'^{\natural},
\]
and it need not equal $C_{\Gamma'^{\natural}}(G)$. Nor is $\Gamma$ determined as a normalizer by $\Gamma^{\natural}$. For the maximal choice $\Gamma=C_{\Gamma'}(G)$ the proposed equality holds, by the displayed intersection identity.
\end{claim}
\passage{25 bis; 157 verso; 312}

\begin{proof}
Both identities follow by testing separately the two defining conditions: membership in $\Gamma$ and pointwise stabilization of each critical datum. To refute recovery, take the first row of E.23. The choices $\Gamma=1$ and $\Gamma=C_2$ have the same pointwise subgroup $1$. Its normalizer inside the centralizer $C_2$ is $C_2$, so the procedure cannot recover the first choice. For the proposed equality without the maximal-choice hypothesis, use the second row and take $\Gamma=1$: its pointwise subgroup is $1$, whereas the centralizer pointwise kernel is $C_2$. All groups and their kernels are included in the finite check. Thus the failure persists even in the specifically requested discrete cases and has a complete intersection formula as its replacement.
\end{proof}

\section{F. Statements usable in the later sections}

We use $\Gamma_{g,\nu}$ for the extended geometric group permitting puncture permutations and $\Gamma_{g,\nu}^{!+}$ for the pure oriented group, following \S19, folio 102, scans 204--205. In each formula below the chosen variant is stated explicitly. The full profinite loop-outer group is $B=\Out_{\rm lac}(P)$; the closure $H=\overline M$ of a chosen geometric subgroup $M$ is a subgroup of $B$. We do not identify that closure with the abstract profinite completion without the congruence assertion required for that identification.

\begin{claim}{F.1}{Proved}{LM4-F.a: fixed discretizations and prediscretizations}
\label{LM4-F.a}
For $M=\Gamma_{g,\nu}$ and $H=\overline M\leq B$, the sets of outer discretizations and prediscretizations are $B/M$ and $B/H$. For a given finite $G\leq B$ their fixed sets are exactly D.16's transporter sets and centralizer-orbit decompositions. In the three types of D.8, every fixed prediscretization has a fixed outer discretization above it. The same lifting conclusion holds for a finite solvable $G$ and genus at most two for the orientation-preserving geometric group covered by Boggi's Theorem C. Over a fixed coset represented by $b$, the fibre is
\[
 (H/M)^{b^{-1}Gb},
\]
with its full labelled-embedding decomposition as in D.16.
\end{claim}
\passage{28; 178--181; 352--358}

\begin{proof}
The identifications and the fixed-set computation are proved in D.16, including their inner-automorphism ambiguities. A fixed coset $bH$ has $K=b^{-1}Gb\leq H$. Its fibre consists of $bhM$, $h\in H$; this coset is fixed precisely when $h^{-1}Kh\leq M$, proving the formula. In the three low types D.8 conjugates every finite $K\leq H$ into $M$, giving a point in that fibre. In genus at most two and for solvable $K$, Boggi's Theorem C \cite{R-Boggi} gives the same conjugation into the oriented geometric group. The congruence subgroup property in these genera identifies its completion with $H$, so the theorem applies to the displayed subgroup. We retain the labelled embeddings in the fibre decomposition because subgroup conjugacy alone does not fix an automorphism of $K$.

The actual based version has $A(P)/J$ in place of $B/M$, as D.16 proves. A chosen genuine lift of $G$ acts on it; an outer subgroup alone acts on its quotient by $P$. These are the exact maps through which the fixed-set assertions of \S28, folios 178--181, scans 352--358 are used. For a finite group in the full $B$ not known to stabilize a prediscretization, the same formula still determines the fixed set as a transporter; its nonemptiness is exactly the unresolved recognition question D.12, not an additional conclusion of this lemma.
\end{proof}

\begin{claim}{F.2}{Refuted}{LM4-F.b, finiteness and equality of the two centralizers}
\label{LM4-F.b}
The centralizer of a finite subgroup $G\leq\Gamma^+_{g,\nu}$ need not be finite, either discretely or in the geometric completion. In genus at most two, for finite solvable $G$, the normalizer comparison and the open-kernel argument below give
\[
 C_{\widehat{\Gamma^+_{g,\nu}}}(G)
       =\overline{C_{\Gamma^+_{g,\nu}}(G)}.
\]
The corresponding equality with the full loop-outer group in place of the geometric completion is false already for $(g,\nu)=(0,3)$ and $G=S_3$.
\end{claim}
\passage{30; 194--196; 383--387}

\begin{proof}
On a closed genus-two surface let $G$ be generated by the hyperelliptic involution. This involution is central in $\Gamma_2^+$, so its centralizer is the entire group. A Dehn twist about a nonseparating curve has infinite order: its action on integral homology sends a transverse vector $b$ to $b+a$, and its $n$th power sends it to $b+na$. This also gives infinite order in the congruence image, since some reduction modulo an integer detects every nonzero $n$. Thus the completed centralizer is infinite as well. For finite solvable $G$, the normalizer comparison, with the source induction and its topology input audited in F.7, gives $N_{\widehat M}(G)=\overline{N_M(G)}$ for $M=\Gamma^+_{g,\nu}$. The continuous conjugation homomorphism from this profinite normalizer to the finite group $\Aut(G)$ has open kernel $C_{\widehat M}(G)$. Intersecting the dense subgroup $N_M(G)$ with that open kernel gives a dense subgroup of the kernel: every open subset of the kernel is open in the normalizer and meets $N_M(G)$. The intersection is exactly $C_M(G)$, proving the displayed centralizer equality. The map of the completion to the loop-outer group is injective in these genera by the congruence subgroup property. The cited sources state the equality for every finite $G$; F.7 records the precise missing intersection step for their nonsolvable assertion. The passage from any such equality to finiteness would still require the discrete centralizer to be finite, contradicted by the hyperelliptic example.

For the smallest hyperbolic type, $\Gamma_{0,3}^+=S_3$, and $C_{S_3}(S_3)=1$. Harbater--Schneps identify $C_B(S_3)=GT_0$. This group is infinite: its multiplier contains the cyclotomic image $\widehat\Z^\times$ of $G_\Q$. These are actual homomorphisms on the same group $\widehat F_2$, with the geometric $S_3$ supplied by rational M\"obius transformations, so the commutation with the Galois image is the naturality of the fundamental-group action. The two centralizers are therefore $1$ and an infinite group. With orientation reversal included the geometric centralizer becomes $C_2$, and the same failure remains.

The statement denoted (c) on scans 383--385, \S30, folios 194--195, concerns the centralizer of an arithmetic germ $\Gamma^{\natural}$. It cannot be replaced merely by a finite $G$: scan 385 itself discusses infinite centralizers of finite groups. The counterexamples above compute the result of the requested replacement explicitly and explain which datum that replacement removes.
\end{proof}

\begin{claim}{F.3}{Proved}{LM4-F.b, the intrinsic open-image statement}
Let $M\leq B$ be an oriented geometric subgroup, $H=\overline M$, $N=N_B(H)$, and let $i:G\hookrightarrow H$ be a fixed finite embedding. Conjugation defines an action of $N/H$ on the set of $H$-conjugacy classes of embeddings $G\hookrightarrow H$. The image of the continuous homomorphism
\[
 C_N(i(G))\longrightarrow N/H
\]
is exactly the stabilizer of $[i]$, and its kernel is $C_H(i(G))$. When the orbit of $[i]$ is finite, this image is open. The orbit is finite for solvable finite $G$ in genus at most two and for every finite $G$ in the three low types of D.8.
\end{claim}
\passage{30; 195 verso--196; 386--387}

\begin{proof}
An element $nH$ stabilizes $[i]$ exactly when there is $h\in H$ with
\[
 ni(g)n^{-1}=hi(g)h^{-1}\quad(g\in G).
\]
This is equivalent to $h^{-1}n\in C_N(i(G))$. It proves the image assertion, and the kernel is obtained by intersecting the centralizer with $H$. Because $H$ is compact, the stabilizer is the image of a compact closed transporter and is closed in $N/H$. A closed subgroup of finite index in a profinite group is open: its finitely many closed cosets have open complements, so each coset is open.

For the claimed cases, finite subgroup conjugacy classes in $H$ come from the finite list of geometric classes, by D.8 or Boggi's solvable statement. There are finitely many geometric classes for each fixed finite group: a finite action has an orbifold quotient of bounded signature by the Euler characteristic formula; there are finitely many epimorphisms from each finitely generated orbifold group to the fixed finite group. Passing from subgroup classes to labelled-embedding classes adds at most $|\Aut(G)|$ choices per subgroup class. Thus the orbit is finite and the image is open. F.6 constructs the arithmetic quotient directly from universal monodromy and proves the open-image assertion in that quotient for every finite geometric subgroup and every hyperbolic type. The inclusion of that quotient into $N/H$ is retained there as the exact comparison map.
\end{proof}

\begin{claim}{F.4}{Proved}{LM4-F.c: the exact canonical lift of a section stabilizer}
\label{LM4-F.c}
Let $\pi$ be a centre-free discrete surface group or its full profinite completion, and put
\[
 A=\Aut_{\rm lac}(\pi),\qquad B=\Out_{\rm lac}(\pi).
\]
Let $G\leq B$ be finite and let $s:G\to A$ be a specified lift with $C_\pi(s(G))=1$. The natural map
\[
 C_A(s(G))\longrightarrow C_B(G)
\]
is injective. Its image is the stabilizer of the $\pi$-conjugacy class $[s]$ in the set of lifts of the labelled inclusion $G\hookrightarrow B$. Its inverse is the canonical homomorphism lifting that stabilizer to $A$ and centralizing $s(G)$. In the profinite case it is continuous. For geometrically realized actions with a nontrivial rotation the section-class set is finite, so this stabilizer has finite index and is open in the profinite case.
\end{claim}
\passage{31; 197 verso--199; 390--393}

\begin{proof}
The kernel is the inner automorphisms centralizing $s(G)$, identified, by centre-freeness, with $C_\pi(s(G))$. For $b\in C_B(G)$ choose an automorphism representative $v$. The conjugate homomorphism $g\mapsto vs(g)v^{-1}$ still lifts the labelled inclusion, because $b$ centralizes each $g$. It represents $[s]$ exactly when some $p\in\pi$ makes $\Inn(p)^{-1}v$ commute with every $s(g)$. This proves that the image equals the claimed stabilizer. The kernel computation proves uniqueness. The inverse is a homomorphism by uniqueness, since products of commuting lifts still commute with $s(G)$.

For a finitely generated profinite $\pi$, $A$ and $B$ are profinite, and the centralizer is closed. The map is a continuous injection from a compact space into a Hausdorff space, so its inverse on its image is continuous. In the discrete geometric situation Part B proves finiteness of the lift classes; D.3 proves their equality with the profinite classes. The orbit of a class is therefore finite. Its stabilizer is closed, as the image of the compact centralizer, hence open. This proves all topology assertions as well as the group assertion.

For the notation of scans 390--393, \S31, folios 197--199, the target $\mathfrak S(\pi_0)^G$ is the centralizer of the chosen genuine lift inside the loop automorphism group. The source is the stabilizer of the corresponding section class inside $\mathcal T(\pi_0)^G=C_{\Out_{\rm lac}(\pi_0)}(G)$. A point $P$ determines that section class after a base path; changing the path conjugates the resulting lift by $\pi_0$, exactly as displayed. Thus the canonicity belongs to the chosen section, and the geometric point supplies its conjugacy class.
\end{proof}

\begin{claim}{F.5}{Refuted}{The unrestricted uniqueness assertion for reflection ovals}
For a reflection fixing an oval, the lifting map in F.4 has kernel $\Z$ in the discrete case and $\widehat\Z$ after completion. Consequently it is not injective and its lifts are not unique without further data.
\end{claim}
\passage{31; 197 verso--199; 390--393; 33; 222 verso--225; 440--445}

\begin{proof}
Part B and D.4 identify $C_\pi(s(G))$ with the primitive oval subgroup and its closure. Every element of that subgroup induces an inner automorphism centralizing the chosen reflection lift and maps to the identity outer class. The kernel computation in F.4 therefore gives the displayed groups. A nonidentity primitive oval element is an explicit second lift of the identity outer element. The corrected general statement is the exact sequence with this kernel and with image the same section stabilizer, proved by the same representative-and-conjugator calculation. No splitting of that sequence is supplied by the point class alone.
\end{proof}

\begin{claim}{F.6}{Proved}{LM4-F.b: open arithmetic image for every finite geometric subgroup}
Let $2g-2+\nu>0$, let $M=\Gamma^+_{g,\nu}$ permit permutations
of the $\nu$ punctures, and retain
\[
 P=\widehat\pi_{g,\nu},\qquad
 B=\Out_{\rm lac}(P),\qquad H=\overline M\leq B.
\]
Write $A_{\rm ar}\leq B$ for the image of universal \emph{arithmetic}
\'etale monodromy of the moduli stack of smooth genus-$g$ curves
with an unordered reduced divisor of degree $\nu$, over $\Q$.
Then $H\triangleleft A_{\rm ar}$, and put
$\Gamma_{\rm ar}=A_{\rm ar}/H$.
For every finite subgroup $G\leq M$, the continuous homomorphism
\[
 C_{A_{\rm ar}}(G)\longrightarrow\Gamma_{\rm ar}
                                                        \tag{F.6.1}
\]
has open image.  Its kernel is $C_H(G)$.
In the notation $N=N_B(H)$ and $q:N\to N/H$, this is exactly
the assertion that
\[
 q(C_N(G))\cap\Gamma_{\rm ar}
       =q(C_{A_{\rm ar}}(G))
       \quad\hbox{is open in }\Gamma_{\rm ar}.           \tag{F.6.2}
\]
The same assertion holds with ordered punctures and the pure
oriented group throughout.  No identification of $H$ with the
abstract completion of $M$ is used.
\end{claim}
\passage{30; 195--196; 385--387}

\begin{proof}
We first type the arithmetic group and its comparison maps.
Denote the indicated moduli stack over $\Q$ by $\mathscr M$,
and choose a geometric base point and an orientation-preserving
topological marking of its complex fiber.  Its universal
punctured curve supplies the outer monodromy homomorphism
\[
 R:\Pi_{\rm ar}=\pi_1^{\rm et}(\mathscr M)
           \longrightarrow B.
\]
The geometric comparison and the arithmetic exact sequence are
\[
 1\longrightarrow\widehat M
  \longrightarrow\Pi_{\rm ar}
  \longrightarrow G_{\Q}\longrightarrow1,
 \qquad R(\widehat M)=H.                               \tag{F.6.3}
\]
These are the constructions in \cite[\S\S7.1--7.2,
equation (29) and the universal \'etale monodromy
representation]{R-BoggiAut2026}.  For the geometric comparison,
finite \'etale covers analytify to finite topological covers;
the inverse comparison is compatible on an \'etale atlas and
its double and triple overlaps, so descends to the stack.
The marked complex moduli stack has contractible Teichm\"uller
space as its covering, with group $M$, which identifies its
topological fundamental group.  The universal curve then
identifies its monodromy with the completed geometric action
on $P$, rather than introducing a separate representation.
All these maps preserve peripheral conjugacy classes, since
the removed divisor is carried to itself in the universal
family.  Their base-point ambiguities are inner automorphisms
of $P$, already removed in $B$; a change of topological marking
conjugates the displayed subgroups by the corresponding element
of $M$.

Set $A_{\rm ar}=R(\Pi_{\rm ar})$.  It is compact and hence
closed in $B$.  Normality in (F.6.3) gives
$H\triangleleft A_{\rm ar}\leq N$.  The quotient map induces
a continuous surjection
\[
 \overline R:G_{\Q}\longrightarrow\Gamma_{\rm ar}.     \tag{F.6.4}
\]
The proof uses this quotient of $G_{\Q}$ as defined.  Its kernel
and the kernel of $\widehat M\to H$ need not be computed.

We next construct a number-field curve with the specified
\emph{labelled topological} $G$-action.  Finite realization
C.5 supplies a compact Riemann surface $X$, a reduced marked
divisor $D$, and a faithful holomorphic action of $G$ preserving
$D$ whose marked mapping classes are the given inclusion
$G\hookrightarrow M$.  Choose projective equations for $X$,
equations for $D$, and formulas for the finitely many
automorphisms and their inverses.  Their finitely many
coefficients lie in a finitely generated $\Q$-subalgebra
$T\subset\C$.  Adjoin the coefficients needed for their
identities and invert finitely many nonzero elements.  This
gives a proper smooth family of curves $\mathcal X\to V$,
a finite \'etale reduced divisor $\mathcal D\to V$, and
the same labelled $G$-action over a nonempty smooth open
subscheme $V\subset\operatorname{Spec}T$.  Smoothness and
disjointness are open conditions.  Faithfulness can be kept
on this open set as well: for each nonidentity $h\in G$,
the equalizer of $h$ and the identity is a proper closed
subset of the generic fiber, and this property persists
after shrinking.  The original inclusion $T\hookrightarrow\C$
defines a point of every such nonempty open, since it has
generic underlying point.

Base change to $\overline\Q$ and take the connected component
whose complex base change contains that point.  After a
finite extension this component and the family are defined
over a number field.  It is a smooth connected variety; its
complex points form a connected, locally path-connected
complex manifold and hence a path-connected space.  It
contains an $\overline\Q$-point, by taking a closed point
in any nonempty affine open.  Choose a piecewise smooth
complex path from the original fiber to this fiber.

For precision, parallel transport along this path preserves
the labelled action.  A smooth horizontal splitting of the
relative tangent sequence can be chosen to agree with the
tangent splitting on $\mathcal D$: locally the divisor is
a disjoint union of sections, and a partition of unity
extends their prescribed splittings.  Average this splitting
over the finite group $G$.  Averaging retains its right-inverse
property and its prescribed values on the divisor, and
makes it $G$-equivariant.  The horizontal lift along the
compact path has complete transport, because the family
is proper and its fibers are compact.  Its transport is an
orientation-preserving diffeomorphism of pairs commuting
with each labelled element of $G$.  Thus the endpoint curve
has exactly the required topological action.  This constructs
a deformation and its comparison diffeomorphism; it does not
assert that the original complex curve is isomorphic to a
curve over $\overline\Q$.

All coefficients of the endpoint curve, divisor, and the
finitely many labelled automorphisms lie in one finite
extension $K/\Q$.  Let $x_K:\operatorname{Spec}K\to\mathscr M$
be the resulting moduli point.  Choose the geometric path
to $x_K$ supplied by the transport just constructed and the
finite-cover comparison.  The morphism $x_K$ then supplies
a continuous homomorphism
\[
 s_K:G_K\longrightarrow\Pi_{\rm ar}
\]
whose composite with $\Pi_{\rm ar}\to G_{\Q}$ is the
inclusion $G_K\hookrightarrow G_{\Q}$.  Under the transported
marking, $R\circ s_K$ is the Galois outer action on the
punctured endpoint curve.

Every $h\in G$ is an automorphism over $K$ of that curve and
its marked divisor.  Naturality of the fiber functor on finite
\'etale covers gives
\[
 (R\circ s_K)(\sigma)\,h
       =h\,(R\circ s_K)(\sigma)
 \quad\hbox{in }B\qquad(\sigma\in G_K).                \tag{F.6.5}
\]
Indeed, pullback by $h$ commutes with scalar conjugation by
$\sigma$; choosing paths for the individual base points can
alter the induced automorphisms only by inner automorphisms
of $P$.  Passing to $B$ gives precisely the equality written.
The equivariant transport aligns these $h$ with the original
labelled inclusion, so (F.6.5) places $R(s_K(G_K))$ in
$C_{A_{\rm ar}}(G)$ itself.

The image of (F.6.1) therefore contains
$\overline R(G_K)$.  The latter has index at most
$[K:\Q]$ in $\Gamma_{\rm ar}$ by the surjectivity of
(F.6.4).  It is compact and hence closed; a closed finite-index
subgroup of a profinite group is open, as proved in F.3.
Every subgroup containing it is a union of its open cosets.
This proves openness in (F.6.1), with kernel $C_H(G)$ by
intersection with $H$.

Finally $q^{-1}(\Gamma_{\rm ar})=A_{\rm ar}$, since
$H\subset A_{\rm ar}$ and $\Gamma_{\rm ar}=q(A_{\rm ar})$.
An element of $C_N(G)$ whose quotient lies in $\Gamma_{\rm ar}$
therefore already belongs to $C_{A_{\rm ar}}(G)$.
This proves both inclusions in (F.6.2).  Ordering the divisor
and using the pure moduli stack changes none of the maps or
arguments when $G$ belongs to the pure oriented group.

For \S30, folios 195--196, scans 385--387, formula (F.6.2)
is the arithmetic open-image assertion with its exact
quotient made explicit.  F.3 supplies the centralizer image
in the entire quotient $N/H$; F.6 supplies its openness
after intersection with the arithmetic image for every
finite geometric $G$.  The inclusion
$\Gamma_{\rm ar}\hookrightarrow N/H$ is the comparison
map between the two statements.  Its image has not been
identified with all of $N/H$ here, so no larger openness
assertion is obtained by suppressing that inclusion.
\end{proof}

\begin{claim}{F.7}{Open}{The nonsolvable centralizer comparison and its checked finite image}
For an oriented mapping class group $M$ of genus at most two,
the normalizer and centralizer comparisons for finite solvable
subgroups have the complete proof below. For a geometric
$A_5\le M=\Gamma^+_{0,\nu}$, the equality
\[
                 C_{\widehat M}(A_5)
                    =\overline{C_M(A_5)}                 \tag{F.7.1}
\]
is stated in \cite[Theorem~C(iii)]{R-Boggi} and
\cite[Theorem~5.1(iii)]{R-BZ}. The displayed proofs contain
an additional closure--intersection equality for which we
have not verified a justification. We do not assert that
(F.7.1) is false.

There is nevertheless equality of their images in the
finite puncture-permutation quotient. Write the invariant
puncture set as a union of some of the exceptional orbits
of sizes $12,20,30$ and $m$ free orbits of size $60$.
Both images fix the exceptional punctures individually
and induce the full group $A_5^m\rtimes S_m$ on the
free orbits. In particular, with only exceptional orbits
the unresolved part of (F.7.1) lies in the pure profinite
mapping class group; its discrete right side is trivial.
\end{claim}
\passage{\S{}30; folios 194--196; scans 383--387; Part D(b) and Part F(b): nonsolvable completion centralizers}

\begin{proof}[Source audit and proved comparisons]
The supplied preprint snapshot of \cite{R-Boggi}
has since appeared in \emph{S\~ao Paulo Journal of
Mathematical Sciences}, volume~20, article~35,
published 23 September 2026, DOI
\texttt{10.1007/s40863-026-00558-y}.
The publisher's full text was checked as well:
Proposition~2.4 retains the same displayed
closure--intersection step, and Lemma~3.2
still invokes that proposition.
The audit below applies to both inspected versions;
it is a proof-step audit and makes no claim based
on the paper's publication status.

For every finite subgroup $J\le M$, there are canonical
continuous maps
\[
 \widehat{C_M(J)}\longrightarrow
       \overline{C_M(J)}\lhook\joinrel\longrightarrow
       C_{\widehat M}(J).                              \tag{F.7.2}
\]
The first is onto by density and compactness. Its
injectivity states that the topology induced from $M$
on the discrete centralizer is its full profinite
topology. Surjectivity of the second arrow is the
ambient centralizer comparison. These are two specific
properties of the displayed maps; injectivity of the
first arrow does not supply surjectivity of the second.

In \cite[Proposition~2.4]{R-Boggi}, choose the displayed
generators $x_1,\ldots,x_k$ of the finite group, each of
prime-power order. The prime-power centralizer theorem
does prove
\[
 C_{\widehat M}(J)
   =\bigcap_i C_{\widehat M}(x_i)
   =\bigcap_i\overline{C_M(x_i)}.                      \tag{F.7.3}
\]
The source then identifies this last intersection with
$\overline{\bigcap_i C_M(x_i)}$. The same final equality
occurs in the proof of
\cite[Theorem~5.1(iii)]{R-BZ}, there indexed by all
$x\in J$. Neither displayed argument supplies a
closure--intersection hypothesis. This exact step is
what is left unverified here. The later virtually
compact special application in \cite[\S6.2]{R-BZ}
has a separate relatively quasiconvex product-separability
input; no such input has been proved here for the
punctured-sphere mapping class group.

\emph{A noncircular solvable proof.}
We first record the prime-power inputs used in the
induction. Goodness and finite virtual cohomological
type of mapping class groups of genus at most two,
and their finite-index subgroups, permit the direct
application of \cite[Theorem~A]{R-BZ}. Thus for a
finite $p$-subgroup $A$ we have
\[
 N_{\widehat M}(A)=\overline{N_M(A)},\qquad
 C_{\widehat M}(A)=\overline{C_M(A)}.                   \tag{F.7.4}
\]
The proof of \cite[Theorem~A and Corollary~B]{R-Boggi}
also supplies the full induced profinite topology on
$C_M(A)$ and hence on $N_M(A)$ in this prime-power case.
Here is why this use avoids the preceding gap. Its
free-action argument is the restriction-injectivity
argument on the finite covering surface group. Its
ramification-removal argument uses Proposition~2.4
only for the specified group $A$, now a $p$-group,
in the good configuration-group extension. For that
application the second equality in (F.7.4), directly
from \cite[Theorem~A]{R-BZ}, replaces Proposition~2.4's
intersection argument. Thus the two required inputs
are justified without using any nonsolvable assertion.
The passage from centralizer to normalizer topology
uses its finite-index kernel under conjugation into
$\Aut(A)$; this kernel is open in the induced topology,
as is detected by any finite quotient in which $A$
injects.

We prove the solvable normalizer comparison, together
with its full induced topology, by induction on the
order of the finite solvable group, simultaneously for
all finite-index subgroups of these mapping class
groups. Finite-index passage preserves the comparison:
intersect the completed normalizer with the corresponding
open subgroup and use density in an open set.

For a nontrivial finite solvable $J$, take a nontrivial
$p$-primary subgroup $A$ of the last nonzero term of
its derived series. That term is abelian, so its
$p$-primary subgroup is characteristic in $J$.
Consequently every normalizer of $J$ normalizes $A$.
The discrete covering normalizer sequence is
\[
 1\longrightarrow A\longrightarrow N_M(A)
                 \longrightarrow Q\longrightarrow1,   \tag{F.7.5}
\]
where $Q$ is a finite-index subgroup of the mapping
class group of the quotient surface with its branch
and puncture decorations. Its genus is again at most
two. The exact quotient map is induced by descent of
the commuting covering diagram; the decorations
record the original punctures, the removed branch
points and the covering subgroup. This is the discrete
normalizer sequence in \cite[\S2.1]{R-Boggi}.

By the prime-power topology input, the completed
normalizer in (F.7.4) is $\widehat{N_M(A)}$.
Completing (F.7.5) therefore gives its quotient by
the finite embedded kernel $A$ as $\widehat Q$.
The subgroup $J/A\le Q$ is smaller and solvable.
The induction hypothesis applies to its normalizer
in $\widehat Q$. The inverse image of that normalizer
under the quotient in (F.7.5), on either side, is
exactly the normalizer of $J$: $J$ is the full inverse
image of $J/A$ because it contains $A$. Pulling back
the dense discrete normalizer proves
\[
                  N_{\widehat M}(J)=\overline{N_M(J)}.
                                                               \tag{F.7.6}
\]
It also proves the topology assertion. Indeed the
completion of this inverse image maps to
$\widehat{N_M(A)}$; on the quotient by $A$ it is
injective by induction, and on $A$ it is the specified
inclusion. Its kernel is therefore trivial.
This completes the induction.

Conjugation gives a continuous homomorphism
\[
       N_{\widehat M}(J)\longrightarrow\Aut(J)
\]
to a finite group. Its kernel is the ambient centralizer
and is open. The dense subgroup $N_M(J)$ from (F.7.6)
meets that open kernel densely, and its intersection
there is exactly $C_M(J)$. This proves
\[
                  C_{\widehat M}(J)=\overline{C_M(J)}
\]
for every finite solvable $J$, without taking an
unjustified intersection of closures.

\emph{The finite image for $A_5$.}
Realize the geometric $A_5$ as the icosahedral rotation
group. Its exceptional point stabilizers have orders
$2,3,5$, and the corresponding transitive orbits are
$A_5/C_2$, $A_5/C_3$, $A_5/C_5$. Every permutation
commuting with this transitive action is right
multiplication by an element of $N_{A_5}(C_d)/C_d$.
This quotient has order two in each of the three
cases; its nonidentity permutation exchanges the
two poles of each relevant rotation axis.

An orientation-preserving homeomorphism commuting
with an odd-order nonidentity rotation fixes its
two poles individually. The local oriented rotation
numbers at the two poles are opposite, and an
orientation-preserving commuting homeomorphism
preserves this local rotation number. This excludes
the nonidentity permutation on the $12$- and
$20$-point orbits for the discrete cyclic centralizers.
Their completed centralizers have the same finite
permutation images by (F.7.4). The full $A_5$
centralizer lies in these cyclic centralizers, so
it fixes these exceptional punctures individually.

For an order-two stabilizer choose a tetrahedral
subgroup $A_4\le A_5$ containing it. An
orientation-preserving homeomorphism commuting
with $A_4$ fixes every ramification point of the
$A_4$ quotient. To see this, its quotient fixes the
branch point of order two. It also fixes each of the
two branch points of order three, rather than
exchanging them, by the preceding local rotation
argument. Isotope the quotient map to the identity
relative to these three branch points and lift the
isotopy. The lift stays in the centralizer; its
endpoint is a deck transformation in $Z(A_4)=1$.
The ramification set is finite and invariant along
the isotopy, so every one of its points is fixed.
The usual discrete covering-centralizer comparison
provides a commuting representative for a centralizing
mapping class. Since $A_4$ is solvable, the proved
centralizer comparison transfers this fixed-puncture
property to $C_{\widehat M}(A_4)$. It therefore holds
for $C_{\widehat M}(A_5)$, excluding the last possible
nonidentity permutation on the $30$-point orbit.

On $m$ free orbits the finite commutant is
$A_5^m\rtimes S_m$. Every element of this commutant
is already realized by the discrete centralizer.
For a direct construction, remove the three branch
points from the quotient sphere. Its unramified
$A_5$-cover has surjective monodromy onto $A_5$.
Move one of the $m$ quotient punctures around a
loop, keeping all the others fixed, and lift the
isotopy starting at the identity. The lift commutes
with $A_5$ and acts on that one free orbit by its
specified monodromy, while fixing the other orbits
and all ramification points. These loops realize
each factor $A_5$. Half-twists interchanging the
quotient punctures realize $S_m$ modulo these factors.
Thus the displayed wreath product is the exact
finite permutation image on both sides of (F.7.1).

The program
\path{code/part_a/a5_centralizer_permutations.py}
constructs $A_5$ as the sixty even permutations of
five letters. It checks its order-two and order-three
generating pair, the three coset actions, normalizer
orders $4,6,10$, and commutants of order two. The
fixed-pole constraints leave only the identity in
each exceptional block; the order-two block uses
the six ramification points for the specified $A_4$.

For completeness, the residual intersection question
has an exact finite-quotient formulation. Take
$u=(1\ 2)(3\ 4)$ and $v=(1\ 3\ 5)$ in $A_5$, which
generate it, and put $D_u=C_M(u)$, $D_v=C_M(v)$.
For open normal subgroups $V\subset U$ of $M$, let
$D_u(V),D_v(V)$ denote their images in $M/V$.
Compactness of coset fibres proves
\[
 \operatorname{im}\bigl(C_{\widehat M}(A_5)\to M/U\bigr)
 =\bigcap_{V\subset U}
       \operatorname{im}\bigl(D_u(V)\cap D_v(V)\to M/U\bigr).
                                                               \tag{F.7.7}
\]
Indeed choose a lift $t\in\widehat M$ of a point
in the right side. For each $V\subset U$, the set
\[
 (\overline{D_u}\,\overline V)\cap
 (\overline{D_v}\,\overline V)\cap t\overline U
\]
is nonempty and compact. These sets are directed
under refinement, so their finite intersections
are nonempty. Compactness produces a single element
in their intersection, which is
$\overline{D_u}\cap\overline{D_v}\cap t\overline U$.
Conversely any such element supplies every finite
coordinate. Equations (F.7.3) and (F.7.7) identify
the remaining assertion precisely: its right side
must equal the image of $C_M(A_5)$ for every $U$.
The finite permutation computation proves that
equality at the puncture-permutation level. It does
not establish it at every refinement. With only
exceptional orbits the discrete centralizer is
trivial, since the quotient sphere has just the
three differently ordered branch points and
$Z(A_5)=1$; the still unproved assertion is then
that the pure completed centralizer is trivial.

\emph{The minimal twelve-puncture case.}
There is a further exact localization of this remaining group.
Write $P=P\Gamma_{0,12}$ and choose $A=A_4\leq A_5$.
The action of $A$ on the twelve vertices is free and transitive.
Its quotient has one free puncture and three branch points of
orders $2,3,3$. The labelled commuting homeomorphism group on
the filled sphere is identified, by descent and the unique
commuting lift, with the identity homeomorphism component of
the quotient fixing those three branch points. It is
contractible. Evaluating the free quotient puncture consequently
identifies the mapping-class centralizer with the free group
$\pi_1(S^2\setminus\{3\text{ branch points}\})=F_2$.
Its permutation on the free fibre is the surjective monodromy
onto $A_4$. The pure centralizer is therefore the exact kernel
\[
 1\longrightarrow C_P(A)=F_{13}\longrightarrow F_2
                    \longrightarrow A_4\longrightarrow1.
                                                               \tag{F.7.8}
\]
Indeed it is the fundamental group of
$X=S^2\setminus R_A$, where the $A_4$ ramification set $R_A$
has $6+4+4=14$ points; its rank is also
$1+12(2-1)=13$ by Schreier. Label the vertices by $g\in A_4$
using a chosen vertex $z_0$. The inclusion is the configuration
map $z\mapsto(gz)_{g\in A_4}$, followed by the quotient by
M\"obius transformations. The proved solvable comparison and
full topology identify its completion with $C_{\widehat P}(A)$.
In particular the unknown group $C_{\widehat M}(A_5)$ lies in
this specified embedded $\widehat F_{13}$.

Each of the six subgroups $C_5\leq A_5$ has two fixed vertices
and two free five-cycles. Retain either one fixed vertex and
one five-cycle, giving twenty-four six-point sets, or retain
both fixed vertices and one five-cycle, giving twelve
seven-point sets. For each retained set $I$, the continuous
forgetting map $r_I$ is defined on $\widehat P$ and is equivariant
for the specified $C_5$. In either target the quotient has
exactly its two branch points and one free puncture. Its
discrete $C_5$ centralizer is therefore exactly the deck group
$C_5$: the quotient mapping class group fixing these three
decorated points is trivial. Prime-power comparison gives the
same completed centralizer. Its intersection with the pure
group is trivial, since the deck group permutes the five-cycle.
It follows that
\[
 C_{\widehat M}(A_5)\subset
 C_{\widehat P}(A)\cap
       \bigcap_{I\text{ among these }36\text{ sets}}\ker(r_I).
                                                               \tag{F.7.9}
\]

The six-point maps in (F.7.9) have a proved common kernel even
on the discrete $F_{13}$. The antipodal permutation on the
twelve vertices commutes with $A_5$. In the regular $A_4$
labelling it is right multiplication by an involution $h\in A_4$.
Every retained six-point set contains exactly one member of
each antipodal pair. Collisions of two entries $g_i z,g_j z$
at a ramification point with stabilizer $H\leq A_4$ occur
precisely when $g_j^{-1}g_i\in H$; these are the right cosets
$gH$. Thus there is no collision among the retained entries
at either fixed point of $h$. Each six-point configuration map
extends across both of these deleted points of $X$. Their
nontrivial peripheral meridians in $F_{13}$ map to the identity
under all twenty-four maps. Completion retains these meridians
as nontrivial primitive procyclic subgroups. Consequently the
joint six-point map on the embedded $\widehat F_{13}$ is
noninjective by these explicit elements.

Every seven-point set retains a collision at each of the
fourteen ramification points: seven choices from six pairs
force a pair, and seven choices from four triples force a
triple collision. The program verifies all twelve sets and
the preceding common invisible axis for the six-point sets,
using right cosets throughout. It also computes the induced
maps on $H_1(F_{13},\F_{101})$. A meridian at an $A_4$
ramification point maps to the sum of the pair generators
within its collision clusters. In the target pure sphere
group the pair generators satisfy the vertex-sum and total-sum
relations. The two fixed points of each cyclic stabilizer give
the same homology column. Each seven-point map and their joint
map have rank six, with seven-dimensional kernel; the joint
six-point map has rank five. These are exact calculations of
the indicated abelian quotients. They neither produce a
nontrivial $A_5$ centralizer element nor prove injectivity of
the joint seven-point map on $\widehat F_{13}$.

F.8 proves that all twelve seven-point kernels in (F.7.9)
are the same, both discretely and after completion. It further
proves that the discrete kernel contains no nonidentity reducible
mapping class and that the completed kernel meets every conjugate
of every full source cusp subgroup trivially. The exact four-point
characters detecting those cusps and all eighty-four finite
subset--inertia checks are supplied there. F.9 subsequently
constructs an explicit nontrivial common kernel element with
a full procyclic closure, and proves that all twenty-four
six-point restrictions factor through these seven-point maps.
The thirty-six-fold kernel is therefore nontrivial. The full
centralizer comparison and geometric conjugacy remain unproved;
no converse to the inclusion (F.7.9) is inferred.

This audit concerns the comparison requested for
\S30, folios 194--196, scans 383--387. F.6's
arithmetic open-image proof applies to $A_5$ without
using (F.7.1), so that proved arithmetic assertion
does not depend on the unresolved intersection.
\end{proof}

\begin{claim}{F.8}{Proved}{The seven-point kernel and its peripheral restrictions}
In the regular $A=A_4$ labelling of (F.7.8), all twelve
seven-point subsets in (F.7.9) form one free orbit for
left multiplication by $A$.  Consequently their kernels
on the specified $F_{13}$ are equal; their kernels on
its profinite completion are equal as well.
The completed kernel meets every conjugate of every
full source cusp subgroup $\widehat\Z$ trivially.
The discrete kernel contains no nonidentity reducible
mapping class under the specified index-twelve
inclusion $F_{13}\le P\Gamma_{0,4}$.
\end{claim}
\passage{\S{}30; folios 194--196; scans 383--387; Part F: nonsolvable finite centralizer comparison}

\begin{proof}

Here is the actual comparison of maps.  Regard the target
as $M_{0,I}$ with its seven labels $I\subset A$, and write
\[
 f_I:X\longrightarrow M_{0,I},\qquad
             z\longmapsto[(gz)_{g\in I}].
\]
For $h\in A$, let
$r_h:M_{0,I}\xrightarrow{\sim}M_{0,hI}$ be the relabelling
which assigns to the label $hg$ the former coordinate
with label $g$.  The configuration defining $f_{hI}(z)$
has coordinate $h(gz)$ at that label.  Applying the same
M\"obius transformation $h^{-1}$ to every coordinate
therefore gives the exact identity of maps
\[
                         f_{hI}=r_h\circ f_I.
\]
This identity also identifies the base configurations;
it induces an isomorphism of the target fundamental
groups intertwining the two homomorphisms from
$\pi_1(X)$.  Functoriality of profinite completion gives
the same commuting diagram for their completions.
Since $r_h$ and its completion are isomorphisms, the
two kernel equalities follow.

The exact label certificate is
\path{code/part_a/seven_point_same_kernel.py}.
It constructs the twelve subsets in the action
$A_5/C_5$ and gives all twelve left multipliers and
the corresponding seven-label permutations.  With
the certificate's zero-based vertex numbering, one
subset is $I=\{0,1,2,3,4,7,10\}$.  Its intersections
with the three cosets of the normal subgroup
$V_4\le A_4$ are respectively
\[
 \begin{array}{c|c}
 \text{coset}&I\text{ retained in that coset}\\ \hline
 \{0,2,7,11\}&\{0,2,7\}\\
 \{1,4,5,9\}&\{1,4\}\\
 \{3,6,8,10\}&\{3,10\}.
 \end{array}
\]
The program checks that none of the twelve subsets
contains a full $V_4$ coset.  Thus the four-label
$V_4$-orbit cross-ratio is not one of the forgetful
projections of any of these seven-point maps.

The intersection of the twelve kernels in (F.7.9)
is therefore the kernel of any one of these named
maps.  F.9 below proves it nontrivial by an exact braid certificate.
This discrete kernel contains no nonidentity reducible
mapping class in the source subgroup of
$P\Gamma_{0,4}$.  To prove this, use the exact
$A_4$ covering
\[
 X=S^2\setminus R_A\longrightarrow
 S^2\setminus\{b_2,b_3,b'_3\}
\]
and its index-twelve injection on fundamental groups
in (F.7.8).  On the four-punctured quotient sphere,
every essential curve separates the moving puncture
$q$ and one branch point $b$ from the other two
branch points.  Let $e=2$ or $3$ be the order at $b$.
The preimage of the disk containing $q,b$ consists
of $12/e$ disks, each containing $e$ of the twelve
points above $q$.  These assertions hold for every
such curve: its monodromy subgroup is the conjugate
of the inertia subgroup at $b$, so its connected
components have degree $e$, and each is a disk
branched only at the one point above $b$.

The lift of a nonzero twist power belonging to the
index-twelve subgroup is a product of twists about
these boundary curves.  More precisely, its exponent
is divisible by $e$, and each lifted twist has exponent
equal to the original exponent divided by $e$.  After
forgetting five points, at least one of these boundaries
remains essential: seven retained points distributed
among six pairs or four triples force a disk containing
at least two retained points, and its complement
contains at least four.  The surviving twists all have
the same nonzero sign; parallel essential components
add their exponents.  Independence of twists about
pairwise disjoint nonisotopic essential curves makes
the resulting mapping class nonidentity.  Finally,
each nonidentity reducible pure mapping class of a
four-punctured sphere is a nonzero power of one such
twist: its invariant multicurve has just one component,
and cutting along it gives two three-punctured spheres.
This proves the asserted exclusion of reducible
elements from the discrete kernel.

There is also an injective detector for each full
profinite cusp.  Let $c$ be a primitive peripheral
generator at any one of the fourteen punctures of
$X$.  Its image in $P\Gamma_{0,4}$ is a conjugate
of $T^e$, where $T$ is a twist as above.  Thus its
lift is the indicated product of twists of exponent
one.  Choose two retained labels in an essential
lifted disk and two outside that disk.  Forgetting
all other labels gives a specified map to
$P\Gamma_{0,4}$.  The inside disk contributes one
twist $U$.  Among the other disks, another twist
survives exactly when both outside labels lie in
one disk; it then gives the same isotopy class and
the same sign.  Hence the image of $c$ is $U^d$
with $d\in\{1,2\}$, up to conjugation.

There is a character
$\ell:P\Gamma_{0,4}\to\Z$ with $\ell(U)=1$.
For example the three peripheral twists in the
free group $P\Gamma_{0,4}=F_2$ have abelianization
vectors $(1,0),(0,1),(-1,-1)$, and an essential
twist is conjugate to one of these under the pure
group.  Choose the corresponding integral linear
functional.  Completing this character and the
two forgetting maps gives a continuous homomorphism
\[
 \widehat F_{13}\longrightarrow
 \widehat{P\Gamma}_{0,7}\longrightarrow
 \widehat{P\Gamma}_{0,4}\longrightarrow\widehat\Z
\]
whose restriction to $\overline{\langle c\rangle}$
is $a\mapsto da$.  Multiplication by $1$ or $2$
is injective on $\widehat\Z$.  The cusp is a full
embedded $\widehat\Z$ in $\widehat F_{13}$, since
its primitive generator is a member of a free
basis after choosing the omitted sphere peripheral.
The same final character detects every conjugate
$g c^a g^{-1}$ with $g\in\widehat F_{13}$, because
its target is abelian.  This proves the asserted
trivial intersection of the completed kernel with
all conjugate cusp groups.  The program gives the
four retained labels and $d$ in each of the
$12\cdot7=84$ subset--inertia-partition checks;
the two fixed points of each inertia subgroup
have the same partition.

The exclusion just proved confines nonidentity discrete
kernel elements to pseudo-Anosov mapping classes. F.9 supplies
an explicit such element by an exact braid certificate.
No assertion about a Nielsen--Thurston type for profinite
elements is used.

The preceding rank-six calculation is a calculation
of its induced abelian map, and does not identify
the kernel of the fundamental-group homomorphism.
The separate file
\path{code/part_a/seven_point_braid_explore.py}
contains a numerical crossing exploration; it is
explicitly exploratory and is not used in this
map comparison. F.9 replaces the discovery computation
by a separate exact rational certificate of a kernel relation.
\end{proof}

\begin{claim}{F.9}{Proved}{An explicit common kernel element for the seven-point maps}
Every one of the twelve seven-point maps in F.7--F.8 has
a nontrivial kernel, both on the specified discrete
$F_{13}$ and on its full profinite completion.  With the
source coordinate, basepoint, and meridians specified
below, the same element
\[
 w=m_{-1}m_i m_1^{-1}m_i^{-1}                         \tag{F.9.1}
\]
lies in all twelve kernels.  A source winding character
sends $w$ to $1$, so its procyclic closure is a full
$\widehat\Z$ contained in their common completed kernel.
The joint kernel of all thirty-six six- and seven-point
maps in (F.7.9) is exactly this same kernel.
This refutes injectivity of these particular restriction
maps; the residual $A_5$ centralizer comparison in F.7
remains unresolved.
\end{claim}
\passage{Section 30, folios 194-196, scans 383-387; Part F, nonsolvable A5 centralizer comparison on the twelve-punctured sphere}

\begin{proof}
\textbf{The specified source and target.}
Use the faithful tetrahedral action generated by
\[
 z\longmapsto-z,\qquad z\longmapsto\frac{z+i}{z-i}.
\]
Its fourteen ramification points are infinity and the
zeros of
\[
 D(z)=z(z^4-1)(z^8+14z^4+1).
\]
Thus $X=\mathbf P^1\setminus R_{A_4}$ has
$\pi_1(X,b)=F_{13}$ for $b=5+4i$.
The seven labels $I=\{0,1,2,3,4,7,10\}$ of F.8,
in that order, have the following coordinate functions:
\[
 \begin{split}
 (m_0(z),\ldots,m_6(z))=
 \left(z,\frac{i(z+1)}{z-1},-\frac1z,
 -\frac{z-i}{z+i},\frac{i(z-1)}{z+1},-z,
 \frac{z-i}{z+i}\right).
 \end{split}                                                   \tag{F.9.2}
\]
Here the subscripts $0,\ldots,6$ denote positions in
the displayed seven-set, not its original labels.
The exact finite part of the certificate below verifies
all $144$ matrix products in this $A_4$, its correspondence
with the labels in F.8, and the ramification polynomial
$D$ by multiplying the fixed-point polynomials of its
seven nontrivial cyclic subgroups.

For $z\in X$, let
\[
 L_z(t)=\frac{(t-m_0(z))(m_1(z)-m_6(z))}
              {(t-m_6(z))(m_1(z)-m_0(z))}.
\]
The seven orbit points are distinct, so this is the
unique M\"obius transformation carrying the first,
second, and seventh points to $0,1,\infty$.
Define the ordered configuration map
\[
 \Phi:X\longrightarrow\operatorname{Conf}_6(\mathbf C),
 \qquad z\longmapsto(R_0(z),\ldots,R_5(z)),
 \quad R_j(z)=L_z(m_j(z)).                              \tag{F.9.3}
\]
The rational functions extend across any removable
singularities in their displayed expressions; their
only poles are among $R_{A_4}$.  In particular
$R_0=0$ and $R_1=1$.
Appending the seventh point $\infty$ and taking the
labelled M\"obius quotient gives a map
\[
 q:\operatorname{Conf}_6(\mathbf C)\longrightarrow M_{0,I}
 \quad\hbox{with}\quad q\circ\Phi=f_I.                 \tag{F.9.4}
\]
These are actual continuous maps, with the displayed
base configurations.  The induced map of fundamental
groups therefore factors through the pure braid group
$P_6\le B_6$.  The inclusion here is induced by the
cover from ordered to unordered configurations and is
injective by path lifting.  We shall prove
$\Phi_*(w)=1$ in $B_6$, which proves $f_{I,*}(w)=1$
by this factorization.

\textbf{The three based source meridians.}
Put $r=27/1000$.  For $c=-1,i,1$, set
\[
 t_c=(\operatorname{Re}c-r)+4i,
 \quad b_c=c-r-ri,
 \quad d_c=c+r-ri,
 \quad e_c=c+r+ri,
 \quad a_c=c-r+ri.
\]
Let $m_c$ be the loop obtained by traversing the
straight polygonal segments with successive vertices
\[
 b,\ t_c,\ b_c,\ d_c,\ e_c,\ a_c,\ b_c,\ t_c,\ b.
                                                               \tag{F.9.5}
\]
The square is counterclockwise; its stem is followed
by its reverse.  Each square contains the single
ramification point $c$.  This also follows directly
from the above polynomial: the other eight finite
points have real and imaginary coordinates with
absolute values $(\sqrt3-1)/2$ or $(\sqrt3+1)/2$,
whereas the remaining finite points are
$0,\pm1,\pm i$.
The exact segment check below also verifies that every
stem and every square edge avoids all zeros of $D$.
Source loop products in (F.9.1) are chronological
concatenation.

\textbf{Exact extraction of the three braids.}
Multiply every coordinate in (F.9.3) by
$\eta=1+(173/1000)i$ and project to its real part.
This multiplication is an orientation-preserving
homeomorphism of the plane.  At the basepoint the
six coordinate positions, in increasing projected
order, are
\[
                      (0,5,3,1,4,2).
\]
Identify this chamber with the standard real ordered
base configuration by a path within the chamber.
The chamber of any fixed strict real ordering is
convex, so the identification is unique up to
homotopy.  Read crossings chronologically.  Write
$j$ for $\sigma_j$, and $-j$ for $\sigma_j^{-1}$;
the positive sign means that the strand on the left
immediately before the crossing has larger imaginary
coordinate.  This fixes both the numbering and the
crossing convention.

The executed certificate
\path{code/part_a/seven_point_kernel_f9.py}
uses only rational and Gaussian-rational arithmetic.
For a segment $z=p+(q-p)t$, $0\le t\le1$, and each
pair of normalized coordinates, write their projected
complex difference before taking real parts as
$A(t)/B(t)$ with $A,B\in\mathbf Q(i)[t]$.
It verifies that $A$, $B$, every coordinate denominator,
and $D(p+(q-p)t)$ have no complex zero on the real
segment: the real and imaginary parts have a greatest
common divisor with no real root there.
The possible crossings are precisely the roots in
$(0,1)$ of
\[
       P(t)=\operatorname{Re}(A(t)\overline{B(t)}).
                                                               \tag{F.9.6}
\]
Clearing rational denominators preserves its roots.
Sturm isolation over $\mathbf Q$ enumerates all of
them.  Their rational isolating intervals are pairwise
disjoint, their multiplicities are one, and
$\operatorname{Im}(A\overline B)$ has a certified
constant nonzero sign on each interval.  No endpoint
is a crossing.  The two strands are checked to be
adjacent in the current order at every crossing,
and the final order is checked against direct exact
evaluation at the endpoint.

There are $18$ different segments up to reversal,
$270$ pair checks, and $64$ isolated crossings in
these segments; the largest crossing-polynomial
degree is $12$.  Reversed segments give the inverse
words.  Between successive crossings the path lies
in a convex ordering chamber.  At a crossing only
the named adjacent pair changes order, with its
certified nonzero imaginary separation.  Straightening
these intervals and crossing neighborhoods gives the
usual elementary braid pieces.  Thus the recorded
words represent the actual monodromies, not just a
linear representation of them.  This is also the
projection-isotopy argument of Artin's Theorem~3
\cite{A-Artin1947}.

The freely reduced words for $m_{-1},m_i,m_1$ are,
respectively, $A,B,C$ below; a line break only separates
consecutive parts of the same word:
\[
\begin{aligned}
 A={}&(1,4,5,2,1,3,2,3,2,1,2,-5,-4,-5,\\
     &\hspace{9mm}-5,4,5,-2,-1,-2,-3,-2,-3,-1,-2,-5,-4,-1),\\
 B={}&(1,4,5,2,3,1,4,2,3,4,-3,1,2,3,4,5,\\
     &\hspace{9mm}-2,5,4,3,2,1,-4,-3,-2,-4,-1,-3,-2,-5,-4,-1),\\
 C={}&(1,4,5,2,3,4,1,2,3,4,5,-1,3,2,-5,-4,\\
     &\hspace{9mm}-4,5,-2,-3,1,-5,-4,-3,-2,-1,-4,-3,-2,-5,-4,-1).
\end{aligned}                                                   \tag{F.9.7}
\]

\textbf{The braid identity and the source detector.}
Use the faithful Artin substitution on
$F_6=\langle X_1,\ldots,X_6\rangle$:
\[
 \sigma_j(X_j)=X_jX_{j+1}X_j^{-1},\qquad
 \sigma_j(X_{j+1})=X_j,
 \qquad \sigma_j(X_k)=X_k\ (k\ne j,j+1).              \tag{F.9.8}
\]
Its inverse is used for a negative letter.  Products
of substitutions are composed from the displayed
word order; the certificate implements composition
by substitution and cancellation of adjacent inverse
letters.  Artin's Theorem~14, with the inverse
generator formula on p.~113, proves faithfulness
\cite{A-Artin1947}.  Simultaneously reversing all
generator signs is an automorphism of the braid
group, and reversing the stacking convention is an
anti-automorphism; neither changes the identity test.

For the exact words (F.9.7), substitution gives
\[
       \operatorname{Artin}(AB)=\operatorname{Artin}(BC),
 \qquad
       \operatorname{Artin}(ABC^{-1}B^{-1})(X_j)=X_j
       \quad(1\le j\le6).                              \tag{F.9.9}
\]
This is an integer free-word calculation.  The output
records all six freely reduced common images for the
first equality and all six single-letter images for
the second.  Faithfulness gives
$AB=BC$ in $B_6$.  By the certified identification of
the three words with the three source monodromies,
(F.9.1) lies in $\ker\Phi_*$, and hence in
$\ker f_{I,*}$ by (F.9.4).

The map $X\to\mathbf C^*$, $z\mapsto z+1$, gives
the integral winding character
\[
 \ell:\pi_1(X,b)\longrightarrow\Z,
 \qquad
 \ell(m_{-1})=1,\quad \ell(m_i)=\ell(m_1)=0.
\]
Therefore $\ell(w)=1$ and $w\ne1$.
Completing this exact character gives
$\widehat\ell:\widehat F_{13}\to\widehat\Z$ with
$\widehat\ell(w)=1$; this proves nontriviality in the
full completion without a detection assumption.
The continuous power map $\widehat\Z\to\widehat F_{13}$,
$a\mapsto w^a$, has composite $a\mapsto a$ with
$\widehat\ell$, so it is injective.  Its image lies
in the closed completed kernel.

Finally F.8 gives the actual map identities
$f_{hI}=r_h\circ f_I$ for the twelve left translates
$hI$, with $r_h$ a labelled target isomorphism.
They carry this same source word, and all its
profinite powers, into the identity in every target.

The twenty-four six-point conditions in (F.7.9)
follow from these twelve conditions by specified
forgetting maps.  For one of the six subgroups
$C_5\le A_5$, let $p_+,p_-$ be its fixed poles and
let $O$ be either of its two five-cycles.  Put
\[
 K=O\cup\{p_+,p_-\},\qquad J_\pm=O\cup\{p_\pm\}.
\]
Forgetting the other pole gives the labelled morphism
$d_\pm:M_{0,K}\to M_{0,J_\pm}$ and the exact identity
$f_{J_\pm}=d_\pm\circ f_K$ on the same based source.
The twelve choices of $(C_5,O)$ give all twelve
seven-sets, and their two pole choices give all
twenty-four six-sets.  The separate executed certificate
\path{code/part_a/seven_point_six_factorization_f9.py}
lists all twenty-four containments, the omitted pole,
and the six positions retained in each seven-tuple.
Functoriality of fundamental groups and of full
profinite completion proves on both sources that
\[
 \bigcap_{I\text{ among all }36}\ker f_{I,*}
       =\bigcap_{K\text{ among the }12\text{ seven-sets}}
                                      \ker f_{K,*}.     \tag{F.9.10}
\]
For the completed equality the starred maps in
(F.9.10) mean their completed homomorphisms.
Consequently (F.9.1) and its full procyclic closure
also satisfy every six-point condition.

The common-kernel injectivity route in F.7 therefore
fails.  The necessary seven-point restrictions of
an element of the residual $A_5$ centralizer still
vanish, but the converse implication has not been
proved.  No centralizer element is constructed here.

The discovery computation used numerical crossings
and a finite Burau image to find the short candidate.
Neither computation is an input to this proof.
The file
\path{code/part_a/seven_point_kernel_f9_output.json}
contains the deterministic exact certificate:
the coordinate functions, Gaussian-rational path
vertices, every crossing polynomial, every rational
root interval and sign, the three words, and the
Artin identities.  The companion stdout file records
the successful execution.
\end{proof}

\clearpage\bibliographystyle{plain}\bibliography{LM4_FINITE_GROUP_ACTIONS_references}
\end{document}
